\documentclass[11pt,oneside]{article}

\usepackage[a4paper,margin=27mm,headheight=15pt]{geometry}
\usepackage[T1]{fontenc}
\usepackage[utf8]{inputenc}
\IfFileExists{libertinus.sty}{\usepackage{libertinus}}{\usepackage{lmodern}}
\usepackage{microtype}
\usepackage{amsmath,amssymb,amsthm,mathtools,mathrsfs}
\usepackage{bm}
\usepackage{booktabs,longtable,array,tabularx}
\usepackage{graphicx}
\usepackage{pgfplots}
\usepackage{xcolor}
\usepackage{enumitem}
\usepackage{fancyhdr}
\usepackage{titlesec}
\usepackage{listings}
\usepackage{xurl}
\usepackage{hyperref}
\usepackage[nameinlink,noabbrev]{cleveref}
\usepackage{aliascnt}

\definecolor{linkblue}{RGB}{25,75,135}
\definecolor{shadegray}{RGB}{246,247,249}
\definecolor{rulegray}{RGB}{195,201,208}
\pgfplotsset{compat=1.18}
\hypersetup{
  colorlinks=true,
  linkcolor=linkblue,
  citecolor=linkblue,
  urlcolor=linkblue,
  pdftitle={Semi-Formalized Research Frontiers for the Fabius Function},
  pdfsubject={Six deduplicated thematic syntheses from eleven research notebooks, crosswalked against the Lean corpus},
  pdfkeywords={Fabius function, Rvachev up-function, research frontiers, conjectures, Lean formalization}
}

\pagestyle{fancy}
\fancyhf{}
\fancyhead[L]{\small Fabius research frontiers}
\fancyhead[R]{\small\nouppercase{\leftmark}}
\fancyfoot[C]{\thepage}
% Canonical project-wide mathematical notation.
% Canonical mathematical notation for active FabiusFunction LaTeX documents.
%
% This file is intentionally a plain TeX input rather than a LaTeX package.
% Every standalone document loads it by an explicit path relative to that
% document's directory; included fragments inherit the definitions from their
% root document.  Keep this file independent of document layout and metadata.
%
% Contract:
%   * one command has one mathematical meaning throughout the active corpus;
%   * names encode semantic roles rather than a document-local choice of letter;
%   * visually similar but differently normalized objects have different print;
%   * every command below is defined normatively in the notation catalogue.
%
% The active corpus excludes docs/archive/.  Historical sources may retain
% their original notation only inside an explicitly labelled quotation.

\makeatletter
\@ifundefined{mathbb}{%
  \PackageError{fabius-notation}{The amssymb package must be loaded first}%
    {Load amsmath and amssymb before inputting fabius-notation.tex.}%
}{}
\@ifundefined{operatorname}{%
  \PackageError{fabius-notation}{The amsmath package must be loaded first}%
    {Load amsmath before inputting fabius-notation.tex.}%
}{}
\makeatother

% Version stamp used by the catalogue and by static audits.
\newcommand{\FabiusNotationVersion}{2026-08-31}

% Number systems.  NaturalNumbers is explicitly the nonnegative convention.
\newcommand{\NaturalNumbers}{\mathbb{N}_{0}}
\newcommand{\PositiveIntegers}{\mathbb{N}_{>0}}
\newcommand{\IntegerNumbers}{\mathbb{Z}}
\newcommand{\RationalNumbers}{\mathbb{Q}}
\newcommand{\RealNumbers}{\mathbb{R}}
\newcommand{\ComplexNumbers}{\mathbb{C}}
\newcommand{\ExtendedRealNumbers}{\overline{\mathbb{R}}}

% The three principal Fabius--Rvachev functions.  Subscripts are deliberate:
% a bare F and a calligraphic F have several unrelated uses in the corpus.
\newcommand{\FabiusBounded}{F_{\mathrm{bd}}}
\newcommand{\FabiusGlobal}{F_{\mathrm{gl}}}
\newcommand{\FabiusQuantile}{Q_{F}}
\newcommand{\FabiusClampedQuantile}{Q_{F,\mathrm{cl}}}
\newcommand{\RvachevUp}{\operatorname{up}}

% Constants, elementary operators, and delimiters.
\newcommand{\EulerE}{\mathrm{e}}
\newcommand{\ImaginaryUnit}{\mathrm{i}}
\newcommand{\Differential}{\mathop{}\!\mathrm{d}}
\newcommand{\AbsoluteValue}[1]{\left\lvert #1\right\rvert}
\newcommand{\Norm}[1]{\left\lVert #1\right\rVert}
\newcommand{\Floor}[1]{\left\lfloor #1\right\rfloor}
\newcommand{\Ceiling}[1]{\left\lceil #1\right\rceil}
\newcommand{\FractionalPart}[1]{\left\{#1\right\}}
\newcommand{\PositivePart}[1]{\left(#1\right)_{+}}
\newcommand{\IndicatorOf}[1]{\mathbf{1}_{#1}}
\newcommand{\SetOf}[1]{\left\{#1\right\}}
\newcommand{\InnerProduct}[1]{\left\langle #1\right\rangle}
\newcommand{\CardinalityOf}[1]{\left\lvert #1\right\rvert}
\newcommand{\DefinitionEquals}{\mathrel{\coloneqq}}
\newcommand{\Sign}{\operatorname{sgn}}
\newcommand{\IdentityOperator}{\operatorname{id}}
\newcommand{\SpanOperator}{\operatorname{span}}
\newcommand{\TraceOperator}{\operatorname{tr}}
\newcommand{\RankOperator}{\operatorname{rank}}
\newcommand{\DiagonalOperator}{\operatorname{diag}}
\newcommand{\ResidueOperator}{\operatorname*{Res}}
\newcommand{\LeastCommonMultiple}{\operatorname{lcm}}
\newcommand{\OrderOperator}{\operatorname{ord}}
\newcommand{\PrincipalLogarithm}{\operatorname{Log}}
\newcommand{\FallingFactorial}[2]{\left(#1\right)^{\underline{#2}}}
\newcommand{\RisingFactorial}[2]{\left(#1\right)^{\overline{#2}}}

% Support, topology, convolution, and metric notation.
\newcommand{\SupportOperator}{\operatorname{supp}}
\newcommand{\SupportOf}[1]{\SupportOperator\!\left(#1\right)}
\newcommand{\TopologicalSupportOf}[1]{\operatorname{tsupp}\!\left(#1\right)}
\newcommand{\Convolution}{\mathbin{*}}
\newcommand{\MetricDistance}{\operatorname{dist}}
\newcommand{\TotalVariationDistance}{d_{\mathrm{TV}}}

% Probability.  Equality and convergence in law are not metric distance.
\newcommand{\Probability}{\mathbb{P}}
\newcommand{\Expectation}{\mathbb{E}}
\newcommand{\Variance}{\operatorname{Var}}
\newcommand{\Covariance}{\operatorname{Cov}}
\newcommand{\LawOperator}{\operatorname{Law}}
\newcommand{\LawOf}[1]{\LawOperator\!\left(#1\right)}
\newcommand{\EqualInLaw}{\mathrel{\overset{\mathrm{law}}{=}}}
\newcommand{\ConvergesInLaw}{\mathrel{\xrightarrow{\mathrm{law}}}}

% Fourier and integral transforms.  The printed labels expose normalization.
% FourierTwoPi uses exp(-2*pi*i*x*xi); FourierAngular uses exp(-i*x*omega).
\newcommand{\FourierTwoPi}[1]{\widehat{#1}_{\,2\pi}}
\newcommand{\FourierAngular}[1]{\widehat{#1}_{\,\mathrm{ang}}}
\newcommand{\LaplaceTransformSymbol}{\mathcal{L}}
\newcommand{\MellinTransformSymbol}{\mathcal{M}}
\newcommand{\LaplaceTransformOf}[1]{\LaplaceTransformSymbol\!\left[#1\right]}
\newcommand{\MellinTransformOf}[1]{\MellinTransformSymbol\!\left[#1\right]}
\newcommand{\MomentGeneratingFunction}[1]{M_{#1}}
\newcommand{\CharacteristicFunction}[1]{\varphi_{#1}}

% The two standard sinc functions must never share one symbol.
\newcommand{\SincRad}{\operatorname{sinc}_{\mathrm{rad}}}
\newcommand{\SincPi}{\operatorname{sinc}_{\pi}}

% Binary arithmetic and Thue--Morse notation.
\newcommand{\BinaryDigitSum}{\operatorname{wt}_{2}}
\newcommand{\ThueMorseSignSymbol}{\varepsilon}
\newcommand{\ThueMorseSign}[1]{\varepsilon_{#1}}
\newcommand{\PadicValuation}[1]{\nu_{#1}}
\newcommand{\TwoAdicValuation}{\nu_{2}}

% q-series notation.  Every base is an explicit argument.
\newcommand{\QPochhammer}[3]{\left(#1;#2\right)_{#3}}
\newcommand{\GaussianBinomial}[3]{%
  \genfrac{[}{]}{0pt}{}{#1}{#2}_{#3}%
}
\newcommand{\QInteger}[2]{\left[#1\right]_{#2}}

% Named special functions.
\newcommand{\Polylogarithm}{\operatorname{Li}}
\newcommand{\LambertW}{W}
\newcommand{\GammaFunction}{\Gamma}
\newcommand{\RiemannZeta}{\zeta}

% Asymptotics.  BigO and LittleO are operator tokens for existing formula
% grammar; prefer the At forms so that the limiting regime is printed.
\newcommand{\BigO}{\mathcal{O}}
\newcommand{\LittleO}{o}
\newcommand{\BigOAt}[2]{\mathcal{O}_{#2}\!\left(#1\right)}
\newcommand{\LittleOAt}[2]{o_{#2}\!\left(#1\right)}

\renewcommand{\headrulewidth}{0.4pt}

\titleformat{\section}{\Large\bfseries}{\thesection}{0.7em}{}
\titleformat{\subsection}{\large\bfseries}{\thesubsection}{0.7em}{}
\titleformat{\subsubsection}{\normalsize\bfseries}{\thesubsubsection}{0.7em}{}
\setcounter{secnumdepth}{3}
\setcounter{tocdepth}{2}
\setlist[itemize]{leftmargin=2em,itemsep=0.25em,topsep=0.4em}
\setlist[enumerate]{leftmargin=2.3em,itemsep=0.25em,topsep=0.4em}

% Long unbreakable inline mathematics occasionally leaves a line with no legal
% break point.  A small emergency stretch lets TeX loosen such a line rather
% than let it protrude into the margin; it applies only to lines that would
% otherwise be overfull.
\setlength{\emergencystretch}{2em}

% ed.: all of these share the `theorem' counter so that numbering runs in one
% ed.: sequence.  Declared as \newtheorem{X}[theorem]{...} they also report the
% ed.: counter name `theorem' to cleveref, so every \Cref to a proposition,
% ed.: lemma, corollary, conjecture, definition, algorithm, example, remark,
% ed.: warning, candidate statement or formalization obligation printed the
% ed.: word ``Theorem''.  In this volume that is not merely untidy: a reference
% ed.: to a formalization obligation or a candidate statement claiming to be a
% ed.: theorem inverts the quarantine distinction the document exists to keep.
% ed.: aliascnt gives each environment its own counter name aliased to
% ed.: `theorem', so the shared numbering is preserved exactly while cleveref
% ed.: sees the true type.
\newtheorem{theorem}{Theorem}[section]

\newaliascnt{proposition}{theorem}
\newtheorem{proposition}[proposition]{Proposition}
\aliascntresetthe{proposition}

\newaliascnt{lemma}{theorem}
\newtheorem{lemma}[lemma]{Lemma}
\aliascntresetthe{lemma}

\newaliascnt{corollary}{theorem}
\newtheorem{corollary}[corollary]{Corollary}
\aliascntresetthe{corollary}

\newaliascnt{conjecture}{theorem}
\newtheorem{conjecture}[conjecture]{Conjecture}
\aliascntresetthe{conjecture}

\theoremstyle{definition}
\newaliascnt{definition}{theorem}
\newtheorem{definition}[definition]{Definition}
\aliascntresetthe{definition}

\newaliascnt{algorithm}{theorem}
\newtheorem{algorithm}[algorithm]{Algorithm}
\aliascntresetthe{algorithm}

\newaliascnt{example}{theorem}
\newtheorem{example}[example]{Example}
\aliascntresetthe{example}

\theoremstyle{remark}
\newaliascnt{remark}{theorem}
\newtheorem{remark}[remark]{Remark}
\aliascntresetthe{remark}

\newaliascnt{warning}{theorem}
\newtheorem{warning}[warning]{Warning}
\aliascntresetthe{warning}

\newaliascnt{candidate}{theorem}
\newtheorem{candidate}[candidate]{Candidate statement}
\aliascntresetthe{candidate}

\newaliascnt{obligation}{theorem}
\newtheorem{obligation}[obligation]{Formalization obligation}
\aliascntresetthe{obligation}

\crefname{proposition}{proposition}{propositions}
\Crefname{proposition}{Proposition}{Propositions}
\crefname{lemma}{lemma}{lemmas}
\Crefname{lemma}{Lemma}{Lemmas}
\crefname{corollary}{corollary}{corollaries}
\Crefname{corollary}{Corollary}{Corollaries}
\crefname{conjecture}{conjecture}{conjectures}
\Crefname{conjecture}{Conjecture}{Conjectures}
\crefname{definition}{definition}{definitions}
\Crefname{definition}{Definition}{Definitions}
\crefname{algorithm}{algorithm}{algorithms}
\Crefname{algorithm}{Algorithm}{Algorithms}
\crefname{example}{example}{examples}
\Crefname{example}{Example}{Examples}
\crefname{remark}{remark}{remarks}
\Crefname{remark}{Remark}{Remarks}
\crefname{warning}{warning}{warnings}
\Crefname{warning}{Warning}{Warnings}
\crefname{candidate}{candidate statement}{candidate statements}
\Crefname{candidate}{Candidate statement}{Candidate statements}
\crefname{obligation}{formalization obligation}{formalization obligations}
\Crefname{obligation}{Formalization obligation}{Formalization obligations}

\newcommand{\repo}{\nolinkurl{Analysis/FabiusFunction}}
\newcommand{\lean}[1]{%
  \ifmmode\text{\texttt{\detokenize{#1}}}\else\nolinkurl{#1}\fi}

\newenvironment{proofidea}{\par\smallskip\noindent\textit{Proof architecture.}\ }{\hfill$\square$\par\smallskip}
\newenvironment{boxedremark}{%
  \par\smallskip\noindent\begin{center}\begin{minipage}{0.94\textwidth}%
  \color{black}\hrule\vspace{0.6em}\small
}{%
  \vspace{0.6em}\hrule\end{minipage}\end{center}\smallskip
}

\lstdefinestyle{pseudo}{
  basicstyle=\ttfamily\small,
  backgroundcolor=\color{shadegray},
  frame=single,
  rulecolor=\color{rulegray},
  columns=fullflexible,
  keepspaces=true,
  showstringspaces=false,
  xleftmargin=0.7em,
  xrightmargin=0.7em,
  aboveskip=0.8em,
  belowskip=0.8em
}

\title{\bfseries Semi-Formalized Research Frontiers for the Fabius Function\\[0.35em]
\large Six thematic syntheses, crosswalked against the Lean corpus}
\author{Boundary layer between prose and proof in the \texttt{ProveIt} Fabius development\\
\small Eleven source notebooks in six syntheses plus a post-audit gap register}
\date{Base research snapshot: 25 August 2026; consolidated repository updated
28 August 2026; formal-status audit updated 1 September 2026\\
\small affine-law checkpoint \nolinkurl{d312c0603};
fractional and inverse-layer-cake checkpoint \nolinkurl{37eab3c56}}

\begin{document}
% Compatibility macros used by the consolidated source notebooks.  They are
% document-local aliases, placed after the verbatim house preamble.
\newcommand{\coef}[2]{[X^{#1}]\,#2}
\newcommand{\Dq}{\mathfrak D}
\newcommand{\shiftop}{\mathsf L}
\newcommand{\eps}{\varepsilon}
\newcommand{\fracpart}[1]{\left\{#1\right\}}
\newcommand{\half}{\tfrac12}
\newcommand{\Bcal}{\mathcal B}
\newcommand{\Toperator}{\mathfrak T}
\newcommand{\Qoperator}{\mathfrak Q}
\newcommand{\anloc}{\mathcal A}
\newcommand{\Gcal}{\mathcal G}
\newcommand{\Pcal}{\mathcal P}
\newcommand{\Wcal}{\mathcal W}
\newcommand{\interior}{\operatorname{int}}
\newcommand{\clamp}{\operatorname{clamp}_{[0,1]}}
\newcommand{\Gauss}{\mathfrak G}
\newcommand{\Comp}{\mathcal C}
\newcommand{\Ham}{H}
\newcommand{\Formalized}{\textsc{Lean}}
\newcommand{\Derived}{\textsc{Derived}}
\newcommand{\Main}{\mathcal M}
\newcommand{\fract}[1]{\left\{#1\right\}}
\newcommand{\Deltah}{\Delta_h}
\newcommand{\Btri}[1]{\binom{#1}{2}}
\newcommand{\Unif}{\operatorname{Unif}}
\newcommand{\Mmain}{\mathcal M}
\newcommand{\CoeffMass}{\mathcal C}
\newcommand{\Prefix}{\mathcal S}
\newcommand{\M}{\mathfrak m}
\newcommand{\Lam}{\Lambda}
\newcommand{\epsTM}{\varepsilon}
\newcommand{\conv}{*}
\newcommand{\Binom}[2]{\binom{#1}{#2}}
\newcommand{\iu}{\mathrm i}
\newcommand{\D}{\mathrm D}
\definecolor{codegreen}{RGB}{40,105,72}
\definecolor{codepurple}{RGB}{105,55,125}
\lstdefinestyle{reportcode}{
  backgroundcolor=\color{shadegray},
  basicstyle=\ttfamily\small,
  keywordstyle=\color{linkblue}\bfseries,
  commentstyle=\color{codegreen},
  stringstyle=\color{codepurple},
  numbers=left,
  numberstyle=\tiny\color{gray},
  numbersep=8pt,
  frame=single,
  rulecolor=\color{rulegray},
  breaklines=true,
  showstringspaces=false,
  tabsize=4,
  columns=fullflexible
}

% The merged volume has three-digit global section numbers.  Article's
% default 1.5em number box is too narrow for them, so widen only that local
% ToC box while preserving the house style's spacing and typography.
\makeatletter
\renewcommand*\l@section[2]{%
  \ifnum \c@tocdepth >\z@
    \addpenalty\@secpenalty
    \addvspace{1.0em \@plus\p@}%
    \setlength\@tempdima{2.3em}%
    \begingroup
      \parindent \z@ \rightskip \@pnumwidth
      \parfillskip -\@pnumwidth
      \leavevmode \bfseries
      \advance\leftskip\@tempdima
      \hskip -\leftskip
      #1\nobreak\hfil
      \nobreak\hb@xt@\@pnumwidth{\hss #2\kern-\p@\kern\p@}\par
    \endgroup
  \fi}
\renewcommand*\l@subsection{\@dottedtocline{2}{1.5em}{3.2em}}
\makeatother

\maketitle

\begin{abstract}
This volume consolidates every mathematical document formerly stored in a subdirectory of
\path{Analysis/FabiusFunction/docs/semi-formalized-research-frontiers/}.  Eleven source
notebooks containing natural-language proofs, symbolic derivations, numerical checks,
conjectural extensions, and explanations built on top of the current Lean library are
organized into six thematic syntheses followed by a post-audit gap register.  They remain
quarantined so that no unformalized claim can be mistaken for part of the verified primary
exposition or lost while awaiting a machine-checked proof.
The post-consolidation audit also records the newly formalized generic naturality,
conditional affine-recursion, and exact-support API for weighted uniform series and the named
geometric-\(q\) weights, series, probability law, recurrence, and reflection.  Every
norm-summable real weighted law with a nonzero coefficient is absolutely continuous with
respect to Lebesgue measure and has null singletons.  For every \(|q|<1\) the geometric law
has those properties and support equal to its exact series range, and for \(0\le q<1\) that
support is \([0,1]\).  The formal API now also names its CDF, proves CDF continuity and
reflection for \(|q|<1\), the exterior tails for \(0\le q<1\), and, for \(0<q<1\), the
conditioning identities, an explicit compactly supported Radon--Nikodym density, and
real-space \(C^\infty\) regularity.  Its CDF and density have exact half-base bridges to the
bounded Fabius function and the affine Rvachev bump.  A generic characteristic-function
argument now proves uniqueness for contractive affine independent-copy equations in complete
second-countable real inner-product spaces; in particular, the named geometric law is the
unique probability solution of its refinement equation for \(|q|<1\), including \(q=0\)
and negative \(q\), with no support, density, or moment hypothesis.  The selected total
density formula is now diagnosed exactly on the nonpositive side: it is identically zero
at \(q=0\) and nonpositive for \(q<0\), so its \texttt{ofReal}-clamped
\texttt{withDensity} measure is zero and cannot equal the geometric-uniform probability law
when \(|q|<1\) and \(q\le0\).  Lean formalization of the paper-level corrected
signed-density replacement for that parameter range remains open.  Write
\(S_q(z)=\operatorname{geometricSincProduct}(q,z)\),
\(G_q(z)=\operatorname{geometricReciprocalGamma}(q,z)\), and
\(z_q(t)=(1-q)t/(2\pi)\).  The formal API now proves, for every real
\(|q|<1\) and \(t\in\RealNumbers\), including \(q=0\) and negative \(q\),
\[
 \phi_q(t)=e^{\ImaginaryUnit t/2}S_q(z_q(t))
 =e^{\ImaginaryUnit t/2}G_q(z_q(t))G_q(-z_q(t)).
\]
It also gives locally uniform convergence on \(\ComplexNumbers\) of the defining products for \(S_q\),
their pointwise multipliability and exact \texttt{HasProd} identity, and entireness of
\(S_q\).  For every fixed complex \(\lVert q\rVert<1\), the single symbol
\(a\mapsto(a;q)_\infty\) now also has a locally uniform defining product, is
entire, vanishes exactly at a displayed factor (equivalently at
\(a=(q^j)^{-1}\) when \(q\ne0\)), and has analytic order one at every zero.
For every real \(|q|<1\), the full phase-bearing prefixes now converge locally
uniformly on the real frequency line, hence uniformly on every compact real-frequency
set, including at \(q=0\) and for negative \(q\).  A corresponding complex-frequency
locally uniform theorem for the full phase-bearing prefix sequence, local-uniform or normal
convergence of the outer spectral Pochhammer product, general-\(q\) rapid decay
and Fourier inversion, and the further centered/MGF, shape, explicit
Bernoulli-cumulant/Bell-moment, transform, derivative-hierarchy, and endpoint-flatness
layers remain in the frontier on the precise boundary recorded below.

Closely overlapping presentations have been merged rather than concatenated.  Each topic
uses the strongest or most complete source as its backbone, imports independent proofs,
sharper constants, examples, warnings, and open obligations from the companion source, and
removes repeated definitions and theorem statements.  The original paths, historically
recorded SHA-256 values, and repository blob identifiers remain in provenance banners and
the repository ledger.  Where a printed SHA-256 value cannot be reproduced from a reachable
blob, it is marked unverified rather than presented as authenticated provenance.  Status
distinctions, citations, tables, algorithms, plots, and genuine alternative proof mechanisms
are retained.
\end{abstract}

\begin{boxedremark}
\textbf{Formalization boundary.}
Nothing is promoted merely by appearing in this volume.  A claim may enter
\path{docs/Fabius_Function_and_Rvachev_Up/Fabius_Function_and_Rvachev_Up.tex}
only after an exact Lean theorem has been checked for the full hypotheses and conclusion,
the declaration and module have been recorded, and the corresponding frontier obligation
has been dispositioned.  A related definition, a weaker theorem, numerical agreement, or a
paper citation is not sufficient.  Some sections quote formalized inputs; that does not
formalize their later deductions.
\end{boxedremark}

\begin{boxedremark}
\textbf{Current source and retained PDF.}
The retained PDF contains the 31 August 2026 Legendre Gaunt--Wigner-square
closed-form overlay.  This source additionally contains the 1 September 2026
five-theorem fixed-nome q-Pochhammer entire-function overlay and the later
zero-definition/six-theorem generalized Rvachev identifiability overlay
described in the audit map below, as well as the six-theorem prime-power
companion-row and zero-definition/three-theorem outer Pochhammer
normal-convergence overlays, the effective-inverse closure, and the latest
expanded q-series/q-calculus union, including the finite root-of-unity,
q-Catalan, Newton-interpolation, q-beta, upper-index Gaussian,
q-Pfaff--Saalsch\"utz, quantum-multinomial, and Gaussian-bound layers, together
with the completed Gaussian linear-coefficient classifier, the
eight-declaration rational-gap inverse module, and the parity-selected Rvachev
superconvergence leaf and the subsequent binary-word, box-partition, and
telescoping-certificate leaves, as well as the exact three-module Lambert
branch-pairing, gap-bijection, and symmetric-identity tranche (respectively
0+7, 4+16, and 0+9), its subsequent zero-definition/five-theorem
Bernoulli branch-gap companion crosswalking
\nolinkurl{eq:pair-Bernoulli-general} and, under the standard removable-origin
convention, \nolinkurl{eq:bernoulli-gen}; those four Lambert modules total four
definitions and 37 theorems, 41 declarations; the exact inverse-recursion
certificates, the Jacobi two-square counting and Lagrange--Rvachev Matrix
leaves, the weighted/geometric barycenter pair, and the
three-definition/seven-theorem geometric Richardson generating-function
module; the three-theorem second-moment strengthening of the existing
Gaussian-binomial cumulant module; the one-definition/seventeen-theorem
arbitrary-space geometric-uniform realization module; the two-definition/
one-theorem regular central q-binomial sum leaf; the zero-definition/
ten-theorem effective fixed-column Gaussian-rate leaf; and the one-definition/
fourteen-theorem Rvachev--Appell Hasse leaf; the zero-definition/two-theorem
greater-than-one Gaussian asymptotics leaf; the one-definition/fourteen-theorem
nodes-only Rvachev amplitude leaf; and the one-definition/eight-theorem
geometric-uniform moment-polynomial algebra leaf; and the zero-definition/
three-theorem Thue--Morse GammaLog differential leaf; and the one-theorem
scaled-geometric polynomial-evaluation strengthening of the existing
\path{GeometricResidualMoments.lean} module; and the zero-definition/one-
theorem real-MGF bridge in
\path{GeometricUniformMomentPolynomialBridge.lean}; the affine-Prouhet
strengthening of the existing zero-definition/sixteen-theorem
\path{FinitePolynomialFunctional.lean} module; the one-definition/four-
theorem local corner-integral leaf \path{ThueMorseCornerIntegral.lean}; and
the zero-definition/three-theorem central Legendre cancellation leaf
\path{RvachevLegendreCentralSum.lean}; and the one-definition/two-theorem
complex moment-product leaf
\path{GeometricUniformComplexMomentProduct.lean}; the zero-definition/one-
theorem half-base root-simplicity leaf
\path{HalfQBinomialRootSimplicity.lean}; the one-definition/two-theorem
exterior reciprocal-germ leaf
\path{GeometricUniformExteriorComplexMomentGerm.lean}; and the zero-definition/
three-theorem sharp coefficient-and-degree leaf
\path{GeometricUniformMomentPolynomialDegree.lean}; the one-definition/six-
theorem \path{RvachevLaurentLeading.lean} leaf; and the eleven-definition/
seventeen-theorem \path{FinitePrefixAppellRecovery.lean} leaf; and the
one-definition/four-theorem
\path{GeometricUniformMomentRatFunc.lean} leaf; the two-theorem extension of
the existing \path{ProbabilityLaplaceMoments.lean} module; and the exhaustive
one-definition/one-theorem
\path{RvachevLegendreBiorthogonality.lean} leaf.  The live source census is
932 facade-reachable modules and 11,688 public declarations, with zero missing
module headers and zero missing declaration comments.  The
Legendre-biorthogonality 925/11,619 census remains an explicit historical
checkpoint.
The sharp
leaf makes \nolinkurl{prop:qF-P-degree-sharp} Exact and supplies the formerly
missing leading/odd-degree clause, so the represented
\nolinkurl{p7:thm:Pn} is Exact.  The new RatFunc leaf packages a single global
rational coefficient, proves its q-factorial clearing, supplies safe inner and
exterior analytic evaluations, and handles the removable value at \(q=1\);
hence \nolinkurl{thm:qF-moment-polynomial} is Exact.  The actual-MGF
normalization remains Exact for real \(\lvert q\rvert<1\), while no analytic
value at a genuine unit-root pole is asserted.  The incoming bridge branch recorded the
successive historical checkpoints 903/11,447,
903/11,448, 904/11,457, and 905/11,458.  In the merged chronology the
real-MGF bridge first reached
915/11,556; the affine theorem, local corner leaf, and central Legendre leaf
then give the successive checkpoints 915/11,557, 916/11,562, and 917/11,565;
the complex product leaf gave the historical 918/11,568 checkpoint.  The
half-base root-simplicity leaf gave the actual merged-main pre-local checkpoint
919/11,569, the exterior leaf gave 920/11,572, and the sharp leaf gave the
historical 921/11,575 checkpoint.  The Laurent leaf then gave 922/11,582, and
the finite-prefix leaf gave the historical 923/11,610 checkpoint.  The RatFunc
leaf gives the historical 924/11,615 checkpoint, the two probability theorems
give 924/11,617, and the Legendre biorthogonality leaf gives the historical
925/11,619 checkpoint.  The subsequent merged source union gives the live
932-module/11,688-public-declaration census with zero documentation gaps.  The incoming-branch
906/11,461 and earlier exterior-branch 907/11,464 censuses remain explicitly
historical.
The retained
257-page, 2,438,299-byte A4 PDF, SHA-256
\nolinkurl{3766761aac90247061f5c955dc84a0feb8567454e10839f1508b9431797ee980},
predates these source-only overlays and the expanded ledger and is therefore a
historical render, not a rendering of the current TeX.
\end{boxedremark}

\tableofcontents
\clearpage

\markboth{Consolidation map and reading conventions}{}
\section*{Consolidation map and reading conventions}
\addcontentsline{toc}{section}{Consolidation map and reading conventions}
Each part below is a thematic synthesis.  Its provenance banner names every source notebook
that contributed to it and records the historically supplied SHA-256 metadata and repository
blob evidence available for that source.  The final register records claims removed by the
primary exposition's exact-counterpart audit.  Namespaced labels and bibliography keys are
retained for stable cross-references, but a prefix no longer means that the surrounding prose
is an untouched source block.

Notation has been standardized at the collisions that affect mathematical meaning.  In
particular, \(\operatorname{App}_n\) denotes the dyadic Appell family while \(A_j\) denotes
the saddle logarithmic coefficients, and cardinal weights use \(\pi_{n,k}\) rather than the
\(\lambda\)-symbol reserved for lattice centers and the lower-Lambert phase.  The final
saddle part keeps its historically local jet symbols \(p_n,d_n,E_m\), but states their
crosswalk explicitly so that its \(d_n\) cannot be confused with the half moments in the
probability and dyadic parts.  Each Fourier chapter states its exponential convention before
formulas involving \(2\pi\).

\begin{longtable}{@{}>{\raggedright\arraybackslash}p{0.21\textwidth}
                         >{\raggedright\arraybackslash}p{0.33\textwidth}
                         >{\raggedright\arraybackslash}p{0.40\textwidth}@{}}
\toprule
Part & Principal subject & Source integration \\
\midrule
\endhead
Probability and generic engines & Probability/Laplace bridges, analytic-density obstruction, weighted paths, Toeplitz convergence, formal exp/log, and Gaussian estimates & Unique Integration material retained; its duplicate inverse and Fabius-specific saddle chapters are cross-referenced to the final part. \\
Repeated integration & Finite box splines, exact derivative thresholds, error kernels, extremizers, and weak-test consequences & Sharper second notebook is the backbone; historical provenance, operator, \(L^p\), weak-error, and comparison results are imported from the first. \\
Exact dyadic web & Moments, Taylor reduction, Appell compression, determinant, transforms, q-binomial rows, stabilization, and Richardson limits & Later Dyadic Web is the backbone; unique Appell, simplex, determinant, prefix-block, filter, and Fourier routes are imported from q-Connections. \\
Dyadic endpoint asymptotics & Exact negative-Laplace decomposition, Lambert phase, vertical saddle, fixed-numerator rays, and arithmetic consequences & Later formulae-to-asymptotics notebook is the backbone; Mellin, coefficient-extraction, Bromwich, squeeze, and numerical material is imported from the bridge. \\
Thue--Morse convergence & Exact discrete bridge, deconvolution, sharp/all-order rates, finite certificates, endpoint identities, and acceleration & Sharper plural-rate notebook is the backbone; finite elementary bounds, zero runs, binary phase, and endpoint formulas are imported from the earlier report. \\
Inverse and small-argument saddle & Total inverse, exact regularity and non-elementarity, endpoint inversion, closed jets, corrections, rigorous all orders, Fourier amplitudes, and numerics & Inverse/Saddle is the backbone; phase-polynomial and right-endpoint transfers, exact coefficient means, higher explicit coefficients, Fourier amplitudes, numerical checks, traps, and the open ledger are imported from Small Argument. \\
Primary-exposition gap register & Claims removed by exact-counterpart audits & Post-consolidation disposition record; every candidate remains open until an exact current Lean theorem is audited. \\
\bottomrule
\end{longtable}

\markboth{Lean object crosswalk}{}
\section*{Lean object crosswalk---inputs, not certification}
\addcontentsline{toc}{section}{Lean object crosswalk---inputs, not certification}
\begin{longtable}{@{}p{0.24\textwidth}p{0.34\textwidth}p{0.34\textwidth}@{}}
\toprule
Object & Representative Lean name & Principal module family \\
\midrule
\endhead
Bounded Fabius function & \nolinkurl{fabiusReal}, \nolinkurl{IsFabius} & \nolinkurl{Basic.lean}, \nolinkurl{PaperStatements.lean} \\
Rvachev up-function & \nolinkurl{rvachevUp} & \nolinkurl{Basic.lean}, \nolinkurl{Differential.lean} \\
Signed global extension & \nolinkurl{extendedFabius} & \nolinkurl{GlobalExtension.lean} \\
Fractional Volterra calculus &
\nolinkurl{fractionalVolterra},
\nolinkurl{Fabius.intervalIntegral_eq_integral_min_of_eq_zero},
\nolinkurl{Fabius.fractionalVolterra_eq_intervalIntegral_min_of_eq_zero},
\nolinkurl{fractionalVolterra_add},
\nolinkurl{fractionalVolterra_affine}, \nolinkurl{fractionalVolterra_add_one_deriv} &
\nolinkurl{FractionalVolterra.lean}, \nolinkurl{FractionalVolterraSemigroup.lean},
\nolinkurl{FractionalVolterraCalculus.lean} \\
Fractional causal Rvachev primitives and Fabius--Rvachev shifts &
\nolinkurl{Fabius.rvachevFractionalPrimitive},
\nolinkurl{Fabius.rvachevFractionalPrimitive_eq_intervalIntegral_min},
\nolinkurl{Fabius.rvachevFractionalPrimitive_nat_succ},
\nolinkurl{Fabius.rvachevFractionalPrimitive_add},
\nolinkurl{fractionalVolterra_add_one_extendedFabius_of_nonneg},
\nolinkurl{fractionalVolterra_add_one_fabiusReal},
\nolinkurl{fractionalVolterra_add_one_rvachevUp} &
\nolinkurl{FabiusFractionalVolterra.lean} \\
Pointwise analytic finite filters &
\nolinkurl{finiteAnalyticSeriesFilter},
\nolinkurl{geometricAnalyticSeriesFilter},
\nolinkurl{geometricLagrangeAnalyticSeriesFilter_shifted} &
\nolinkurl{AnalyticSeriesFilter.lean} \\
Exact dyadic evaluator & \nolinkurl{fabiusDyadicValue} & \nolinkurl{Arithmetic.lean}, \nolinkurl{PaperStatements.lean} \\
Binary and q-binomial series & \nolinkurl{fabiusReductionSum} and global q-series declarations & \nolinkurl{TaylorReduction.lean}, \nolinkurl{FabiusGlobalQBinomialSeries.lean} \\
Thue--Morse grid samples & \nolinkurl{correctedPrefixGridSample} & \nolinkurl{ThueMorseApproximation.lean} \\
Finite complex Thue--Morse block transforms &
\nolinkurl{sum_range_two_pow_binaryWeight_cexp},
\nolinkurl{sum_range_two_pow_binaryWeight_cexp_affine},
\nolinkurl{sum_range_two_pow_thueMorseSign_cexp},
\nolinkurl{sum_range_two_pow_thueMorseSign_cexp_affine} &
\nolinkurl{ThueMorseComplexExponential.lean} \\
Centered mixed differences and symmetric Thue--Morse blocks (exactly two
definitions and eleven theorems).  The arbitrary-half-step layer uses
additive commutative groups and an additive action, permits repeated values,
and gives the empty, singleton, insert, and powerset laws. &
\nolinkurl{symmetricMixedDifference},
\nolinkurl{symmetricMixedDifference_empty},
\nolinkurl{symmetricMixedDifference_singleton},
\nolinkurl{symmetricMixedDifference_insert},
\nolinkurl{symmetricMixedDifference_eq_sum_powerset_smul} &
\nolinkurl{ThueMorseSymmetricDifference.lean} \\
\emph{Same surface, polynomial layer.}  Over a commutative ring, for a
polynomial of degree at most \(|s|\), it extracts the degree-\(|s|\)
coefficient with sign-free factor
\(|s|!\prod_{i\in s}(2a_i)\), annihilates lower degree including the
zero-polynomial/empty-support boundary, and evaluates the first surviving
power. &
\nolinkurl{symmetricMixedDifference_polynomial_eq_coeff_card},
\nolinkurl{symmetricMixedDifference_polynomial_of_degree_lt},
\nolinkurl{symmetricMixedDifference_pow_card} &
\emph{same module} \\
\emph{Same surface, dyadic layer and report boundary.}  The dyadic block is
group-valued.  Its increasing-grid wrappers use a characteristic-zero field:
the all-order form, including \(m=0\), has points
\(x-(1-2^{-m})+2n2^{-m}\).  At positive order \(N=m+1\), the full point is
\(x-(1-2^{-(m+1)})+n/2^m\); its varying term is \(n/2^m\), its mesh spacing
is \(1/2^m\), and its complement sign is \((-1)^N\).  This closes the
Boolean-cube, polynomial, and affine-grid clauses of the continuous-chaos
report's Thue--Morse corner theorem. &
\nolinkurl{symmetricDyadicMixedDifference},
\nolinkurl{symmetricDyadicMixedDifference_zero},
\nolinkurl{symmetricDyadicMixedDifference_eq_sum_thueMorseSign_smul},
\nolinkurl{symmetricDyadicMixedDifference_inv_two_pow_eq_sum_thueMorseSign_smul},
\nolinkurl{symmetricDyadicMixedDifference_inv_two_pow_succ_eq_sum_thueMorseSign_smul} &
\emph{same module} \\
Local centered-box integral (one definition and four theorems).  For
nonnegative half-steps it proves the repeated-integral identity under exactly
\(\lean{IsOpen}\,I\), \(\lean{OrdConnected}\,I\),
\(\lean{ContDiffOn}\,\RealNumbers\,N\,g\,I\), and containment in \(I\) of the
full closed symmetric segment.  The result is local, real-valued, fixes the
recursive nesting order, and includes zero steps and \(N=0\).  Together with
the preceding algebra it makes \nolinkurl{thm:TM-corner} Exact/Complete exactly by
composition.  It does not prove the following Walsh conditional-expectation
corollary or its \(2^{-N}\) normalization. &
\nolinkurl{centeredBoxIntegral},
\nolinkurl{centeredBoxIntegral_zero},
\nolinkurl{centeredBoxIntegral_succ},
\nolinkurl{symmetricMixedDifference_range_eq_centeredBoxIntegral},
\nolinkurl{symmetricMixedDifference_univ_eq_centeredBoxIntegral} &
\nolinkurl{ThueMorseCornerIntegral.lean} \\
Prime-power Pascal companion rows &
\nolinkurl{primePowerChoose_padicValNat_add},
\nolinkurl{primePowerChoose_padicValNat},
\nolinkurl{primePowerSubOneChoose_padicValNat},
\nolinkurl{primePowerSubTwoChoose_padicValNat},
\nolinkurl{twoPowChoose_padicValNat},
\nolinkurl{twoPowSubTwoChoose_padicValNat} &
\nolinkurl{PrimePowerBinomialValuation.lean}; zero definitions and exactly
six theorems.  Row \(p^m-1\) is a \(p\)-adic unit row for \(j<p^m\), and
\(\nu_p\binom{p^m-2}{j-1}=\nu_p(j)\) holds for exactly \(0<j<p^m\), with a
dyadic wrapper.  The upper companion boundary is excluded because its
binomial coefficient is zero. \\
Generic weighted uniform series and laws &
\nolinkurl{uniformProduct_map_of_ae_head_tail_recursion},
\nolinkurl{weightedUniformSeries_map},
\nolinkurl{weightedUniformDistribution_split_of_shift_eq_map},
\nolinkurl{weightedUniformDistribution_absolutelyContinuous},
\nolinkurl{weightedUniformDistribution_nullSingletonClass},
\nolinkurl{weightedUniformDistribution_support_eq_range},
\nolinkurl{weightedUniformDistribution_support_eq_Icc_min_max} &
\nolinkurl{WeightedUniformSeries.lean}, \nolinkurl{WeightedUniformDistribution.lean},
\nolinkurl{WeightedUniformSupport.lean}, \nolinkurl{ProbabilityRepresentation.lean} \\
Contractive affine independent-copy laws &
\nolinkurl{affineIndependentCopyLaw},
\nolinkurl{charFun_affineIndependentCopyLaw},
\nolinkurl{charFun_iterate_of_affineIndependentCopy_fixedPoint},
\nolinkurl{affineIndependentCopyLaw_fixedPoint_unique},
\nolinkurl{affineIndependentCopy_map_fixedPoint_unique} &
\nolinkurl{AffineIndependentCopy.lean} \\
Named geometric uniform law &
\nolinkurl{geometricUniformWeight},
\nolinkurl{geometricUniformDistribution_selfSimilar},
\nolinkurl{geometricUniformDistribution_absolutelyContinuous},
\nolinkurl{geometricUniformDistribution_nullSingletonClass},
\nolinkurl{geometricUniformDistribution_support_eq_range},
\nolinkurl{geometricUniformDistribution_support_eq_Icc},
\nolinkurl{weightedSumDistribution_support_eq_Icc},
\nolinkurl{eq_geometricUniformDistribution_of_selfSimilar} &
\nolinkurl{GeometricUniformLaw.lean}, \nolinkurl{GeometricUniformUniqueness.lean},
\nolinkurl{ProbabilityRepresentation.lean} \\
Geometric sinc-prefix factorization &
Closed digit factor, exact finite residual factorization, pointwise phase-bearing prefix
limit, and dyadic weighted-sum wrappers &
\nolinkurl{GeometricUniformDictionary.lean},
\nolinkurl{GeometricSincFactorization.lean} \\
Geometric sinc-product convergence &
Locally uniform convergence, pointwise multipliability and exact \nolinkurl{HasProd}, and
entireness of the geometric sinc product &
\nolinkurl{GeometricReciprocalGamma.lean} \\
Generalized Rvachev exponent identifiability &
\nolinkurl{analyticOrderAt_generalizedRvachevProduct_two_pow},
\nolinkurl{exponent_succ_eq_toNat_analyticOrderAt_generalizedRvachevProduct},
\nolinkurl{generalizedRvachevProduct_eq_iff} &
\nolinkurl{GeneralizedRvachevIdentifiability.lean}; zero definitions and
exactly six theorems \\
Geometric characteristic/sinc/Gamma bridges &
Exact phase-bearing characteristic-function identities through the rescaled geometric sinc
product and paired reciprocal-Gamma factors, together with locally uniform real-frequency
convergence of the full phase-bearing prefixes and its compact-set specialization, for every
real \(\lvert q\rvert<1\) &
\nolinkurl{GeometricSincCharacteristicFunction.lean} \\
Fixed-nome entire complex q-Pochhammer symbol &
\nolinkurl{hasProdLocallyUniformly_complexQPochhammerInf},
\nolinkurl{complexQPochhammerInf_differentiable},
\nolinkurl{complexQPochhammerInf_eq_zero_iff},
\nolinkurl{complexQPochhammerInf_eq_zero_iff_eq_inv_pow},
\nolinkurl{analyticOrderAt_complexQPochhammerInf_of_eq_zero} &
\nolinkurl{QPochhammerEntire.lean}; zero definitions and exactly five theorems,
under exactly \(\lVert q\rVert<1\); the division-free factor-zero theorem
includes \(q=0\), while the reciprocal-power form assumes \(q\ne0\) \\
Outer geometric spectral Pochhammer normal convergence &
\nolinkurl{hasProdLocallyUniformly_geometricSincProduct_complexQPochhammerInf},
\nolinkurl{hasProdLocallyUniformly_rvachevFourierProduct_complexQPochhammerInf},
\nolinkurl{hasProdLocallyUniformly_rvachevFourier_complexQPochhammerInf} &
\nolinkurl{GeometricPochhammerNormalConvergence.lean}; zero definitions and
exactly three theorems.  The general theorem assumes exactly
\(\lVert q\rVert<1\), includes \(q=0\), and the last two are the dyadic
Rvachev/Fabius specializations.  Centered/MGF, outside-disk reciprocal,
pole-divisor, and zero--pole clauses remain outside Lean. \\
Finite q-Pochhammer dissection &
\nolinkurl{finiteQPochhammerIn_dissection},
\nolinkurl{finiteQPochhammerIn_dissection_remainder} &
\nolinkurl{QPochhammerDissection.lean}; zero definitions and exactly two
theorems over every commutative ring, with the remainder theorem assuming
exactly \(u\le r\) \\
Generic infinite q-Pochhammer product &
\nolinkurl{qPochhammerInfIn}; all twenty-nine public theorem names and their
hypothesis boundaries are listed in the exact crosswalk below &
\nolinkurl{QPochhammerInfinite.lean}; one definition and exactly twenty-nine
theorems, distinct from the complex analytic-order wrapper above \\
Infinite q-binomial theorem and Euler expansions &
\nolinkurl{gaussianMajorant}, \nolinkurl{hasSum_qBinomial_theorem},
\nolinkurl{hasSum_euler_product}, \nolinkurl{hasSum_euler_reciprocal} &
\nolinkurl{QBinomialTheoremInfinite.lean}; one definition and twenty-two
theorems, inventoried below \\
Effective fixed-column Gaussian limit and rates &
\nolinkurl{tendsto_gaussianBinomial_add_atTop},
\nolinkurl{gaussianBinomial_fixedColumn_relativeError_isBigO},
\nolinkurl{gaussianBinomial_shifted_fixedColumn_relativeError_isBigO} &
\nolinkurl{GaussianBinomialFixedColumnRate.lean}; zero definitions and ten
theorems, inventoried below.  Together with the existing fixed-column limit,
this makes \nolinkurl{thm:fixed-column-limit} exact, including \(q=0\). \\
Jacobi triple product and pentagonal sums &
\nolinkurl{thetaExponent}, \nolinkurl{pentagonalExponent},
\nolinkurl{hasSum_jacobi_triple_product}, \nolinkurl{hasSum_pentagonal_pairs} &
\nolinkurl{JacobiTripleProduct.lean}; two definitions and twenty-five
theorems, inventoried below \\
Finite q-Pascal summations &
\nolinkurl{sum_gaussianBinomial_succ_mul},
\nolinkurl{sum_gaussianBinomial_succ_mul'},
\nolinkurl{Commute.gaussianBinomial_left},
\nolinkurl{Commute.gaussianBinomial_right} &
\nolinkurl{QPascalSummation.lean}; zero definitions and four theorems \\
Gaussian-binomial continuity and classical limit &
\nolinkurl{continuous_gaussianBinomial},
\nolinkurl{tendsto_gaussianBinomial_nhds_one},
\nolinkurl{gaussianBinomial_eq_finiteQPochhammerIn_div} &
\nolinkurl{GaussianBinomialContinuity.lean}; zero definitions and three theorems \\
Generic Gaussian-polynomial degree, palindromicity, mean, and linear coefficient &
\nolinkurl{reflect_add_of_natDegree_le}, \nolinkurl{reflect_one'},
\nolinkurl{gaussianBinomial_natDegree_le},
\nolinkurl{gaussianBinomial_zero_left}, \nolinkurl{gaussianBinomial_diag'},
\nolinkurl{reflect_gaussianBinomial},
\nolinkurl{coeff_gaussianBinomial_reflect},
\nolinkurl{coeff_gaussianBinomial_zero},
\nolinkurl{coeff_gaussianBinomial_top},
\nolinkurl{gaussianBinomial_natDegree}, \nolinkurl{gaussianBinomial_monic},
\nolinkurl{two_mul_derivative_gaussianBinomial_eval_one},
\nolinkurl{coeff_gaussianBinomial_one_of_pos_of_lt},
\nolinkurl{coeff_gaussianBinomial_one} &
\nolinkurl{GaussianBinomialPalindromic.lean}; zero definitions and fourteen theorems \\
Universal Gaussian-polynomial degree and palindromicity &
\nolinkurl{natDegree_gaussianBinomial_universal},
\nolinkurl{gaussianBinomial_universal_monic},
\nolinkurl{coeff_zero_gaussianBinomial_universal},
\nolinkurl{gaussianBinomial_universal_reflect},
\nolinkurl{coeff_gaussianBinomial_universal_symm} &
\nolinkurl{GaussianBinomialPolynomialStructure.lean}; zero definitions and five theorems \\
Primitive-root Gaussian block &
\nolinkurl{gaussianBinomial_isPrimitiveRoot_eq_zero},
\nolinkurl{neg_one_pow_mul_pow_choose_two},
\nolinkurl{finiteQPochhammerIn_isPrimitiveRoot} &
\nolinkurl{PrimitiveRootBlock.lean}; zero definitions and three theorems \\
The q-Lucas theorem &
\nolinkurl{add_mul_add_sub_one},
\nolinkurl{choose_two_add}, \nolinkurl{coeff_finiteQPochhammerIn_neg_X},
\nolinkurl{finiteQPochhammerIn_neg_X_block},
\nolinkurl{coeff_block_pow_mul}, \nolinkurl{pow_choose_two_add_mul_eq},
\nolinkurl{gaussianBinomial_q_lucas} &
\nolinkurl{QLucas.lean}; zero definitions and seven theorems; its local
\nolinkurl{two_mul_choose_two} is private, while the public theorem belongs to
\nolinkurl{QChuVandermonde.lean} \\
Terminating two-phi-one reversal &
\nolinkurl{twoPhiOneFinite}, \nolinkurl{twoPhiOneReflection},
\nolinkurl{twoPhiOneFinite_reversal}, \nolinkurl{twoPhiOne_reversal},
\nolinkurl{twoPhiOne_reversal_twice} &
\nolinkurl{TwoPhiOneReversal.lean} 2+12; the actual-tsum reversal and its
involutive parameter map retain their exact displayed nonvanishing hypotheses \\
The two q-Chu--Vandermonde sums &
\nolinkurl{finiteQPochhammerIn_div_eq_sum_chu},
\nolinkurl{finiteQPochhammerIn_div_eq_sum_chu_second},
\nolinkurl{q_chu_vandermonde_first},
\nolinkurl{q_chu_vandermonde_second},
\nolinkurl{q_chu_vandermonde_second_by_reversal},
\nolinkurl{twoPhiOne_q_chu_vandermonde_first},
\nolinkurl{twoPhiOne_q_chu_vandermonde_second} &
\nolinkurl{QChuVandermonde.lean} 0+10; both displayed formulas are exact on
their full rational domain, while the named reversal proof retains two
additional nonvanishing hypotheses \\
Jacobi two-square count and Lambert forms &
\nolinkurl{sumSqRep_two_eq_four_mul_twoSquareDivisorSum},
\nolinkurl{sumSqRep_two_eq_four_mul_prod},
\nolinkurl{theta_sq_eq_chi4_lambert},
\nolinkurl{theta_sq_eq_odd_lambert} &
\nolinkurl{JacobiTwoSquareCount.lean}; zero definitions and four theorems;
the product form retains its valuation hypothesis and both Lambert forms are
unconditional in a complete normed field under \(\lVert q\rVert<1\) \\
Cyclotomic carry criterion and primitive-root multiple value &
\nolinkurl{cyclotomic_exponent_eq_one_iff},
\nolinkurl{cyclotomic_dvd_gaussianBinomial_iff},
\nolinkurl{gaussianBinomial_mul_isPrimitiveRoot} &
\nolinkurl{CyclotomicDivisibility.lean}; zero definitions and three theorems \\
MacMahon q-Catalan polynomial &
\nolinkurl{qCatalan}; \nolinkurl{map_qInt}, \nolinkurl{qInt_X_monic},
\nolinkurl{qInt_X_natDegree}, \nolinkurl{X_sub_one_mul_qInt},
\nolinkurl{qInt_X_eq_prod_cyclotomic},
\nolinkurl{qInt_X_dvd_gaussianBinomial_rat},
\nolinkurl{qInt_X_dvd_gaussianBinomial_int},
\nolinkurl{qInt_X_mul_qCatalan}, \nolinkurl{qCatalan_natDegree},
\nolinkurl{qCatalan_eval_one_mul}, \nolinkurl{qCatalan_eval_one} &
\nolinkurl{QCatalan.lean}; one definition and eleven theorems \\
Newton interpolation and geometric divided differences &
\nolinkurl{newtonCoeff}, \nolinkurl{nodeNewtonPoly},
\nolinkurl{newtonInterpolant};
\nolinkurl{newtonCoeff_eq}, \nolinkurl{newtonCoeff_zero},
\nolinkurl{newtonCoeff_mul_prod}, \nolinkurl{newtonPoly_succ},
\nolinkurl{eval_newtonPoly}, \nolinkurl{degree_newtonPoly_lt},
\nolinkurl{newtonPoly_eq_interpolate},
\nolinkurl{eq_newtonPoly_of_eval_eq},
\nolinkurl{coeff_newtonPoly_self}, \nolinkurl{newtonCoeff_eq_sum},
\nolinkurl{nodal_range_pow}, \nolinkurl{prod_erase_pow_sub_pow},
\nolinkurl{newtonCoeff_pow_eq_sum}, \nolinkurl{newtonPoly_succ},
\nolinkurl{eval_newtonPoly}, \nolinkurl{degree_newtonPoly_lt},
\nolinkurl{newtonPoly_eq_interpolate},
\nolinkurl{eq_newtonPoly_of_eval_eq}, \nolinkurl{coeff_newtonPoly_self} &
\nolinkurl{NewtonInterpolation.lean}; three definitions and nineteen theorems \\
Jackson q-beta integral &
\nolinkurl{qBeta}; \nolinkurl{qNumber_pos}, \nolinkurl{qBeta_term_eq},
\nolinkurl{qBeta_eq_prod}, \nolinkurl{qBeta_eq_qGamma},
\nolinkurl{qBeta_comm}, \nolinkurl{qBeta_pos},
\nolinkurl{qBeta_add_one_left}, \nolinkurl{qBeta_add_one_right} &
\nolinkurl{QBetaIntegral.lean}; one definition and eight theorems \\
Quantum binomial theorem &
\nolinkurl{quantumPlane_mul_pow}, \nolinkurl{quantum_binomial} &
\nolinkurl{QuantumBinomial.lean}; zero definitions and two noncommutative theorems \\
Rogers--Szeg\H{o} polynomials &
\nolinkurl{rogersSzego}, \nolinkurl{rogersSzego_succ},
\nolinkurl{rogersSzego_three_term},
\nolinkurl{hasSum_rogersSzego_generating} &
\nolinkurl{RogersSzegoPolynomial.lean}; one definition and nine theorems \\
Generic continuous-CDF bridge &
\nolinkurl{continuous_cdf_of_nullSingleton},
\nolinkurl{cdf_reflection_sub},
\nolinkurl{measure_eq_withDensity_of_cdf_hasDerivAt} &
\nolinkurl{ContinuousCDF.lean} \\
Geometric uniform CDF and density &
\nolinkurl{geometricUniformCDF},
\nolinkurl{geometricUniformDensity},
\nolinkurl{geometricUniformCDF_eq_intervalIntegral},
\nolinkurl{geometricUniformDistribution_eq_withDensity},
\nolinkurl{contDiff_geometricUniformCDF},
\nolinkurl{contDiff_geometricUniformDensity} &
\nolinkurl{GeometricUniformCDF.lean} \\
Half-base CDF and density bridges &
\nolinkurl{weightedSumCDF_eq_geometricUniformCDF_one_half},
\nolinkurl{geometricUniformCDF_one_half_eq_fabiusReal},
\nolinkurl{geometricUniformDensity_one_half_eq_rvachevUp} &
\nolinkurl{ProbabilityRepresentation.lean} \\
Generic finite-measure Cauchy calculus &
\nolinkurl{measureCauchyDomain}, \nolinkurl{measureCauchyTransform},
\nolinkurl{measureCauchyPower}, \nolinkurl{measureCauchyTransform_apply},
\nolinkurl{measureCauchyPower_one}, \nolinkurl{measureCauchyPower_map_affine},
\nolinkurl{measureCauchyTransform_map_affine}, \nolinkurl{isOpen_measureCauchyDomain},
\nolinkurl{hasDerivAt_measureCauchyTransform},
\nolinkurl{analyticOn_measureCauchyTransform},
\nolinkurl{measureCauchyPower_succ_of_uniformAffineFixedPoint},
\nolinkurl{hasDerivAt_measureCauchyTransform_of_uniformAffineFixedPoint} &
\nolinkurl{MeasureCauchyTransform.lean} \\
Generic bounded-measure centered moment/Laurent calculus
(one definition, fifteen theorems) &
\nolinkurl{measureCauchyMoment}, \nolinkurl{inv_sub_eq_sum_range_add},
\nolinkurl{measureCauchyMoment_zero}, \nolinkurl{measurable_inv_sub_rclike},
\nolinkurl{integrable_centered_pow_div_sub_of_ae_norm_sub_le},
\nolinkurl{integrable_inv_sub_of_ae_norm_sub_le},
\nolinkurl{norm_measureCauchyMoment_le},
\nolinkurl{integral_inv_sub_eq_sum_range_measureCauchyMoment_add},
\nolinkurl{norm_integral_inv_sub_sub_sum_range_measureCauchyMoment_le},
\nolinkurl{summable_measureCauchyMoment_laurent},
\nolinkurl{integral_inv_sub_eq_tsum_measureCauchyMoment},
\nolinkurl{hasSum_measureCauchyMoment_laurent},
\nolinkurl{measureCauchyTransform_eq_sum_range_measureCauchyMoment_add},
\nolinkurl{norm_measureCauchyTransform_sub_sum_range_measureCauchyMoment_le},
\nolinkurl{measureCauchyTransform_eq_tsum_measureCauchyMoment},
\nolinkurl{hasSum_measureCauchyTransform_measureCauchyMoment_laurent} &
\nolinkurl{MeasureCauchyMomentLaurent.lean} \\
Geometric-uniform Cauchy hierarchy &
\nolinkurl{geometricUniformStieltjesDomain},
\nolinkurl{geometricUniformStieltjesTransform},
\nolinkurl{geometricUniformStieltjesPower},
\nolinkurl{geometricUniformStieltjesTransform_apply},
\nolinkurl{geometricUniformStieltjesPower_one},
\nolinkurl{isOpen_geometricUniformStieltjesDomain},
\nolinkurl{analyticOn_geometricUniformStieltjesTransform},
\nolinkurl{hasDerivAt_geometricUniformStieltjesTransform_refinement},
\nolinkurl{geometricUniformStieltjesPower_succ} &
\nolinkurl{GeometricUniformCauchy.lean} \\
Canonical Cauchy--Stieltjes powers and refinement &
\nolinkurl{rvachevCauchyDomain}, \nolinkurl{fabiusStieltjesDomain},
\nolinkurl{rvachevCauchyTransform}, \nolinkurl{fabiusStieltjesTransform},
\nolinkurl{rvachevCauchyPower}, \nolinkurl{fabiusStieltjesPower},
\nolinkurl{rvachevCauchyTransform_apply},
\nolinkurl{fabiusStieltjesTransform_apply},
\nolinkurl{rvachevCauchyPower_one}, \nolinkurl{fabiusStieltjesPower_one},
\nolinkurl{rvachevCauchyTransform_eq_integral_rvachevUp},
\nolinkurl{isOpen_rvachevCauchyDomain}, \nolinkurl{isOpen_fabiusStieltjesDomain},
\nolinkurl{hasDerivAt_rvachevCauchyTransform},
\nolinkurl{hasDerivAt_fabiusStieltjesTransform},
\nolinkurl{analyticOn_rvachevCauchyTransform},
\nolinkurl{analyticOn_fabiusStieltjesTransform},
\nolinkurl{fabiusStieltjesTransform_eq_two_mul_rvachevCauchyTransform},
\nolinkurl{rvachevCauchyTransform_eq_inv_two_mul_fabiusStieltjesTransform},
\nolinkurl{fabiusStieltjesPower_eq_two_pow_mul_rvachevCauchyPower},
\nolinkurl{hasDerivAt_fabiusStieltjesTransform_refinement},
\nolinkurl{fabiusStieltjesPower_succ}, \nolinkurl{rvachevCauchyPower_succ} &
\nolinkurl{CauchyTransform.lean} \\
Radius-one, center-zero up-measure Laurent specialization &
\nolinkurl{ae_norm_sub_zero_le_one_rvachevMeasure},
\nolinkurl{measureCauchyMoment_rvachevMeasure_zero},
\nolinkurl{integral_inv_sub_eq_sum_upMoment_add},
\nolinkurl{abs_integral_inv_sub_sub_sum_le},
\nolinkurl{integral_inv_sub_eq_tsum_upMoment},
\nolinkurl{hasSum_upMoment_laurent},
\nolinkurl{integral_inv_sub_complex_eq_sum_upMoment_add},
\nolinkurl{norm_integral_inv_sub_complex_sub_sum_le},
\nolinkurl{integral_inv_sub_complex_eq_tsum_upMoment},
\nolinkurl{hasSum_upMoment_laurent_complex},
\nolinkurl{rvachevCauchyTransform_eq_tsum_upMoment} &
\nolinkurl{StieltjesMomentLaurent.lean},
\nolinkurl{StieltjesCauchyTransform.lean} \\
Weighted survival, lower-CDF, fractional, and inverse layer cake &
\nolinkurl{integral_smul_setIntegral_subgraph},
\nolinkurl{integral_smul_setIntegral_supergraph},
\nolinkurl{intervalIntegral_survival_smul_eq_integral_clamp},
\nolinkurl{intervalIntegral_cdf_smul_eq_integral_clamp},
\nolinkurl{fractionalVolterra_measureReal_Iic_eq_integral_posPart},
\nolinkurl{fractionalVolterra_fabiusReal_eq_integral_posPart},
\nolinkurl{galoisConnection_Icc_restrict_of_lt_iff_lt},
\nolinkurl{intervalIntegral_smul_comp_fabiusInv},
\nolinkurl{intervalIntegral_mul_comp_sub_of_lt_iff_lt_of_absolutelyContinuousOnInterval},
\nolinkurl{intervalIntegral_mul_comp_of_lt_iff_lt_of_absolutelyContinuousOnInterval},
\nolinkurl{intervalIntegral_mul_comp_fabiusInv_of_absolutelyContinuousOnInterval},
\nolinkurl{intervalIntegral_fabiusInv_eq_one_half} &
\nolinkurl{SubgraphFubini.lean}, \nolinkurl{SurvivalLayerCake.lean},
\nolinkurl{CDFLayerCake.lean}, \nolinkurl{FractionalCDFLayerCake.lean},
\nolinkurl{FabiusFractionalIntegral.lean},
\nolinkurl{InverseLayerCake.lean}, \nolinkurl{ProbabilityLaplaceMoments.lean} \\
Fractional affine/order-raising calculus, causal Rvachev primitives, and shifts &
\nolinkurl{fractionalVolterra_affine},
\nolinkurl{fractionalVolterra_comp_affine},
\nolinkurl{fractionalVolterra_add_one_deriv},
\nolinkurl{Fabius.rvachevFractionalPrimitive},
\nolinkurl{Fabius.rvachevFractionalPrimitive_eq_intervalIntegral_min},
\nolinkurl{Fabius.rvachevFractionalPrimitive_nat_succ},
\nolinkurl{Fabius.rvachevFractionalPrimitive_add},
\nolinkurl{fractionalVolterra_add_one_extendedFabius_of_nonneg},
\nolinkurl{fractionalVolterra_add_one_fabiusReal},
\nolinkurl{fractionalVolterra_add_one_rvachevUp} &
\nolinkurl{FractionalVolterraCalculus.lean},
\nolinkurl{FabiusFractionalVolterra.lean} \\
Absolutely-continuous inverse-pair calculus &
\nolinkurl{intervalIntegral_deriv_mul_comp_add_comp_mul_deriv_of_lt_iff_lt},
\nolinkurl{intervalIntegral_deriv_mul_fabiusReal_add_fabiusInv_mul_deriv},
\nolinkurl{intervalIntegral_fabiusReal_add_fabiusInv_eq_one} &
\nolinkurl{InversePairIntegral.lean} \\
Finite-node and geometric formal-series filters &
\nolinkurl{finitePowerSeriesFilter},
\nolinkurl{coeff_finitePowerSeriesFilter},
\nolinkurl{finitePowerSeriesFilter_rescale},
\nolinkurl{map_finitePowerSeriesFilter},
\nolinkurl{coeff_finitePowerSeriesFilter_of_exact_of_le},
\nolinkurl{geometricSeriesFilter},
\nolinkurl{coeff_geometricSeriesFilter_of_exact},
\nolinkurl{geometricSeriesFilter_eq_residual_mk},
\nolinkurl{geometricLagrangeSeriesFilter_eq_residual_mk} &
\nolinkurl{FinitePowerSeriesFilter.lean},
\nolinkurl{GeometricPowerSeriesFilter.lean} \\
Unconditionally summable analytic-series filters &
\nolinkurl{finiteAnalyticSeriesFilter},
\nolinkurl{geometricAnalyticSeriesFilter},
\nolinkurl{finiteAnalyticSeriesFilter_eq_tsum},
\nolinkurl{finiteAnalyticSeriesFilter_eq_head_add_tail_of_exact},
\nolinkurl{geometricAnalyticSeriesFilter_eq_constant_add_gaussian_tsum},
\nolinkurl{geometricLagrangeAnalyticSeriesFilter_shifted} &
\nolinkurl{AnalyticSeriesFilter.lean} \\
Formal implicit roots and nonmonic Hensel lifting &
\nolinkurl{FormalImplicitRoot.exists_isRoot_sub_mem},
\nolinkurl{FormalImplicitRoot.eq_of_isRoot_of_sub_mem},
\nolinkurl{FormalImplicitRoot.existsUnique_isRoot_sub_mem},
\nolinkurl{PowerSeries.Implicit.existsUnique_isRoot_constantCoeff},
\nolinkurl{PowerSeries.Implicit.existsUnique_zeroConstant_root},
\nolinkurl{PowerSeries.Implicit.root},
\nolinkurl{PowerSeries.Implicit.constantCoeff_root},
\nolinkurl{PowerSeries.Implicit.eval_root},
\nolinkurl{PowerSeries.Implicit.eq_root} &
\nolinkurl{ImplicitPowerSeries.lean} \\
Totalized inverse & \nolinkurl{fabiusInv}, \nolinkurl{deriv_deriv_fabiusInv_eq_zero_iff}, \nolinkurl{tendsto_deriv_fabiusInv_atTop_at_zero_right}, \nolinkurl{fabiusInv_differentiableAt_iff}, \nolinkurl{fabiusInv_contDiffAt_infty_iff}, \nolinkurl{fabiusInv_analyticAt_iff} & \nolinkurl{FabiusInverse.lean}, \nolinkurl{InverseNotElementary.lean} \\
Effective dyadic inverse continuity (two definitions and twelve theorems).
With
\(D_{\mathrm{fac}}(r)=2^{\binom{r+1}{2}}(r+1)!\) and
\(D_{\Delta}(r)=2^{\binom{r+1}{2}}(r+1)^{r+1}\), the positive zeroth
recurrence term gives
\(\bigl(2^{\binom r2}(r+1)!(2^r-1)\bigr)^{-1}\le F(2^{-r})\) for \(r>0\).
For every \(r\), both reciprocal denominator bounds follow,
\(D_{\mathrm{fac}}(r)\le D_{\Delta}(r)\), both denominators are primitive
recursive, and each denominator has both strict and closed inverse moduli:
strict input error gives strict inverse error below \(2^{-r}\), while closed
input error gives inverse error at most \(2^{-r}\).  The factorial witness at \(r=n\) proves
\nolinkurl{EffectivelyUniformContinuous}; it does not supply
\nolinkurl{SequentiallyComputable}, a combined
\nolinkurl{IsComputableRealFunction} theorem, a logarithmic \(r(n)\), an exact
least/ceiling denominator, tolerant bisection, or input-bit asymptotics. &
\nolinkurl{inverseFabiusFactorialDenominator},
\nolinkurl{inverseFabiusDeltaDenominator};
\nolinkurl{inverseFabiusFactorialDenominator_eq},
\nolinkurl{inverseFabiusFactorialDenominator_primrec},
\nolinkurl{inverseFabiusDeltaDenominator_primrec},
\nolinkurl{fabiusReal_inverse_two_pow_one_term_lower_bound},
\nolinkurl{inv_inverseFabiusFactorialDenominator_le_fabiusReal},
\nolinkurl{inverseFabiusFactorialDenominator_le_deltaDenominator},
\nolinkurl{inv_inverseFabiusDeltaDenominator_le_fabiusReal},
\nolinkurl{abs_fabiusInv_sub_lt_inverse_two_pow_of_lt_factorialDenominator},
\nolinkurl{abs_fabiusInv_sub_le_inverse_two_pow_of_le_factorialDenominator},
\nolinkurl{abs_fabiusInv_sub_lt_inverse_two_pow_of_lt_deltaDenominator},
\nolinkurl{abs_fabiusInv_sub_le_inverse_two_pow_of_le_deltaDenominator},
\nolinkurl{fabiusInv_effectivelyUniformContinuous} &
\nolinkurl{FabiusInverseEffectiveContinuity.lean} \\
Logarithmic reciprocal inverse continuity (three definitions and fifteen
theorems).  The primitive-recursive selector is
\(r_{\log}(n)=\operatorname{Nat.log2}(n)+1\); for \(n>0\) it is the least
natural \(r\) with \(n<2^r\) and satisfies \(r_{\log}(n)\le n\).  The
factorial and Delta denominators take value \(1\) at zero, compose their
dyadic predecessors with \(r_{\log}\) at positive inputs, are primitive
recursive, and retain the factorial-to-Delta comparison.  For either
denominator, strict and closed input thresholds both give strict inverse error
below \(1/n\), and each denominator separately witnesses
\nolinkurl{EffectivelyUniformContinuous}.  The exact endpoint-mass ceiling
denominator, tolerant bisection, \nolinkurl{SequentiallyComputable}, the
combined \nolinkurl{IsComputableRealFunction} theorem, and input-bit
asymptotics remain outside this logarithmic-modulus leaf; the next two rows
record the downstream realizer closure. &
\nolinkurl{inverseFabiusLogarithmicOrder},
\nolinkurl{inverseFabiusLogarithmicFactorialDenominator},
\nolinkurl{inverseFabiusLogarithmicDeltaDenominator};
\nolinkurl{inverseFabiusLogarithmicOrder_eq_succ_log2},
\nolinkurl{inverseFabiusLogarithmicOrder_primrec},
\nolinkurl{inverseFabiusLogarithmicOrder_isLeast},
\nolinkurl{inverseFabiusLogarithmicOrder_le_self},
\nolinkurl{inverseFabiusLogarithmicFactorialDenominator_of_pos},
\nolinkurl{inverseFabiusLogarithmicDeltaDenominator_of_pos},
\nolinkurl{inverseFabiusLogarithmicFactorialDenominator_primrec},
\nolinkurl{inverseFabiusLogarithmicDeltaDenominator_primrec},
\nolinkurl{inverseFabiusLogarithmicFactorialDenominator_le_deltaDenominator},
\nolinkurl{abs_fabiusInv_sub_lt_inv_nat_of_lt_logarithmicFactorialDenominator},
\nolinkurl{abs_fabiusInv_sub_lt_inv_nat_of_le_logarithmicFactorialDenominator},
\nolinkurl{abs_fabiusInv_sub_lt_inv_nat_of_lt_logarithmicDeltaDenominator},
\nolinkurl{abs_fabiusInv_sub_lt_inv_nat_of_le_logarithmicDeltaDenominator},
\nolinkurl{fabiusInv_effectivelyUniformContinuous_logarithmic},
\nolinkurl{fabiusInv_effectivelyUniformContinuous_logarithmicDelta} &
\nolinkurl{FabiusInverseLogarithmicModulus.lean} \\
Sharp exact dyadic inverse modulus (two definitions and ten theorems).
For fixed \(r\), the natural ceiling of \(1/F(2^{-r})\) is positive, has
reciprocal at most the exact rational endpoint mass, and is the least positive
integer with that property.  It gives the strict output target \(2^{-r}\);
every smaller positive denominator fails at the endpoint pair
\(0,F(2^{-r})\), so it is least among strict integer moduli for this fixed
dyadic target.  The logarithmic wrapper takes value \(1\) at zero and uses the
fixed-target ceiling at \(r_{\log}(n)\) for \(n>0\), giving a strict \(1/n\)
output witness.  This does not prove leastness for the weaker \(1/n\) target,
assert a modulus at the convention-only zero input, or claim an input-bit
running-time bound.  Both denominator maps are primitive recursive. &
\nolinkurl{inverseFabiusExactDyadicDenominator},
\nolinkurl{inverseFabiusExactLogarithmicDenominator};
\nolinkurl{inverseFabiusExactDyadicDenominator_primrec},
\nolinkurl{inverseFabiusExactDyadicDenominator_pos},
\nolinkurl{inv_inverseFabiusExactDyadicDenominator_le_fabiusAtInverseTwoPow},
\nolinkurl{inverseFabiusExactDyadicDenominator_isLeast},
\nolinkurl{abs_fabiusInv_sub_lt_inverse_two_pow_of_lt_exactDyadicDenominator},
\nolinkurl{exists_fabiusInv_gap_of_lt_exactDyadicDenominator},
\nolinkurl{inverseFabiusExactDyadicDenominator_isLeast_strictModulus},
\nolinkurl{inverseFabiusExactLogarithmicDenominator_primrec},
\nolinkurl{inverseFabiusExactLogarithmicDenominator_of_pos},
\nolinkurl{abs_fabiusInv_sub_lt_inv_nat_of_lt_exactLogarithmicDenominator} &
\nolinkurl{FabiusInverseExactDyadicModulus.lean} \\
Fixed-depth effective monotone inversion (two definitions and six theorems).
At requested output precision \(p\), a natural-recursive controller performs
exactly \(p\) dyadic updates.  At depth \(s\), its bracket index \(j\)
represents \([j2^{-s},(j+1)2^{-s}]\), and it compares forward and target
dyadic names at precision \(p+d(p)+3+s\) with a four-numerator-unit margin.
A certified sign selects one half; an inconclusive comparison stores the
midpoint and doubles that stored numerator on every remaining step, preserving
its real value while advancing to denominator \(2^p\).  The generic theorem
assumes a positive computable denominator and its strict inverse-modulus
implication; the following gap module constructs such a certificate from
supplied positive rational dyadic-gap lower bounds. &
\nolinkurl{SequentiallyComputableOn},
\nolinkurl{unitClamp};
\nolinkurl{unitClamp_sequentiallyComputable},
\nolinkurl{tolerantDifference_error},
\nolinkurl{tolerantDifference_safe_updates},
\nolinkurl{tolerantDifference_inconclusive},
\nolinkurl{tolerantBisection_correct},
\nolinkurl{effectiveInversionOn_Icc} &
\nolinkurl{EffectiveMonotoneInverse.lean} \\
Effective inversion from computable positive rational dyadic gaps (exactly
four definitions and four theorems).  A computable sequence of positive natural
numerators and denominators gives a positive computable reciprocal denominator.
For a strict increasing inverse pair on \([0,1]\), the lower bound
\(\alpha(p)\le f(x+2^{-p})-f(x)\) for every
\(x\in[0,1-2^{-p}]\) yields the inverse modulus, sequential computability,
and effective uniform continuity.  The total result is precisely the clamped
extension \(x\mapsto g(\lean{unitClamp}(x))\), not a claim about the
unclamped values of \(g\) outside \([0,1]\). &
\nolinkurl{EffectivelyUniformContinuousOn},
\nolinkurl{ComputablePositiveRationalSequence},
\nolinkurl{ComputablePositiveRationalSequence.value},
\nolinkurl{ComputablePositiveRationalSequence.reciprocalDenominator},
\nolinkurl{ComputablePositiveRationalSequence.reciprocalDenominator_spec},
\nolinkurl{inverseModulus_of_positiveRationalGap},
\nolinkurl{effectiveInversionOn_Icc_of_computablePositiveRationalGap},
\nolinkurl{clampedEffectiveInversion_of_computablePositiveRationalGap} &
\nolinkurl{EffectiveGapInverse.lean} \\
Computability of the totalized Fabius inverse (zero definitions and one
theorem).  The unit-interval realizer uses the computable centered-spline
oracle and the Delta denominator; the computable unit clamp totalizes it, and
the logarithmic-Delta modulus supplies effective uniform continuity. &
\nolinkurl{fabiusInv_isComputableRealFunction} &
\nolinkurl{FabiusInverseComputable.lean} \\
Lower-Lambert coordinate & \nolinkurl{fabiusLambertPhase} & \nolinkurl{FabiusLambertSaddle.lean} \\
Exact real-branch pairing (zero definitions and seven theorems).  For
\(x\in(-e^{-1},0)\) and
\(\delta=W_0(x)-W_{-1}(x)>0\),
\(W_0=-\delta/(e^\delta-1)\),
\(W_{-1}=-\delta e^\delta/(e^\delta-1)
=-\delta/(1-e^{-\delta})\), and
\(x=W_0e^{W_0}\) with this displayed value of \(W_0\).  Division of the two
branch equations gives \(W_{-1}=W_0e^\delta\), and subtraction gives all
three formulas. &
\nolinkurl{principalLambertW_sub_lowerLambertW_pos},
\nolinkurl{lowerLambertW_eq_principalLambertW_mul_exp_gap},
\nolinkurl{principalLambertW_eq_neg_gap_div},
\nolinkurl{lowerLambertW_eq_neg_gap_mul_exp_div},
\nolinkurl{lowerLambertW_eq_neg_gap_div_one_sub_exp_neg},
\nolinkurl{eq_neg_gap_div_mul_exp},
\nolinkurl{principalLambertW_lowerLambertW_eq_of_exp_gap} &
\nolinkurl{LambertWBranchPairing.lean}; exhaustive public surface, on the
strict common branch domain only. \\
The branch gap as a global coordinate (four definitions and sixteen
theorems).  For every \(\delta>0\), the same formulas define
\(-1<A(\delta)<0\), \(B(\delta)<-1\), and
\(X(\delta)=A(\delta)e^{A(\delta)}\).  Since
\(B=A-\delta\) and \(Be^B=Ae^A\), branch uniqueness proves that
\(X\) and \(x\mapsto W_0(x)-W_{-1}(x)\) are inverse bijections between
\((0,\infty)\) and \((-e^{-1},0)\).  With \(t=e^\delta>1\), the two branch
values are \(-\log t/(t-1)\), \(-t\log t/(t-1)\), and the argument is
\(-\log t\,t^{-1/(t-1)}/(t-1)\). &
\nolinkurl{gapPrincipal}, \nolinkurl{gapLower}, \nolinkurl{gapArg},
\nolinkurl{branchGap};
\nolinkurl{gap_denominator_pos}, \nolinkurl{gapPrincipal_mem_Ioo},
\nolinkurl{gapLower_eq_mul_exp}, \nolinkurl{gapLower_eq_sub},
\nolinkurl{gapLower_lt_neg_one}, \nolinkurl{gapLower_mul_exp},
\nolinkurl{gapArg_mem_Ioo}, \nolinkurl{principalLambertW_gapArg},
\nolinkurl{lowerLambertW_gapArg}, \nolinkurl{branchGap_gapArg},
\nolinkurl{gapArg_branchGap}, \nolinkurl{branchGap_invOn},
\nolinkurl{branchGap_bijOn}, \nolinkurl{principalLambertW_gapArg_log},
\nolinkurl{lowerLambertW_gapArg_log}, \nolinkurl{gapArg_log} &
\nolinkurl{LambertWGapBijection.lean}; exhaustive public surface.  The
strict tangent inequalities for \(e^\delta\) put the two constructed values
in their unique branch ranges. \\
Symmetric branch laws (zero definitions and nine theorems).  On the same
strict domain,
\(W_{-1}/W_0=e^\delta\),
\(W_0+W_{-1}=-\delta\coth(\delta/2)<-2\), and
\(W_0W_{-1}=\delta^2/(4\sinh^2(\delta/2))\in(0,1)\).
The bounds are \(\sinh y>y\) and \(y\coth y>1\) at \(y=\delta/2\). &
\nolinkurl{lowerLambertW_div_principalLambertW_eq_exp_branchGap},
\nolinkurl{principalLambertW_add_lowerLambertW_eq_exp_branchGap},
\nolinkurl{principalLambertW_add_lowerLambertW_eq_cosh_div_sinh_branchGap},
\nolinkurl{principalLambertW_mul_lowerLambertW_eq_exp_branchGap},
\nolinkurl{principalLambertW_mul_lowerLambertW_eq_sinh_sq_branchGap},
\nolinkurl{principalLambertW_add_lowerLambertW_lt_neg_two},
\nolinkurl{principalLambertW_mul_lowerLambertW_pos},
\nolinkurl{principalLambertW_mul_lowerLambertW_lt_one},
\nolinkurl{principalLambertW_mul_lowerLambertW_mem_Ioo} &
\nolinkurl{LambertWBranchSymmetry.lean}; exhaustive public surface.  At the
branch point the gap is zero and the limiting sum/product equal \(-2,1\);
zero is the singular lower endpoint.  Neither endpoint is in these theorems,
and this symmetry leaf itself asserts no Bernoulli-gap series or higher
Puiseux/logarithmic expansion. \\
The Bernoulli branch-gap series and complex generating identity
(zero definitions and five theorems), crosswalking exactly
\nolinkurl{eq:pair-Bernoulli-general} and, under the explicit standard
removable-origin convention, \nolinkurl{eq:bernoulli-gen}.  The complex
Bernoulli exponential generating series is summable if and only if
\(\lVert z\rVert<2\pi\), so it diverges on the boundary and throughout the
exterior; precisely on that disk it has sum
\(\lean{complexExpm1Div}(z)^{-1}\), which is \(1\) at zero and rewrites to
\(z/(e^z-1)\) away from zero.  On
\(x\in(-e^{-1},0)\) with positive gap \(\delta<2\pi\), the two sums are
\(W_0=-\sum_{n\geq0}B_n\delta^n/n!\) and
\(W_{-1}=-\delta-\sum_{n\geq0}B_n\delta^n/n!\). &
\nolinkurl{summable_norm_bernoulli_mul_pow_div_factorial},
\nolinkurl{summable_bernoulli_mul_pow_div_factorial_iff},
\nolinkurl{hasSum_bernoulli_mul_pow_div_factorial},
\nolinkurl{hasSum_bernoulli_mul_pow_div_factorial_complex_iff},
\nolinkurl{principalLambertW_lowerLambertW_eq_bernoulliSeries} &
\nolinkurl{LambertWBranchGapBernoulli.lean}; exhaustive public surface.
No equality with Lean's literal totalized quotient at zero is claimed.  The
compiled exact-radius proof uses even Bernoulli--zeta coefficient bounds and
a positive even subsequence, not the Guide's nearest-nonzero-zero route.  Both
labelled identities are Exact under the stated convention; all higher or
convergent Puiseux claims remain open. \\
Raw real Lambert curvature (the first 13 declarations of the exact 36-name
curvature inventory).  For \(z>-e^{-1}\), the principal branch has the
nonsingular formula
\(W''=-e^{-2W}(W+2)/(W+1)^3\), hence \(W_0''(0)=-2\) and strict concavity
on the closed domain \([-e^{-1},\infty)\).  Lower-branch calculus and all
three sign equivalences assume exactly \(z\in(-e^{-1},0)\); the unique zero is
\(-2e^{-2}\), with strict convexity on
\([-e^{-1},-2e^{-2}]\) and strict concavity on
\([-2e^{-2},0)\).  The endpoint-inclusive shape results do not assert a
branch-point derivative, vertical-tangent limit, or Puiseux law; the two
later branch-point leaves in the next rows supply the vertical geometry and
leading square-root law. &
\nolinkurl{deriv_principalLambertW},
\nolinkurl{deriv_principalLambertW_hasDerivAt},
\nolinkurl{deriv_deriv_principalLambertW},
\nolinkurl{deriv_deriv_principalLambertW_zero},
\nolinkurl{deriv_deriv_principalLambertW_neg},
\nolinkurl{strictConcaveOn_principalLambertW},
\nolinkurl{deriv_lowerLambertW_hasDerivAt},
\nolinkurl{deriv_deriv_lowerLambertW},
\nolinkurl{deriv_deriv_lowerLambertW_pos_iff},
\nolinkurl{deriv_deriv_lowerLambertW_neg_iff},
\nolinkurl{deriv_deriv_lowerLambertW_eq_zero_iff},
\nolinkurl{strictConvexOn_lowerLambertW_left},
\nolinkurl{strictConcaveOn_lowerLambertW_right} &
\nolinkurl{LambertWCurvature.lean} \\
Lambert branch-point vertical geometry (eight theorems).  Along the
right-hand filter at \(b=-\exp(-1)\), the principal derivative and endpoint
secant slope tend to \(+\infty\), while their lower-branch counterparts tend
to \(-\infty\).  Consequently neither branch has a finite right derivative
on \(\operatorname{Ioi}(b)\), and neither totalized real branch is
differentiable at \(b\). &
\nolinkurl{tendsto_deriv_principalLambertW_branchPoint_atTop},
\nolinkurl{tendsto_deriv_lowerLambertW_branchPoint_atBot},
\nolinkurl{tendsto_principalLambertW_secantSlope_branchPoint_atTop},
\nolinkurl{tendsto_lowerLambertW_secantSlope_branchPoint_atBot},
\nolinkurl{principalLambertW_not_differentiableWithinAt_branchPoint},
\nolinkurl{lowerLambertW_not_differentiableWithinAt_branchPoint},
\nolinkurl{principalLambertW_not_differentiableAt_branchPoint},
\nolinkurl{lowerLambertW_not_differentiableAt_branchPoint} &
\nolinkurl{LambertWBranchPointGeometry.lean} \\
Lambert leading branch-point asymptotics (one noncomputable definition and
eight theorems).  Put \(b=-\exp(-1)\).  The intrinsic scale
\(\mathcal S(z)=\sqrt{2\exp(1)(z-b)}\) is positive to the right of \(b\) and
has the stated exact square.  Both squared displacement ratios tend
to \(2\exp(1)\), their squared forms are asymptotically equivalent, and the
signs select
\(W_0(z)+1\sim\mathcal S(z)\) and
\(W_{-1}(z)+1\sim-\mathcal S(z)\) along the same right-hand filter.  These are
only the leading square-root terms: an \(O(z-b)\) remainder, higher
coefficients, and a convergent or full Puiseux expansion remain open. &
\nolinkurl{lambertWBranchPointScale},
\nolinkurl{lambertWBranchPointScale_pos},
\nolinkurl{lambertWBranchPointScale_sq},
\nolinkurl{tendsto_principalLambertW_add_one_sq_div_branchPoint},
\nolinkurl{tendsto_lowerLambertW_add_one_sq_div_branchPoint},
\nolinkurl{principalLambertW_add_one_sq_isEquivalent_branchPoint},
\nolinkurl{lowerLambertW_add_one_sq_isEquivalent_branchPoint},
\nolinkurl{principalLambertW_add_one_isEquivalent_branchPoint},
\nolinkurl{lowerLambertW_add_one_isEquivalent_branchPoint} &
\nolinkurl{LambertWBranchPointAsymptotics.lean} \\
Generic scaled power--exponential curvature (four declarations).  Write \(P\)
for the power--exponential peak at parameters \(m,A,\beta\).
If \(m\in\NaturalNumbers\setminus\{0\}\), \(A>0\), \(\beta>0\), and \(0<x<P\), both
inverse branches satisfy
\(\lambda''=\lambda\bigl(m-(m-\beta\lambda)^2\bigr)/
\bigl(x^2(m-\beta\lambda)^3\bigr)\).  No generic square-root threshold,
sign, inflection, or global shape theorem is claimed. &
\nolinkurl{deriv_principalPowerExponentialPhase_hasDerivAt},
\nolinkurl{deriv_deriv_principalPowerExponentialPhase},
\nolinkurl{deriv_lowerPowerExponentialPhase_hasDerivAt},
\nolinkurl{deriv_deriv_lowerPowerExponentialPhase} &
\nolinkurl{PowerExponentialLambertCurvature.lean} \\
Fabius Lambert curvature (one definition and eighteen theorems, completing
the 36-name inventory).  Put \(L=\log2\),
\(x_*=2e^{-2}/L\), and \(P=e^{-1}/L\).  On \(0<x<P\), the lower phase has
\(\lambda(x_*)=2/L\),
\(\lambda''=L\lambda^2(2-L\lambda)/
\bigl(x^2(1-L\lambda)^3\bigr)\), sign equivalences on either side of \(x_*\),
and its unique zero there; its exact shape domains are
\(\operatorname{Ioc}(0,x_*)\) and \(\operatorname{Icc}(x_*,P)\).
The principal affine-\(W_0\) identity is global; continuity and strict
convexity hold on \(\operatorname{Iic}(P)\), while first and second calculus
and positivity assume \(x<P\), with \(\lambda_0''(0)=2L\).  This
specialization adds no branch-point derivative, vertical-tangent limit, or
Puiseux expansion. &
\nolinkurl{fabiusLambertInflectionInput},
\nolinkurl{fabiusLambertInflectionInput_mem_Ioo},
\nolinkurl{fabiusLambertPhase_inflectionInput},
\nolinkurl{deriv_deriv_fabiusLambertPhase},
\nolinkurl{deriv_deriv_fabiusLambertPhase_pos_iff},
\nolinkurl{deriv_deriv_fabiusLambertPhase_neg_iff},
\nolinkurl{deriv_deriv_fabiusLambertPhase_eq_zero_iff},
\nolinkurl{deriv_deriv_fabiusLambertPhase_inflectionInput},
\nolinkurl{strictConvexOn_fabiusLambertPhase_left},
\nolinkurl{strictConcaveOn_fabiusLambertPhase_right},
\nolinkurl{fabiusPrincipalLambertPhase_eq_principalLambertW},
\nolinkurl{fabiusPrincipalLambertPhase_continuousOn_Iic},
\nolinkurl{fabiusPrincipalLambertPhase_hasDerivAt},
\nolinkurl{deriv_fabiusPrincipalLambertPhase},
\nolinkurl{deriv_fabiusPrincipalLambertPhase_hasDerivAt},
\nolinkurl{deriv_deriv_fabiusPrincipalLambertPhase},
\nolinkurl{deriv_deriv_fabiusPrincipalLambertPhase_pos},
\nolinkurl{deriv_deriv_fabiusPrincipalLambertPhase_zero},
\nolinkurl{strictConvexOn_fabiusPrincipalLambertPhase} &
\nolinkurl{PowerExponentialLambertFabiusCurvature.lean} \\
All-orders saddle expansion & saddle mass/log coefficient declarations & \nolinkurl{FabiusFullAsymptoticExpansion.lean} and related modules \\
\bottomrule
\end{longtable}

\begin{boxedremark}
\textbf{Explicit nonclaims.}
The crosswalk identifies formal inputs and nearby APIs; it is not a claim that every formula
in the corresponding part is formalized.  In particular, the sharp quantitative
Thue--Morse rates, repeated-integration norm identities, inverse endpoint asymptotic,
general ordered-composition coefficient formulas, Fourier-amplitude estimates, and any
claim marked derived, numerical, conjectural, or open remain in this frontier until an exact
Lean declaration is audited.  The named weights \((1-q)q^n\), their law, measure recurrence,
reflection, support equal to the series range for \(|q|<1\), and exact support \([0,1]\) for
\(0\le q<1\), together with absolute continuity and null singletons throughout \(|q|<1\),
are now exact formal inputs.  So are the named CDF, its continuity and reflection for
\(|q|<1\), its tails for \(0\le q<1\), and, for \(0<q<1\), its refinement integrals,
explicit compactly supported density, measure-level Radon--Nikodym identity, and real-space
\(C^\infty\) regularity.  The generic affine-law operator now has formal probability
preservation, product-map and characteristic-function factorizations, finite recurrence
iteration, and probability-fixed-point uniqueness for \(|q|<1\); the named geometric law is
therefore characterized by its affine equation, including at \(q=0\) and negative \(q\).
The scaled digit now has a closed characteristic function, and for every \(|q|<1\) the
phase-collapsed finite residual formula and pointwise convergence of its phase-bearing sinc
prefixes are formal, with exact dyadic wrappers.  The defining products for
\texttt{geometricSincProduct} now converge locally uniformly on \(\ComplexNumbers\), have named
\texttt{Multipliable} and \texttt{HasProd} theorems, and define an entire function.  For
\(z_q(t)=(1-q)t/(2\pi)\), the exact named bridges give
\(\phi_q(t)=e^{\ImaginaryUnit t/2}S_q(z_q(t))
=e^{\ImaginaryUnit t/2}G_q(z_q(t))G_q(-z_q(t))\) throughout \(|q|<1\), including zero
and negative ratios.  The selected total density formula is formally zero at \(q=0\),
nonpositive for negative ratios, and gives the zero \texttt{withDensity} measure throughout
\(q\le0\), hence not the probability law on the nonpositive contraction range.  The full
phase-bearing prefixes now converge locally uniformly on the real frequency line and
uniformly on every compact real-frequency set throughout \(|q|<1\), including zero and
negative ratios.  What remains here is Lean formalization of the paper-level corrected
signed-density construction for \(q\le0\), a corresponding complex-frequency locally
uniform theorem for the full phase-bearing prefix sequence, rapid decay and Fourier
inversion, named centered/MGF packaging, a closed
higher-derivative hierarchy and endpoint-flatness wrappers, shape theorems, further
transforms, and explicit Bernoulli-cumulant/Bell-moment formulas and asymptotics.  Finite
characteristic-function, MGF, and CGF tail factorizations are already part of the geometric
tail dictionary.
The generic and geometric Cauchy rows above additionally make the transform DDE and
adjacent-power hierarchy direct named theorems, including negative nonzero ratios.  Here
``closed higher-derivative hierarchy'' means the report's real-space density derivatives
and endpoint wrappers; it does not deny the proved Cauchy-power recurrence.
\end{boxedremark}


% BEGIN SYNTHESIZED TOPIC: probability, discrete, and generic asymptotic engines
% Historically recorded SHA-256 (not reproduced by a reachable source blob):
% 21222ae5a8c64cf556dac562fd66943ae0b6ed881408e23e37edfa9113bdecbd
% Authoritative archived standalone Git blob: c450f8c7ab2764b02759c5295a7ae1f909c43fee
\clearpage
\part{Probability, Discrete, and Generic Asymptotic Engines}
\label{part:integration}
\begin{boxedremark}
\textbf{Source provenance and status.} Former path: \nolinkurl{Fabius_Integration_Research_Frontiers/Fabius_Integration_Research_Frontiers.tex}.  This broad mixed-status input supplied probability bridges, discrete engines, generic saddle machinery, and an earlier inverse presentation.  The inverse and Fabius-specific saddle material now live only in \cref{part:inverse}; this part retains the reusable engines and its source-specific proofs.\par\smallskip Historically recorded source SHA-256: \nolinkurl{21222ae5a8c64cf556dac562fd66943ae0b6ed881408e23e37edfa9113bdecbd}.  The forensic audit did not reproduce that value from a reachable whole-file blob; the authoritative archived standalone source is Git blob \nolinkurl{c450f8c7ab2764b02759c5295a7ae1f909c43fee}.
\end{boxedremark}
\section*{Topic abstract}
The long article \emph{The Fabius Function and Rvachev's Up-Function} gives a broad
human-readable account of the formal development in \texttt{ProveIt}.  Its breadth is also
its unavoidable weakness: several new branches of the Lean library postdate the article,
and several generic proof engines are named only in the module guide because they have no
counterpart in the exposition.  This companion supplies those missing parts.

The opening audit records a formalized redundancy in the original compact-support
characterization: positivity of its dilation parameter follows from the remaining
hypotheses.  The analytic-locus chapter then develops the generic elementary-function
obstruction and applies it to the forward Fabius family and its total inverse.  The
probability chapters supply the measure-theoretic bridges between the product law, its
survival function, tilted Laplace moments, finite atomic laws, and half-weighted step
approximants, together with the post-snapshot continuous-linear naturality, conditional
operator/scalar affine-recursion, and real absolute-continuity APIs for arbitrary weighted
uniform laws and the named geometric-\(q\) specialization.  Two reusable
discrete engines follow: weighted path sums for triangular
recurrences and a finite-row Toeplitz convergence theorem.  Finally, the part exposes the
generic machinery used by the all-orders saddle proof: parameter-dependent Poincar\'e
expansions, recursive exponential and logarithmic coefficients, exact finite polynomial
remainders, Gaussian contraction, and uniform polynomial tails.  The inverse theory,
closed Fabius jets, explicit corrections, and contour assembly are cross-referenced to their
single canonical treatment in \cref{part:inverse}.

The purpose is complementarity rather than repetition.  Definitions and headline theorems
already proved in detail in the main article are recalled only when they are needed as input.
Every substantial assertion below is either a direct human-readable rendering of a current
Lean declaration, a proof-level explanation of a generic module that the main article
explicitly omits, or a clearly identified elementary corollary of those declarations.


\begin{boxedremark}
\textbf{Formalization boundary.}
This is a research-frontier document.  It preserves the source draft used during integration,
including exact Lean-backed results, proof sketches, ordinary-analysis corollaries, stronger
generalizations, and formulas that still lack exact declarations.  Its presence under
\path{semi-formalized-research-frontiers} is intentional: no mathematical assertion in this
notebook should be attributed to the formal library merely because a related definition or
weaker theorem exists.  The primary exposition contains the audited, declaration-backed
subset.  Among the obligations retained here are generic moment-matrix consequences,
arbitrary-family Toeplitz wording, generic Gaussian coefficient bounds, and the geometric-
\(q\) Fourier, centered-family, higher-derivative, endpoint-flatness, shape,
transform, and moment/cumulant layers beyond the formalized CDF and positive-ratio
smooth-density API.  The same API now proves that the selected total density formula is
zero at \(q=0\), nonpositive for negative \(q\), and has zero clamped
\texttt{withDensity} measure for \(q\le0\), which cannot represent \(\mu_q\) on the
nonpositive contraction range.  Only formalization of the paper-level sign-corrected
density remains open.  The
probability-law fixed point, by contrast, is now
unique under the exact contractive hypotheses recorded below, the phase-bearing sinc
prefix has a formal pointwise limit, and the oriented Cauchy-transform DDE and
adjacent-power hierarchy are formal for nonzero contractions.  The latter should not be
confused with the remaining Fourier and real-space density-derivative layers.  The inverse
endpoint and Fabius-specific coefficient obligations are tracked in \cref{part:inverse}.
\end{boxedremark}

\section{Scope, conventions, and the audit boundary}
\label{integration:sec:scope}

Let $F:\RealNumbers\to[0,1]$ denote the bounded Fabius function, totalized in the usual way by
$F(x)=0$ for $x\le0$ and $F(x)=1$ for $x\ge1$.  Let
\begin{equation}
 \RvachevUp(x)=F(1-|x|)
 \label{integration:eq:up-fold-supp}
\end{equation}
be Rvachev's compactly supported up-function, and let $\FabiusGlobal$ denote the signed
global extension used in the Lean development.  On $[0,1]$ the basic identities are
\begin{equation}
 F(1-x)=1-F(x),
 \qquad
 F'(x)=2\RvachevUp(2x-1),
 \qquad
 F'(x)=2F(2x)\quad(0<x<\tfrac12).
 \label{integration:eq:basic-inputs}
\end{equation}
These are inputs here, not topics to be reproved.

The current formal library also packages their exact curvature consequence:
\begin{equation}
 F''(x)=8\bigl(\RvachevUp(4x-1)-\RvachevUp(4x-3)\bigr),
 \label{integration:eq:forward-second-derivative}
\end{equation}
with \(F''>0\) exactly on \((0,1/2)\), \(F''<0\) exactly on \((1/2,1)\), and
\(F''=0\) exactly off \((0,1)\) or at \(1/2\).  These are
\path{deriv_deriv_fabiusReal}, \path{deriv_deriv_fabiusReal_pos_iff},
\path{deriv_deriv_fabiusReal_neg_iff}, and
\path{deriv_deriv_fabiusReal_eq_zero_iff}.

The repository snapshot underlying this supplement is the following commit on
the \texttt{codex/fabius-both-papers} branch, dated 25 August 2026.
\begin{center}
 \texttt{eaea1b9afe2b7ff0b7dd880bd1710e303dec6d80}.
\end{center}
An earlier version of the main article left four clusters to its module guide:
Gaussian contraction and Gaussian tails, the formal saddle exponential and
logarithm algebra, the abstract Toeplitz row lemma, and the generic weighted-path solver.
Those four clusters receive complete mathematical chapters below.  The non-elementarity
development is also expanded here, while the inverse API is handed to \cref{part:inverse}
instead of being duplicated.  The current snapshot also proves that positivity of the
dilation parameter in the original compact-support characterization is redundant.  Finally,
several bridges are present only as one- or two-sentence transitions; they are expanded here
into proof-level arguments.

\subsection{Three levels of coverage}

To keep the audit precise, each topic belongs to one of the following levels.
\begin{longtable}{@{}p{0.19\textwidth}p{0.70\textwidth}@{}}
\toprule
Status & Meaning in this companion\\
\midrule
\endhead
\textbf{Expanded formal input} & A current Lean module whose exact declarations are inputs to
a longer proof-oriented or frontier discussion here: principally
\path{ElementaryFunction.lean}, \path{NotElementary.lean},
\path{InverseNotElementary.lean}, \path{ProbabilityLaplaceMoments.lean}, and
\path{FractionalVolterra.lean}, together with
\path{FractionalVolterraSemigroup.lean}, \path{FractionalVolterraCalculus.lean},
\path{FabiusFractionalVolterra.lean}, \path{CDFLayerCake.lean},
\path{FractionalCDFLayerCake.lean}, \path{FabiusFractionalIntegral.lean},
\path{InversePairIntegral.lean}, \path{FinitePowerSeriesFilter.lean}, and
\path{GeometricPowerSeriesFilter.lean}, together with
\path{AnalyticSeriesFilter.lean}, \path{ThueMorseComplexProductBridge.lean},
\path{MeasureCauchyTransform.lean}, \path{MeasureCauchyMomentLaurent.lean},
\path{GeometricUniformCauchy.lean},
\path{CauchyTransform.lean}, \path{CauchyCDF.lean},
\path{CauchySurvival.lean}, \path{CauchyHigherPowers.lean}, and
\path{CauchyRenormalization.lean}, together with the real logarithmic and
generalized-order, complex Laurent/Herglotz, integrated Stieltjes--Perron, and initial
Jacobi-data modules, and
\path{ContinuousCDF.lean}, \path{GeometricUniformCDF.lean}, and
\path{ImplicitPowerSeries.lean}.\\
\textbf{Missing engine} & A generic theorem family that the main article deliberately leaves in
the module guide: path sums, Toeplitz rows, formal coefficient recurrences, finite
remainders, and Gaussian polynomial estimates.\\
\textbf{Compressed bridge} & A statement whose endpoints occur in the main article but whose
intermediate measure theory, topology, or asymptotic algebra is suppressed.\\
\bottomrule
\end{longtable}

\subsection{Notation that must not collide}

The exact half moments are denoted by $d_n$, following the main article.  The bounded saddle
exponent jets are instead denoted by
\begin{equation}
 Z_n(t),
 \label{integration:eq:Z-notation}
\end{equation}
although their Lean declaration is
\lean{negativeLaplaceBoundedExponentJet}.  This avoids using the same letter for two wholly
different sequences.  We write $L=\log 2$, $\Psi$ for the smooth one-periodic correction in
the negative-Laplace exponent, and $\lambda(x)$ for the exact lower-Lambert phase at small
positive $x$.

\begin{boxedremark}
\textbf{Formal boundary convention.}  When a statement is a theorem in the current Lean
library, the corresponding declaration is named in the margin of the discussion or in the
concordance in \cref{integration:app:concordance}.  When a statement is an immediate classical calculus
corollary but is not presently packaged as a declaration, it is called a \emph{corollary in
ordinary analysis}.  The forward and inverse second-derivative profiles, midpoint inverse
jet, strict closed-half shape, and endpoint interior-derivative limits are now packaged
formally; the distinction remains relevant for general higher inverse
derivatives and for the transfer of the highest-jet coefficient from mass coefficients to
logarithmic coefficients.
\end{boxedremark}

\subsection{The positivity hypothesis in the original characterization is redundant}
\label{integration:sec:scale-positive}

The compact-support characterization used in the first Arias de Reyna paper
\cite{integration:arias05442} assumes a smooth function $\phi:\RealNumbers\to\RealNumbers$, a real dilation
parameter $k$, and
\begin{equation}
 \operatorname{tsupp}\phi=[-1,1],\qquad
 \phi(x)>0\quad(-1<x<1),\qquad
 \phi(0)=1,
 \label{integration:eq:original-characterization-data}
\end{equation}
together with the derivative law
\begin{equation}
 \phi'(x)=k\bigl(\phi(2x+1)-\phi(2x-1)\bigr).
 \label{integration:eq:original-characterization-derivative}
\end{equation}
The paper states $k>0$ as a hypothesis.  In the current formal development that field is retained
for source fidelity, but a smart constructor proves that its positivity follows from the other
assumptions.

\begin{theorem}[Redundancy of scale positivity]
\label{integration:thm:scale-positive-redundant}
Suppose $\phi\in C^\infty(\RealNumbers)$ satisfies
\eqref{integration:eq:original-characterization-data} and
\eqref{integration:eq:original-characterization-derivative} for some $k\in\RealNumbers$.  Then $k>0$.
Consequently these data determine an instance of the full original-characterization predicate
without a separate positivity proof.
\end{theorem}

\begin{proof}
The support identity and continuity imply $\phi(-1)=0$: the function vanishes on
$(-\infty,-1)$, and $-1$ is a limit point of that interval.  The mean-value theorem on
$[-1,0]$ therefore gives a point $c\in(-1,0)$ such that
\begin{equation}
 \phi'(c)=\frac{\phi(0)-\phi(-1)}{0-(-1)}=1.
 \label{integration:eq:original-characterization-mvt}
\end{equation}
Now $2c-1<-1$, so the support condition gives $\phi(2c-1)=0$, whereas
$-1<2c+1<1$, so interior positivity gives $\phi(2c+1)>0$.  Evaluating
\eqref{integration:eq:original-characterization-derivative} at $c$ yields
\begin{equation}
 1=k\phi(2c+1).
 \label{integration:eq:original-characterization-positive-product}
\end{equation}
The second factor is positive, hence $k>0$.
\end{proof}

This is \lean{IsOriginalFabius.mk\_of\_derivative\_law}.  The proof is small, but it clarifies
the logical content of the original theorem: the sign of the scale is forced before the later
Fourier argument identifies its exact value $k=2$.

\section{The elementary-function obstruction and inverse handoff}

\subsection{Canonical inverse treatment}
\label{integration:sec:inverse}
\label{integration:sec:inverse-steepness}
The totalized inverse---written \(I\) in this source's audit tables and \(\FabiusClampedQuantile\) in the
canonical chapter---is treated once, in the combined inverse-and-saddle synthesis
of \cref{part:inverse}.  That treatment includes the order isomorphism on \([0,1]\),
the implementation through \nolinkurl{fabiusOrderIso}, \nolinkurl{fabiusInv}, and
\nolinkurl{Set.IccExtend}, the four-way comparison calculus, reflection and exact dyadic
inversion, transported flatness bounds, endpoint steepness, reciprocal first and second
derivatives, strict closed-half curvature, and the exact differentiability,
\(C^\infty\), and analytic loci.  Keeping those results in a single later chapter avoids
presenting the same inverse package twice; the present part retains only the generic
analytic-locus obstruction that is reused for several function classes.

\subsection{Elementary functions through their analytic locus}
\label{integration:sec:elementary}

The non-elementarity proof does not attempt differential algebra or Liouville theory.  It
uses a topological-analytic invariant that is both stronger and easier to formalize: every
one-variable elementary function is analytic on a dense open set, whereas the Fabius family
is nonanalytic at every point of the relevant interval.

\subsubsection{The exact formal class}

\begin{definition}[The class $\mathsf{Elem}$]
\label{integration:def:elementary-class}
Let $\mathsf{Elem}$ be the smallest class of total functions $\RealNumbers\to\RealNumbers$ containing all
constant functions and the identity and closed under
\begin{equation}
 +,\quad\cdot,\quad-,\quad(\cdot)^{-1},\quad
 f\mapsto f^r\ (r\in\RealNumbers),\quad
 \exp,\ \log,\ \sin,\ \cos,\ \arcsin,\ \arctan,
 \label{integration:eq:elementary-generators}
\end{equation}
where every unary operation is applied pointwise.
\end{definition}

This is the inductive predicate \lean{IsElementary}.  Composition is deliberately not a
constructor.

\begin{proposition}[Closure derived from the generators]
\label{integration:prop:elementary-derived}
The class $\mathsf{Elem}$ is closed under composition.  It is consequently closed under
subtraction, division, natural powers, polynomial and rational functions, square roots,
absolute values, tangent, hyperbolic functions, $\arccos$, $\operatorname{arsinh}$, and
signed odd roots.
\end{proposition}

\begin{proof}
Induct on the construction of the outer function $g$.  Precomposing a constant or the
identity with an elementary $f$ is elementary.  Every remaining constructor commutes with
precomposition: for example $(g_1+g_2)\circ f=(g_1\circ f)+(g_2\circ f)$, and similarly for
the unary generators.  This proves closure under composition.  The remaining operations
are identities in the generated algebra.  For example
\[
 |f|=\sqrt{f^2},
 \qquad
 \tan f=\frac{\sin f}{\cos f},
 \qquad
 \sinh f=\frac{\EulerE^f-\EulerE^{-f}}2.
\]
For an odd positive integer $n$, the signed root may be written
\[
 \frac{f}{|f|}|f|^{1/n},
\]
with value zero at $f=0$ under the total reciprocal convention.
\end{proof}

\begin{remark}[Totalization makes the theorem stronger]
\label{integration:rem:elementary-totalization}
Mathlib treats reciprocal, logarithm, real power, and inverse trigonometric functions as
total maps on $\RealNumbers$; outside their classical domains they receive specified extension values.
Thus \cref{integration:def:elementary-class} is at least as large as the class represented by classical
expressions on open domains.  Proving that $F$ does not belong to this larger class is a
stronger statement, not a domain-convention loophole.
\end{remark}

\subsubsection{Analytic loci}

For a function $f:\RealNumbers\to\RealNumbers$, define
\begin{equation}
 \anloc(f)=\{x\in\RealNumbers:f\text{ is real analytic at }x\}.
 \label{integration:eq:analytic-locus}
\end{equation}
Analyticity is local, hence $\anloc(f)$ is open for every $f$.  The substantive theorem is
density for elementary functions.

The only difficult induction step is composition with a unary generator that has finitely
many exceptional values.  The following lemma isolates it.

\begin{lemma}[Finite-bad-value composition principle]
\label{integration:lem:finite-bad-composition}
Let $f,g:\RealNumbers\to\RealNumbers$.  Suppose that $\anloc(f)$ is dense and that there is a finite set
$B\subset\RealNumbers$ such that $g$ is analytic at every point of $\RealNumbers\setminus B$.  Then
$\anloc(g\circ f)$ is dense.
\end{lemma}

\begin{proof}
Let $U$ be a nonempty open interval.  Because $\anloc(f)$ is dense and open, choose
$x_0\in U\cap\anloc(f)$.  We prove that every punctured neighborhood of $x_0$ contains a
point where $g\circ f$ is analytic.

Fix $c\in B$.  The analytic function $f-c$ near $x_0$ has two possibilities.
\begin{enumerate}
 \item It vanishes identically on a neighborhood of $x_0$.  Then $f$ is locally constant
 equal to $c$, so $g\circ f$ is locally constant and therefore analytic at $x_0$, regardless
 of whether $g$ is analytic at $c$.
 \item It is not locally zero.  The isolated-zero theorem gives a punctured neighborhood of
 $x_0$ on which $f(x)\ne c$.
\end{enumerate}
If the first alternative occurs for any $c\in B$, we are done.  Otherwise intersect the
finitely many punctured neighborhoods.  The intersection still contains points arbitrarily
close to $x_0$; intersect once more with the open set $U\cap\anloc(f)$.  At any resulting
point $x$, the value $f(x)$ avoids $B$, so $g$ is analytic at $f(x)$ and $f$ is analytic at
$x$.  The analytic chain rule makes $g\circ f$ analytic at $x$.  Thus every nonempty open
$U$ meets $\anloc(g\circ f)$.
\end{proof}

\begin{theorem}[Dense open analytic locus]
\label{integration:thm:elementary-dense-analytic}
If $f\in\mathsf{Elem}$, then $\anloc(f)$ is dense and open.  Equivalently, the nonanalytic
locus
\begin{equation}
 \RealNumbers\setminus\anloc(f)
 \label{integration:eq:elementary-exceptional-set}
\end{equation}
is closed and nowhere dense.
\end{theorem}

\begin{proof}
Openness is true for every analytic locus.  Density follows by induction on the elementary
construction.

Constants and the identity are analytic everywhere.  For sums and products, the
intersection of the two dense open analytic loci is dense and each point in the intersection
is analytic for the result.  Negation is immediate.  For a unary generator, apply
\cref{integration:lem:finite-bad-composition} with the following exceptional sets:
\begin{center}
\begin{tabular}{@{}lll@{}}
\toprule
Generator & Exceptional values & Analytic elsewhere\\
\midrule
$t\mapsto t^{-1}$ & $\{0\}$ & yes\\
$t\mapsto t^r$ & $\{0\}$ & yes for every fixed $r\in\RealNumbers$\\
$\log t$ & $\{0\}$ & yes under the total real-log convention\\
$\arcsin t$ & $\{-1,1\}$ & yes\\
$\exp t,\sin t,\cos t,\arctan t$ & $\varnothing$ & everywhere\\
\bottomrule
\end{tabular}
\end{center}
The complement of a dense open set is closed with empty interior, hence nowhere dense.
\end{proof}

\subsubsection{The reusable obstruction}

\begin{theorem}[No agreement on an interior]
\label{integration:thm:elementary-obstruction}
Let $h:\RealNumbers\to\RealNumbers$ be nonanalytic at every point of $\interior(S)$, where
$S\subset\RealNumbers$ has nonempty interior.  No elementary function $g$ can agree with $h$ on all of
$S$.
\end{theorem}

\begin{proof}
If $g$ were elementary, \cref{integration:thm:elementary-dense-analytic} would give a point
$x\in\interior(S)$ at which $g$ is analytic.  Equality on $S$ implies equality on a whole
neighborhood of $x$ contained in $S$, so analyticity transfers from $g$ to $h$ at $x$, a
contradiction.
\end{proof}

The theorem is stronger than saying that $h$ itself is not elementary: it rules out any
piecewise repair outside $S$ and any elementary surrogate on a substantial subset.

\subsection{Strong non-elementarity of the Fabius family}
\label{integration:sec:not-elementary}

The main article proves nowhere analyticity of $F$ on $[0,1]$, of $\RvachevUp$ on $[-1,1]$, and of
the signed extension on its active interval.  Combining those results with
\cref{integration:thm:elementary-obstruction} gives the new non-elementarity branch almost for free.

\begin{theorem}[Strong local non-elementarity of $F$]
\label{integration:thm:fabius-strong-non-elementary}
Let $U\subset[0,1]$ have nonempty interior.  There is no $g\in\mathsf{Elem}$ such that
\begin{equation}
 g(x)=F(x)\qquad(x\in U).
 \label{integration:eq:elementary-F-agreement}
\end{equation}
In particular, no elementary function agrees with $F$ on $(0,1)$, and $F$ is not an
elementary function on $\RealNumbers$.
\end{theorem}

\begin{proof}
The function $F$ is nonanalytic at every point of $[0,1]$, hence at every point of
$\interior(U)$.  Apply \cref{integration:thm:elementary-obstruction}.
\end{proof}

\begin{theorem}[Strong local non-elementarity of $\RvachevUp$]
\label{integration:thm:up-strong-non-elementary}
If $U\subset[-1,1]$ has nonempty interior, no elementary $g$ agrees with $\RvachevUp$ on $U$.
Consequently $\RvachevUp$ is not elementary.
\end{theorem}

\begin{theorem}[Strong local non-elementarity of the signed extension]
\label{integration:thm:global-strong-non-elementary}
If $U\subset[0,2)$ has nonempty interior, no elementary $g$ agrees with the signed global
extension $\FabiusGlobal$ on $U$.  Consequently $\FabiusGlobal$ is not elementary.
\end{theorem}

\begin{proof}[Proof of the last two theorems]
The corresponding nowhere-analyticity theorems supply the hypothesis of
\cref{integration:thm:elementary-obstruction}.  The global conclusions follow by taking a nonempty open
subinterval of the active support and observing that an elementary representation on all of
$\RealNumbers$ would in particular agree there.
\end{proof}

The same analytic-density obstruction applies to the totalized inverse.  To avoid a second
statement and proof of the inverse theorem here, \cref{inverse:sec:inverse-calculus} records
its exact analytic locus, the resulting local non-elementarity declarations, and the
additional branch-domain hypotheses needed for the enlarged elementary-or-inverse class.

\begin{remark}[Why compact support alone is not the argument]
A compactly supported nonzero function cannot be globally real analytic, but elementary
functions need not be analytic everywhere: $|x|$, $x^{-1}$, and totalized logarithms already
have exceptional points.  What defeats elementarity is not compact support by itself; it is
nonanalyticity on an entire interval, while elementary exceptional sets are nowhere dense.
\end{remark}

\begin{remark}[Why this is sharper than a differential-algebra slogan]
The conclusion does not depend on choosing a particular syntax for ``closed form'' beyond
the explicit inductive class.  More importantly, it is local: changing the function outside
a small open interval cannot make it elementary on that interval.  Thus the theorem rules
out not only a global elementary formula for $F$, but also an elementary formula valid on
any nontrivial open patch of its natural domain.
\end{remark}

\section{Probability laws and the bridges the headline formulas conceal}

\subsection{The product probability law in measure-theoretic form}
\label{integration:sec:product-law}

The random-series representation in the main article is concise because it speaks in the
usual probabilistic shorthand.  The formal construction must establish continuity,
measurability, product independence, pushforward identities, and support before it may use
the same shorthand.  This section records that chain explicitly.

\subsubsection{Sample space and uniform convergence}

Let
\begin{equation}
 \Omega=[0,1]^{\NaturalNumbers}
 \label{integration:eq:sample-space}
\end{equation}
with its product topology and product Borel probability measure
\begin{equation}
 \mathbf P=\bigotimes_{n\ge0}\operatorname{Unif}[0,1].
 \label{integration:eq:uniform-product}
\end{equation}
Write $\omega_n$ for the $n$th coordinate and define
\begin{equation}
 X(\omega)=\sum_{n=0}^{\infty}\frac{\omega_n}{2^{n+1}}.
 \label{integration:eq:weighted-coordinate-sum}
\end{equation}
For every $\omega$,
\begin{equation}
 0\le \frac{\omega_n}{2^{n+1}}\le2^{-n-1},
 \qquad
 \sum_{n\ge0}2^{-n-1}=1.
 \label{integration:eq:coordinate-majorant}
\end{equation}
The Weierstrass test therefore gives uniform convergence on $\Omega$.  Every partial sum is
continuous, so $X$ is continuous and hence measurable.  Termwise comparison in
\eqref{integration:eq:coordinate-majorant} gives the pointwise support statement
\begin{equation}
 0\le X(\omega)\le1.
 \label{integration:eq:X-support}
\end{equation}

Let
\begin{equation}
 \mu=\mathbf P\circ X^{-1}
 \label{integration:eq:weighted-sum-law}
\end{equation}
be the pushforward law.  It is a probability measure and
\begin{equation}
 \mu([0,1])=1,
 \qquad
 \mu(\RealNumbers\setminus[0,1])=0,
 \qquad
 \SupportOperator(\mu)=[0,1].
 \label{integration:eq:law-support}
\end{equation}
The Lean declarations are \lean{weightedCoordinateSum},
\lean{continuous\_weightedCoordinateSum}, \lean{weightedSumDistribution},
\lean{weightedSumDistribution_Icc}, \lean{weightedSumDistribution_compl_Icc},
\lean{ae_weightedSumDistribution_mem_Icc},
\lean{weightedSumDistribution_restrict_Icc}, and the exact topological-support theorem
\lean{weightedSumDistribution_support_eq_Icc}.  The almost-everywhere and restriction
declarations are owned upstream by \path{ProbabilityRepresentation.lean}.

\subsubsection{Deleting the head coordinate}

Define the left shift
\begin{equation}
 (T\omega)_n=\omega_{n+1}.
 \label{integration:eq:tail-map}
\end{equation}
The product measure is invariant under this deletion in the sense that $T_\#\mathbf P=\mathbf
P$, and the head coordinate $\omega_0$ is independent of the entire tail $T\omega$.  Reindexing
the series gives the exact pointwise split
\begin{equation}
 X(\omega)=\frac{\omega_0+X(T\omega)}2.
 \label{integration:eq:X-head-tail-split}
\end{equation}
Consequently the pair $(\omega_0,X(T\omega))$ has law
\begin{equation}
 \operatorname{Unif}[0,1]\otimes\mu.
 \label{integration:eq:head-tail-law}
\end{equation}
Pushing \eqref{integration:eq:X-head-tail-split} through this joint law gives the self-similarity identity
\begin{equation}
 \mu=\left(\operatorname{Unif}[0,1]\otimes\mu\right)
       \circ\left[(u,x)\mapsto\frac{u+x}{2}\right]^{-1}.
 \label{integration:eq:law-self-similarity}
\end{equation}

\subsubsection{Post-snapshot generic naturality and affine recursion}
\label{integration:sec:weighted-naturality}

The formal mechanism is now independent of the dyadic weights.  For weights \(w:\NaturalNumbers\to E\)
write
\[
 S_w(\omega)=\sum_{n\ge0}\omega_nw_n,
 \qquad
 \mu_w=(S_w)_*\mathbf P,
 \qquad
 w_n^+=w_{n+1}.
\]
First, let \(G,H\) be arbitrary measurable spaces, let \(g:\Omega\to G\) be measurable,
and let \(\Phi:[0,1]\times G\to H\) be measurable.  An almost-everywhere recursion
\[
 f(\omega)=\Phi\bigl(\omega_0,g(T\omega)\bigr)
 \quad\text{for \(\mathbf P\)-almost every \(\omega\)}
\]
implies
\begin{equation}
 f_*\mathbf P
 =\Phi_*\bigl(\operatorname{Unif}[0,1]\otimes g_*\mathbf P\bigr).
 \label{integration:eq:generic-ae-recursion}
\end{equation}
No separate measurability hypothesis on \(f\) is required; a pointwise recursion is an
immediate specialization.

Now let \(E\) be a real Banach space, let \(F\) be a real normed space, let
\(L:E\to F\) be continuous linear, and suppose \(\sum_n\lVert w_n\rVert<\infty\).
The target \(F\) need not be complete.  Then
\begin{equation}
 L(S_w(\omega))=S_{Lw}(\omega),\qquad L_*\mu_w=\mu_{Lw},
 \label{integration:eq:weighted-naturality}
\end{equation}
where the law identity uses the Borel measurable structures on \(E\) and \(F\).
For every real scalar \(c\), the stronger pointwise identity
\(S_{cw}=cS_w\) holds without completeness or summability.  Its law-level version retains
source completeness, Borel structure, and norm summability because the underlying
\(\operatorname{Measure.map}\) definition is totalized.

If \(A:E\to E\) is continuous linear and \(w_{n+1}=A(w_n)\), then
\begin{align}
 A_*\mu_w&=\mu_{w^+},\\
 S_w(\omega)&=\omega_0w_0+A(S_w(T\omega)),\\
 \mu_w&=\bigl((u,z)\mapsto u w_0+A z\bigr)_*
   \bigl(\operatorname{Unif}[0,1]\otimes\mu_w\bigr).
 \label{integration:eq:weighted-operator-recursion}
\end{align}
The scalar case \(A=c\,\mathrm{id}\) requires no sign or contraction condition on \(c\)
beyond the explicit norm-summability hypothesis on the displayed weights.

For real laws there is also an absolute-continuity smoothing principle.  If \(\nu\) is any
s-finite measure on \(\RealNumbers\) and \(a\ne0\), then the image of
\(\operatorname{Unif}[0,1]\otimes\nu\) under \((u,x)\mapsto ua+x\) is absolutely continuous
with respect to Lebesgue measure.  Consequently a norm-summable real weighted law \(\mu_w\)
is absolutely continuous whenever \(w_n\ne0\) for at least one index \(n\); the nonzero
coefficient need not be the head.  Under the same hypotheses every singleton is
\(\mu_w\)-null.

The eleven public declarations in this tranche are
\lean{ProbabilityRepresentation.uniformProduct_map_of_ae_head_tail_recursion},
\lean{ProbabilityRepresentation.uniformProduct_map_of_head_tail_recursion},
\lean{ProbabilityRepresentation.weightedUniformSeries_map},
\lean{ProbabilityRepresentation.weightedUniformSeries_smul},
\lean{ProbabilityRepresentation.weightedUniformDistribution_map},
\lean{ProbabilityRepresentation.weightedUniformDistribution_smul},
\lean{ProbabilityRepresentation.weightedUniformDistribution_map_eq_shift},
\lean{ProbabilityRepresentation.weightedUniformSeries_split_of_shift_eq_map},
\lean{ProbabilityRepresentation.weightedUniformDistribution_split_of_shift_eq_map},
\lean{ProbabilityRepresentation.weightedUniformSeries_geometric_split}, and
\lean{ProbabilityRepresentation.weightedUniformDistribution_geometric_split} in
\path{WeightedUniformSeries.lean}.  The separate dyadic bridge
\lean{ProbabilityRepresentation.weightedSumDistribution_eq_weightedUniformDistribution}
in \path{ProbabilityRepresentation.lean} identifies the law \(\mu\) above with
\(\mu_w\) at \(w_n=1/2/2^n\).

The two generic declarations containing the word \emph{geometric} are conditional
scalar-shift theorems.  The compatibility surface additionally exports
\lean{ProbabilityRepresentation.weightedUniformSeries_smul_weights},
\lean{ProbabilityRepresentation.uniformProduct_map_head_tail_function},
\lean{ProbabilityRepresentation.weightedUniformDistribution_isProbabilityMeasure},
\lean{ProbabilityRepresentation.weightedUniformDistribution_smul_weights},
\lean{ProbabilityRepresentation.uniformProduct_map_head_tail_weightedUniformSeries},
\lean{ProbabilityRepresentation.weightedUniformDistribution_unitInterval}, and
\lean{ProbabilityRepresentation.weightedUniformDistribution_compl_unitInterval} through
\path{WeightedUniformSeries.lean} and \path{WeightedUniformDistribution.lean}.  The
law-level scalar wrapper is oriented as \(\mu_{cw}=(x\mapsto cx)_*\mu_w\).

The same compatibility module supplies the exact Bochner barycenter
\[
 \int_E x\,\Differential\mu_w(x)=\frac12\mathbin{\cdot}\sum_{n=0}^{\infty}w_n
\]
for norm-summable weights in a complete Borel real normed space.  This is
\lean{ProbabilityRepresentation.integral_id_weightedUniformDistribution}; it uses no sign,
order, independence expansion, or termwise-integration hypothesis.

The four-declaration absolute-continuity tranche in
\path{WeightedUniformDistribution.lean} consists of
\lean{ProbabilityRepresentation.uniformScaledAdd_absolutelyContinuous},
\lean{ProbabilityRepresentation.weightedUniformDistribution_absolutelyContinuous_of_head_ne_zero},
\lean{ProbabilityRepresentation.weightedUniformDistribution_absolutelyContinuous}, and
\lean{ProbabilityRepresentation.weightedUniformDistribution_nullSingletonClass}.  The first
allows an arbitrary s-finite tail law, the second assumes \(w_0\ne0\), and the sharp final
pair assumes only \(\exists n,\ w_n\ne0\); the last theorem returns Mathlib's
\texttt{NullSingletonClass} explicitly.
Together with its six compatibility theorems and the barycenter theorem,
\path{WeightedUniformDistribution.lean} has exactly eleven public theorems and no public
definitions.

The separate exact-support layer has ten public declarations in
\path{WeightedUniformSupport.lean}:
\lean{ProbabilityRepresentation.isOpenPosMeasure_infinitePi},
\lean{ProbabilityRepresentation.uniformProduct_isOpenPosMeasure},
\lean{ProbabilityRepresentation.support_map_eq_closure_range_of_continuous},
\lean{ProbabilityRepresentation.weightedUniformSeries_constCoordinates},
\lean{ProbabilityRepresentation.weightedUniformDistribution_support_eq_range},
\lean{ProbabilityRepresentation.range_weightedUniformSeries_eq_Icc_min_max},
\lean{ProbabilityRepresentation.weightedUniformDistribution_support_eq_Icc_min_max},
\lean{ProbabilityRepresentation.range_weightedUniformSeries_eq_Icc},
\lean{ProbabilityRepresentation.weightedUniformDistribution_support_eq_Icc}, and
\lean{ProbabilityRepresentation.weightedUniformDistribution_support_eq_unitInterval}.
Thus an arbitrary norm-summable weighted law in a complete Borel real normed space has
support exactly the range of its continuous series map.  For arbitrary signed real weights,
both range and support are exactly
\[
 \left[\sum_n\min\{w_n,0\},\ \sum_n\max\{w_n,0\}\right];
\]
cancellation creates no gaps.  Constant coordinate paths give the nonnegative specialization
\([0,\sum_n w_n]\), including the degenerate zero-total-weight case.

There is now also a named specialization.  For
\[
 w_n(q)=(1-q)q^n,\qquad S_q=S_{w(q)},\qquad \mu_q=\mu_{w(q)},
\]
the zero and successor coefficients are exact for every real \(q\).  Under \(|q|<1\),
the weights are norm-summable with sum one, \(S_q\) is continuous and measurable, and
\(\mu_q\) is a probability measure satisfying
\[
 S_q(\omega)=(1-q)\omega_0+qS_q(T\omega),\qquad
 \mu_q=\bigl((u,z)\mapsto(1-q)u+qz\bigr)_*
   \bigl(\operatorname{Unif}[0,1]\otimes\mu_q\bigr),
\]
together with pointwise and law-level reflection.  Under \(0\le q<1\), the weights are
nonnegative, the series lies in \([0,1]\), and the law gives mass one to that interval and
zero to its complement.  For every \(|q|<1\), its topological support is exactly the series
range, its law is absolutely continuous with respect to Lebesgue measure, and all singletons
are null; neither of the latter two statements needs \(q\ge0\).  Under \(0\le q<1\), both
the series range and the support are exactly \([0,1]\).  Under the weaker sharp hypothesis
\(|q|<1\), including \(q=0\) and negative contractions, its exact mean is
\[
 \int_{\RealNumbers}x\,\Differential\mu_q(x)=\frac12;
\]
no positivity hypothesis is used.

The twenty-four declarations are
\lean{ProbabilityRepresentation.geometricUniformWeight},
\lean{ProbabilityRepresentation.geometricUniformWeight_zero},
\lean{ProbabilityRepresentation.geometricUniformWeight_succ},
\lean{ProbabilityRepresentation.hasSum_geometricUniformWeight},
\lean{ProbabilityRepresentation.summable_norm_geometricUniformWeight},
\lean{ProbabilityRepresentation.geometricUniformWeight_nonneg},
\lean{ProbabilityRepresentation.geometricUniformSeries},
\lean{ProbabilityRepresentation.continuous_geometricUniformSeries},
\lean{ProbabilityRepresentation.measurable_geometricUniformSeries},
\lean{ProbabilityRepresentation.geometricUniformSeries_split},
\lean{ProbabilityRepresentation.geometricUniformSeries_reflect},
\lean{ProbabilityRepresentation.geometricUniformSeries_mem_Icc},
\lean{ProbabilityRepresentation.geometricUniformDistribution},
\lean{ProbabilityRepresentation.geometricUniformDistribution_isProbabilityMeasure},
\lean{ProbabilityRepresentation.geometricUniformDistribution_absolutelyContinuous},
\lean{ProbabilityRepresentation.geometricUniformDistribution_nullSingletonClass},
\lean{ProbabilityRepresentation.uniformProduct_map_head_tail_geometricUniformSeries},
\lean{ProbabilityRepresentation.geometricUniformDistribution_selfSimilar},
\lean{ProbabilityRepresentation.geometricUniformDistribution_reflection},
\lean{ProbabilityRepresentation.integral_id_geometricUniformDistribution_eq_one_half},
\lean{ProbabilityRepresentation.geometricUniformDistribution_Icc},
\lean{ProbabilityRepresentation.geometricUniformDistribution_compl_Icc},
\lean{ProbabilityRepresentation.geometricUniformDistribution_support_eq_range}, and
\lean{ProbabilityRepresentation.geometricUniformDistribution_support_eq_Icc} in
\path{GeometricUniformLaw.lean}.  At \(q=1/2\),
\lean{ProbabilityRepresentation.weightedCoordinateSum_eq_geometricUniformSeries_one_half}
and
\lean{ProbabilityRepresentation.weightedSumDistribution_eq_geometricUniformDistribution_one_half}
identify the dyadic series and law with this specialization, while
\lean{ProbabilityRepresentation.weightedSumDistribution_support_eq_Icc} records the exact
support of the dyadic law.

\paragraph{Contractive affine independent-copy laws.}
Compiler-validated checkpoint \texttt{d312c0603} adds a generic layer which does not depend
on geometric weights, real-valued digits, or uniformity.  Let \(D\) be any measurable
space, let \(E\) be a complete second-countable real inner-product space with its Borel
measurable structure, let \(\rho\) be a probability measure on \(D\), let \(d:D\to E\) be
measurable, and let \(q\in\RealNumbers\).  For a probability measure \(\lambda\) on \(E\), set
\begin{equation}
 \mathcal A_{\rho,d,q}(\lambda)
  =(d_*\rho)*\bigl((x\mapsto qx)_*\lambda\bigr).
 \label{integration:eq:affine-independent-copy-law}
\end{equation}
This is again a probability measure, and convolution agrees with the direct independent-copy
description
\begin{equation}
 \mathcal A_{\rho,d,q}(\lambda)
 =\bigl((u,x)\mapsto d(u)+qx\bigr)_*(\rho\otimes\lambda).
 \label{integration:eq:affine-independent-copy-map}
\end{equation}
Consequently, writing \(\CharacteristicFunction{\lambda}\) for the characteristic
function of \(\lambda\),
\begin{equation}
 \CharacteristicFunction{\mathcal A_{\rho,d,q}(\lambda)}(t)
 =\CharacteristicFunction{d_*\rho}(t)\,\CharacteristicFunction{\lambda}(qt).
 \label{integration:eq:affine-independent-copy-charfun}
\end{equation}

\begin{theorem}[Affine recurrence, finite iteration, and uniqueness]
\label{integration:thm:affine-independent-copy-unique}
If \(\lambda\) is a fixed point of \(\mathcal A_{\rho,d,q}\), then
\begin{align}
 \CharacteristicFunction{\lambda}(t)
   &=\CharacteristicFunction{d_*\rho}(t)\,\CharacteristicFunction{\lambda}(qt),
      \label{integration:eq:affine-charfun-recurrence}\\
 \CharacteristicFunction{\lambda}(t)
   &=\left(\prod_{k=0}^{n-1}\CharacteristicFunction{d_*\rho}(q^kt)\right)
      \CharacteristicFunction{\lambda}(q^nt) \qquad(n\in\NaturalNumbers).
      \label{integration:eq:affine-charfun-iterate}
\end{align}
If \(|q|<1\), any two probability fixed points are equal.  The same conclusion holds when
the hypotheses are presented directly in the product-map form
\eqref{integration:eq:affine-independent-copy-map}.
\end{theorem}

\begin{proofidea}
For two fixed points, subtract the one-step recurrences.  Since every characteristic
function has modulus at most one, the common digit factor gives
\[
 \left|\CharacteristicFunction{\lambda}(t)-\CharacteristicFunction{\nu}(t)\right|
 \le
 \left|\CharacteristicFunction{\lambda}(qt)-\CharacteristicFunction{\nu}(qt)\right|.
\]
Iterating sends \(t\) to \(q^nt\to0\).  Continuity at the origin and
\(\CharacteristicFunction{\lambda}(0)=\CharacteristicFunction{\nu}(0)=1\) make the
right-hand side tend to zero, after which
characteristic-function extensionality identifies the measures.  No infinite product is
used: \eqref{integration:eq:affine-charfun-iterate} is a finite identity.  No support,
density, absolute-continuity, or moment hypothesis occurs, and both \(q=0\) and negative
\(q\) are included.
\end{proofidea}

Taking \(D=[0,1]\), \(\rho=\operatorname{Unif}[0,1]\), \(E=\RealNumbers\), and
\(d(u)=(1-q)u\) yields the named characterization
\begin{equation}
 \lambda
 =\bigl((u,x)\mapsto(1-q)u+qx\bigr)_*
   \bigl(\operatorname{Unif}[0,1]\otimes\lambda\bigr)
 \quad\Longrightarrow\quad
 \lambda=\mu_q                                             \label{integration:eq:geometric-law-unique}
\end{equation}
for every \(|q|<1\).  Thus the geometric refinement equation is a uniqueness
characterization even at \(q=0\) and for negative \(q\), with no support, density, or moment
assumption on the competing probability law.

\paragraph{Exact affine-law API crosswalk.}
The definition and eight theorems in \path{AffineIndependentCopy.lean}, all in namespace
\path{ProbabilityTheory}, are
\lean{affineIndependentCopyLaw},
\lean{affineIndependentCopyLaw_isProbabilityMeasure},
\lean{affineIndependentCopyLaw_eq_map_prod},
\lean{charFun_affineIndependentCopyLaw},
\lean{charFun_eq_mul_charFun_of_affineIndependentCopy_fixedPoint},
\lean{charFun_iterate_of_affineIndependentCopy_fixedPoint},
\lean{eq_of_charFun_affine_recurrence},
\lean{affineIndependentCopyLaw_fixedPoint_unique}, and
\lean{affineIndependentCopy_map_fixedPoint_unique}.
The last two expose respectively the operator-fixed-point and product-map formulations.
The specialization \eqref{integration:eq:geometric-law-unique} is exactly
\lean{ProbabilityRepresentation.eq_geometricUniformDistribution_of_selfSimilar} in
\path{GeometricUniformUniqueness.lean}.

\paragraph{Exact geometric sinc-prefix API crosswalk.}
The predecessor \path{GeometricUniformDictionary.lean} gives finite measure,
characteristic-function, MGF, and CGF tail factorizations.  The six public declarations in
\path{GeometricSincFactorization.lean} are
\lean{charFun_geometricUniformDigit},
\lean{charFun_geometricUniformDistribution_prefix},
\lean{charFun_geometricUniformDistribution_prefix_sinc},
\lean{tendsto_prefix_sinc_charFun},
\lean{charFun_weightedSumDistribution_prefix_sinc}, and
\lean{tendsto_prefix_sinc_charFun_weightedSumDistribution}.
The scaled-digit identity holds for every real \(q\).  Writing
\(\phi_q(t)=\operatorname{charFun}(\mu_q)(t)\), the contraction hypothesis
\(|q|<1\) gives, for every \(m\in\NaturalNumbers\) and \(t\in\RealNumbers\),
\begin{equation}
 \phi_q(t)=e^{\ImaginaryUnit(1-q^m)t/2}
 \left(\prod_{k=0}^{m-1}
   \SincRad\!\left(\frac{(1-q)q^kt}{2}\right)\right)
 \phi_q(q^m t),                                           \label{integration:eq:geometric-sinc-residual}
\end{equation}
and the phase-bearing prefix obtained by removing \(\phi_q(q^m t)\) tends pointwise to
\(\phi_q(t)\).  This includes \(q=0\) and negative \(q\); the last two declarations are
the exact \(q=1/2\) weighted-sum versions.  This is a sequence-of-prefixes
\texttt{Tendsto} theorem.

The four new public theorems in \path{GeometricReciprocalGamma.lean} are
\lean{hasProdLocallyUniformly_geometricSincProduct},
\lean{geometricSincProductFactors_multipliable},
\lean{hasProd_geometricSincProduct}, and
\lean{geometricSincProduct_differentiable}.  For every complex \(q\) with
\(\lVert q\rVert<1\), they prove locally uniform convergence on \(\ComplexNumbers\) of the defining
products for \(S_q(z)=\operatorname{geometricSincProduct}(q,z)\), pointwise
multipliability and the exact \texttt{HasProd} identity, and entireness of \(S_q\).
The zero-definition, four-theorem module
\path{GeometricSincCharacteristicFunction.lean} consists exactly of
\lean{charFun_geometricUniformDistribution_eq_phase_mul_geometricSincProduct} and
\lean{charFun_geometricUniformDistribution_eq_phase_mul_geometricReciprocalGamma},
followed by
\lean{tendstoLocallyUniformly_prefix_sinc_charFun} and
\lean{tendstoUniformlyOn_prefix_sinc_charFun}.
Writing \(G_q(z)=\operatorname{geometricReciprocalGamma}(q,z)\) and
\(z_q(t)=(1-q)t/(2\pi)\), these declarations give, for every real \(|q|<1\) and
\(t\in\RealNumbers\),
\begin{align}
 \phi_q(t)
 &=e^{\ImaginaryUnit t/2}S_q\bigl(z_q(t)\bigr),
 \label{integration:eq:geometric-sinc-characteristic}\\
 &=e^{\ImaginaryUnit t/2}G_q\bigl(z_q(t)\bigr)G_q\bigl(-z_q(t)\bigr).
 \label{integration:eq:geometric-reciprocal-gamma-characteristic}
\end{align}
The scope again includes \(q=0\) and negative \(q\).  If
\[
 P_m(t)=e^{\ImaginaryUnit(1-q^m)t/2}
 \prod_{k=0}^{m-1}\SincRad\!\left(\frac{(1-q)q^kt}{2}\right),
\]
the last two theorems state that \(P_m\to\phi_q\) locally uniformly on the real frequency
line and, explicitly, uniformly on every compact \(K\subset\RealNumbers\).  This is stronger than the
predecessor's pointwise limit, while remaining a real-frequency result; no declaration yet
packages locally uniform convergence of the full phase-bearing prefix sequence on the
complex frequency plane.  These are frequency-side convergence statements and assert no
density identity or density regularity consequence.

The adjacent one-definition, nine-theorem module
\path{RvachevPochhammerFactorization.lean} has the exhaustive surface
\lean{complexQPochhammerInf},
\lean{complexQPochhammerInf_eq_tprod},
\lean{multipliable_one_sub_mul_pow_complex},
\lean{hasProd_complexQPochhammerInf},
\lean{tendsto_finiteQPochhammerIn_complex},
\lean{summable_norm_sineTerm_qpow_pair},
\lean{geometricSincProduct_eq_tprod_pair},
\lean{geometricSincProduct_eq_tprod_complexQPochhammerInf},
\lean{rvachevFourierProduct_eq_tprod_complexQPochhammerInf}, and
\lean{rvachevFourier_eq_tprod_complexQPochhammerInf}.
For complex \(q,z\) with \(\lVert q\rVert<1\), the middle three new declarations justify
the scale--zero index exchange and prove globally
\begin{equation}
 S_q(z)=\prod_{k\ge0}
 \left(\frac{z^2}{(k+1)^2};q^2\right)_\infty .
 \label{integration:eq:geometric-sinc-pochhammer}
\end{equation}
This includes \(q=0\) and individual zero factors; the last two declarations are its
dyadic Rvachev-product and bounded-Fabius Fourier specializations.  The compound frontier
still lacks a named centered/MGF wrapper, the outside-disk reciprocal formula,
the pole divisor, and the zero--pole exchange.  The separate
zero-definition/three-theorem module
\path{GeometricPochhammerNormalConvergence.lean} now supplies the formerly
missing outer-product convergence statement.  Its declarations are
\lean{hasProdLocallyUniformly_geometricSincProduct_complexQPochhammerInf},
\lean{hasProdLocallyUniformly_rvachevFourierProduct_complexQPochhammerInf}, and
\lean{hasProdLocallyUniformly_rvachevFourier_complexQPochhammerInf}.  The first
holds on the whole complex plane for every \(\lVert q\rVert<1\), including
\(q=0\); the last two are the nome-\(1/4\) Rvachev-product and bounded-Fabius
specializations.  Hence the compound spectral theorem remains partial even
though its locally-uniform/normal-convergence clause is now exact.

The downstream zero-definition, five-theorem leaf
\path{QPochhammerEntire.lean} has the exhaustive public surface
\lean{hasProdLocallyUniformly_complexQPochhammerInf},
\lean{complexQPochhammerInf_differentiable},
\lean{complexQPochhammerInf_eq_zero_iff},
\lean{complexQPochhammerInf_eq_zero_iff_eq_inv_pow}, and
\lean{analyticOrderAt_complexQPochhammerInf_of_eq_zero}.
For each fixed complex \(q\) satisfying exactly \(\lVert q\rVert<1\), the
factors \(a\mapsto1-aq^j\) converge locally uniformly on the whole complex
\(a\)-plane to \(\lean{complexQPochhammerInf}\), which is entire in \(a\).
The third theorem gives the exact raw locus
\[
 \lean{complexQPochhammerInf}(a,q)=0
 \quad\Longleftrightarrow\quad
 \exists j\in\NaturalNumbers,\;1-aq^j=0,
\]
the fourth rewrites this as the reciprocal-power lattice when \(q\ne0\), and
the fifth gives analytic order one at every such zero.  The division-free form
includes \(q=0\), with its unique zero \(a=1\).  There is no joint holomorphy
theorem in \(q\) or global entire-function growth, order, type, or asymptotic
claim; the outer local-uniform theorem belongs to the separate module above.

The finite and general-ring companion layers are separate.
\path{QPochhammerDissection.lean} has zero definitions and exactly two
theorems, \lean{finiteQPochhammerIn_dissection} and
\lean{finiteQPochhammerIn_dissection_remainder}.  Over every commutative ring,
they give the exact residue-class split at length \(rn\) and, under exactly
\(u\le r\), the partial-period split with lengths \(n+1\) for \(s<u\) and
\(n\) for \(u\le s<r\).  They assume nothing about the nome and use no
topology, division, cancellation, or nonvanishing premise.

\begin{theorem}[Entire fixed-nome Pochhammer symbol and simple zero lattice]
\label{integration:thm:qpochhammer-entire}
Let \(q\in\ComplexNumbers\) satisfy \(\lVert q\rVert<1\), and write
\[
 P_q(a)=(a;q)_\infty=\prod_{j=0}^{\infty}(1-aq^j).
\]
The finite prefixes converge locally uniformly in \(a\) on all of
\(\ComplexNumbers\), and \(P_q\) is entire.  Its zeros satisfy
\begin{align}
 P_q(a)=0
 &\Longleftrightarrow \exists j\in\NaturalNumbers:\ 1-aq^j=0,
 \label{integration:eq:qpochhammer-factor-zero}\\
 &\Longleftrightarrow \exists j\in\NaturalNumbers:\ a=(q^j)^{-1}
 \qquad(q\ne0).                                             \label{integration:eq:qpochhammer-inverse-lattice}
\end{align}
Every zero has complex analytic order one.  The first zero description is the
total one: it includes \(q=0\), when the only vanishing factor has index zero,
the only zero is \(a=1\), and it is simple.
\end{theorem}

\begin{proofidea}
Use the generic product-at-summable-scales theorem with \(f(w)=1-w\) and
scales \(q^j\).  The geometric estimate
\(\sum_j\lVert q^j\rVert<\infty\), the identity
\(f(w)-1=-w=O(w)\), and entireness of \(f\) give the locally uniform product
and entireness of \(P_q\).  The generic no-hidden-zero theorem gives
\eqref{integration:eq:qpochhammer-factor-zero}.  If \(q\ne0\), then \(q^j\ne0\),
so \(1-aq^j=0\) is equivalent to \(a=(q^j)^{-1}\), proving
\eqref{integration:eq:qpochhammer-inverse-lattice}.

For simplicity, first show that a zero index is unique.  If \(q=0\), the
equation \(aq^j=1\) forces \(j=0\).  If \(q\ne0\), two zero indices \(j,k\)
give \(a\ne0\) and \(q^j=q^k\); taking norms and using
\(0<\lVert q\rVert<1\) forces \(j=k\).  Splitting the product after the unique
zero index gives the global factorization
\[
 P_q(z)=(z-a)U(z),\qquad
 U(z)=-q^j(z;q)_jP_q(zq^{j+1}).
\]
The cofactor is entire.  At \(z=a\), its constant \(q^j\) is nonzero (also in
the \(q=0\) case, where \(j=0\)), its finite prefix contains no zero by index
uniqueness, and its tail cannot vanish because a tail zero would shift to a
second zero factor with index \(j+1+k\).  Thus \(U(a)\ne0\), and the displayed
linear factorization gives analytic order one.
\end{proofidea}

\paragraph{Exact fixed-nome Pochhammer API crosswalk.}
The new \path{QPochhammerEntire.lean} leaf has no public definitions and
exactly five public theorems:
\lean{hasProdLocallyUniformly_complexQPochhammerInf},
\lean{complexQPochhammerInf_differentiable},
\lean{complexQPochhammerInf_eq_zero_iff},
\lean{complexQPochhammerInf_eq_zero_iff_eq_inv_pow}, and
\lean{analyticOrderAt_complexQPochhammerInf_of_eq_zero}.
They are fixed-\(q\), single-symbol results.  They prove no joint holomorphy in
\((a,q)\), outside-disk reciprocal formula, or centered
characteristic-function/MGF package.  The separate three-theorem module above
proves local-uniform convergence of the outer product in
\eqref{integration:eq:geometric-sinc-pochhammer}, but none of those remaining
compound spectral clauses.  This is a source-only overlay: the retained
257-page canonical PDF predates both these five theorems and the outer-product
convergence tranche and remains a historical render until a three-pass
Libertinus rebuild is performed.

\paragraph{Exact generic infinite-product and dissection crosswalk.}
The newer generic layer is complementary to, and not a renaming of, the
complex wrapper just described.  The complete public surface of
\path{QPochhammerDissection.lean} is the two theorems
\lean{finiteQPochhammerIn_dissection} and
\lean{finiteQPochhammerIn_dissection_remainder}.  Over every commutative ring,
the first proves \((a;q)_{rn}=\prod_{s<r}(aq^s;q^r)_n\) for all natural
\(r,n\); the second proves the two-block formula at length \(rn+u\) under
exactly \(u\le r\), including \(u=r\).  Neither theorem uses a topology, a
norm, a nonvanishing hypothesis, or any restriction on \(q\).

The exhaustive public surface of \path{QPochhammerInfinite.lean} is the
definition \lean{qPochhammerInfIn} and twenty-nine theorems.  In an arbitrary
topological commutative ring the definition is the total
\(\prod'_{j\ge0}(1-aq^j)\), and \lean{qPochhammerInfIn_eq_tprod} records this
defining equality; no convergence is implied outside a multipliable range.
For a normed ring with \(\lVert1\rVert=1\), the elementary inputs are
\lean{summable_norm_mul_pow},
\lean{one_sub_ne_zero_of_norm_lt_one}, and
\lean{norm_mul_pow_self_lt_one}.  The finite result
\lean{finiteQPochhammerIn_self_ne_zero} additionally assumes a commutative
ring with no zero divisors.

For a complete normed commutative ring with \(\lVert1\rVert=1\), under exactly
\(\lVert q\rVert<1\), the seven convergence and factorization declarations are
\lean{multipliable_one_sub_mul_pow_of_norm_lt_one},
\lean{hasProd_qPochhammerInfIn},
\lean{tendsto_finiteQPochhammerIn_qPochhammerInfIn},
\lean{qPochhammerInfIn_eq_finite_mul_shift},
\lean{qPochhammerInfIn_succ_shift},
\lean{qPochhammerInfIn_eq_factor_mul}, and
\lean{qPochhammerInfIn_dissection}.  They give multipliability, the named
\texttt{HasProd}, convergence of finite prefixes, arbitrary-prefix
concatenation, the first-factor shift, removal of an arbitrary factor, and
the infinite residue-class identity; only the last also assumes \(0<r\).
Adding a multiplicative norm gives the three exact zero declarations
\lean{qPochhammerInfIn_ne_zero},
\lean{qPochhammerInfIn_eq_zero_iff}, and
\lean{qPochhammerInfIn_self_ne_zero}; the middle statement is the
division-free equivalence with \(aq^j=1\), so it still includes \(q=0\).

The unconditional field support layer consists of
\lean{qPochhammerInfIn_eq_tprod_smul} and
\lean{isBigO_one_sub_sub_one} over a normed field;
\lean{summable_norm_pow_of_norm_lt_one} additionally assumes
\(\lVert q\rVert<1\), followed by
\lean{differentiable_finiteQPochhammerIn} over a nontrivially normed field.
For a complete normed field, \(\lVert q\rVert<1\), and \(q\ne0\),
\lean{qPochhammerInfIn_eq_zero_iff_exists_inv_pow} rewrites the zero set as
\(a=q^{-j}\).  The locally uniform and continuity statements
\lean{hasProdLocallyUniformly_qPochhammerInfIn} and
\lean{continuous_qPochhammerInfIn} require the same strict contraction and
local compactness.
The algebraic cofactor normalization
\lean{pow_sq_mul_finiteQPochhammerIn_inv_pow_self} needs only a field and
\(q\ne0\).

The remaining six theorems are complex.  Under \(\lVert q\rVert<1\),
\lean{differentiable_qPochhammerInfIn} proves entireness and
\lean{hasDerivAt_qPochhammerInfIn_of_mul_pow_eq_one} gives the derivative from
the raw equation \(a_0q^j=1\), including \(q=0,j=0\).  With \(q\ne0\),
\lean{hasDerivAt_qPochhammerInfIn_inv_pow} gives
\[
 \frac{\Differential}{\Differential a}(a;q)_\infty\bigg|_{a=q^{-j}}
 =(-1)^{j+1}(q^{\binom j2})^{-1}(q;q)_j(q;q)_\infty,
\]
and \lean{deriv_qPochhammerInfIn_inv_pow_ne_zero} proves the displayed
coefficient nonzero.  Its literal conclusion is coefficient nonvanishing;
the preceding \texttt{HasDerivAt} declaration supplies the derivative
identification.  The division-free theorem
\lean{deriv_qPochhammerInfIn_ne_zero_of_mul_pow_eq_one} strengthens this to
nonvanishing of the actual derivative at every raw factor zero, including
\(q=0\), and \lean{analyticOrderAt_qPochhammerInfIn_of_eq_zero} concludes that
every zero of the generic product has analytic order exactly one.  The
separate public name \lean{complexQPochhammerInf} has the corresponding
compatibility theorem \lean{analyticOrderAt_complexQPochhammerInf_of_eq_zero},
also including \(q=0\); conversely the five-theorem wrapper contains no generic
ring-valued concatenation or dissection theorem.  Neither API proves joint
holomorphy in \(q\) or an outside-disk reciprocal theory.

\paragraph{Exact latest q-series tranche.}
The downstream \path{GaussianBinomialPalindromic.lean} has no public
definitions and exactly fourteen public theorems:
\lean{reflect_add_of_natDegree_le}, \lean{reflect_one'},
\lean{gaussianBinomial_natDegree_le},
\lean{gaussianBinomial_zero_left}, \lean{gaussianBinomial_diag'},
\lean{reflect_gaussianBinomial},
\lean{coeff_gaussianBinomial_reflect},
\lean{coeff_gaussianBinomial_zero},
\lean{coeff_gaussianBinomial_top},
\lean{gaussianBinomial_natDegree}, \lean{gaussianBinomial_monic},
\lean{two_mul_derivative_gaussianBinomial_eval_one},
\lean{coeff_gaussianBinomial_one_of_pos_of_lt}, and
\lean{coeff_gaussianBinomial_one}.  Over every commutative
semiring these give generic reflection helpers, the Gaussian degree bound,
zero and diagonal values, reflection in degree \(k(n-k)\), constant and top
coefficients one, bounded-index coefficient reversal, and the division-free
mean identity.  The interior theorem gives linear coefficient one under
exactly \(0<k<n\), while the total theorem gives one in that case and zero
otherwise, explicitly covering \(k=0\), \(k=n\), and \(n<k\).  Both hold over
every commutative semiring; exact degree and monicity alone require
nontriviality.

The companion \path{GaussianBinomialPolynomialStructure.lean} has no public
definitions and exactly five public theorems:
\lean{natDegree_gaussianBinomial_universal},
\lean{gaussianBinomial_universal_monic},
\lean{coeff_zero_gaussianBinomial_universal},
\lean{gaussianBinomial_universal_reflect}, and
\lean{coeff_gaussianBinomial_universal_symm}.  For \(k\le n\), the universal
natural-coefficient Gaussian polynomial has exact natural degree \(k(n-k)\),
is monic, has constant coefficient one, and equals its reflection in that
degree.  If \(d\le k(n-k)\), its degree-\(d\) and degree-\(k(n-k)-d\)
coefficients agree.  With the existing row symmetry and reciprocal identities,
this closes the canonical forward q-binomial structure theorem.  The adjacent
general-semiring theorem \lean{coeff_gaussianBinomial_one} supplies the
complete linear-coefficient classifier.

The source-only \path{CentralQBinomialReduction.lean} has no definitions and
six theorems: \lean{finiteQPochhammerIn_mul_neg},
\lean{finiteQPochhammerIn_two_mul},
\lean{finiteQPochhammerIn_map_ringHom},
\lean{central_gaussianBinomial_sq_mul_int},
\lean{central_gaussianBinomial_sq_mul}, and
\lean{central_gaussianBinomial_sq_div}.  They give sign pairing, even--odd
dissection, naturality, the integral certificate, the commutative-ring
division-free central squared-base reduction, and its field quotient under two
nonzero-denominator hypotheses.  \path{CyclotomicFactorization.lean} has no
definitions and seven theorems: \lean{div_add_div_le_div},
\lean{div_le_div_add_div_add_one}, \lean{mem_range_and_mem_divisors_iff},
\lean{finiteQPochhammerIn_X_eq_prod_cyclotomic},
\lean{finiteQPochhammerIn_X_eq_gaussianBinomial_mul},
\lean{prod_cyclotomic_pow_div_extend}, and
\lean{gaussianBinomial_X_eq_prod_cyclotomic}.  They give exponent bounds,
divisor reindexing, finite-product identities over commutative rings, and the
Gaussian cyclotomic factorization over an integral domain.

The four-module root-of-unity/q-Catalan tranche is also exhaustive.
\path{PrimitiveRootBlock.lean} is 0+3:
\lean{gaussianBinomial_isPrimitiveRoot_eq_zero},
\lean{neg_one_pow_mul_pow_choose_two}, and
\lean{finiteQPochhammerIn_isPrimitiveRoot}.  In a commutative integral domain,
a primitive \(d\)-th root kills \(\GaussianBinomial{d}{k}{\zeta}\) for
\(0<k<d\); under \(d>0\), the top phase is \,\(-1\) and
\((y;\zeta)_d=1-y^d\).

\path{QLucas.lean} is 0+7: \lean{add_mul_add_sub_one}, \lean{choose_two_add},
\lean{coeff_finiteQPochhammerIn_neg_X},
\lean{finiteQPochhammerIn_neg_X_block}, \lean{coeff_block_pow_mul},
\lean{pow_choose_two_add_mul_eq}, and \lean{gaussianBinomial_q_lucas}.
For \(d>0\), a primitive \(d\)-th root in a commutative integral domain, and
\(b,s<d\), the conclusion is
\(\GaussianBinomial{ad+b}{rd+s}{\zeta}=\binom ar
\GaussianBinomial{b}{s}{\zeta}\).
Its local \path{two_mul_choose_two} is private; the unique public theorem of
that name belongs to \path{QChuVandermonde.lean}.

\path{CyclotomicDivisibility.lean} is 0+3:
\lean{cyclotomic_exponent_eq_one_iff},
\lean{cyclotomic_dvd_gaussianBinomial_iff}, and
\lean{gaussianBinomial_mul_isPrimitiveRoot}.  For \(k\le n\), \(d>0\), the
exponent-one and rational-polynomial divisibility criteria are exactly
\(n\bmod d<k\bmod d\); a positive-order primitive root in a commutative
integral domain gives \(\GaussianBinomial{an}{bn}{\zeta}=\binom ab\).

\path{QCatalan.lean} is 1+11: the definition \lean{qCatalan}, followed by
\lean{map_qInt}, \lean{qInt_X_monic}, \lean{qInt_X_natDegree},
\lean{X_sub_one_mul_qInt}, \lean{qInt_X_eq_prod_cyclotomic},
\lean{qInt_X_dvd_gaussianBinomial_rat},
\lean{qInt_X_dvd_gaussianBinomial_int}, \lean{qInt_X_mul_qCatalan},
\lean{qCatalan_natDegree}, \lean{qCatalan_eval_one_mul}, and
\lean{qCatalan_eval_one}.  The integral quotient exists because
\([n+1]_X\mid\GaussianBinomial{2n}{n}{X}\), has degree \(n(n-1)\), and
evaluates at one to the ordinary Catalan number.

\path{NewtonInterpolation.lean} is 3+19: definitions \lean{newtonCoeff},
\lean{nodeNewtonPoly}, and \lean{newtonInterpolant}; theorems \lean{newtonCoeff_eq},
\lean{newtonCoeff_zero}, \lean{newtonCoeff_mul_prod},
\lean{nodeNewtonPoly_succ}, \lean{eval_nodeNewtonPoly},
\lean{degree_nodeNewtonPoly_lt}, \lean{nodeNewtonPoly_eq_interpolate},
\lean{eq_nodeNewtonPoly_of_eval_eq}, \lean{coeff_nodeNewtonPoly_self},
\lean{newtonCoeff_eq_sum}, \lean{nodal_range_pow},
\lean{prod_erase_pow_sub_pow}, \lean{newtonCoeff_pow_eq_sum},
\lean{newtonPoly_succ}, \lean{eval_newtonPoly},
\lean{degree_newtonPoly_lt}, \lean{newtonPoly_eq_interpolate},
\lean{eq_newtonPoly_of_eval_eq}, and \lean{coeff_newtonPoly_self}.  The final
six are compatibility wrappers for the incoming \lean{newtonInterpolant}
name; the scalar \lean{newtonPoly} from
\path{NewtonBasisGeneratingFunction.lean} remains a different object.
\path{QBetaIntegral.lean} is 1+8: definition \lean{qBeta}; theorems
\lean{qNumber_pos}, \lean{qBeta_term_eq}, \lean{qBeta_eq_prod},
\lean{qBeta_eq_qGamma}, \lean{qBeta_comm}, \lean{qBeta_pos},
\lean{qBeta_add_one_left}, and \lean{qBeta_add_one_right}.  The first module
retains its field, distinct-node, nonzero-product, and geometric-grid
hypotheses; the second retains \(0<q<1\) and its positive-argument hypotheses.

Six further facade-reachable modules contribute exactly four definitions and
sixty-five theorems.  The small public inventories are
\lean{continuous_gaussianBinomial},
\lean{tendsto_gaussianBinomial_nhds_one}, and
\lean{gaussianBinomial_eq_finiteQPochhammerIn_div} in the zero-definition,
three-theorem \path{GaussianBinomialContinuity.lean};
\lean{sum_gaussianBinomial_succ_mul},
\lean{sum_gaussianBinomial_succ_mul'},
\lean{Commute.gaussianBinomial_left}, and
\lean{Commute.gaussianBinomial_right} in the zero-definition, four-theorem
\path{QPascalSummation.lean}; and \lean{quantumPlane_mul_pow} with
\lean{quantum_binomial} in the zero-definition, two-theorem
\path{QuantumBinomial.lean}.

The one-definition/twenty-two-theorem
\path{QBinomialTheoremInfinite.lean} surface is
\lean{gaussianMajorant}, \lean{hasSum_of_tendsto_of_dominated},
\lean{one_le_qPochhammerInfIn_neg},
\lean{finiteQPochhammerIn_neg_le_qPochhammerInfIn},
\lean{qPochhammerInfIn_nonneg},
\lean{qPochhammerInfIn_le_finiteQPochhammerIn},
\lean{qPochhammerInfIn_pos_of_lt_one},
\lean{qPochhammerInfIn_zero_left},
\lean{norm_finiteQPochhammerIn_le_neg_norm},
\lean{finiteQPochhammerIn_norm_le_norm},
\lean{norm_finiteQPochhammerIn_le},
\lean{qPochhammerInfIn_norm_le_norm_finiteQPochhammerIn},
\lean{finiteQPochhammerIn_ne_zero_of_norm_lt_one},
\lean{gaussianBinomial_eq_div}, \lean{gaussianMajorant_nonneg},
\lean{norm_gaussianBinomial_le},
\lean{tendsto_finiteQPochhammerIn_pow_atTop},
\lean{tendsto_gaussianBinomial_atTop},
\lean{summable_pow_choose_two_mul_pow}, \lean{hasSum_euler_product},
\lean{qPochhammerInfIn_ne_zero_of_norm_lt_one},
\lean{hasSum_qBinomial_theorem}, and \lean{hasSum_euler_reciprocal}.

The subsequent zero-definition/ten-theorem
\path{GaussianBinomialFixedColumnRate.lean} surface is
\lean{norm_finiteQPochhammerIn_pow_sub_one_le_exp},
\lean{norm_finiteQPochhammerIn_pow_sub_one_le},
\lean{norm_finiteQPochhammerIn_self_mul_gaussianBinomial_sub_one_le},
\lean{norm_gaussianBinomial_sub_inv_finiteQPochhammerIn_le},
\lean{norm_gaussianBinomial_add_sub_inv_finiteQPochhammerIn_le},
\lean{tendsto_gaussianBinomial_add_atTop},
\lean{gaussianBinomial_fixedColumn_relativeError_isBigO},
\lean{gaussianBinomial_shifted_fixedColumn_relativeError_isBigO},
\lean{gaussianBinomial_fixedColumn_error_isBigO}, and
\lean{gaussianBinomial_shifted_fixedColumn_error_isBigO}.
For a normed commutative ring with normalized multiplicative norm and
\(\lVert q\rVert\le1\), the first two bound
\(\lVert(q^m;q)_k-1\rVert\) by
\(\exp(k\lVert q\rVert^m)-1\) and
\(k e^k\lVert q\rVert^m\); the third gives
the denominator-free relative Gaussian bound for \(k\le n\), still meaningful
at roots of unity.  Over any normed field, \(\lVert q\rVert<1\) gives the
fixed and shifted nonasymptotic additive estimates, shifted convergence, and
the four relative/additive \texttt{IsBigO} laws at rates \(q^{n-k+1}\) and
\(q^{n+1}\).  No completeness or \(q\ne0\) premise is present and no sharpness
is claimed for the elementary constant.  With
\lean{tendsto_gaussianBinomial_atTop}, this proves all four clauses of
\nolinkurl{thm:fixed-column-limit}; the relative pair is the two displayed
multiplicative \(1+O(\cdot)\) statements.

The one-definition/fourteen-theorem
\path{RvachevAppellHasse.lean} surface is headed by
\lean{Appell.polynomialTransform}; its theorems are
\lean{Appell.polynomialTransform_apply},
\lean{Appell.polynomialTransform_monomial},
\lean{Appell.polynomialTransform_eq_sum_hasseDeriv_of_natDegree_lt},
\lean{Appell.polynomialTransform_eq_sum_hasseDeriv},
\lean{rvachevReciprocalMomentRat_odd},
\lean{rvachevDeconvolutionLinearMap_eq_appellPolynomialTransform},
\lean{rvachevDeconvolvedPolynomial_eq_sum_even_hasseDeriv},
\lean{eval_hasseDeriv_prod_X_sub_C_eq_elementarySymmetricEval},
\lean{eval_rvachevDeconvolvedPolynomial_prod_X_sub_C},
\lean{eval_rvachevDeconvolvedPolynomial_qFallingPower},
\lean{lagrangeBasis_eq_nodalWeight_mul_prod_X_sub_C},
\lean{lagrangeRvachevDecoder_eq_nodalWeight_mul_sum},
\lean{geometric_nodalWeight_eq_geometricQPochhammer}, and
\lean{geometric_lagrangeRvachevDecoder_eq}.  The generic commutative-semiring
layer is the finite Appell--Hasse formula; the Rvachev specialization proves
odd reciprocal centered moments vanish and reduces deconvolution to its even
Hasse terms.  Taylor--Vieta expansion of root products then gives the
elementary-symmetric, q-falling, nodal-weight, and full geometric-decoder
formulas.  Consequently \nolinkurl{gq:prop:q-Appell-falling} is exact after
composition with
\lean{normalized_sum_Ioo_rvachevDeconvolvedPolynomial_mul_shifted_rvachevUp},
and \nolinkurl{gq:thm:gaussian-Appell-decoder} is exact after composition with
\lean{normalized_sum_Ioo_lagrangeRvachevDecoder_mul_shifted_rvachevUp}; the
separate synthesis theorems provide the atom reconstruction.  Here
\(\gamma\) is the formal sequence \lean{rvachevReciprocalMomentRat}; no
analytic convergence of a reciprocal moment-generating function is claimed.
The algebraic formulas remain total at zero or colliding nodes through
totalized field inversion, whereas their cardinal use has \(c>0\),
\(0<q<1\), and the stated mesh, interval, and degree hypotheses.  Nothing
here proves a larger matrix right inverse or decoder optimality.

The two-definition/twenty-five-theorem \path{JacobiTripleProduct.lean}
inventory is \lean{thetaExponent}, \lean{pentagonalExponent},
\lean{two_dvd_mul_sub_one}, \lean{mul_sub_one_nonneg},
\lean{two_mul_thetaExponent}, \lean{thetaExponent_natCast},
\lean{thetaExponent_neg_natCast},
\lean{choose_two_add_choose_two_succ_eq},
\lean{sum_mul_sum_eq_sum_antidiagonal}, \lean{thetaCoeff_pair_eq},
\lean{sum_antidiagonal_thetaCoeff}, \lean{finite_triple_product_poly},
\lean{pow_mul_finiteQPochhammerIn_div},
\lean{neg_one_zpow_natCast_sub}, \lean{finite_triple_product},
\lean{norm_gaussianBinomialInt_le},
\lean{tendsto_gaussianBinomialInt_central},
\lean{summable_pow_thetaExponent_mul_zpow},
\lean{hasSum_jacobi_triple_product},
\lean{hasSum_jacobi_triple_product'},
\lean{mul_three_mul_sub_one_nonneg},
\lean{two_dvd_mul_three_mul_sub_one},
\lean{two_mul_pentagonalExponent}, \lean{pentagonalExponent_eq},
\lean{pentagonalExponent_pos}, \lean{hasSum_pentagonal}, and
\lean{hasSum_pentagonal_pairs}.

Finally, the one-definition/nine-theorem
\path{RogersSzegoPolynomial.lean} surface is \lean{rogersSzego},
\lean{rogersSzego_zero}, \lean{rogersSzego_one},
\lean{rogersSzego_succ},
\lean{one_sub_pow_succ_mul_gaussianBinomial_succ_succ},
\lean{rogersSzego_dilation}, \lean{rogersSzego_three_term},
\lean{summable_norm_pow_div_finiteQPochhammerIn_self},
\lean{sum_antidiagonal_euler_mul_euler}, and
\lean{hasSum_rogersSzego_generating}.  This is a source inventory, not a claim
that the retained 257-page PDF has been rebuilt from the merged source.

\begin{theorem}[Dyadic orders identify generalized Rvachev exponents]
\label{integration:thm:generalized-rvachev-identifiability}
Let \(a:\NaturalNumbers\to\NaturalNumbers\) satisfy
\(\sum_{h\ge0}a(h)/2^h<\infty\), and define
\[
 \Phi_a(z)=\prod_{h\ge0}
   \SincRad\!\left(\frac{\pi z}{2^h}\right)^{a(h)}.
\]
Then every dyadic analytic order is the corresponding inclusive prefix:
\begin{equation}
 \OrderOperator_{z=2^n}\Phi_a(z)=\sum_{h=0}^{n}a(h).
 \label{integration:eq:generalized-rvachev-dyadic-order}
\end{equation}
In particular,
\begin{align}
 a(0)&=\operatorname{ENat.toNat}\!\left(
   \OrderOperator_{z=1}\Phi_a(z)\right),\nonumber\\
 a(n+1)&=\operatorname{ENat.toNat}\!\left(
   \OrderOperator_{z=2^{n+1}}\Phi_a(z)\right)
 -\operatorname{ENat.toNat}\!\left(
   \OrderOperator_{z=2^n}\Phi_a(z)\right).
 \label{integration:eq:generalized-rvachev-exponent-recovery}
\end{align}
For two sequences \(a,c\) satisfying their respective summability
hypotheses,
\begin{equation}
 \Phi_a=\Phi_c\quad\Longleftrightarrow\quad a=c.
 \label{integration:eq:generalized-rvachev-product-equality}
\end{equation}
\end{theorem}

\begin{proofidea}
The complete integer zero-divisor theorem says that the order at a positive
integer \(r\) is the weighted scale multiplicity
\(\sum_{h=0}^{\TwoAdicValuation(r)}a(h)\).  Since \(\TwoAdicValuation(2^n)=n\), this is exactly
\eqref{integration:eq:generalized-rvachev-dyadic-order}.  The initial sample
\(n=0\) is at \(2^0=1\), not at zero: it is \(a(0)\), whereas every
\(\Phi_a\) has value one at the origin.  The prefix identity
\(P_a(n+1)=P_a(n)+a(n+1)\) gives the second formula by taking consecutive
differences.  Lean first applies \(\operatorname{ENat.toNat}\); this loses no
information because the order formula has already identified each order with
a finite natural cast.

If \(\Phi_a=\Phi_c\) as functions, applying the order functional at each
dyadic point gives equal prefix sequences.  The zeroth prefix recovers the
head, and additive cancellation of consecutive prefixes recovers every
successor exponent.  The converse follows by substitution.  Both summability
hypotheses are necessary for this analytic formulation because the infinite
product is totalized outside the admissible class.
\end{proofidea}

\paragraph{Exact identifiability API and boundary.}
The zero-definition, six-theorem
\path{GeneralizedRvachevIdentifiability.lean} surface consists of
\lean{weightSequence_eq_of_weightedScaleMultiplicity_base_pow_eq},
\lean{analyticOrderAt_generalizedRvachevProduct_two_pow},
\lean{exponent_zero_eq_toNat_analyticOrderAt_generalizedRvachevProduct},
\lean{exponent_succ_eq_toNat_analyticOrderAt_generalizedRvachevProduct},
\lean{exponentSequence_eq_of_analyticOrderAt_two_pow_eq}, and
\lean{generalizedRvachevProduct_eq_iff}.  The first is the reusable arithmetic
core for every base \(b>1\) and every additive cancellative commutative weight
monoid; the remaining five specialize the complete generalized zero-divisor
API to admissible dyadic products.  These statements identify the analytic
divisor, meaning zero locations with multiplicities.  The unweighted zero set
alone does not retain exponent sizes, and sampling product values at dyadic
zeros records only vanishing.  No spectral-zeta reconstruction,
cumulant-sample reconstruction, or underlying probability-law
identifiability is formalized; those remain open.  This one-module,
six-declaration overlay raised its historical census from 629/8,546 to
630/8,552.  With later concurrent integration, including the six-theorem
prime-power companion row, zero-definition/three-theorem outer Pochhammer
normal-convergence leaf, two-module fixed-depth effective-inverse closure, the
eight-declaration rational-gap inverse leaf, the expanded q-series/q-calculus
union (including the completed Gaussian linear-coefficient classifier,
root-of-unity/q-Catalan tranche, Newton-interpolation/q-beta tranche,
upper-index Gaussian, q-Pfaff--Saalsch\"utz, quantum-multinomial, and
Gaussian-bound layers), and the parity-selected Rvachev synthesis leaf, the
intermediate source census is 672/8,875.  The three later
\path{BinaryWordInversions.lean}, \path{BoxPartitions.lean}, and
\path{TelescopingCertificate.lean} leaves add thirty-four declarations and
bring the historical source checkpoint to 675/8,909.  The subsequent source
union includes the exact 36-declaration Lambert branch-coordinate tranche,
exact inverse and Jacobi two-square leaves, the Lagrange--Rvachev Matrix and
weighted/geometric barycenter additions, the twelve-declaration terminating-
basic-hypergeometric closure, and the three-definition/seven-theorem
\path{GeometricRichardsonGenerating.lean} leaf.  Three later theorems in the
existing \path{GaussianBinomialCumulants.lean} module add the
characteristic-zero second-derivative formula and the two division-free raw-
second-moment and variance-numerator identities over every commutative
semiring.  The later \path{LambertWBranchGapBernoulli.lean} (0+5),
\path{GeometricUniformRealization.lean} (1+17), and
\path{RegularCentralQBinomialSum.lean} (2+1), and
\path{GaussianBinomialFixedColumnRate.lean} (0+10), and
\path{RvachevAppellHasse.lean} (1+14) leaves bring the historical source census to
the historical 910/11,525 checkpoint.  The later
\path{GaussianBinomialGreaterOneAsymptotics.lean} (0+2),
\path{RvachevLagrangeNodesOnly.lean} (1+14), and
\path{GeometricUniformMomentPolynomial.lean} (1+8) leaves bring the next
source checkpoint to 913/11,551.  The subsequent
\path{ThueMorseGammaTowerDifferential.lean} (0+3) leaf brings the live
source checkpoint to 914/11,554.  The subsequent one-theorem strengthening
of the existing \path{GeometricResidualMoments.lean} module leaves the module
count unchanged and brings the next source checkpoint to 914/11,555.  The
zero-definition/one-theorem
\path{GeometricUniformMomentPolynomialBridge.lean} leaf brings the next
historical source checkpoint to 915/11,556.  The one-theorem affine
strengthening of the existing \path{FinitePolynomialFunctional.lean} module
leaves the module count unchanged and brings the next checkpoint to
915/11,557.  The \path{ThueMorseCornerIntegral.lean} (1+4) and
\path{RvachevLegendreCentralSum.lean} (0+3) leaves then bring the successive
checkpoints to 916/11,562 and 917/11,565; the complex moment-product leaf gave
the historical 918/11,568 checkpoint.  The zero-definition/one-theorem
\path{HalfQBinomialRootSimplicity.lean} leaf gave the actual merged-main
pre-local checkpoint 919/11,569; the exterior reciprocal-germ leaf gave
920/11,572; and the zero-definition/three-theorem sharp-degree leaf gave the
historical 921/11,575 checkpoint.  The
\path{RvachevLaurentLeading.lean} (1+6) leaf then gave 922/11,582, and
\path{FinitePrefixAppellRecovery.lean} (11+17) gave the historical 923/11,610
checkpoint.  The \path{GeometricUniformMomentRatFunc.lean} (1+4) leaf gives
the historical 924/11,615 checkpoint.  The two-theorem
\path{ProbabilityLaplaceMoments.lean} extension then gives 924/11,617, and
the \path{RvachevLegendreBiorthogonality.lean} (1+1) leaf gives the historical
925/11,619 checkpoint.  The subsequent merged source union gives the live
932-module/11,688-public-declaration census, with zero missing module headers
and zero missing declaration comments.

The Laurent leaf makes \nolinkurl{is:p2:thm:Laurent-leading} Exact through
the manuscript-normalized punctured-neighborhood limit and its Fourier
coordinate, odd-core, nonvanishing, generic-cofactor, and general integer-pole
companions.  Puncturing is essential because inversion is totalized at the
pole; no lower Laurent coefficient or later Appell-coefficient asymptotic is
promoted.  It first raised the inverse concordance to 52 Lean-proved rows.
The finite-prefix leaf then makes
\nolinkurl{is:p2:thm:finite-prefix-expansion} and
\nolinkurl{is:p2:thm:exact-recovery} Exact for every \(N,n\), including
\(N=0\), with bases \(1/2\) and \(1/4\) and respectively \(n+1\) and
\(\Floor{n/2}+1\) consecutive prefixes.  Its exact degrees \(n\) and
\(\Floor{n/2}\) are outer degrees in
\(\lean{Polynomial}\,(\lean{Polynomial}\,\RationalNumbers)\); a fixed-inner-\(x\)
centered specialization can drop degree, for example for odd \(n\) at
\(x=0\).  Its prefix moments are an algebraic finite-convolution model, not a
random-variable, \lean{HasLaw}, or analytic-MGF realization.  At that
checkpoint the inverse package had 54 Lean-proved and 91 human-proved
frontier rows.

The exhaustive public surface of
\path{FinitePrefixThueMorseCollapse.lean} is zero definitions and eight
theorems:
\lean{Appell.sum_thueMorseSign_mul_eval_poly},
\lean{Fabius.sum_thueMorseSign_mul_uncenteredDyadicPrefixAppellPolynomialRat},
\lean{Fabius.sum_thueMorseSign_mul_uncenteredDyadicPrefixAppellPolynomialRat_of_lt},
\lean{Fabius.sum_thueMorseSign_mul_uncenteredDyadicPrefixAppellPolynomialRat_self},
\lean{Fabius.sum_thueMorseSign_mul_centeredDyadicPrefixAppellPolynomialRat},
\lean{Fabius.sum_thueMorseSign_mul_centeredDyadicPrefixAppellPolynomialRat_succ},
\lean{Fabius.sum_thueMorseSign_mul_centeredDyadicPrefixAppellPolynomialRat_of_lt},
and
\lean{Fabius.sum_thueMorseSign_mul_centeredDyadicPrefixAppellPolynomialRat_self}.
Writing \(A^{\mathrm{unc}}_{N,n}\) and \(A^{\mathrm{cen}}_{N,n}\) for the
two rational prefix Appell polynomials and \(s_N=1-2^{-N}\), their exact
complete signed responses are
\[
 \sum_{k<2^N}\varepsilon_k
   A^{\mathrm{unc}}_{N,n}(x+k/2^N)
 =(-1)^N2^{-\binom{N+1}{2}}n^{\underline{N}}x^{n-N}
\]
and
\[
 \sum_{k<2^N}\varepsilon_k
   A^{\mathrm{cen}}_{N,n}(x+s_N-2k/2^N)
 =2^{-\binom{N}{2}}n^{\underline{N}}x^{n-N}.
\]
The uncentered formula has the manuscript's sign and
\(2^{-\binom{N+1}{2}}\) factor, whereas the centered formula is sign-free
with \(2^{-\binom{N}{2}}\).  When \(N=m+1\), the successor theorem uses the
literal manuscript grid \(x+s_{m+1}-k/2^m\).  The two \texttt{of\_lt}
theorems supply Prouhet cancellation below the block depth, and the two
\texttt{self} theorems give the first nonzero constants.  Consequently
\nolinkurl{is:p2:thm:TM-uncentered},
\nolinkurl{is:p2:cor:Prouhet-canonical}, and
\nolinkurl{is:p2:thm:TM-centered} are Exact by composition.  Both total main
theorems include \(N=0\), strengthening the positive-depth statements.  The
leaf is rational coefficient algebra only: it creates no random variable or
\lean{HasLaw}, proves no analytic MGF, and supplies no Barnes-function
identification.  These three promotions put the 194-row inverse concordance
at 57 Lean-proved, 88 human-proved, 10 conjectural, 15 open, and 24
nonassertoric rows.

The existing \path{ProbabilityLaplaceMoments.lean} module adds exactly the
two theorems
\lean{weightedSumDistribution_real_Ici_eq_rvachevUp_of_nonneg} and
\lean{integral_pow_weightedSumDistribution_eq_mul_intervalIntegral_rvachevUp}.
The first uses atomlessness of the weighted-sum law to replace the established
strict tail by the manuscript's closed event \(X\geq t\) for every \(t\geq0\).
The second gives, for every natural \(n\geq1\),
\[
 \int_{\RealNumbers}x^n\Differential\mu_X(x)
 =n\int_0^1t^{n-1}\RvachevUp(t)\Differential t,
\]
with the expectation over the full law and compact support consumed by the
existing survival layer-cake theorem.  Together with the global up/Fabius
identities these make \nolinkurl{prop:up-tail} and
\nolinkurl{cor:up-moments} Exact.  No density theorem or degree-zero extension
is asserted.

The exhaustive public surface of
\path{RvachevLegendreBiorthogonality.lean} is the definition
\lean{rvachevLegendreAnalysisKernel} and the theorem
\lean{rvachevLegendreBiorthogonality}.  For every
\(F:\lean{BoundedFabius}\) satisfying \(\lean{IsFabius}\,F\), natural
\(M\ne0\), arbitrary \(\ell,m\), and
\(\ell\leq\lean{padicValNat}\,2\,M\), the theorem is the literal normalized
open-block identity
\[
 \frac1M\sum_{-2M<k<2M}Q_\ell^-(k/M)\Lambda_m(k/M)
 =\begin{cases}1,&m=\ell,\\0,&m\ne\ell.\end{cases}
\]
The analysis degree \(m\) is unrestricted.  The proof integrates the existing
exact degree-\(\ell\) finite Rvachev synthesis on \([-1,1]\) against
\((2m+1)P_m/2\), interchanges the finite sum and interval integral, and applies
Legendre orthogonality.  This makes \nolinkurl{thm:leg-biorthogonality}
Exact.  It does not prove rationality of \(\Lambda_m\), the reverse finite
projector, rank or trace, or Cauchy--Binet; hence
\nolinkurl{cor:leg-biorthogonal-matrices} remains Partial.

Finally,
\lean{halfQBinomial_sum_rootMultiplicity_two_pow}, composed with the existing
rational root classifier and Gaussian/half-q coefficient identification,
makes \nolinkurl{cor:halfbase-root-locus} Exact for the canonical rational
polynomial.  Injective scalar extension preserves the multiplicities, but
this does not classify all roots over every extension field.  The q forward
ledger is 181 Exact / 78 Partial / 15 None / 8 interface rows; its source
concordance is 94 Lean-proved / 384 human-proved frontier /
60 non-applicable / 9 conjectures.

The nodes-only leaf makes
\nolinkurl{cor:lag-nodes-only} Exact/Complete by composition, with rationality
only at rational evaluation points and a formal coefficient-level Bell
identity; it does not promote the larger right-inverse row.  Within this

The exhaustive nodes-only surface is the definition
\lean{rvachevDeconvolvedPolynomialRat} and the fourteen theorems
\lean{map_rvachevDeconvolvedPolynomialRat},
\lean{rvachevDeconvolvedPolynomial_eq_sum_appell},
\lean{eval_rvachevDeconvolvedPolynomial_eq_sum_even_iterateDerivative},
\lean{rvachevDeconvolvedPolynomial_prod_X_sub_C_eq_sum_appell},
\lean{eval_rvachevDeconvolvedPolynomial_lagrangeBasis_eq_sum_even_iterateDerivative},
\lean{eval_rvachevDeconvolvedPolynomial_lagrangeBasis_eq_nodalWeight_mul_sum_appell},
\lean{lagrangeRvachevDecoder_eq_nodalWeight_mul_sum_appell},
\lean{map_lagrangeBasis_ratCast},
\lean{map_rvachevDeconvolvedPolynomialRat_lagrangeBasis},
\lean{lagrangeRvachevDecoder_eq_ratCast},
\lean{rvachevRawMomentRat_eq_centeredRvachevFullMoment},
\lean{momentCumulant_rvachevRawMomentRat_eq_centeredRvachevFullCumulant},
\lean{momentCumulant_rvachevRawMomentRat_even_eq_bernoulliMersenne}, and
\lean{rvachevReciprocalMomentRat_eq_completeBellPolynomial_neg_centeredCumulant}.
Its Exact/Complete status is compositional rather than a claim of one wrapper:
rationality holds at rational evaluation points, and the Bell identity is
formal coefficient algebra rather than analytic reciprocal-MGF convergence.
Independently, the existing \path{LagrangeRvachevSynthesis.lean} theorems
\lean{normalized_sum_Ioo_lagrangeRvachevDecoder_mul_shifted_rvachevUp} and
\lean{sum_Ioo_lagrangeRvachevAtomCoefficient_mul_shifted_rvachevUp} make
\nolinkurl{thm:lag-cardinal} Exact/Complete by assembly.  The still-Partial
\nolinkurl{thm:lag-right-inverse} and decoder optimality are not promoted.

\path{GaussianBinomialGreaterOneAsymptotics.lean} consists exactly of
\lean{gaussianBinomial_gt_one_fixedColumn_relativeError_isBigO} and
\lean{gaussianBinomial_gt_one_central_isEquivalent}.  For real \(q>1\), they
give the printed fixed-column relative-error rate
\((q^{-1})^{n-k+1}\) and central equivalence
\([2m,m]_q\sim q^{m^2}(q^{-1};q^{-1})_\infty^{-1}\).
Together with evaluated reciprocity they make \nolinkurl{cor:qgreaterone}
Exact.  Natural subtraction is total and reciprocity is used only eventually
when \(k\le n\); no shifted-central, non-real, or wider-nome result is claimed.

\path{ThueMorseGammaTowerDifferential.lean} consists exactly of
\lean{hasDerivAt_mellin_mellinKernel_parameter},
\lean{hasDerivAt_thueMorseGammaLog_succ}, and
\lean{iteratedDeriv_thueMorseGammaLog}.  They prove the Mellin parameter-shift
derivative for every complex spectral exponent, the successor GammaLog law,
and its falling-factorial iteration for \(k\le r\), all under \(a>0\).
Hence \nolinkurl{p2:thm:gamma-tower} is Exact when its logarithm is read as
the existing chosen GammaLog coordinate; no principal-complex-log identity
or nonpositive-parameter differential law is asserted.

\path{ThueMorseCornerIntegral.lean} has the one definition
\lean{centeredBoxIntegral} and four theorems
\lean{centeredBoxIntegral_zero}, \lean{centeredBoxIntegral_succ},
\lean{symmetricMixedDifference_range_eq_centeredBoxIntegral}, and
\lean{symmetricMixedDifference_univ_eq_centeredBoxIntegral}.  Composed with
the existing symmetric-difference and report-grid algebra, this makes
\nolinkurl{thm:TM-corner} Exact/Complete.  The analytic theorem assumes nonnegative
half-steps, \(\lean{IsOpen}\,I\), \(\lean{OrdConnected}\,I\),
\(\lean{ContDiffOn}\,\RealNumbers\,N\,g\,I\), and containment of the full
closed symmetric segment in \(I\), rather than global smoothness.  It includes
zero steps and \(N=0\), is real-valued, and fixes the recursive integration
order.  It proves neither the following Walsh conditional-expectation
construction nor its \(2^{-N}\) normalization.

\path{RvachevLegendreCentralSum.lean} has zero definitions and the three
theorems \lean{eval_legendrePolynomial_even_zero},
\lean{eval_rvachevLegendreDeconvolutionPolynomial_even}, and
\lean{rvachevLegendreCentralSum}.  For \(F:\lean{BoundedFabius}\) satisfying
\(\lean{IsFabius}\,F\) and every \(n\), including \(n=0\), the final theorem
is the literal finite central-binomial cancellation at mesh \(4^n\), obtained
from central Legendre evaluation, evenness, support truncation, and pairing of
opposite indices.  It makes only \nolinkurl{cor:leg-central-sum} Exact; no
Jacobi decoder, reverse spectral closure, or larger Lagrange right-inverse row
is promoted.

\path{RvachevLegendreBiorthogonality.lean} has exactly the one definition
\lean{rvachevLegendreAnalysisKernel} and the one theorem
\lean{rvachevLegendreBiorthogonality}.  The definition is precisely the
translated normalized analysis kernel
\[
 \Lambda_m(c)=\frac{2m+1}{2}\int_{-1}^{1}
   \RvachevUp(x-c)P_m(x)\,\Differential x,
\]
so \nolinkurl{def:leg-Lambda} is Exact.  For every nonzero natural \(M\), all
\(\ell,m\), and \(\ell\leq\TwoAdicValuation(M)\), the theorem proves
\[
 \frac1M\sum_{-2M<k<2M}Q_\ell(k/M)\Lambda_m(k/M)
 =\begin{cases}1,&m=\ell,\\0,&m\ne\ell.\end{cases}
\]
Its integer index set is the literal open block \(\lvert k\rvert<2M\), and
degree zero is included.  Hence \nolinkurl{thm:leg-biorthogonality} is
Exact/Complete.  This leaf does not establish the support, smoothness, parity,
origin-value, and Fourier--Bessel clauses of \nolinkurl{thm:leg-Lambda}, nor
the finite-matrix projector corollary; those remain outside the promotion.

Within this
source chronology, the initial zero-definition/
four-theorem \path{LambertWBranchGapBernoulli.lean} tranche first brought the
exact-radius checkpoint to 903/11,447, and its fifth theorem raised that
historical checkpoint to 903/11,448.  The two new declarations are
\lean{summable_bernoulli_mul_pow_div_factorial_iff} and
\lean{hasSum_bernoulli_mul_pow_div_factorial_complex_iff}; the latter gives the
complex \(\operatorname{HasSum}\) value
\(\lean{complexExpm1Div}(z)^{-1}\) exactly on the open disk.  Thus the leaf crosswalks
\nolinkurl{eq:pair-Bernoulli-general} and, under the explicit standard
removable-origin convention, the whole \nolinkurl{eq:bernoulli-gen}.  This
does not claim the literal Lean quotient at zero or formalize the Guide's
nearest-zero proof route; Puiseux claims remain open.  The subsequent
one-definition/eight-theorem
\path{GeometricUniformMomentPolynomial.lean} leaf raised the next historical
source census to 904/11,457.  It supplies the recursively defined rational
polynomial, its residual-product recurrence, triangular degree bound,
reciprocal-factorial value at zero, and the cases \(P_0,\ldots,P_4\).  Its
zero-definition/one-theorem companion
\path{GeometricUniformMomentPolynomialBridge.lean} raised its incoming branch
checkpoint to 905/11,458.  Its exhaustive public surface is
\lean{geometricUniformMomentPolynomial_eval₂_eq_mgf_taylorCoefficient}, which
for every real \(\lvert q\rvert<1\) identifies that polynomial with the
finite-q-Pochhammer normalization of the Taylor coefficient of the actual
geometric-uniform MGF.  Hence \nolinkurl{p7:eq:Pn-def} is Exact in this real
contraction regime, including signed \(q\) and \(q=0\).  The subsequent
\path{GeometricUniformComplexMomentProduct.lean} leaf has the exhaustive
public surface one definition and two theorems:
\lean{geometricUniformComplexMomentProduct},
\lean{hasProdLocallyUniformly_geometricUniformComplexMomentProduct}, and
\lean{geometricUniformMomentPolynomial_eval₂_eq_complexMomentProduct_taylorCoefficient}.
For every complex \(q\) with \(\lVert q\rVert<1\), these construct the actual
infinite product, prove its locally uniform convergence on the complex plane,
and give the finite-q-Pochhammer normalization of its Taylor coefficient.  The
result is a complex analytic analogue, not a complex probability-moment
extension of \nolinkurl{p7:eq:Pn-def}.  The merged-main pre-local
\path{HalfQBinomialRootSimplicity.lean} leaf has no public definitions and the
single theorem \lean{halfQBinomial_sum_rootMultiplicity_two_pow}.  For every
\(j<n\), it proves over \(\RationalNumbers\) that the exact half-base
q-binomial coefficient polynomial has root multiplicity one at \(2^j\).
Together with the existing complete rational root classification this makes
every root in that locus simple, without an arbitrary-base, arbitrary-field,
or general cyclotomic claim.  The subsequent
\path{GeometricUniformExteriorComplexMomentGerm.lean} leaf also has the
exhaustive public surface one definition and two theorems:
\lean{geometricUniformExteriorComplexMomentGerm},
\lean{analyticAt_geometricUniformExteriorComplexMomentGerm}, and
\lean{geometricUniformMomentPolynomial_eval₂_eq_exteriorComplexMomentGerm_taylorCoefficient}.
For every complex \(q\) with \(1<\lVert q\rVert\), they define the actual
reciprocal germ \((\mathcal A_{q^{-1}}(-z))^{-1}\), prove its analyticity at
zero, and give the same finite-q-Pochhammer normalization of its Taylor
coefficient.  No boundary statement at \(\lVert q\rVert=1\) is made by this
analytic leaf.  The
zero-definition/three-theorem
\path{GeometricUniformMomentPolynomialDegree.lean} leaf consists exhaustively
of \lean{coeff_geometricUniformMomentPolynomial_choose_two},
\lean{coeff_geometricUniformMomentPolynomial_choose_two_sub_one}, and
\lean{geometricUniformMomentPolynomial_natDegree_eq}.  It identifies the
coefficient at \(\binom n2\) with \(B_n(1)/n!=(-1)^nB_n/n!\), gives the
displayed subleading coefficient for \(n\ge2\), and proves exact degree
\(\binom n2\) for \(n=1\) and even \(n\), including \(n=0\), and
\(\binom n2-1\) for odd \(n>1\).  Thus
\nolinkurl{prop:qF-P-degree-sharp} and the represented
\nolinkurl{p7:thm:Pn} are Exact.  The final
\path{GeometricUniformMomentRatFunc.lean} leaf has the exhaustive public
surface one definition and four theorems:
\lean{geometricUniformMomentRatFunc},
\lean{qFactorial_mul_geometricUniformMomentRatFunc},
\lean{eval_geometricUniformMomentRatFunc_eq_complexMomentProduct_taylorCoefficient},
\lean{eval_geometricUniformMomentRatFunc_eq_exteriorComplexMomentGerm_taylorCoefficient},
and \lean{eval_geometricUniformMomentRatFunc_one}.  It packages one rational
\(a_n=d_n/n!\), proves \([n]_X!a_n=\mathcal P_n\) globally in
\texttt{RatFunc}, identifies safe evaluation with the actual inner and
exterior coefficients under the two strict norm hypotheses, and specializes
at \(q=1\) via \([n]_1!=n!\).  Consequently
\nolinkurl{thm:qF-moment-polynomial} is Exact by composition.  The theorem
does not evaluate through genuine unit-root poles or assert unit-circle
analytic continuation, so the broader MGF and germ claims remain Partial.
The inner complex \(1+2\) leaf produced the historical 906/11,461
checkpoint, and the earlier exterior branch had the historical 907/11,464
checkpoint.  In the merged chronology the complex product gave the historical
918/11,568 checkpoint, the half-base root leaf gave the actual merged-main
pre-local 919/11,569 checkpoint, the exterior leaf gave 920/11,572, and the
sharp-degree leaf gave the historical 921/11,575 checkpoint.  The Laurent leaf
then gave 922/11,582, and the finite-prefix leaf gave the historical
923/11,610 checkpoint.  The RatFunc leaf gives the historical 924/11,615
checkpoint, the two probability theorems give 924/11,617, and the Legendre
biorthogonality leaf gives the historical 925/11,619 checkpoint.  The
subsequent merged source union gives the live census of 932 modules and
11,688 public declarations, with zero documentation gaps.
The existing
\path{GeometricResidualMoments.lean} module now has zero definitions and nine
public theorems.  Its
\lean{sum_geometricLagrangeWeight_mul_scaled_geometric_pow_of_pos} theorem
supplies the displayed positive-degree scaled-geometric moment formula, while
the new \lean{sum_geometricLagrangeWeight_mul_eval_scaled_geometric} theorem
supplies exact evaluation at zero for every polynomial whose degree is at most
the interpolation order.  Together they make
\nolinkurl{cor:scaled-geometric-moments} Exact by composition.  Both theorems
hold over an arbitrary field under injectivity of \(k\mapsto q^k\) on
\(\operatorname{range}(p+1)\); the polynomial theorem permits every scale
\(c\), including zero, and thus subsumes the manuscript's \(c\ne0\) case.
No larger interpolation or node-collision claim is made.  The theorem
\lean{sum_weight_mul_eval_affine_of_topCoeff_extractor} in the existing
zero-definition/sixteen-theorem \path{FinitePolynomialFunctional.lean} module
transports a same-ring top-coefficient extractor across \(a+bX\) over every
commutative semiring.  Composed with
\lean{halfQBinomial_negativeDyadic_polynomial_sum_eq_mersenne}, it makes
\nolinkurl{cor:geometric-prouhet-affine} Exact under the established
rational-polynomial half-base convention.  It needs neither \(b\ne0\) nor
distinct transformed nodes and includes \(b=0\) and \(n=0\); it does not
extend the concrete half-base extractor to arbitrary coefficient rings.
The q forward ledger is 181 Exact / 78 Partial / 15 None / 8 interface rows:
\nolinkurl{p7:thm:Pn} moves Partial-to-Exact and
\nolinkurl{prop:qF-P-degree-sharp} moves None-to-Exact, while the RatFunc
assembly moves \nolinkurl{thm:qF-moment-polynomial} Partial-to-Exact and the
probability extension moves \nolinkurl{prop:up-tail} and
\nolinkurl{cor:up-moments} Partial-to-Exact.
These
additions are source-only, and the retained 257-page PDF
remains explicitly historical; no PDF was built for this update.
The newest six module counts are
\path{QPochhammerInfiniteBounds.lean} (0+5),
\path{HeineTransformation.lean} (2+5),
\path{QGaussSummation.lean} (0+2),
\path{QPochhammerComplexOrder.lean} (1+4),
\path{BasicHypergeometricSeries.lean} (2+5), and
\path{QMultinomial.lean} (1+9).
The four subsequent module counts are
\path{GaussianBinomialPalindromic.lean} (0+14),
\path{JacksonIntegral.lean} (1+7),
\path{QExponential.lean} (3+8), and
\path{ThetaQuasiPeriodicity.lean} (1+6).  Their five definitions and
thirty-five theorems add Gaussian-polynomial palindromicity and the total
linear-coefficient classifier, q-exponential
and q-derivative laws, Jackson integration, and bilateral-theta
quasi-periodicity and zeros, preserving the displayed analytic hypotheses.
The four newest module counts are
\path{GaussianBinomialPolynomialStructure.lean} (0+5),
\path{JacobiCubic.lean} (0+2),
\path{QPochhammerLogDerivative.lean} (0+10), and
\path{QPochhammerOrderDerivative.lean} (0+3).  Their twenty theorems add
universal Gaussian polynomial structure, Jacobi's cubic identity, the
q-Pochhammer logarithmic derivative and Lambert-series form, and differentiation
with respect to complex order under the displayed hypotheses.
The final two module counts are
\path{CentralQBinomialReduction.lean} (0+6) and
\path{CyclotomicFactorization.lean} (0+7).  Their thirteen theorems add
finite-symbol sign pairing and even--odd dissection, the central Gaussian
reduction, and cyclotomic factorizations, with the quotient-denominator and
integral-domain assumptions retained explicitly.

The four root-of-unity and q-Catalan module counts are
\path{CyclotomicDivisibility.lean} (0+3),
\path{PrimitiveRootBlock.lean} (0+3),
\path{QCatalan.lean} (1+11), and
\path{QLucas.lean} (0+7).  Their one definition and twenty-four theorems add
the cyclotomic carry criterion, complete primitive-root q-Pochhammer blocks,
q-Lucas over integral domains, and the integral q-Catalan polynomial with its
exact degree and Catalan specialization.

The newest two module counts are
\path{QBetaIntegral.lean} (1+8) and
\path{NewtonInterpolation.lean} (3+19).  Their four definitions and
twenty-seven theorems evaluate the Jackson q-beta integral as both an
infinite-product quotient and a q-gamma quotient for \(0<q<1\) and positive
arguments, and construct Newton coefficients and interpolants with evaluation,
uniqueness, divided differences, and the explicit geometric-grid formula.  The
node-polynomial and incoming interpolant names coexist as proved compatibility
aliases, while the older scalar \path{NewtonBasisGeneratingFunction.lean}
\nolinkurl{newtonPoly} remains distinct.  No endpoint or \(q\to1\) limit theorem
is included.

The latest three q-series module counts are
\path{GaussianBinomialInteger.lean} (1+10),
\path{GaussianBinomialComplexOrder.lean} (1+5), and
\path{QPfaffSaalschutz.lean} (0+3).  Their two definitions and eighteen
theorems extend Gaussian coefficients to integer and principal-branch complex
upper parameters, prove both integer q-Pascal laws and negative-index
reflection, derive the reciprocal finite and generalized complex-order
q-binomial series, and establish the terminating balanced \({}_3\phi_2\)
q-Pfaff--Saalsch\"utz summation over a field.  Nonzero nomes and parameters,
strict-contraction bounds, and every finite-product nonvanishing hypothesis
remain explicit.

The final zero-definition/five-theorem module
\path{QuantumMultinomial.lean} decomposes natural tuple antidiagonals,
transports Gaussian symmetry to every semiring, proves q-multinomial
coefficient commutation, and establishes the ordered noncommutative
q-multinomial expansion.  It assumes exactly that the nome commute with each
variable and that \(x_jx_i=q(x_ix_j)\) for \(i<j\); the finite identity is
division-free and needs no ambient commutativity or convergence hypothesis.

The latest Gaussian-bound module \path{GaussianBinomialBounds.lean} is 0+6.
It imports the stronger \lean{finiteQPochhammerIn\_self\_pos} from
\path{GeneralQConditionNumber.lean} rather than recounting it, evaluates field
reciprocity, proves the constant-term lower and infinite-q-Pochhammer upper
bounds for nonnegative real strict contractions, and transfers them to
dimension-dominant two-sided bounds for every real \(Q>1\).  Its \(k\le n\),
nonzero, order, and strict-contraction premises remain explicit.  The
parity-selected synthesis module
\path{RvachevSuperconvergentSynthesis.lean} is 1+8: it names the selected
dyadic phase representatives and proves the corresponding monomial,
polynomial, deconvolution, and Appell sampling identities.  It asserts neither
a complete phase classification nor maximality beyond those representatives.

The final terminating-basic-hypergeometric inventory is
\path{TwoPhiOneReversal.lean} (2+12) and
\path{QChuVandermonde.lean} (0+10).  It proves the finite-to-tsum termination
bridge, reflected-parameter involution, one- and two-step reversal formulas,
both q-Chu--Vandermonde evaluations, and full-domain actual-tsum wrappers.
Accordingly the two q-Chu formulas and the reversal lemma are exact.  The
stronger provenance statement that reversal proves the second formula on the
whole displayed domain remains partial: the named by-reversal theorem needs
\(C\ne0\) and \((A;q)_n\ne0\), while the full-domain theorem instead uses a
direct denominator-cleared finite q-Cauchy proof.  No rational-continuation or
cleared commutative-ring extension is formalized.

The earlier compiler-validated checkpoint \texttt{e6f535e5d} supplied the two modules that
promote the
real-space CDF and positive-ratio density layer.  Write
\begin{equation}
 F_q(x)=\operatorname{cdf}(\mu_q)(x),
 \qquad
 f_q(x)=\frac{F_q(x/q)-F_q((x-(1-q))/q)}{1-q}.
 \label{integration:eq:geometric-cdf-density}
\end{equation}
The definitions are total for real \(q\).  The CDF is monotone, measurable, nonnegative, and
at most one for every parameter.  When \(|q|<1\), it is continuous and satisfies
\begin{equation}
 F_q(1-x)=1-F_q(x),
 \qquad F_q(1/2)=1/2.                                      \label{integration:eq:geometric-cdf-reflection}
\end{equation}
For \(0\le q<1\), including \(q=0\), it has the exact tails
\begin{equation}
 F_q(x)=0\quad(x\le0),
 \qquad F_q(x)=1\quad(x\ge1).                              \label{integration:eq:geometric-cdf-tails}
\end{equation}
For \(0<q<1\), the conditioning and ordinary interval-integral identities are
\begin{align}
 F_q(x)&=\int_0^1F_q\!\left(\frac{x-(1-q)u}{q}\right)\Differential u,
                                                                  \label{integration:eq:geometric-cdf-conditioning}\\
 &=\frac{q}{1-q}\int_{(x-(1-q))/q}^{x/q}F_q(t)\Differential t.             \label{integration:eq:geometric-cdf-ordinary-integral}
\end{align}
On the same strict range, \(F_q'(x)=f_q(x)\) at every real point.  The density is continuous,
nonnegative and symmetric about \(1/2\), vanishes on both exterior half-lines, has ordinary
support contained in \((0,1)\), topological support contained in \([0,1]\), and compact
support.  The exact measure identity and smoothness conclusions are
\begin{equation}
 \mu_q=(\operatorname{Leb}:\operatorname{Measure}(\RealNumbers))
   \mathbin{\mathrm{withDensity}}
   \bigl(x\mapsto\operatorname{ofReal}(f_q(x))\bigr),
 \qquad F_q,f_q\in C^\infty(\RealNumbers).                            \label{integration:eq:geometric-density-with-density}
\end{equation}
The proof of \(C^\infty\) bootstraps the real-space affine derivative refinement; it does not
formalize the report's Fourier-product argument.

\paragraph{Exact CDF/density API crosswalk.}
The generic module \path{ContinuousCDF.lean} exports
\lean{continuous_cdf_of_nullSingleton},
\lean{cdf_reflection_sub}, and
\lean{measure_eq_withDensity_of_cdf_hasDerivAt}.  The first converts null singletons into
CDF continuity, the second converts affine reflection invariance into CDF reflection, and the
third identifies an everywhere pointwise CDF derivative with the measure's Lebesgue density.

All of the following names lie in \path{GeometricUniformCDF.lean} and the namespace
\path{Fabius.ProbabilityRepresentation}.  The total definition and universal CDF surface are
\lean{geometricUniformCDF},
\lean{monotone_geometricUniformCDF},
\lean{geometricUniformCDF_nonneg},
\lean{geometricUniformCDF_le_one}, and
\lean{measurable_geometricUniformCDF}.  The \(|q|<1\) continuity, reflection, and midpoint
statements are
\lean{continuous_geometricUniformCDF},
\lean{geometricUniformCDF_reflection}, and
\lean{geometricUniformCDF_one_half}.  The \(0\le q<1\) tails are
\lean{geometricUniformCDF_zero_of_nonpos} and
\lean{geometricUniformCDF_one_of_one_le}.

The four nonpositive-parameter diagnostic declarations are
\lean{geometricUniformDensity_zero},
\lean{geometricUniformDensity_nonpos_of_neg},
\lean{volume_withDensity_geometricUniformDensity_eq_zero_of_nonpos}, and
\lean{geometricUniformDistribution_ne_withDensity_geometricUniformDensity_of_nonpos}.
They state respectively that the selected total formula is identically zero at \(q=0\), is
nonpositive for every \(q<0\), produces the zero \texttt{withDensity} measure for every
\(q\le0\), and differs from \(\mu_q\) whenever additionally \(|q|<1\).  These are
diagnostic theorems about the selected formula, not a Lean formalization of the paper-level
corrected signed-density replacement for nonpositive parameters.

For \(0<q<1\), the refinement, density, and derivative declarations are
\lean{geometricUniformCDF_eq_integral},
\lean{geometricUniformCDF_eq_intervalIntegral},
\lean{geometricUniformDensity},
\lean{geometricUniformCDF_hasDerivAt},
\lean{deriv_geometricUniformCDF},
\lean{continuous_geometricUniformDensity}, and
\lean{geometricUniformDensity_nonneg}.  Exterior vanishing and support are
\lean{geometricUniformDensity_zero_of_nonpos},
\lean{geometricUniformDensity_zero_of_one_le},
\lean{support_geometricUniformDensity_subset_Ioo},
\lean{support_geometricUniformDensity_subset_Icc},
\lean{tsupport_geometricUniformDensity_subset_Icc}, and
\lean{geometricUniformDensity_hasCompactSupport}.  Finally,
\lean{geometricUniformDensity_reflection},
\lean{geometricUniformDistribution_eq_withDensity},
\lean{contDiff_geometricUniformCDF}, and
\lean{contDiff_geometricUniformDensity} give density symmetry, the Radon--Nikodym
identification, and smoothness of every finite order.

The half-base crosswalk in \path{ProbabilityRepresentation.lean} is exact at every real
argument:
\lean{weightedSumCDF_eq_geometricUniformCDF_one_half} identifies the two CDF definitions,
\lean{geometricUniformCDF_one_half_eq_fabiusReal} identifies that CDF with every bounded
Fabius representative, and
\lean{geometricUniformDensity_one_half_eq_rvachevUp} proves
\begin{equation}
 f_{1/2}(x)=2\RvachevUp(2x-1).                                    \label{integration:eq:geometric-density-half-bridge}
\end{equation}

\begin{boxedremark}
\textbf{Remaining geometric-\(q\) frontier.}
The promotion is deliberately not read past its hypotheses.  The law \(\mu_q\) is absolutely
continuous for every \(|q|<1\), but the selected total function \(f_q\) is proved to be its
Radon--Nikodym density only when \(0<q<1\).  Its behavior for \(q\le0\) is not open:
\(f_0=0\), \(f_q\le0\) for \(q<0\), and
\(\operatorname{Leb}.\mathrm{withDensity}(\operatorname{ofReal}\circ f_q)=0\), which is
not \(\mu_q\) on the nonpositive contraction range.  Lean formalization of the paper-level
corrected signed-density construction for \(q=0\) and negative \(q\) remains in the research
quarantine, together with:
\begin{itemize}
\item a complex-frequency locally uniform theorem for the full phase-bearing prefix
      sequence, beyond the formal locally uniform real-frequency limit, its compact-set
      specialization, and the locally uniform convergence of the unphased complex sinc
      product; quantitative rapid decay and the Fourier-inversion construction of smooth
      densities;
\item separately named centered variables, MGF wrappers, and densities and their affine
      transport package;
\item a closed all-order derivative-refinement hierarchy and explicit endpoint-flatness
      wrappers, beyond the formal existence of all finite derivatives;
\item log-concavity, strictness, central plateaux and other parameter-dependent shape claims;
\item further transform formulas, explicit Bernoulli cumulants and Bell moments, and related
      asymptotic or Gaussian-limit statements.
\end{itemize}
The finite factorization and iteration
\eqref{integration:eq:affine-independent-copy-charfun}--%
\eqref{integration:eq:affine-charfun-iterate}, and the probability-law uniqueness theorem,
are formal; \eqref{integration:eq:geometric-sinc-residual} and its pointwise prefix limit
now supply the finite geometric sinc specialization, while
\eqref{integration:eq:geometric-sinc-characteristic}--%
\eqref{integration:eq:geometric-reciprocal-gamma-characteristic} supply the exact infinite
sinc and reciprocal-Gamma characteristic-function bridges.  The locally uniform
\texttt{HasProd} theorem controls \(S_q\) as a complex-variable product, while
\lean{tendstoLocallyUniformly_prefix_sinc_charFun} and
\lean{tendstoUniformlyOn_prefix_sinc_charFun} control the full phase-bearing prefixes on the
real frequency line and on each compact real-frequency set.  These declarations do not
supply the corresponding complex-frequency full-prefix theorem or the Fourier estimates in
the first item.
The affine measure fixed point is also the direct theorem
\lean{eq_geometricUniformDistribution_of_selfSimilar}.  The transform DDE and adjacent
Cauchy-power hierarchy are direct named results in \path{GeometricUniformCauchy.lean}; the
transform bullets above refer to the remaining Fourier/Laplace and moment formulas.
\end{boxedremark}

\begin{proposition}[The smoothing equation for the CDF]
\label{integration:prop:CDF-smoothing-detailed}
Let $H(y)=\mu(({-\infty},y])$.  Then for every $y\in\RealNumbers$,
\begin{equation}
 H(y)=\int_0^1 H(2y-u)\Differential u.
 \label{integration:eq:CDF-smoothing-detailed}
\end{equation}
For $0\le y\le1/2$ this reduces to
\begin{equation}
 H(y)=\int_0^{2y}H(t)\Differential t.
 \label{integration:eq:CDF-left-detailed}
\end{equation}
\end{proposition}

\begin{proof}
Condition on the head coordinate in \eqref{integration:eq:X-head-tail-split}.  For a fixed $u$,
\[
 \Probability\!\left(\frac{u+X'}2\le y\right)=\Probability(X'\le2y-u)=H(2y-u),
\]
where $X'$ has law $\mu$ and is independent of $u$.  Integrate over $u\in[0,1]$.  If
$0\le y\le1/2$, support implies $H(2y-u)=0$ when $u>2y$; substitute $t=2y-u$ and use the
reflection identity established below, or reverse the interval directly.
\end{proof}

\subsubsection{Coordinate reflection}

Let $R:\Omega\to\Omega$ be
\begin{equation}
 (R\omega)_n=1-\omega_n.
 \label{integration:eq:coordinate-reflection}
\end{equation}
Lebesgue measure on $[0,1]$ is invariant under $u\mapsto1-u$, so the infinite product measure
is invariant under $R$.  Moreover
\begin{equation}
 X(R\omega)=\sum_{n\ge0}\frac{1-\omega_n}{2^{n+1}}
 =1-X(\omega).
 \label{integration:eq:X-reflection}
\end{equation}
Thus the law is reflection invariant:
\begin{equation}
 (x\mapsto1-x)_\#\mu=\mu.
 \label{integration:eq:law-reflection}
\end{equation}
Conditioning on the tail \(X'\), the head coordinate \(U_0\) is uniform, so the law of
\((U_0+X')/2\) is absolutely continuous with respect to Lebesgue measure.  This smoothing
step is now formalized by
\lean{ProbabilityRepresentation.geometricUniformDistribution_absolutelyContinuous} at
\(q=1/2\);
\lean{ProbabilityRepresentation.geometricUniformDistribution_nullSingletonClass} packages
the resulting null singletons.  Thus \(H\) is continuous.  Reflection invariance then gives
\begin{equation}
 H(1-y)=1-H(y).
 \label{integration:eq:CDF-reflection}
\end{equation}
The fixed-point uniqueness theorem for the bounded Fabius equation now identifies
\begin{equation}
 H=F
 \quad\text{on all of }\RealNumbers.
 \label{integration:eq:CDF-is-F}
\end{equation}
The phrase ``on all of $\RealNumbers$'' includes the two constant tails; it is not a reference to the
signed extension $\FabiusGlobal$.

\begin{theorem}[Probability representation, global form]
\label{integration:thm:probability-representation-global}
For every $y\in\RealNumbers$,
\begin{equation}
 F(y)=\Probability\!\left(\sum_{n\ge0}\frac{U_n}{2^{n+1}}\le y\right),
 \label{integration:eq:F-as-CDF-global}
\end{equation}
where the $U_n$ are independent uniform variables on $[0,1]$.  Equivalently,
\begin{equation}
 \RvachevUp(x)=\mu(({-\infty},1-|x|])
 \label{integration:eq:up-as-CDF-global}
\end{equation}
for every $x\in\RealNumbers$.
\end{theorem}

\subsection{The survival-function Fubini identity}
\label{integration:sec:survival-fubini}

A particularly useful bridge identifies integration against the Fabius law
with an ordinary integral against $\RvachevUp$.  It has now been promoted in the
main article from the scalar differentiable bridge to a Banach-valued
\(L^1\) family.  The formal development contains interval-general
finite-measure full and clipped survival-kernel identities, together with
support-free clamped and interval-supported partial/full layer-cake forms.
This section retains the Fubini mechanism and records the exact stronger API
below.

For $t\ge0$, support, atomlessness, and
\eqref{integration:eq:CDF-is-F} give
\begin{equation}
 \mu([t,\infty))=\mu((t,\infty))=1-F(t)=\RvachevUp(t).
 \label{integration:eq:survival-is-up}
\end{equation}
The equality remains correct for $t>1$, when both sides are zero.

\begin{theorem}[Banach-valued survival-kernel identity]
\label{integration:thm:survival-kernel-Banach}
Let $E$ be a real Banach space and let $k:[0,1]\to E$ be Bochner integrable.
Then
\begin{equation}
 \int_{[0,1]}\left(\int_0^x k(t)\Differential t\right)\Differential\mu(x)
 =\int_0^1 \RvachevUp(t)\,k(t)\Differential t.
 \label{integration:eq:survival-kernel-Banach}
\end{equation}
In particular, for scalar $k$ and
$K(x)=\int_0^xk(t)\Differential t$,
\[
 \int_0^1k(t)\RvachevUp(t)\Differential t=\Expectation[K(X)].
\]
Taking \(k(t)=nt^{n-1}\) gives, for every natural \(n\geq1\),
\[
 \int_{\RealNumbers}x^n\Differential\mu(x)
 =n\int_0^1t^{n-1}\RvachevUp(t)\Differential t.
\]
The choice $E=\ComplexNumbers$, regarded as a real Banach space, gives the same identity
for complex kernels.
\end{theorem}

\begin{proof}
Put
\[
 A=\{(t,x)\in[0,1]\times\RealNumbers:t<x\},
 \qquad
 H(t,x)=\IndicatorOf{A}(t,x)k(t).
\]
Integrability of $k$ and finiteness of the probability law make $H$
Bochner integrable.  Integrating first in $t$ gives
$\int_0^xk(t)\Differential t$ for $x\in[0,1]$; integrating first in $x$ gives
$\mu((t,\infty))k(t)=\RvachevUp(t)k(t)$.  Bochner--Fubini equates the two
iterated integrals.
\end{proof}

\begin{corollary}[Expectation from the survival function]
\label{integration:thm:survival-Fubini}
Let $g,g':\RealNumbers\to\RealNumbers$ be continuous and suppose
\begin{equation}
 g(x)-g(0)=\int_0^x g'(t)\Differential t
 \qquad(0\le x\le1).
 \label{integration:eq:g-FTC-hypothesis}
\end{equation}
Then
\begin{equation}
 \int_{[0,1]}g(x)\Differential\mu(x)
 =g(0)+\int_0^1 g'(t)\RvachevUp(t)\Differential t.
 \label{integration:eq:survival-Fubini}
\end{equation}
\end{corollary}

\begin{proof}
Apply \cref{integration:thm:survival-kernel-Banach} to $k=g'$, use
\eqref{integration:eq:g-FTC-hypothesis}, and add the constant $g(0)$.
\end{proof}

\paragraph{Formalization status.}
The interval-general Banach theorem is
\lean{integral_Icc_intervalIntegral_eq_intervalIntegral_smul_survival}.
Its finite-measure clipped companion is
\lean{integral_Icc_intervalIntegral_min_eq_intervalIntegral_smul_survival}.
The Fabius full-endpoint specializations are
\lean{integral_Icc_intervalIntegral_eq_intervalIntegral_smul_rvachevUp} and
\lean{intervalIntegral_mul_rvachevUp_eq_integral_Icc_intervalIntegral};
the corresponding clipped declarations proving the full
\(c\in[0,1]\) family are
\lean{integral_Icc_intervalIntegral_min_eq_intervalIntegral_smul_rvachevUp}
and
\lean{intervalIntegral_mul_rvachevUp_eq_integral_Icc_intervalIntegral_min};
the differentiable corollary is
\lean{integral_unit_eq_zero_add_integral_deriv_mul_rvachevUp}.  The two newest
theorems in \path{ProbabilityLaplaceMoments.lean} are
\lean{weightedSumDistribution_real_Ici_eq_rvachevUp_of_nonneg}, the exact
closed-tail form above, and
\lean{integral_pow_weightedSumDistribution_eq_mul_intervalIntegral_rvachevUp},
the full-law moment identity for every natural \(n\geq1\).  Together with the
existing global up/Fabius bridges they make \nolinkurl{prop:up-tail} and
\nolinkurl{cor:up-moments} Exact.  No density theorem or degree-zero extension
is asserted.

\begin{remark}[Endpoint conventions disappear under integration]
The Fubini kernel naturally uses the open ray $(t,\infty)$, while the
probabilistic tail proposition uses $[t,\infty)$.  Atomlessness makes those
events equimeasurable.  The CDF uses $(-\infty,t]$; its complement is exact,
and no endpoint atom enters because $\mu$ has a continuous CDF.
Likewise, replacing $(0,1)$ by any of the standard half-open versions changes no Lebesgue
integral.  The formal proof nevertheless records these null-set conversions explicitly.
\end{remark}

\begin{theorem}[Banach-valued survival layer cake]
\label{integration:thm:survival-layer-cake}
Let \(E\) be a real Banach space and \(0\le c\le1\).  If \(k:\RealNumbers\to E\) is
interval-integrable on \([0,c]\), then
\begin{equation}
 \int_0^c \RvachevUp(t)k(t)\Differential t
 =\int_{\RealNumbers}\left(\int_0^{\min\{z,c\}}k(t)\Differential t\right)\Differential\mu(z).
 \label{integration:eq:survival-layer-cake-partial}
\end{equation}
If \(k\) is interval-integrable on \([0,1]\), then
\begin{equation}
 \int_0^1 \RvachevUp(t)k(t)\Differential t
 =\int_{\RealNumbers}\left(\int_0^z k(t)\Differential t\right)\Differential\mu(z).
 \label{integration:eq:survival-layer-cake-full}
\end{equation}
No positivity hypothesis is imposed on \(k\); in particular \(E=\ComplexNumbers\) gives the
complex-kernel form.
\end{theorem}

The exact specializations are
\lean{ProbabilityRepresentation.intervalIntegral_rvachevUp_smul_eq_integral_min}
and \lean{ProbabilityRepresentation.intervalIntegral_rvachevUp_smul_eq_integral} in
\path{ProbabilityLaplaceMoments.lean}.  Their reusable finite-measure versions---including
the support-free clamped identity---are the three
\lean{intervalIntegral_survival_smul_eq_integral_clamp},
\lean{intervalIntegral_survival_smul_eq_integral_min_of_ae_mem_Icc}, and
\lean{intervalIntegral_survival_smul_eq_integral_of_ae_mem_Icc} declarations in
\path{SurvivalLayerCake.lean}, all derived from
\lean{integral_smul_setIntegral_subgraph} in \path{SubgraphFubini.lean}.

The later two-theorem extension of \path{ProbabilityLaplaceMoments.lean} is
exhaustively
\lean{weightedSumDistribution_real_Ici_eq_rvachevUp_of_nonneg} and
\lean{integral_pow_weightedSumDistribution_eq_mul_intervalIntegral_rvachevUp}.
The first uses the established singleton-null law to identify
\(\mu_X([t,\infty))=\RvachevUp(t)\) for every \(t\geq0\).  Composed with the
existing global reflection and nonnegative complement formulas, it makes
\nolinkurl{prop:up-tail} Exact with the manuscript's closed-tail
\(\Probability(X\geq t)\) convention.  The second proves
\[
 \int_{\RealNumbers}x^n\,\Differential\mu_X(x)
 =n\int_0^1t^{n-1}\RvachevUp(t)\,\Differential t
 \qquad(n\in\NaturalNumbers,\ n\geq1),
\]
and makes \nolinkurl{cor:up-moments} Exact.  Its expectation is the full-line
integral against the canonical weighted-sum law; degree zero is intentionally
excluded because the displayed right side would be zero rather than the
law's total mass one.  No broader MGF, fractional-moment, or complex-moment
row is promoted.

\begin{theorem}[Atom-preserving lower-CDF and fractional Fabius layer cake]
\label{integration:thm:fractional-cdf-layer-cake}
For a finite measure \(\nu\) on \(\RealNumbers\), \(a\le b\), and a Banach-valued map \(k\)
interval-integrable on \([a,b]\),
\begin{equation}
 \int_a^b \nu(({-\infty},t])k(t)\Differential t
 =\int_{\RealNumbers}\left(
   \int_{\min\{\max\{z,a\},b\}}^b k(t)\Differential t\right)\Differential\nu(z).
 \label{integration:eq:cdf-layer-cake-clamp}
\end{equation}
For the partial identity on \([c,b]\), only \(c\le b\) and the one-sided
almost-everywhere bound \(z\le b\) are needed.  Compact support on \([a,b]\) gives the
convenient partial corollary and the full terminal-primitive identity.  No atomlessness
hypothesis is needed.
For the Fabius law \(\mu\), every \(\alpha>0\), and every \(x\in\RealNumbers\),
\begin{equation}
 I_{0+}^{\alpha}F(x)
 =\frac{1}{\Gamma(\alpha+1)}\int_{\RealNumbers}(x-z)_+^\alpha\Differential\mu(z)
 =\frac{1}{\Gamma(\alpha+1)}
   \int_{\Omega}(x-X(\omega))_+^\alpha\Differential\mathbf P(\omega).
 \label{integration:eq:fabius-fractional-stopped-power}
\end{equation}
At \(x=1\), reflection of \(\mu\) turns the stopped power into the ordinary moment:
\begin{equation}
 I_{0+}^{\alpha}F(1)
 =\frac{1}{\Gamma(\alpha+1)}\int_{\RealNumbers}z^\alpha\Differential\mu(z)
 =\frac{1}{\Gamma(\alpha+1)}
   \int_{\Omega}X(\omega)^\alpha\Differential\mathbf P(\omega).
 \label{integration:eq:fabius-fractional-endpoint-moment}
\end{equation}
\end{theorem}

The closed-supergraph engine is \lean{integral_smul_setIntegral_supergraph}.  Its four
finite-measure forms in \path{CDFLayerCake.lean} are
\lean{intervalIntegral_cdf_smul_eq_integral_clamp},
\lean{intervalIntegral_cdf_smul_eq_integral_max_of_ae_le},
\lean{intervalIntegral_cdf_smul_eq_integral_max_of_ae_mem_Icc}, and
\lean{intervalIntegral_cdf_smul_eq_integral_of_ae_mem_Icc}.  The support-free clamped,
lower-supported positive-part, probability-CDF positive-part, and compactly supported
endpoint-power consequences in
\path{FractionalCDFLayerCake.lean} are
\lean{fractionalVolterra_measureReal_Iic_eq_integral_clamp},
\lean{fractionalVolterra_measureReal_Iic_eq_integral_posPart},
\lean{fractionalVolterra_cdf_eq_integral_posPart}, and
\lean{fractionalVolterra_measureReal_Iic_eq_integral_rpow_of_ae_mem_Icc}.  Finally,
\lean{ProbabilityRepresentation.fractionalVolterra_fabiusReal_eq_integral_posPart},
\lean{ProbabilityRepresentation.fractionalVolterra_fabiusReal_eq_uniformProduct_integral_posPart},
\lean{ProbabilityRepresentation.fractionalVolterra_fabiusReal_one_eq_integral_rpow}, and
\lean{ProbabilityRepresentation.fractionalVolterra_fabiusReal_one_eq_uniformProduct_integral_rpow}
in \path{FabiusFractionalIntegral.lean} are exactly the two Fabius displays above.  These
are positive-real integral identities, not complex-order continuations, fractional
derivatives, Caputo formulas, or folded-Rvachev, quantile, or inverse-function
specializations.

\begin{remark}[Formalized inverse layer-cake boundary]
Suppose the two clocks preserve ordered compact intervals and satisfy
\(C(x)<u\Longleftrightarrow x<Q(u)\) throughout their rectangle.  The theorem
\lean{galoisConnection_Icc_restrict_of_lt_iff_lt} says that their interval
restrictions form a Galois connection and hence are monotone.  The generic
explicit-primitive and pointwise-\(C^1\) identities
\lean{intervalIntegral_smul_intervalIntegral_of_lt_iff_lt} and
\lean{intervalIntegral_smul_comp_of_lt_iff_lt} therefore construct the measurable
clamped representative needed by the subgraph engine instead of assuming either clock
measurable.  They allow real or complex weights and Banach-valued primitives.  Their
Fabius forms are \lean{intervalIntegral_smul_intervalIntegral_fabiusInv} and
\lean{intervalIntegral_smul_comp_fabiusInv}.

For a real-valued function \(\Psi\) absolutely continuous on \([a,b]\) and a real or
complex integrable weight \(\phi\), the stronger regularity interface is
\begin{equation}
 \int_c^d\phi(u)\bigl(\Psi(Q(u))-\Psi(a)\bigr)\Differential u
 =\int_a^b\Psi'(x)\left(\int_{C(x)}^d\phi(u)\Differential u\right)\Differential x.
 \label{integration:eq:ac-inverse-layer-cake}
\end{equation}
This is
\lean{intervalIntegral_mul_comp_sub_of_lt_iff_lt_of_absolutelyContinuousOnInterval};
\lean{intervalIntegral_mul_comp_of_lt_iff_lt_of_absolutelyContinuousOnInterval}
removes the subtraction under \(\Psi(a)=0\), and
\lean{intervalIntegral_mul_comp_fabiusInv_of_absolutelyContinuousOnInterval} is the
exact unit-interval Fabius wrapper.  Absolute continuity supplies derivative
integrability, so there is no separate integrability premise on the right.  Degenerate
ordered intervals are included.

The real weighted-area specialization is \lean{intervalIntegral_mul_fabiusInv_eq}.
The module also exports
\lean{intervalIntegral_fabiusInv_eq_intervalIntegral_rvachevUp} and
\lean{intervalIntegral_fabiusInv_eq_one_half}, which identify the inverse area with the
unit-interval $\RvachevUp$ area and with \(1/2\).  These eleven declarations exhaust the public
theorems of \path{InverseLayerCake.lean}.  Separately,
\lean{intervalIntegral_deriv_mul_comp_add_comp_mul_deriv_of_lt_iff_lt} in
\path{InversePairIntegral.lean} packages arbitrary real absolutely-continuous outer
functions on an arbitrary ordered compact rectangle whose two clocks preserve the displayed
intervals and satisfy the strict inverse-order relation.  Clock measurability is derived,
not assumed, and degenerate ordered intervals are included.  For every \(y\in[a,b]\), its
variable-upper companion
\lean{intervalIntegral_deriv_mul_comp_add_comp_mul_deriv_to_of_lt_iff_lt} restricts the
identity to \([a,y]\times[c,C(y)]\).  The complete and variable-upper Fabius wrappers are
\lean{intervalIntegral_deriv_mul_fabiusReal_add_fabiusInv_mul_deriv} and
\lean{intervalIntegral_deriv_mul_fabiusReal_add_fabiusInv_mul_deriv_to}; the two
identity-function specializations are
\lean{intervalIntegral_fabiusReal_add_fabiusInv_eq_one} and
\lean{intervalIntegral_fabiusInv_to_eq_mul_sub_intervalIntegral_fabiusReal}.  These six
declarations exhaust the public theorems of \path{InversePairIntegral.lean}.  They still do
not package reversed endpoints, improper singular-endpoint extensions, complex-valued
arbitrary-AC primitives, higher or fractional inverse primitives, or complex Mellin
continuation.  In particular, accepting complex weights does not by itself export the
report's named Fourier or Mellin transform corollaries.
\end{remark}

\begin{remark}[Formalized positive-real fractional Volterra boundary]
\label{integration:rem:fractional-volterra-boundary}
At compiled full-facade checkpoint \texttt{149332f9d}, the definition
\lean{fractionalVolterra} in \path{FractionalVolterra.lean} is a total
Gamma-normalized oriented interval integral for a real order, arbitrary real base and
endpoint, and an input valued in any real normed space.  The unconditional declarations
\lean{fractionalVolterra_self}, \lean{fractionalVolterra_congr},
\lean{fractionalVolterra_smul}, \lean{fractionalVolterra_one}, and
\lean{fractionalVolterra_nat_succ} give the base-point value, locality on the closed
unoriented endpoint interval, real-scalar compatibility, the ordinary order-one primitive,
and exact agreement with \lean{normalizedVolterra} at every positive natural order.
In the classical regime \(0<\alpha\) and \(a\le x\), a continuous input has an
interval-integrable fractional kernel by
\lean{intervalIntegrable_fractionalVolterra_kernel}, and
\lean{fractionalVolterra_add_input} gives additivity for two such continuous inputs.
If \(a\le b\) and \(a\le x\), the operator-independent theorem
\lean{intervalIntegral_eq_integral_min_of_eq_zero} cuts an interval-integrable kernel off
at \(\min(x,b)\) whenever it vanishes on the open right tail \((b,x)\).
Under the same endpoint bounds, \(0<\alpha\), and continuity on \([a,x]\), its fractional
specialization is
\lean{fractionalVolterra_eq_intervalIntegral_min_of_eq_zero}:
\begin{equation}
 I_{a+}^{\alpha}f(x)
 =\int_a^{\min(x,b)}
   \frac{(x-t)^{\alpha-1}}{\Gamma(\alpha)}f(t)\Differential t.
 \label{integration:eq:fractional-volterra-support-cutoff}
\end{equation}
For \(0<\alpha\), \(0<\beta\), and \(s\le x\),
\lean{intervalIntegrable_fractionalVolterra_betaKernel} in
\path{FractionalVolterraSemigroup.lean} proves interval integrability of the raw shifted
beta kernel, including the degenerate case \(s=x\).  When \(s<x\),
\lean{intervalIntegral_fractionalVolterra_betaKernel} proves
\begin{equation}
 \int_s^x(x-u)^{\alpha-1}(u-s)^{\beta-1}\Differential u
 =(x-s)^{\alpha+\beta-1}
   \frac{\Gamma(\alpha)\Gamma(\beta)}{\Gamma(\alpha+\beta)}.
 \label{integration:eq:fractional-volterra-beta-kernel}
\end{equation}
and \lean{intervalIntegral_fractionalVolterra_normalizedBetaKernel} proves
\begin{equation}
 \int_s^x
 \frac{(x-u)^{\alpha-1}}{\Gamma(\alpha)}
 \frac{(u-s)^{\beta-1}}{\Gamma(\beta)}\Differential u
 =\frac{(x-s)^{\alpha+\beta-1}}{\Gamma(\alpha+\beta)}.
 \label{integration:eq:fractional-volterra-normalized-beta-kernel}
\end{equation}
If the target is complete, \(a\le x\), and \(f\) is continuous on \([a,x]\), then
\lean{fractionalVolterra_add} gives
\begin{equation}
 I_{a+}^{\alpha+\beta}f(x)
 =I_{a+}^{\alpha}\!\left[u\mapsto I_{a+}^{\beta}f(u)\right](x).
 \label{integration:eq:fractional-volterra-semigroup}
\end{equation}
For a complete target, \(v\in E\), and \(a<x\), the same module proves
\begin{align}
 I_{a+}^{\alpha}\!\left[t\mapsto
   \frac{(t-a)^{\beta-1}}{\Gamma(\beta)}v\right](x)
 &=\frac{(x-a)^{\alpha+\beta-1}}{\Gamma(\alpha+\beta)}v,
 &&\alpha,\beta>0,
 \label{integration:eq:fractional-volterra-normalized-rpow}\\
 I_{a+}^{\alpha}\!\left[t\mapsto(t-a)^\rho v\right](x)
 &=\frac{\Gamma(\rho+1)(x-a)^{\alpha+\rho}}
        {\Gamma(\alpha+\rho+1)}v,
 &&\alpha>0,\ \rho>-1.
 \label{integration:eq:fractional-volterra-rpow}
\end{align}
The constant specialization includes \(a=x\):
\begin{equation}
 I_{a+}^{\alpha}[t\mapsto v](x)
 =\frac{(x-a)^\alpha}{\Gamma(\alpha+1)}v
 \qquad(\alpha>0,\ a\le x).
 \label{integration:eq:fractional-volterra-const}
\end{equation}
The exact declarations are \lean{fractionalVolterra_normalized_rpow_smul},
\lean{fractionalVolterra_rpow_smul}, and \lean{fractionalVolterra_const}.
The companion \path{FractionalVolterraCalculus.lean} proves increasing-affine covariance
for arbitrary real order, \(0<c\), and \(a\le x\), without regularity, integrability, or
completeness:
\begin{equation}
 I_{(ca+d)+}^{\alpha}f(cx+d)
 =c^\alpha I_{a+}^{\alpha}[t\mapsto f(ct+d)](x).
 \label{integration:eq:fractional-volterra-affine}
\end{equation}
Its inverse-smul form is \lean{fractionalVolterra_comp_affine}.  For a Banach-valued
continuous \(g\) with right-sided interior derivative \(g'\), assuming \(g'\) is
interval-integrable, \lean{fractionalVolterra_add_one_deriv} gives
\begin{equation}
 I_{a+}^{\alpha+1}g'(x)
 =I_{a+}^{\alpha}g(x)
  -\frac{(x-a)^\alpha}{\Gamma(\alpha+1)}g(a)
 \qquad(\alpha>0,\ a\le x),
 \label{integration:eq:fractional-volterra-add-one-deriv}
\end{equation}
including \(a=x\); \lean{fractionalVolterra_add_one_deriv_of_eq_zero} removes the
boundary term when \(g(a)=0\).  The total definition
\lean{rvachevFractionalPrimitive} in \path{FabiusFractionalVolterra.lean} sets
\(U_\beta(x)=I_{(-1)+}^{\beta}\RvachevUp(x)\).  Its three causal declarations prove
\begin{align}
 U_\beta(x)
 &=\frac1{\Gamma(\beta)}
   \int_{-1}^{\min(x,1)}(x-t)^{\beta-1}\RvachevUp(t)\Differential t,
 &&\beta>0,\ -1\le x,
 \label{integration:eq:rvachev-fractional-primitive-min}\\
 U_{n+1}(x)&=J_{-1}^{n+1}\RvachevUp(x),
 &&n\in\NaturalNumbers,\ x\in\RealNumbers,
 \label{integration:eq:rvachev-fractional-primitive-natural}\\
 U_{\alpha+\beta}(x)&=I_{(-1)+}^{\alpha}U_\beta(x),
 &&\alpha,\beta>0,\ -1\le x.
 \label{integration:eq:rvachev-fractional-primitive-semigroup}
\end{align}
They are \lean{rvachevFractionalPrimitive_eq_intervalIntegral_min},
\lean{rvachevFractionalPrimitive_nat_succ}, and
\lean{rvachevFractionalPrimitive_add}.  The module's other three declarations prove
\begin{align}
 I_{0+}^{\alpha+1}\FabiusGlobal(x)
 &=2^\alpha I_{0+}^{\alpha}\FabiusGlobal(x/2), &&\alpha>0,\ x\ge0,
 \label{integration:eq:fractional-shift-global-Fabius}\\
 I_{0+}^{\alpha+1}F(x)
 &=2^\alpha I_{0+}^{\alpha}F(x/2), &&\alpha>0,\ 0\le x\le1,
 \label{integration:eq:fractional-shift-bounded-Fabius}\\
 I_{-1+}^{\alpha+1}\RvachevUp(x)
 &=2^\alpha I_{0+}^{\alpha}F((x+1)/2), &&\alpha>0,\ x\ge-1.
 \label{integration:eq:fractional-shift-Up}
\end{align}
They are \lean{fractionalVolterra_add_one_extendedFabius_of_nonneg},
\lean{fractionalVolterra_add_one_fabiusReal}, and
\lean{fractionalVolterra_add_one_rvachevUp}.  No order-zero operator law,
reversed-endpoint fractional interpretation or covariance, nonpositive-scale covariance,
semigroup theorem
for merely interval-integrable inputs or noncomplete targets, Caputo or other fractional
derivative, complex-order theorem, iterated fractional shift, shifted dyadic-lattice or
endpoint-moment wrapper, transform or tail-series theorem, or inverse-function fractional
specialization is included.  The positive-real Fabius-CDF
specialization is recorded separately in
\cref{integration:thm:fractional-cdf-layer-cake}.
\end{remark}

\begin{remark}[Formalized unconditional analytic-filter boundary]
\label{integration:rem:analytic-series-filter-boundary}
At the same compiled checkpoint, \path{AnalyticSeriesFilter.lean} turns the finite moment
algebra of a row \((\xi_i,w_i)\) into an exact pointwise series theorem.  If each weighted
sampled series
\(
  \sum_{m\ge0}w_i\bigl((\xi_i z)^m a_m\bigr)
\)
is unconditionally summable, then \lean{finiteAnalyticSeriesFilter_eq_tsum} proves
\begin{equation}
 \sum_i\sum_{m\ge0}w_i\bigl((\xi_i z)^m a_m\bigr)
 =\sum_{m\ge0}\left(\sum_iw_i\xi_i^m\right)z^m a_m.
 \label{integration:eq:analytic-series-filter-diagonal}
\end{equation}
The generic statement works over a commutative semiring acting on a normed additive group,
and a zero weight needs no convergence at its node.  Exactness through degree \(p\) splits
off the reproduced finite head; at target zero the head is exactly \(a_0\).  For a geometric
row over a commutative ring with continuous constant scalar multiplication, the finite
row-moment factor in degree \(p+1+r\) is
\begin{equation}
 (-1)^p q^{\binom{p+1}{2}}\GaussianBinomial{p+r}{p}{q}.
 \label{integration:eq:analytic-series-filter-gaussian-tail}
\end{equation}
This is \lean{geometricAnalyticSeriesFilter_eq_constant_add_gaussian_tsum}.  Assuming
\(j\mapsto q^j\) is injective on \(\{0,\ldots,p\}\), the geometric-Lagrange and
shifted-scale wrappers over a scalar field impose sample summability only where the
corresponding weight is nonzero.  These are Mathlib \lean{Summable}/\lean{tsum} identities
and hence unconditional and pointwise.  They do not cover conditionally-only convergent
boundary series, prove radius or uniform convergence or quantitative norm/error bounds,
determine residual signs or asymptotic acceleration rates, instantiate an actual sinc
product, or supply a formal-\lean{PowerSeries}-evaluation bridge.
\end{remark}

\subsection{Raw, endpoint, and tilted Laplace moments}
\label{integration:sec:probability-laplace}

The main article develops half moments and the negative generating function from the
analytic side.  The new probability--Laplace module proves that those quantities are exactly
the moments of the product law.

\subsubsection{Three moment families}

For a probability measure $\nu$ supported on $[0,1]$, define
\begin{align}
 U_n(\nu)&=\int_{[0,1]}(1-x)^n\Differential\nu(x),                  \label{integration:eq:endpoint-moment-def}\\
 L_k(\nu;s)&=\int_{[0,1]}\EulerE^{-sx}x^k\Differential\nu(x),           \label{integration:eq:raw-laplace-def}\\
 J_k(s)&=\int_0^1t^k\RvachevUp(t)\EulerE^{-st}\Differential t.                 \label{integration:eq:survival-tilt-def}
\end{align}
For real $s$, the integrand in \eqref{integration:eq:raw-laplace-def} is nonnegative, so
\begin{equation}
 L_k(\nu;s)\ge0.
 \label{integration:eq:laplace-moment-nonnegative}
\end{equation}

Apply \cref{integration:thm:survival-Fubini} to $g(x)=\EulerE^{-sx}x^k$.  For $k=0$ one obtains
\begin{equation}
 L_0(\mu;s)=1-sJ_0(s).
 \label{integration:eq:laplace-zero-survival}
\end{equation}
For $k\ge0$, applying it to $g(x)=\EulerE^{-sx}x^{k+1}$ gives
\begin{equation}
 L_{k+1}(\mu;s)=(k+1)J_k(s)-sJ_{k+1}(s).
 \label{integration:eq:laplace-successor-survival}
\end{equation}
These equations are exactly the defining clauses of
\lean{fabiusLaplaceMoment}.

\begin{theorem}[Law moments equal analytic Fabius moments]
\label{integration:thm:law-moments-identification}
For every $k\in\NaturalNumbers$ and $s\in\RealNumbers$,
\begin{equation}
 L_k(\mu;s)=M_k(s),
 \label{integration:eq:law-M-identification}
\end{equation}
where
\begin{align}
 M_0(s)&=G(-s),                                            \label{integration:eq:M-zero}\\
 M_{k+1}(s)&=(k+1)J_k(s)-sJ_{k+1}(s)                      \label{integration:eq:M-successor}
\end{align}
are the tilted Fabius moments and $G$ is the generating function of the main article.
In particular $M_k(s)\ge0$ for real $s$.
\end{theorem}

\begin{proof}
The case $k=0$ is the integration-by-parts identity
\eqref{integration:eq:laplace-zero-survival}, which is also the negative generating-function formula.
The successor case is \eqref{integration:eq:laplace-successor-survival}.
\end{proof}

\subsubsection{Differential recurrence}

Differentiation under the compact integral in \eqref{integration:eq:survival-tilt-def} is justified by a
uniform exponential majorant in a neighborhood of each real $s$.  It gives
\begin{equation}
 J_k'(s)=-J_{k+1}(s).
 \label{integration:eq:J-differential-recurrence}
\end{equation}
Substituting into \eqref{integration:eq:M-zero}--\eqref{integration:eq:M-successor} yields the cleaner recurrence
\begin{equation}
 M_k'(s)=-M_{k+1}(s).
 \label{integration:eq:M-differential-recurrence}
\end{equation}
Iteration gives
\begin{equation}
 M_k^{(n)}(s)=(-1)^nM_{k+n}(s).
 \label{integration:eq:M-iterated-derivative}
\end{equation}
Thus all $M_k$ are $C^\infty$ on $\RealNumbers$, and the negative generating function is completely
monotone:
\begin{equation}
 (-1)^n\frac{\Differential^n}{\Differential s^n}G(-s)=M_n(s)\ge0.
 \label{integration:eq:complete-monotonicity}
\end{equation}

\subsubsection{Reflection and the half moments}

Reflection invariance \eqref{integration:eq:law-reflection} gives
\begin{equation}
 \int(1-x)^n\Differential\mu(x)=\int x^n\Differential\mu(x).
 \label{integration:eq:endpoint-raw-reflection}
\end{equation}
At zero tilt the analytic moment normalization is the exact rational half-moment sequence
$d_n$:
\begin{equation}
 M_n(0)=d_n.
 \label{integration:eq:M-zero-half-moment}
\end{equation}
Combining the identifications yields the full probability interpretation.

\begin{theorem}[Probability meaning of the half moments]
\label{integration:thm:half-moments-probability}
For every $n\in\NaturalNumbers$,
\begin{equation}
 d_n=\Expectation[X^n]=\Expectation[(1-X)^n]
     =\int_{[0,1]}x^n\Differential\mu(x)
     =\int_{[0,1]}(1-x)^n\Differential\mu(x).
 \label{integration:eq:half-moment-probability}
\end{equation}
Moreover
\begin{equation}
 \left.\frac{\Differential^n}{\Differential s^n}G(-s)\right|_{s=0}=(-1)^nd_n.
 \label{integration:eq:generating-jet-half-moments}
\end{equation}
\end{theorem}

\begin{proof}
Use \eqref{integration:eq:law-M-identification} at $s=0$, then
\eqref{integration:eq:M-zero-half-moment} and \eqref{integration:eq:endpoint-raw-reflection}.  The derivative identity
is \eqref{integration:eq:M-iterated-derivative} with $k=0$ and $s=0$.
\end{proof}

\begin{boxedremark}
The equality $d_n=\Expectation[X^n]$ is not merely an interpretation of a previously known rational
sequence.  It is the bridge that licenses positivity, log-convexity, moment-matrix
positivity, and Laplace comparison arguments without rebuilding them from the recurrence for
$d_n$.  The formal library keeps the bridge in a separate module precisely so the
probabilistic and saddle branches can import it independently.
\end{boxedremark}

The probability chapter hands its zero-th tilted moment directly to the endpoint
asymptotic synthesis.  At \(k=0\),
\eqref{integration:eq:law-M-identification} identifies the zero-th tilted law moment with
\(M_0(s)=G(-s)\).  For \(s>0\), the notation of \cref{part:dyadicasym} calls this same
negative Laplace transform \(B(s)\); its independent-coordinate product is already recorded
in \eqref{dyadicasym:eq:B-product}.  Thus the exact negative-Laplace product decomposed there
is the Laplace transform of the same product-space law, rather than a second independent
model.

\subsection{From atomic laws to half-endpoint histograms}
\label{integration:sec:step-measure-bridge}

The main article introduces the polynomial step approximants $\Wcal_n$ and then moves quickly
to their convergence.  The formal proof has a nontrivial middle layer: the coefficients first
define an atomic probability measure; each atom is smeared across a centered cell; interval
integrals of the resulting histogram are sandwiched between masses of slightly shrunken and
enlarged intervals.

\subsubsection{Half-endpoint indicators are ordinary indicators almost everywhere}

For $a<b$, define
\begin{equation}
 \IndicatorOf{[a,b]}^{*}(x)=
 \begin{cases}
  1,&a<x<b,\\
  \tfrac12,&x=a\text{ or }x=b,\\
  0,&\text{otherwise}.
 \end{cases}
 \label{integration:eq:half-endpoint-indicator}
\end{equation}
The two endpoint values are essential for pointwise gluing of adjacent cells, but irrelevant
to Lebesgue integration.

\begin{lemma}[Almost-everywhere replacement]
\label{integration:lem:half-indicator-ae}
For every $a,b\in\RealNumbers$,
\begin{equation}
 \IndicatorOf{[a,b]}^{*}=\IndicatorOf{(a,b]}
 \quad\text{Lebesgue-almost everywhere}.
 \label{integration:eq:half-indicator-ae}
\end{equation}
Consequently
\begin{equation}
 \int_\RealNumbers\IndicatorOf{[a,b]}^{*}(x)\Differential x=\max(b-a,0),
 \label{integration:eq:half-indicator-total-integral}
\end{equation}
and, for $c\le d$,
\begin{equation}
 \int_c^d\IndicatorOf{[a,b]}^{*}(x)\Differential x
 =\max\bigl(\min(d,b)-\max(c,a),0\bigr).
 \label{integration:eq:half-indicator-overlap}
\end{equation}
\end{lemma}

\begin{proof}
The functions differ only at $a$ and $b$, a null set.  The first integral is the length of
$(a,b]$.  The second is the length of
$(c,d]\cap(a,b]=(\max(c,a),\min(d,b)]$, with truncation at zero when the interval is empty.
\end{proof}

\subsubsection{Atomic measure, cell geometry, and exact normalization}

At level $n$, let $x_{n,m}$ be the location of the $m$th polynomial atom and let
$p_{n,m}\ge0$ be its normalized mass, with
\begin{equation}
 \sum_m p_{n,m}=1.
 \label{integration:eq:atomic-mass-one}
\end{equation}
The corresponding atomic law is
\begin{equation}
 \nu_n=\sum_m p_{n,m}\delta_{x_{n,m}}.
 \label{integration:eq:polynomial-atomic-law}
\end{equation}
Let
\begin{equation}
 h_n=2^{-n-1}
 \label{integration:eq:step-half-width}
\end{equation}
be the half-width of a cell and
\begin{equation}
 C_{n,m}=[x_{n,m}-h_n,x_{n,m}+h_n].
 \label{integration:eq:step-cell}
\end{equation}
Its full width is $2h_n=2^{-n}$.  Smearing the atom uniformly over its cell produces the
histogram
\begin{equation}
 \Wcal_n(x)=2^n\sum_m p_{n,m}\IndicatorOf{C_{n,m}}^{*}(x).
 \label{integration:eq:histogram-from-atoms}
\end{equation}
By \eqref{integration:eq:half-indicator-total-integral}, each term has integral $p_{n,m}$, hence
\begin{equation}
 \int_\RealNumbers\Wcal_n(x)\Differential x=1.
 \label{integration:eq:histogram-mass-one}
\end{equation}
The pointwise half weights ensure that when two cells share an endpoint, the two half
contributions add to the correct plateau value.

\subsubsection{The overlap sandwich}

For a fixed interval $[a,b]$ with $a\le b$, let
\begin{equation}
 \ell_{n,m}(a,b)=|[a,b]\cap C_{n,m}|
 =\max\bigl(\min(b,x_{n,m}+h_n)-\max(a,x_{n,m}-h_n),0\bigr).
 \label{integration:eq:cell-overlap-length}
\end{equation}
Then
\begin{equation}
 \int_a^b\Wcal_n(x)\Differential x
 =\sum_m p_{n,m}\,2^n\ell_{n,m}(a,b).
 \label{integration:eq:histogram-overlap-sum}
\end{equation}
The normalized overlap $2^n\ell_{n,m}$ lies in $[0,1]$.

If $x_{n,m}\in(a+h_n,b-h_n]$, then the entire cell lies inside $[a,b]$, so the normalized
overlap is $1$.  If the cell meets $[a,b]$, then its center lies in
$(a-h_n,b+h_n]$.  Therefore, term by term,
\begin{equation}
 \IndicatorOf{(a+h_n,b-h_n]}(x_{n,m})
 \le2^n\ell_{n,m}(a,b)
 \le\IndicatorOf{(a-h_n,b+h_n]}(x_{n,m}).
 \label{integration:eq:overlap-indicator-sandwich}
\end{equation}
Multiplication by $p_{n,m}$ and summation gives the central bridge.

\begin{theorem}[Atomic--histogram interval sandwich]
\label{integration:thm:atomic-histogram-sandwich}
For $a\le b$,
\begin{equation}
 \nu_n((a+h_n,b-h_n])
 \le\int_a^b\Wcal_n(x)\Differential x
 \le\nu_n((a-h_n,b+h_n]).
 \label{integration:eq:atomic-histogram-sandwich}
\end{equation}
\end{theorem}

This is the exact content of the inner- and outer-overlap lemmas in
\path{StepMeasureBridge.lean}.  Its importance is conceptual: it converts weak convergence
of measures into convergence of integrals of discontinuous histograms without claiming
pointwise convergence of densities from weak convergence alone.

\subsubsection{Portmanteau plus a moving-boundary squeeze}

The atomic laws $\nu_n$ coincide with the finite convolution laws and converge weakly to the
absolutely continuous law with density $\RvachevUp$.  Because that limit gives zero mass to every
singleton, the frontier of $(a,b]$ is null.  Portmanteau therefore gives
\begin{equation}
 \nu_n((a,b])\longrightarrow\int_a^b\RvachevUp(x)\Differential x.
 \label{integration:eq:portmanteau-Ioc}
\end{equation}
The intervals in \eqref{integration:eq:atomic-histogram-sandwich} move with $n$, so one more squeeze is
needed.  Fix $\varepsilon>0$ and choose $\delta>0$ with $4\delta<\varepsilon$.  Eventually
$h_n\le\delta$, and hence
\begin{align}
 (a+\delta,b-\delta]&\subset(a+h_n,b-h_n],                 \label{integration:eq:fixed-inner-subset}\\
 (a-h_n,b+h_n]&\subset(a-\delta,b+\delta].                 \label{integration:eq:moving-outer-subset}
\end{align}
The bound $0\le\RvachevUp\le1$ gives
\begin{align}
 0&\le\int_a^b\RvachevUp-\int_{a+\delta}^{b-\delta}\RvachevUp\le2\delta, \label{integration:eq:up-shrink-error}\\
 0&\le\int_{a-\delta}^{b+\delta}\RvachevUp-\int_a^b\RvachevUp\le2\delta. \label{integration:eq:up-expand-error}
\end{align}
Apply \eqref{integration:eq:portmanteau-Ioc} to the two fixed comparison intervals and combine with
\eqref{integration:eq:atomic-histogram-sandwich}.

\begin{theorem}[Interval-integral convergence of the step approximants]
\label{integration:thm:step-integral-convergence}
For every $a,b\in\RealNumbers$,
\begin{equation}
 \int_a^b\Wcal_n(x)\Differential x
 \longrightarrow
 \int_a^b\RvachevUp(x)\Differential x.
 \label{integration:eq:step-integral-convergence}
\end{equation}
\end{theorem}

\begin{proof}
The preceding fixed-$\delta$ sandwich proves the result for $a\le b$.  If $b<a$, reverse the
orientation of both interval integrals.
\end{proof}

\begin{warning}[The logical order cannot be shortened]
Weak convergence of $\nu_n$ does not by itself imply pointwise convergence of the
histograms $\Wcal_n$, nor even convergence of their values at a single point.  The formal
route is
\[
\begin{aligned}
 \text{weak convergence}
 &\Longrightarrow \text{null-boundary interval masses}\\
 &\Longrightarrow \text{moving-cell sandwich}
 \Longrightarrow \text{interval-integral convergence}.
\end{aligned}
\]
Other parts of the development then use shape, refinement, Fourier decay, or difference
quotients to obtain stronger function-level convergence.
\end{warning}

\section{The reusable discrete engines}

\subsection{Weighted paths solve triangular recurrences}
\label{integration:sec:path-sums}

The exact inverse-power recurrence for $F(2^{-n})$ is triangular.  The main article states its
composition formula, but the Lean proof first establishes a generic solver for every
semiring-valued triangular recurrence.  That abstraction explains both the form of the
answer and the uniqueness argument.

\subsubsection{Compositions as increasing paths}

A composition of $n$ is an ordered list of positive integers
\begin{equation}
 c=(r_1,\ldots,r_m),
 \qquad
 r_1+\cdots+r_m=n.
 \label{integration:eq:composition}
\end{equation}
Its partial sums
\begin{equation}
 s_0=0,
 \qquad
 s_j=r_1+\cdots+r_j
 \quad(1\le j\le m)
 \label{integration:eq:composition-partial-sums}
\end{equation}
form a strictly increasing path
\begin{equation}
 0=s_0<s_1<\cdots<s_m=n.
 \label{integration:eq:increasing-path}
\end{equation}
Conversely, every such path determines the positive increments $r_j=s_j-s_{j-1}$.

Let $R$ be a semiring and let
\begin{equation}
 w:\NaturalNumbers\times\NaturalNumbers\to R
 \label{integration:eq:generic-edge-weight}
\end{equation}
be an edge-weight function.  Define
\begin{align}
 W_w(c)&=\prod_{j=1}^{m}w(s_{j-1},s_j),                    \label{integration:eq:path-weight}\\
 P_w(n)&=\sum_{c\models n}W_w(c).                          \label{integration:eq:path-sum}
\end{align}
For $n=0$, there is one empty composition and its empty product is $1$, so
\begin{equation}
 P_w(0)=1.
 \label{integration:eq:path-sum-zero}
\end{equation}

\subsubsection{Last-edge decomposition}

Every nonempty path to $n$ has a unique last predecessor $k<n$.  Deleting its final edge
gives a path to $k$, and appending the edge $k\to n$ reverses the operation.  Therefore

\begin{theorem}[Last-edge decomposition]
\label{integration:thm:path-last-edge}
For $n>0$,
\begin{equation}
 P_w(n)=\sum_{k=0}^{n-1}P_w(k)w(k,n).
 \label{integration:eq:path-last-edge}
\end{equation}
\end{theorem}

\begin{proof}
Partition the finite set of compositions of $n$ by the sum $k$ of all blocks except the last.
For fixed $k$, deleting the final block is a bijection with compositions of $k$.  Path-weight
multiplicativity gives
\[
 W_w(c\mathbin{\|}(n-k))=W_w(c)w(k,n).
\]
Summing first over compositions of $k$ and then over $k<n$ gives
\eqref{integration:eq:path-last-edge}.
\end{proof}

One can instead partition by the number of edges.  Let $P_{w,r}(n)$ be the contribution of
compositions with exactly $r$ blocks.  Uniformly in $n$,
\begin{equation}
 P_w(n)=\sum_{r=0}^{n}P_{w,r}(n),
 \label{integration:eq:path-length-uniform}
\end{equation}
and for $n>0$ the zero-edge term vanishes, so the range may be written $1\le r\le n$.

\subsubsection{The generic triangular solver}

\begin{theorem}[Path-sum solution of a triangular recurrence]
\label{integration:thm:triangular-path-solver}
Let $x:\NaturalNumbers\to R$ satisfy
\begin{align}
 x_0&=b,                                                    \label{integration:eq:triangular-initial}\\
 x_n&=\sum_{k=0}^{n-1}x_k w(k,n)\qquad(n>0).               \label{integration:eq:triangular-recurrence}
\end{align}
Then
\begin{equation}
 x_n=bP_w(n)
 \qquad(n\ge0).
 \label{integration:eq:triangular-path-solution}
\end{equation}
Consequently the recurrence and initial value determine $x$ uniquely.
\end{theorem}

\begin{proof}
Use strong induction on $n$.  At $n=0$, \eqref{integration:eq:path-sum-zero} gives
$bP_w(0)=b=x_0$.  For $n>0$, insert the induction hypotheses into
\eqref{integration:eq:triangular-recurrence}:
\[
 x_n=\sum_{k<n}bP_w(k)w(k,n)
 =b\sum_{k<n}P_w(k)w(k,n)
 =bP_w(n),
\]
where the last equality is \cref{integration:thm:path-last-edge}.  Any two solutions therefore equal the
same path sum at each index.
\end{proof}

\begin{remark}
No subtraction, division, topology, or order is used.  The theorem works in an arbitrary
semiring.  This is why the formal development separates it from the rational Fabius
specialization.
\end{remark}

\subsection{Specialization to the inverse-dyadic recurrence}
\label{integration:sec:fabius-composition-formula}

The Fabius specialization is stated only once, as
\eqref{dyadicweb:eq:ordered-composition} in the exact dyadic part.  There it receives both
the infinite-product/geometric-simplex derivation and the independent recurrence-path
derivation obtained by specializing \cref{integration:thm:triangular-path-solver}.  Keeping
the generic solver here preserves the reusable semiring argument without repeating the
same inverse-dyadic theorem and its small numerical cases.

\subsection{A finite-row Toeplitz convergence theorem}
\label{integration:sec:toeplitz}

The complex-shift and discrete-limit formulas involve a row of coefficients depending on
$n$, not a fixed summability kernel.  The main article verifies the Fabius specialization but
omits the abstract theorem.  The needed result is a finite-row version of the Toeplitz
summability principle in a normed ring.

\subsubsection{Abstract statement}

Let $K$ be a normed ring.  For each $n$, let $J_n$ be a finite index set, let
$w_{n,j}\in K$, and let $q(n,j)\in\NaturalNumbers$ be the sampled sequence index.  Assume:
\begin{align}
 \sum_{j\in J_n}w_{n,j}&=1,                                \label{integration:eq:toeplitz-row-mass}\\
 \sum_{j\in J_n}\|w_{n,j}\|&\le C                         \label{integration:eq:toeplitz-variation}
\end{align}
for one constant $C\ge0$, and assume that the sampled indices move uniformly into every
tail:
\begin{equation}
 \forall N\ \exists n_0\ \forall n\ge n_0\ \forall j\in J_n,
 \qquad q(n,j)\ge N.
 \label{integration:eq:toeplitz-tail-index}
\end{equation}

\begin{theorem}[Finite-row Toeplitz convergence]
\label{integration:thm:finite-row-toeplitz}
If $H_m\to L$ in $K$, then
\begin{equation}
 \sum_{j\in J_n}w_{n,j}H_{q(n,j)}\longrightarrow L.
 \label{integration:eq:finite-row-toeplitz}
\end{equation}
\end{theorem}

\begin{proof}
Fix $\varepsilon>0$.  Since $H_m\to L$, choose $N$ such that
\begin{equation}
 \|H_m-L\|<\frac{\varepsilon}{C+1}
 \qquad(m\ge N).
 \label{integration:eq:toeplitz-tail-error}
\end{equation}
Choose $n_0$ from \eqref{integration:eq:toeplitz-tail-index}.  For $n\ge n_0$, use the row-mass identity
to subtract $L$:
\begin{align*}
 \sum_{j\in J_n}w_{n,j}H_{q(n,j)}-L
 &=\sum_{j\in J_n}w_{n,j}\bigl(H_{q(n,j)}-L\bigr).
\end{align*}
Therefore
\begin{align*}
 \left\|\sum_{j\in J_n}w_{n,j}H_{q(n,j)}-L\right\|
 &\le\sum_{j\in J_n}\|w_{n,j}\|\,\|H_{q(n,j)}-L\|\\
 &<\frac{\varepsilon}{C+1}\sum_{j\in J_n}\|w_{n,j}\|\\
 &\le\varepsilon\frac{C}{C+1}<\varepsilon.
\end{align*}
This proves convergence.
\end{proof}

\begin{remark}
The weights need not be positive, the rows need not be nested, and no entrywise limit of
$w_{n,j}$ is required.  The theorem uses exactly three facts: row mass one, uniformly bounded
total variation, and uniform migration of all sampled indices into the tail.
\end{remark}

\subsection{The Fabius specialization and its canonical exact treatment}
\label{integration:sec:fabius-toeplitz}
The abstract theorem above is the reusable normed-ring engine.  Its Fabius specialization
is developed only once in \cref{part:dyadicweb}: the half-base \(q\)-binomial rows have
mass one, uniformly bounded total variation, and sampled indices \(2n-j\) escaping every
finite prefix.  The shifted-spline definition there also records the indispensable
half-cell cutoff
\[
 \Floor{2^p x-\frac12}+1
 =\PositivePart{\Floor{2^p x+\frac12}},
\]
which converts an inclusive upper index into a natural range length without creating a
spurious summand below the first cell.  The exact-dyadic chapter then strengthens the old
coarse convergence statement to an explicit \(O_Q(4^{-n})\) estimate, exact stabilization
at dyadic inputs, and the synchronized \(q\)-Richardson limits.  Those sharper statements,
rather than a second copy of the row calculation, are the canonical specialization of
\cref{integration:thm:finite-row-toeplitz} in this synthesis.

\section{The formal and analytic machinery of the all-orders saddle expansion}

\subsection{Parameter-dependent Poincar\'e expansions}
\label{integration:sec:poincare-api}

The endpoint asymptotic is not an ordinary expansion with constant coefficients.  Its
coefficients retain a bounded one-periodic dependence on the exact Lambert phase.  The
formal library therefore uses an asymptotic API in which coefficient families may depend on
the asymptotic parameter, provided that dependence remains bounded.

Let $\ell$ be a filter on a parameter space $A$, let $\sigma:A\to\mathbb K$ be a scale tending
to zero, and let $f:A\to E$ take values in a normed vector space over the normed field
$\mathbb K$.  For coefficient families $c_k:A\to E$, put
\begin{equation}
 S_N[f](a)=\sum_{k=0}^{N-1}\sigma(a)^k c_k(a).
 \label{integration:eq:parameter-partial-sum}
\end{equation}

\begin{definition}[Full parameter-dependent expansion]
\label{integration:def:parameter-expansion}
We write
\begin{equation}
 f(a)\sim_{\ell}\sum_{k\ge0}\sigma(a)^kc_k(a)
 \label{integration:eq:parameter-expansion-symbol}
\end{equation}
when
\begin{enumerate}
 \item every $c_k$ is bounded along $\ell$, equivalently $c_k=O_{\ell}(1)$;
 \item for every $N\in\NaturalNumbers$,
 \begin{equation}
  f-S_N[f]=O_{\ell}(\sigma^N).
  \label{integration:eq:parameter-expansion-remainder}
 \end{equation}
\end{enumerate}
\end{definition}

This is \lean{HasAsymptoticExpansion}.  Boundedness of the coefficients is not redundant: if
arbitrary unbounded parameter dependence were allowed, powers of $\sigma$ could be moved
between the ``coefficient'' and the scale, destroying the meaning of order.

\subsubsection{Elementary algebra of full expansions}

Suppose $f$ and $g$ have expansions with the same scale.  Direct manipulation of finite
partial sums proves:
\begin{align}
 f+g&\sim\sum_{k\ge0}\sigma^k(c_k+d_k),                    \label{integration:eq:expansion-add}\\
 f-g&\sim\sum_{k\ge0}\sigma^k(c_k-d_k),                    \label{integration:eq:expansion-sub}\\
 \alpha f&\sim\sum_{k\ge0}\sigma^k(\alpha c_k).            \label{integration:eq:expansion-smul}
\end{align}
The API also supports transport across eventual equality and pullback along a map tending to
the source filter.  These apparently routine lemmas are what let the final proof switch
between complex and real saddle ratios, replace exact identities after a threshold, and move
from the logarithmic coordinate to $x\downarrow0$ without reopening every remainder
estimate.

\subsubsection{Flat functions}

\begin{definition}[Flat relative to a scale]
\label{integration:def:flat-scale}
A function $r$ is flat relative to $\sigma$ along $\ell$ if
\begin{equation}
 r=O_{\ell}(\sigma^N)
 \qquad\text{for every }N\in\NaturalNumbers.
 \label{integration:eq:flat-all-orders}
\end{equation}
\end{definition}

\begin{theorem}[Flat remainders are invisible to all coefficients]
\label{integration:thm:flat-invisible}
A flat $r$ has the full expansion with every coefficient zero.  Consequently, if
\begin{equation}
 f\sim\sum_{k\ge0}\sigma^kc_k,
 \label{integration:eq:f-exp-before-flat}
\end{equation}
then
\begin{equation}
 f+r\sim\sum_{k\ge0}\sigma^kc_k,
 \qquad
 f-r\sim\sum_{k\ge0}\sigma^kc_k.
 \label{integration:eq:flat-preserves-expansion}
\end{equation}
The implications hold in both directions.
\end{theorem}

\begin{proof}
For the zero-coefficient expansion, the $N$th partial sum is identically zero, so the
remainder condition is exactly \eqref{integration:eq:flat-all-orders}.  Add or subtract this expansion
coefficientwise from the expansion of $f$.  Reversibility follows by moving $r$ to the other
side.
\end{proof}

This theorem is the formal reason that the forward negative-Laplace tail can be restored
after the central saddle expansion without changing any algebraic-order coefficient.

\subsubsection{Exact logarithmic coordinates}

Define
\begin{equation}
 x(t)=2^{-t},
 \qquad
 t(x)=-\log_2x
 \quad(x>0).
 \label{integration:eq:small-log-coordinates}
\end{equation}
These maps are exact inverses on $x>0$, and
\begin{equation}
 t\to+\infty\iff x(t)\downarrow0,
 \qquad
 x\downarrow0\iff t(x)\to+\infty.
 \label{integration:eq:coordinate-filter-equivalence}
\end{equation}

\begin{theorem}[Full-expansion coordinate equivalence]
\label{integration:thm:small-log-expansion-equivalence}
For functions $f,\sigma,c_k$ on the positive half-line,
\begin{equation}
 f(2^{-t})\sim_{t\to\infty}
 \sum_{k\ge0}\sigma(2^{-t})^kc_k(2^{-t})
 \label{integration:eq:log-coordinate-expansion}
\end{equation}
if and only if
\begin{equation}
 f(x)\sim_{x\downarrow0}
 \sum_{k\ge0}\sigma(x)^kc_k(x).
 \label{integration:eq:small-coordinate-expansion}
\end{equation}
\end{theorem}

\begin{proof}
Pull back the full expansion along either coordinate map.  The two composites are eventually
identical on their respective filters by \eqref{integration:eq:small-log-coordinates}; boundedness and
every remainder estimate therefore transfer in both directions.
\end{proof}

\subsection{Recursive coefficients of a formal exponential}
\label{integration:sec:exp-coeff}

The saddle integrand contains the exponential of a formal perturbation.  Rather than invoke
Bell polynomials or repeatedly expand a universal power series, the Lean development uses a
finite recurrence whose $n$th output depends only on the first $n$ input coefficients.

Let $R$ be a commutative $\RationalNumbers$-algebra, and let
\begin{equation}
 E(u)=\sum_{j\ge1}E_j u^j
 \label{integration:eq:formal-exponent}
\end{equation}
be a formal exponent with zero constant term.  Define $a_n\in R$ recursively by
\begin{align}
 a_0&=1,                                                    \label{integration:eq:expcoeff-zero}\\
 a_{n+1}&=\frac1{n+1}\sum_{j=0}^{n}(j+1)E_{j+1}a_{n-j}.
                                                               \label{integration:eq:expcoeff-recurrence}
\end{align}

\begin{theorem}[Correctness and uniqueness of the exponential recurrence]
\label{integration:thm:expcoeff-correctness}
The power series
\begin{equation}
 A(u)=\sum_{n\ge0}a_nu^n
 \label{integration:eq:formal-mass-series}
\end{equation}
satisfies
\begin{equation}
 A'(u)=E'(u)A(u),
 \qquad
 A(0)=1.
 \label{integration:eq:formal-exp-ode}
\end{equation}
It is the unique formal power series with these two properties, and therefore
\begin{equation}
 A(u)=\exp(E(u)).
 \label{integration:eq:formal-exp-identification}
\end{equation}
\end{theorem}

\begin{proof}
The coefficient of $u^n$ in $A'$ is $(n+1)a_{n+1}$.  The coefficient of $u^n$ in $E'A$ is
\[
 \sum_{j=0}^{n}(j+1)E_{j+1}a_{n-j},
\]
so \eqref{integration:eq:formal-exp-ode} is coefficientwise equivalent to
\eqref{integration:eq:expcoeff-recurrence}.  Conversely the differential equation and initial value
determine $a_0$, then $a_1$, then $a_2$, and so on, proving uniqueness.  The universal
formal exponential satisfies the same equation and initial value.
\end{proof}

The first coefficients are
\begin{align}
 a_1={}&E_1,                                                  \label{integration:eq:a1}\\
 a_2={}&E_2+\frac12E_1^2,                                    \label{integration:eq:a2}\\
 a_3={}&E_3+E_1E_2+\frac16E_1^3,                             \label{integration:eq:a3}\\
 a_4={}&E_4+E_1E_3+\frac12E_2^2
       +\frac12E_1^2E_2+\frac1{24}E_1^4.                    \label{integration:eq:a4}
\end{align}
Equation \eqref{integration:eq:a4} is the exact finite identity used for the second saddle correction.

\subsubsection{Structural laws}

The recurrence is preferable to a black-box exponential because it exposes several laws
needed downstream.

\begin{proposition}[Finite dependence]
\label{integration:prop:expcoeff-finite-dependence}
The coefficient $a_n$ depends only on $E_1,\ldots,E_n$.  If two exponent families agree
through order $n$, their $n$th exponential coefficients agree.
\end{proposition}

\begin{proposition}[Naturality]
\label{integration:prop:expcoeff-naturality}
For every morphism of commutative $\RationalNumbers$-algebras $\varphi:R\to S$,
\begin{equation}
 \varphi(a_n(E))=a_n(\varphi\circ E).
 \label{integration:eq:expcoeff-map}
\end{equation}
In particular, evaluation of polynomial- or function-valued coefficients commutes with the
recurrence.
\end{proposition}

\begin{proposition}[Rescaling, sums, and parity]
\label{integration:prop:expcoeff-structure}
Let $c\in R$ be central.
\begin{enumerate}
 \item If $\widetilde E_j=c^jE_j$, then
 \begin{equation}
  a_n(\widetilde E)=c^na_n(E).
  \label{integration:eq:expcoeff-rescale}
 \end{equation}
 \item If $E$ and $F$ are two exponent families, then
 \begin{equation}
  a_n(E+F)=\sum_{r+s=n}a_r(E)a_s(F).
  \label{integration:eq:expcoeff-add-convolution}
 \end{equation}
 \item If $E_j(-v)=(-1)^jE_j(v)$, then
 \begin{equation}
  a_n(-v)=(-1)^na_n(v).
  \label{integration:eq:expcoeff-parity}
 \end{equation}
\end{enumerate}
\end{proposition}

\begin{proofidea}
Naturality is induction through finite sums and rational scalar multiplication.
Rescaling follows from the recurrence because every summand has total degree
$(j+1)+(n-j)=n+1$.  The addition law is the coefficient form of
$\exp(E+F)=\exp(E)\exp(F)$.  Parity is the specialization $c=-1$ of rescaling after
pointwise reflection in $v$.
\end{proofidea}

\subsection{Recursive coefficients of a formal logarithm}
\label{integration:sec:log-coeff}

Let
\begin{equation}
 A(u)=1+\sum_{n\ge1}a_nu^n.
 \label{integration:eq:unit-mass-series}
\end{equation}
We seek
\begin{equation}
 B(u)=\log A(u)=\sum_{n\ge1}b_nu^n.
 \label{integration:eq:formal-log-series}
\end{equation}
Instead of expanding the universal logarithm, use
\begin{equation}
 A(u)B'(u)=A'(u).
 \label{integration:eq:formal-log-ode}
\end{equation}
Solving the coefficient of $u^n$ for $b_{n+1}$ gives
\begin{align}
 b_0&=0,                                                     \label{integration:eq:logcoeff-zero}\\
 b_{n+1}&=a_{n+1}-\frac1{n+1}
 \sum_{j=0}^{n-1}(n-j)b_{n-j}a_{j+1}.                      \label{integration:eq:logcoeff-recurrence}
\end{align}
The empty sum at $n=0$ gives $b_1=a_1$.

\begin{theorem}[Correctness and uniqueness of the logarithm recurrence]
\label{integration:thm:logcoeff-correctness}
If $a_0=1$, the series generated by \eqref{integration:eq:logcoeff-recurrence} is the unique series with
zero constant term satisfying \eqref{integration:eq:formal-log-ode}; hence it is the formal logarithm of
$A$.
\end{theorem}

\begin{proof}
The coefficient of $u^n$ in $AB'$ contains the term $(n+1)b_{n+1}$ from $a_0b'$, plus terms
involving only $b_1,\ldots,b_n$.  Equating it with $(n+1)a_{n+1}$ and dividing by $n+1$
gives \eqref{integration:eq:logcoeff-recurrence}.  This triangular relation proves uniqueness.  The
formal logarithm satisfies $B(0)=0$ and $B'A=A'$.
\end{proof}

The first coefficients are
\begin{align}
 b_1={}&a_1,                                                  \label{integration:eq:b1}\\
 b_2={}&a_2-\frac12a_1^2,                                   \label{integration:eq:b2}\\
 b_3={}&a_3-a_1a_2+\frac13a_1^3,                            \label{integration:eq:b3}\\
 b_4={}&a_4-a_1a_3-\frac12a_2^2+a_1^2a_2-\frac14a_1^4.     \label{integration:eq:b4}
\end{align}
As with the exponential recurrence, $b_n$ depends only on $a_1,\ldots,a_n$, commutes with
$\RationalNumbers$-algebra morphisms and evaluation, obeys the rescaling law
\begin{equation}
 a_j\mapsto c^ja_j\quad\Longrightarrow\quad b_n\mapsto c^nb_n,
 \label{integration:eq:logcoeff-rescale}
\end{equation}
and preserves parity.

\begin{remark}[Why $a_0$ is a hypothesis, not an input]
The recursive definition \eqref{integration:eq:logcoeff-recurrence} never reads $a_0$; it is meaningful
for every sequence.  Correctness as a logarithm, however, requires $a_0=1$.  Separating the
recurrence from its normalization makes finite coefficient calculations easier while keeping
the theorem honest.
\end{remark}

\subsubsection{From mass expansions to logarithmic expansions}

Suppose a positive real function has a full expansion
\begin{equation}
 R(a)\sim1+\sum_{n\ge1}\sigma(a)^na_n(a),
 \qquad \sigma(a)\to0,
 \label{integration:eq:mass-expansion-before-log}
\end{equation}
with bounded coefficient families.  Finite polynomial remainders, discussed in the next
section, show that for each $N$,
\begin{equation}
 \log R(a)=\sum_{n=1}^{N-1}\sigma(a)^nb_n(a)+O(\sigma(a)^N),
 \label{integration:eq:log-asymptotic-transfer}
\end{equation}
where $b_n$ is generated by \eqref{integration:eq:logcoeff-recurrence}.  The positivity hypothesis fixes
the real branch and is eventually automatic when $R\to1$.

\subsection{Exact finite remainders from formal coefficient identities}
\label{integration:sec:finite-remainder}

Formal power-series equality says that two expressions have the same coefficients below a
chosen order.  An asymptotic proof needs an evaluated identity with an explicit factor of the
first omitted power.  The module \path{SaddleExpansionFiniteRemainder.lean} provides exactly
that bridge without invoking an infinite series.

Fix $L\ge1$ and truncate the exponent:
\begin{equation}
 E_{<L}(X)=\sum_{k=1}^{L-1}E_kX^k.
 \label{integration:eq:exponent-trunc-polynomial}
\end{equation}
Truncate the scalar exponential polynomial at the same order and substitute:
\begin{equation}
 T_L(X)=\sum_{q=0}^{L-1}\frac{E_{<L}(X)^q}{q!}.
 \label{integration:eq:finite-exp-substitution}
\end{equation}
Let
\begin{equation}
 A_{<L}(X)=\sum_{k=0}^{L-1}a_kX^k
 \label{integration:eq:expcoeff-trunc-polynomial}
\end{equation}
be the polynomial formed from the recursive exponential coefficients.

\begin{theorem}[Exact divisibility of the finite defect]
\label{integration:thm:finite-defect-divisibility}
If $E_0=0$, then
\begin{equation}
 X^L\mid\bigl(T_L(X)-A_{<L}(X)\bigr).
 \label{integration:eq:defect-divisibility}
\end{equation}
Equivalently, there is a canonical polynomial $Q_L(X)$ such that
\begin{equation}
 T_L(X)=A_{<L}(X)+X^LQ_L(X).
 \label{integration:eq:exact-finite-remainder-polynomial}
\end{equation}
\end{theorem}

\begin{proof}
The coefficient of $X^k$ in $T_L$ agrees with that of the full formal substitution
$\exp(E(X))$ for every $k<L$: terms of scalar exponential degree $q\ge L$ cannot contribute
below degree $L$ because $E$ has zero constant term, and truncating $E$ itself does not alter
coefficients below $L$.  By \cref{integration:thm:expcoeff-correctness}, that common coefficient is
$a_k$.  Hence the first $L$ coefficients of the defect vanish, which is equivalent to
divisibility by $X^L$.  The formal construction takes $Q_L$ to be the result of applying
polynomial division by $X$ exactly $L$ times.
\end{proof}

Evaluating at $x\in R$ gives the identity actually used in estimates:
\begin{equation}
 \sum_{q=0}^{L-1}\frac1{q!}
 \left(\sum_{k=1}^{L-1}E_kx^k\right)^q
 =\sum_{k=0}^{L-1}a_kx^k+x^LQ_L(x).
 \label{integration:eq:finite-remainder-evaluated}
\end{equation}

\begin{boxedremark}
The quotient $Q_L$ is a polynomial in $x$ whose coefficients are algebraic expressions in
the finitely many parameters $E_1,\ldots,E_{L-1}$.  If those parameter families are uniformly
bounded, then $Q_L$ is uniformly bounded on a fixed compact $x$-window.  Thus the exact
factorization \eqref{integration:eq:finite-remainder-evaluated} turns formal coefficient agreement into a
uniform $O(x^L)$ estimate with no appeal to convergence of an infinite formal series.
\end{boxedremark}

\subsection{Gaussian contraction: integration becomes polynomial algebra}
\label{integration:sec:gaussian-contraction}

After the saddle variable is scaled by $v=\sqrt{\lambda}\,\theta$, the leading kernel is the
standard Gaussian
\begin{equation}
 \gamma(v)=\EulerE^{-v^2/2}.
 \label{integration:eq:standard-gaussian-kernel}
\end{equation}
Every finite reference term is $\gamma(v)$ times a complex polynomial.  The generic
contraction module turns its whole-line integral into an algebraic linear functional.

\subsubsection{Gaussian moments}

Define the unnormalized moments
\begin{equation}
 M_n^{\mathrm G}=\int_{\RealNumbers}\EulerE^{-v^2/2}v^n\Differential v.
 \label{integration:eq:real-gaussian-moment}
\end{equation}
Every moment is absolutely integrable.  Symmetry gives
\begin{equation}
 M_{2j+1}^{\mathrm G}=0.
 \label{integration:eq:gaussian-odd-zero}
\end{equation}
Integration by parts, with $u=v^{n+1}$ and
$\Differential w=-v\EulerE^{-v^2/2}\Differential v$, gives
\begin{equation}
 M_{n+2}^{\mathrm G}=(n+1)M_n^{\mathrm G}.
 \label{integration:eq:gaussian-moment-recurrence}
\end{equation}
Since
\begin{equation}
 M_0^{\mathrm G}=\sqrt{2\pi},
 \label{integration:eq:gaussian-moment-zero}
\end{equation}
the normalized moments
\begin{equation}
 \mu_n=\frac{M_n^{\mathrm G}}{\sqrt{2\pi}}
 \label{integration:eq:normalized-gaussian-moment}
\end{equation}
satisfy
\begin{equation}
 \mu_0=1,
 \quad
 \mu_1=0,
 \quad
 \mu_{n+2}=(n+1)\mu_n.
 \label{integration:eq:normalized-gaussian-recurrence}
\end{equation}
Thus
\begin{equation}
 \mu_{2j}=(2j-1)!!,
 \qquad
 \mu_{2j+1}=0.
 \label{integration:eq:normalized-gaussian-closed}
\end{equation}
The first even values needed below are
\begin{equation}
 \mu_4=3,
 \quad
 \mu_6=15,
 \quad
 \mu_8=105,
 \quad
 \mu_{10}=945,
 \quad
 \mu_{12}=10395.
 \label{integration:eq:gaussian-moments-second-correction}
\end{equation}

\subsubsection{The contraction functional}

For $p(X)=\sum_kp_kX^k\in\ComplexNumbers[X]$, define
\begin{equation}
 \Gcal[p]=\sum_kp_k\mu_k.
 \label{integration:eq:gaussian-contraction-def}
\end{equation}
This is a $\ComplexNumbers$-linear map $\ComplexNumbers[X]\to\ComplexNumbers$.

\begin{theorem}[Gaussian polynomial contraction]
\label{integration:thm:gaussian-contraction}
For every complex polynomial $p$,
\begin{equation}
 \int_{\RealNumbers}\EulerE^{-v^2/2}p(v)\Differential v
 =\sqrt{2\pi}\,\Gcal[p].
 \label{integration:eq:gaussian-contraction-integral}
\end{equation}
In particular
\begin{equation}
 \Gcal[X^{2j}]=(2j-1)!!,
 \qquad
 \Gcal[X^{2j+1}]=0.
 \label{integration:eq:gaussian-contraction-monomials}
\end{equation}
\end{theorem}

\begin{proof}
Expand $p$ as a finite sum of monomials, interchange that finite sum with the integral, and
use \eqref{integration:eq:normalized-gaussian-moment}.  Integrability follows coefficientwise from the
absolute Gaussian moments.
\end{proof}

The same coefficientwise expansion yields the norm estimate
\begin{equation}
 \left|\Gcal[p]\right|
 \le\sum_k|p_k|\mu_{2\Ceiling{k/2}},
 \label{integration:eq:gaussian-contraction-norm-bound}
\end{equation}
with odd coefficients contributing zero in the exact functional.  In the formal proof a
slightly more uniform support sum is used so it composes smoothly with polynomial families.

\begin{remark}[Parity is an algebraic annihilator]
The vanishing of odd Gaussian moments is not merely a simplification after integration.  It
is the mechanism that removes every odd power of the small parameter $\varepsilon=
\lambda^{-1/2}$ from the mass expansion.  The exponent coefficient of order $m$ has parity
$m$ in $v$, the formal exponential recurrence preserves total parity, and
\eqref{integration:eq:gaussian-contraction-monomials} kills every odd result.  The final expansion is
therefore in integer powers of $\lambda^{-1}$.
\end{remark}

\subsection{Uniform Gaussian tails for polynomial references}
\label{integration:sec:gaussian-tails}

Whole-line contraction is useful only after the central-window integral has been replaced by
a whole-line integral with a negligible error.  The omitted Gaussian-tail modules provide a
single coefficient weight that controls every polynomial reference.

\subsubsection{A factorial coefficient weight}

For $p(X)=\sum_kp_kX^k$, define
\begin{equation}
 \mathfrak W(p)=\sum_{k\in\SupportOperator p}|p_k|\,8k!.
 \label{integration:eq:gaussian-tail-weight}
\end{equation}
It is finite and nonnegative.

\begin{lemma}[Monomial Gaussian tail]
\label{integration:lem:monomial-gaussian-tail}
For $A\ge4$ and $k\in\NaturalNumbers$,
\begin{equation}
 \int_{|v|>A}\EulerE^{-v^2/2}|v|^k\Differential v
 \le 8k!\,\frac{\EulerE^{-A^2/4}}{A}.
 \label{integration:eq:monomial-gaussian-tail}
\end{equation}
\end{lemma}

\begin{proofidea}
Split the kernel as
$\EulerE^{-v^2/2}=\EulerE^{-v^2/4}\EulerE^{-v^2/4}$.  The factor
$|v|^k\EulerE^{-v^2/4}$ is controlled uniformly by a factorial bound, while the remaining
Gaussian tail is estimated by integrating the derivative of $\EulerE^{-v^2/4}$ and using
$|v|\ge A$.  The constants are chosen so the result holds simultaneously for every
$k\ge0$ once $A\ge4$; no optimization is needed in the saddle application.
\end{proofidea}

\begin{theorem}[Polynomial Gaussian tail]
\label{integration:thm:polynomial-gaussian-tail}
For every $p\in\ComplexNumbers[X]$ and $A\ge4$,
\begin{equation}
 \int_{|v|>A}\left|\EulerE^{-v^2/2}p(v)\right|\Differential v
 \le \mathfrak W(p)\frac{\EulerE^{-A^2/4}}{A}.
 \label{integration:eq:polynomial-gaussian-tail}
\end{equation}
\end{theorem}

\begin{proof}
Use
$|p(v)|\le\sum_{k\in\SupportOperator p}|p_k||v|^k$, integrate the finite sum, and apply
\cref{integration:lem:monomial-gaussian-tail} coefficientwise.
\end{proof}

\subsubsection{The all-orders central radius}

For expansion order $N$ and scale parameter $b\ge1$, set
\begin{equation}
 A_N(b)=\sqrt{32(N+1)\log b}.
 \label{integration:eq:order-central-radius}
\end{equation}
Then
\begin{equation}
 \EulerE^{-A_N(b)^2/4}=b^{-8(N+1)}.
 \label{integration:eq:order-radius-exponential}
\end{equation}
The exponent $8(N+1)$ is much stronger than the requested algebraic order $N+1$, leaving
ample room for the factor $1/A_N(b)$ and for bounded coefficient weights.

\begin{theorem}[All-orders polynomial tail]
\label{integration:thm:all-orders-polynomial-tail}
Let $b(a)\to+\infty$ along a filter and let $p_a\in\ComplexNumbers[X]$ satisfy
\begin{equation}
 \mathfrak W(p_a)=O(1).
 \label{integration:eq:bounded-polynomial-tail-weight}
\end{equation}
Then for every fixed $N$,
\begin{equation}
 \int_{|v|>A_N(b(a))}
 \left|\EulerE^{-v^2/2}p_a(v)\right|\Differential v
 =O\bigl(b(a)^{-(N+1)}\bigr).
 \label{integration:eq:all-orders-polynomial-tail}
\end{equation}
\end{theorem}

\begin{proof}
Eventually $A_N(b(a))\ge4$ and $b(a)\ge1$.  Apply
\cref{integration:thm:polynomial-gaussian-tail}, substitute
\eqref{integration:eq:order-radius-exponential}, and use
\[
 b^{-8(N+1)}A_N(b)^{-1}\le b^{-(N+1)}.
\]
Multiplication by the bounded family \eqref{integration:eq:bounded-polynomial-tail-weight} preserves the
big-$O$ rate.
\end{proof}

\begin{remark}
The theorem is formulated for a moving polynomial family.  This is essential: the saddle
reference polynomial depends periodically on the exact phase.  Periodicity and smoothness
make its finitely many coefficient functions bounded, hence its factorial weight is bounded.
\end{remark}

\subsection{Handoff to the Fabius-specific saddle synthesis}
\label{integration:sec:saddle-jets}
\label{integration:sec:saddle-coeff-engine}
\label{integration:sec:second-correction}
\label{integration:sec:full-assembly}
\label{integration:sec:exact-phase-boundary}
The preceding Poincar\'e, formal exponential/logarithm, finite-remainder, and Gaussian
modules are deliberately generic.  Their Fabius-specific instantiation is assembled once
in \cref{part:inverse}, after the exact lower-Lambert phase and periodic negative-Laplace
decomposition of \cref{part:dyadicasym}.  There the closed shifted-Stirling jets feed the
two-monomial exponent coefficients, the formal exponential produces the saddle-mass
coefficients, Gaussian contraction removes the local variable, and the recursive logarithm
produces the periodic corrections.  The same chapter gives the explicit first and second
corrections and the rigorous all-orders central/minor-arc/flat-tail assembly.

This separation also fixes the claim boundary.  The mass top-jet theorem is formalized;
the transfer to the logarithmic highest-\(\Psi\)-derivative coefficient is retained as a
derived statement.  Likewise, the expansion is stated in the exact phase: replacing that
phase by a truncated elementary approximation inside periodic coefficients requires a
separate composition theorem.  The full formulas, status ledger, and remaining obligations
are therefore kept with the canonical application rather than repeated here.

\clearpage
\subsection*{Supplementary material from this synthesis}
\addcontentsline{toc}{section}{Supplementary material from this synthesis}
\subsection{Executable recurrences in mathematical pseudocode}
\label{integration:app:algorithms}

The following algorithms are not substitutes for the proofs above.  They are compact
implementations of the finite algebra that appears repeatedly in the formal development.
All arrays are zero-indexed.

\begin{algorithm}[Weighted path sum]
\label{integration:alg:path-sum}
The dynamic program below computes $P_w(n)$ without enumerating compositions.  It is the
triangular recurrence proved equivalent to the composition sum in
\cref{integration:thm:triangular-path-solver}.
\end{algorithm}
\begin{lstlisting}[style=pseudo]
function PathSums(N, weight):
    P[0] := 1
    for n := 1 to N:
        P[n] := 0
        for k := 0 to n - 1:
            P[n] := P[n] + P[k] * weight(k, n)
    return P

function FabiusInverseDyadic(N):
    weight(k, n) := 1 / ((2^n - 1) * factorial(n - k + 1))
    a := PathSums(N, weight)
    for n := 0 to N:
        value[n] := 2^(-binomial(n, 2)) * a[n]
    return value
\end{lstlisting}

\begin{algorithm}[Formal exponential coefficients]
\label{integration:alg:expcoeff}
Given $E[1],\ldots,E[N]$ in a commutative $\RationalNumbers$-algebra, this returns the first $N$
coefficients of $\exp(\sum_{j\ge1}E[j]u^j)$.
\end{algorithm}
\begin{lstlisting}[style=pseudo]
function ExpCoeff(E, N):
    a[0] := 1
    for n := 0 to N - 1:
        s := 0
        for j := 0 to n:
            s := s + (j + 1) * E[j + 1] * a[n - j]
        a[n + 1] := s / (n + 1)
    return a
\end{lstlisting}

\begin{algorithm}[Formal logarithm coefficients]
\label{integration:alg:logcoeff}
The input is a unit series $1+\sum_{j\ge1}a[j]u^j$.  The output satisfies
$\log(1+\sum a[j]u^j)=\sum b[j]u^j$.
\end{algorithm}
\begin{lstlisting}[style=pseudo]
function LogCoeff(a, N):
    b[0] := 0
    for n := 0 to N - 1:
        s := 0
        for j := 0 to n - 1:
            s := s + (n - j) * b[n - j] * a[j + 1]
        b[n + 1] := a[n + 1] - s / (n + 1)
    return b
\end{lstlisting}

\begin{algorithm}[Normalized Gaussian contraction]
\label{integration:alg:gaussian-contraction}
The moment table can be built once to the degree of the input polynomial.
\end{algorithm}
\begin{lstlisting}[style=pseudo]
function GaussianMoments(D):
    mu[0] := 1
    if D >= 1: mu[1] := 0
    for n := 0 to D - 2:
        mu[n + 2] := (n + 1) * mu[n]
    return mu

function GaussianContract(coeff, D):
    mu := GaussianMoments(D)
    result := 0
    for k := 0 to D:
        result := result + coeff[k] * mu[k]
    return result
\end{lstlisting}

\begin{algorithm}[Saddle coefficient pipeline]
\label{integration:alg:saddle-pipeline}
This is the finite algebra at fixed logarithmic order $N$.  The analytic proof is what
justifies applying it to the central integral with a uniform remainder.
\end{algorithm}
\begin{lstlisting}[style=pseudo]
function SaddleCoefficients(N, phase):
    # 1. Build bounded exponent jets Z[0 .. 2N - 1].
    Z := ClosedJets(2*N - 1, phase)

    # 2. Build exponent polynomials E[1 .. 2N].
    for m := 1 to 2*N:
        E[m] := i^m * (Z[m - 1] * X^m / factorial(m)
                       + (-1)^(m + 1) * X^(m + 2) / (m + 2))

    # 3. Exponentiate formally in epsilon.
    a := ExpCoeff(E, 2*N)

    # 4. Contract even epsilon orders against the normalized Gaussian.
    B[0] := 1
    for j := 1 to N:
        B[j] := GaussianContract(coefficients(a[2*j]), degree(a[2*j]))

    # 5. Convert mass coefficients to logarithmic coefficients.
    A := LogCoeff(B, N)
    return (B, A)
\end{lstlisting}

\subsection{Lean declaration and module concordance}
\label{integration:app:concordance}

File names below are relative to
\begin{center}
 \repo/\texttt{Lean/FabiusFunction/}.
\end{center}
The table lists the principal declarations, not every supporting lemma.

\begingroup\small
\begin{longtable}{@{}>{\raggedright\arraybackslash\scriptsize}p{0.36\textwidth}
                         >{\raggedright\arraybackslash}p{0.57\textwidth}@{}}
\toprule
Module & Content explained in this companion\\
\midrule
\endhead
\path{OriginalCharacterization.lean} &
\lean{IsOriginalFabius.mk\_of\_derivative\_law}: the support/mean-value argument showing
that scale positivity follows from the other original-characterization assumptions;
\cref{integration:sec:scale-positive}.\\
\path{FabiusInverse.lean} &
The complete order, calculus, curvature, exact-locus, and endpoint-steepness interface is
consolidated in \cref{part:inverse}; \cref{integration:sec:inverse} is the handoff.\\
\path{ElementaryFunction.lean} &
The inductive predicate \lean{IsElementary}, composition closure, derived elementary
operations, analytic loci, the finite-bad-value composition theorem, and the dense-open
analytic-locus theorem; \cref{integration:sec:elementary}.\\
\path{NotElementary.lean} &
No elementary function agrees with a bounded Fabius function, $\RvachevUp$, or the signed extension
on a set with nonempty interior; \cref{integration:sec:not-elementary}.\\
\path{InverseNotElementary.lean} &
Exact analytic locus of the total inverse and exclusion from both the elementary class and
the class enlarged by admissible inverse branches; \cref{integration:sec:not-elementary}.\\
\path{ProbabilityRepresentation.lean} &
The product sample space, uniformly convergent coordinate sum, pushforward law, the exact
dyadic-to-generic bridge
\lean{ProbabilityRepresentation.weightedSumDistribution_eq_weightedUniformDistribution},
the two half-base geometric bridges, head--tail split, coordinate reflection,
the exact support theorem \lean{ProbabilityRepresentation.weightedSumDistribution_support_eq_Icc},
and identification of the CDF with $F$;
\cref{integration:sec:product-law}.\\
\path{WeightedUniformSeries.lean} &
The eleven-declaration generic tranche: almost-everywhere and pointwise head--tail map
recursion, continuous-linear and scalar naturality, shifted-law transport, and
operator/scalar affine recursions; \cref{integration:sec:weighted-naturality}.\\
\path{WeightedUniformDistribution.lean} &
The exact zero-definition/eleven-theorem compatibility surface: head--tail, probability,
the half-total-weight barycenter, scalar-law and unit-interval results, followed by the
four-declaration s-finite uniform-smoothing, sharp weighted-law absolute-continuity, and
null-singleton tranche;
\cref{integration:sec:weighted-naturality}.\\
\path{WeightedUniformSupport.lean} &
The ten-declaration exact-support tranche: open positivity for arbitrary products,
continuous-pushforward support, support equal to the series range, exact signed-real
min/max intervals, and the nonnegative and unit-mass specializations;
\cref{integration:sec:weighted-naturality}.\\
\path{GeometricUniformLaw.lean} &
The twenty-four-declaration named geometric weight, series, and law API: summability and
mass, continuity and measurability, affine splitting, reflection, exact mean \(1/2\),
absolute continuity, null singletons, and support equal to the series range for \(|q|<1\),
including zero and negative ratios, with exact support \([0,1]\)
for \(0\le q<1\);
\cref{integration:sec:weighted-naturality}.\\
\path{AffineIndependentCopy.lean} &
The one-definition/eight-theorem generic affine-law tranche: probability preservation,
convolution/product-map equality, characteristic-function factorization and finite
iteration, a recurrence-level uniqueness engine, and operator/product-map fixed-point
uniqueness for \(|q|<1\) in complete second-countable real inner-product targets;
\cref{integration:thm:affine-independent-copy-unique}.\\
\path{GeometricUniformUniqueness.lean} &
Uniqueness of the named geometric law among all probability solutions of its affine
independent-copy equation for \(|q|<1\), including zero and negative parameters and without
support, density, or moment hypotheses;
\cref{integration:eq:geometric-law-unique}.\\
\path{GeometricSincFactorization.lean} &
The unconditional closed scaled-digit characteristic function and, for \(|q|<1\), the
closed-factor and phase-collapsed finite residual formulas, pointwise phase-bearing sinc
prefix convergence, and dyadic weighted-sum wrappers;
\cref{integration:eq:geometric-sinc-residual}.\\
\path{GeometricReciprocalGamma.lean} &
Locally uniform convergence on \(\ComplexNumbers\), pointwise multipliability and exact
\texttt{HasProd}, and entireness of the geometric sinc product, together with its
reciprocal-Gamma reflection factorization;
\cref{integration:eq:geometric-sinc-characteristic}.\\
\path{GeometricSincCharacteristicFunction.lean} &
The zero-definition, four-theorem interface: the two exact bridges from the geometric-uniform
characteristic function to the rescaled geometric sinc product and paired reciprocal-Gamma
factors, followed by locally uniform convergence of the full phase-bearing prefixes on the
real frequency line and its uniform-on-compact-sets specialization, for every real
\(|q|<1\), including zero and negative ratios;
\cref{integration:eq:geometric-sinc-characteristic,integration:eq:geometric-reciprocal-gamma-characteristic}.\\
\path{RvachevPochhammerFactorization.lean} &
The one-definition, nine-theorem complex Pochhammer interface: strict-contraction product
convergence, absolute paired perturbation summability, the global general-\(q\)
\(S_q(z)=\prod_{k\ge0}(z^2/(k+1)^2;q^2)_\infty\) factorization, and its dyadic
Rvachev/Fabius specializations;
\cref{integration:eq:geometric-sinc-pochhammer}.\\
\path{PrimePowerBinomialValuation.lean} &
The exhaustive zero-definition, six-theorem prime-power row interface: the
additive/subtraction formulas for row \(p^m\), the unit row \(p^m-1\), the
exact \(p^m-2\) companion valuation on \(0<j<p^m\), and both dyadic wrappers.
The strict companion upper bound is essential, and no valuation-histogram
claim is made.\\
\path{GeometricPochhammerNormalConvergence.lean} &
The zero-definition, three-theorem outer-product interface: locally uniform
convergence on \(\ComplexNumbers\) to \(S_q\) for every complex
\(\lVert q\rVert<1\), including \(q=0\), followed by the nome-\(1/4\)
Rvachev-product and bounded-Fabius specializations.  The centered/MGF,
outside-disk reciprocal, pole-divisor, and zero--pole clauses of the compound
spectral theorem remain open;
\cref{integration:eq:geometric-sinc-pochhammer}.\\
\path{GeneralizedRvachevIdentifiability.lean} &
The zero-definition, six-theorem recovery interface: dyadic analytic orders
are inclusive exponent prefixes, consecutive finite-order differences decode
the sequence, and two admissible generalized Rvachev products are equal
exactly when their exponent sequences are equal;
\cref{integration:thm:generalized-rvachev-identifiability}.\\
\path{QPochhammerEntire.lean} &
The zero-definition, five-theorem fixed-nome interface: locally uniform
convergence in the symbol variable, entireness, the total factor-zero criterion,
the reciprocal-power lattice for nonzero \(q\), and analytic order one at every
zero, including the \(q=0\) boundary;
\cref{integration:thm:qpochhammer-entire}.\\
\path{QPochhammerDissection.lean} &
The exhaustive zero-definition, two-theorem finite residue-class interface:
exact dissection over every commutative ring for all natural \(r,n\), and the
remainder form under \(u\le r\), including the endpoint \(u=r\).\\
\path{QPochhammerInfinite.lean} &
The exhaustive one-definition, twenty-seven-theorem generic infinite-product
interface: topological-ring totalization; strict-contraction convergence,
concatenation, factor removal and positive-modulus dissection in complete
normed commutative rings; the factor-zero locus under a multiplicative norm;
locally compact field continuity; and complex entireness with an explicit
nonzero reciprocal-zero derivative when \(q\ne0\).  It supplies no
\texttt{analyticOrderAt} theorem and does not replace the preceding wrapper.\\
\path{GaussianBinomialAtNegOneDerivative.lean} &
The zero-definition, four-theorem first-jet interface at \(q=-1\): two exact
Gaussian derivative evaluations and simple-root multiplicity for the
even-row/odd-column locus, first over \(\IntegerNumbers\) and then over
characteristic-zero commutative rings under the strict interior hypothesis.\\
\path{QBinomialTheoremInfinite.lean} &
The exhaustive one-definition, twenty-two-theorem dominated-limit and
infinite q-binomial interface, including Euler's product and reciprocal sums.\\
\path{GaussianBinomialFixedColumnRate.lean} &
The exhaustive zero-definition, ten-theorem effective fixed-column interface:
finite-product defect bounds, a denominator-free relative Gaussian bound,
fixed and shifted nonasymptotic additive estimates, shifted convergence, and
four relative/additive geometric \texttt{IsBigO} results.  The ring layer uses
\(\lVert q\rVert\le1\); the normed-field layer uses only
\(\lVert q\rVert<1\), without completeness or (q\ne0\).\\
\path{JacobiTripleProduct.lean} &
The exhaustive two-definition, twenty-five-theorem finite/infinite triple-product
and pentagonal-sum interface.\\
\path{QPascalSummation.lean} &
The exhaustive zero-definition, four-theorem finite q-Pascal summation and
commuting Gaussian-binomial interface.\\
\path{GaussianBinomialContinuity.lean} &
The exhaustive zero-definition, three-theorem continuity, \(q\to1\), and
finite-Pochhammer quotient interface.\\
\path{GaussianBinomialPalindromic.lean} &
The exhaustive zero-definition, fourteen-theorem generic reflection, degree,
boundary coefficient, palindromicity, exact-degree/monicity, mean, and total
linear-coefficient interface.\\
\path{GaussianBinomialPolynomialStructure.lean} &
The exhaustive zero-definition, five-theorem universal degree, monicity,
constant-coefficient, reflection, and bounded-index coefficient-symmetry interface.\\
\path{CentralQBinomialReduction.lean} &
The exhaustive zero-definition, six-theorem central squared-base reduction interface.\\
\path{CyclotomicFactorization.lean} &
The exhaustive zero-definition, seven-theorem finite-product and Gaussian cyclotomic-factorization interface.\\
\path{PrimitiveRootBlock.lean} &
The exhaustive zero-definition, three-theorem primitive-root Gaussian-vanishing, phase, and complete-block interface.\\
\path{QLucas.lean} &
The exhaustive zero-definition, seven-theorem natural quadratic, block-coefficient, phase, and q-Lucas interface; its similarly named local helper is private, and the public \path{two_mul_choose_two} belongs to \path{QChuVandermonde.lean}.\\
\path{TwoPhiOneReversal.lean} &
The exhaustive two-definition, twelve-theorem reflected-parameter, finite-to-tsum, involution, and terminating two-phi-one reversal interface.\\
\path{QChuVandermonde.lean} &
The exhaustive zero-definition, ten-theorem q-Chu--Vandermonde interface, including both denominator-cleared identities, both quotient and actual-tsum evaluations, and the separately named partial-domain reversal proof of the second.\\
\path{JacobiTwoSquareCount.lean} &
The exhaustive zero-definition, four-theorem nonzero two-square count, conditional prime product, and unconditional chi-four and alternating-odd Lambert interface.\\
\path{CyclotomicDivisibility.lean} &
The exhaustive zero-definition, three-theorem cyclotomic carry, rational divisibility, and primitive-root multiple-value interface.\\
\path{QCatalan.lean} &
The exhaustive one-definition, eleven-theorem integral q-Catalan divisibility, degree, and value-at-one interface.\\
\path{NewtonInterpolation.lean} &
The exhaustive three-definition, nineteen-theorem finite-node and geometric-grid Newton interface, including the collision-free node-polynomial API and incoming interpolant compatibility wrappers.\\
\path{QBetaIntegral.lean} &
The exhaustive one-definition, eight-theorem Jackson q-beta and q-Gamma interface.\\
\path{QuantumBinomial.lean} &
The exhaustive zero-definition, two-theorem quantum-plane and noncommutative
q-binomial interface.\\
\path{RogersSzegoPolynomial.lean} &
The exhaustive one-definition, nine-theorem Rogers--Szeg\H{o} polynomial,
recurrence, and generating-series interface.\\
\path{ContinuousCDF.lean} &
The three reusable real-CDF bridges: continuity from null singletons, CDF reflection from
affine reflection invariance, and identification of an everywhere pointwise CDF derivative
as a Lebesgue Radon--Nikodym density;
\cref{integration:sec:weighted-naturality}.\\
\path{GeometricUniformCDF.lean} &
The named geometric CDF and density API: continuity and reflection for \(|q|<1\), tails for
\(0\le q<1\), the exact zero/nonpositive/zero-measure/nonidentity diagnosis for the selected
formula when \(q\le0\), and, for \(0<q<1\), conditioning and interval-integral refinement,
everywhere differentiation, nonnegativity, symmetry, compact support, the exact
\texttt{withDensity} identity, and real-space \(C^\infty\) regularity;
\cref{integration:sec:weighted-naturality}.\\
\path{SubgraphFubini.lean} &
\lean{integral_smul_setIntegral_subgraph} and
\lean{integral_smul_setIntegral_supergraph}, the real- or complex-weighted Banach-valued
open-subgraph and atom-preserving closed-supergraph exchanges over arbitrary s-finite
measures; \cref{integration:sec:survival-fubini}.\\
\path{SurvivalLayerCake.lean} &
The support-free clamped and interval-supported partial/full Banach-valued survival
identities; \cref{integration:thm:survival-layer-cake}.\\
\path{CDFLayerCake.lean}, \path{FractionalCDFLayerCake.lean} &
The support-free clamped, one-sided-upper-supported partial, compactly supported partial,
and full Banach-valued lower-CDF identities, followed by the positive-real clamped,
lower-supported positive-part, probability-CDF, and compact endpoint-power fractional forms;
\cref{integration:thm:fractional-cdf-layer-cake}.\\
\path{FabiusFractionalIntegral.lean} &
The distribution-law and product-sample-space positive-real stopped-power formulas for
every real endpoint, together with both unit-endpoint ordinary-moment forms;
\cref{integration:thm:fractional-cdf-layer-cake}.\\
\path{InverseLayerCake.lean} &
The compact-interval Galois connection; measurable-free generic and Fabius
explicit-primitive/pointwise-\(C^1\) layer cake; the real-primitive, real- or
complex-weighted arbitrary-absolutely-continuous subtraction and normalized forms; the
full weighted inverse-area identity, equality with the $\RvachevUp$ area, and the inverse mean
\(1/2\); eleven public theorems in total;
\cref{integration:sec:survival-fubini}.\\
\path{InversePairIntegral.lean} &
The full and variable-upper real two-function absolutely-continuous inverse-pair identities
on arbitrary ordered compact rectangles, with clock monotonicity and measurable clamped
representatives derived from interval preservation and the strict order equivalence; their
Fabius wrappers, the full area sum, and the partial inverse-area identity; six public
theorems in total;
\cref{integration:sec:survival-fubini}.\\
\path{MeasureCauchyTransform.lean}, \path{MeasureCauchyMomentLaurent.lean},
\path{GeometricUniformCauchy.lean},
\path{CauchyTransform.lean}, \path{CauchyCDF.lean},
\path{CauchySurvival.lean}, \path{CauchyHigherPowers.lean},
\path{CauchyRenormalization.lean} &
The generic finite-measure affine Cauchy calculus; the arbitrary-center
\(\mathrm{RCLike}\) moment/Laurent expansion under an almost-everywhere radius-\(R\) ball
and \(R<\lVert z-c\rVert\), with exact remainder, finite-mass geometric bound,
\path{Summable}/\path{tsum}/\path{HasSum} forms, and complex named-transform wrappers;
the exact-support geometric-\(q\) transform DDE and adjacent-power hierarchy for every
nonzero \(|q|<1\), including
negative ratios; the canonical unit/centered all-power bridge, named unit Stieltjes DDE,
and both complex-slit adjacent-order recurrences; atom-exact CDF and survival integration by
parts at every positive kernel power; and the centered DDE and all-order finite
Thue--Morse derivative orbit.  The generic fixed-point theorems need no probability
normalization or topological/measurability assumptions on their invariant carrier.\\
\path{StieltjesLogFixedPoint.lean}, \path{StieltjesResolventHierarchy.lean},
\path{StieltjesGeneralizedOrder.lean}, \path{StieltjesMomentLaurent.lean},
\path{StieltjesCauchyTransform.lean}, \path{StieltjesHerglotz.lean},
\path{StieltjesInversion.lean}, \path{StieltjesPerron.lean},
\path{OrthogonalPolynomialJacobi.lean} &
The real logarithmic fixed point and integer hierarchy for \(z>1\), real order lowering for
\(\alpha>1\), the radius-one center-zero up-measure Laurent specialization through
\nolinkurl{ae_norm_sub_zero_le_one_rvachevMeasure} and
\nolinkurl{measureCauchyMoment_rvachevMeasure_zero}, its retained real and complex
exterior \nolinkurl{upMoment} expansions, finite-height
Herglotz--Poisson formulas, integrated interval Stieltjes--Perron inversion, and initial
exact Jacobi data.  No complex logarithmic continuation, complex order, pointwise or
nontangential Sokhotski--Plemelj/Hilbert-principal-value formula, or complete
J-fraction/Pad\'e theory is asserted.\\
\path{GramStieltjesNaturality.lean}, \path{RvachevRationalJacobi.lean} &
The first module has zero public definitions and six public theorems:
\lean{momentPairing_map},
\lean{map_gramStieltjesNumerator},
\lean{map_gramStieltjesPolynomial},
\lean{map_gramStieltjesNorm},
\lean{map_gramStieltjesJacobiDiagonal}, and
\lean{map_gramStieltjesJacobiSubdiagonal}.  They prove scalar naturality of the finite
pairing, fraction-free numerator, normalized polynomial, norm, and Jacobi coefficients at
their commutative-semiring, commutative-ring, and field boundaries.  The second module has
four public definitions,
\lean{rvachevHankelDetRat},
\lean{rvachevOrthoPolynomialRat},
\lean{rvachevOrthoNormRat}, and
\lean{rvachevJacobiSubdiagonalRat}, and thirteen public theorems:
\lean{upMoment_eq_rvachevRawMomentRat_cast},
\lean{rvachevHankelDetRat_cast},
\lean{rvachevHankelDetRat_pos},
\lean{rvachevOrthoPolynomialRat_cast},
\lean{rvachevOrthoPolynomialRat_isMonicOfDegree},
\lean{momentPairing_rvachevOrthoPolynomialRat_eq_zero},
\lean{rvachevOrthoNormRat_cast},
\lean{rvachevOrthoNormRat_pos},
\lean{gramStieltjesJacobiDiagonal_rvachevRawMomentRat_eq_zero},
\lean{rvachevJacobiSubdiagonalRat_cast},
\lean{rvachevJacobiSubdiagonalRat_pos},
\lean{rvachevJacobiSubdiagonalRat_eq_det_ratio}, and
\lean{rvachevOrthoPolynomialRat_three_term}.  This gives positive rational determinants,
monic rational orthogonal polynomials, positive rational norms, a zero diagonal, positive
rational subdiagonals, their real-cast comparisons, the determinant cross-ratio, and the
centered rational recurrence in every degree.  The zero-based subdiagonal at index $n$ is
the conventional $\beta_{n+1}$.  These four rational definitions are noncomputable finite
objects, not native-code evaluators, and do not themselves evaluate $H_4$ or $\beta_4$;
the downstream executable tranche and its values leaf do so without changing this
all-degree interface.  This module proves no root, Christoffel, quadrature,
continued-fraction, Pad\'e, or asymptotic theorem.\\
\path{PolynomialMomentGramDeterminant.lean} &
Exactly two public definitions,
\lean{polynomialCoefficientMatrix} and \lean{polynomialMomentGramMatrix}, and seven public
theorems:
\lean{polynomialCoefficientMatrix_apply},
\lean{polynomialMomentGramMatrix_apply},
\lean{polynomialMomentGramMatrix_eq_transpose_mul_hankel_mul},
\lean{polynomialMomentGramMatrix_det_eq_coefficient_det_sq_mul},
\lean{polynomialCoefficientMatrix_det_eq_prod_coeff},
\lean{polynomialMomentGramMatrix_det_eq_prod_coeff_sq_mul}, and
\lean{gramStieltjesJacobiSubdiagonal_eq_polynomialMomentGramMatrix_det_ratio}.
For a coherent family $p_j$, coefficient expansion proves
$G_n=C_n^{\mathsf T}H_nC_n$; the degree bound makes $C_n$ upper triangular with diagonal
$[X^j]p_j$, so taking determinants gives
$\det G_n=(\prod_{j<n}[X^j]p_j)^2\det H_n$.  A nonzero diagonal then transports the
zero-based finite Jacobi subdiagonal to the polynomial-Gram determinant cross-ratio.  This
last Lean equality imposes no Hankel-nonvanishing hypothesis because field division is
total: a zero middle Hankel determinant makes both cross-ratios zero, not a genuine
nonsingular Jacobi recurrence.  This is
finite ring/field algebra with no measure, positivity, entry-rationality,
Gaunt/Wigner/$3j$, Christoffel, or infinite-Jacobi assertion.\\
\path{LegendrePolynomialRational.lean} &
Exactly two public definitions, the executable scalar coefficient
\lean{legendrePolynomialCoeffRat} and the noncomputable polynomial wrapper
\lean{legendrePolynomialRat}, and six public theorems:
\lean{legendrePolynomialRat_cast},
\lean{coeff_legendrePolynomialRat},
\lean{natDegree_legendrePolynomialRat},
\lean{coeff_legendrePolynomialRat_self},
\lean{coeff_legendrePolynomialRat_self_ne_zero}, and
\lean{coeff_legendrePolynomialRat_self_div_succ}.
The bounded rational coefficient sum agrees coefficientwise with the wrapper; the wrapper
casts to the real Rodrigues-normalized ordinary Legendre polynomial, has degree \(n\),
leading coefficient \(2^{-n}\binom{2n}{n}\), nonzero top coefficient, and consecutive-top
quotient \((n+1)/(2n+1)\).  This is finite coefficient algebra, with no root, Christoffel,
quadrature, orthogonality, or asymptotic assertion.\\
\path{LegendreGaunt.lean} &
Exactly four public definitions,
\lean{legendreLebesgueMomentRat},
\lean{legendreGauntRat},
\lean{legendreGaunt}, and
\lean{legendreProductLinearizationCoeffRat}, and twelve public theorems:
\lean{legendreLebesgueMomentRat_even},
\lean{legendreLebesgueMomentRat_odd},
\lean{legendreLebesgueMomentRat_cast},
\lean{legendreGauntRat_eq_momentFunctional},
\lean{legendreGauntRat_cast},
\lean{legendreGauntRat_swap_left},
\lean{legendreGauntRat_swap_right},
\lean{legendrePolynomial_mul_eq_sum_gaunt},
\lean{legendrePolynomialRat_mul_eq_sum_gaunt},
\lean{legendreGauntRat_eq_zero_of_odd_sum},
\lean{legendreGauntRat_eq_zero_of_add_lt}, and
\lean{legendreGauntRat_eq_zero_of_triangle_violation}.\\
\emph{LegendreGaunt semantics} &
The executable rational monomial moments on \([-1,1]\) feed a bounded triple
coefficient sum, whose cast is the real triple-Legendre interval integral.  The two swap
theorems generate the index symmetries, and the normalized coefficient gives exact finite
product linearization over rational and real scalars.  Odd total degree and any strict
triangle violation force vanishing; these are necessary support conditions only,
with no converse or nonvanishing theorem in this module.  The downstream closed-form
module supplies the exact total squared zero-row Wigner \(3j\) identification, factorial
forms, sharp support, and nonnegativity results; none of those conclusions should be read
back into this predecessor module.\\
\path{LegendreGauntClosedForm.lean} &
Exactly two public definitions,
\lean{legendreGauntAdmissible} and \lean{legendreWignerThreeJZeroSqRat}, and twenty-five
public theorems:
\lean{legendreGauntAdmissible_iff_exists_pairwise_add},
\lean{legendreGauntAdmissible_pairwise_add},
\lean{legendreWignerThreeJZeroSqRat_pairwise_add},
\lean{legendreWignerThreeJZeroSqRat_pairwise_add_factorial},
\lean{legendreWignerThreeJZeroSqRat_eq_factorial_of_halfSum},
\lean{legendreWignerThreeJZeroSqRat_eq_zero_of_not_admissible},
\lean{legendreGauntRat_add_boundary},
\lean{legendreGauntRat_add_boundary_eq_two_mul_wignerThreeJZeroSqRat},
\lean{legendreGauntRat_zero_left},
\lean{legendreGauntRat_zero_left_eq_two_mul_wignerThreeJZeroSqRat},
\lean{legendreGauntRat_pairwise_add_eq_two_mul_wignerThreeJZeroSqRat},
\lean{legendreGauntRat_eq_zero_of_not_admissible},
\lean{legendreGauntRat_eq_two_mul_wignerThreeJZeroSqRat},
\lean{legendreGaunt_eq_two_mul_wignerThreeJZeroSqRat},
\lean{legendreWignerThreeJZeroSqRat_pos_iff_admissible},
\lean{legendreWignerThreeJZeroSqRat_nonneg},
\lean{legendreWignerThreeJZeroSqRat_eq_zero_iff_not_admissible},
\lean{legendreGauntRat_pos_iff_admissible},
\lean{legendreGauntRat_eq_zero_iff_not_admissible},
\lean{legendreGaunt_pos_iff_admissible},
\lean{legendreGaunt_eq_zero_iff_not_admissible},
\lean{legendreGauntRat_nonneg},
\lean{legendreGaunt_nonneg},
\lean{legendreProductLinearizationCoeffRat_eq_mul_wignerThreeJZeroSqRat}, and
\lean{legendreProductLinearizationCoeffRat_pos_iff_admissible}.
Admissibility is even total degree together with the three weak triangle inequalities,
equivalently \(i=b+c\), \(j=a+c\), \(k=a+b\).  With \(s=a+b+c\), the total rational
square datum \(W_{ijk}^{2}=\lean{legendreWignerThreeJZeroSqRat}(i,j,k)\) is
\[
 W_{ijk}^{2}=
 \frac{\binom{2a}{a}\binom{2b}{b}\binom{2c}{c}}
      {(2s+1)\binom{2s}{s}},
\]
with equivalent pairwise-coordinate and half-sum factorial formulas.  At arbitrary
natural indices the global identities are
\[
 G_{ijk}^{\RationalNumbers}=2W_{ijk}^{2},
 \qquad
 G_{ijk}=2\bigl(W_{ijk}^{2}:\RealNumbers\bigr).
\]
The datum is positive exactly on the admissible support and zero otherwise; the rational
and real Gaunt coefficients have the same sharp support and are nonnegative.  The
coefficient of \(P_k\) in \(P_iP_j\) is \((2k+1)W_{ijk}^{2}\) and is positive exactly on
that support.  Separate boundary theorems cover \(k=i+j\) and \(i=0\).  This is
deliberately the total square of the integer-index zero-row datum: no signed symbol or
Condon--Shortley phase is chosen, and half-integer indices, nonzero magnetic indices,
general \(3j/6j/9j\) theory, Wigner orthogonality, and recoupling identities remain open.\\
\path{FabiusLegendreHankelDeterminant.lean} &
Exactly two public definitions, \lean{upLegendreGramMatrix} and
\lean{upLegendreGramDet}, and seven public theorems:
\lean{upLegendreGramMatrix_apply_eq_integral},
\lean{upLegendreGramDet_eq_prod_leadingCoeff_sq_mul_hankelDet},
\lean{upLegendreGramDet_zero}, \lean{upLegendreGramDet_pos},
\lean{coeff_legendrePolynomial_self_div_succ},
\lean{gramStieltjesJacobiSubdiagonal_upMoment_eq_upLegendreGramDet_ratio}, and
\lean{rvachevJacobiSubdiagonalRat_cast_eq_upLegendreGramDet_ratio}.
For arbitrary $F:\lean{BoundedFabius}$, let $D_n$ be the first-$n$ Legendre Gram
determinant in its up-moment pairing and let $L_j=2^{-j}\binom{2j}{j}$.  The determinant
identity $D_n=(\prod_{j<n}L_j^2)\det H_n$, the leading-coefficient quotient, and the real
Gram cross-ratio hold for every such $F$.  Here $D_0=1$ is the determinant of the empty
$0\mathbin{\times}0$ Gram matrix, extending the paper's prior $N\geq1$ definition.  The
index $n$ is zero-based---the conventional $\beta_{n+1}$---and
\[
 \beta_{n+1}=\left(\frac{n+1}{2n+1}\right)^2
 \frac{D_{n+2}D_n}{D_{n+1}^2},
\]
with totalized division in the arbitrary-$F$ theorem.  Identifying Gram entries with
integrals, proving $D_n>0$, and identifying this ratio with the real cast of the rational
Rvachev subdiagonal require $h_F:\lean{IsFabius}\,F$; positivity then makes the ratio a
genuine nonsingular Jacobi formula.  This module itself proves no finite Gaunt entry
formula or entry rationality by that route; the two downstream Gaunt modules below do.
It proves no Wigner/$3j$ or factorial closed form, Christoffel reconstruction, or infinite
product.  The shifted-atom kernel \lean{rvachevTranslateGram} is a distinct
unweighted interval Gram construction.\\
\path{FabiusLegendreRationalGram.lean} &
Exactly three executable public definitions,
\lean{rvachevLegendreGramEntryRat},
\lean{rvachevLegendreGramMatrixRat}, and
\lean{rvachevLegendreGramDetRat}, and eleven public theorems:
\lean{rvachevLegendreGramEntryRat_eq_momentPairing},
\lean{rvachevLegendreGramMatrixRat_apply},
\lean{rvachevLegendreGramMatrixRat_eq_polynomialMomentGramMatrix},
\lean{rvachevLegendreGramEntryRat_cast},
\lean{rvachevLegendreGramMatrixRat_cast},
\lean{rvachevLegendreGramDetRat_cast},
\lean{rvachevLegendreGramDetRat_eq_prod_leadingCoeff_sq_mul_rvachevHankelDetRat},
\lean{rvachevLegendreGramDetRat_zero},
\lean{rvachevLegendreGramDetRat_pos},
\lean{rvachevOrthoNormRat_eq_rvachevLegendreGramDetRat_ratio}, and
\lean{rvachevJacobiSubdiagonalRat_eq_rvachevLegendreGramDetRat_ratio}.
The bounded rational entry sum and its finite matrix agree with the abstract moment-pairing
objects.  For \(F:\lean{BoundedFabius}\) satisfying \(\lean{IsFabius}\,F\), their entry,
matrix, and determinant cast to the real up-law objects.  The rational determinant is the
Hankel determinant times the squared-leading-coefficient product, equals one in order zero,
and is positive; its exact ratios give the rational norm and zero-based
\(\beta_{n+1}\), with the latter prefactor \(((n+1)/(2n+1))^2\).  The next module
supplies the finite Gaunt entry formulas, and its downstream closed-form module converts
them to finite total squared zero-row Wigner-$3j$ sums.  The rational Gram module itself
proves neither formula.  Signed symbols and phases, half-integer or
nonzero-magnetic-index theory, and general recoupling and Wigner orthogonality,
Christoffel reconstruction, roots or quadrature, infinite Jacobi
products/continued fractions, and asymptotics remain open.\\
\path{FabiusLegendreGaunt.lean} &
Exactly one public definition, \lean{canonicalRvachevFullLegendreCoefficientRat}, and
eight public theorems:
\lean{canonicalRvachevFullLegendreCoefficientRat_even},
\lean{canonicalRvachevFullLegendreCoefficientRat_odd},
\lean{canonicalRvachevFullLegendreCoefficientRat_cast},
\lean{canonicalRvachevFullLegendreCoefficientRat_eq_normalized_moment},
\lean{rvachevLegendreGramEntryRat_eq_sum_full_gaunt},
\lean{rvachevLegendreGramEntryRat_eq_sum_gaunt},
\lean{rvachevLegendreGramMatrixRat_apply_eq_sum_gaunt}, and
\lean{upLegendreGramMatrix_apply_eq_sum_gaunt}.
The definition equals \lean{canonicalRvachevLegendreCoefficientRat}\((k/2)\) when
\(2\mid k\) and zero otherwise, so its even restriction is the canonical sequence and its
odd entries vanish.  Write \(c_k^{\mathrm{full}}\) for this definition and
\(M_{\RationalNumbers}\) for the rational moment functional.  Unconditionally in \(k\),
\[
 c_k^{\mathrm{full}}=\frac{2k+1}{2}\,M_{\RationalNumbers}(P_k).
\]
For all natural \(i,j\), write \(H_{ij}^{\RationalNumbers}\) for the rational Gram entry,
\(c_r\) for the canonical even coefficient, and \(G_{ijk}^{\RationalNumbers}\) for the
rational Gaunt coefficient.  Then
\[
 H_{ij}^{\RationalNumbers}=\sum_{k=0}^{i+j}c_k^{\mathrm{full}}G_{ijk}^{\RationalNumbers},
 \qquad
 H_{ij}^{\RationalNumbers}=\sum_{r=0}^{\Floor{(i+j)/2}}c_rG_{ij,2r}^{\RationalNumbers}.
\]
The rational matrix theorem gives the even formula for every \(i,j:\lean{Fin}\,n\).
Only the cast theorem and real matrix-entry formula require
\(F:\lean{BoundedFabius}\) with \(h_F:\lean{IsFabius}\,F\): they identify the full
rational coefficient with \lean{rvachevFullLegendreCoefficient}\((F,k)\) and the real
entry with the same even finite sum using
\lean{rvachevLegendreCoefficient}\((F,r)\) and
\lean{legendreGaunt}\((i,j,2r)\).  The other six theorems are unconditional.  These are
finite polynomial and finite-sum identities.  Together with \path{LegendreGaunt.lean},
they close the finite Gaunt-integral/product-linearization layer and the finite Rvachev
Gram-entry expansion.  This module itself proves no Wigner/$3j$ identification or
factorial formula and no nonnegativity by that route; the two downstream closed-form
modules supply exactly the total squared integer zero-row results and finite Gram sums.
No signed-symbol or phase convention, half-integer or nonzero-magnetic-index theory,
infinite Legendre-series/integral interchange, Christoffel reconstruction, or named values
for the displayed \(G_3\) entries are proved here.\\
\path{FabiusLegendreGauntClosedForm.lean} &
Zero public definitions and exactly three public theorems:
\lean{rvachevLegendreGramEntryRat_eq_two_mul_sum_wignerThreeJZeroSqRat},
\lean{rvachevLegendreGramMatrixRat_apply_eq_two_mul_sum_wignerThreeJZeroSqRat}, and
\lean{upLegendreGramMatrix_apply_eq_two_mul_sum_wignerThreeJZeroSqRat}.
They substitute the exact identity
\(G_{ijk}^{\RationalNumbers}=2W_{ijk}^{2}\) into the finite Gaunt expansions.  Thus,
unconditionally for natural \(i,j\),
\[
 H_{ij}^{\RationalNumbers}
 =2\sum_{r=0}^{\Floor{(i+j)/2}}c_r W_{ij,2r}^{2},
\]
and the same formula holds entrywise for the executable rational Gram matrix.  The real
up-law matrix formula replaces \(c_r\) by
\lean{rvachevLegendreCoefficient}\((F,r)\) and casts the square datum to
\(\RealNumbers\); it requires exactly \(F:\lean{BoundedFabius}\) and
\(h_F:\lean{IsFabius}\,F\).  These are finite sums and inherit the square-datum boundary:
no signed Wigner symbol, phase convention, half-integer or nonzero-magnetic indices,
general \(3j/6j/9j\) theory, orthogonality, recoupling, or infinite-series assertion.\\
\path{FabiusLegendreRationalGramValues.lean} &
Zero public definitions and exactly eleven public theorems:
\lean{moment_four},
\lean{rvachevLegendreGramDetRat_one},
\lean{rvachevLegendreGramDetRat_two},
\lean{rvachevLegendreGramDetRat_three},
\lean{rvachevLegendreGramDetRat_four},
\lean{rvachevLegendreGramDetRat_five},
\lean{rvachevOrthoNormRat_four},
\lean{rvachevJacobiSubdiagonalRat_three},
\lean{hankelRatio_four},
\lean{integral_sq_upOrthoPolynomial_four}, and
\lean{hankelRatio_four_div_three}.
They compute the raw eighth moment as $132809/32531625$ and the order-one through
order-five rational Legendre Gram determinants as
$1$, $1/9$, $8/2025$, $39616/602791875$, and
$16544275456/27453718922765625$.  They then give
$H_4=\lean{rvachevOrthoNormRat}(4)=26727424/55791736875$ and, with the Lean
subdiagonal index zero-based,
$\beta_4=\lean{rvachevJacobiSubdiagonalRat}(3)=835232/4640643$.
These first eight computations are unconditional over $\RationalNumbers$.  The final three theorems
transport the norm to the real Hankel ratio, its squared orthogonal-polynomial integral,
and the quotient of the fourth and third real Hankel ratios only for
$F:\lean{BoundedFabius}$ together with $h_F:\lean{IsFabius}\,F$.
Together with the preceding modules, this values leaf completes an eleven-module
finite-moment/Legendre/Gaunt inventory of twenty public definitions and 109 public
theorems, hence 129 public declarations.  The former nine-module checkpoint was eighteen
definitions and eighty-one theorems; the two closed-form leaves add two definitions and
twenty-eight theorems.  The finite Gaunt--zero-row-Wigner-square identification, its
factorial form, sharp support and positivity, product-linearization coefficients, and
finite rational and real up-law Gram-entry square sums are therefore exact.  The values
leaf still gives no named theorem for the displayed \(G_3\) entries.  Signed or general
Wigner symbols and phase conventions, half-integer or nonzero-magnetic-index symbols,
general recoupling and Wigner orthogonality, Christoffel reconstruction, roots or
quadrature, infinite Jacobi products/continued fractions, and asymptotics remain open.\\
\path{FractionalVolterra.lean}, \path{FractionalVolterraSemigroup.lean} &
The total real-order Gamma-normalized operator, its unconditional endpoint/locality/scalar,
order-one, and positive-natural-order rules, and its positive-order ordered-endpoint
continuous-kernel integrability and input additivity; raw shifted-beta-kernel integrability
for positive orders and ordered endpoints \(s\le x\); the raw and normalized values under a
strict endpoint inequality; and additive composition of positive orders for continuous
Banach-valued inputs on ordered compact intervals.  A generic open-tail cutoff and its
fractional specialization truncate compact-right-support integrals at \(\min(x,b)\).
Normalized shifted powers, arbitrary shifted real powers in the sharp range \(\rho>-1\),
and constants have their exact Gamma-quotient values.  No corresponding semigroup for
merely interval-integrable inputs or noncomplete targets, fractional derivative, complex
order, order-zero law, or reversed-endpoint fractional interpretation;
\cref{integration:rem:fractional-volterra-boundary}.\\
\path{FractionalVolterraCalculus.lean}, \path{FabiusFractionalVolterra.lean} &
Increasing-affine covariance for arbitrary real order, positive scale, and ordered endpoints;
positive-order Banach-valued integration by parts across an interval-integrable right
interior derivative with its exact
left-boundary term; a total causal Rvachev fractional primitive with its support-truncated
formula, positive-natural bridge, and positive-order semigroup; and the three positive-real
single-step order shifts for the signed and bounded Fabius functions and \(\RvachevUp\).  These
modules export four theorems and one definition with six theorems, respectively, but no
nonpositive-scale/reversed covariance, fractional derivative or Caputo operator,
complex-order continuation, shifted dyadic-lattice or endpoint-moment wrapper,
transform/tail series, or inverse specialization;
\cref{integration:rem:fractional-volterra-boundary}.\\
\path{ProbabilityLaplaceMoments.lean} &
The survival function, the finite-measure Banach-valued $L^1$ full and clipped
survival-kernel identities, their partial/full Rvachev layer-cake specializations,
the differentiable compact-support Fubini identity, raw and endpoint moment
identifications, zero-tilt normalization, and positivity.  Its later exhaustive
two-theorem extension consists of
\lean{weightedSumDistribution_real_Ici_eq_rvachevUp_of_nonneg} and
\lean{integral_pow_weightedSumDistribution_eq_mul_intervalIntegral_rvachevUp},
the atomless closed-tail identity on \(t\geq0\) and the full-law
positive-natural-degree moment identity for \(n\geq1\), respectively.  They
make \nolinkurl{prop:up-tail} and \nolinkurl{cor:up-moments} Exact; the dyadic almost-sure
support and restriction declarations
are owned upstream by \path{ProbabilityRepresentation.lean};
\cref{integration:sec:survival-fubini,integration:sec:probability-laplace}.\\
\path{LaplaceMoments.lean} &
\lean{tiltedSurvivalMoment}, \lean{fabiusLaplaceMoment}, the differential recurrence,
iterated derivatives, smoothness, and evaluation at zero tilt; \cref{integration:sec:probability-laplace}.\\
\path{StepMeasureBridge.lean} &
Almost-everywhere replacement of half-endpoint indicators, exact overlap integrals, atomic
measure comparisons, the shrunken/enlarged interval sandwich, portmanteau, and convergence of
every interval integral; \cref{integration:sec:step-measure-bridge}.\\
\path{FabiusInverseDyadicClosedForm.lean} &
Compositions, path weights, last-edge decomposition, generic triangular solver, uniqueness,
and the exact composition formula for $F(2^{-n})$; \cref{integration:sec:path-sums,integration:sec:fabius-composition-formula}.\\
\path{FabiusDiscreteLimitToeplitz.lean} &
Corrected prefix length, exact $q$-binomial row mass, total-variation bound $16$, the generic
finite-row Toeplitz theorem, and the specialization $q(n,j)=2n-j$;
\cref{integration:sec:toeplitz,integration:sec:fabius-toeplitz}.\\
\path{FinitePowerSeriesFilter.lean}, \path{GeometricPowerSeriesFilter.lean} &
The coefficientwise diagonal action of arbitrary finite-node weighted rescale filters,
rescaling and coefficient-ring naturality, exactness through the prescribed order, and the
denominator-free Gaussian residual multiplier for geometric filters, including the
distinct-node Lagrange specialization.  These are formal power-series identities only and
assert no analytic convergence, norm bound, or remainder estimate.\\
\path{AnalyticSeriesFilter.lean} &
The definitions \lean{finiteAnalyticSeriesFilter} and
\lean{geometricAnalyticSeriesFilter}; diagonal summability and evaluation
\lean{summable_finiteAnalyticSeriesFilter_diagonal} and
\lean{finiteAnalyticSeriesFilter_eq_tsum}; the exact head/tail decompositions
\lean{finiteAnalyticSeriesFilter_eq_head_add_tail_of_exact} and
\lean{finiteAnalyticSeriesFilter_eq_constant_add_tail_of_exact_zero}; and the exact Gaussian
tails \lean{geometricAnalyticSeriesFilter_eq_constant_add_gaussian_tsum},
\lean{geometricLagrangeAnalyticSeriesFilter_eq_constant_add_gaussian_tsum}, and
\lean{geometricLagrangeAnalyticSeriesFilter_shifted}.  Every sampled-series assumption is
unconditional and explicit; no conditional/boundary convergence, radius or uniform
convergence, error/norm or asymptotic estimate, positivity, or Fabius-specific acceleration
claim is made.\\
\path{ImplicitPowerSeries.lean} &
One distinguished root definition and eight nonmonic formal implicit-root theorems.  The
three \path{FormalImplicitRoot} results give existence, uniqueness, and unique existence in
a complete adic commutative ring from an approximate root and a derivative unit modulo the
ideal (standalone uniqueness uses the stated Jacobson-radical bound).  The six
\path{PowerSeries.Implicit} declarations give arbitrary-base and zero-constant lifting, the
distinguished root, its constant coefficient and evaluation equation, and uniqueness over
an arbitrary commutative coefficient ring.  No concrete inverse-Fabius germ, analytic
convergence, plateau localization, flat-remainder, or quantile theorem is asserted.\\
\path{ThueMorseComplexProductBridge.lean} &
Total complex finite Thue--Morse sinc and negative-Laplace block identities at every level,
including the removable origin, their quotient forms away from zero, origin simplification
laws, and the exact finite Fourier--Laplace rotation.  No infinite-product convergence,
analytic logarithm, tail estimate, affine lacunary cosine companion, or measure-specific
transform theorem is asserted.\\
\path{SaddleExpansionAlgebra.lean} &
\lean{expCoeff}, its differential equation, uniqueness, naturality, rescaling, addition,
parity, and \lean{HasAsymptoticExpansion}; \cref{integration:sec:poincare-api,integration:sec:exp-coeff}.\\
\path{SaddleLogExpansionAlgebra.lean} &
\lean{logCoeff}, the identity $AB'=A'$, low-order formulas, uniqueness, naturality,
rescaling, and parity; \cref{integration:sec:log-coeff}.\\
\path{SaddleExpansionFiniteRemainder.lean} &
Finite exponent substitution, equality of coefficients below the truncation order,
divisibility of the defect by $X^L$, and the evaluated exact remainder;
\cref{integration:sec:finite-remainder}.\\
\path{SaddleExpansionSmallArgument.lean} &
Exact transfer of full expansions between $t\to\infty$ and $x=2^{-t}\downarrow0$;
\cref{integration:sec:poincare-api}.\\
\path{SaddleExpansionFlat.lean} &
Zero-coefficient expansion of all-orders-flat errors and invariance of every coefficient under
addition or subtraction of such an error; \cref{integration:sec:poincare-api,integration:sec:full-assembly}.\\
\path{GaussianPolynomialContraction.lean} &
Gaussian integrability, odd moments, integration-by-parts recurrence, normalized moments,
linear polynomial contraction, and the whole-line integral identity;
\cref{integration:sec:gaussian-contraction}.\\
\path{GaussianPolynomialTail.lean} &
The factorial coefficient weight and explicit complementary-interval tail bound;
\cref{integration:sec:gaussian-tails}.\\
\path{GaussianPolynomialTailAllOrders.lean} &
The radius $\sqrt{32(N+1)\log b}$ and arbitrary algebraic tail order for bounded polynomial
families; \cref{integration:sec:gaussian-tails}.\\
\path{FabiusSaddle*.lean} &
The closed jets, shifted-Stirling conversion, explicit low-order corrections, top mass jet,
and the central/tail/all-orders Fabius application are mapped in the canonical concordance
of \cref{inverse:sec:module-concordance}; this part supplies only the generic coefficient,
remainder, and Gaussian engines used there.\\
\bottomrule
\end{longtable}
\endgroup

\subsection{Audit matrix against the main article}
\label{integration:app:audit-matrix}

\begingroup\small
\begin{longtable}{@{}p{0.28\textwidth}p{0.18\textwidth}p{0.46\textwidth}@{}}
\toprule
Topic & Status before this companion & Added here\\
\midrule
\endhead
Original scale positivity & Assumed in the source statement & Mean-value proof that the dilation
parameter is automatically positive before its exact value $k=2$ is derived.\\
Total inverse $I$ & Absent & Canonical treatment consolidated in \cref{part:inverse}; this
part retains only the analytic-density obstruction shared with the forward functions.\\
Elementary function class & Absent & Exact inductive closure, totalization convention,
finite-bad-value composition, and dense-open analytic locus.\\
Non-elementarity & Absent & No local agreement on any set with nonempty interior for $F$,
$\RvachevUp$, the signed extension, or the total inverse, including the enlarged admissible
inverse-branch class.\\
Product probability law & Headline proof only & Uniform convergence and measurability,
head--tail pushforward identity, coordinate reflection, and exact dyadic support \([0,1]\).\\
Generic weighted-law naturality, absolute continuity, and support & Not recorded in the pinned companion snapshot &
Almost-everywhere and pointwise head--tail recursion, continuous-linear and scalar
naturality, shifted-law transport, conditional operator/scalar affine fixed-point laws, and
the seven-name compatibility surface; a four-name tranche gives s-finite uniform smoothing,
sharp nontrivial weighted-law absolute continuity, and null singletons; a separate ten-name
tranche gives support equal to the series range, exact signed-real min/max intervals, and
nonnegative specializations.\\
Named geometric-\(q\) law & Not recorded in the pinned companion snapshot &
Exact weights, summability and mass, series and probability law, affine self-similarity,
reflection, absolute continuity, null singletons, and support equal to the series range for
\(|q|<1\), exact support \([0,1]\) for \(0\le q<1\), a named CDF continuous and symmetric
for \(|q|<1\), its exact nonnegative-range tails, and, for \(0<q<1\), refinement integrals,
an explicit compactly supported Radon--Nikodym density, and real-space \(C^\infty\)
regularity.  Series, law, CDF, and density half-base bridges are formal, and
\texttt{geometricSincProduct} is now backed by locally uniform convergence on \(\ComplexNumbers\),
pointwise multipliability and an exact \texttt{HasProd} theorem, and entireness.  For each
fixed complex \(\lVert q\rVert<1\), the single symbol \(a\mapsto(a;q)_\infty\)
likewise has a locally uniform defining product, is entire, has exactly the
factor-zero locus (the reciprocal-power lattice for nonzero \(q\)), and has
only simple zeros.  With
\(z_q(t)=(1-q)t/(2\pi)\), exact named bridges prove
\(\phi_q(t)=e^{\ImaginaryUnit t/2}S_q(z_q(t))
=e^{\ImaginaryUnit t/2}G_q(z_q(t))G_q(-z_q(t))\) throughout \(|q|<1\), including
zero and negative ratios.  The same module proves locally uniform convergence of the full
phase-bearing prefixes on the real frequency line and uniform convergence on every compact
real-frequency set.  Failure of the selected total density formula on \(q\le0\) is formal;
the outer spectral Pochhammer product now also converges locally uniformly on
the complex plane, including at \(q=0\), with exact dyadic Rvachev/Fabius
specializations.  However,
Lean formalization of the paper-level corrected signed-density replacement, a corresponding
complex-frequency locally uniform theorem for the full phase-bearing prefix sequence,
general-\(q\) rapid decay and Fourier inversion, centered/MGF package, closed
higher-derivative and endpoint-flatness wrappers, shape, further transforms, and explicit
Bernoulli-cumulant/Bell-moment formulas and asymptotics remain outside the formal boundary.
In contrast, the finite residual sinc formula, pointwise characteristic-prefix limit,
locally uniform complex sinc product, exact infinite characteristic/Gamma bridges, generic
affine-law factorization, finite characteristic-function iteration, and fixed-point
uniqueness are formal, and specialize throughout \(|q|<1\).\\
The all-power half-base bridge and the transform DDE/adjacent-power hierarchy are likewise
formal for every nonzero \(|q|<1\), including negative ratios; these are separate from the
remaining Fourier and real-space density-derivative layers.\\
Probability--Laplace bridge & Absent & Generic Banach-valued partial and full survival layer
cake, inverse layer cake and exact inverse area, identification of raw, endpoint, and tilted
moments, complete monotonicity, and half-moment interpretation.\\
Atomic-to-step bridge & Compressed & Half-endpoint null-set conversion, exact cell overlap,
shrunken/enlarged interval sandwich, portmanteau, and moving-boundary squeeze.\\
Weighted path solver & Explicitly omitted & Generic semiring theorem, uniqueness, edge-count
decomposition, Fabius specialization, and worked values.\\
Finite-row Toeplitz theorem & Explicitly omitted & Abstract normed-ring proof, exact Fabius
row mass, uniform variation bound, tail-index condition, and corrected prefix rounding.\\
Parameter-dependent full expansions & Only summarized & Exact API, coefficient boundedness,
closure, flat errors, and coordinate transfer.\\
Formal exponential/logarithm & Explicitly omitted & Recurrences, formal differential equations,
uniqueness, low orders, naturality, rescaling, addition, and parity.\\
Finite formal-to-analytic remainder & Absent & Exact divisibility by the first omitted power
and evaluated polynomial remainder.\\
Pointwise analytic finite filters & Absent & Unconditional diagonalization under weighted
sample-series summability, exact finite heads and tails, and Gaussian residual tails for
geometric and geometric-Lagrange rows.  Conditional boundary series, uniform estimates,
signs/rates, and sinc-product instantiation remain outside the formal result.\\
Gaussian contraction/tails & Explicitly omitted & Moment recurrence, contraction functional,
factorial coefficient weight, explicit tail, and all-orders radius.\\
Fabius-specific saddle application & Split across overlapping notebooks & Consolidated once
in \cref{part:inverse}, including the closed jets, two explicit corrections, rigorous
all-orders assembly, exact-phase warning, higher symbolic coefficients, and numerical
claim boundary.\\
\bottomrule
\end{longtable}
\endgroup

The following substantial topics are intentionally not repeated because the main article
already treats them at proof level: construction and uniqueness of $F$ and $\RvachevUp$; the original
Rvachev characterization apart from the redundancy isolated in \cref{integration:sec:scale-positive}; the infinite sinc product and Fourier decay; support, positivity,
shape, and nowhere analyticity; exact moment recurrences and denominator arithmetic; dyadic
rational evaluation, $q$-Pochhammer and $q$-binomial formulas; Thue--Morse interpolation and
iterated difference identities; Legendre expansions, coefficient Parseval and exact tails, the
complete orthogonality of the polynomial-form Legendre blocks, the $A_2$ Legendre energy
identity, exact rational energy approximants and their casted convergence, the Fourier and
dyadic sinc-product energy integrals, least-squares approximation, Rvachev reciprocal-moment
Appell deconvolution, exact shifted-up polynomial synthesis, and literal finite Legendre
translate blocks with their atom-Gram expansion; and the computational API.  Their use as
input above is always signposted.

At compiled source checkpoint \texttt{ca66a1247}, the exact concordance for the newly closed
Legendre energy input is
\path{FabiusLegendreEnergy.lean}.  Its two public definitions are
\lean{rvachevLegendreBlock} and \lean{fabiusSquareEnergy}.  Its eight public theorems are
\lean{intervalIntegral_rvachevLegendreBlock_mul},
\lean{integral_sq_eval_rvachevLegendrePartialSumPolynomial},
\lean{hasSum_rvachevLegendreCoefficient_energy},
\lean{hasSum_rvachevLegendreCoefficient_energy_tail},
\lean{rvachevLegendreSquaredError_partialSum_eq_tsum_tail},
\lean{integral_sq_rvachevUp_eq_two_mul_fabiusSquareEnergy},
\lean{hasSum_fabiusSquareEnergy_legendre}, and
\lean{fabiusSquareEnergy_eq_tsum_legendre}.  This closes complete orthogonality for the
polynomial-form blocks, coefficient Parseval and its exact shifted tail, and the real-variable
Legendre series for $A_2=\int_0^1F(t)^2\Differential t$.

At compiled source checkpoint \texttt{9d5f41c2c}, the exact concordance for the rational
energy layer is \path{FabiusLegendreRationalEnergy.lean}.  Its three executable definitions are
\lean{canonicalRvachevLegendreCoefficientRat}, \lean{fabiusSquareEnergyTermRat}, and
\lean{fabiusSquareEnergyPartialSumRat}.  Its fifteen public theorems are
\lean{canonicalRvachevLegendreCoefficientRat_cast},
\lean{canonicalRvachevLegendreCoefficientRat_zero},
\lean{fabiusSquareEnergyTermRat_cast}, \lean{fabiusSquareEnergyTermRat_nonneg},
\lean{fabiusSquareEnergyTermRat_zero}, \lean{fabiusSquareEnergyPartialSumRat_cast},
\lean{monotone_fabiusSquareEnergyPartialSumRat},
\lean{fabiusSquareEnergyPartialSumRat_pos},
\lean{fabiusSquareEnergyPartialSumRat_zero},
\lean{fabiusSquareEnergyPartialSumRat_one},
\lean{fabiusSquareEnergyPartialSumRat_two},
\lean{fabiusSquareEnergyPartialSumRat_three},
\lean{hasSum_fabiusSquareEnergy_ratCast}, \lean{fabiusSquareEnergy_eq_tsum_ratCast}, and
\lean{tendsto_fabiusSquareEnergyPartialSumRat_cast}.  Consequently every finite
Legendre-energy partial sum has a positive rational representative, those representatives are
monotone, and their real casts converge to $A_2$.  The report's four displayed cutoffs
$N=0,1,2,3$ are certified exactly as $1/4$, $7/18$, $3271/8100$, and
$3246043/8037225$.

At compiled source checkpoint \texttt{b9b240bc0}, the Fourier-energy concordance is
\path{FabiusSquareEnergyFourier.lean}.  It exports no definitions and exactly four theorems:
\lean{integral_norm_sq_rvachevFourier_eq_two_mul_fabiusSquareEnergy},
\lean{fabiusSquareEnergy_eq_integral_Ioi_norm_sq_rvachevFourier},
\lean{fabiusSquareEnergy_eq_integral_Ioi_tprod_sinc_sq}, and
\lean{fabiusSquareEnergy_eq_scaled_integral_Ioi_tprod_sinc_sq}.  They identify the full
real-axis squared Fourier mass with $2A_2$, its positive-half-line integral with $A_2$, and
that half-line integral with the Self-Reconstruction report's unscaled
$\prod_{j\ge0}\SincRad^2(\pi\xi/2^j)$ formula and its $1/(2\pi)$-scaled sine-quotient formula
indexed from one.

The current Rvachev moment--Appell source concordance is
\path{RvachevMomentAppell.lean}.  Its six public definitions are
\lean{rvachevRawMomentRat}, \lean{rvachevReciprocalMomentRat},
\lean{rvachevAppellPolynomialRat}, \lean{rvachevAppellPolynomial}, and
\lean{rvachevDeconvolvedPolynomial}, together with
\lean{rvachevDeconvolutionLinearMap}.  Its thirty-three public theorems are
\lean{rvachevRawMomentRat_zero}, \lean{rvachevRawMomentRat_even},
\lean{rvachevRawMomentRat_odd}, \lean{rvachevReciprocalMomentRat_zero},
\lean{binomialConv_rvachevRawMomentRat_reciprocal},
\lean{rvachevReciprocalMomentRat_eq_completeBellPolynomial},
\lean{monic_rvachevAppellPolynomialRat},
\lean{natDegree_rvachevAppellPolynomialRat},
\lean{rvachevAppellPolynomial_eq_poly_cast},
\lean{monic_rvachevAppellPolynomial},
\lean{natDegree_rvachevAppellPolynomial},
\lean{eval_rvachevAppellPolynomial_add},
\lean{integral_pow_mul_rvachev_eq_rvachevRawMomentRat_cast},
\lean{integral_eval_rvachevAppellPolynomial_add_mul_rvachev},
\lean{rvachevDeconvolutionLinearMap_apply},
\lean{rvachevDeconvolvedPolynomial_zero},
\lean{rvachevDeconvolvedPolynomial_add},
\lean{rvachevDeconvolvedPolynomial_smul},
\lean{rvachevDeconvolvedPolynomial_finsetSum},
\lean{rvachevDeconvolvedPolynomial_C_mul},
\lean{rvachevDeconvolvedPolynomial_monomial},
\lean{rvachevDeconvolvedPolynomial_X_pow},
\lean{coeff_rvachevDeconvolvedPolynomial_natDegree},
\lean{natDegree_rvachevDeconvolvedPolynomial_le},
\lean{natDegree_rvachevDeconvolvedPolynomial},
\lean{leadingCoeff_rvachevDeconvolvedPolynomial},
\lean{rvachevDeconvolvedPolynomial_eq_zero_iff},
\lean{rvachevDeconvolutionLinearMap_injective},
\lean{rvachevDeconvolvedPolynomial_injective}, and
\lean{integral_eval_rvachevDeconvolvedPolynomial_add_mul_rvachev},
\lean{integral_eval_rvachevDeconvolvedPolynomial_sub_mul_rvachev},
\lean{integral_eval_rvachevAppellPolynomial_sub_mul_rvachev}, and
\lean{integral_eval_rvachevAppellPolynomial_mul_rvachev_eq_zero}.  They package the full
rational raw-moment sequence, its formal binomial-convolution reciprocal and complete-Bell
description, rational and real monic Appell families of exact degree, and coefficientwise
polynomial deconvolution as a real linear map preserving zero, addition, scalar
multiplication, finite sums, and constant-polynomial multiplication.  It sends a monomial
and $X^n$ to the correspondingly scaled and unscaled Rvachev--Appell polynomial.  Its
triangular top term is unchanged, so it preserves exact natural degree and leading
coefficient, has trivial kernel, and is injective as both the packaged linear map and the
underlying operation.  Besides additive smoothing, reflection invariance of the density gives
the exact centered convolution of every deconvolved polynomial, the corresponding Appell
identity recovers $x^n$, and every positive-degree Rvachev--Appell polynomial has zero mean.
This reciprocal-moment
Appell family is not the dyadic
$\operatorname{App}_n$ family fixed in the global notation convention above.

The current \path{RvachevPolynomialSynthesis.lean} source exports no public definitions and exactly five public
theorems: \lean{tsum_rvachevDeconvolvedPolynomial_mul_shifted_rvachevUp},
\lean{normalized_tsum_rvachevDeconvolvedPolynomial_mul_shifted_rvachevUp},
\lean{sum_Ioo_rvachevDeconvolvedPolynomial_mul_shifted_rvachevUp}, and
\lean{normalized_sum_Ioo_rvachevDeconvolvedPolynomial_mul_shifted_rvachevUp}, together with
\lean{normalized_tsum_shifted_rvachevDeconvolvedPolynomial_mul_rvachevUp}.  For every
nonzero natural mesh $M$ and polynomial of degree at most $\TwoAdicValuation(M)$, these give the global
raw and normalized \texttt{tsum} identities and, on $[-1,1]$, their exact finite
$k\in(-2M,2M)$ forms.  The fifth theorem gives arbitrary real phase $\theta$:
the normalized sum samples
$\mathscr D_{\RvachevUp}P\bigl(x-(\theta+k)/M\bigr)$ against
$\RvachevUp\bigl((\theta+k)/M\bigr)$ and recovers $P(x)$.  Its exact range is
$\deg P\leq \TwoAdicValuation(M)$ at arbitrary phase.

The compiled \path{RvachevSuperconvergentSynthesis.lean} module exports exactly
one definition and eight theorems.  The definition
\lean{IsRvachevSuperconvergentPhase} selects $0,1/2$ when
$\TwoAdicValuation(M)+1$ is odd and $1/4,3/4$ when it is even.
\lean{isRvachevSuperconvergentPhase_two_pow_iff} gives the dyadic form: even
$N$ selects endpoint phases and odd $N$ quarter phases.  The other seven
theorems are
\lean{tsum_quarter_monomial_eq_integral_of_even_deg},
\lean{tsum_three_quarters_monomial_eq_integral_of_even_deg},
\lean{tsum_shifted_monomial_eq_integral_superconvergent},
\lean{tsum_shifted_polynomial_eq_integral_superconvergent},
\lean{integral_polynomial_mul_rvachevUp_eq_normalized_tsum_superconvergent},
\lean{normalized_tsum_shifted_rvachevDeconvolvedPolynomial_mul_rvachevUp_superconvergent},
and
\lean{normalized_tsum_shifted_rvachevAppellPolynomial_mul_rvachevUp_superconvergent}.
For every nonzero natural $M$, they prove selected-phase exactness through
degree $\TwoAdicValuation(M)+1$, physical quadrature, deconvolved-polynomial
reconstruction, and the Appell monomial specialization.  This arbitrary-$M$
result is stronger than the dyadic-only manuscript claim.  It names exact
phase representatives, not congruence classes or a complete classification,
and proves no maximality, positivity, or rationality theorem.

The downstream \path{LagrangeRvachevMatrix.lean} leaf has the exhaustive
four-definition/six-theorem public surface.  Its definitions and abbreviation are
\lean{rvachevAtomIndexSet}, \lean{RvachevAtomIndex},
\lean{lagrangeRvachevEncoderMatrix}, and \lean{lagrangeRvachevDecoderMatrix}; its theorems
are \lean{lagrangeRvachevEncoderMatrix_nonneg},
\lean{sum_lagrangeRvachevEncoderMatrix_row_eq_one},
\lean{sum_lagrangeRvachevDecoderMatrix_row_eq_one},
\lean{lagrangeRvachevEncoderMatrix_mul_decoderMatrix},
\lean{exists_neg_entry_of_rightInverse_of_row_overlap}, and
\lean{exists_lagrangeRvachevDecoderMatrix_entry_neg_of_row_overlap}.  On the exact finite
atom block \(\{k\in\IntegerNumbers:-2M<k<2M\}\), it packages
\(A_{ik}=M^{-1}\RvachevUp(v_i-k/M)\) and the sampled Lagrange--Rvachev decoder.  Encoder
nonnegativity requires no distinctness or nonzero mesh; encoder row unitality requires
\lean{IsFabius}, (M\ne0), and the selected node in \([-1,1]\); decoder row unitality
requires a nonempty distinct family but no Fabius or mesh-admissibility hypothesis.
Distinct in-range nodes, (M\ne0), and
\(\lvert s\rvert-1\le\TwoAdicValuation(M)\) give the typed equation \(AD=I\).

Generically, entrywise \(A\ge0\), row-unital \(D\), \(AD=I\), and a column strictly positive
in two distinct rows force a negative entry of \(D\); encoder row normalization is not used.
The specialization remains conditional on an actual overlap.  No theorem asserts that all
admissible nodes overlap, or gives the range, kernel, rank, trace, characteristic polynomial,
or Cauchy--Binet data of \(DA\), a geometric closed form, or decoder optimality.

At the focused-build checkpoint of 30 August 2026,
\path{CompositeMeshSharpness.lean} exports one public definition,
\lean{rvachevCombExactThrough}, and seven public theorems:
\lean{exists_shift_tsum_shifted_monomial_ne_integral_nat_real},
\lean{rvachevCombExactThrough_iff_padicValNat},
\lean{rvachevCombExactThrough_iff_pow_two_dvd},
\mbox{\texttt{rvachevCombExactThrough\_two\_pow}},
\lean{two_pow_le_of_rvachevCombExactThrough},
\lean{isLeast_rvachevCombExactThrough}, and
\lean{isLeast_rvachevCombExactThrough_even}.  Universal shifted-comb exactness through
degree $d$ holds for a nonzero natural mesh $M$ exactly when $d\le \TwoAdicValuation(M)$, equivalently
when $2^d\mid M$.  The monomial of degree $\TwoAdicValuation(M)+1$ fails at some real shift, $2^d$ is
the least universally exact mesh, and $4^N$ is the least mesh for the complete polynomial
space through degree $2N$.  These statements do not assert minimality for an individual
Legendre polynomial, one fixed $S_N$, or a target-adapted mesh.

At compiled source checkpoint \texttt{a3854643d},
\path{FabiusLegendreTranslateBlocks.lean} exports six public
definitions: \lean{rvachevLegendreDeconvolutionPolynomial},
\lean{rvachevLegendreScale}, \lean{rvachevLegendreIndexSet},
\lean{rvachevLegendreAtomCoefficient}, \lean{rvachevLegendreTranslateBlock}, and
\lean{rvachevTranslateGram}.  Its seven public theorems are
\lean{natDegree_rvachevLegendreDeconvolutionPolynomial},
\lean{leadingCoeff_rvachevLegendreDeconvolutionPolynomial},
\lean{eval_legendrePolynomial_eq_sum_rvachevUp},
\lean{eval_legendrePolynomial_even_eq_sum_rvachevUp},
\lean{rvachevLegendreTranslateBlock_eq_rvachevLegendreBlock},
\lean{intervalIntegral_rvachevLegendreTranslateBlock_mul}, and
\lean{sum_rvachevLegendreAtomCoefficient_mul_gram}.  For
$Q_d=\mathscr D_{\RvachevUp}P_d$, the first two prove exact natural degree $d$ and the explicit unchanged
leading coefficient $2^{-d}\binom{2d}{d}$.  The remaining theorems specialize the exact finite
synthesis to mesh $2^d$, then to $4^n$ for the even mode, identify each literal finite
translate block with the existing polynomial block on $[-1,1]$, and prove both its complete
orthogonality and the expanded finite atom-Gram identity.

The later \path{RvachevLegendreBiorthogonality.lean} leaf has the exhaustive
one-definition/one-theorem public surface
\lean{rvachevLegendreAnalysisKernel} and
\lean{rvachevLegendreBiorthogonality}.  The former is the canonical translated
Legendre analysis kernel, and the latter gives its exact finite Kronecker
pairing with the deconvolved modes for every nonzero natural mesh satisfying
\(\ell\leq\TwoAdicValuation(M)\), over the literal open block
\(-2M<k<2M\).  It makes \nolinkurl{def:leg-Lambda} and
\nolinkurl{thm:leg-biorthogonality} Exact/Complete, but proves no
Fourier--Bessel representation or matrix-projector identity.

At compiled source checkpoint \texttt{a3854643d}, the focused-build
\path{FabiusLegendreTranslateSeries.lean} exports five public definitions:
\lean{rvachevLegendrePartialSumDeconvolutionPolynomial},
\lean{rvachevLegendrePartialSumAtomCoefficient},
\lean{rvachevLegendrePartialSumTranslateBlock},
\lean{rvachevLegendreTranslateBlockOnInterval}, and
\lean{rvachevLegendrePartialSumTranslateBlockOnInterval}.  Its twenty-five public theorems are
\lean{natDegree_rvachevLegendrePartialSumDeconvolutionPolynomial},
\lean{leadingCoeff_rvachevLegendrePartialSumDeconvolutionPolynomial},
\lean{summable_norm_rvachevLegendreTranslateBlock},
\lean{summable_rvachevLegendreTranslateBlock},
\lean{hasSum_rvachevLegendreTranslateBlock},
\lean{tsum_rvachevLegendreTranslateBlock},
\lean{rvachevLegendrePartialSumDeconvolutionPolynomial_eq_sum},
\lean{rvachevLegendrePartialSumAtomCoefficient_eq_sum},
\lean{eval_rvachevLegendrePartialSumPolynomial_eq_tsum_rvachevUp},
\lean{eval_rvachevLegendrePartialSumPolynomial_eq_sum_rvachevUp},
\lean{rvachevLegendrePartialSumTranslateBlock_eq_eval_partialSumPolynomial},
\lean{rvachevLegendrePartialSumTranslateBlock_eq_sum_translateBlock},
\lean{rvachevLegendreTranslateBlockOnInterval_apply},
\lean{rvachevLegendreTranslateBlockOnInterval_eq_smul},
\lean{summable_norm_rvachevLegendreTranslateBlockOnInterval},
\lean{summable_rvachevLegendreTranslateBlockOnInterval},
\lean{hasSum_rvachevLegendreTranslateBlock_uniform},
\lean{tsum_rvachevLegendreTranslateBlock_uniform},
\lean{rvachevLegendrePartialSumTranslateBlockOnInterval_apply},
\lean{rvachevLegendrePartialSumTranslateBlockOnInterval_eq_eval_partialSumPolynomial},
\lean{rvachevLegendrePartialSumTranslateBlockOnInterval_eq_sum},
\lean{tendsto_rvachevLegendrePartialSumTranslateBlockOnInterval},
\lean{rvachevLegendrePartialSumTranslateBlock_tendstoUniformlyOn},
\lean{tendsto_norm_rvachevLegendrePartialSumTranslateBlockOnInterval_sub}, and
\lean{tendsto_rvachevLegendrePartialSumTranslateBlock}.  They close the literal outer
translate-block series with pointwise and supremum-norm absolute summability and pointwise
and uniform \texttt{HasSum}/\texttt{tsum} forms.  They also identify
$C_N=\mathscr D_{\RvachevUp}S_N$ with the finite sum of deconvolved modes, prove that $C_N$ has the
same natural degree and leading coefficient as $S_N$ (including degenerate visible-degree
drops), expand its common-mesh
coefficients, prove global and finite synthesis at mesh $4^N$, and identify the resulting
function with both the polynomial partial sum and the sum of the separately scaled blocks.
The bundled common-mesh partial trains converge to $\RvachevUp$ in $C([-1,1])$, hence in the
interval supremum norm; raw \texttt{TendstoUniformlyOn}, norm-error-to-zero, and pointwise
forms are exported as well.  No rate, coefficientwise limit, or uniform convergence outside
$[-1,1]$ is asserted.

The nine modules \path{RvachevMomentAppell.lean},
\path{RvachevPolynomialSynthesis.lean},
\path{RvachevSuperconvergentSynthesis.lean},
\path{LagrangeRvachevSynthesis.lean},
\path{LagrangeRvachevMatrix.lean},
\path{CompositeMeshSharpness.lean},
\path{FabiusLegendreTranslateBlocks.lean},
\path{FabiusLegendreTranslateSeries.lean}, and
\path{RvachevLegendreBiorthogonality.lean} have public definition/theorem
inventories $6/33$, $0/5$, $1/8$, $2/7$, $4/6$, $1/7$, $6/7$, $5/25$, and
$1/1$, respectively, for exactly $125$ public declarations.  Universal
whole-space mesh sharpness is now proved, but minimality for an individual Legendre mode,
block, or partial sum is not.  Their APIs also do not prove an analytic
reciprocal-MGF or Appell generating-series identity, a literal differential-series
realization of deconvolution, the displayed low reciprocal coefficients, parity or the
displayed closed formulas for the deconvolved Legendre family, rationality of the atom rows,
range/kernel/rank/trace/characteristic-polynomial or Cauchy--Binet results for the
coefficient-space projector,
equality of the fixed-scale and separately scaled coefficient vectors, an unconditional
$\operatorname{natDegree}(S_N)=2N$ theorem or nonvanishing of the top Legendre coefficient, or the later
refinement, projector, and asymptotic layers.

\subsection{Notation crosswalk}
\label{integration:app:notation}

\begin{longtable}{@{}p{0.18\textwidth}p{0.29\textwidth}>{\raggedright\arraybackslash\footnotesize}p{0.43\textwidth}@{}}
\toprule
This companion & Main-article notation & Principal Lean declaration\\
\midrule
\endhead
$I(y)$ & inverse of bounded $F$ & \lean{fabiusInv}\\
$\mu$ & law of the weighted uniform sum & \lean{weightedSumDistribution}\\
$J_k(s)$ & tilted survival moment & \lean{tiltedSurvivalMoment}\\
$M_k(s)$ & tilted probability/Fabius moment & \lean{fabiusLaplaceMoment}\\
$d_n$ & exact half moment & \lean{halfMoment}\\
$P_w(n)$ & weighted composition/path sum & generic path-sum definition in
\path{FabiusInverseDyadicClosedForm.lean}\\
$\sigma$ & abstract small scale & first argument of \lean{HasAsymptoticExpansion}\\
$a_n$ & formal exponential coefficient & \lean{SaddleExpansion.expCoeff}\\
$b_n$ & formal logarithm coefficient & \lean{SaddleExpansion.logCoeff}\\
$\Gcal[p]$ & normalized Gaussian contraction &
\lean{SaddleExpansion.}\newline\lean{gaussianPolynomialContraction}\\
$Z_n(t)$ & bounded exponent jet $\widetilde{\mathcal Z}_n$ &
\lean{negativeLaplaceBoundedExponentJet}\\
$B_j(t)$ & normalized saddle mass coefficient & \lean{fabiusSaddleMassCoefficient}\\
$A_j(t)$ & logarithmic saddle coefficient & \lean{fabiusSaddleLogCoefficient}\\
$\lambda(x)$ & exact lower-Lambert phase & \lean{fabiusLambertPhase}\\
$\mathcal M(x)$ & corrected sharp Lambert main & \lean{fabiusSharpLambertMain}\\
$\mathcal S(x)$ & normalized real saddle mass/ratio & \lean{fabiusSaddleRatio} and the
kernel-mass bridge\\
\bottomrule
\end{longtable}

\subsection{A compact dependency diagram}
\label{integration:app:dependency}

The proof dependencies of the material unique to this companion can be read from the
following diagram.  An arrow means that the target uses the source as a mathematical input,
not merely as a Lean import.

\begin{center}
\begin{minipage}{0.96\textwidth}\footnotesize
\begin{align*}
 &\text{support + interior positivity + mean-value theorem + derivative law}
   \longrightarrow k>0,\\[0.35ex]
 &\text{strict monotonicity + surjectivity + flatness}
   \longrightarrow \text{total inverse}
   \longrightarrow \text{endpoint steepness},\\[0.35ex]
 &\text{nowhere analyticity + dense elementary analytic locus}
   \longrightarrow \text{strong non-elementarity},\\[0.35ex]
 &\text{product law}
   \longrightarrow \text{survival identity}
   \longrightarrow \text{Laplace/raw/endpoint moment bridge},\\[0.35ex]
 &\text{weak convergence + null boundaries + cell geometry}
   \longrightarrow \text{step-integral convergence},\\[0.35ex]
 &\text{triangular recurrence}
   \longrightarrow \text{weighted paths}
   \longrightarrow \text{composition formula},\\[0.35ex]
 &\text{$q$-binomial row mass + variation + tail indices}
   \longrightarrow \text{discrete Toeplitz limit},\\[0.35ex]
 &\text{closed jets}
   \longrightarrow \text{exponent coefficients}
   \longrightarrow \text{formal exponential}\\[-0.1ex]
 &\hspace{8em}\longrightarrow \text{Gaussian contraction}
   \longrightarrow \text{mass coefficients},\\[0.35ex]
 &\text{mass coefficients}
   \longrightarrow \text{formal logarithm}
   \longrightarrow \text{log coefficients},\\[0.35ex]
 &\text{finite remainders + central Taylor + Gaussian tails + minor arcs}\\[-0.1ex]
 &\hspace{8em}\longrightarrow \text{all-orders saddle mass},\\[0.35ex]
 &\text{all-orders mass + exact reduction + flat tail + coordinate transfer}\\[-0.1ex]
 &\hspace{8em}\longrightarrow \text{full exact-phase expansion}.
\end{align*}
\end{minipage}
\end{center}

\nocite{integration:fabius1966,integration:arias1982,integration:arias05442,integration:arias2018,integration:mainarticle,integration:repo}
\renewcommand{\refname}{Topic references}
\begin{thebibliography}{99}

\bibitem{integration:fabius1966}
J.~Fabius,
\newblock A probabilistic example of a nowhere analytic $C^\infty$-function,
\newblock \emph{Zeitschrift f\"ur Wahrscheinlichkeitstheorie und Verwandte Gebiete}
\textbf{5} (1966), 173--174.

\bibitem{integration:arias1982}
J.~Arias de Reyna,
\newblock On the Fabius function,
\newblock \emph{The American Mathematical Monthly} \textbf{89} (1982), 291--298.

\bibitem{integration:arias05442}
J.~Arias de Reyna,
\newblock An infinitely differentiable function with compact support: Definition and properties,
\newblock arXiv:1702.05442.

\bibitem{integration:arias2018}
J.~Arias de Reyna,
\newblock Arithmetic of the Fabius function,
\newblock arXiv:1702.06487, version 3 (2017).

\bibitem{integration:mainarticle}
V.~Reshetnikov,
\newblock \emph{The Fabius Function and Rvachev's Up-Function: Exact Arithmetic,
Probability, Fourier Analysis, and Corrected Sharp Asymptotics},
\newblock \texttt{ProveIt},
\path{Analysis/FabiusFunction/docs/Fabius_Function_and_Rvachev_Up/}.

\bibitem{integration:repo}
V.~Reshetnikov,
\newblock \texttt{ProveIt}: formal development of the Fabius function and Rvachev's
up-function,
\newblock GitHub repository, snapshot
\texttt{eaea1b9afe2b7ff0b7dd880bd1710e303dec6d80}, 25 August 2026,
\newblock \url{https://github.com/VladimirReshetnikov/ProveIt/tree/eaea1b9afe2b7ff0b7dd880bd1710e303dec6d80/Analysis/FabiusFunction}.

\end{thebibliography}


\renewcommand{\arraystretch}{1}
% END SYNTHESIZED TOPIC: probability, discrete, and generic asymptotic engines


% BEGIN SYNTHESIZED TOPIC: repeated integration and Rvachev's up-function
% Former source SHA-256 values:
% eadcb4b414ac93723a91acbd5062d44340f78134a2cecbc18f7d7ba67eb2c9be
% fad16072df30d9e6eb5df03da57dd217769180730581f1dfd19fa5be0d16b262
\clearpage
\part{Repeated Integration and Rvachev's Up-Function}
\label{part:repeated2}
\begin{boxedremark}
\textbf{Synthesized provenance and status.} This part merges the former
\nolinkurl{Repeated_Integration_and_Rvachev_Up/Repeated_Integration_and_Rvachev_Up.tex}
and
\nolinkurl{Rvachev_Up_from_Repeated_Integration-2/Rvachev_Up_from_Repeated_Integration.tex}.
The second source supplies the canonical sharp threshold \(n\ge r+2\), the first
admissible-stage argument, the complete comparison-kernel data, and the rigorous
all-orders proof.  The first source contributes the 2012 Stack Exchange provenance,
the weak-test and \(L^p\) corollaries, the operator interpretation, and the comparison
with ordinary truncation.  The resulting sharp identities and expansions remain
natural-language proofs and formal obligations.\par\smallskip
Former source SHA-256 values:
\nolinkurl{eadcb4b414ac93723a91acbd5062d44340f78134a2cecbc18f7d7ba67eb2c9be} and
\nolinkurl{fad16072df30d9e6eb5df03da57dd217769180730581f1dfd19fa5be0d16b262}.
\end{boxedremark}
\section*{Topic abstract}
A Math Stack Exchange question asks for a closed form for the sequence
\[
 f_0(x)=\frac12\IndicatorOf{(-1,1)}(x),\qquad
 f_n(x)=2\int_{-1}^{x}\bigl(f_{n-1}(2t+1)-f_{n-1}(2t-1)\bigr)\Differential t.
\]
The first integration produces the piecewise linear tent
$f_1(x)=\PositivePart{1-|x|}$, and the subsequent iterates visually converge very rapidly to
Rvachev's compactly supported $C^\infty$ up-function.  This report gives a complete
finite-stage description and a sharp convergence theorem.

The integral operator is exactly a probabilistic averaging operator: if $g$ is the density
of $Y$ and $U$ is uniform on $[-1,1]$, then the next density is that of $(Y+U)/2$.
Consequently, for $n\ge1$,
\[
 f_n=u_{2^{-1}}\conv u_{2^{-2}}\conv\cdots\conv u_{2^{-n}}\conv u_{2^{-n}},
 \qquad
 u_a(x)=\frac1{2a}\IndicatorOf{[-a,a]}(x).
\]
This identifies $f_n$ as a nonuniform box spline.  Collapsing its general
inclusion--exclusion formula with the Thue--Morse generating polynomial yields, for
$n\ge1$,
\[
 f_n(x)=\frac{2^{n(n+1)/2-1}}{n!}
 \sum_{j=0}^{2^n}\bigl(\tau_n(j)-\tau_n(j-1)\bigr)
 \PositivePart{x+1-\frac{j}{2^{n-1}}}^{n},
\]
where $\tau_n(j)=(-1)^{\BinaryDigitSum(j)}$ on $0\le j<2^n$ and is zero outside that range.
The $n$th derivative is therefore a rescaled Thue--Morse step function.

The principal quantitative result derived here is exact, rather than merely asymptotic.  For every
fixed $r\ge0$ and every $n\ge r+2$,
\[
 \boxed{\displaystyle
 \Norm{f_n^{(r)}-\RvachevUp^{(r)}}_\infty
 =\frac{2^{\binom{r+3}{2}}}{9\,4^n}.}
\]
Thus, once the $r$th derivative has passed its first admissible spline stage, every further
integration divides its uniform error by exactly four.  In the indexing that starts with the
piecewise linear tent $g_0=f_1$, this becomes
\[
 \Norm{g_m-\RvachevUp}_\infty=\frac{2}{9\,4^m}\qquad(m\ge1),
\]
while the exceptional initial error is $\Norm{g_0-\RvachevUp}_\infty=13/72$.
The proof is based on an exact Peano-kernel identity.  The comparison kernel between a
uniform random variable and the up-distributed random variable has both integral and
maximum equal to $1/9$; the finite partial density has disjoint extremal plateaux whose
height is known exactly.  Their combination gives both the upper bound and matching
extremizers.


\section{The problem, the indexing, and the answer}
\label{repeated2:sec:introduction}

This construction was posed independently on Mathematics Stack Exchange in 2012
\cite{repeated2:mse-question}.  Nathan Grigg found the exact finite-stage
Heaviside/Thue--Morse expression \cite{repeated2:mse-answer}; the probability and
box-spline viewpoints below explain its coefficients and, more importantly, expose the
sharp convergence law.  The iteration is not ordinary truncation of the dyadic random
series.  It replaces the omitted infinite tail by one uniform box of exactly the same
support radius, preserving support, mass, symmetry, and mean while leaving a controlled
variance defect.

The recurrence considered in the Math Stack Exchange question
\cite{repeated2:mse-question,repeated2:mse-answer} is
\begin{equation}
 f_0(x)=\frac12\IndicatorOf{(-1,1)}(x),
 \qquad
 f_n(x)=2\int_{-1}^{x}
 \bigl(f_{n-1}(2t+1)-f_{n-1}(2t-1)\bigr)\Differential t.
 \label{repeated2:eq:original-recurrence}
\end{equation}
Changing the values of $f_0$ at $x=\pm1$ changes neither any later iterate nor any of the
norm statements below.  We therefore use whichever endpoint convention is convenient.

There is a small indexing issue in the informal description of the construction.  The
literal seed $f_0$ in \eqref{repeated2:eq:original-recurrence} is piecewise constant; the first iterate
is the piecewise linear function
\begin{equation}
 f_1(x)=\PositivePart{1-|x|}.
 \label{repeated2:eq:tent}
\end{equation}
When one says that the up-function arises by ``repeatedly integrating a piecewise linear
function,'' the natural indexing is therefore
\begin{equation}
 g_m=f_{m+1},\qquad g_0(x)=\PositivePart{1-|x|}.
 \label{repeated2:eq:linear-indexing}
\end{equation}
Both indexings will be retained: $f_n$ matches the question, while $g_m$ matches the
piecewise-linear wording.

\subsection{Main finite-stage formula}

Write $\BinaryDigitSum(j)$ for the number of $1$'s in the binary expansion of $j$.  For $n\ge0$ define
the zero-extended finite Thue--Morse word
\begin{equation}
 \tau_n(j)=
 \begin{cases}
  (-1)^{\BinaryDigitSum(j)},&0\le j<2^n,\\
  0,&\text{otherwise},
 \end{cases}
 \qquad
 \Delta\tau_n(j)=\tau_n(j)-\tau_n(j-1).
 \label{repeated2:eq:tau-definition}
\end{equation}
The accepted Stack Exchange answer gives a formula equivalent to the following one; the
probabilistic and spline derivation in \cref{repeated2:sec:exact-spline} explains why its coefficients
are first differences of the Thue--Morse word.

\begin{theorem}[Exact $n$th iterate]
\label{repeated2:thm:exact-iterate-intro}
For every $n\ge1$ and $x\in\RealNumbers$,
\begin{equation}
 \boxed{\displaystyle
 f_n(x)=\frac{2^{n(n+1)/2-1}}{n!}
 \sum_{j=0}^{2^n}\Delta\tau_n(j)
 \PositivePart{x+1-\frac{j}{2^{n-1}}}^{n}.}
 \label{repeated2:eq:exact-iterate-intro}
\end{equation}
Equivalently,
\begin{equation}
 f_n(x)=\frac{1}{2n!\,2^{n(n-1)/2}}
 \sum_{j=0}^{2^n}\Delta\tau_n(j)
 \PositivePart{2^nx+2^n-2j}^{n}.
 \label{repeated2:eq:exact-iterate-mse-scaling}
\end{equation}
The knots are
\begin{equation}
 -1,-1+2^{1-n},\ldots,1-2^{1-n},1.
 \label{repeated2:eq:iterate-knots}
\end{equation}
On every open knot interval $f_n$ is a polynomial of degree $n$, and globally
$f_n\in C^{n-1}(\RealNumbers)$.
\end{theorem}

In the tent indexing \eqref{repeated2:eq:linear-indexing}, \cref{repeated2:thm:exact-iterate-intro} reads
\begin{equation}
 \boxed{\displaystyle
 g_m(x)=\frac{2^{m(m+3)/2}}{(m+1)!}
 \sum_{j=0}^{2^{m+1}}\Delta\tau_{m+1}(j)
 \PositivePart{x+1-\frac{j}{2^m}}^{m+1}.}
 \label{repeated2:eq:exact-tent-indexing}
\end{equation}
Thus the $m$th repeated integration of the tent is a degree-$(m+1)$ spline on the dyadic
mesh of width $2^{-m}$.

\subsection{Main sharp convergence theorem}

For a bounded function $h$ let $\Norm{h}_\infty=\sup_{x\in\RealNumbers}|h(x)|$.  When $h\in C^r$,
we use the max version of the $C^r$ norm,
\begin{equation}
 \Norm{h}_{C^r}=\max_{0\le j\le r}\Norm{h^{(j)}}_\infty.
 \label{repeated2:eq:Cr-norm}
\end{equation}

\begin{theorem}[Exact derivative error]
\label{repeated2:thm:main-exact-rate-intro}
Let $r\ge0$.
The first iterate belonging to $C^r$ is $f_{r+1}$, and its $r$th-derivative error is
\begin{equation}
 \Norm{f_{r+1}^{(r)}-\RvachevUp^{(r)}}_\infty
 =\frac{13}{72}\,2^{\binom{r+1}{2}}.
 \label{repeated2:eq:first-admissible-error}
\end{equation}
For every $n\ge r+2$,
\begin{equation}
 \boxed{\displaystyle
 \Norm{f_n^{(r)}-\RvachevUp^{(r)}}_\infty
 =\frac{2^{\binom{r+3}{2}}}{9\,4^n}.}
 \label{repeated2:eq:main-exact-rate}
\end{equation}
All of these bounds are attained at the same $2^{r+2}$ points
\begin{equation}
 c_{r+2,j}=-1+\frac{2j+1}{2^{r+2}},
 \qquad 0\le j<2^{r+2}.
 \label{repeated2:eq:error-extremizer-grid}
\end{equation}
At the first admissible stage,
\begin{equation}
 f_{r+1}^{(r)}(c_{r+2,j})-\RvachevUp^{(r)}(c_{r+2,j})
 =(-1)^{\BinaryDigitSum(j)}\frac{13}{72}2^{\binom{r+1}{2}},
 \label{repeated2:eq:signed-first-stage-error}
\end{equation}
and for $n\ge r+2$,
\begin{equation}
 f_n^{(r)}(c_{r+2,j})-\RvachevUp^{(r)}(c_{r+2,j})
 =(-1)^{\BinaryDigitSum(j)}\frac{2^{\binom{r+3}{2}}}{9\,4^n}.
 \label{repeated2:eq:signed-extremizer-error}
\end{equation}
\end{theorem}

The exact quartering begins immediately after the exceptional first $C^r$ spline:
\begin{equation}
 \frac{\Norm{f_{r+2}^{(r)}-\RvachevUp^{(r)}}_\infty}
 {\Norm{f_{r+1}^{(r)}-\RvachevUp^{(r)}}_\infty}=\frac4{13},
 \qquad
 \frac{\Norm{f_{n+1}^{(r)}-\RvachevUp^{(r)}}_\infty}
 {\Norm{f_n^{(r)}-\RvachevUp^{(r)}}_\infty}=\frac14
 \quad(n\ge r+2).
 \label{repeated2:eq:exact-error-ratios}
\end{equation}
For the tent indexing this yields the particularly simple uniform result
\begin{equation}
 \Norm{g_0-\RvachevUp}_\infty=\frac{13}{72},
 \qquad
 \boxed{\displaystyle
 \Norm{g_m-\RvachevUp}_\infty=\frac{2}{9\,4^m}\quad(m\ge1).}
 \label{repeated2:eq:tent-uniform-rate}
\end{equation}
Thus every integration after the first one gains exactly two binary digits in the sharp
uniform error.

\begin{boxedremark}
The approximation in this report is distinct from the centered finite-CDF spline already
discussed in the companion article.  Here $f_n$ is itself a probability \emph{density}, and
the last unresolved tail is replaced by a scaled uniform variable.  The true tail and the
uniform replacement have equal mass and equal mean.  The first surviving discrepancy is
therefore quadratic in the tail scale $2^{-n}$, which explains the factor $4^{-n}$; the
comparison-kernel argument below determines the coefficient exactly.
\end{boxedremark}

\section{The integration operator is a probability operator}
\label{repeated2:sec:probability-operator}

For $a>0$, let
\begin{equation}
 u_a(x)=\frac1{2a}\IndicatorOf{[-a,a]}(x)
 \label{repeated2:eq:uniform-density}
\end{equation}
be the density of a random variable uniform on $[-a,a]$.  Let $T$ denote the operator in
\eqref{repeated2:eq:original-recurrence}:
\begin{equation}
 (Th)(x)=2\int_{-1}^{x}\bigl(h(2t+1)-h(2t-1)\bigr)\Differential t.
 \label{repeated2:eq:T-definition}
\end{equation}

\begin{lemma}[Moving-window and convolution forms]
\label{repeated2:lem:T-convolution}
Suppose $h$ is integrable and supported in $[-1,1]$.  Then
\begin{equation}
 (Th)(x)=\int_{2x-1}^{2x+1}h(s)\Differential s
 =2(h\conv u_1)(2x).
 \label{repeated2:eq:T-moving-window}
\end{equation}
\end{lemma}

\begin{proof}
In the first term of \eqref{repeated2:eq:T-definition}, put $s=2t+1$; in the second, put
$s=2t-1$.  This gives
\[
 (Th)(x)=\int_{-1}^{2x+1}h(s)\Differential s-\int_{-3}^{2x-1}h(s)\Differential s.
\]
The support assumption allows the lower endpoint $-3$ in the second integral to be
replaced by $-1$, after which the difference is the integral over $[2x-1,2x+1]$.
The convolution identity follows directly from the definition of $u_1$.
\end{proof}

\begin{proposition}[Probabilistic meaning]
\label{repeated2:prop:T-probability}
If $h$ is the density of a random variable $Y$, and $U$ is independent and uniform on
$[-1,1]$, then $Th$ is the density of
\begin{equation}
 \frac{Y+U}{2}.
 \label{repeated2:eq:T-random-map}
\end{equation}
In particular, $T$ preserves nonnegativity, total mass one, evenness, and support in
$[-1,1]$.
\end{proposition}

\begin{proof}
The density of $Y+U$ is $h\conv u_1$.  Scaling a random variable by $1/2$ multiplies its
density by $2$ and evaluates it at $2x$, giving exactly \eqref{repeated2:eq:T-moving-window}.
The listed properties are immediate.
\end{proof}

Let $U_1,U_2,\ldots$ be independent random variables, each uniform on $[-1,1]$.
Start with $X_0=U_1$ and define recursively
\begin{equation}
 X_n=\frac{X_{n-1}+U_{n+1}}{2}.
 \label{repeated2:eq:X-recursion}
\end{equation}
Then $f_n$ is the density of $X_n$.

\begin{theorem}[Finite random-sum representation]
\label{repeated2:thm:finite-random-sum}
For every $n\ge1$,
\begin{equation}
 X_n\ \stackrel{d}{=}\
 \sum_{k=1}^{n}2^{-k}U_k+2^{-n}U_{n+1}.
 \label{repeated2:eq:Xn-random-sum}
\end{equation}
Consequently,
\begin{equation}
 f_n=u_{2^{-1}}\conv u_{2^{-2}}\conv\cdots\conv
 u_{2^{-n}}\conv u_{2^{-n}}.
 \label{repeated2:eq:fn-box-convolution}
\end{equation}
\end{theorem}

\begin{proof}
Expanding \eqref{repeated2:eq:X-recursion} gives weights
$2^{-n},2^{-n},2^{-(n-1)},\ldots,2^{-1}$ on $n+1$ independent copies of $U$.
Because the variables are identically distributed, the two copies of $2^{-n}$ may be
placed last, giving \eqref{repeated2:eq:Xn-random-sum}.  Convolution of the corresponding scaled
uniform densities gives \eqref{repeated2:eq:fn-box-convolution}.
\end{proof}

\begin{remark}[The natural Markov chain need not converge pathwise]
The almost-sure convergence obtained by reordering the summands into the nested series is a
coupling of the \emph{laws}.  The forward recursion \eqref{repeated2:eq:X-recursion}
continually inserts fresh noise of size $1/2$, so it should instead be viewed as a Markov
chain converging to its stationary law.  Reordering the i.i.d. summands produces the nested
series coupling; it does not assert pathwise convergence of the original forward chain.
\end{remark}

Several facts that are less transparent in the integral recursion are now immediate:
\begin{equation}
 f_n\ge0,
 \qquad \int_{\RealNumbers}f_n(x)\Differential x=1,
 \qquad f_n(-x)=f_n(x),
 \qquad \SupportOperator f_n=[-1,1].
 \label{repeated2:eq:basic-fn-properties}
\end{equation}
Moreover, for $n\ge1$,
\begin{equation}
 f_n(0)=\int_{-1}^{1}f_{n-1}(s)\Differential s=1.
 \label{repeated2:eq:fn-zero-one}
\end{equation}

\section{The limit is Rvachev's up-function}
\label{repeated2:sec:limit-identification}

Consider the almost surely convergent random series
\begin{equation}
 X=\sum_{k=1}^{\infty}2^{-k}U_k.
 \label{repeated2:eq:up-random-series}
\end{equation}
Its support is $[-1,1]$.  Denote its density by $\RvachevUp$.  Existence and smoothness also follow
from the Fourier product below; they are classical properties of the Rvachev--Arias de
Reyna construction \cite{repeated2:arias1982}.

\subsection{Self-similarity and the differential equation}

Splitting off the first summand in \eqref{repeated2:eq:up-random-series} gives the distributional
fixed-point identity
\begin{equation}
 X\ \stackrel{d}{=}\ \frac{U+X'}2,
 \label{repeated2:eq:up-fixed-point-random}
\end{equation}
where $X'$ is an independent copy of $X$.  By \cref{repeated2:prop:T-probability}, its density is a
fixed point of $T$:
\begin{equation}
 \RvachevUp(x)=\int_{2x-1}^{2x+1}\RvachevUp(s)\Differential s.
 \label{repeated2:eq:up-fixed-point-integral}
\end{equation}
Differentiating gives
\begin{equation}
 \RvachevUp'(x)=2\bigl(\RvachevUp(2x+1)-\RvachevUp(2x-1)\bigr).
 \label{repeated2:eq:up-differential-equation}
\end{equation}
Also $\RvachevUp(0)=\int_{-1}^{1}\RvachevUp=1$.  The even compactly supported smooth solution of
\eqref{repeated2:eq:up-differential-equation} normalized by $\RvachevUp(0)=1$ is Rvachev's up-function.
Thus the repeated-integration limit is not merely analogous to $\RvachevUp$; it is exactly $\RvachevUp$
with the normalization used in the companion article.

\subsection{Fourier product}

Use the non-unitary characteristic Fourier transform
\begin{equation}
 \mathscr F h(t)=\int_{\RealNumbers}h(x)\EulerE^{-\iu tx}\Differential x,
 \qquad \SincRad z=\begin{cases}\sin z/z,&z\ne0,\\1,&z=0.\end{cases}
 \label{repeated2:eq:fourier-convention}
\end{equation}
Then $\mathscr F u_a(t)=\SincRad(at)$.  From \eqref{repeated2:eq:fn-box-convolution},
\begin{equation}
 \mathscr F f_n(t)
 =\SincRad(2^{-n}t)\prod_{k=1}^{n}\SincRad(2^{-k}t)
 =\left(\prod_{k=1}^{n-1}\SincRad(2^{-k}t)\right)\SincRad^2(2^{-n}t).
 \label{repeated2:eq:finite-fourier-product}
\end{equation}
Under the symmetric convention used in the Stack Exchange question, the right side is
multiplied by $(2\pi)^{-1/2}$, exactly reproducing the conjectured formula there.
Taking the limit gives
\begin{equation}
 \mathscr F\RvachevUp(t)=\prod_{k=1}^{\infty}\SincRad(2^{-k}t).
 \label{repeated2:eq:up-fourier-product}
\end{equation}
The product is rapidly decreasing together with all its derivatives; Fourier inversion
therefore gives a $C^\infty$ density.  The product also makes the finite-depth correction
transparent: $f_n$ has the first $n$ factors of the limiting product, with the last factor
repeated once.

\subsection{Relation with the bounded Fabius function}

Let $\xi_k=(U_k+1)/2$, so that the $\xi_k$ are independent uniform variables on $[0,1]$.
Then
\begin{equation}
 Y=\sum_{k=1}^{\infty}2^{-k}\xi_k=\frac{X+1}{2}.
 \label{repeated2:eq:Y-X-relation}
\end{equation}
The cumulative distribution function of $Y$ is the bounded Fabius function $F$.  Hence
\begin{equation}
 \Probability(X\le x)=F\left(\frac{x+1}{2}\right),
 \label{repeated2:eq:X-cdf-F}
\end{equation}
and, using $F(1-y)=1-F(y)$,
\begin{equation}
 \Probability(X>t)=F\left(\frac{1-t}{2}\right),\qquad 0\le t\le1.
 \label{repeated2:eq:X-tail-F}
\end{equation}
On the two halves of the support,
\begin{equation}
 \RvachevUp(x)=
 \begin{cases}
 F(x+1),&-1\le x\le0,\\
 F(1-x),&0\le x\le1.
 \end{cases}
 \label{repeated2:eq:up-F-relation}
\end{equation}
This is exactly the normalization used in the existing exposition.

\section{Exact finite splines and the Thue--Morse collapse}
\label{repeated2:sec:exact-spline}

The convolution formula \eqref{repeated2:eq:fn-box-convolution} already gives a closed form in the
language of box splines.  To obtain the explicit truncated-power expression, we first
record the general one-dimensional inclusion--exclusion formula.

\subsection{A general nonuniform box-spline formula}

\begin{lemma}[Density of a sum of centered uniforms]
\label{repeated2:lem:general-box-spline}
Let $V_i$ be independent and uniform on $[-a_i,a_i]$, where $a_i>0$ and
$1\le i\le m$.  Put $A=a_1+\cdots+a_m$.  The density $b_{\bm a}$ of
$V_1+\cdots+V_m$ is
\begin{equation}
 b_{\bm a}(x)=
 \frac{1}{2^m(m-1)!\prod_{i=1}^{m}a_i}
 \sum_{\varepsilon\in\{0,1\}^m}(-1)^{|\varepsilon|}
 \PositivePart{x+A-2\sum_{i=1}^{m}\varepsilon_i a_i}^{m-1}.
 \label{repeated2:eq:general-box-spline}
\end{equation}
\end{lemma}

\begin{proof}
Translate $V_i$ to $W_i=V_i+a_i$, which is uniform on $[0,2a_i]$.  The distributional
derivative of the density of $W_i$ is
$(\delta_0-\delta_{2a_i})/(2a_i)$.  After convolving the $m$ factors and integrating
$m-1$ times from the left endpoint, one obtains the usual inclusion--exclusion formula
for the density of $W_1+\cdots+W_m$.  Translating the sum back by $A$ gives
\eqref{repeated2:eq:general-box-spline}.  Equivalently, one may verify the formula by taking
Laplace transforms.
\end{proof}

For $f_n$, the $m=n+1$ half-widths are
\begin{equation}
 2^{-1},2^{-2},\ldots,2^{-(n-1)},2^{-n},2^{-n}.
 \label{repeated2:eq:fn-halfwidths}
\end{equation}
Their sum is $1$, and their product is $2^{-n(n+3)/2}$.  Thus the scalar prefactor in
\eqref{repeated2:eq:general-box-spline} becomes
\begin{equation}
 \frac{2^{n(n+1)/2-1}}{n!}.
 \label{repeated2:eq:box-prefactor}
\end{equation}
It remains to combine subsets that produce the same knot.

\subsection{The subset-sum polynomial}

After multiplication by $2^{n-1}$, the possible increments $2a_i$ in
\eqref{repeated2:eq:general-box-spline} are
\begin{equation}
 2^{n-1},2^{n-2},\ldots,2,1,1.
 \label{repeated2:eq:integer-increments}
\end{equation}
The signed subset-sum generating polynomial is therefore
\begin{align}
 B_n(z)
 &=(1-z)^2\prod_{r=1}^{n-1}(1-z^{2^r}) \notag\\
 &=(1-z)\prod_{r=0}^{n-1}(1-z^{2^r}).
 \label{repeated2:eq:subset-polynomial}
\end{align}
The classical finite Thue--Morse identity is
\begin{equation}
 \prod_{r=0}^{n-1}(1-z^{2^r})
 =\sum_{j=0}^{2^n-1}(-1)^{\BinaryDigitSum(j)}z^j
 =\sum_{j\in\IntegerNumbers}\tau_n(j)z^j.
 \label{repeated2:eq:thue-morse-product}
\end{equation}

\begin{theorem}[Finite complex Boolean-cube and Thue--Morse transforms; \Formalized]
\label{repeated2:thm:complex-thue-morse-transform}
Let $m\in\NaturalNumbers$ and $u,a,z\in\ComplexNumbers$, and write
$\varepsilon_n=(-1)^{\BinaryDigitSum(n)}$.  Then
\begin{align}
 \sum_{n=0}^{2^m-1}u^{\BinaryDigitSum(n)}\exp(nz)
 &=\prod_{j=0}^{m-1}\left(1+u\exp(2^jz)\right),
 \label{repeated2:eq:weighted-complex-block}\\
 \sum_{n=0}^{2^m-1}u^{\BinaryDigitSum(n)}\exp(a+nz)
 &=\exp(a)\prod_{j=0}^{m-1}\left(1+u\exp(2^jz)\right),
 \label{repeated2:eq:weighted-complex-block-affine}\\
 \sum_{n=0}^{2^m-1}\varepsilon_n\exp(nz)
 &=\prod_{j=0}^{m-1}\left(1-\exp(2^jz)\right),
 \label{repeated2:eq:signed-complex-block}\\
 \sum_{n=0}^{2^m-1}\varepsilon_n\exp(a+nz)
 &=\exp(a)\prod_{j=0}^{m-1}\left(1-\exp(2^jz)\right).
 \label{repeated2:eq:signed-complex-block-affine}
\end{align}
The statement includes $m=0$: the sum then contains the unique zero-bit word $n=0$ and
the product is empty.  The unshifted identities read $1=1$, while the affine identities
read $\exp(a)=\exp(a)$.
\end{theorem}

\begin{proof}
The four identities are, in order, the Lean declarations
\lean{sum_range_two_pow_binaryWeight_cexp},
\lean{sum_range_two_pow_binaryWeight_cexp_affine},
\lean{sum_range_two_pow_thueMorseSign_cexp}, and
\lean{sum_range_two_pow_thueMorseSign_cexp_affine} in
\path{ThueMorseComplexExponential.lean}.  Conceptually, the binary expansion of
$n<2^m$ identifies $n$ with a subset of $\{0,\ldots,m-1\}$; multiplying the independent
choices produces \eqref{repeated2:eq:weighted-complex-block}.  The affine identity factors
out $\exp(a)$, and the two signed identities replace the Hamming-weight marker by the
Thue--Morse sign.
\end{proof}

\begin{boxedremark}
\textbf{Exact formal boundary.}
These are finite identities valid at every complex argument, so they require no convergence
or reality hypothesis.  The cited declarations are the raw exponential layer; the complex
half-angle conversion and total sinc and negative-Laplace prefixes, including their removable
origins, are separately formal in \path{ThueMorseComplexProductBridge.lean}.  Neither finite
module proves an infinite product or exponential-generating-function convergence statement,
nor do they cover a parity word with position-dependent digit weights.
\end{boxedremark}

Multiplication by $1-z$ takes first differences of the coefficient sequence:
\begin{equation}
 B_n(z)=\sum_{j=0}^{2^n}\Delta\tau_n(j)z^j.
 \label{repeated2:eq:first-difference-generating}
\end{equation}
This explains the otherwise rather mysterious coefficient
$c_n(j)-c_n(j-1)$ in the Stack Exchange answer: the duplicated smallest uniform interval
contributes the extra factor $1-z$, turning the Thue--Morse word into its first difference.

\begin{proof}[Proof of \cref{repeated2:thm:exact-iterate-intro}]
Insert the widths \eqref{repeated2:eq:fn-halfwidths} into \cref{repeated2:lem:general-box-spline}.  The
subset sum indexed by $\varepsilon$ equals $j/2^{n-1}$ after grouping terms with the same
integer $j$; by \eqref{repeated2:eq:first-difference-generating}, the total sign of that group is
$\Delta\tau_n(j)$.  The scalar is \eqref{repeated2:eq:box-prefactor}.  This gives
\eqref{repeated2:eq:exact-iterate-intro}.  Factoring $2^{-n}$ from each truncated power gives
\eqref{repeated2:eq:exact-iterate-mse-scaling}.
\end{proof}

\subsection{Derivatives, knots, and exact smoothness}

For $0\le r\le n-1$, differentiating \eqref{repeated2:eq:exact-iterate-intro} gives
\begin{equation}
 f_n^{(r)}(x)=
 \frac{2^{n(n+1)/2-1}}{(n-r)!}
 \sum_{j=0}^{2^n}\Delta\tau_n(j)
 \PositivePart{x+1-\frac{j}{2^{n-1}}}^{n-r}.
 \label{repeated2:eq:finite-spline-derivatives}
\end{equation}
The formula is also valid away from the knots for $r=n$.  If
\begin{equation}
 -1+\frac{j}{2^{n-1}}<x<-1+\frac{j+1}{2^{n-1}},
 \qquad 0\le j<2^n,
 \label{repeated2:eq:open-cell}
\end{equation}
then the first-difference sum telescopes and gives
\begin{equation}
 \boxed{\displaystyle
 f_n^{(n)}(x)=2^{n(n+1)/2-1}\tau_n(j).}
 \label{repeated2:eq:top-derivative-thue-morse}
\end{equation}
Thus the highest ordinary derivative is literally a Thue--Morse step word.  Because this
step function has jumps at the active knots, $f_n$ is $C^{n-1}$ but not $C^n$ globally.  Each
application of $T$ adds exactly one degree and exactly one order of classical smoothness.

\subsection{The first iterates}

The first three functions are
\begin{align}
 f_0(x)&=\frac12\IndicatorOf{(-1,1)}(x),
 \label{repeated2:eq:f0-explicit}\\
 f_1(x)&=\PositivePart{1-|x|},
 \label{repeated2:eq:f1-explicit}\\
 f_2(x)&=
 \begin{cases}
 1-2x^2,&|x|\le\tfrac12,\\
 2(1-|x|)^2,&\tfrac12\le|x|\le1,\\
 0,&|x|\ge1.
 \end{cases}
 \label{repeated2:eq:f2-explicit}
\end{align}
The next one already displays the full Thue--Morse organization compactly:
\begin{equation}
 f_3(x)=\frac{16}{3}\sum_{j=0}^{8}\Delta\tau_3(j)
 \PositivePart{x+1-\frac j4}^{3}.
 \label{repeated2:eq:f3-truncated-power}
\end{equation}

\begin{figure}[htbp]
\centering
\begin{tikzpicture}
\begin{axis}[
  width=0.91\textwidth,
  height=0.42\textwidth,
  xmin=-1.12,xmax=1.12,
  ymin=-0.04,ymax=1.08,
  axis lines=middle,
  xlabel={$x$},
  ylabel={$f_n(x)$},
  xtick={-1,-0.5,0,0.5,1},
  ytick={0,0.5,1},
  grid=both,
  major grid style={draw=rulegray},
  minor grid style={draw=rulegray!45},
  legend style={at={(0.5,-0.18)},anchor=north,legend columns=3,draw=none},
  samples=501
]
\addplot[thick,gray,densely dotted,domain=-1.1:1.1]
  {abs(x)<1 ? 0.5 : 0};
\addlegendentry{$f_0$}
\addplot[thick,dashed,domain=-1.1:1.1]
  {max(1-abs(x),0)};
\addlegendentry{$f_1$}
\addplot[thick,linkblue,domain=-1.1:1.1]
  {abs(x)<=0.5 ? 1-2*x^2 : (abs(x)<=1 ? 2*(1-abs(x))^2 : 0)};
\addlegendentry{$f_2$}
\end{axis}
\end{tikzpicture}
\caption{The rectangle, the piecewise linear tent, and the quadratic spline.  Later
iterates add one polynomial degree and one smoothness order at each step.}
\label{repeated2:fig:first-iterates}
\end{figure}

\begin{figure}[htbp]
\centering
\begin{tikzpicture}
\begin{axis}[
  width=0.90\textwidth,
  height=0.47\textwidth,
  xmin=-1,xmax=1,ymin=0,ymax=1.05,
  xlabel={$x$},ylabel={density},
  xtick={-1,-0.75,-0.5,-0.25,0,0.25,0.5,0.75,1},
  ytick={0,0.25,0.5,0.75,1},
  grid=both,
  minor tick num=1,
  tick label style={font=\small},
  label style={font=\small},
  legend style={font=\small,at={(0.5,-0.20)},anchor=north,legend columns=4},
]
\addplot+[mark=none,line width=1.15pt] coordinates {
(-1.00000000,0.0000000000)
(-0.98437500,0.0000000021)
(-0.96875000,0.0000005727)
(-0.95312500,0.0000104681)
(-0.93750000,0.0000689621)
(-0.92187500,0.0002681716)
(-0.90625000,0.0007601198)
(-0.89062500,0.0017496732)
(-0.87500000,0.0034722156)
(-0.85937500,0.0061713187)
(-0.84375000,0.0100905539)
(-0.82812500,0.0154653054)
(-0.81250000,0.0225004393)
(-0.79687500,0.0313479791)
(-0.78125000,0.0421000415)
(-0.76562500,0.0547959004)
(-0.75000000,0.0694443120)
(-0.73437500,0.0860458408)
(-0.71875000,0.1045999223)
(-0.70312500,0.1250978002)
(-0.68750000,0.1475002009)
(-0.67187500,0.1717150073)
(-0.65625000,0.1975901963)
(-0.64062500,0.2249209014)
(-0.62500000,0.2534717388)
(-0.60937500,0.2829991368)
(-0.59375000,0.3132595238)
(-0.57812500,0.3440175159)
(-0.56250000,0.3750682468)
(-0.54687500,0.4062596932)
(-0.53125000,0.4374997382)
(-0.51562500,0.4687491080)
(-0.50000000,0.4999990463)
(-0.48437500,0.5312489847)
(-0.46875000,0.5624983544)
(-0.45312500,0.5937383994)
(-0.43750000,0.6249298458)
(-0.42187500,0.6559805767)
(-0.40625000,0.6867385689)
(-0.39062500,0.7169989559)
(-0.37500000,0.7465263539)
(-0.35937500,0.7750771912)
(-0.34375000,0.8024078964)
(-0.32812500,0.8282830853)
(-0.31250000,0.8524978917)
(-0.29687500,0.8749002924)
(-0.28125000,0.8953981704)
(-0.26562500,0.9139522519)
(-0.25000000,0.9305537806)
(-0.23437500,0.9452021923)
(-0.21875000,0.9578980512)
(-0.20312500,0.9686501136)
(-0.18750000,0.9774976533)
(-0.17187500,0.9845327873)
(-0.15625000,0.9899075387)
(-0.14062500,0.9938267740)
(-0.12500000,0.9965258770)
(-0.10937500,0.9982484195)
(-0.09375000,0.9992379728)
(-0.07812500,0.9997299211)
(-0.06250000,0.9999291306)
(-0.04687500,0.9999876246)
(-0.03125000,0.9999975200)
(-0.01562500,0.9999980906)
(0.00000000,0.9999980927)
(0.01562500,0.9999980906)
(0.03125000,0.9999975200)
(0.04687500,0.9999876246)
(0.06250000,0.9999291306)
(0.07812500,0.9997299211)
(0.09375000,0.9992379728)
(0.10937500,0.9982484195)
(0.12500000,0.9965258770)
(0.14062500,0.9938267740)
(0.15625000,0.9899075387)
(0.17187500,0.9845327873)
(0.18750000,0.9774976533)
(0.20312500,0.9686501136)
(0.21875000,0.9578980512)
(0.23437500,0.9452021923)
(0.25000000,0.9305537806)
(0.26562500,0.9139522519)
(0.28125000,0.8953981704)
(0.29687500,0.8749002924)
(0.31250000,0.8524978917)
(0.32812500,0.8282830853)
(0.34375000,0.8024078964)
(0.35937500,0.7750771912)
(0.37500000,0.7465263539)
(0.39062500,0.7169989559)
(0.40625000,0.6867385689)
(0.42187500,0.6559805767)
(0.43750000,0.6249298458)
(0.45312500,0.5937383994)
(0.46875000,0.5624983544)
(0.48437500,0.5312489847)
(0.50000000,0.4999990463)
(0.51562500,0.4687491080)
(0.53125000,0.4374997382)
(0.54687500,0.4062596932)
(0.56250000,0.3750682468)
(0.57812500,0.3440175159)
(0.59375000,0.3132595238)
(0.60937500,0.2829991368)
(0.62500000,0.2534717388)
(0.64062500,0.2249209014)
(0.65625000,0.1975901963)
(0.67187500,0.1717150073)
(0.68750000,0.1475002009)
(0.70312500,0.1250978002)
(0.71875000,0.1045999223)
(0.73437500,0.0860458408)
(0.75000000,0.0694443120)
(0.76562500,0.0547959004)
(0.78125000,0.0421000415)
(0.79687500,0.0313479791)
(0.81250000,0.0225004393)
(0.82812500,0.0154653054)
(0.84375000,0.0100905539)
(0.85937500,0.0061713187)
(0.87500000,0.0034722156)
(0.89062500,0.0017496732)
(0.90625000,0.0007601198)
(0.92187500,0.0002681716)
(0.93750000,0.0000689621)
(0.95312500,0.0000104681)
(0.96875000,0.0000005727)
(0.98437500,0.0000000021)
(1.00000000,0.0000000000)
};
\addlegendentry{$\RvachevUp$}
\addplot+[mark=none,densely dashed] coordinates {
(-1.00000000,0.0000000000)
(-0.98437500,0.0156240463)
(-0.96875000,0.0312490463)
(-0.95312500,0.0468740463)
(-0.93750000,0.0624990463)
(-0.92187500,0.0781240463)
(-0.90625000,0.0937490463)
(-0.89062500,0.1093740463)
(-0.87500000,0.1249990463)
(-0.85937500,0.1406240463)
(-0.84375000,0.1562490463)
(-0.82812500,0.1718740463)
(-0.81250000,0.1874990463)
(-0.79687500,0.2031240463)
(-0.78125000,0.2187490463)
(-0.76562500,0.2343740463)
(-0.75000000,0.2499990463)
(-0.73437500,0.2656240463)
(-0.71875000,0.2812490463)
(-0.70312500,0.2968740463)
(-0.68750000,0.3124990463)
(-0.67187500,0.3281240463)
(-0.65625000,0.3437490463)
(-0.64062500,0.3593740463)
(-0.62500000,0.3749990463)
(-0.60937500,0.3906240463)
(-0.59375000,0.4062490463)
(-0.57812500,0.4218740463)
(-0.56250000,0.4374990463)
(-0.54687500,0.4531240463)
(-0.53125000,0.4687490463)
(-0.51562500,0.4843740463)
(-0.50000000,0.4999990463)
(-0.48437500,0.5156240463)
(-0.46875000,0.5312490463)
(-0.45312500,0.5468740463)
(-0.43750000,0.5624990463)
(-0.42187500,0.5781240463)
(-0.40625000,0.5937490463)
(-0.39062500,0.6093740463)
(-0.37500000,0.6249990463)
(-0.35937500,0.6406240463)
(-0.34375000,0.6562490463)
(-0.32812500,0.6718740463)
(-0.31250000,0.6874990463)
(-0.29687500,0.7031240463)
(-0.28125000,0.7187490463)
(-0.26562500,0.7343740463)
(-0.25000000,0.7499990463)
(-0.23437500,0.7656240463)
(-0.21875000,0.7812490463)
(-0.20312500,0.7968740463)
(-0.18750000,0.8124990463)
(-0.17187500,0.8281240463)
(-0.15625000,0.8437490463)
(-0.14062500,0.8593740463)
(-0.12500000,0.8749990463)
(-0.10937500,0.8906240463)
(-0.09375000,0.9062490463)
(-0.07812500,0.9218740463)
(-0.06250000,0.9374990463)
(-0.04687500,0.9531240463)
(-0.03125000,0.9687490463)
(-0.01562500,0.9843740463)
(0.00000000,0.9999980927)
(0.01562500,0.9843740463)
(0.03125000,0.9687490463)
(0.04687500,0.9531240463)
(0.06250000,0.9374990463)
(0.07812500,0.9218740463)
(0.09375000,0.9062490463)
(0.10937500,0.8906240463)
(0.12500000,0.8749990463)
(0.14062500,0.8593740463)
(0.15625000,0.8437490463)
(0.17187500,0.8281240463)
(0.18750000,0.8124990463)
(0.20312500,0.7968740463)
(0.21875000,0.7812490463)
(0.23437500,0.7656240463)
(0.25000000,0.7499990463)
(0.26562500,0.7343740463)
(0.28125000,0.7187490463)
(0.29687500,0.7031240463)
(0.31250000,0.6874990463)
(0.32812500,0.6718740463)
(0.34375000,0.6562490463)
(0.35937500,0.6406240463)
(0.37500000,0.6249990463)
(0.39062500,0.6093740463)
(0.40625000,0.5937490463)
(0.42187500,0.5781240463)
(0.43750000,0.5624990463)
(0.45312500,0.5468740463)
(0.46875000,0.5312490463)
(0.48437500,0.5156240463)
(0.50000000,0.4999990463)
(0.51562500,0.4843740463)
(0.53125000,0.4687490463)
(0.54687500,0.4531240463)
(0.56250000,0.4374990463)
(0.57812500,0.4218740463)
(0.59375000,0.4062490463)
(0.60937500,0.3906240463)
(0.62500000,0.3749990463)
(0.64062500,0.3593740463)
(0.65625000,0.3437490463)
(0.67187500,0.3281240463)
(0.68750000,0.3124990463)
(0.70312500,0.2968740463)
(0.71875000,0.2812490463)
(0.73437500,0.2656240463)
(0.75000000,0.2499990463)
(0.76562500,0.2343740463)
(0.78125000,0.2187490463)
(0.79687500,0.2031240463)
(0.81250000,0.1874990463)
(0.82812500,0.1718740463)
(0.84375000,0.1562490463)
(0.85937500,0.1406240463)
(0.87500000,0.1249990463)
(0.89062500,0.1093740463)
(0.90625000,0.0937490463)
(0.92187500,0.0781240463)
(0.93750000,0.0624990463)
(0.95312500,0.0468740463)
(0.96875000,0.0312490463)
(0.98437500,0.0156240463)
(1.00000000,0.0000000000)
};
\addlegendentry{$f_1$}
\addplot+[mark=none,dashdotted] coordinates {
(-1.00000000,0.0000000000)
(-0.98437500,0.0004882515)
(-0.96875000,0.0019530654)
(-0.95312500,0.0043944418)
(-0.93750000,0.0078123808)
(-0.92187500,0.0122068822)
(-0.90625000,0.0175779462)
(-0.89062500,0.0239255726)
(-0.87500000,0.0312497616)
(-0.85937500,0.0395505130)
(-0.84375000,0.0488278270)
(-0.82812500,0.0590817034)
(-0.81250000,0.0703121424)
(-0.79687500,0.0825191438)
(-0.78125000,0.0957027078)
(-0.76562500,0.1098628342)
(-0.75000000,0.1249995232)
(-0.73437500,0.1411127746)
(-0.71875000,0.1582025886)
(-0.70312500,0.1762689650)
(-0.68750000,0.1953119040)
(-0.67187500,0.2153314054)
(-0.65625000,0.2363274694)
(-0.64062500,0.2583000958)
(-0.62500000,0.2812492847)
(-0.60937500,0.3051750362)
(-0.59375000,0.3300773501)
(-0.57812500,0.3559562266)
(-0.56250000,0.3828116655)
(-0.54687500,0.4106436670)
(-0.53125000,0.4394522309)
(-0.51562500,0.4692373574)
(-0.50000000,0.4999990463)
(-0.48437500,0.5307607353)
(-0.46875000,0.5605458617)
(-0.45312500,0.5893544257)
(-0.43750000,0.6171864271)
(-0.42187500,0.6440418661)
(-0.40625000,0.6699207425)
(-0.39062500,0.6948230565)
(-0.37500000,0.7187488079)
(-0.35937500,0.7416979969)
(-0.34375000,0.7636706233)
(-0.32812500,0.7846666872)
(-0.31250000,0.8046861887)
(-0.29687500,0.8237291276)
(-0.28125000,0.8417955041)
(-0.26562500,0.8588853180)
(-0.25000000,0.8749985695)
(-0.23437500,0.8901352584)
(-0.21875000,0.9042953849)
(-0.20312500,0.9174789488)
(-0.18750000,0.9296859503)
(-0.17187500,0.9409163892)
(-0.15625000,0.9511702657)
(-0.14062500,0.9604475796)
(-0.12500000,0.9687483311)
(-0.10937500,0.9760725200)
(-0.09375000,0.9824201465)
(-0.07812500,0.9877912104)
(-0.06250000,0.9921857119)
(-0.04687500,0.9956036508)
(-0.03125000,0.9980450273)
(-0.01562500,0.9995098412)
(0.00000000,0.9999980927)
(0.01562500,0.9995098412)
(0.03125000,0.9980450273)
(0.04687500,0.9956036508)
(0.06250000,0.9921857119)
(0.07812500,0.9877912104)
(0.09375000,0.9824201465)
(0.10937500,0.9760725200)
(0.12500000,0.9687483311)
(0.14062500,0.9604475796)
(0.15625000,0.9511702657)
(0.17187500,0.9409163892)
(0.18750000,0.9296859503)
(0.20312500,0.9174789488)
(0.21875000,0.9042953849)
(0.23437500,0.8901352584)
(0.25000000,0.8749985695)
(0.26562500,0.8588853180)
(0.28125000,0.8417955041)
(0.29687500,0.8237291276)
(0.31250000,0.8046861887)
(0.32812500,0.7846666872)
(0.34375000,0.7636706233)
(0.35937500,0.7416979969)
(0.37500000,0.7187488079)
(0.39062500,0.6948230565)
(0.40625000,0.6699207425)
(0.42187500,0.6440418661)
(0.43750000,0.6171864271)
(0.45312500,0.5893544257)
(0.46875000,0.5605458617)
(0.48437500,0.5307607353)
(0.50000000,0.4999990463)
(0.51562500,0.4692373574)
(0.53125000,0.4394522309)
(0.54687500,0.4106436670)
(0.56250000,0.3828116655)
(0.57812500,0.3559562266)
(0.59375000,0.3300773501)
(0.60937500,0.3051750362)
(0.62500000,0.2812492847)
(0.64062500,0.2583000958)
(0.65625000,0.2363274694)
(0.67187500,0.2153314054)
(0.68750000,0.1953119040)
(0.70312500,0.1762689650)
(0.71875000,0.1582025886)
(0.73437500,0.1411127746)
(0.75000000,0.1249995232)
(0.76562500,0.1098628342)
(0.78125000,0.0957027078)
(0.79687500,0.0825191438)
(0.81250000,0.0703121424)
(0.82812500,0.0590817034)
(0.84375000,0.0488278270)
(0.85937500,0.0395505130)
(0.87500000,0.0312497616)
(0.89062500,0.0239255726)
(0.90625000,0.0175779462)
(0.92187500,0.0122068822)
(0.93750000,0.0078123808)
(0.95312500,0.0043944418)
(0.96875000,0.0019530654)
(0.98437500,0.0004882515)
(1.00000000,0.0000000000)
};
\addlegendentry{$f_2$}
\addplot+[mark=none,dotted,line width=1.1pt] coordinates {
(-1.00000000,0.0000000000)
(-0.98437500,0.0000012715)
(-0.96875000,0.0000203447)
(-0.95312500,0.0001029958)
(-0.93750000,0.0003255184)
(-0.92187500,0.0007947237)
(-0.90625000,0.0016479408)
(-0.89062500,0.0030530161)
(-0.87500000,0.0052083135)
(-0.85937500,0.0083401715)
(-0.84375000,0.0126749287)
(-0.82812500,0.0184109489)
(-0.81250000,0.0257160788)
(-0.79687500,0.0347276471)
(-0.78125000,0.0455524651)
(-0.76562500,0.0582668266)
(-0.75000000,0.0729165077)
(-0.73437500,0.0895167670)
(-0.71875000,0.1080523459)
(-0.70312500,0.1284774683)
(-0.68750000,0.1507158404)
(-0.67187500,0.1746606509)
(-0.65625000,0.2001745710)
(-0.64062500,0.2270897543)
(-0.62500000,0.2552078366)
(-0.60937500,0.2843024796)
(-0.59375000,0.3141473448)
(-0.57812500,0.3445440681)
(-0.56250000,0.3753248031)
(-0.54687500,0.4063522209)
(-0.53125000,0.4375195103)
(-0.51562500,0.4687503775)
(-0.50000000,0.4999990463)
(-0.48437500,0.5312477152)
(-0.46875000,0.5624785824)
(-0.45312500,0.5936458717)
(-0.43750000,0.6246732896)
(-0.42187500,0.6554540246)
(-0.40625000,0.6858507479)
(-0.39062500,0.7156956130)
(-0.37500000,0.7447902560)
(-0.35937500,0.7729083384)
(-0.34375000,0.7998235216)
(-0.32812500,0.8253374417)
(-0.31250000,0.8492822523)
(-0.29687500,0.8715206244)
(-0.28125000,0.8919457467)
(-0.26562500,0.9104813256)
(-0.25000000,0.9270815849)
(-0.23437500,0.9417312660)
(-0.21875000,0.9544456275)
(-0.20312500,0.9652704456)
(-0.18750000,0.9742820139)
(-0.17187500,0.9815871437)
(-0.15625000,0.9873231640)
(-0.14062500,0.9916579211)
(-0.12500000,0.9947897792)
(-0.10937500,0.9969450766)
(-0.09375000,0.9983501518)
(-0.07812500,0.9992033689)
(-0.06250000,0.9996725743)
(-0.04687500,0.9998950969)
(-0.03125000,0.9999777479)
(-0.01562500,0.9999968211)
(0.00000000,0.9999980927)
(0.01562500,0.9999968211)
(0.03125000,0.9999777479)
(0.04687500,0.9998950969)
(0.06250000,0.9996725743)
(0.07812500,0.9992033689)
(0.09375000,0.9983501518)
(0.10937500,0.9969450766)
(0.12500000,0.9947897792)
(0.14062500,0.9916579211)
(0.15625000,0.9873231640)
(0.17187500,0.9815871437)
(0.18750000,0.9742820139)
(0.20312500,0.9652704456)
(0.21875000,0.9544456275)
(0.23437500,0.9417312660)
(0.25000000,0.9270815849)
(0.26562500,0.9104813256)
(0.28125000,0.8919457467)
(0.29687500,0.8715206244)
(0.31250000,0.8492822523)
(0.32812500,0.8253374417)
(0.34375000,0.7998235216)
(0.35937500,0.7729083384)
(0.37500000,0.7447902560)
(0.39062500,0.7156956130)
(0.40625000,0.6858507479)
(0.42187500,0.6554540246)
(0.43750000,0.6246732896)
(0.45312500,0.5936458717)
(0.46875000,0.5624785824)
(0.48437500,0.5312477152)
(0.50000000,0.4999990463)
(0.51562500,0.4687503775)
(0.53125000,0.4375195103)
(0.54687500,0.4063522209)
(0.56250000,0.3753248031)
(0.57812500,0.3445440681)
(0.59375000,0.3141473448)
(0.60937500,0.2843024796)
(0.62500000,0.2552078366)
(0.64062500,0.2270897543)
(0.65625000,0.2001745710)
(0.67187500,0.1746606509)
(0.68750000,0.1507158404)
(0.70312500,0.1284774683)
(0.71875000,0.1080523459)
(0.73437500,0.0895167670)
(0.75000000,0.0729165077)
(0.76562500,0.0582668266)
(0.78125000,0.0455524651)
(0.79687500,0.0347276471)
(0.81250000,0.0257160788)
(0.82812500,0.0184109489)
(0.84375000,0.0126749287)
(0.85937500,0.0083401715)
(0.87500000,0.0052083135)
(0.89062500,0.0030530161)
(0.90625000,0.0016479408)
(0.92187500,0.0007947237)
(0.93750000,0.0003255184)
(0.95312500,0.0001029958)
(0.96875000,0.0000203447)
(0.98437500,0.0000012715)
(1.00000000,0.0000000000)
};
\addlegendentry{$f_4$}
\end{axis}
\end{tikzpicture}
\caption{The first smooth finite box splines and their limit.  The plotted up-function is a
high-stage numerical evaluation used only for illustration.}
\label{repeated2:fig:iterates-comparison}
\end{figure}


\section{Replacing the last tail: the exact error decomposition}
\label{repeated2:sec:error-decomposition}

Put
\begin{equation}
 S_n=\sum_{k=1}^{n}2^{-k}U_k,
 \qquad a_n=2^{-n},
 \label{repeated2:eq:Sn-an}
\end{equation}
and let $p_n$ be the density of $S_n$.  Let $X'$ be an independent copy of the limiting
random variable $X$.  The finite and infinite laws split as
\begin{equation}
 X_n\stackrel d=S_n+a_nU,
 \qquad
 X\stackrel d=S_n+a_nX'.
 \label{repeated2:eq:tail-replacement}
\end{equation}
If
\begin{equation}
 \rho_a(x)=a^{-1}\RvachevUp(x/a)
 \label{repeated2:eq:scaled-up-density}
\end{equation}
then \eqref{repeated2:eq:tail-replacement} is exactly
\begin{equation}
 f_n=p_n\conv u_{a_n},
 \qquad
 \RvachevUp=p_n\conv\rho_{a_n}.
 \label{repeated2:eq:finite-infinite-convolution-split}
\end{equation}
The repeated-integration approximation therefore does one very specific thing: it keeps the
first $n$ dyadic uniform coordinates exact and replaces the remaining self-similar tail
$a_nX'$ by a single uniform coordinate $a_nU$.

\subsection{The comparison kernel}

\begin{definition}[Uniform-versus-up Peano kernel]
\label{repeated2:def:comparison-kernel}
For $t\in\RealNumbers$, define
\begin{equation}
 K(t)=\Expectation\PositivePart{U-t}-\Expectation\PositivePart{X-t}.
 \label{repeated2:eq:K-definition}
\end{equation}
\end{definition}

The kernel contains every constant needed for the sharp error theorem.  We first calculate
its basic invariants.

\begin{lemma}[Second and first absolute moments]
\label{repeated2:lem:up-moments}
The limiting random variable $X$ satisfies
\begin{equation}
 \Expectation X^2=\frac19,
 \qquad
 \Expectation|X|=\frac5{18}.
 \label{repeated2:eq:X-moments}
\end{equation}
Consequently,
\begin{equation}
 F\left(\frac14\right)=\frac5{72}.
 \label{repeated2:eq:F-quarter}
\end{equation}
\end{lemma}

\begin{proof}
Independence and centering in \eqref{repeated2:eq:up-random-series} give
\[
 \Expectation X^2=\frac13\sum_{k=1}^{\infty}4^{-k}=\frac19.
\]
For fixed $x\in[-1,1]$ and $U$ uniform on $[-1,1]$,
\[
 \Expectation\bigl(|x+U|\mid x\bigr)=\frac12\int_{-1}^{1}|x+u|\Differential u
 =\frac{1+x^2}{2}.
\]
Using $X\stackrel d=(X'+U)/2$ gives
\[
 \Expectation|X|=\frac12\Expectation|X'+U|
 =\frac14\bigl(1+\Expectation X^2\bigr)=\frac5{18}.
\]
By symmetry, $\Expectation X_+=\frac12\Expectation|X|=5/36$.  On the other hand,
\eqref{repeated2:eq:X-tail-F} and $F'(y)=2F(2y)$ imply
\[
 \Expectation X_+=\int_0^1F\left(\frac{1-t}{2}\right)\Differential t
 =2\int_0^{1/2}F(y)\Differential y=2F\left(\frac14\right),
\]
which proves \eqref{repeated2:eq:F-quarter}.
\end{proof}

\begin{theorem}[Exact kernel data]
\label{repeated2:thm:kernel-data}
The kernel $K$ is even, continuous, supported in $[-1,1]$, and nonnegative.  For
$0\le t\le1$,
\begin{equation}
 K(t)=\frac{(1-t)^2}{4}-2F\left(\frac{1-t}{4}\right).
 \label{repeated2:eq:K-explicit}
\end{equation}
Moreover,
\begin{equation}
 \boxed{\displaystyle
 \int_{-1}^{1}K(t)\Differential t=\frac19,
 \qquad
 \Norm{K}_\infty=K(0)=\frac19,}
 \label{repeated2:eq:K-mass-and-max}
\end{equation}
while
\begin{equation}
 \Norm{K'}_\infty=\frac{13}{72}.
 \label{repeated2:eq:Kprime-max}
\end{equation}
\end{theorem}

\begin{proof}
Evenness follows from symmetry and the fact that $U$ and $X$ have the same mean.  Both
variables are supported in $[-1,1]$, so $K$ vanishes outside that interval.

For nonnegativity, let $\phi$ be convex.  From the fixed-point identity
$X\stackrel d=(X'+U)/2$ and convexity,
\[
 \Expectation\phi(X)\le\frac12\Expectation\phi(X')+\frac12\Expectation\phi(U).
\]
Since $X'$ has the same law as $X$, this rearranges to
$\Expectation\phi(X)\le\Expectation\phi(U)$.  Thus $X$ is below $U$ in convex order.  Applying this to
$\phi_t(y)=\PositivePart{y-t}$ gives $K(t)\ge0$.

For $0\le t\le1$,
\[
 \Expectation\PositivePart{U-t}=\frac{(1-t)^2}{4}.
\]
Also, by \eqref{repeated2:eq:X-tail-F},
\begin{align*}
 \Expectation\PositivePart{X-t}
 &=\int_t^1\Probability(X>s)\Differential s
 =\int_t^1F\left(\frac{1-s}{2}\right)\Differential s\\
 &=2\int_0^{(1-t)/2}F(y)\Differential y
 =2F\left(\frac{1-t}{4}\right),
\end{align*}
which proves \eqref{repeated2:eq:K-explicit}.

The Peano identity proved in \cref{repeated2:lem:peano-identity} below, applied to
$h(y)=y^2/2$, gives
\[
 \int K(t)\Differential t=\frac12\bigl(\Expectation U^2-\Expectation X^2\bigr)
 =\frac12\left(\frac13-\frac19\right)=\frac19.
\]
At the origin, \cref{repeated2:lem:up-moments} gives
\[
 K(0)=\Expectation U_+-\Expectation X_+=\frac14-\frac5{36}=\frac19.
\]
It remains to see that this is the maximum.  Differentiating
\eqref{repeated2:eq:K-explicit} gives
\begin{equation}
 K'(t)=F\left(\frac{1-t}{2}\right)-\frac{1-t}{2}.
 \label{repeated2:eq:Kprime-F}
\end{equation}
On $[0,1/2]$, the Fabius function is convex: $F'(y)=2F(2y)$ is increasing.  Since
$F(0)=0$ and $F(1/2)=1/2$, convexity places its graph below the chord, so
$F(y)\le y$.  Therefore $K'(t)\le0$ for $t\ge0$, and $K(0)$ is the maximum.

Finally, put $y=(1-t)/2$ in \eqref{repeated2:eq:Kprime-F}.  The magnitude of $K'$ is the maximum of
$y-F(y)$ on $[0,1/2]$.  Its derivative is $1-2F(2y)$, which vanishes exactly at
$y=1/4$ because $F(1/2)=1/2$.  Thus
\[
 \Norm{K'}_\infty=\frac14-F\left(\frac14\right)
 =\frac14-\frac5{72}=\frac{13}{72}.
\]
\end{proof}

The equality between the mass and the maximum in \eqref{repeated2:eq:K-mass-and-max} is the key
numerical coincidence behind the clean final rate.  The mass controls all sufficiently
smooth stages; the maximum controls the borderline stage at which a distributional
box-spline derivative becomes an ordinary continuous error function.

\subsection{Peano identity and the pointwise error formula}

\begin{lemma}[Peano identity]
\label{repeated2:lem:peano-identity}
Let $h'$ be absolutely continuous on $[-1,1]$.  Then
\begin{equation}
 \Expectation h(U)-\Expectation h(X)=\int_{-1}^{1}K(t)h''(t)\Differential t.
 \label{repeated2:eq:peano-identity}
\end{equation}
\end{lemma}

\begin{proof}
For $y\in[-1,1]$,
\[
 h(y)=h(-1)+h'(-1)(y+1)+\int_{-1}^{1}\PositivePart{y-t}h''(t)\Differential t.
\]
Take expectations first with $y=U$ and then with $y=X$.  The constant terms cancel, and
the linear terms cancel because both random variables have mean zero.  Fubini's theorem
then gives \eqref{repeated2:eq:peano-identity}.
\end{proof}

Apply the lemma with $h(y)=p_n^{(r)}(x-a_ny)$.  Whenever $n\ge r+3$, the needed weak
second derivative is bounded and the result is an ordinary integral.

\begin{proposition}[Exact pointwise error identity]
\label{repeated2:prop:pointwise-error}
Let $r\ge0$ and $n\ge r+3$.  Then for every $x\in\RealNumbers$,
\begin{equation}
 f_n^{(r)}(x)-\RvachevUp^{(r)}(x)
 =a_n^2\int_{-1}^{1}K(t)
 p_n^{(r+2)}(x-a_nt)\Differential t.
 \label{repeated2:eq:exact-pointwise-error}
\end{equation}
\end{proposition}

\begin{proof}
From \eqref{repeated2:eq:finite-infinite-convolution-split},
\[
 f_n^{(r)}(x)=\Expectation p_n^{(r)}(x-a_nU),
 \qquad
 \RvachevUp^{(r)}(x)=\Expectation p_n^{(r)}(x-a_nX).
\]
Use \cref{repeated2:lem:peano-identity}; differentiating twice with respect to its argument contributes
the factor $a_n^2$.
\end{proof}

There is also a distributional form that remains valid at the borderline stage.  Define
\begin{equation}
 G_a(x)=aK(x/a).
 \label{repeated2:eq:Ga-definition}
\end{equation}
Since $K''$ is the difference between the densities of $U$ and $X$ in the distributional
sense,
\begin{equation}
 u_a-\rho_a=\D^2G_a.
 \label{repeated2:eq:kernel-distribution-identity}
\end{equation}
Consequently,
\begin{equation}
 f_n-\RvachevUp=p_n\conv\D^2G_{a_n}.
 \label{repeated2:eq:error-distribution-convolution}
\end{equation}
This identity will handle the first two smoothness thresholds without any appeal to a
classical second derivative of $p_n$.

It also gives a sharp weak-test estimate that is useful even when a uniform density norm
is not the natural metric.

\begin{corollary}[Sharp weak error]
\label{repeated2:cor:weak-error}
For every $\phi\in C^2(\RealNumbers)$ with bounded second derivative,
\begin{equation}
 \left|\Expectation\phi(X_n)-\Expectation\phi(X)\right|
 \le \frac{4^{-n}}9\Norm{\phi''}_\infty.
 \label{repeated2:eq:weak-error}
\end{equation}
Before taking absolute values the difference is nonnegative for convex $\phi$.
The constant is sharp: equality holds for $\phi(x)=x^2$.
\end{corollary}

\begin{proof}
Pair \eqref{repeated2:eq:error-distribution-convolution} with $\phi$ and move the
two distributional derivatives onto $\phi$.  The prefix density has mass one, while
$G_{a_n}$ is nonnegative and has mass $4^{-n}/9$ by
\eqref{repeated2:eq:K-mass-and-max}.  Convexity gives the sign.  For $\phi(x)=x^2$,
the identity is exactly the variance difference
$\Expectation X_n^2-\Expectation X^2=2\,4^{-n}/9$.
\end{proof}

\begin{lemma}[Summatory Thue--Morse cancellation]
\label{repeated2:lem:TM-prefix}
Let $\eps_j=(-1)^{\BinaryDigitSum(j)}$.  Then for every $J\ge0$,
\begin{equation}
 \sum_{j=0}^{J}\eps_j=
 \begin{cases}
  \eps_{J/2},&J\text{ even},\\
  0,&J\text{ odd}.
 \end{cases}                                                \label{repeated2:eq:TM-prefix}
\end{equation}
In particular, every prefix sum has absolute value at most one.
\end{lemma}

\begin{proof}
The identities $\eps_{2q}=\eps_q$ and $\eps_{2q+1}=-\eps_q$ make every complete adjacent pair
cancel.  An even endpoint leaves the last even term unpaired.
\end{proof}

\section{Sharp derivative geometry of the partial densities}
\label{repeated2:sec:partial-density-geometry}

The pointwise identity \eqref{repeated2:eq:exact-pointwise-error} becomes sharp because the derivatives
of $p_n$ are disjoint signed translates of a single tail density, each having a flat top of
known width and height.

For $m\ge0$, define the dyadic centers
\begin{equation}
 c_{m,j}=-1+\frac{2j+1}{2^m},
 \qquad 0\le j<2^m.
 \label{repeated2:eq:cmj-definition}
\end{equation}
For $m=0$, this is the singleton $\{0\}$.  At the top piecewise-polynomial derivative
order, derivatives below are understood through their bounded almost-everywhere
representatives; values at the finitely many knots may be chosen arbitrarily between the
adjacent limits and do not affect the $L^\infty$ norm used in this section.

\begin{theorem}[Derivative decomposition and exact plateaux]
\label{repeated2:thm:partial-derivative-geometry}
Let $0\le m\le n-1$.  Let $q_{m,n}$ be the density of the tail sum
\begin{equation}
 R_{m,n}=\sum_{k=m+1}^{n}2^{-k}U_k.
 \label{repeated2:eq:Rmn}
\end{equation}
Then, almost everywhere,
\begin{equation}
 p_n^{(m)}(x)=2^{\binom m2}
 \sum_{j=0}^{2^m-1}(-1)^{\BinaryDigitSum(j)}q_{m,n}(x-c_{m,j}).
 \label{repeated2:eq:pnm-decomposition}
\end{equation}
The translated supports in this sum have disjoint interiors.  Moreover,
\begin{equation}
 \Norm{p_n^{(m)}}_\infty=2^{\binom{m+1}{2}},
 \label{repeated2:eq:pnm-norm}
\end{equation}
and on the central interval of each translated tail,
\begin{equation}
 p_n^{(m)}(x)=(-1)^{\BinaryDigitSum(j)}2^{\binom{m+1}{2}}
 \quad\text{whenever}\quad |x-c_{m,j}|<2^{-n}.
 \label{repeated2:eq:pnm-plateau}
\end{equation}
\end{theorem}

\begin{proof}
In the distributional sense,
\begin{equation}
 \D u_{2^{-k}}=2^{k-1}\bigl(\delta_{-2^{-k}}-\delta_{2^{-k}}\bigr).
 \label{repeated2:eq:uniform-derivative-delta}
\end{equation}
Differentiate the first $m$ factors in
$p_n=u_{2^{-1}}\conv\cdots\conv u_{2^{-n}}$.  The product of the scalar factors is
\[
 \prod_{k=1}^{m}2^{k-1}=2^{\binom m2}.
\]
The $2^m$ signed sums of the locations $\pm2^{-k}$ are precisely the centers
$c_{m,j}$, in increasing order, and the sign is $(-1)^{\BinaryDigitSum(j)}$.  Convolution with the
undifferentiated tail gives \eqref{repeated2:eq:pnm-decomposition}.

Adjacent centers are separated by $2^{1-m}$.  The support radius of $q_{m,n}$ is
\[
 \sum_{k=m+1}^{n}2^{-k}=2^{-m}-2^{-n},
\]
so its support diameter is $2^{1-m}-2^{1-n}$, strictly smaller than the center spacing.
Hence the translated interiors are disjoint.

Scaling shows
\begin{equation}
 q_{m,n}(x)=2^m p_{n-m}(2^mx).
 \label{repeated2:eq:qmn-scaling}
\end{equation}
The density $p_s$ has supremum one: its first factor $u_{1/2}$ has height one, and
convolution with a probability density cannot increase the supremum.  It actually equals
one on $|x|<2^{-s}$, because the remaining factors have total support radius
$1/2-2^{-s}$.  Thus $q_{m,n}$ has height $2^m$ and is constant at that height on
$|x|<2^{-n}$.  Multiplying by $2^{\binom m2}$ proves
\eqref{repeated2:eq:pnm-norm} and \eqref{repeated2:eq:pnm-plateau}.
\end{proof}

At the missing endpoint $m=n$, the same calculation gives a signed Dirac comb:
\begin{equation}
 \D^np_n=2^{\binom n2}
 \sum_{j=0}^{2^n-1}(-1)^{\BinaryDigitSum(j)}\delta_{c_{n,j}}.
 \label{repeated2:eq:pn-top-delta-comb}
\end{equation}
The centers are spaced by $2^{1-n}=2a_n$.

\subsection{A sharp derivative norm for the up-function}

The finite plateau geometry also recovers a useful sharp fact about the limit itself.

\begin{corollary}[Exact derivative peaks of $\RvachevUp$]
\label{repeated2:cor:up-derivative-norm}
For every $m\ge0$,
\begin{equation}
 \boxed{\displaystyle
 \Norm{\RvachevUp^{(m)}}_\infty=2^{\binom{m+1}{2}}.}
 \label{repeated2:eq:up-derivative-norm}
\end{equation}
More precisely,
\begin{equation}
 \RvachevUp^{(m)}(c_{m,j})=(-1)^{\BinaryDigitSum(j)}2^{\binom{m+1}{2}}
 \qquad(0\le j<2^m).
 \label{repeated2:eq:up-derivative-peaks}
\end{equation}
\end{corollary}

\begin{proof}
Take $n\ge m+1$.  Since $\RvachevUp=p_n\conv\rho_{a_n}$,
\[
 \RvachevUp^{(m)}(x)=\int p_n^{(m)}(x-y)\rho_{a_n}(y)\Differential y.
\]
Convolution with a probability density cannot exceed
$\Norm{p_n^{(m)}}_\infty$, giving the upper bound.  At $x=c_{m,j}$, the whole support
$[-a_n,a_n]$ of $\rho_{a_n}$ lies in the constant plateau
\eqref{repeated2:eq:pnm-plateau}, so the convolution equals the signed plateau height.  This proves
both statements.
\end{proof}

\begin{corollary}[Higher-order flatness at the peak grid]
\label{repeated2:cor:higher-flatness-grid}
For every $q\ge0$, every $0\le j<2^q$, and every $k>q$,
\begin{equation}
 \RvachevUp^{(k)}(c_{q,j})=0.
 \label{repeated2:eq:higher-flatness-grid}
\end{equation}
\end{corollary}

\begin{proof}
At $q=0$, the only point is $c_{0,0}=0$.  On the right of zero,
$\RvachevUp(x)=F(1-x)$, and the smooth bounded Fabius function is constant on $[1,\infty)$;
therefore all positive-order derivatives of $\RvachevUp$ vanish at zero.

Let $q\ge1$, put $a=2^{-q}$, and write $\RvachevUp=p_q\conv\rho_a$.  For
$\ell=k-q\ge1$,
\[
 \RvachevUp^{(k)}=(\D^q p_q)\conv\rho_a^{(\ell)}.
\]
Insert the signed Dirac comb \eqref{repeated2:eq:pn-top-delta-comb}.  At $x=c_{q,j}$ every
off-diagonal translate is evaluated at distance at least $2a$, outside the support
$[-a,a]$ of $\rho_a$.  The diagonal translate is
$\rho_a^{(\ell)}(0)=a^{-\ell-1}\RvachevUp^{(\ell)}(0)=0$ by the first paragraph.
Thus every term vanishes.
\end{proof}

\section{The exact convergence rate}
\label{repeated2:sec:exact-rate}

We now combine the nonnegative kernel of \cref{repeated2:sec:error-decomposition} with the exact
plateaux of \cref{repeated2:sec:partial-density-geometry}.

\subsection{All stages after the first admissible one}

\begin{theorem}[Sharp $L^\infty$ error for every derivative]
\label{repeated2:thm:sharp-derivative-error}
Let $r\ge0$ and $n\ge r+2$.  Then
\begin{equation}
 \Norm{f_n^{(r)}-\RvachevUp^{(r)}}_\infty
 =\frac{4^{-n}}{9}\Norm{\RvachevUp^{(r+2)}}_\infty
 =\frac{2^{\binom{r+3}{2}}}{9\,4^n}.
 \label{repeated2:eq:sharp-derivative-error}
\end{equation}
At the grid \eqref{repeated2:eq:error-extremizer-grid}, the signed equality
\eqref{repeated2:eq:signed-extremizer-error} holds.
\end{theorem}

\begin{proof}
Put $m=r+2$.

First suppose $n\ge m+1$.  From \eqref{repeated2:eq:exact-pointwise-error}, the nonnegativity and
mass of $K$, and \eqref{repeated2:eq:pnm-norm},
\begin{align*}
 |f_n^{(r)}(x)-\RvachevUp^{(r)}(x)|
 &\le a_n^2\int K(t)\Differential t\,\Norm{p_n^{(m)}}_\infty\\
 &=4^{-n}\cdot\frac19\cdot2^{\binom{m+1}{2}}.
\end{align*}
At $x=c_{m,j}$, every point $c_{m,j}-a_nt$ with $-1<t<1$ lies in the constant plateau
\eqref{repeated2:eq:pnm-plateau}.  Hence no inequality remains:
\[
 f_n^{(r)}(c_{m,j})-\RvachevUp^{(r)}(c_{m,j})
 =(-1)^{\BinaryDigitSum(j)}4^{-n}2^{\binom{m+1}{2}}\int K(t)\Differential t.
\]
This proves the result for $n\ge m+1$.

It remains to treat the borderline case $n=m$.  Use
\eqref{repeated2:eq:error-distribution-convolution} and move all $m=r+2$ derivatives to $p_m$:
\[
 f_m^{(r)}-\RvachevUp^{(r)}=(\D^mp_m)\conv G_{a_m}.
\]
By \eqref{repeated2:eq:pn-top-delta-comb},
\begin{equation}
 f_m^{(r)}(x)-\RvachevUp^{(r)}(x)
 =2^{\binom m2}a_m
 \sum_{j=0}^{2^m-1}(-1)^{\BinaryDigitSum(j)}
 K\left(\frac{x-c_{m,j}}{a_m}\right).
 \label{repeated2:eq:borderline-error-kernel-copies}
\end{equation}
The copies of $K$ have support radius $a_m$ and their centers are $2a_m$ apart, so their
interiors are disjoint.  Therefore
\[
 \Norm{f_m^{(r)}-\RvachevUp^{(r)}}_\infty
 =2^{\binom m2}a_m\Norm K_\infty
 =2^{\binom m2-m}\frac19.
\]
Since
$\binom m2-m=\binom{m+1}{2}-2m$, this is exactly the right side of
\eqref{repeated2:eq:sharp-derivative-error} with $n=m$.  At $x=c_{m,j}$, use $K(0)=1/9$ to obtain
the signed formula.
\end{proof}

The two cases of the proof use different geometric data:
\begin{itemize}
\item for $n\ge r+3$, the error averages a constant extremal plateau and uses
$\int K=1/9$;
\item for $n=r+2$, the error is a disjoint sum of rescaled kernel copies and uses
$\Norm K_\infty=1/9$.
\end{itemize}
The equality of those two kernel constants is exactly what makes the same formula valid at
the threshold and beyond it.

\subsection{The first \texorpdfstring{$C^r$}{C-r} spline}

The first iterate with $r$ continuous derivatives is $f_{r+1}$.  Its error has a different
constant, controlled by $K'$ rather than by $K$.

\begin{theorem}[Exceptional first stage]
\label{repeated2:thm:first-admissible-stage}
For every $r\ge0$,
\begin{equation}
 \Norm{f_{r+1}^{(r)}-\RvachevUp^{(r)}}_\infty
 =2^{\binom{r+1}{2}}\Norm{K'}_\infty
 =\frac{13}{72}2^{\binom{r+1}{2}}.
 \label{repeated2:eq:first-stage-full}
\end{equation}
The supremum is attained at every point $c_{r+2,k}$, and the signed value is
\eqref{repeated2:eq:signed-first-stage-error}.
\end{theorem}

\begin{proof}
Set $n=r+1$.  In \eqref{repeated2:eq:error-distribution-convolution}, the total derivative order is
$r+2=n+1$.  Move $n$ derivatives to $p_n$ and one to $G_{a_n}$.  Since
$G_{a_n}'(x)=K'(x/a_n)$, \eqref{repeated2:eq:pn-top-delta-comb} gives
\begin{equation}
 f_n^{(r)}(x)-\RvachevUp^{(r)}(x)
 =2^{\binom n2}\sum_{j=0}^{2^n-1}(-1)^{\BinaryDigitSum(j)}
 K'\left(\frac{x-c_{n,j}}{a_n}\right).
 \label{repeated2:eq:first-stage-Kprime-copies}
\end{equation}
Again the supports have disjoint interiors, so the supremum is
$2^{\binom n2}\Norm{K'}_\infty$.  Substitute $n=r+1$ and
\eqref{repeated2:eq:Kprime-max}.

For the signed values, \eqref{repeated2:eq:Kprime-F} gives
$K'(-1/2)=13/72$ and $K'(1/2)=-13/72$.  The two points
$c_{n,j}\mp a_n/2$ are respectively $c_{n+1,2j}$ and $c_{n+1,2j+1}$; using
$\tau_{n+1}(2j)=\tau_n(j)$ and
$\tau_{n+1}(2j+1)=-\tau_n(j)$ in
\eqref{repeated2:eq:first-stage-Kprime-copies} proves
\eqref{repeated2:eq:signed-first-stage-error}.
\end{proof}

\subsection{Exact \texorpdfstring{$C^r$}{C-r} convergence}

Because the constants in \eqref{repeated2:eq:sharp-derivative-error} increase strictly with the
derivative order, the highest derivative controls the max norm \eqref{repeated2:eq:Cr-norm}.

\begin{corollary}[Exact $C^r$ norm]
\label{repeated2:cor:exact-Cr-rate}
For every $r\ge0$,
\begin{equation}
 \Norm{f_{r+1}-\RvachevUp}_{C^r}
 =\frac{13}{72}2^{\binom{r+1}{2}},
 \label{repeated2:eq:first-Cr-error}
\end{equation}
and for every $n\ge r+2$,
\begin{equation}
 \boxed{\displaystyle
 \Norm{f_n-\RvachevUp}_{C^r}
 =\frac{2^{\binom{r+3}{2}}}{9\,4^n}.}
 \label{repeated2:eq:exact-Cr-rate}
\end{equation}
Thus, for every fixed $r$, the tail of the sequence lies in $C^r(\RealNumbers)$ and converges to
$\RvachevUp$ in the $C^r$ norm; the error is exactly geometric from the second admissible stage
onward.  No finite $f_n$ is itself $C^\infty$, so this is deliberately stated as eventual
convergence in every finite smoothness norm rather than literally as convergence of a
sequence inside the Fr\'echet space $C^\infty(\RealNumbers)$.
\end{corollary}

\begin{proof}
For $n\ge r+2$, the constants in \eqref{repeated2:eq:sharp-derivative-error} grow strictly with the
derivative order, since the ratio of the order-$(j+1)$ constant to the order-$j$ constant
is $2^{j+3}$.  Hence the $r$th derivative controls the maximum in
\eqref{repeated2:eq:Cr-norm}.  At the exceptional stage $n=r+1$, the assertion is immediate for
$r=0$.  For $r\ge1$, the largest lower-order error is the order-$(r-1)$ error, and its
ratio to the order-$r$ error in \eqref{repeated2:eq:first-stage-full} is
\[
 \frac{2^{\binom{r+2}{2}}/(9\,4^{r+1})}
 {(13/72)2^{\binom{r+1}{2}}}
 =\frac{2^{2-r}}{13}<1.
\]
Thus the highest available derivative controls the max norm at the first stage as well.
\end{proof}

For the tent indexing $g_m=f_{m+1}$, this becomes the following direct answer to the
piecewise-linear formulation.

\begin{corollary}[Repeated integration starting from the tent]
\label{repeated2:cor:tent-indexed-rate}
Let $g_0(x)=\PositivePart{1-|x|}$ and $g_{m+1}=Tg_m$.  For every $r\ge0$,
\begin{equation}
 \Norm{g_r-\RvachevUp}_{C^r}
 =\frac{13}{72}2^{\binom{r+1}{2}},
 \label{repeated2:eq:tent-first-Cr}
\end{equation}
and, for every $m\ge r+1$,
\begin{equation}
 \boxed{\displaystyle
 \Norm{g_m-\RvachevUp}_{C^r}
 =\frac{2^{\binom{r+3}{2}-2}}{9\,4^m}.}
 \label{repeated2:eq:tent-Cr-rate}
\end{equation}
In particular, \eqref{repeated2:eq:tent-uniform-rate} holds.
\end{corollary}

\subsection{The uniform errors and exact extremal values}

For $r=0$, the complete startup behavior is
\begin{equation}
 \Norm{f_0-\RvachevUp}_\infty=\frac12,
 \qquad
 \Norm{f_1-\RvachevUp}_\infty=\frac{13}{72},
 \qquad
 \Norm{f_n-\RvachevUp}_\infty=\frac{8}{9\,4^n}\quad(n\ge2).
 \label{repeated2:eq:all-uniform-errors}
\end{equation}
The value for $f_0$ is immediate from $\RvachevUp(0)=1$ and $0\le\RvachevUp\le1$.  For $f_1$, put
$y=1-|x|$.  By \eqref{repeated2:eq:up-F-relation},
\begin{equation}
 f_1(x)-\RvachevUp(x)=y-F(y),
 \label{repeated2:eq:f1-error-F}
\end{equation}
whose absolute maximum is attained at $y=1/4$ and equals
$1/4-F(1/4)=13/72$.

The four extremizers from \eqref{repeated2:eq:error-extremizer-grid} are
$\pm1/4,\pm3/4$.  Since
\begin{equation}
 \RvachevUp(\pm3/4)=F(1/4)=\frac5{72},
 \qquad
 \RvachevUp(\pm1/4)=1-F(1/4)=\frac{67}{72},
 \label{repeated2:eq:up-quarter-values}
\end{equation}
we obtain exact finite-stage values for every $n\ge2$:
\begin{align}
 f_n(\pm3/4)&=\frac5{72}+\frac{8}{9\,4^n},
 \label{repeated2:eq:fn-threequarter}\\
 f_n(\pm1/4)&=\frac{67}{72}-\frac{8}{9\,4^n}.
 \label{repeated2:eq:fn-onequarter}
\end{align}
The first few are shown in \cref{repeated2:tab:exact-errors}.

\begin{table}[htbp]
\centering
\caption{Sharp uniform errors and two extremal values.}
\label{repeated2:tab:exact-errors}
\begin{tabular}{c@{\qquad}c@{\qquad}c@{\qquad}c}
\toprule
$n$&$\Norm{f_n-\RvachevUp}_\infty$&$f_n(3/4)$&$f_n(1/4)$\\
\midrule
$0$&$1/2$&--&--\\
$1$&$13/72$&$1/4$&$3/4$\\
$2$&$1/18$&$1/8$&$7/8$\\
$3$&$1/72$&$1/12$&$11/12$\\
$4$&$1/288$&$7/96$&$89/96$\\
$5$&$1/1152$&$9/128$&$119/128$\\
$6$&$1/4608$&$107/1536$&$1429/1536$\\
\bottomrule
\end{tabular}
\end{table}

\begin{figure}[htbp]
\centering
\begin{tikzpicture}
\begin{axis}[
  width=0.90\textwidth,
  height=0.47\textwidth,
  xmin=-1,xmax=1,ymin=-0.98,ymax=0.98,
  xlabel={$x$},ylabel={$4^n(f_n(x)-\RvachevUp(x))$},
  xtick={-1,-0.75,-0.5,-0.25,0,0.25,0.5,0.75,1},
  ytick={-0.8888889,-0.5,0,0.5,0.8888889},
  yticklabels={$-8/9$,$-1/2$,$0$,$1/2$,$8/9$},
  grid=both,
  minor tick num=1,
  tick label style={font=\small},
  label style={font=\small},
  legend style={font=\small,at={(0.5,-0.20)},anchor=north,legend columns=3},
]
\addplot+[mark=none,densely dashed] coordinates {
(-1.00000000,0.0000000000)
(-0.98437500,0.0078119904)
(-0.96875000,0.0312398836)
(-0.95312500,0.0701435800)
(-0.93750000,0.1238946997)
(-0.92187500,0.1910193703)
(-0.90625000,0.2690852216)
(-0.89062500,0.3548143910)
(-0.87500000,0.4444407357)
(-0.85937500,0.5340671101)
(-0.84375000,0.6197963689)
(-0.82812500,0.6978623692)
(-0.81250000,0.7649872484)
(-0.79687500,0.8187386364)
(-0.78125000,0.8576426606)
(-0.76562500,0.8810709413)
(-0.75000000,0.8888833786)
(-0.73437500,0.8810709413)
(-0.71875000,0.8576426606)
(-0.70312500,0.8187386364)
(-0.68750000,0.7649872484)
(-0.67187500,0.6978623692)
(-0.65625000,0.6197963689)
(-0.64062500,0.5340671101)
(-0.62500000,0.4444407357)
(-0.60937500,0.3548143910)
(-0.59375000,0.2690852216)
(-0.57812500,0.1910193703)
(-0.56250000,0.1238946997)
(-0.54687500,0.0701435800)
(-0.53125000,0.0312398836)
(-0.51562500,0.0078119904)
(-0.50000000,0.0000000000)
(-0.48437500,-0.0078119904)
(-0.46875000,-0.0312398836)
(-0.45312500,-0.0701435800)
(-0.43750000,-0.1238946997)
(-0.42187500,-0.1910193703)
(-0.40625000,-0.2690852216)
(-0.39062500,-0.3548143910)
(-0.37500000,-0.4444407357)
(-0.35937500,-0.5340671101)
(-0.34375000,-0.6197963689)
(-0.32812500,-0.6978623692)
(-0.31250000,-0.7649872484)
(-0.29687500,-0.8187386363)
(-0.28125000,-0.8576426606)
(-0.26562500,-0.8810709412)
(-0.25000000,-0.8888833786)
(-0.23437500,-0.8810709412)
(-0.21875000,-0.8576426606)
(-0.20312500,-0.8187386363)
(-0.18750000,-0.7649872484)
(-0.17187500,-0.6978623692)
(-0.15625000,-0.6197963689)
(-0.14062500,-0.5340671101)
(-0.12500000,-0.4444407357)
(-0.10937500,-0.3548143910)
(-0.09375000,-0.2690852216)
(-0.07812500,-0.1910193703)
(-0.06250000,-0.1238946997)
(-0.04687500,-0.0701435800)
(-0.03125000,-0.0312398836)
(-0.01562500,-0.0078119904)
(0.00000000,0.0000000000)
(0.01562500,-0.0078119904)
(0.03125000,-0.0312398836)
(0.04687500,-0.0701435800)
(0.06250000,-0.1238946997)
(0.07812500,-0.1910193703)
(0.09375000,-0.2690852216)
(0.10937500,-0.3548143910)
(0.12500000,-0.4444407357)
(0.14062500,-0.5340671101)
(0.15625000,-0.6197963689)
(0.17187500,-0.6978623692)
(0.18750000,-0.7649872484)
(0.20312500,-0.8187386363)
(0.21875000,-0.8576426606)
(0.23437500,-0.8810709412)
(0.25000000,-0.8888833786)
(0.26562500,-0.8810709412)
(0.28125000,-0.8576426606)
(0.29687500,-0.8187386363)
(0.31250000,-0.7649872484)
(0.32812500,-0.6978623692)
(0.34375000,-0.6197963689)
(0.35937500,-0.5340671101)
(0.37500000,-0.4444407357)
(0.39062500,-0.3548143910)
(0.40625000,-0.2690852216)
(0.42187500,-0.1910193703)
(0.43750000,-0.1238946997)
(0.45312500,-0.0701435800)
(0.46875000,-0.0312398836)
(0.48437500,-0.0078119904)
(0.50000000,0.0000000000)
(0.51562500,0.0078119904)
(0.53125000,0.0312398836)
(0.54687500,0.0701435800)
(0.56250000,0.1238946997)
(0.57812500,0.1910193703)
(0.59375000,0.2690852216)
(0.60937500,0.3548143910)
(0.62500000,0.4444407357)
(0.64062500,0.5340671101)
(0.65625000,0.6197963689)
(0.67187500,0.6978623692)
(0.68750000,0.7649872484)
(0.70312500,0.8187386364)
(0.71875000,0.8576426606)
(0.73437500,0.8810709413)
(0.75000000,0.8888833786)
(0.76562500,0.8810709413)
(0.78125000,0.8576426606)
(0.79687500,0.8187386364)
(0.81250000,0.7649872484)
(0.82812500,0.6978623692)
(0.84375000,0.6197963689)
(0.85937500,0.5340671101)
(0.87500000,0.4444407357)
(0.89062500,0.3548143910)
(0.90625000,0.2690852216)
(0.92187500,0.1910193703)
(0.93750000,0.1238946997)
(0.95312500,0.0701435800)
(0.96875000,0.0312398836)
(0.98437500,0.0078119904)
(1.00000000,0.0000000000)
};
\addlegendentry{$n=2$}
\addplot+[mark=none,dashdotted] coordinates {
(-1.00000000,0.0000000000)
(-0.98437500,0.0013018924)
(-0.96875000,0.0103797772)
(-0.95312500,0.0344857556)
(-0.93750000,0.0789188080)
(-0.92187500,0.1455959445)
(-0.90625000,0.2326001847)
(-0.89062500,0.3346325783)
(-0.87500000,0.4444410536)
(-0.85937500,0.5542495288)
(-0.84375000,0.6562819224)
(-0.82812500,0.7432861626)
(-0.81250000,0.8099632992)
(-0.79687500,0.8543963515)
(-0.78125000,0.8785023299)
(-0.76562500,0.8875802147)
(-0.75000000,0.8888821071)
(-0.73437500,0.8875802147)
(-0.71875000,0.8785023299)
(-0.70312500,0.8543963515)
(-0.68750000,0.8099632992)
(-0.67187500,0.7432861626)
(-0.65625000,0.6562819224)
(-0.64062500,0.5542495288)
(-0.62500000,0.4444410536)
(-0.60937500,0.3346325783)
(-0.59375000,0.2326001847)
(-0.57812500,0.1455959445)
(-0.56250000,0.0789188079)
(-0.54687500,0.0344857556)
(-0.53125000,0.0103797771)
(-0.51562500,0.0013018924)
(-0.50000000,-0.0000000000)
(-0.48437500,-0.0013018925)
(-0.46875000,-0.0103797772)
(-0.45312500,-0.0344857557)
(-0.43750000,-0.0789188080)
(-0.42187500,-0.1455959446)
(-0.40625000,-0.2326001848)
(-0.39062500,-0.3346325784)
(-0.37500000,-0.4444410536)
(-0.35937500,-0.5542495289)
(-0.34375000,-0.6562819225)
(-0.32812500,-0.7432861627)
(-0.31250000,-0.8099632993)
(-0.29687500,-0.8543963516)
(-0.28125000,-0.8785023301)
(-0.26562500,-0.8875802148)
(-0.25000000,-0.8888821072)
(-0.23437500,-0.8875802148)
(-0.21875000,-0.8785023301)
(-0.20312500,-0.8543963516)
(-0.18750000,-0.8099632993)
(-0.17187500,-0.7432861627)
(-0.15625000,-0.6562819225)
(-0.14062500,-0.5542495289)
(-0.12500000,-0.4444410536)
(-0.10937500,-0.3346325784)
(-0.09375000,-0.2326001848)
(-0.07812500,-0.1455959446)
(-0.06250000,-0.0789188080)
(-0.04687500,-0.0344857557)
(-0.03125000,-0.0103797772)
(-0.01562500,-0.0013018925)
(0.00000000,-0.0000000000)
(0.01562500,-0.0013018925)
(0.03125000,-0.0103797772)
(0.04687500,-0.0344857556)
(0.06250000,-0.0789188080)
(0.07812500,-0.1455959445)
(0.09375000,-0.2326001848)
(0.10937500,-0.3346325783)
(0.12500000,-0.4444410536)
(0.14062500,-0.5542495288)
(0.15625000,-0.6562819224)
(0.17187500,-0.7432861626)
(0.18750000,-0.8099632992)
(0.20312500,-0.8543963515)
(0.21875000,-0.8785023300)
(0.23437500,-0.8875802147)
(0.25000000,-0.8888821071)
(0.26562500,-0.8875802147)
(0.28125000,-0.8785023300)
(0.29687500,-0.8543963515)
(0.31250000,-0.8099632992)
(0.32812500,-0.7432861626)
(0.34375000,-0.6562819224)
(0.35937500,-0.5542495288)
(0.37500000,-0.4444410536)
(0.39062500,-0.3346325783)
(0.40625000,-0.2326001848)
(0.42187500,-0.1455959445)
(0.43750000,-0.0789188080)
(0.45312500,-0.0344857556)
(0.46875000,-0.0103797772)
(0.48437500,-0.0013018924)
(0.50000000,0.0000000000)
(0.51562500,0.0013018925)
(0.53125000,0.0103797772)
(0.54687500,0.0344857557)
(0.56250000,0.0789188080)
(0.57812500,0.1455959446)
(0.59375000,0.2326001848)
(0.60937500,0.3346325784)
(0.62500000,0.4444410536)
(0.64062500,0.5542495289)
(0.65625000,0.6562819225)
(0.67187500,0.7432861627)
(0.68750000,0.8099632992)
(0.70312500,0.8543963516)
(0.71875000,0.8785023300)
(0.73437500,0.8875802148)
(0.75000000,0.8888821072)
(0.76562500,0.8875802148)
(0.78125000,0.8785023300)
(0.79687500,0.8543963516)
(0.81250000,0.8099632992)
(0.82812500,0.7432861627)
(0.84375000,0.6562819224)
(0.85937500,0.5542495289)
(0.87500000,0.4444410536)
(0.89062500,0.3346325783)
(0.90625000,0.2326001848)
(0.92187500,0.1455959446)
(0.93750000,0.0789188080)
(0.95312500,0.0344857557)
(0.96875000,0.0103797772)
(0.98437500,0.0013018924)
(1.00000000,0.0000000000)
};
\addlegendentry{$n=3$}
\addplot+[mark=none,line width=1.1pt] coordinates {
(-1.00000000,0.0000000000)
(-0.98437500,0.0000783989)
(-0.96875000,0.0032098493)
(-0.95312500,0.0202300827)
(-0.93750000,0.0619748303)
(-0.92187500,0.1313403460)
(-0.90625000,0.2254303760)
(-0.89062500,0.3334091890)
(-0.87500000,0.4444410534)
(-0.85937500,0.5554729179)
(-0.84375000,0.6634517308)
(-0.82812500,0.7575417608)
(-0.81250000,0.8269072765)
(-0.79687500,0.8686520241)
(-0.78125000,0.8856722574)
(-0.76562500,0.8888037078)
(-0.75000000,0.8888821066)
(-0.73437500,0.8888037077)
(-0.71875000,0.8856722572)
(-0.70312500,0.8686520238)
(-0.68750000,0.8269072761)
(-0.67187500,0.7575417603)
(-0.65625000,0.6634517302)
(-0.64062500,0.5554729172)
(-0.62500000,0.4444410527)
(-0.60937500,0.3334091881)
(-0.59375000,0.2254303752)
(-0.57812500,0.1313403451)
(-0.56250000,0.0619748293)
(-0.54687500,0.0202300816)
(-0.53125000,0.0032098483)
(-0.51562500,0.0000783978)
(-0.50000000,-0.0000000012)
(-0.48437500,-0.0000784001)
(-0.46875000,-0.0032098505)
(-0.45312500,-0.0202300839)
(-0.43750000,-0.0619748316)
(-0.42187500,-0.1313403474)
(-0.40625000,-0.2254303774)
(-0.39062500,-0.3334091904)
(-0.37500000,-0.4444410549)
(-0.35937500,-0.5554729194)
(-0.34375000,-0.6634517323)
(-0.32812500,-0.7575417623)
(-0.31250000,-0.8269072781)
(-0.29687500,-0.8686520257)
(-0.28125000,-0.8856722591)
(-0.26562500,-0.8888037095)
(-0.25000000,-0.8888821084)
(-0.23437500,-0.8888037094)
(-0.21875000,-0.8856722590)
(-0.20312500,-0.8686520256)
(-0.18750000,-0.8269072779)
(-0.17187500,-0.7575417622)
(-0.15625000,-0.6634517321)
(-0.14062500,-0.5554729192)
(-0.12500000,-0.4444410547)
(-0.10937500,-0.3334091902)
(-0.09375000,-0.2254303773)
(-0.07812500,-0.1313403472)
(-0.06250000,-0.0619748315)
(-0.04687500,-0.0202300838)
(-0.03125000,-0.0032098504)
(-0.01562500,-0.0000784000)
(0.00000000,-0.0000000011)
(0.01562500,-0.0000784000)
(0.03125000,-0.0032098504)
(0.04687500,-0.0202300838)
(0.06250000,-0.0619748314)
(0.07812500,-0.1313403472)
(0.09375000,-0.2254303772)
(0.10937500,-0.3334091901)
(0.12500000,-0.4444410546)
(0.14062500,-0.5554729190)
(0.15625000,-0.6634517319)
(0.17187500,-0.7575417619)
(0.18750000,-0.8269072776)
(0.20312500,-0.8686520252)
(0.21875000,-0.8856722585)
(0.23437500,-0.8888037089)
(0.25000000,-0.8888821077)
(0.26562500,-0.8888037088)
(0.28125000,-0.8856722583)
(0.29687500,-0.8686520249)
(0.31250000,-0.8269072772)
(0.32812500,-0.7575417614)
(0.34375000,-0.6634517313)
(0.35937500,-0.5554729183)
(0.37500000,-0.4444410538)
(0.39062500,-0.3334091892)
(0.40625000,-0.2254303763)
(0.42187500,-0.1313403462)
(0.43750000,-0.0619748304)
(0.45312500,-0.0202300828)
(0.46875000,-0.0032098494)
(0.48437500,-0.0000783989)
(0.50000000,0.0000000000)
(0.51562500,0.0000783990)
(0.53125000,0.0032098494)
(0.54687500,0.0202300828)
(0.56250000,0.0619748305)
(0.57812500,0.1313403463)
(0.59375000,0.2254303763)
(0.60937500,0.3334091893)
(0.62500000,0.4444410538)
(0.64062500,0.5554729183)
(0.65625000,0.6634517312)
(0.67187500,0.7575417612)
(0.68750000,0.8269072769)
(0.70312500,0.8686520246)
(0.71875000,0.8856722579)
(0.73437500,0.8888037084)
(0.75000000,0.8888821072)
(0.76562500,0.8888037083)
(0.78125000,0.8856722579)
(0.79687500,0.8686520245)
(0.81250000,0.8269072768)
(0.82812500,0.7575417610)
(0.84375000,0.6634517310)
(0.85937500,0.5554729181)
(0.87500000,0.4444410536)
(0.89062500,0.3334091891)
(0.90625000,0.2254303761)
(0.92187500,0.1313403461)
(0.93750000,0.0619748304)
(0.95312500,0.0202300827)
(0.96875000,0.0032098493)
(0.98437500,0.0000783989)
(1.00000000,0.0000000000)
};
\addlegendentry{$n=6$}
\addplot+[mark=none,dotted,domain=-1:1,samples=2] {8/9};
\addplot+[mark=none,dotted,domain=-1:1,samples=2] {-8/9};
\end{axis}
\end{tikzpicture}
\caption{The errors after multiplication by $4^n$.  Their shape approaches $\RvachevUp''/9$, while
the exact extrema remain fixed at $\pm8/9$ from the second stage onward.  Numerical curves are
shown only as an illustration of the proved identities.}
\label{repeated2:fig:scaled-errors-supplement}
\end{figure}


\section{The asymptotic error profile}
\label{repeated2:sec:asymptotic-profile}

The exact norm formula gives the rate, but the Peano identity also identifies the whole
limiting error shape.

\begin{theorem}[Uniform leading profile]
\label{repeated2:thm:leading-profile}
For every fixed $r\ge0$,
\begin{equation}
 4^n\bigl(f_n^{(r)}-\RvachevUp^{(r)}\bigr)
 \longrightarrow \frac19\RvachevUp^{(r+2)}
 \qquad\text{uniformly on }\RealNumbers.
 \label{repeated2:eq:leading-profile}
\end{equation}
\end{theorem}

\begin{proof}
Put $m=r+2$.  From \eqref{repeated2:eq:exact-pointwise-error},
\[
 4^n\bigl(f_n^{(r)}(x)-\RvachevUp^{(r)}(x)\bigr)
 =\int K(t)p_n^{(m)}(x-a_nt)\Differential t.
\]
For fixed $m$, $p_n^{(m)}\to\RvachevUp^{(m)}$ uniformly.  One direct estimate follows from
$\RvachevUp=p_n\conv\rho_{a_n}$ and the mean-value theorem:
\[
 \Norm{p_n^{(m)}-\RvachevUp^{(m)}}_\infty
 \le a_n\Expectation|X|\,\Norm{p_n^{(m+1)}}_\infty,
\]
for all sufficiently large $n$.  The same derivative bound shows that replacing
$p_n^{(m)}(x-a_nt)$ by $p_n^{(m)}(x)$ costs $O(a_n)$ uniformly.  Finally use
$\int K=1/9$.
\end{proof}

Combining \eqref{repeated2:eq:leading-profile} with \eqref{repeated2:eq:up-derivative-norm} predicts the sharp
constant in \eqref{repeated2:eq:sharp-derivative-error}.  The exact theorem says more: the norm of the
leading profile is already the exact norm of the full error for every $n\ge r+2$.

There is an operator-theoretic reason for the factor $1/4$.  Differentiating the fixed-point
identity $T\RvachevUp=\RvachevUp$ gives
\begin{equation}
 T\bigl(\RvachevUp^{(r)}\bigr)=2^{-r}\RvachevUp^{(r)}.
 \label{repeated2:eq:T-eigenfunctions}
\end{equation}
Thus $\RvachevUp''$ is an eigenfunction with eigenvalue $1/4$.  Matching the tail mass and mean
annihilates the modes of orders zero and one, so the second-derivative mode is the first one
that survives.

\begin{corollary}[$L^p$ leading-profile asymptotics]
\label{repeated2:cor:Lp}
For every fixed $r\ge0$ and $1\le p<\infty$,
\begin{equation}
 \lim_{n\to\infty}4^n
 \Norm{f_n^{(r)}-\RvachevUp^{(r)}}_{L^p}
 =\frac19\Norm{\RvachevUp^{(r+2)}}_{L^p}.
 \label{repeated2:eq:Lp-asymptotic}
\end{equation}
For $p=\infty$, the corresponding norm equality is already exact for every
$n\ge r+2$.
\end{corollary}

\begin{proof}
Use the uniform convergence in \eqref{repeated2:eq:leading-profile} and the common compact
support for finite $p$; use \eqref{repeated2:eq:sharp-derivative-error} for $p=\infty$.
\end{proof}

\subsection{A full even-derivative expansion}

The Fourier product provides an all-orders refinement.  Put
\begin{equation}
 \Phi(z)=\prod_{k=1}^{\infty}\SincRad(z/2^k),
 \qquad
 A(z)=\frac{\SincRad z}{\Phi(z)}.
 \label{repeated2:eq:A-ratio}
\end{equation}
The quotient is analytic near the origin (indeed on $|z|<2\pi$) and meromorphic
globally.  Wherever $\Phi(2^{-n}t)\ne0$ one has
\begin{equation}
 \mathscr F f_n(t)=\mathscr F\RvachevUp(t)A(2^{-n}t),
 \label{repeated2:eq:fourier-ratio}
\end{equation}
and at the exceptional frequencies the product on the right is interpreted by its
continuous finite-product extension.  Only the Taylor germ at zero is used below.  It is
convenient to encode that germ by
\begin{equation}
 B(w)=A(\iu w)
 =\frac{\sinh w/w}{\displaystyle\prod_{k=1}^{\infty}
 \frac{\sinh(w/2^k)}{w/2^k}}
 =\sum_{m=0}^{\infty}\beta_mw^{2m}.
 \label{repeated2:eq:B-generating}
\end{equation}
The first coefficients are
\begin{equation}
 \beta_0=1,
 \qquad \beta_1=\frac19,
 \qquad \beta_2=\frac{2}{2025},
 \qquad \beta_3=-\frac1{2679075},
 \qquad \beta_4=\frac{58}{10247461875}.
 \label{repeated2:eq:beta-first}
\end{equation}

\begin{proposition}[All-orders $C^r$ expansion]
\label{repeated2:prop:all-orders-expansion}
For fixed integers $r\ge0$ and $M\ge1$,
\begin{equation}
 f_n^{(r)}
 =\sum_{m=0}^{M-1}\beta_m4^{-mn}\RvachevUp^{(r+2m)}
 +O_{r,M}(4^{-Mn})
 \label{repeated2:eq:all-orders-expansion}
\end{equation}
uniformly on $\RealNumbers$ as $n\to\infty$.
\end{proposition}

\begin{proof}
Choose $R$ with $1<R<2\pi$ and put $a=2^{-n}$.  On $|z|\le R$, Taylor's theorem gives
\begin{equation}
 A(z)=\sum_{m=0}^{M-1}\beta_m(-1)^m z^{2m}+O_{M,R}(|z|^{2M}).
 \label{repeated2:eq:A-Taylor-remainder}
\end{equation}
On the low-frequency region $|t|\le R/a$, multiply this by
$(\iu t)^r\mathscr F\RvachevUp(t)$.  Since $\mathscr F\RvachevUp$ is a Schwartz function, the
$L^1$ norm of the resulting remainder is at most
\[
 C_{r,M,R}a^{2M}\int_{\RealNumbers}|t|^{r+2M}|\mathscr F\RvachevUp(t)|\Differential t
 =O_{r,M}(a^{2M}).
\]
The sign in \eqref{repeated2:eq:A-Taylor-remainder} is exactly the sign in
$(\iu t)^{2m}=(-1)^mt^{2m}$, so Fourier inversion produces the displayed derivatives of
$\RvachevUp$.

It remains to control $|t|>R/a$.  From \eqref{repeated2:eq:finite-fourier-product} and
$|\SincRad y|\le\min(1,|y|^{-1})$, for $|t|\ge R2^n$,
\begin{equation}
 |\mathscr F f_n(t)|
 \le 2^{n(n+1)/2+n}|t|^{-(n+1)}.
 \label{repeated2:eq:finite-product-high-frequency}
\end{equation}
Hence, for fixed $r$ and large $n$,
\[
 \int_{|t|>R2^n}|t|^r|\mathscr F f_n(t)|\Differential t
 \le C_r R^{r-n}2^{-n^2/2+(r+3/2)n},
\]
which is smaller than $a^{2M}$ for all sufficiently large $n$.  The high-frequency tails
of the finitely many $\RvachevUp^{(r+2m)}$ terms are also $O_{r,M}(a^{2M})$, again because
$\mathscr F\RvachevUp$ is Schwartz.  Fourier inversion now proves the uniform remainder in
\eqref{repeated2:eq:all-orders-expansion}.
\end{proof}

For explicit coefficient generation, write
\begin{equation}
 \log B(w)=\sum_{m=1}^{\infty}\lambda_mw^{2m}.
 \label{repeated2:eq:lambda-log}
\end{equation}
With the standard Bernoulli numbers $B_{2m}$,
\begin{equation}
 \lambda_m=
 \frac{2^{2m-1}B_{2m}}{m(2m)!}
 \frac{4^m-2}{4^m-1}.
 \label{repeated2:eq:lambda-formula}
\end{equation}
Indeed,
\[
 \log\frac{\sinh w}{w}
 =\sum_{m\ge1}\frac{2^{2m-1}B_{2m}}{m(2m)!}w^{2m},
 \qquad
 \sum_{k\ge1}4^{-mk}=\frac1{4^m-1};
\]
subtracting the logarithms of the denominator factors in \eqref{repeated2:eq:B-generating} gives
\eqref{repeated2:eq:lambda-formula}.  The coefficients $\beta_m$ are then obtained recursively from
\begin{equation}
 \beta_0=1,
 \qquad
 \beta_m=\frac1m\sum_{j=1}^{m}j\lambda_j\beta_{m-j}.
 \label{repeated2:eq:beta-recurrence}
\end{equation}
Thus the first terms of \eqref{repeated2:eq:all-orders-expansion} are
\begin{equation}
 f_n^{(r)}=\RvachevUp^{(r)}
 +\frac{\RvachevUp^{(r+2)}}{9\,4^n}
 +\frac{2\RvachevUp^{(r+4)}}{2025\,16^n}
 -\frac{\RvachevUp^{(r+6)}}{2679075\,64^n}
 +O_r(256^{-n}).
 \label{repeated2:eq:first-asymptotic-terms}
\end{equation}
At the extremal grid \eqref{repeated2:eq:error-extremizer-grid}, all higher derivative terms vanish
by \cref{repeated2:cor:higher-flatness-grid}; there the leading term is not merely asymptotic but
exactly equal to the full error, as \eqref{repeated2:eq:signed-extremizer-error} states.

\section{Comparison with ordinary truncation}
\label{repeated2:sec:ordinary-truncation}

The ordinary prefix density $p_n$ and the repeated-integration density $f_n$ retain the
same first $n$ dyadic boxes, but they handle the remaining tail differently.
\begin{center}
\begin{tabularx}{0.96\textwidth}{@{}p{0.23\textwidth}XX@{}}
\toprule
 & Ordinary truncation $p_n$ & Repeated integration $f_n$\\
\midrule
Random variable & $S_n$ & $S_n+2^{-n}U$\\
Number of boxes & $n$ & $n+1$\\
Support radius & $1-2^{-n}$ & $1$ exactly\\
Tail model & omitted & uniform on the exact tail support\\
Mean error & zero by symmetry & zero by symmetry\\
Leading smooth error & uncorrected tail scale & variance defect, scale $4^{-n}$\\
Sharp uniform law & not the law studied here & $8/(9\,4^n)$ for $n\ge2$\\
\bottomrule
\end{tabularx}
\end{center}
The additional uniform box is therefore a support-and-moment correction rather than a
negligible finite-product artifact.  If one could also replace the tail by a positive,
compactly supported law with matching variance, the $\RvachevUp''$ term would disappear and the
next formal mode suggests a $16^{-n}$ scheme.  This variance-matched construction is a
conjectural direction; no such positive scheme is asserted here.

\section{Exact computation and numerical use}
\label{repeated2:sec:computation}

A piecewise-polynomial representation also admits a simple primitive recursion.  If
$A_{n-1}$ is a continuous primitive of $f_{n-1}$, then
\begin{equation}
 f_n(x)=A_{n-1}(2x+1)-A_{n-1}(2x-1).
 \label{repeated2:eq:primitive-algorithm}
\end{equation}
It preserves exact rational coefficients and raises the polynomial degree by one, making
it convenient when every cell, rather than one point, is required.

\subsection{A direct exact-arithmetic evaluator}

Formula \eqref{repeated2:eq:exact-iterate-intro} immediately gives an exact rational evaluator at
rational arguments.  The following compact Python implementation uses arbitrary-precision
fractions.  It is intended as executable notation rather than as the fastest large-$n$
algorithm.

\begin{lstlisting}[style=reportcode,language=Python,caption={Exact evaluation of the finite iterates.},label={repeated2:lst:exact-evaluator}]
from fractions import Fraction
from math import factorial


def tau(n: int, j: int) -> int:
    if j < 0 or j >= (1 << n):
        return 0
    return -1 if (j.bit_count() & 1) else 1


def positive_power(x: Fraction, n: int) -> Fraction:
    return x**n if x > 0 else Fraction(0)


def f(n: int, x: Fraction) -> Fraction:
    x = Fraction(x)
    if n == 0:
        return Fraction(1, 2) if abs(x) < 1 else Fraction(0)

    coefficient = Fraction(
        1 << (n * (n + 1) // 2 - 1), factorial(n)
    )
    mesh = Fraction(1, 1 << (n - 1))

    return coefficient * sum(
        (tau(n, j) - tau(n, j - 1))
        * positive_power(x + 1 - j * mesh, n)
        for j in range((1 << n) + 1)
    )
\end{lstlisting}

Only indices satisfying $j<2^{n-1}(x+1)$ contribute, so the loop may be truncated at the
active cell.  Derivatives are obtained by replacing $n!$ with $(n-r)!$ and the power $n$
with $n-r$, as in \eqref{repeated2:eq:finite-spline-derivatives}.

\subsection{Conditioning and stable evaluation}

The truncated-power formula is exact but can exhibit severe cancellation at large $n$:
individual terms are much larger than their sum.  Three alternatives are preferable in
floating-point work.

\begin{enumerate}[label=\textup{(\roman*)}]
\item Use the convolution description \eqref{repeated2:eq:fn-box-convolution} and propagate a
piecewise-polynomial representation from one level to the next.
\item Use a B-spline/de Boor representation on the dyadic knot sequence rather than the
truncated-power basis.
\item In Fourier space, evaluate the finite sinc product \eqref{repeated2:eq:finite-fourier-product}
and invert it with a quadrature adapted to a compactly supported smooth target.
\end{enumerate}

For certified approximation of $\RvachevUp$, the sharp error gives an immediate stopping rule.  To
obtain uniform error at most $\varepsilon$ from the tent-indexed sequence, it is necessary
and sufficient that
\begin{equation}
 m\ge
 \Ceiling{\frac12\log_2\left(\frac{2}{9\varepsilon}\right)}
 \qquad(m\ge1).
 \label{repeated2:eq:stopping-rule}
\end{equation}
For a fixed derivative order $r$, replace the numerator $2/9$ by
$2^{\binom{r+3}{2}-2}/9$ and also require $m\ge r+1$.

\section{How this fits the existing ProveIt development}
\label{repeated2:sec:formalization-roadmap}

The present approximation is closely related to, but not identical with, two families
already developed in the repository.

\begin{itemize}
\item The finite convolution and polynomial machinery in
\nolinkurl{FabiusFunction/EarlyApproximants.lean} provides reusable measure-theoretic and
coefficient infrastructure.
\item The Thue--Morse spline machinery in
\nolinkurl{FabiusFunction/FabiusUniformSpline.lean} already contains many of the finite-power
sum and block-translation lemmas needed to identify the top derivative pattern.
\item The centered finite-CDF spline in the main article approximates the bounded Fabius
function.  The sequence here approximates the up-density itself and has a different tail
replacement and a different sharp rate.
\end{itemize}

A natural formalization could be organized as follows.

\begin{enumerate}[label=\textup{(\arabic*)}]
\item Define the probability operator $T$ and prove
$Th(x)=\int_{2x-1}^{2x+1}h$ for supported integrable $h$.
\item Define $X_n$ recursively and prove the finite random-sum identity
\eqref{repeated2:eq:Xn-random-sum}; derive \eqref{repeated2:eq:fn-box-convolution} at the level of measures.
\item Introduce the nonuniform box spline and prove the general inclusion--exclusion lemma
\eqref{repeated2:eq:general-box-spline}.
\item Reuse the finite Thue--Morse product to prove
\eqref{repeated2:eq:first-difference-generating}, then obtain the exact iterate and derivative formulas.
\item Define $K$ by \eqref{repeated2:eq:K-definition}.  Prove convex domination with the one-line
Jensen argument, then formalize its mass, maximum, and derivative maximum.
\item Prove the derivative decomposition \eqref{repeated2:eq:pnm-decomposition}, support separation,
and the exact plateaux.
\item Combine the Peano identity with the plateau lemma, splitting the proof into the
ordinary case $n\ge r+3$ and the distributional threshold $n=r+2$.
\end{enumerate}

The analytic core of the sharp theorem is therefore modest.  The technically delicate
parts for Lean are likely to be bookkeeping for weak derivatives of interval densities and
the transfer between measure convolution and classical derivatives.  One can avoid most
of that friction by proving the borderline formulas directly from finite signed Dirac
combs and using ordinary absolutely continuous derivatives only after the threshold.

\section{Conclusions}
\label{repeated2:sec:conclusion}

The repeated-integration construction has three mutually reinforcing exact descriptions:

\begin{enumerate}[label=\textup{(\roman*)}]
\item an averaging recursion $X\mapsto(X+U)/2$;
\item a convolution of dyadically scaled uniform densities;
\item a dyadic polynomial spline whose top derivative is a Thue--Morse word.
\end{enumerate}

These descriptions identify the limit as Rvachev's up-function and reduce the finite-stage
closed form to a single first-difference Thue--Morse sum.  More unexpectedly, the same
probabilistic representation yields a completely sharp convergence theorem.  The finite
stage and the limit differ only in the last scaled tail: a uniform variable replaces an
up-distributed variable.  Their mass and mean agree, so the error starts at second order in
$2^{-n}$.  The comparison kernel quantifies this replacement exactly, while the partial
box spline supplies disjoint extremal plateaux.  The outcome is the exact identity
\[
 \Norm{f_n^{(r)}-\RvachevUp^{(r)}}_\infty
 =\frac{2^{\binom{r+3}{2}}}{9\,4^n}
 \qquad(n\ge r+2),
\]
not just an upper bound or an asymptotic equivalent.

In the most literal piecewise-linear formulation, starting from the tent and integrating
again and again, the answer is especially crisp: after the initial tent, the uniform error
is exactly $2/(9\,4^m)$.  Every additional integration divides the worst-case error by four,
and the points at which that error is attained are fixed dyadic points rather than moving
with the iteration.

\renewcommand{\refname}{Topic references}
\begin{thebibliography}{99}

\bibitem{repeated2:mse-question}
kram1032,
\emph{Recursive Integration over Piecewise Polynomials: Closed form?},
Mathematics Stack Exchange, question 240687, 19 November 2012,
\url{https://math.stackexchange.com/questions/240687/recursive-integration-over-piecewise-polynomials-closed-form}.

\bibitem{repeated2:mse-answer}
Nathan Grigg,
\emph{Answer to ``Recursive Integration over Piecewise Polynomials: Closed form?''},
Mathematics Stack Exchange, 22 November 2012,
\url{https://math.stackexchange.com/a/240805}.

\bibitem{repeated2:arias1982}
Juan Arias de Reyna,
\emph{An infinitely differentiable function with compact support: Definition and
properties},
English translation, arXiv:1702.05442, 2017; original Spanish version in
\emph{Revista de la Real Academia de Ciencias Exactas, Físicas y Naturales de Madrid}
\textbf{76} (1982), 21--38.

\bibitem{repeated2:arias2018}
Juan Arias de Reyna,
\emph{Arithmetic of the Fabius function},
arXiv:1702.06487; \emph{Integers} \textbf{18} (2018), A51.

\bibitem{repeated2:fabius1966}
Jaap Fabius,
\emph{A probabilistic example of a nowhere analytic $C^\infty$-function},
\emph{Zeitschrift für Wahrscheinlichkeitstheorie und Verwandte Gebiete}
\textbf{5} (1966), 173--174.

\bibitem{repeated2:jessen-wintner1935}
Børge Jessen and Aurel Wintner,
\emph{Distribution functions and the Riemann zeta function},
\emph{Transactions of the American Mathematical Society}
\textbf{38} (1935), 48--88.

\bibitem{repeated2:proveit}
Vladimir Reshetnikov,
\emph{ProveIt: Analysis/FabiusFunction},
\url{https://github.com/VladimirReshetnikov/ProveIt/tree/main/Analysis/FabiusFunction},
repository snapshot consulted 25 August 2026.

\end{thebibliography}


\renewcommand{\arraystretch}{1}
% END SYNTHESIZED TOPIC: repeated integration and Rvachev's up-function




% BEGIN SYNTHESIZED TOPIC: exact dyadic calculus and alternative representations
% Former source SHA-256 values:
% 81a0a911ac7ab28e12c6b87ccaf76ffccaf32e48f873abff18fb7a2d01bcf3e5
% 462276b10fcd32b0446deb7cfedc4ec07c2ae55dbd333d5ff9b1d98f07df89e1
\clearpage
\part{The Dyadic Web of the Fabius Function}
\label{part:dyadicweb}
\begin{boxedremark}
\textbf{Synthesized provenance and status.} This part merges the former
\nolinkurl{Fabius_Dyadic_q_Connections/Fabius_Dyadic_q_Connections.tex} and
\nolinkurl{Fabius_Dyadic_Formulae_and_Alternative_Representations/Fabius_Dyadic_Formulae_and_Alternative_Representations.tex}.
The Dyadic Web supplies the canonical master-coefficient, exact-stabilization, and
$q$-Richardson framework.  The earlier source contributes the Appell compression,
ordered-composition and Hessenberg formulas, binary prefix rearrangement, limiting filter,
and integral-transform views.  The former coarse $O(2^{-n})$ discrete estimate is omitted
in favor of the explicit $O_Q(4^{-n})$ theorem below.  Claims without an exact audited Lean
declaration remain frontier obligations.\par\smallskip Former source SHA-256 values:
\nolinkurl{81a0a911ac7ab28e12c6b87ccaf76ffccaf32e48f873abff18fb7a2d01bcf3e5} and
\nolinkurl{462276b10fcd32b0446deb7cfedc4ec07c2ae55dbd333d5ff9b1d98f07df89e1}.
\end{boxedremark}
\section*{Topic abstract}
Exact values of the Fabius function at dyadic rationals appear in several forms that look
unrelated: a Thue--Morse sum corrected by the even moments of Rvachev's up-function, a finite
half-base $q$-binomial expression with a seemingly arbitrary complex translation, a
bit-recursive Taylor reduction, an absolutely convergent binary series at arbitrary argument,
and a finite $q$-weighted discrete approximant whose limit is the signed global Fabius
function.  This companion article identifies the common algebra behind these formulae and
makes their connections explicit.

The central object is the coefficient
\[
 [X^n]\bigl(B_m(X)\mathscr A(X)\bigr),
 \qquad
 B_m(X)=\sum_{h<m}(-1)^{\BinaryDigitSum(h)}\EulerE^{(m-1-h)X},
\]
where $\mathscr A$ is the exponential moment series of the Fabius probability law.  Direct
expansion gives the classical moment formula; probabilistic interpretation gives an
expectation formula; Cauchy's formula gives a contour representation; and geometric divided
differences through $-1,-2,\ldots,-2^n$ give the $q$-binomial formula.  The same coefficient
also produces new finite identities: an exact moment-corrected expansion in centered splines,
its reciprocal deconvolution, and a residue formula for the Gaussian interpolation row.

At arbitrary arguments the two infinite constructions have different architectures.  The
binary series is a genuine telescoping sum whose summands are individually independent of the
complex shift.  The discrete formula is instead a Toeplitz transform of shifted splines; its
finite rows generally depend on the shift, although their limit does not.  Writing the row
generating polynomial as
\[
 W_n(z)=\frac{(z/2;1/2)_n}{(1/2;1/2)_n}
\]
reveals it as a $q$-Richardson filter: it preserves constants and annihilates the geometric
error modes $2^{-p},2^{-2p},\ldots,2^{-np}$.  This yields an explicit uniform
$\BigO(4^{-n})$ error bound.  Most importantly, at a dyadic input $m/2^s$ the discrete
approximant stabilizes exactly for every row $n\ge s$; thus the ``limit formula'' literally
becomes the exact dyadic $q$-formula after the denominator has been resolved.  A synchronized
nearest-dyadic version provides a second, termwise translation-invariant discrete limit.


\section{Scope, conventions, and the formula map}
\label{dyadicweb:sec:scope}

This article is a companion rather than a replacement for the comprehensive exposition
\cite{dyadicweb:mainarticle}.  It assumes the construction and basic analytic theory developed there
and concentrates on one cluster of results: exact dyadic evaluation, binary reduction, finite
half-base $q$-calculus, and discrete limits.  The proofs formalized in Lean under \repo
\cite{dyadicweb:proveit} are the source of the exact identities used below.  Several consequences in
this article are then derived directly from those identities; a formalization map is given in
\cref{dyadicweb:sec:formal-map}.

We use three related functions, with the same convention as the companion article:
\begin{itemize}
\item $F$ is the bounded Fabius distribution function, equal to $0$ on $(-\infty,0]$ and to
      $1$ on $[1,\infty)$;
\item $\RvachevUp$ is the even compactly supported Rvachev density on $[-1,1]$;
\item $\FabiusGlobal$ is the signed global extension satisfying
      $\FabiusGlobal'(x)=2\FabiusGlobal(2x)$ on $\RealNumbers$.
\end{itemize}
They agree on $[0,1]$.  Put
\[
 \ThueMorseSignSymbol_r=(-1)^{\BinaryDigitSum(r)},\qquad
 c_k=\int_{-1}^{1}u^{2k}\RvachevUp(u)\Differential u,
\]
and let $d_n$ denote the half moments normalized by
\begin{equation}
 F(2^{-n})=2^{-\binom n2}\frac{d_n}{n!}.                 \label{dyadicweb:eq:inverse-dyadic-input}
\end{equation}
The normalized moment series is
\begin{equation}
 \mathscr A(X)=\sum_{n\ge0}d_n\frac{X^n}{n!}.             \label{dyadicweb:eq:A-def}
\end{equation}
We reserve $q$ for the fixed base $q=1/2$ and use $Q\in\ComplexNumbers$ for a freely chosen translation.
This distinction is essential: in the expressions $(a;q)_n$ and
$(r-2^k+Q)^{n+k}$ the two symbols play completely different roles.

\subsection{Five presentations of the same value}

The identities studied below fall into five families.  Their superficial differences are
summarized in \cref{dyadicweb:tab:formula-map}; the rest of the article explains the arrows between
rows.

\begin{table}[ht]
\centering
\small
\begin{tabularx}{\textwidth}{@{}>{\raggedright\arraybackslash}p{0.19\textwidth}
>{\raggedright\arraybackslash}p{0.18\textwidth}
>{\raggedright\arraybackslash}p{0.22\textwidth}X@{}}
\toprule
Presentation & Range and status & Principal coefficients & Mechanism \\
\midrule
Moment/Thue--Morse formula & finite and exact at $m/2^n$ & $c_k$ and $\ThueMorseSignSymbol_h$ & coefficient expansion in the centered probability law \\
Half-base $q$-formula & finite and exact at $m/2^n$ & $\GaussianBinomial nk{1/2}$ and $(1/2;1/2)_n$ & geometric divided difference after block cancellation \\
Binary Taylor series & infinite and absolutely convergent for $x\ge0$; finite at dyadics & binary prefixes and $d_j$ & exact residual telescope along the binary expansion \\
Global $q$-binary series & same outer series as the preceding row & a finite $q$-identity inside each scale & termwise replacement of a Taylor polynomial \\
Discrete $q$-limit & one finite row for each $n$, convergent for all $x\in\RealNumbers$ & normalized Gaussian row $\omega_{n,j}$ & Toeplitz transform of shifted finite splines \\
\bottomrule
\end{tabularx}
\caption{The formula families.  The two infinite constructions are not partial-sum versions
of one another.}
\label{dyadicweb:tab:formula-map}
\end{table}

The main structural conclusions can be stated before any calculation.

\begin{boxedremark}
\textbf{The common core.}
The exact dyadic formulae are all coefficient extractors applied to the same entire function
$B_m\mathscr A$.  The binary series repeatedly reuses its Taylor coefficients at successive
binary scales.  The discrete formula applies a normalized $q$-Pochhammer polynomial to the
sequence of finite-spline degrees.  At a dyadic input these two arbitrary-argument mechanisms
both become finite: the binary series terminates, while the discrete rows stabilize exactly.
\end{boxedremark}

\subsection{The exact dyadic formula taken as input}

For $0\le m\le2^n$ the moment formula proved in the formal development is
\begin{equation}
 \boxed{
 F\!\left(\frac{m}{2^n}\right)
 =\frac{2^{-\binom{n+1}{2}}}{n!}
 \sum_{h=0}^{m-1}\ThueMorseSignSymbol_h
 \sum_{k=0}^{\Floor{n/2}}
 \binom n{2k}c_k(2m-2h-1)^{n-2k}.}                        \label{dyadicweb:eq:input-moment-dyadic}
\end{equation}
The same finite expression, interpreted with $\FabiusGlobal$ on the left, is valid for every
$m,n\in\NaturalNumbers$.  The $q$-binomial version is
\begin{align}
 \FabiusGlobal(m/2^n)
 &=\frac{1}{2^{n^2}(1/2;1/2)_n}
 \sum_{k=0}^{n}
 \frac{\GaussianBinomial nk{1/2}}{4^{\binom k2}(n+k)!}             \notag\\
 &\qquad\times
 \sum_{r=0}^{m2^k-1}\ThueMorseSignSymbol_r(r-m2^k+Q)^{n+k},               \label{dyadicweb:eq:input-q-dyadic}
\end{align}
valid for every $Q\in\ComplexNumbers$ and independent of $Q$ as a complete finite expression.  A central
goal below is to explain why \eqref{dyadicweb:eq:input-moment-dyadic} and
\eqref{dyadicweb:eq:input-q-dyadic} are two readings of one coefficient.

\section{The master coefficient and its probabilistic meaning}
\label{dyadicweb:sec:master}

\subsection{The probability law behind the moment series}

Let $U_1,U_2,\ldots$ be independent random variables, each uniformly distributed on
$[0,1]$, and set
\begin{equation}
 Y=\sum_{j=1}^{\infty}\frac{U_j}{2^j},
 \qquad Z=2Y-1.                                             \label{dyadicweb:eq:Y-Z-def}
\end{equation}
Then $Y\in[0,1]$, its distribution function is $F$, and $Z$ has density $\RvachevUp$ on $[-1,1]$.
Consequently
\begin{equation}
 d_n=\Expectation(Y^n),\qquad c_k=\Expectation(Z^{2k}),\qquad \Expectation(Z^{2k+1})=0. \label{dyadicweb:eq:prob-moments}
\end{equation}
The moment generating function of one uniform variable is
\[
 E_0(X)=\frac{\EulerE^X-1}{X}=\int_0^1\EulerE^{Xu}\Differential u,
\]
with the removable value $E_0(0)=1$.  Independence gives the locally uniformly convergent
entire product
\begin{equation}
 \mathscr A(X)=\Expectation(\EulerE^{XY})
 =\prod_{j=1}^{\infty}E_0(X/2^j).                          \label{dyadicweb:eq:A-product}
\end{equation}
Separating the first factor, or merely replacing $X$ by $2X$ in the product, gives the
refinement equation
\begin{equation}
 \mathscr A(2X)=E_0(X)\mathscr A(X).                       \label{dyadicweb:eq:A-refinement}
\end{equation}
The centered moment series
\begin{equation}
 \mathscr C(T)=\Expectation(\EulerE^{TZ})
 =\sum_{k\ge0}c_k\frac{T^{2k}}{(2k)!}                     \label{dyadicweb:eq:C-def}
\end{equation}
satisfies the exact change of coordinates
\begin{equation}
 \boxed{\mathscr A(X)=\EulerE^{X/2}\mathscr C(X/2).}           \label{dyadicweb:eq:A-C-transform}
\end{equation}
The familiar hyperbolic-sine form follows by writing
$E_0(t)=\EulerE^{t/2}\sinh(t/2)/(t/2)$:
\begin{align}
 \mathscr A(X)
 &=\EulerE^{X/2}\prod_{j=1}^{\infty}
   \frac{\sinh(X/2^{j+1})}{X/2^{j+1}},                    \label{dyadicweb:eq:A-sinh-product}\\
 \mathscr C(T)
 &=\prod_{j=1}^{\infty}\frac{\sinh(T/2^j)}{T/2^j}.        \label{dyadicweb:eq:C-sinh-product}
\end{align}
Thus the moment series used in exact arithmetic is the real-exponential counterpart of the
sinc product in Fourier analysis.

\subsection{The Thue--Morse prefix series}

For $m\in\NaturalNumbers$ define
\begin{equation}
 B_m(X)=\sum_{h=0}^{m-1}\ThueMorseSignSymbol_h\EulerE^{(m-1-h)X},
 \qquad B_0(X)=0.                                         \label{dyadicweb:eq:Bm-def}
\end{equation}
This finite exponential polynomial records the numerator $m$ and nothing else.  Its two
factors have distinct roles:
\begin{itemize}
\item $B_m$ carries the finite Thue--Morse prefix and therefore the binary numerator;
\item $\mathscr A$ carries the universal Fabius moments and therefore the analytic law.
\end{itemize}
Their product is the promised master object.

\begin{theorem}[Master coefficient identity]
\label{dyadicweb:thm:master-coefficient}
For every $m,n\in\NaturalNumbers$,
\begin{equation}
 \boxed{
 \FabiusGlobal\!\left(\frac{m}{2^n}\right)
 =2^{-\binom n2}[X^n]\bigl(B_m(X)\mathscr A(X)\bigr).}    \label{dyadicweb:eq:master-coefficient}
\end{equation}
If $0\le m\le2^n$, the left side is the bounded value $F(m/2^n)$.
\end{theorem}

\begin{proof}
Put $t_h=2m-2h-1$.  By \eqref{dyadicweb:eq:C-def},
\[
 \sum_{k=0}^{\Floor{n/2}}\binom n{2k}c_kt_h^{n-2k}
 =n![W^n]\bigl(\EulerE^{t_hW}\mathscr C(W)\bigr).
\]
Using \eqref{dyadicweb:eq:A-C-transform} at $X=2W$ gives
$\mathscr C(W)=\EulerE^{-W}\mathscr A(2W)$.  Since $t_h-1=2(m-1-h)$,
\begin{align*}
 n![W^n]\bigl(\EulerE^{t_hW}\mathscr C(W)\bigr)
 &=n![W^n]\bigl(\EulerE^{2(m-1-h)W}\mathscr A(2W)\bigr)\\
 &=n!2^n[X^n]\bigl(\EulerE^{(m-1-h)X}\mathscr A(X)\bigr).
\end{align*}
Multiply by $\ThueMorseSignSymbol_h$, sum over $h<m$, and insert the result in
\eqref{dyadicweb:eq:input-moment-dyadic}.  The power of two becomes
$2^{n-\binom{n+1}{2}}=2^{-\binom n2}$.
\end{proof}

The theorem compresses the double sum in \eqref{dyadicweb:eq:input-moment-dyadic} into one coefficient,
but it also expands naturally in several different directions.

\begin{corollary}[Expectation form]
\label{dyadicweb:cor:expectation-dyadic}
For every $m,n\in\NaturalNumbers$,
\begin{align}
 \FabiusGlobal(m/2^n)
 &=\frac{2^{-\binom n2}}{n!}
   \sum_{h=0}^{m-1}\ThueMorseSignSymbol_h\,
   \Expectation\bigl(m-1-h+Y\bigr)^n                              \label{dyadicweb:eq:expectation-Y}\\
 &=\frac{2^{-\binom{n+1}{2}}}{n!}
   \sum_{h=0}^{m-1}\ThueMorseSignSymbol_h\,
   \Expectation\bigl(2m-2h-1+Z\bigr)^n.                            \label{dyadicweb:eq:expectation-Z}
\end{align}
\end{corollary}

\begin{proof}
Because $\mathscr A(X)=\Expectation(\EulerE^{XY})$,
\[
 [X^n]\bigl(\EulerE^{aX}\mathscr A(X)\bigr)
 =\frac1{n!}\Expectation(a+Y)^n.
\]
Apply this termwise to $B_m$.  The second form follows from
$m-1-h+Y=(2m-2h-1+Z)/2$.
\end{proof}

Expanding the centered expectation in \eqref{dyadicweb:eq:expectation-Z}, using symmetry to discard the
odd moments, recovers \eqref{dyadicweb:eq:input-moment-dyadic} immediately.  In this sense the
Thue--Morse/moment formula is simply the binomial theorem applied inside one probability
expectation.

\begin{corollary}[Cauchy coefficient representation]
\label{dyadicweb:cor:cauchy-X}
For any $\rho>0$,
\begin{equation}
 \FabiusGlobal(m/2^n)
 =\frac{2^{-\binom n2}}{2\pi\ImaginaryUnit}
  \int_{|X|=\rho}\frac{B_m(X)\mathscr A(X)}{X^{n+1}}\Differential X. \label{dyadicweb:eq:cauchy-X}
\end{equation}
\end{corollary}

This contour formula is elementary, but useful conceptually: the $q$-binomial formula will
turn out to be a second residue computation, now in an auxiliary geometric node variable.

\subsection{The numerator and the law separate completely}

Expanding \eqref{dyadicweb:eq:master-coefficient} without centering gives another exact finite form:
\begin{equation}
 \FabiusGlobal(m/2^n)
 =\frac{2^{-\binom n2}}{n!}
 \sum_{h=0}^{m-1}\ThueMorseSignSymbol_h
 \sum_{j=0}^{n}\binom njd_j(m-1-h)^{n-j}.                 \label{dyadicweb:eq:d-moment-prefix}
\end{equation}
The $d_j$ are universal and may be precomputed once.  The numerator enters only through the
finite power sums
\(
 \sum_{h<m}\ThueMorseSignSymbol_h(m-1-h)^r.
\)
This is often a more transparent bridge to the Taylor reduction series than the centered
$c_k$ formula, because the reduction polynomials themselves are written in the $d_j$.

\begin{remark}[Two independent cancellations]
The Thue--Morse prefix in \eqref{dyadicweb:eq:d-moment-prefix} may annihilate low ordinary powers when
$m$ is a complete binary block.  The $q$-binomial formula introduces a second cancellation
across the geometric scale index.  These cancellations live in different variables and should
not be conflated: one acts on $h$ or $r$ inside a binary block, the other on the scale node
$-2^k$.
\end{remark}

\section{Moment coordinates and alternative exact representations}
\label{dyadicweb:sec:moment-alternatives}

\subsection[The triangular transforms between the c and d sequences]{The triangular transforms between $c_k$ and $d_n$}

The affine relation $Z=2Y-1$ gives both directions of the moment change of coordinates.

\begin{proposition}[Moment transforms]
\label{dyadicweb:prop:moment-transforms}
For every $n,k\in\NaturalNumbers$,
\begin{align}
 d_n
 &=2^{-n}\sum_{k=0}^{\Floor{n/2}}\binom n{2k}c_k,
                                                               \label{dyadicweb:eq:d-from-c-new}\\
 c_k
 &=\sum_{j=0}^{2k}(-1)^j2^j\binom{2k}{j}d_j.              \label{dyadicweb:eq:c-from-d-new}
\end{align}
\end{proposition}

\begin{proof}
Since $Y=(1+Z)/2$, expand $Y^n$ and use the vanishing odd moments of $Z$.  Conversely expand
$(2Y-1)^{2k}$; because $2k$ is even, the sign of the term of degree $j$ is $(-1)^j$.
\end{proof}

The first formula explains the odd-looking shift $2m-2h-1$ in
\eqref{dyadicweb:eq:input-moment-dyadic}: it is precisely the centered coordinate obtained after
converting the half moments to the even moments of $\RvachevUp$.

\subsection{An infinite-cube composition formula}

Each factor in \eqref{dyadicweb:eq:A-product} has the expansion
\[
 E_0(X/2^j)=\sum_{r\ge0}\frac{X^r}{2^{jr}(r+1)!}.
\]
Taking the coefficient of $X^n$ gives an exact positive series over finite-support weak
compositions.

\begin{proposition}[Composition representation of the half moments]
\label{dyadicweb:prop:composition-d}
For every $n\in\NaturalNumbers$,
\begin{equation}
 d_n=n!\!
 \sum_{\substack{r_1,r_2,\ldots\ge0\\r_1+r_2+\cdots=n}}
 \prod_{j\ge1}\frac{2^{-jr_j}}{(r_j+1)!}.                 \label{dyadicweb:eq:composition-d}
\end{equation}
Only finitely many $r_j$ are nonzero in each summand, and the series converges absolutely.
Equivalently,
\begin{equation}
 d_n=\lim_{J\to\infty}\int_{[0,1]^J}
       \left(\sum_{j=1}^{J}\frac{u_j}{2^j}\right)^n
       \Differential u_1\cdots\Differential u_J.                              \label{dyadicweb:eq:cube-d}
\end{equation}
\end{proposition}

This gives a $q$-free, manifestly positive representation of every inverse-dyadic value:
\begin{equation}
 F(2^{-n})=\frac{2^{-\binom n2}}{n!}d_n
 =2^{-\binom n2}
 \sum_{r_1+r_2+\cdots=n}
 \prod_{j\ge1}\frac{2^{-jr_j}}{(r_j+1)!}.                \label{dyadicweb:eq:inverse-composition}
\end{equation}

Grouping the nonzero entries of the weak composition gives a finite ordered-composition
form.  If $\rho=(r_1,\ldots,r_\ell)\vDash n$ and
$s_j=r_1+\cdots+r_j$, then
\begin{equation}
 \boxed{
 F(2^{-n})=2^{-\binom n2}
 \sum_{\rho\vDash n}\prod_{j=1}^{\ell}
 \frac1{(2^{s_j}-1)(r_j+1)!}.}
 \label{dyadicweb:eq:ordered-composition}
\end{equation}
There is also an independent recurrence-path proof.  Apply
\cref{integration:thm:triangular-path-solver} to the normalized recurrence with
\[
 w(k,n)=\frac1{(2^n-1)(n-k+1)!},\qquad 0\leq k<n.
\]
A path \(0=s_0<s_1<\cdots<s_\ell=n\) corresponds to the composition
\(r_j=s_j-s_{j-1}\); its chronological edge product is exactly the product in
\eqref{dyadicweb:eq:ordered-composition}.  Thus the geometric-simplex calculation below
and the generic triangular solver give two genuinely different derivations of the one
canonical formula.
Indeed, choosing distinct active scales $1\le j_1<\cdots<j_\ell$ produces the
geometric simplex sum
\begin{equation}
 \sum_{1\le j_1<\cdots<j_\ell}2^{-\sum_i j_ir_i}
 =\prod_{h=1}^{\ell}\frac1{2^{t_h}-1},\qquad
 t_h=r_h+\cdots+r_\ell.
 \label{dyadicweb:eq:geometric-simplex}
\end{equation}
To prove it, write $j_i=i+g_1+\cdots+g_i$ with $g_h\ge0$.  The independent
geometric series contribute $(1-2^{-t_h})^{-1}$, and the initial powers cancel
because $\sum_ht_h=\sum_iir_i$; reversing the composition changes the tail sums
$t_h$ into the prefix sums $s_j$.

The same triangular recurrence has a compact determinant form.  Put
$a_n=d_n/n!$, $a_0=1$, and
\begin{equation}
 \beta_{n,k}=\frac1{(2^n-1)(n-k+1)!},\qquad 0\le k<n.
 \label{dyadicweb:eq:Hessenberg-beta}
\end{equation}
Let $H_n$ be the lower-Hessenberg matrix whose last-row pattern is
$(\beta_{n,0},\ldots,\beta_{n,n-1})$, whose entries below the superdiagonal are
$\beta_{i,j}$, and whose superdiagonal entries are $-1$.  Then
\begin{equation}
 \boxed{F(2^{-n})=2^{-\binom n2}\det H_n.}
 \label{dyadicweb:eq:Hessenberg-determinant}
\end{equation}
Expansion along the last row gives
$\det H_n=\sum_{k<n}\beta_{n,k}\det H_k$, exactly the recurrence for $a_n$.
Expanding the determinant by Hessenberg paths recovers
\eqref{dyadicweb:eq:ordered-composition}.

\subsection{Cumulants, Bernoulli numbers, and Bell polynomials}

The infinite product also gives a compact partition formula.  With Bernoulli numbers in the
convention $B_2=1/6$, $B_4=-1/30$, one has
\begin{equation}
 \log\frac{\EulerE^X-1}{X}
 =\frac X2+\sum_{r\ge1}\frac{B_{2r}}{2r(2r)!}X^{2r}.      \label{dyadicweb:eq:log-exprel}
\end{equation}
Summing this identity over the dyadic scales in \eqref{dyadicweb:eq:A-product} yields
\begin{equation}
 \log\mathscr A(X)
 =\frac X2+
  \sum_{r\ge1}
  \frac{B_{2r}}{2r(2r)!(2^{2r}-1)}X^{2r}.                 \label{dyadicweb:eq:log-A}
\end{equation}
Thus the cumulants $\kappa_j$ of $Y$ are
\begin{equation}
 \kappa_1=\frac12,\qquad
 \kappa_{2r}=\frac{B_{2r}}{2r(2^{2r}-1)},\qquad
 \kappa_{2r+1}=0\quad(r\ge1).                            \label{dyadicweb:eq:Y-cumulants}
\end{equation}

\begin{theorem}[Bell-polynomial representation]
\label{dyadicweb:thm:bell-d}
Let $\mathcal B_n$ be the complete exponential Bell polynomial.  Then
\begin{equation}
 d_n=\mathcal B_n(\kappa_1,\ldots,\kappa_n),               \label{dyadicweb:eq:bell-d}
\end{equation}
and explicitly
\begin{equation}
 d_n=n!\!
 \sum_{\substack{m_1,\ldots,m_n\ge0\\
                  1m_1+2m_2+\cdots+nm_n=n}}
 \prod_{r=1}^{n}\frac1{m_r!}
 \left(\frac{\kappa_r}{r!}\right)^{m_r}.                 \label{dyadicweb:eq:partition-d}
\end{equation}
Consequently
\begin{equation}
 \boxed{
 F(2^{-n})=\frac{2^{-\binom n2}}{n!}
 \mathcal B_n(\kappa_1,\ldots,\kappa_n).}                 \label{dyadicweb:eq:inverse-bell}
\end{equation}
\end{theorem}

\begin{proof}
Equation \eqref{dyadicweb:eq:log-A} is the standard cumulant expansion
$\log\mathscr A(X)=\sum_{j\ge1}\kappa_jX^j/j!$.  Exponentiating and comparing coefficients
is exactly the defining generating identity for the complete Bell polynomials.
\end{proof}

The representation is useful when Bernoulli arithmetic is already available.  It also gives
the moment recurrence
\begin{equation}
 d_{n+1}=\sum_{j=0}^{n}\binom nj\kappa_{j+1}d_{n-j},       \label{dyadicweb:eq:cumulant-moment-recurrence}
\end{equation}
while \eqref{dyadicweb:eq:A-refinement} gives the rational recurrence
\begin{equation}
 (2^n-1)d_n
 =\sum_{r=1}^{n}\binom nr\frac{d_{n-r}}{r+1}
 \qquad(n\ge1).                                           \label{dyadicweb:eq:d-recurrence-companion}
\end{equation}

\begin{table}[ht]
\centering
\small
\begin{tabular}{@{}c c c c@{}}
\toprule
$n$ & $d_n$ & $c_n$ & $F(2^{-n})$ \\
\midrule
0 & $1$ & $1$ & $1$ \\
1 & $1/2$ & $1/9$ & $1/2$ \\
2 & $5/18$ & $19/675$ & $5/72$ \\
3 & $1/6$ & $583/59535$ & $1/288$ \\
4 & $143/1350$ & $132809/32531625$ & $143/2073600$ \\
5 & $19/270$ & $46840699/24405225075$ & $19/33177600$ \\
\bottomrule
\end{tabular}
\caption{The columns use different index conventions: $c_n$ is the $2n$th centered moment,
whereas $d_n$ is the $n$th half moment.}
\label{dyadicweb:tab:moments}
\end{table}

\subsection{Exact moment correction of the matched spline}
\label{dyadicweb:subsec:moment-corrected-spline}

The centered finite spline of degree $p\ge1$ is
\begin{equation}
 S_p(x)=\frac{(-1)^p}{2^{\binom p2}p!}
 \sum_{r=0}^{N_p(x)-1}\ThueMorseSignSymbol_r\left(r-2^px+\frac12\right)^p,
 \qquad
 N_p(x)=\max\!\left(0,\Floor{2^px+\frac12}\right). \label{dyadicweb:eq:centered-spline-local}
\end{equation}
At degree zero set
\begin{equation}
 S_0(x)=\sum_{r=0}^{N_0(x)-1}\ThueMorseSignSymbol_r.                        \label{dyadicweb:eq:S0-local}
\end{equation}
The companion article proves the global estimate
\begin{equation}
 \sup_{x\in\RealNumbers}|S_p(x)-\FabiusGlobal(x)|\le2^{-p}.             \label{dyadicweb:eq:centered-global-error-input}
\end{equation}
At the matched grid point $m/2^p$ the cutoff is $m$ and
\begin{equation}
 S_p(m/2^p)=\frac{1}{2^{\binom{p+1}{2}}p!}
 \sum_{h=0}^{m-1}\ThueMorseSignSymbol_h(2m-2h-1)^p.                       \label{dyadicweb:eq:matched-spline}
\end{equation}
The $k=0$ part of \eqref{dyadicweb:eq:input-moment-dyadic} is exactly this matched spline.  The remaining
moments supply a finite correction hierarchy.

\begin{theorem}[Moment-corrected matched-spline identity]
\label{dyadicweb:thm:moment-corrected-spline}
For all $m,n\in\NaturalNumbers$,
\begin{equation}
 \boxed{
 \FabiusGlobal(m/2^n)=
 \sum_{k=0}^{\Floor{n/2}}
 \frac{c_k}{(2k)!\,2^{k(2n-2k+1)}}
 S_{n-2k}\!\left(\frac{m}{2^{n-2k}}\right).}              \label{dyadicweb:eq:moment-corrected-spline}
\end{equation}
For $0\le m\le2^n$ the left side is $F(m/2^n)$.  The spline arguments on the right are
understood globally and may lie outside $[0,1]$.
\end{theorem}

\begin{proof}
In \eqref{dyadicweb:eq:input-moment-dyadic} put $p=n-2k$.  By
\eqref{dyadicweb:eq:matched-spline},
\[
 \sum_{h=0}^{m-1}\ThueMorseSignSymbol_h(2m-2h-1)^p
 =2^{\binom{p+1}{2}}p!\,S_p(m/2^p).
\]
Moreover
\[
 \frac1{n!}\binom n{2k}p!=\frac1{(2k)!},
\]
and
\[
 \binom{n+1}{2}-\binom{p+1}{2}=k(2n-2k+1).
\]
Substitution gives \eqref{dyadicweb:eq:moment-corrected-spline}.
\end{proof}

The formula has three noteworthy features.
\begin{enumerate}
\item Its first term is the obvious finite-spline approximation $S_n(m/2^n)$.
\item The correction drops the spline degree by two at each step, reflecting the evenness of
      $\RvachevUp$ and the disappearance of odd centered moments.
\item The numerator $m$ remains fixed, while the denominator is coarsened from $2^n$ to
      $2^{n-2k}$.  Thus the formula couples degree reduction to dyadic rescaling.
\end{enumerate}

\begin{corollary}[Exact defect of the matched spline]
\label{dyadicweb:cor:matched-defect}
For every $m,n\in\NaturalNumbers$,
\begin{equation}
 \FabiusGlobal(m/2^n)-S_n(m/2^n)
 =\sum_{k=1}^{\Floor{n/2}}
 \frac{c_k}{(2k)!\,2^{k(2n-2k+1)}}
 S_{n-2k}\!\left(\frac{m}{2^{n-2k}}\right).               \label{dyadicweb:eq:matched-defect}
\end{equation}
\end{corollary}

This is stronger than an error estimate: it identifies the error exactly as a finite sequence
of lower-degree global splines.  The small factor in the first correction is
$c_1/(2\cdot2^{2n-1})=1/(9\cdot2^{2n})$.

\subsection{The reciprocal correction transform}

The previous relation is triangular and can be inverted.  Define reciprocal centered moments
$\gamma_k$ by
\begin{equation}
 \frac1{\mathscr C(T)}=
 \sum_{k\ge0}\gamma_k\frac{T^{2k}}{(2k)!}.                \label{dyadicweb:eq:gamma-def}
\end{equation}
Thus
\begin{equation}
 \gamma_0=1,
 \qquad
 \gamma_n=-\sum_{k=0}^{n-1}\binom{2n}{2k}\gamma_kc_{n-k}
 \quad(n\ge1),                                             \label{dyadicweb:eq:gamma-recurrence}
\end{equation}
and the first values are
\begin{equation}
 1,\quad-\frac19,\quad\frac{31}{675},\quad
 -\frac{2347}{59535},\quad\frac{1904369}{32531625},\ldots. \label{dyadicweb:eq:gamma-values}
\end{equation}

\begin{theorem}[Reciprocal moment deconvolution]
\label{dyadicweb:thm:reciprocal-spline}
For every $m,n\in\NaturalNumbers$,
\begin{equation}
 \boxed{
 S_n(m/2^n)=
 \sum_{k=0}^{\Floor{n/2}}
 \frac{\gamma_k}{(2k)!\,2^{k(2n-2k+1)}}
 \FabiusGlobal\!\left(\frac{m}{2^{n-2k}}\right).}            \label{dyadicweb:eq:reciprocal-spline}
\end{equation}
\end{theorem}

\begin{proof}
For a fixed numerator $m$, introduce
\[
 P_p(m)=\sum_{h=0}^{m-1}\ThueMorseSignSymbol_h(2m-2h-1)^p,
 \qquad
 G_p(m)=2^{\binom{p+1}{2}}p!\,\FabiusGlobal(m/2^p).
\]
The exact moment formula says, in exponential generating functions,
\[
 \sum_{p\ge0}G_p(m)\frac{T^p}{p!}
 =\mathscr C(T)\sum_{p\ge0}P_p(m)\frac{T^p}{p!}.
\]
Multiply by $1/\mathscr C(T)$, compare coefficients, and then use
$P_p(m)=2^{\binom{p+1}{2}}p!S_p(m/2^p)$.  The same triangular exponent simplification as in
\cref{dyadicweb:thm:moment-corrected-spline} gives \eqref{dyadicweb:eq:reciprocal-spline}.
\end{proof}

Equations \eqref{dyadicweb:eq:moment-corrected-spline} and \eqref{dyadicweb:eq:reciprocal-spline} are a finite
convolution/deconvolution pair.  They show that the exact dyadic values and the matched centered
splines carry precisely the same information once the universal moment series $\mathscr C$ is
known.

\begin{boxedremark}
\textbf{First bridge to the discrete theory.}
The centered spline is not merely an approximation from which exact dyadic arithmetic is
separate.  At a matched grid point, the exact value is obtained by a finite universal moment
correction of that spline, and the correction is invertible.  Later, the half-base
$q$-Pochhammer row will give a second exact correction, using consecutive spline degrees at
the \emph{same} argument rather than parity-separated degrees at rescaled arguments.
\end{boxedremark}

\section[The half-base q-formula as geometric coefficient extraction]{The half-base $q$-formula as geometric coefficient extraction}
\label{dyadicweb:sec:q-extraction}

\subsection{The Gaussian row is a divided difference}

For a general base $q$, write
\begin{equation}
 (a;q)_n=\prod_{j=0}^{n-1}(1-aq^j),                       \label{dyadicweb:eq:qpoch-def-local}
\end{equation}
and let $\GaussianBinomial nkq$ denote the Gaussian binomial coefficient.  At $q=1/2$ the finite
$q$-binomial theorem is
\begin{equation}
 (z;1/2)_n=\sum_{k=0}^{n}(-1)^k2^{-\binom k2}
 \GaussianBinomial nk{1/2}z^k.                                      \label{dyadicweb:eq:finite-q-binomial-local}
\end{equation}
Put
\begin{equation}
 x_k=-2^k,
 \qquad
 w_{n,k}=(-1)^k2^{-\binom k2}\GaussianBinomial nk{1/2},
 \qquad
 \mathfrak m_n=\prod_{\ell=1}^{n}(2^\ell-1).              \label{dyadicweb:eq:q-nodes-weights}
\end{equation}
Then
\begin{equation}
 \sum_{k=0}^{n}w_{n,k}x_k^d=
 \begin{cases}
 0,&0\le d<n,\\
 \mathfrak m_n,&d=n,
 \end{cases}                                               \label{dyadicweb:eq:q-node-moments}
\end{equation}
and
\begin{equation}
 \frac{w_{n,k}}{\mathfrak m_n}
 =\frac1{\prod_{\ell\ne k}(x_k-x_\ell)},
 \qquad
 (1/2;1/2)_n=2^{-\binom{n+1}{2}}\mathfrak m_n.             \label{dyadicweb:eq:q-barycentric-normalization}
\end{equation}
Equivalently, the half-base coefficients contain only integral Gaussian arithmetic and an
explicit power of two:
\begin{equation}
 \GaussianBinomial nk{1/2}=2^{-k(n-k)}\GaussianBinomial nk2.
 \label{dyadicweb:eq:qbinomial-inversion}
\end{equation}
Thus the functional
\begin{equation}
 \Dq_n[P]=\frac1{\mathfrak m_n}\sum_{k=0}^{n}w_{n,k}P(-2^k) \label{dyadicweb:eq:Dq-def}
\end{equation}
is the divided difference $[-1,-2,\ldots,-2^n]P$.  For every polynomial of degree at most
$n$, it returns the leading coefficient.

\begin{proposition}[Residue form of the geometric extractor]
\label{dyadicweb:prop:q-residue-extractor}
Let
\[
 \Pi_n(z)=\prod_{k=0}^{n}(z+2^k),
\]
and let $P$ be a polynomial of degree at most $n$.  If $\Gamma$ encloses all nodes
$-1,-2,\ldots,-2^n$, then
\begin{equation}
 [z^n]P(z)=\Dq_n[P]
 =\frac{1}{2\pi\ImaginaryUnit}\int_\Gamma\frac{P(z)}{\Pi_n(z)}\Differential z. \label{dyadicweb:eq:q-residue-extractor}
\end{equation}
\end{proposition}

\begin{proof}
The residues at the simple poles $x_k=-2^k$ are
$P(x_k)/\Pi_n'(x_k)=P(x_k)/\prod_{\ell\ne k}(x_k-x_\ell)$, giving the barycentric sum in
\eqref{dyadicweb:eq:Dq-def}.  Lagrange interpolation identifies that sum with the leading coefficient.
\end{proof}

\begin{remark}[The same extractor at a general geometric base]
For $0<q<1$, let the nodes be $1,q^{-1},\ldots,q^{-n}$.  Lagrange interpolation gives,
for every polynomial $P$ of degree at most $n$,
\begin{equation}
 [t^n]P(t)
 =\frac{(-1)^nq^{\binom{n+1}{2}}}{(q;q)_n}
  \sum_{k=0}^{n}(-1)^kq^{\binom k2}\GaussianBinomial nkq P(q^{-k}).
 \label{dyadicq:eq:q-leading-extractor-general}
\end{equation}
The corresponding barycentric denominator is
\begin{equation}
 \frac1{\displaystyle\prod_{j\ne k}(q^{-k}-q^{-j})}
 =\frac{(-1)^{n-k}q^{\binom{n+1}{2}+\binom k2}}
 {(q;q)_n}\GaussianBinomial nkq.                                   \label{dyadicq:eq:barycentric-denominator}
\end{equation}
Equivalently, the finite $q$-binomial theorem kills every lower monomial after the
specialization $z=q^{-m}$, while the surviving top term is normalized by
\begin{equation}
 (q^{-n};q)_n=(-1)^nq^{-\binom{n+1}{2}}(q;q)_n.           \label{dyadicq:eq:q-survivor}
\end{equation}
Thus the half-base stencil above is the $q=1/2$ instance of a generic geometric divided
difference, not an isolated Gaussian-binomial coincidence.
\end{remark}

\begin{remark}[Division-free scaled characteristic and \(q\)-difference functional]
\label{dyadicq:rem:scaled-qdiff}
There is a stronger denominator-free row, valid without ordering or
distinctness assumptions.  Over a commutative ring put
\[
 K^{-1}_{q,s}(n,k)
 =(-s)^{n-k}q^{\binom{n-k}{2}}\GaussianBinomial{n}{k}{q}.
\]
Then, for arbitrary \(q,s,z\) and every \(n\ge0\),
\begin{equation}
 \sum_{k=0}^{n}K^{-1}_{q,s}(n,k)z^k
 =\prod_{j=0}^{n-1}(z-sq^j).
 \label{dyadicq:eq:scaled-qdiff-characteristic}
\end{equation}
This is the scaled characteristic polynomial.  Its \(s=1,\ z=q^d\)
specialization gives every monomial moment,
\begin{equation}
 \sum_{k=0}^{n}K^{-1}_{q,1}(n,k)(q^k)^d
 =\prod_{j=0}^{n-1}(q^d-q^j),
 \label{dyadicq:eq:qdiff-moment}
\end{equation}
and therefore vanishes whenever \(d<n\).

The exact polynomial statement survives arbitrary scalar extension.  If
\(A\) is a semiring, \(R\) a commutative ring,
\(\varphi\colon A\to R\) a semiring homomorphism, and \(P\in A[X]\) has
degree at most \(n\), then
\begin{equation}
 \sum_{k=0}^{n}K^{-1}_{q,1}(n,k)\,P^\varphi(q^k)
 =\varphi([X^n]P)\prod_{j=0}^{n-1}(q^n-q^j),
 \label{dyadicq:eq:qdiff-top-coeff}
\end{equation}
where \(P^\varphi\) is obtained by mapping coefficients through \(\varphi\).
If \(\deg P<n\), the left side is exactly zero.  This includes \(n=0\) and
the zero polynomial.  The nodes may collide and the surviving product in
\eqref{dyadicq:eq:qdiff-top-coeff} may itself vanish; there is no division,
nonzero or invertible base, domain, characteristic, topology, or convergence
hypothesis.

These four claims are exactly the four public theorems of
\path{QDifferenceAnnihilation.lean}:
\lean{Fabius.sum_scaledGaussianBinomialInverseKernel_mul_pow},
\lean{Fabius.sum_gaussianBinomialInverseKernel_mul_geometric_pow},
\texttt{Fabius.qDifference\_sum\_eval}\textsubscript{2}\allowbreak
\texttt{\_eq\_map\_coeff\_mul}, and
\texttt{Fabius.qDifference\_sum\_eval}\textsubscript{2}\allowbreak
\texttt{\_eq\_zero\_of\_degree\_lt}.
The theorem
\lean{Fabius.geometricQBinomialWeightNumerator_eq_scaledGaussianBinomialInverseKernel},
now in \path{GeometricQBinomialLagrange.lean}, identifies the
denominator-free geometric Lagrange numerator with
\(K^{-1}_{q,q}(n,k)\) at every pair of natural indices, including above the
diagonal.  Thus its generating polynomial is the \(s=q\) specialization of
\eqref{dyadicq:eq:scaled-qdiff-characteristic}.  This global numerator
identity does not weaken the field, finite-node injectivity, or in-range
assumptions in the normalized quotient formula
\eqref{dyadicq:eq:q-leading-extractor-general}.

For completeness, \path{QBinomialInversionSpecializations.lean} now has
exactly two public definitions and four public theorems:
\lean{Fabius.qGaussianResidualCoeff},
\lean{Fabius.qGaussianReconstructionCoeff},
\lean{Fabius.qGaussianResidualCoeff_eq},
\lean{Fabius.qGaussianReconstructionCoeff_eq},
\lean{Fabius.qGaussianReconstructionCoeff_residualCoeff_delta}, and
\lean{Fabius.qGaussianResidualCoeff_reconstructionCoeff_delta}.  At base
\(q^2\) and independent scale \(-q\), these give the two rows
\[
 (-q)^{n-k}\GaussianBinomial{n}{k}{q^2},
 \qquad q^{(n-k)^2}\GaussianBinomial{n}{k}{q^2},
\]
The two definitions and these two pointwise closed-form theorems require only
an arbitrary, possibly noncommutative ring.  Exactly the two convolution
theorems require a commutative ring, and prove that both total finite orders
are the Kronecker delta.

\medskip
\noindent The remaining finite geometric \(q\)-layer has the following
exhaustive public surface.
\begin{itemize}
\item\label{dyadicq:item:elementary-qpochhammer}
\path{QPochhammerElementaryIdentities.lean} has exactly 13 public theorems:
\path{Fabius.finiteQPochhammerIn_base_reversal_units},
\path{Fabius.finiteQPochhammerIn_inv_base_reversal_units},
\path{Fabius.finiteQPochhammerIn_base_reversal},
\path{Fabius.finiteQPochhammerIn_inv_base_reversal},
\path{Fabius.prod_pow_sub_pow_eq_finiteQPochhammerIn},
\path{Fabius.pow_mul_finiteQPochhammerIn_inv_pow_eq},
\path{Fabius.finiteQPochhammerIn_inv_pow_eq_self_div},
\path{Fabius.finiteQPochhammerIn_inv_pow_eq_zero_of_lt},
\path{Fabius.one_sub_mul_gaussianBinomial_one},
\path{Fabius.gaussianBinomial_adjacent_mul},
\path{Fabius.gaussianBinomial_row_adjacent_mul},
\path{Fabius.gaussianBinomial_adjacent_div}, and
\path{Fabius.gaussianBinomial_row_adjacent_div}.
The two unit reversals, the root-safe terminating numerator, the first-column
clearer, and both adjacent cross identities hold over every commutative ring.
In particular the polynomial identities remain valid at roots of unity, and
the two adjacent cross identities are total in all natural \(n,k\): zero
extension makes both sides vanish on and above the row boundary.  The two
field reversal wrappers require exactly \(a\ne0\) and \(q\ne0\).  The cleared
terminating formula and the \(k>N\) zero theorem require \(q\ne0\), while the
displayed quotient for \(k\le N\) additionally requires
\((q;q)_{N-k}\ne0\).  The two adjacent quotient theorems instead retain the
strict range \(k<n\) and require exactly their displayed Gaussian and
linear-factor denominators to be nonzero; neither assumes \(q\ne0\).

\item The six theorems of \path{GeometricCompleteHomogeneous.lean} are
\path{Fabius.completeHomogeneousEval_geometric},
\path{Fabius.completeHomogeneousEval_scaled_geometric},
\path{Fabius.completeHomogeneousEvalOn_range_pow_eq_gaussianBinomial},
\path{Fabius.completeHomogeneousEvalOn_range_pow_eq_gaussianBinomial_degree},
\path{Fabius.gaussianBinomial_add_symm}, and
\path{Fabius.gaussianBinomial_symm_via_completeHomogeneous}.  They prove both
principal-specialization orientations, common scaling, and both generic
Gaussian symmetry laws over every commutative semiring, without division,
distinctness, or cancellation.

\item The five theorems of
\path{GeometricLagrangeCompleteHomogeneous.lean} are
\path{Fabius.completeHomogeneousEvalOn_geometric_range},
\path{Fabius.sum_geometricLagrangeWeight_mul_pow_succ_add_eq_gaussianBinomial},
\path{Fabius.geometricLagrangeQMoment_eq_residual_gaussianBinomial},
\path{Fabius.completeHomogeneousEvalOn_geometric_range_eq_qBinomial}, and
\path{Fabius.geometricLagrangeQMoment_eq_residual_qBinomial_via_completeHomogeneous}.
The first is semiring-level; the field statement uses finite-node injectivity;
the rational quotient bridges retain their explicit nonzero-Pochhammer or
\(0<q<1\) assumptions.

\item \path{GeometricLagrangeQMoments.lean} has the one definition
\path{Fabius.geometricLagrangeQMoment} and exactly 37 theorems:
\path{Fabius.geometricLagrangeQMoment_eq_weightPolynomial_eval},
\path{Fabius.geometricLagrangeQMoment_eq_forwardRichardson_eval},
\path{Fabius.geometricRootPolynomial_inv_eval_pow_mul_signedPowers},
\path{Fabius.geometricRootPolynomial_inv_eval_pow_mul_triangular},
\path{Fabius.geometricRootPolynomial_inv_eval_one_mul_triangular},
\path{Fabius.geometricLagrangeQMoment_eq_qPochhammer},
\path{Fabius.geometricLagrangeQMoment_zero},
\path{Fabius.geometricLagrangeQMoment_eq_zero},
\path{Fabius.geometricRootPolynomial_inv_eval_pow_eq_qPochhammer_of_le},
\path{Fabius.geometricLagrangeQMoment_eq_residual_qPochhammer},
\path{Fabius.qPochhammer_self_add},
\path{Fabius.qPochhammer_self_pos_of_pos_of_lt_one},
\path{Fabius.qBinomial_pos_of_pos_of_lt_one},
\path{Fabius.gaussianBinomial_eq_qBinomial_of_pos_of_lt_one},
\path{Fabius.qPochhammer_pow_pos_of_pos_of_lt_one},
\path{Fabius.qPochhammer_tail_div_self_eq_qBinomial},
\path{Fabius.geometricLagrangeQMoment_eq_residual_qBinomial},
\path{Fabius.geometricLagrangeQMoment_firstUncancelled},
\path{Fabius.negOnePow_mul_geometricLagrangeQMoment_eq_positiveResidual},
\path{Fabius.negOnePow_mul_geometricLagrangeQMoment_pos},
\path{Fabius.qPochhammer_self_succ},
\path{Fabius.qBinomial_succ_succ_of_pos_of_lt_one'},
\path{Fabius.qBinomial_succ_succ_of_pos_of_lt_one},
\path{Fabius.qBinomial_theorem_of_pos_of_lt_one},
\path{Fabius.sum_qBinomial_triangular_succ_eq_neg_qPochhammer},
\path{Fabius.abs_geometricLagrangeWeight_eq_qBinomial},
\path{Fabius.abs_geometricLagrangeWeight_eq_sign_mul},
\path{Fabius.abs_geometricLagrangeWeight_complement_eq_qBinomial},
\path{Fabius.sum_abs_geometricLagrangeWeight_eq_qPochhammer_ratio},
\path{Fabius.neg_qPochhammer_div_self_eq_prod},
\path{Fabius.sum_abs_geometricLagrangeWeight_eq_prod},
\path{Fabius.quarterGeometricLagrangeQMoment_eq_qPochhammer},
\path{Fabius.quarterGeometricLagrangeQMoment_eq_zero},
\path{Fabius.quarterGeometricLagrangeQMoment_eq_residual_qPochhammer},
\path{Fabius.quarterGeometricLagrangeQMoment_eq_residual_qBinomial},
\path{Fabius.quarterGeometricLagrangeQMoment_firstUncancelled}, and
\path{Fabius.sum_abs_quarterGeometricLagrangeWeight_eq_qPochhammer_ratio}.
These are finite rational identities.  Quotient statements retain their
nonzero denominators, and positivity, sign, and absolute-value formulas retain
\(0<q<1\); the module makes no analytic convergence or remainder claim.

\item \path{GeometricRichardsonGenerating.lean} has exactly three public
definitions,
\path{Fabius.geometricRichardsonKernel},
\path{Fabius.qPochhammerNormalizedDataSeries}, and
\path{Fabius.geometricRichardsonTransform}, and seven public theorems:
\path{Fabius.coeff_rescale_qPochhammerSeries_eq_geometricRichardsonKernel},
\path{Fabius.coeff_qPochhammerNormalizedDataSeries},
\path{Fabius.geometricRichardsonTransform_generating},
\path{Fabius.geometricRichardsonTransform_eq_sum_lagrange},
\path{Fabius.geometricLagrangeRichardson_generating},
\path{Fabius.hasSum_geometricRichardsonTransform_mul_pow}, and
\path{Fabius.hasSum_geometricLagrangeRichardson_mul_pow}.
The theorem \path{Fabius.geometricLagrangeRichardson_generating} is the exact
formal counterpart of the canonical comb-manuscript label
\path{gq:thm:richardson-generating}.  The underlying convolution factorization
is formal over an arbitrary commutative ring with totalized inverses; its
Lagrange-row identification assumes a field and nonzero nome.  The two
analytic \path{HasSum} theorems additionally assume a complete normed field,
strict nome contraction, and absolute summability of the normalized data
series.  They do not provide convergence without those hypotheses or an
analytic acceleration bound.

\item \path{GeometricUniformMomentPolynomial.lean} has the one public
definition \path{Fabius.geometricUniformMomentPolynomial} and exactly eight
public theorems:
\path{Fabius.geometricUniformMomentPolynomial_zero},
\path{Fabius.geometricUniformMomentPolynomial_succ},
\path{Fabius.geometricUniformMomentPolynomial_natDegree_le},
\path{Fabius.geometricUniformMomentPolynomial_eval_zero},
\path{Fabius.geometricUniformMomentPolynomial_one},
\path{Fabius.geometricUniformMomentPolynomial_two},
\path{Fabius.geometricUniformMomentPolynomial_three}, and
\path{Fabius.geometricUniformMomentPolynomial_four}.  These declarations
define the recursive family directly in \(\RationalNumbers[X]\), prove its
division-free residual-product recurrence, triangular degree bound and value
at zero, and calculate \(P_0,\ldots,P_4\).  The companion
\path{GeometricUniformMomentPolynomialDegree.lean} has no public definitions
and exactly three public theorems:
\path{Fabius.coeff_geometricUniformMomentPolynomial_choose_two},
\path{Fabius.coeff_geometricUniformMomentPolynomial_choose_two_sub_one}, and
\path{Fabius.geometricUniformMomentPolynomial_natDegree_eq}.  They give the
top Bernoulli coefficient for every \(n\), the displayed subleading coefficient
for \(n\ge2\), and exact parity-sensitive degree, including \(n=0\) and
\(n=1\).  Thus \nolinkurl{prop:qF-P-degree-sharp} and the represented
\nolinkurl{p7:thm:Pn} are Exact.  The exhaustive public surface of the
companion \path{GeometricUniformMomentPolynomialBridge.lean} is no
definitions and the single theorem
\path{Fabius.geometricUniformMomentPolynomial_eval₂_eq_mgf_taylorCoefficient}.
For every real \(\lvert q\rvert<1\) and every \(n\), that theorem identifies
the polynomial evaluation with the finite-q-Pochhammer normalization of the
\(n\)-th Taylor coefficient of the genuine geometric-uniform MGF.  Therefore
\nolinkurl{p7:eq:Pn-def} is Exact in this real contraction regime.  The
exhaustive public surface of
\path{GeometricUniformComplexMomentProduct.lean} is the definition
\path{Fabius.geometricUniformComplexMomentProduct} and the two theorems
\path{Fabius.hasProdLocallyUniformly_geometricUniformComplexMomentProduct} and
\path{Fabius.geometricUniformMomentPolynomial_eval₂_eq_complexMomentProduct_taylorCoefficient}.
They construct the actual locally uniform product and its normalized
Taylor-coefficient bridge for every complex \(\lVert q\rVert<1\).  This is an
analytic analogue rather than a complex probability-moment identity.  The
exhaustive public surface of
\path{GeometricUniformExteriorComplexMomentGerm.lean} is the definition
\path{Fabius.geometricUniformExteriorComplexMomentGerm} and the two theorems
\path{Fabius.analyticAt_geometricUniformExteriorComplexMomentGerm} and
\path{Fabius.geometricUniformMomentPolynomial_eval₂_eq_exteriorComplexMomentGerm_taylorCoefficient}.
They construct the actual reciprocal germ, prove its analyticity at zero, and
give its normalized Taylor-coefficient bridge for every complex
\(1<\lVert q\rVert\).  The final companion
\path{GeometricUniformMomentRatFunc.lean} has the one public definition
\path{Fabius.geometricUniformMomentRatFunc} and exactly four public theorems:
\path{Fabius.qFactorial_mul_geometricUniformMomentRatFunc},
\path{Fabius.eval_geometricUniformMomentRatFunc_eq_complexMomentProduct_taylorCoefficient},
\path{Fabius.eval_geometricUniformMomentRatFunc_eq_exteriorComplexMomentGerm_taylorCoefficient},
and \path{Fabius.eval_geometricUniformMomentRatFunc_one}.  These package the
common rational coefficient, prove its global q-factorial clearing, identify
its safe inner and exterior evaluations, and supply the removable \(q=1\)
specialization.  Accordingly \nolinkurl{thm:qF-moment-polynomial} is Exact by
composition.  Neither this algebraic continuation nor either analytic leaf
asserts an analytic value at a genuine unit-root pole or a unit-circle
continuation theorem.

\item The five theorems of \path{FinitePolynomialFilterExactness.lean} are
\path{Fabius.polynomialFilter_response_eq},
\path{Fabius.polynomialFilter_exact},
\path{Fabius.normalizedGeometricRootPolynomial_filter_exact},
\path{Fabius.forwardGeometricRichardsonPolynomial_filter_exact}, and
\path{Fabius.forwardGeometricRichardsonPolynomial_filter_firstUncancelled}.
The first two are arbitrary commutative-semiring response and exactness laws;
the field-valued geometric forms retain their nonzero-base and normalization-
denominator hypotheses and expose the first surviving mode.

\item \path{QBinomialCauchy.lean} has one public definition,
\path{Fabius.qBernsteinBasis}, and exactly five public theorems:
\path{Fabius.finite_qCauchy_identity},
\path{Fabius.finiteQPochhammerIn_mul_eq_sum_gaussianBinomial},
\path{Fabius.finite_qCauchy_identity_reflected},
\path{Fabius.sum_qBernsteinBasis}, and
\path{Fabius.finite_qCauchy_second_identity}.  The second name is the compatibility
spelling of the primary finite $q$-Cauchy identity.  The reflected convolution, the
denominator-free $q$-Bernstein partition of unity, and the second Cauchy identity hold
over every commutative ring, including degree zero, $q=0$, roots of unity, positive
characteristic, and zero divisors, with no quotient or nonvanishing hypothesis.

\item \path{QuarterCatalanGerm.lean} has the two definitions
\path{Fabius.quarterCatalanCoefficient} and
\path{Fabius.quarterCatalanGermSeries}, and the thirteen theorems
\path{Fabius.quarterCatalanCoefficient_zero},
\path{Fabius.quarterCatalanCoefficient_succ_eq_report},
\path{Fabius.quarterCatalanGermSeries_coeff},
\path{Fabius.quarterCatalanGermSeries_coeff_succ},
\path{Fabius.quarterCatalanGermSeries_constantCoeff},
\path{Fabius.quarterCatalanGermSeries_equation},
\path{Fabius.powerSeries_quadratic_injectiveOn_zeroConstant},
\path{Fabius.eq_quarterCatalanGermSeries_of_equation}, and
\path{Fabius.existsUnique_quarterCatalanGermSeries},
\path{Fabius.dyadicGermTwo_functionalEquation},
\path{Fabius.rescale_dyadicGermTwo_eq_quadraticInverse},
\path{Fabius.dyadicGermTwo_eq_rescale_quadraticInverse}, and
\path{Fabius.coeff_dyadicGermTwo_succ}.  The explicit series is the unique
zero-constant solution of \(D+4D^2=(4/9)X\) in \(\RationalNumbers[[X]]\).  Rescaling the
distinguished rational dyadic germ by \(9/4\) identifies it with the Catalan inverse of
\(X+4X^2\), and its positive coefficients are \((4/9)^{m+1}(-4)^m C_m\).
This module is formal power-series algebra and asserts no convergence or equality of real
functions on a neighborhood.

\item \path{FabiusInverseQuarterJet.lean} has zero public definitions and exactly two
theorems:
\path{Fabius.iteratedDeriv_centeredFabiusInv_quarter_eq_quadraticInverse} and
\path{Fabius.iteratedDeriv_fabiusInv_five_seventy_two_succ}.  For every bounded Fabius
solution, the complete centered inverse jet at \(5/72=F(1/4)\) is the factorial-scaled
coefficient sequence of \path{QuadraticInverse.inverse} at parameter \(4\); equivalently,
\(G^{(m+1)}(5/72)=(m+1)!(-4)^m C_m\).  The conclusion is equality of all jets, not
local analytic equality and not a named nonzero flat-remainder theorem.

\item \path{QuarterCatalanRichardson.lean} has the three definitions
\path{Fabius.finiteRescaleFilter},
\path{Fabius.geometricRichardsonPowerSeriesFilter}, and
\path{Fabius.quarterCatalanRichardsonFilter}, and the 15 theorems
\path{Fabius.finiteRescaleFilter_coeff},
\path{Fabius.geometricRichardsonPowerSeriesFilter_coeff},
\path{Fabius.geometricRichardsonPowerSeriesFilter_coeff_zero},
\path{Fabius.geometricRichardsonPowerSeriesFilter_coeff_eq_zero},
\path{Fabius.geometricRichardsonPowerSeriesFilter_coeff_eq_qPochhammer},
\path{Fabius.geometricRichardsonPowerSeriesFilter_coeff_eq_qBinomial},
\path{Fabius.geometricRichardsonPowerSeriesFilter_firstUncancelled_coeff_of_nonzero},
\path{Fabius.geometricRichardsonPowerSeriesFilter_firstUncancelled_coeff},
\path{Fabius.quarterCatalanRichardsonFilter_coeff},
\path{Fabius.quarterCatalanRichardsonFilter_coeff_zero},
\path{Fabius.quarterCatalanRichardsonFilter_coeff_eq_zero},
\path{Fabius.quarterCatalanRichardsonFilter_coeff_eq_zero_of_le},
\path{Fabius.quarterCatalanRichardsonFilter_coeff_eq_qBinomial},
\path{Fabius.quarterCatalanRichardsonFilter_coeff_succ_eq_qBinomial}, and
\path{Fabius.quarterCatalanRichardsonFilter_firstUncancelled_coeff}.  These are
coefficientwise formal-power-series identities only: no convergence, real
error sign, remainder estimate, or analytic acceleration is asserted.
\end{itemize}
The additional theorem
\path{Fabius.sum_lagrangeEvalWeight_mul_pow_card_add_zero} in
\path{LagrangeResidualMoments.lean} gives every higher evaluation-at-zero
moment of a nonempty distinct field-valued node family as the negative signed
nodal product times its complete homogeneous function.  Its nonempty premise
is exactly what removes the exceptional empty-row \(0^0\) term.
\end{remark}

The $q$-symbols in the dyadic formula therefore have a concrete and limited task: they are the
closed form of a divided-difference stencil on a geometric grid.  No infinite
basic-hypergeometric summation is involved.

\subsection{The parametric coefficient polynomial}

Introduce an auxiliary node variable $z$ and define
\begin{equation}
 \mathscr R_m(z;X)
 =\frac{\EulerE^{-X}B_m(zX)\mathscr A(zX)}{\mathscr A(-X)},
 \qquad
 r_{m,n}(z)=[X^n]\mathscr R_m(z;X).                        \label{dyadicweb:eq:Rmn-def}
\end{equation}
Because $\mathscr A(-X)$ is a unit with constant term $1$, this is a well-defined formal
series and an entire analytic germ.  More importantly, $r_{m,n}$ is a polynomial in $z$.

\begin{lemma}[Node polynomial and its leading coefficient]
\label{dyadicweb:lem:Rmn-leading}
For fixed $m,n$, the polynomial $r_{m,n}(z)$ has degree at most $n$ and
\begin{equation}
 [z^n]r_{m,n}(z)=[X^n]\bigl(B_m(X)\mathscr A(X)\bigr)
 =2^{\binom n2}\FabiusGlobal(m/2^n).                         \label{dyadicweb:eq:Rmn-leading}
\end{equation}
\end{lemma}

\begin{proof}
Write
\(
 \EulerE^{-X}/\mathscr A(-X)=\sum_{j\ge0}u_jX^j
\)
with $u_0=1$.  Since
$B_m(zX)\mathscr A(zX)=(B_m\mathscr A)(zX)$, the coefficient of $X^n$ is a polynomial in $z$
of degree at most $n$.  Its $z^n$ term can only use the constant $u_0$, so its leading
coefficient is $[X^n](B_m\mathscr A)$.  Apply \cref{dyadicweb:thm:master-coefficient}.
\end{proof}

The exact dyadic value is now the leading coefficient of a degree-$n$ polynomial sampled at
$-1,-2,\ldots,-2^n$.  Applying \cref{dyadicweb:prop:q-residue-extractor} gives a compact alternative
to the displayed $q$-sum.

\begin{corollary}[Geometric-node contour representation]
\label{dyadicweb:cor:q-node-contour}
For every $m,n\in\NaturalNumbers$,
\begin{equation}
 \boxed{
 \FabiusGlobal(m/2^n)=
 \frac{2^{-\binom n2}}{2\pi\ImaginaryUnit}
 \int_\Gamma
 \frac{r_{m,n}(z)}{\prod_{k=0}^{n}(z+2^k)}\Differential z,}          \label{dyadicweb:eq:q-node-contour}
\end{equation}
where $\Gamma$ encloses the geometric nodes.
\end{corollary}

Together, \eqref{dyadicweb:eq:cauchy-X} and \eqref{dyadicweb:eq:q-node-contour} exhibit two independent Cauchy
extractions:
\begin{center}
\begin{tabular}{c@{\qquad}c}
coefficient in the analytic variable $X$
& leading coefficient in the geometric node $z$ \\
$[X^n](B_m\mathscr A)$
& $[z^n]r_{m,n}(z)$ \\
ordinary Cauchy kernel $X^{-n-1}$
& barycentric kernel $\Pi_n(z)^{-1}$.
\end{tabular}
\end{center}
The finite $q$-binomial formula is what results when the second contour is evaluated by
residues and each node value is converted to a Thue--Morse block sum.

\subsection{Why the node values are Thue--Morse power sums}

For a complete binary block set
\begin{equation}
 C_k(X)=\sum_{r=0}^{2^k-1}\ThueMorseSignSymbol_r\EulerE^{(r-2^k)X}.              \label{dyadicweb:eq:Ck-local}
\end{equation}
The digit factorization
\(
 \sum_{r<2^k}\ThueMorseSignSymbol_rz^r=\prod_{j=0}^{k-1}(1-z^{2^j})
\)
gives
\begin{equation}
 C_k(X)=(-1)^k2^{\binom k2}X^k
 \EulerE^{-X}\prod_{j=0}^{k-1}E_0(-2^jX).                     \label{dyadicweb:eq:Ck-factor-local}
\end{equation}
Hence $C_k$ has a zero of exact order $k$.  The refinement
\eqref{dyadicweb:eq:A-refinement}, iterated $k$ times, gives
\begin{equation}
 \frac{\EulerE^{-X}\mathscr A(-2^kX)}{\mathscr A(-X)}
 =\EulerE^{-X}\prod_{j=0}^{k-1}E_0(-2^jX).                    \label{dyadicweb:eq:A-geometric-bridge}
\end{equation}
Combining \eqref{dyadicweb:eq:Ck-factor-local} and \eqref{dyadicweb:eq:A-geometric-bridge} identifies the node
specialization of \eqref{dyadicweb:eq:Rmn-def} with a normalized complete block.

There is also a useful differential-operator reading of this cancellation.  For a polynomial
or entire function $P$, set
\begin{equation}
 (\Toperator_kP)(c)=\sum_{r=0}^{2^k-1}\ThueMorseSignSymbol_rP(c+r).         \label{dyadicq:eq:TM-operator}
\end{equation}
Writing $D=\Differential/\Differential c$, the binary digits of $r$ give
\begin{equation}
 \Toperator_k=\prod_{j=0}^{k-1}\left(1-e^{2^jD}\right). \label{dyadicq:eq:TM-factorization}
\end{equation}
Every factor starts with $-2^jD$, so degrees below $k$ vanish and the first surviving
moment is
\begin{equation}
 \sum_{r=0}^{2^k-1}\ThueMorseSignSymbol_rr^k=(-1)^k2^{\binom k2}k!.       \label{dyadicq:eq:first-TM-moment}
\end{equation}
More generally,
\begin{equation}
 c\longmapsto\frac1{(n+k)!}\sum_{r<2^k}\ThueMorseSignSymbol_r(r+c)^{n+k} \label{dyadicq:eq:TM-degree-n}
\end{equation}
has degree at most $n$, with leading coefficient
\begin{equation}
 [c^n]\left(\frac1{(n+k)!}\sum_{r<2^k}\ThueMorseSignSymbol_r(r+c)^{n+k}\right)
 =\frac{(-1)^k2^{\binom k2}}{n!}.                         \label{dyadicq:eq:TM-leading-coeff}
\end{equation}
With $\operatorname{sinhc}(z)=\sinh(z)/z$, its midpoint-symmetric factorization is
\begin{equation}
 \Toperator_k=(-1)^k2^{\binom k2}D^k e^{(2^k-1)D/2}
 \prod_{j=0}^{k-1}\operatorname{sinhc}(2^{j-1}D).        \label{dyadicq:eq:TM-sinhc}
\end{equation}
This exposes a $k$th derivative followed by an even differential correction about the
block midpoint, explaining why centered and translated forms arise naturally.

For a general numerator, the long block factors as
\begin{equation}
 \sum_{r=0}^{m2^k-1}\ThueMorseSignSymbol_r\EulerE^{(r-m2^k)X}
 =B_m(-2^kX)C_k(X).                                       \label{dyadicweb:eq:long-block-local}
\end{equation}
After division by the forced factor
$(-1)^k2^{\binom k2}X^k$, the coefficient of $X^n$ is
\begin{equation}
 \frac{(-1)^k2^{-\binom k2}}{(n+k)!}
 \sum_{r=0}^{m2^k-1}\ThueMorseSignSymbol_r(r-m2^k)^{n+k}.                 \label{dyadicweb:eq:node-value-power-sum}
\end{equation}
Multiplication by the interpolation weight $w_{n,k}$ cancels the sign and contributes a
second factor $2^{-\binom k2}$.  Their product is the characteristic
$4^{-\binom k2}$ of \eqref{dyadicweb:eq:input-q-dyadic}.  Finally
\eqref{dyadicweb:eq:q-barycentric-normalization} converts $\mathfrak m_n$ into
$2^{n^2}(1/2;1/2)_n$ after the factor $2^{\binom n2}$ from
\cref{dyadicweb:lem:Rmn-leading} is included.  This recovers the exact $q$-formula with $Q=0$.

\subsection{Translation invariance is lower-degree annihilation}

The shift $Q$ enters a block generating function by multiplication with $\EulerE^{QX}$.  If
\(
 P_{m,k}(X)
\)
denotes the normalized long-block series after its forced $X^k$ has been removed, then the
$q$-numerator is a weighted sum of $[X^n](\EulerE^{QX}P_{m,k}(X))$.  Expanding the exponential,
\begin{equation}
 [X^n](\EulerE^{QX}P_{m,k}(X))
 =\sum_{s=0}^{n}\frac{Q^{n-s}}{(n-s)!}[X^s]P_{m,k}(X).     \label{dyadicweb:eq:Q-expansion}
\end{equation}
For each $s<n$, the Gaussian row sees a polynomial of geometric-node degree at most $s$ and
annihilates it.  Only the $s=n$ term survives; it has no $Q$ factor.

\begin{proposition}[Finite translation invariance; Lean formalized]
\label{dyadicweb:prop:finite-Q-invariance}
For fixed $m,n$, the polynomial
\begin{equation}
 Q\longmapsto
 \sum_{k=0}^{n}
 \frac{\GaussianBinomial nk{1/2}}{4^{\binom k2}(n+k)!}
 \sum_{r=0}^{m2^k-1}\ThueMorseSignSymbol_r(r-m2^k+Q)^{n+k}                 \label{dyadicweb:eq:Q-polynomial}
\end{equation}
is constant.
\end{proposition}

The constant-polynomial identity and its coefficientwise form are the exact declarations
\path{qBinomialThueMorseDyadicTranslatedFormulaPolynomial_eq_const} and
\path{qBinomialThueMorseDyadicTranslatedFormulaPolynomial_coeff} in
\path{FabiusDyadicQBinomialScalar.lean}.

This is stronger than saying that a limit is independent of $Q$: every coefficient of
$Q,Q^2,\ldots$ vanishes in the finite expression itself.  The distinction becomes important
in \cref{dyadicweb:sec:discrete}, where the finite discrete rows generally do depend on $Q$.

\subsection{A useful abstract summary}

The exact $q$-formula is a composition of two annihilators:
\begin{enumerate}
\item the complete Thue--Morse block has a zero of order $k$, removing all inner polynomial
      degrees below $k$;
\item the Gaussian row on $-2^k$ removes all remaining node degrees below $n$ and returns the
      leading degree $n$.
\end{enumerate}
The exponent $n+k$ in the inner power is therefore forced: $k$ degrees are consumed by the
binary block, leaving exactly $n$ degrees for the geometric extractor.  The formula is a
hybrid of automatic-sequence cancellation, ordinary exponential generating functions, and
finite $q$-calculus; it is not a terminating basic-hypergeometric summation in the usual
${}_r\phi_s$ sense \cite{dyadicweb:gasper2004,dyadicweb:dlmf17}.

\section{Binary Taylor reduction and the global series}
\label{dyadicweb:sec:binary}

\subsection{The finite reduction polynomial}

The inverse-dyadic moments assemble into the polynomial
\begin{equation}
 \mathcal T_m(y)=
 \sum_{j=0}^{m}
 2^{\binom{j+1}{2}-\binom{m-j}{2}}
 \frac{d_{m-j}}{(m-j)!}\frac{y^j}{j!}.                    \label{dyadicweb:eq:T-def-local}
\end{equation}
The signed global extension satisfies the exact Taylor-block identity
\begin{equation}
 \FabiusGlobal(2^{-m}+y)=-\FabiusGlobal(y)+\mathcal T_m(y),
 \qquad 0\le y<2^{-m}.                                    \label{dyadicweb:eq:T-reduction-local}
\end{equation}
This is the elementary step behind both the terminating dyadic recursion and the infinite
binary telescope.

\subsection{Appell compression of every Taylor block}

To avoid collision with the later saddle coefficients, write
\begin{equation}
 \operatorname{App}_m(\theta)=\Expectation(Y+\theta)^m
 =\sum_{j=0}^{m}\binom mjd_{m-j}\theta^j.
 \label{dyadicweb:eq:Appell-definition}
\end{equation}
Its exponential generating function is
$\sum_m\operatorname{App}_m(\theta)X^m/m!=e^{\theta X}\mathscr A(X)$;
hence it is an Appell family,
\begin{equation}
 \operatorname{App}_m'(\theta)=m\operatorname{App}_{m-1}(\theta),\qquad
 \operatorname{App}_m(\theta)=(-1)^m\operatorname{App}_m(-1-\theta).
 \label{dyadicweb:eq:Appell-identities}
\end{equation}
The full Appell addition law is
\begin{equation}
 \operatorname{App}_m(\theta+c)
 =\sum_{j=0}^{m}\binom mj\operatorname{App}_{m-j}(\theta)c^j.
 \label{dyadicq:eq:Appell-addition}
\end{equation}
The irregular powers of two in \eqref{dyadicweb:eq:T-def-local} collapse under the natural
tail scale:
\begin{equation}
 \boxed{
 \mathcal T_m(2^{-m}\theta)
 =2^{-\binom m2}\frac{\operatorname{App}_m(\theta)}{m!}.}
 \label{dyadicweb:eq:Appell-compression}
\end{equation}
Equivalently, at an arbitrary scale,
\begin{equation}
 \mathcal T_m(y)=2^{-\binom m2}
 \frac{\operatorname{App}_m(2^my)}{m!}.                  \label{dyadicq:eq:R-Appell-general}
\end{equation}
Indeed, after substituting $y=2^{-m}\theta$, the exponent in every $j$th term is
$\binom{j+1}{2}-\binom{m-j}{2}-mj=-\binom m2$.
For $0\le\theta<1$, support and positivity give
$0\le\operatorname{App}_m(\theta)\le2^m$.

For orientation, the first rows are
\begin{align*}
 \operatorname{App}_0(\theta)&=1,\\
 \operatorname{App}_1(\theta)&=\theta+\frac12,\\
 \operatorname{App}_2(\theta)&=\theta^2+\theta+\frac5{18},\\
 \operatorname{App}_3(\theta)&=\theta^3+\frac32\theta^2+\frac56\theta+\frac16,\\
 \operatorname{App}_4(\theta)&=\theta^4+2\theta^3+\frac53\theta^2
              +\frac23\theta+\frac{143}{1350}.
\end{align*}

For $x\ge0$ set
\begin{equation}
 p_m(x)=\Floor{2^mx},
 \qquad
 y_m(x)=x-\frac{p_m(x)}{2^m},
 \qquad 0\le y_m(x)<2^{-m},                               \label{dyadicweb:eq:prefix-tail-local}
\end{equation}
and put
\begin{equation}
 p^*_{m-1}(x)=
 \begin{cases}
 \Floor{x/2},&m=0,\\
 p_{m-1}(x),&m\ge1.
 \end{cases}                                               \label{dyadicweb:eq:previous-prefix-local}
\end{equation}
Define
\begin{equation}
 \alpha_m(x)=(-1)^{\BinaryDigitSum(p_m(x))}
 \bigl(2p^*_{m-1}(x)-p_m(x)\bigr).                         \label{dyadicweb:eq:alpha-local}
\end{equation}
This coefficient is $0$ unless the newly exposed binary digit is $1$, in which case it is a
sign.

\begin{theorem}[Binary reduction series]
\label{dyadicweb:thm:binary-series-local}
For every $x\ge0$,
\begin{equation}
 \boxed{
 \FabiusGlobal(x)=\sum_{m=0}^{\infty}
 \alpha_m(x)\mathcal T_m(y_m(x)).}                        \label{dyadicweb:eq:binary-series-local}
\end{equation}
The series converges absolutely, and its partial sum through $m=N$ has the uniform remainder
bound
\begin{equation}
 \left|\FabiusGlobal(x)-\sum_{m=0}^{N}
 \alpha_m(x)\mathcal T_m(y_m(x))\right|\le2^{1-N}
 \qquad(N\ge1).                                           \label{dyadicweb:eq:binary-tail-local}
\end{equation}
\end{theorem}

The proof is an exact residual telescope.  If
\(
 R_m(x)=(-1)^{\BinaryDigitSum(p_m(x))}\FabiusGlobal(y_m(x)),
\)
then
\[
 \alpha_m(x)\mathcal T_m(y_m(x))=R_{m-1}(x)-R_m(x)
 \quad(m\ge1),
\]
with a separately correct scale-zero term.  Since $0\le y_m<2^{-m}$ and
$|\FabiusGlobal(y)|\le2y$ near the origin, $|R_m|\le2^{1-m}$.

\subsection{The sparse digit form on the unit interval}

For $0\le x<1$, write the floor-selected binary expansion
\begin{equation}
 x=0.\delta_1\delta_2\delta_3\cdots
 =\sum_{m\ge1}\delta_m2^{-m},
 \qquad \delta_m\in\{0,1\},                               \label{dyadicweb:eq:binary-digits}
\end{equation}
where $p_m$ is the integer represented by the first $m$ digits.  Let
\(
 s_m=\delta_1+\cdots+\delta_m
\)
and let
\begin{equation}
 t_m=\sum_{\ell>m}\delta_\ell2^{-\ell}=y_m(x)             \label{dyadicweb:eq:binary-tail-digits}
\end{equation}
be the unscaled tail.  If $\delta_m=1$, then
$p_m=2p_{m-1}+1$ and
\[
 \alpha_m=(-1)^{s_m}(-1)=(-1)^{s_{m-1}};
\]
if $\delta_m=0$, then $\alpha_m=0$.

\begin{corollary}[Digit-sparse series]
\label{dyadicweb:cor:digit-sparse}
For $0\le x<1$,
\begin{equation}
 \boxed{
 F(x)=\sum_{\substack{m\ge1\\\delta_m=1}}
 (-1)^{s_{m-1}}\mathcal T_m(t_m).}                         \label{dyadicweb:eq:digit-sparse}
\end{equation}
At a dyadic rational the floor convention chooses the terminating binary expansion, so the
sum is finite and contains exactly one term for each set bit of the numerator.
\end{corollary}

If $\theta_m=2^mt_m=\fract{2^mx}$, the Appell compression gives the equivalent
digit-sparse form
\begin{equation}
 F(x)=\sum_{\substack{m\ge1\\\delta_m=1}}
 (-1)^{s_{m-1}}
 \frac{2^{-\binom m2}}{m!}\operatorname{App}_m(\theta_m).
 \label{dyadicweb:eq:digit-Appell}
\end{equation}
It converges uniformly on $0\le x<1$.  If $S_N$ denotes its truncation through
$m=N$, then
\begin{equation}
 \boxed{
 |F(x)-S_N(x)|
 \le2\,\frac{2^{N+1-\binom{N+1}{2}}}{(N+1)!}.}
 \label{dyadicweb:eq:digit-Appell-tail}
\end{equation}
The ratio of consecutive majorants is $2^{1-m}/(m+1)\le1/2$; the digit series
therefore has a superexponential certificate, substantially sharper than the generic
$2^{1-N}$ residual bound needed for the all-real telescope.

This is the most direct connection between the fast bit-recursive evaluator and the arbitrary
argument series: the finite recursion is not a separate algorithm pasted onto the infinite
formula; it is literally the same telescope when the binary tail eventually becomes zero.

\subsection{Direct one-row Appell and nested digit--\texorpdfstring{$q$}{q} forms}

The triangular reconstruction below is not the only way to insert the Gaussian row.  Retaining
the Appell variable throughout the finite identity gives the direct centered formula.
\begin{theorem}[Centered Gaussian-binomial representation of the Appell row]
\label{dyadicq:thm:centered-q-Appell}
For every $n\ge0$ and $\theta\in\ComplexNumbers$,
\begin{equation}
 \boxed{
 \frac{\operatorname{App}_n(\theta)}{n!}
 =\frac{2^{-\binom{n+1}{2}}}{(1/2;1/2)_n}
 \sum_{k=0}^{n}\frac{\GaussianBinomial nk{1/2}}{4^{\binom k2}(n+k)!}
 \sum_{r=0}^{2^k-1}\ThueMorseSignSymbol_r\bigl(r-2^k(1+\theta)\bigr)^{n+k}.}
 \label{dyadicq:eq:centered-q-Appell}
\end{equation}
\end{theorem}
At $\theta=0$, multiplication by $2^{-\binom n2}$ recovers the centered inverse-dyadic
formula.  More generally, a common translation may be inserted in every inner block.
\begin{theorem}[Raw translated Appell representation]
\label{dyadicq:thm:raw-q-Appell}
For every $n\ge0$ and $\theta,\eta\in\ComplexNumbers$,
\begin{equation}
 \boxed{
 \frac{\operatorname{App}_n(\theta)}{n!}
 =\frac{(-1)^n2^{-\binom{n+1}{2}}}{(1/2;1/2)_n}
 \sum_{k=0}^{n}\frac{\GaussianBinomial nk{1/2}}{4^{\binom k2}(n+k)!}
 \sum_{r=0}^{2^k-1}\ThueMorseSignSymbol_r\bigl(r+\eta+2^k\theta\bigr)^{n+k}.}
 \label{dyadicq:eq:raw-q-Appell}
\end{equation}
The right-hand side is independent of $\eta$.
\end{theorem}

For a fixed outer scale $n$, the translation must be common to every $k$-term in that row.
Different outer scales may use different translations $\eta_n$ without changing any term;
this is a genuine translation gauge and can reduce intermediate powers in floating-point
work, but varying the shift within one $k$-row invalidates the cancellation.

Termwise substitution into \eqref{dyadicweb:eq:digit-Appell} gives, for $0\le x<1$,
\begin{align}
 F(x)&=\sum_{\substack{m\ge1\\\delta_m=1}}
 (-1)^{s_{m-1}}
 \frac1{2^{m^2}(1/2;1/2)_m}
 \sum_{k=0}^{m}\frac{\GaussianBinomial mk{1/2}}{4^{\binom k2}(m+k)!}\notag\\[-0.2em]
 &\hspace{5em}\times
 \sum_{r=0}^{2^k-1}\ThueMorseSignSymbol_r\bigl(r-2^k(1+\theta_m)\bigr)^{m+k},
 \label{dyadicq:eq:centered-digit-q}
\end{align}
and, for any fixed $\eta\in\ComplexNumbers$,
\begin{align}
 F(x)&=\sum_{\substack{m\ge1\\\delta_m=1}}
 (-1)^{s_{m-1}}
 \frac1{(-2)^{m^2}(1/2;1/2)_m}
 \sum_{k=0}^{m}\frac{\GaussianBinomial mk{1/2}}{4^{\binom k2}(m+k)!}\notag\\[-0.2em]
 &\hspace{5em}\times
 \sum_{r=0}^{2^k-1}\ThueMorseSignSymbol_r\bigl(r+\eta+2^k\theta_m\bigr)^{m+k}.
 \label{dyadicq:eq:raw-digit-q}
\end{align}
Each $m$th summand equals the corresponding Appell summand before any infinite summation is
interchanged.  At a dyadic input the series terminates, so these are finite exact identities.

\subsection[Replacing each Taylor coefficient by a finite q-identity]{Replacing each Taylor coefficient by a finite $q$-identity}

For $Q\in\ComplexNumbers$ define the finite inverse-block numerator
\begin{equation}
 \mathcal N_n^{(Q)}=
 \sum_{k=0}^{n}
 \frac{\GaussianBinomial nk{1/2}}{4^{\binom k2}(n+k)!}
 \sum_{r=0}^{2^k-1}\ThueMorseSignSymbol_r(r-2^k+Q)^{n+k}.                  \label{dyadicweb:eq:NnQ-local}
\end{equation}
The exact one-block $q$-formula is equivalent to
\begin{equation}
 \frac{\mathcal N_n^{(Q)}}
 {2^{\binom n2}(1/2;1/2)_n}
 =\frac{2^nd_n}{n!},                                      \label{dyadicweb:eq:NnQ-normalization-local}
\end{equation}
which is independent of $Q$.  Insert these finite identities into the polynomial
\eqref{dyadicweb:eq:T-def-local}:
\begin{equation}
 \mathcal Q_m^{(Q)}(y)=
 2^{-\binom{m+1}{2}}
 \sum_{n=0}^{m}
 \frac{(2^{m+1}y)^{m-n}}{(m-n)!}
 \frac{\mathcal N_n^{(Q)}}
 {2^{\binom n2}(1/2;1/2)_n}.                              \label{dyadicweb:eq:QmQ-local}
\end{equation}

\begin{theorem}[$q$-Taylor identity and global $q$-series]
\label{dyadicweb:thm:q-global-local}
For every $m$, $y$, and $Q\in\ComplexNumbers$,
\begin{equation}
 \mathcal Q_m^{(Q)}(y)=\mathcal T_m(y).                   \label{dyadicweb:eq:Qm-equals-Tm-local}
\end{equation}
Consequently, for fixed $Q$ and every $x\ge0$,
\begin{equation}
 \boxed{
 \FabiusGlobal(x)=\sum_{m=0}^{\infty}
 \alpha_m(x)\mathcal Q_m^{(Q)}(y_m(x)),}                 \label{dyadicweb:eq:global-q-series-local}
\end{equation}
with absolute convergence.
\end{theorem}

The key word is \emph{termwise}: every finite polynomial
$\mathcal Q_m^{(Q)}$ already equals $\mathcal T_m$ and is already independent of $Q$ before
the outer infinite sum is considered.  The $q$-symbols do not create the convergence of the
outer series.  They merely provide an alternative finite presentation of each Taylor block.

\subsection{The dyadic triangle of finite representations}

Suppose $x=m/2^s\in[0,1]$ is written with a terminating binary expansion.  Then three different
finite mechanisms evaluate the same rational number:
\begin{enumerate}
\item \eqref{dyadicweb:eq:input-moment-dyadic} sums the Thue--Morse prefix $0\le h<m$ and applies the
      even-moment correction in one shot;
\item \eqref{dyadicweb:eq:input-q-dyadic} uses one geometric row $0\le k\le s$ and two nested
      cancellations;
\item \eqref{dyadicweb:eq:digit-sparse} walks only through the set bits of $m$, using inverse-dyadic
      Taylor data.
\end{enumerate}
The finite sums have no natural term-by-term identification.  Their equality is mediated by
the master coefficient \eqref{dyadicweb:eq:master-coefficient} and by the exact Taylor reduction
\eqref{dyadicweb:eq:T-reduction-local}.  A fourth finite representation, obtained from the discrete
limit, appears in the next section.

The combinatorial bridge between the one-shot and set-bit organizations is an exact
prefix-block decomposition.  Write
$M/2^N=0.\delta_1\ldots\delta_N{}_2$ and
$p_m=\Floor{2^mM/2^N}$.  For any function $P$ on the nonnegative integers and
$k\ge0$,
\begin{align}
 \sum_{r=0}^{M2^k-1}\ThueMorseSignSymbol_rP(r)
 &=\sum_{\substack{1\le m\le N\\\delta_m=1}}
 \ThueMorseSignSymbol_{p_{m-1}}
 \sum_{s=0}^{2^{N-m+k}-1}\ThueMorseSignSymbol_s
 P\!\left(p_{m-1}2^{N-m+k+1}+s\right).
 \label{dyadicweb:eq:prefix-block-decomposition}
\end{align}
The interval $[0,M)$ is the disjoint union of the blocks selected by the set bits.
After scaling by $2^k$, each block begins at a multiple of twice its length, so adding
the local coordinate creates no binary carry and its Thue--Morse sign factors.  This
explains why the one-shot $q$ formula and the sparse Appell formula group the same exact
quantity so differently.

\section[Discrete q-limits, exact stabilization, and extrapolation]{Discrete $q$-limits, exact stabilization, and extrapolation}
\label{dyadicweb:sec:discrete}

\subsection{Shifted splines and the finite row}

For $p\ge1$, $Q\in\ComplexNumbers$, and $x\in\RealNumbers$, define
\begin{equation}
 S_p^{(Q)}(x)=
 \frac{(-1)^p}{2^{\binom p2}p!}
 \sum_{r=0}^{N_p(x)-1}\ThueMorseSignSymbol_r(r-2^px+Q)^p.                  \label{dyadicweb:eq:shifted-spline-local}
\end{equation}
At $Q=1/2$ this is the centered spline $S_p$.  For $x\ge0$, the formal development proves
the exact fixed-cutoff Taylor relation and the audited bound
\begin{equation}
 \bigl|S_p^{(Q)}(x)-S_p(x)\bigr|
 \le2^{-(p-1)}\EulerE^{|Q-1/2|}.                              \label{dyadicweb:eq:shifted-spline-bound-formal}
\end{equation}
This is \path{norm_fabiusComplexShiftSpline_sub_center_le_half_pow_mul_exp} in
\path{FabiusComplexShiftSpline.lean}, with hypotheses $1\le p$ and $0\le x$.  When $x\le0$,
both shifted splines vanish by \path{fabiusComplexShiftSpline_eq_zero_of_nonpos}; an all-real
bound follows by a direct case split, but is not presently packaged as one public theorem.
The sharper estimate used in the frontier calculation below is
\begin{equation}
 \bigl|S_p^{(Q)}(x)-S_p(x)\bigr|
 \le2^{-(p-1)}\bigl(\EulerE^{|Q-1/2|}-1\bigr).                 \label{dyadicweb:eq:shifted-spline-bound-local}
\end{equation}
It is a natural-language candidate, not presently an exact Lean theorem; its formalization
obligation and the $p=2$ edge-case warning are recorded in the Primary Exposition Gap
Register at the end of this volume.  Consequently
\eqref{dyadicweb:eq:shifted-total-error}, \cref{dyadicweb:thm:D-error}, and
\cref{dyadicweb:cor:D-Q-decay} retain frontier status with their printed sharp constants.
Together with \eqref{dyadicweb:eq:centered-global-error-input},
\begin{equation}
 \bigl|S_p^{(Q)}(x)-\FabiusGlobal(x)\bigr|
 \le(2\EulerE^{|Q-1/2|}-1)2^{-p}.                              \label{dyadicweb:eq:shifted-total-error}
\end{equation}

The discrete $q$-approximant is
\begin{align}
 \mathcal D_n^{(Q)}(x)
 &=\frac1{2^{n^2}(1/2;1/2)_n}
 \sum_{k=0}^{n}
 \frac{\GaussianBinomial nk{1/2}}{4^{\binom k2}(n+k)!}             \notag\\
 &\qquad\times
 \sum_{r=0}^{N_{n+k}(x)-1}\ThueMorseSignSymbol_r
       (r-2^{n+k}x+Q)^{n+k}.                              \label{dyadicweb:eq:DnQ-local}
\end{align}
For every fixed $Q$ and every real $x$, the Lean development proves
\begin{equation}
 \mathcal D_n^{(Q)}(x)\longrightarrow\FabiusGlobal(x).       \label{dyadicweb:eq:DnQ-limit-local}
\end{equation}
The notation strongly resembles the exact dyadic formula
\eqref{dyadicweb:eq:input-q-dyadic}.  The following result makes the resemblance literal.

\subsection{At dyadic inputs the limit stabilizes exactly}

\begin{theorem}[Finite stabilization on the dyadic grid]
\label{dyadicweb:thm:dyadic-stabilization}
Let $x=m/2^s$ with $m,s\in\NaturalNumbers$.  Then, for every $n\ge s$ and every $Q\in\ComplexNumbers$,
\begin{equation}
 \boxed{\mathcal D_n^{(Q)}(x)=\FabiusGlobal(x).}              \label{dyadicweb:eq:dyadic-stabilization}
\end{equation}
If $0\le x\le1$, the right side is $F(x)$.
\end{theorem}

\begin{proof}
For $n\ge s$ put $M=m2^{n-s}$, so that $x=M/2^n$.  For every $0\le k\le n$,
\[
 2^{n+k}x=M2^k\in\NaturalNumbers,
 \qquad
 N_{n+k}(x)=\Floor{M2^k+\frac12}=M2^k.
\]
Therefore \eqref{dyadicweb:eq:DnQ-local} becomes exactly the dyadic formula
\eqref{dyadicweb:eq:input-q-dyadic} with numerator $M$ and denominator exponent $n$.  That formula is
valid for unreduced dyadic representations and for every complex translation $Q$.
\end{proof}

This is the sharpest connection between the exact and limiting formulae: the discrete formula
extends the exact dyadic row to arbitrary $x$, and on a dyadic grid it needs no limiting
process once the row depth reaches the denominator depth.

\begin{corollary}[Exact finite shifted-spline extrapolation]
\label{dyadicweb:cor:exact-shifted-extrapolation}
For $x=m/2^s$, $n\ge s$, and every $Q\in\ComplexNumbers$,
\begin{equation}
 \FabiusGlobal(x)=\sum_{j=0}^{n}\omega_{n,j}
 S_{2n-j}^{(Q)}(x),                                       \label{dyadicweb:eq:exact-shifted-extrapolation}
\end{equation}
where the weights $\omega_{n,j}$ are defined in
\eqref{dyadicweb:eq:omega-local} below.
\end{corollary}

Thus every dyadic value admits an exact representation as a finite signed combination of
consecutive spline degrees, all evaluated at the same point.  Each shifted spline may depend
strongly on $Q$, but the complete combination does not.

\subsection{The row polynomial}

Reindexing \eqref{dyadicweb:eq:DnQ-local} by $j=n-k$ gives the exact identity
\begin{equation}
 \mathcal D_n^{(Q)}(x)=
 \sum_{j=0}^{n}\omega_{n,j}S_{2n-j}^{(Q)}(x),              \label{dyadicweb:eq:D-toeplitz-local}
\end{equation}
with
\begin{equation}
 \omega_{n,j}=
 \frac{(-1)^j\GaussianBinomial nj{1/2}2^{-\binom{j+1}{2}}}
 {(1/2;1/2)_n}.                                           \label{dyadicweb:eq:omega-local}
\end{equation}
The full row is encoded by one normalized $q$-Pochhammer polynomial.

\begin{proposition}[Generating polynomial of the Toeplitz row]
\label{dyadicweb:prop:Wn}
Define
\begin{equation}
 W_n(z)=\sum_{j=0}^{n}\omega_{n,j}z^j.                    \label{dyadicweb:eq:Wn-def}
\end{equation}
Then
\begin{equation}
 \boxed{
 W_n(z)=\frac{(z/2;1/2)_n}{(1/2;1/2)_n}
 =\prod_{\ell=1}^{n}\frac{1-z/2^\ell}{1-1/2^\ell}.}      \label{dyadicweb:eq:Wn-product}
\end{equation}
Consequently
\begin{align}
 W_n(1)&=1,                                                \label{dyadicweb:eq:Wn-one}\\
 W_n(2^r)&=0\qquad(1\le r\le n),                         \label{dyadicweb:eq:Wn-roots}\\
 \sum_{j=0}^{n}|\omega_{n,j}|
 &=W_n(-1)=\frac{(-1/2;1/2)_n}{(1/2;1/2)_n},              \label{dyadicweb:eq:Wn-variation}\\
 \sum_{j=0}^{n}|\omega_{n,j}|2^j
 &=W_n(-2)=\frac{(-1;1/2)_n}{(1/2;1/2)_n}.                \label{dyadicweb:eq:Wn-weighted-variation}
\end{align}
\end{proposition}

\begin{proof}
Apply the finite $q$-binomial theorem to $(z/2;1/2)_n$.  Its coefficient of $z^j$ is
\[
 (-1)^j2^{-\binom j2}2^{-j}\GaussianBinomial nj{1/2}
 =(-1)^j2^{-\binom{j+1}{2}}\GaussianBinomial nj{1/2},
\]
which gives \eqref{dyadicweb:eq:Wn-product} after division by $(1/2;1/2)_n$.  The evaluations at
$1,2^r,-1,-2$ follow from the product and from the alternating signs of the coefficients.
\end{proof}

\paragraph{Current formal status.}
For rational $z$, the full identity \eqref{dyadicweb:eq:Wn-product} is
\lean{sum_range_discreteLimitWeight_mul_pow}, and the exact dyadic root locus
\eqref{dyadicweb:eq:Wn-roots} is
\lean{sum_range_discreteLimitWeight_mul_pow_eq_zero_iff}.  The mass and
unweighted variation evaluations are respectively
\lean{sum_range_discreteLimitWeight} and
\lean{sum_abs_discreteLimitWeight}, all in
\path{FabiusDiscreteLimitToeplitz.lean}.  The weighted variation at $-2$,
the logarithmic-derivative moment formulas below, and the limiting filter are
not claimed as named Lean counterparts here.

Logarithmic differentiation at $z=1$ records the row's first signed moment:
\begin{equation}
\sum_{j=0}^{n}j\omega_{n,j}=W_n'(1)
 =-\sum_{\ell=1}^{n}\frac1{2^\ell-1}.
 \label{dyadicweb:eq:row-first-moment}
\end{equation}
Further logarithmic differentiation gives every higher factorial row moment.  The coefficients
also have the finite symmetric form
\begin{equation}
 \omega_{n,j}=\frac{(-1)^j2^{-\binom{j+1}{2}}}
 {(1/2;1/2)_j(1/2;1/2)_{n-j}}.                            \label{dyadicq:eq:weights-symmetric}
\end{equation}
For each fixed $j$ this gives the limit
\begin{equation}
 \omega_{\infty,j}=
 \frac{(-1)^j2^{-\binom{j+1}{2}}}
 {(1/2;1/2)_j(1/2;1/2)_\infty},
 \label{dyadicweb:eq:limit-row-weight}
\end{equation}
with stationary generating filter
\begin{equation}
 W_\infty(z)=\frac{(z/2;1/2)_\infty}{(1/2;1/2)_\infty}.
 \label{dyadicweb:eq:limit-filter}
\end{equation}
It has unit mass and finite total variation.  This limiting filter is a useful diagnostic,
although the finite roots in \eqref{dyadicweb:eq:Wn-roots} are what produce exact
$q$-Richardson cancellation.

The familiar row-sum identity and bounded variation are therefore only two evaluations of a
single polynomial.  The roots $2,4,\ldots,2^n$ contain additional information that the basic
Toeplitz proof does not need.

\subsection[A q-Pochhammer filter in the degree variable]{A $q$-Pochhammer filter in the degree variable}

Fix $x$ and $Q$ and regard $p\mapsto S_p^{(Q)}(x)$ as a sequence.  Let the backward shift
$\shiftop$ act by
\(
 \shiftop S_p^{(Q)}=S_{p-1}^{(Q)}.
\)
Then \eqref{dyadicweb:eq:D-toeplitz-local} and \eqref{dyadicweb:eq:Wn-product} can be written as
\begin{equation}
 \boxed{
 \mathcal D_n^{(Q)}(x)
 =W_n(\shiftop)S_{2n}^{(Q)}(x)
 =\prod_{\ell=1}^{n}
   \frac{1-2^{-\ell}\shiftop}{1-2^{-\ell}}
   S_{2n}^{(Q)}(x).}                                      \label{dyadicweb:eq:q-filter-operator}
\end{equation}
Each factor preserves a constant sequence because its value at $\shiftop=1$ is $1$.
Meanwhile the zero $W_n(2^r)=0$ means that it removes a geometric error mode with ratio
$2^{-r}$ in the degree.

\begin{proposition}[$q$-Richardson exactness]
\label{dyadicweb:prop:q-richardson}
Suppose a rational-valued model sequence has the form
\begin{equation}
 u_p=L+\sum_{r=1}^{n}a_r2^{-rp}.                           \label{dyadicweb:eq:model-errors}
\end{equation}
Then
\begin{equation}
 \sum_{j=0}^{n}\omega_{n,j}u_{2n-j}=L.                    \label{dyadicweb:eq:q-richardson}
\end{equation}
\end{proposition}

\begin{proof}
The constant contributes $L W_n(1)=L$.  The $r$th error contributes
\[
 a_r2^{-2nr}\sum_{j=0}^{n}\omega_{n,j}2^{rj}
 =a_r2^{-2nr}W_n(2^r)=0.
\]
\end{proof}

\paragraph{Current formal status.}
The module-valued theorem
\lean{discreteLimitWeight_qRichardson_exact_module} in
\path{FabiusDiscreteLimitToeplitz.lean} holds for sequences in any additive
commutative monoid carrying a $\RationalNumbers$-module structure.  It needs no target
multiplication, topology, order, or norm: the rational row weights and dyadic
mode coefficients act only by scalar multiplication.  Its single-mode core is
\lean{sum_range_discreteLimitWeight_mul_geometricMode_smul_eq_zero}.
The theorem \lean{discreteLimitWeight_qRichardson_exact} is the $M=\RationalNumbers$
specialization displayed above; its zero-based coefficient $a_r$ multiplies
the mode $2^{-(r+1)p}$.  These theorems transport cancellation at the rational
dyadic nodes; they do not classify zeros at arbitrary scalar arguments.

No claim is needed that the actual spline error has a finite expansion of the form
\eqref{dyadicweb:eq:model-errors}.  The proposition explains the algebraic design of the row: it is a
geometric-grid analogue of Richardson extrapolation, with the normalized $q$-Pochhammer
polynomial as its error-annihilating filter.

\subsection[A quantitative four-to-the-minus-n convergence bound]{A quantitative $4^{-n}$ convergence bound}

The bounded-variation proof of \eqref{dyadicweb:eq:DnQ-limit-local} uses only that all sampled degrees
$2n-j$ tend to infinity.  Retaining their actual sizes gives a stronger quantitative
consequence.

\begin{theorem}[Explicit discrete-row error; frontier sharpening]
\label{dyadicweb:thm:D-error}
For $n\ge1$, $Q\in\ComplexNumbers$, and every $x\in\RealNumbers$, let
\begin{equation}
 H_n=\frac{(-1;1/2)_n}{(1/2;1/2)_n}.                      \label{dyadicweb:eq:Hn-def}
\end{equation}
Then
\begin{equation}
 \boxed{
 |\mathcal D_n^{(Q)}(x)-\FabiusGlobal(x)|
 \le H_n\bigl(2\EulerE^{|Q-1/2|}-1\bigr)4^{-n}.}              \label{dyadicweb:eq:D-error}
\end{equation}
Moreover $H_n\le64$, so the convergence is uniformly $\BigO_Q(4^{-n})$ on the whole real
line.  At the centered shift,
\begin{equation}
 |\mathcal D_n^{(1/2)}(x)-\FabiusGlobal(x)|\le H_n4^{-n}.     \label{dyadicweb:eq:D-centered-error}
\end{equation}
\end{theorem}

\begin{proof}
Apply \eqref{dyadicweb:eq:shifted-total-error} with $p=2n-j$ in
\eqref{dyadicweb:eq:D-toeplitz-local}:
\begin{align*}
 |\mathcal D_n^{(Q)}(x)-\FabiusGlobal(x)|
 &\le\bigl(2\EulerE^{|Q-1/2|}-1\bigr)
 \sum_{j=0}^{n}|\omega_{n,j}|2^{-(2n-j)}\\
 &=\bigl(2\EulerE^{|Q-1/2|}-1\bigr)4^{-n}
 \sum_{j=0}^{n}|\omega_{n,j}|2^j.
\end{align*}
Use \eqref{dyadicweb:eq:Wn-weighted-variation}.  Finally
\[
 H_n
 =\frac{2}{1-2^{-n}}
   \frac{(-1/2;1/2)_{n-1}}{(1/2;1/2)_{n-1}}
 \le4\cdot16=64,
\]
where the last factor is the uniform variation bound proved in the formal development.
\end{proof}

The same argument isolates the finite translation dependence.

\begin{corollary}[Decay of the shift dependence]
\label{dyadicweb:cor:D-Q-decay}
For every $x\in\RealNumbers$ and $n\ge1$,
\begin{equation}
 |\mathcal D_n^{(Q)}(x)-\mathcal D_n^{(1/2)}(x)|
 \le2H_n\bigl(\EulerE^{|Q-1/2|}-1\bigr)4^{-n}.                \label{dyadicweb:eq:D-Q-decay}
\end{equation}
\end{corollary}

Thus the finite rows can depend dramatically on $Q$ at small $n$, but the dependence is
uniformly suppressed at a geometric rate.  This is analytically different from
\cref{dyadicweb:prop:finite-Q-invariance}, where the finite dyadic expression is exactly constant in
$Q$.

\subsection[A synchronized nearest-dyadic q-limit]{A synchronized nearest-dyadic $q$-limit}
\label{dyadicweb:subsec:synchronized}

There is a simpler way to turn the exact dyadic formula into a limit at arbitrary argument.
For $x\ge0$ define the nearest nonnegative dyadic numerator
\begin{equation}
 M_n(x)=\max\!\left(0,\Floor{2^nx+\frac12}\right)=N_n(x),
 \qquad
 \widetilde x_n=\frac{M_n(x)}{2^n}.                        \label{dyadicweb:eq:nearest-dyadic}
\end{equation}
Use this one numerator at every inner scale:
\begin{align}
 \widetilde{\mathcal D}_n^{(Q)}(x)
 &=\frac1{2^{n^2}(1/2;1/2)_n}
 \sum_{k=0}^{n}
 \frac{\GaussianBinomial nk{1/2}}{4^{\binom k2}(n+k)!}             \notag\\
 &\qquad\times
 \sum_{r=0}^{M_n(x)2^k-1}\ThueMorseSignSymbol_r
       (r-M_n(x)2^k+Q)^{n+k}.                             \label{dyadicweb:eq:synchronized-D}
\end{align}
Unlike \eqref{dyadicweb:eq:DnQ-local}, every finite row is now an exact dyadic formula.

\begin{theorem}[Synchronized exact-grid representation]
\label{dyadicweb:thm:synchronized}
For every $n$, $Q\in\ComplexNumbers$, and $x\ge0$,
\begin{equation}
 \boxed{
 \widetilde{\mathcal D}_n^{(Q)}(x)
 =\FabiusGlobal\!\left(\frac{M_n(x)}{2^n}\right).}           \label{dyadicweb:eq:synchronized-exact}
\end{equation}
In particular, the left side is exactly independent of $Q$.  On $0\le x\le1$,
\begin{equation}
 \left|\widetilde{\mathcal D}_n^{(Q)}(x)-F(x)\right|
 \le2^{-n},                                                \label{dyadicweb:eq:synchronized-error}
\end{equation}
and hence
\begin{equation}
 F(x)=\lim_{n\to\infty}\widetilde{\mathcal D}_n^{(Q)}(x). \label{dyadicweb:eq:synchronized-limit}
\end{equation}
\end{theorem}

\begin{proof}
Equation \eqref{dyadicweb:eq:synchronized-exact} is
\eqref{dyadicweb:eq:input-q-dyadic} with numerator $M_n(x)$.  For $x\in[0,1]$,
$|\widetilde x_n-x|\le2^{-n-1}$ and the bounded Fabius function is $2$-Lipschitz, giving
\eqref{dyadicweb:eq:synchronized-error}.
\end{proof}

The same construction applied to \eqref{dyadicweb:eq:input-moment-dyadic} gives a purely
moment/Thue--Morse limit:
\begin{align}
 F(x)=\lim_{n\to\infty}
 &\frac{2^{-\binom{n+1}{2}}}{n!}
 \sum_{h=0}^{M_n(x)-1}\ThueMorseSignSymbol_h                              \notag\\
 &\times\sum_{k=0}^{\Floor{n/2}}
 \binom n{2k}c_k(2M_n(x)-2h-1)^{n-2k},                   \label{dyadicweb:eq:synchronized-moment-limit}
\end{align}
again with error at most $2^{-n}$ on $[0,1]$.  This representation contains no $q$-symbols at
all: it is exact dyadic arithmetic followed by nearest-grid convergence.

The synchronized limit is conceptually simple but converges at the basic grid rate
$\BigO(2^{-n})$.  The asynchronous row \eqref{dyadicweb:eq:DnQ-local} is adapted to the natural spline
degree $n+k$ at each scale and, by \cref{dyadicweb:thm:D-error}, achieves the stronger uniform bound
$\BigO_Q(4^{-n})$.  The normalized $q$-Pochhammer row acts as an acceleration filter.

\subsection{What fails away from a dyadic point}

Let
\begin{equation}
 \theta_n(x)=2^nx-M_n(x),
 \qquad -\frac12\le\theta_n(x)<\frac12.                  \label{dyadicweb:eq:theta-n}
\end{equation}
Then for every $k\ge0$,
\begin{equation}
 N_{n+k}(x)
 =M_n(x)2^k+\Floor{2^k\theta_n(x)+\frac12}, \label{dyadicweb:eq:cutoff-defect}
\end{equation}
and the shifted coordinate in the inner power is
\begin{equation}
 r-2^{n+k}x+Q
 =r-M_n(x)2^k+\bigl(Q-2^k\theta_n(x)\bigr).               \label{dyadicweb:eq:shift-defect}
\end{equation}
Thus the source discrete row differs from the synchronized exact dyadic row in exactly two
ways:
\begin{enumerate}
\item the complete block of length $M_n2^k$ acquires a boundary defect
      $\Floor{2^k\theta_n+1/2}$;
\item the common translation $Q$ is replaced by the scale-dependent translation
      $Q-2^k\theta_n$.
\end{enumerate}
At a resolved dyadic point, $\theta_n=0$, so both defects disappear simultaneously; this is
another proof of \cref{dyadicweb:thm:dyadic-stabilization}.  Away from the dyadic grid the shifts vary
with $k$, so the finite common-translation invariance of
\cref{dyadicweb:prop:finite-Q-invariance} cannot be applied across the row.

\subsection{Two infinite constructions that share a limit but not terms}

The global binary $q$-series \eqref{dyadicweb:eq:global-q-series-local} and the discrete limit
\eqref{dyadicweb:eq:DnQ-limit-local} both evaluate $\FabiusGlobal$, but their relation is not that of an
infinite series and its partial sums.

\begin{table}[ht]
\centering
\small
\begin{tabularx}{\textwidth}{@{}p{0.22\textwidth}XX@{}}
\toprule
Feature & Binary $q$-series & Discrete $q$-limit \\
\midrule
Outer index & binary scale $m$ & approximation row $n$ \\
Finite inner object & $\mathcal Q_m^{(Q)}=\mathcal T_m$ & shifted splines $S_{2n-j}^{(Q)}$ \\
Shift dependence & absent in every summand & generally present in every finite row \\
Convergence source & exact residual telescope & Toeplitz summability of spline convergence \\
At dyadic $x$ & outer series terminates at the last set bit & rows stabilize exactly once $n$ reaches the denominator depth \\
Natural error bound & $2^{1-N}$ after scale $N$ & $H_n(2\EulerE^{|Q-1/2|}-1)4^{-n}$ \\
\bottomrule
\end{tabularx}
\caption{The two arbitrary-argument $q$ constructions.}
\label{dyadicweb:tab:two-infinite}
\end{table}

\begin{warning}
It is therefore incorrect to call $\mathcal D_n^{(Q)}$ the $n$th partial sum of the global
$q$-binomial series.  The two families use the same finite half-base identities, but they
insert them into different analytic frameworks.
\end{warning}

\subsection{The first rows}

For reference, the first normalized Toeplitz rows are
\begin{align*}
 (\omega_{1,j})_{j=0}^{1}&=(2,-1),\\
 (\omega_{2,j})_{j=0}^{2}&=(8/3,-2,1/3),\\
 (\omega_{3,j})_{j=0}^{3}&=(64/21,-8/3,2/3,-1/21),\\
 (\omega_{4,j})_{j=0}^{4}&=(1024/315,-64/21,8/9,-2/21,1/315).
\end{align*}
Their ordinary and $2^j$-weighted variations are shown in
\cref{dyadicweb:tab:row-constants}.

\begin{table}[ht]
\centering
\small
\begin{tabular}{@{}c c c@{}}
\toprule
$n$ & $\sum_j|\omega_{n,j}|$ & $H_n=\sum_j|\omega_{n,j}|2^j$ \\
\midrule
1 & $3$ & $4$ \\
2 & $5$ & $8$ \\
3 & $45/7$ & $80/7$ \\
4 & $51/7$ & $96/7$ \\
5 & $1683/217$ & $3264/217$ \\
\bottomrule
\end{tabular}
\caption{Two evaluations of the row polynomial: $W_n(-1)$ and $W_n(-2)$.}
\label{dyadicweb:tab:row-constants}
\end{table}

\section{A unified dyadic identity and a worked example}
\label{dyadicweb:sec:worked}

\subsection{The complete equality chain}

Let $x=m/2^s\in[0,1]$, use the terminating binary digits
$x=0.\delta_1\cdots\delta_s$, and let $N\ge s$.  Combining the preceding results gives the
following chain of exact finite expressions:
\begin{align}
 F(x)
 &=2^{-\binom s2}[X^s]\bigl(B_m(X)\mathscr A(X)\bigr)     \label{dyadicweb:eq:chain-coefficient}\\
 &=\frac{2^{-\binom{s+1}{2}}}{s!}
   \sum_{h=0}^{m-1}\ThueMorseSignSymbol_h
   \sum_{k=0}^{\Floor{s/2}}
   \binom s{2k}c_k(2m-2h-1)^{s-2k}                       \label{dyadicweb:eq:chain-moment}\\
 &=\frac{1}{2^{s^2}(1/2;1/2)_s}
   \sum_{k=0}^{s}\frac{\GaussianBinomial sk{1/2}}{4^{\binom k2}(s+k)!}
   \sum_{r=0}^{m2^k-1}\ThueMorseSignSymbol_r(r-m2^k+Q)^{s+k}              \label{dyadicweb:eq:chain-q}\\
 &=\sum_{k=0}^{\Floor{s/2}}
   \frac{c_k}{(2k)!2^{k(2s-2k+1)}}
   S_{s-2k}\!\left(\frac{m}{2^{s-2k}}\right)             \label{dyadicweb:eq:chain-moment-spline}\\
 &=\sum_{\substack{1\le j\le s\\\delta_j=1}}
   (-1)^{\delta_1+\cdots+\delta_{j-1}}
   \mathcal T_j\!\left(\sum_{\ell>j}\delta_\ell2^{-\ell}\right) \label{dyadicweb:eq:chain-binary}\\
 &=\sum_{j=0}^{N}\omega_{N,j}S_{2N-j}^{(Q)}(x).           \label{dyadicweb:eq:chain-q-spline}
\end{align}
The letter $k$ has unrelated local meanings in different lines; this is unavoidable because
one line indexes even moments and another indexes geometric scales.  The equalities are not
termwise transformations.  They are different factorizations of the same exact value:
\begin{itemize}
\item \eqref{dyadicweb:eq:chain-coefficient} is the common scalar invariant;
\item \eqref{dyadicweb:eq:chain-moment} expands the probability law in centered moments;
\item \eqref{dyadicweb:eq:chain-q} extracts the same coefficient by geometric interpolation;
\item \eqref{dyadicweb:eq:chain-moment-spline} corrects one matched spline through even moments;
\item \eqref{dyadicweb:eq:chain-binary} telescopes through the set bits;
\item \eqref{dyadicweb:eq:chain-q-spline} extrapolates consecutive spline degrees at the same point.
\end{itemize}

\subsection[The value at 5/16]{The value at $5/16$}

Take
\[
 x=\frac5{16}=0.0101_2,
 \qquad
 F(x)=\frac{305857}{2073600}
 \approx0.1475004822530864.                                \label{dyadicweb:eq:F-five-sixteen}
\]
This one value illustrates every bridge above.

\subsubsection{Master coefficient and moment formula}

The prefix polynomial is
\begin{equation}
 B_5(X)=\EulerE^{4X}-\EulerE^{3X}-\EulerE^{2X}+\EulerE^X-1.                  \label{dyadicweb:eq:B5}
\end{equation}
Using
\(
 d_0=1,d_1=1/2,d_2=5/18,d_3=1/6,d_4=143/1350
\)
in $\mathscr A$ gives
\begin{equation}
 [X^4](B_5\mathscr A)=\frac{305857}{32400}.               \label{dyadicweb:eq:B5-coefficient}
\end{equation}
Since $2^{-\binom42}=1/64$, \eqref{dyadicweb:eq:master-coefficient} gives
\eqref{dyadicweb:eq:F-five-sixteen}.  Equivalently, with
$c_1=1/9$ and $c_2=19/675$,
\begin{equation}
 F(5/16)=\frac{2^{-10}}{4!}
 \sum_{h=0}^{4}\ThueMorseSignSymbol_h
 \left[t_h^4+6c_1t_h^2+c_2\right],
 \qquad t_h=9-2h.                                         \label{dyadicweb:eq:five-sixteen-moment}
\end{equation}

\subsubsection{Moment correction of the matched spline}

The three spline values required by \eqref{dyadicweb:eq:moment-corrected-spline} are
\begin{equation}
 S_4(5/16)=\frac{1205}{8192},
 \qquad
 S_2(5/4)=\frac{15}{16},
 \qquad
 S_0(5)=-1.                                                \label{dyadicweb:eq:five-sixteen-splines}
\end{equation}
Hence
\begin{align}
 F(5/16)
 &=S_4(5/16)+\frac1{2304}S_2(5/4)
   +\frac{19}{16588800}S_0(5)                             \notag\\
 &=\frac{1205}{8192}+\frac5{12288}-\frac{19}{16588800}
 =\frac{305857}{2073600}.                                 \label{dyadicweb:eq:five-sixteen-correction}
\end{align}
The matched degree-four spline is already close, but the exact discrepancy is resolved by two
universal lower-degree corrections.

\subsubsection{The two set bits}

The binary digits equal $1$ only at positions $2$ and $4$.  The reduction polynomial at the
first position is
\begin{equation}
 \mathcal T_2(y)=\frac5{72}+y+4y^2,                        \label{dyadicweb:eq:T2-explicit}
\end{equation}
and
\(
 \mathcal T_4(0)=F(1/16)=143/2073600.
\)
The signs are $+$ at position $2$ and $-$ at position $4$, so
\begin{align}
 F(5/16)
 &=\mathcal T_2(1/16)-\mathcal T_4(0)                     \notag\\
 &=\frac{85}{576}-\frac{143}{2073600}
 =\frac{305857}{2073600}.                                 \label{dyadicweb:eq:five-sixteen-binary}
\end{align}
The prefix formula sums five Thue--Morse terms; the binary formula needs only the two set bits.

\subsubsection[Exact q-extrapolation of centered splines]{Exact $q$-extrapolation of centered splines}

At row $N=4$, \eqref{dyadicweb:eq:chain-q-spline} reads
\begin{align}
 F(5/16)
 ={}&\frac{1024}{315}S_8(5/16)
 -\frac{64}{21}S_7(5/16)
 +\frac89S_6(5/16)                                      \notag\\
 &-\frac2{21}S_5(5/16)
 +\frac1{315}S_4(5/16).                                  \label{dyadicweb:eq:five-sixteen-q-extrap}
\end{align}
The individual spline values are
\begin{equation}
 \begin{split}
 S_4&=\frac{1205}{8192},\qquad
 S_5=\frac{289799}{1966080},\qquad
 S_6=\frac{13917509}{94371840},\\
 S_7&=\frac{74236247}{503316480},\qquad
 S_8=\frac{395939357}{2684354560},
 \end{split}                                               \label{dyadicweb:eq:five-sixteen-S4-S8}
\end{equation}
all evaluated at $5/16$.  None equals $F(5/16)$, but the $q$-Pochhammer filter combines them
exactly.

\subsubsection{Finite shift dependence before stabilization}

The first discrete rows illustrate the difference between exact finite translation invariance
and asymptotic translation independence.

\begin{table}[ht]
\centering
\small
\begin{tabular}{@{}c c c c@{}}
\toprule
$n$ & $\mathcal D_n^{(0)}(5/16)$ & $\mathcal D_n^{(1/2)}(5/16)$ & $\mathcal D_n^{(3)}(5/16)$ \\
\midrule
1 & $0.156250000$ & $0.156250000$ & $3.906250000$ \\
2 & $0.151041667$ & $0.147460938$ & $-0.531250000$ \\
3 & $0.147505374$ & $0.147501983$ & $0.125788484$ \\
4 & $0.147500482$ & $0.147500482$ & $0.147500482$ \\
\bottomrule
\end{tabular}
\caption{The row becomes exactly equal to $F(5/16)$ at $n=4$, independently of $Q$.}
\label{dyadicweb:tab:five-sixteen-D}
\end{table}

The exact dyadic $q$-formula itself is independent of $Q$ already at row four, although its
outer summands vary substantially.  Their decimal contributions are:

\begin{table}[ht]
\centering
\small
\begin{tabular}{@{}c r r r@{}}
\toprule
scale $k$ & $Q=0$ & $Q=1/2$ & $Q=3$ \\
\midrule
0 & $\phantom{-}0.000626240$ & $\phantom{-}0.000466967$ & $\phantom{-}0.000000000$ \\
1 & $-0.016345021$ & $-0.014038037$ & $-0.005184307$ \\
2 & $\phantom{-}0.141657926$ & $\phantom{-}0.131089095$ & $\phantom{-}0.084488329$ \\
3 & $-0.467452954$ & $-0.449506045$ & $-0.365062120$ \\
4 & $\phantom{-}0.489014292$ & $\phantom{-}0.479488502$ & $\phantom{-}0.433258580$ \\
\midrule
sum & $0.147500482$ & $0.147500482$ & $0.147500482$ \\
\bottomrule
\end{tabular}
\caption{Different translated decompositions of the same exact dyadic value.}
\label{dyadicweb:tab:five-sixteen-q-terms}
\end{table}

\section{Fourier representation}
\label{dyadicweb:sec:fourier-representation}

The exact-arithmetic network has a complementary global integral form.  For
$Z=2Y-1$, independence gives
\begin{equation}
 \FourierTwoPi{\RvachevUp}(\xi)=\prod_{j=0}^{\infty}
 \SincRad\!\left(\frac{\pi\xi}{2^j}\right),
 \qquad
 \FourierTwoPi{f}(\xi)=\int_\RealNumbers
 f(t)\EulerE^{-2\pi\ImaginaryUnit\xi t}\Differential t.
 \label{dyadicweb:eq:up-Fourier}
\end{equation}
Integrating the Fourier inversion formula from the symmetry center gives, for
$0<x<1$,
\begin{equation}
 \boxed{
 F(x)=\frac12+\frac1\pi\int_0^\infty
 \frac{\sin(2\pi\xi(2x-1))}{\xi}
 \prod_{j=0}^{\infty}\SincRad\!\left(\frac{\pi\xi}{2^j}\right)\Differential\xi.}
 \label{dyadicweb:eq:F-Fourier}
\end{equation}
The rapid decay of the sinc product makes this a convergent integral.  Its rationality at
dyadic $x$ is a nonlocal cancellation, in contrast with the manifestly finite Appell and
$q$ formulae.

An arbitrary-argument contour representation comes from the same moment product.  If
\begin{equation}
 L_Y(s)=\Expectation e^{-sY}=\mathscr A(-s)
 =\prod_{j=1}^{\infty}\frac{1-e^{-s/2^j}}{s/2^j},         \label{dyadicq:eq:LX}
\end{equation}
then, for $\sigma>0$ and $0<x<1$, the frontier Bromwich formula is
\begin{equation}
 \boxed{
 F(x)=\frac1{2\pi\ImaginaryUnit}\lim_{T\to\infty}
 \int_{\sigma-\ImaginaryUnit T}^{\sigma+\ImaginaryUnit T}\frac{e^{sx}}s
 \prod_{j=1}^{\infty}\frac{1-e^{-s/2^j}}{s/2^j}\Differential s.} \label{dyadicq:eq:Bromwich}
\end{equation}
Thus one product has three complementary readings: coefficient extraction gives moments and
exact dyadic values, logarithmic expansion gives cumulants and Bell-polynomial formulae, and
contour inversion gives $F(x)$ directly.  This arbitrary-$x$ representation is distinct from
the inverse-dyadic Bromwich-to-coefficient collapse in the next part.

\section{Computational and conceptual consequences}
\label{dyadicweb:sec:consequences}

\subsection{Which representation is best for which task?}

The formulae are mathematically equivalent at dyadic inputs, but not computationally
interchangeable.

\begin{table}[ht]
\centering
\small
\begin{tabularx}{\textwidth}{@{}p{0.25\textwidth}p{0.31\textwidth}X@{}}
\toprule
Task & Recommended representation & Reason \\
\midrule
Prove rationality or denominator facts & master coefficient or moment formula & all arithmetic is finite and the universal moments are explicit rationals \\
Evaluate one dyadic with a sparse numerator & binary Taylor reduction & the number of reductions is the number of set bits \\
Evaluate many numerators at one depth & precompute $d_j$ or $c_k$, then use prefix recurrences & universal moment data are reused \\
Explain the $q$-coefficients & geometric divided difference and residue form & Gaussian coefficients are recognized as interpolation weights, not mysterious special-function decorations \\
Approximate arbitrary $x$ with a simple certified error & centered spline $S_p$ & direct global error $2^{-p}$ and one finite Thue--Morse sum \\
Accelerate arbitrary-argument spline convergence & discrete row $\mathcal D_n^{(Q)}$ & exact $q$-Pochhammer filter and uniform $\BigO_Q(4^{-n})$ bound \\
Use only exact dyadic arithmetic in a limiting scheme & synchronized row $\widetilde{\mathcal D}_n^{(Q)}$ or synchronized moment limit & every finite stage is an exact dyadic value and is translation-independent \\
Trace dependence on binary digits & digit-sparse series & one nonzero summand per exposed $1$-bit \\
\bottomrule
\end{tabularx}
\caption{Representation selection.  The $q$-formula is structurally illuminating but is
usually not the cheapest naive exact evaluator.}
\label{dyadicweb:tab:representation-selection}
\end{table}

The exact $q$-formula contains inner blocks of length $m2^k$, so a direct implementation can
be expensive.  Its value is primarily structural: it exposes geometric interpolation,
translation invariance, and the extrapolation filter.  For exact numerical evaluation, the
bit recursion or the master coefficient with precomputed moments is generally preferable.

For an exact dyadic $x=M/2^N$, a sparse-bit evaluator first computes the universal moments
$d_0,\ldots,d_N$ from \eqref{dyadicweb:eq:d-recurrence-companion}, evaluates the required
Appell rows by Horner's rule or \eqref{dyadicweb:eq:Appell-definition}, and sums only over the
set bits in \eqref{dyadicweb:eq:digit-Appell}.  A straightforward implementation uses
$O(N^2)$ rational arithmetic operations; this is an arithmetic-operation count, not a
bit-complexity bound.  It avoids the exponentially long inner Thue--Morse blocks in the closed
$q$ formula.

For a non-dyadic argument, the same digit--Appell series skips zero bits and has the
superexponential certificate \eqref{dyadicweb:eq:digit-Appell-tail}.  The centered nested
$q$-series can be poorly conditioned because large powers cancel.  The common-translation
gauge in \eqref{dyadicq:eq:raw-q-Appell} may reduce those intermediate magnitudes, although
exact rational or ball arithmetic remains preferable; again, the chosen translation must be
common within each finite row.

\subsection{A compact exact-arithmetic scheme}

Equation \eqref{dyadicweb:eq:d-recurrence-companion} computes all half moments with rational arithmetic,
after which \eqref{dyadicweb:eq:d-moment-prefix} evaluates a dyadic value.

\begin{algorithm}[Moment-prefix evaluator]
\label{dyadicweb:alg:moment-prefix}
Given $m,n\in\NaturalNumbers$:
\begin{enumerate}
\item Set $d_0=1$ and, for $1\le j\le n$, compute
\[
 d_j=\frac1{2^j-1}\sum_{r=1}^{j}\binom jr\frac{d_{j-r}}{r+1}.
\]
\item Form
\[
 V=\sum_{h=0}^{m-1}\ThueMorseSignSymbol_h
   \sum_{j=0}^{n}\binom njd_j(m-1-h)^{n-j}.
\]
\item Return $2^{-\binom n2}V/n!$.
\end{enumerate}
The result is $\FabiusGlobal(m/2^n)$ and is $F(m/2^n)$ when $0\le m\le2^n$.
\end{algorithm}

A literal pseudocode version is included for clarity.

\begin{lstlisting}[style=pseudo,caption={Exact dyadic evaluation from half moments.}]
d[0] := 1
for j := 1,...,n do
    s := 0
    for r := 1,...,j do
        s := s + binomial(j,r) * d[j-r] / (r+1)
    end
    d[j] := s / (2^j - 1)
end

V := 0
for h := 0,...,m-1 do
    inner := 0
    for j := 0,...,n do
        inner := inner + binomial(n,j) * d[j] * (m-1-h)^(n-j)
    end
    V := V + (-1)^(binary_weight(h)) * inner
end
return V / (2^(n*(n-1)/2) * n!)
\end{lstlisting}

The centered-moment formula can be evaluated similarly.  Its advantage is parity: only
$\Floor{n/2}+1$ moments occur.  The half-moment formula aligns more directly with the
binary Taylor polynomials.

\subsection{Exact arithmetic checks implied by the connections}

The network of representations supplies a strong verification suite.  For chosen $m,n$ one
can independently check:
\begin{enumerate}
\item the master coefficient against the centered-moment formula;
\item the half-base $q$ expression at two or more translations $Q$;
\item the terminating binary reduction when $m/2^n\in[0,1]$;
\item the moment-corrected matched-spline identity;
\item the exact $q$-extrapolation row at any order $N\ge n$;
\item representation independence under $(m,n)\mapsto(2m,n+1)$.
\end{enumerate}
These checks stress different cancellation patterns.  Agreement is therefore more informative
than recomputing the same formula in a rearranged order.

The last item is now more than a check: \path{qBinomialThueMorseDyadicFormula_refine}
proves it exactly.  The stronger
\path{qBinomialThueMorseDyadicTranslatedFormula_eq_of_rat_eq} allows any two
numerator--exponent pairs representing the same rational dyadic value and permits different
rational translations on the two sides.  The translated refinement and half-shifted variants
are proved in \path{FabiusQBinomialFormula.lean}.  These theorem-level invariances do not
formalize the separate finite-row stabilization question in \cref{dyadicweb:sec:discrete}.

\subsection{Endpoint and convention warnings}

Several small conventions are mathematically active.

\begin{warning}[The scale-zero term]
The global binary series begins at $m=0$.  It happens to vanish for $0\le x<1$, but it equals
$1$ at $x=1$ and controls the signed behavior on integer blocks.  A series beginning at
$m=1$ is valid only on the half-open unit interval.
\end{warning}

\begin{warning}[Integer exponents]
The exponent
$\binom{j+1}{2}-\binom{m-j}{2}$ in \eqref{dyadicweb:eq:T-def-local} is an integer and may be negative.
Replacing integer subtraction by truncated natural subtraction changes the polynomial.
\end{warning}

\begin{warning}[Centered powers are signed]
Expressions such as $(r-m2^k+Q)^{n+k}$ live in a characteristic-zero scalar field.  The
quantity $r-m2^k$ is deliberately negative on much of the block.  Truncated subtraction in
$\NaturalNumbers$ destroys the formula.
\end{warning}

\begin{warning}[Binary expansions at dyadics]
The floor prefixes $p_m=\Floor{2^mx}$ select the terminating expansion at a dyadic
rational.  Replacing it by the alternative expansion ending in infinitely many $1$-digits
changes the termwise digit series even though the represented real number is unchanged.
\end{warning}

\section{Relation to the Lean development}
\label{dyadicweb:sec:formal-map}

The principal source theorems are formalized in the following modules.  Filenames are relative
to \repo.

\begingroup
\small
\setlength\LTleft{0pt}
\setlength\LTright{0pt}
\begin{longtable}{@{}p{0.34\textwidth}p{0.62\textwidth}@{}}
\toprule
Module & Role in the present article \\
\midrule
\endfirsthead
\toprule
Module & Role in the present article \\
\midrule
\endhead
\path{DyadicClosedForm.lean} & Exact Thue--Morse/even-moment formula for arbitrary dyadic numerators. \\
\path{HalfQBinomial.lean} & Finite $q$-Pochhammer and Gaussian-binomial calculus at $q=1/2$, including the interpolation annihilation row. \\
\path{HalfQBinomialRootSimplicity.lean} & Exactly one public theorem and no definitions: \lean{halfQBinomial_sum_rootMultiplicity_two_pow} proves that every dyadic root $2^j$, $j<n$, of the exact rational half-base coefficient polynomial has multiplicity one.  Together with the preceding complete rational root classification, all roots in that locus are simple; no arbitrary-base or arbitrary-field claim is made. \\
\path{QBinomialReciprocity.lean} & Exactly four public theorems: \lean{gaussianBinomial_reciprocity_units}, \lean{gaussianBinomial_reciprocity}, \lean{gaussianBinomial_neg_one_eq_zero_of_odd_degree}, and \lean{gaussianBinomial_neg_one_even_odd_eq_zero}.  They give unit-valued reciprocity over every commutative semiring, its nonzero-element semifield wrapper, odd-degree vanishing at \(q=-1\) over every commutative ring, and the total even-row/odd-column specialization.  The two reciprocity identities are total above the diagonal by zero extension; no quotient, cancellation, domain, characteristic, topology, or convergence assumption is used. \\
\path{GaussianBinomialPalindromic.lean} & Exactly fourteen public theorems and no definitions: generic reflection helpers, Gaussian degree bound, zero/diagonal values, reflection, constant/top coefficients, coefficient symmetry, division-free mean, exact degree/monicity over a nontrivial commutative semiring, and the arbitrary-commutative-semiring linear-coefficient classifier, equal to one exactly for \(0<k<n\) and zero otherwise. \\
\path{GaussianBinomialPolynomialStructure.lean} & Exactly five public theorems and no definitions: exact universal natural degree, monicity, constant coefficient one, reflection at degree \(k(n-k)\), and bounded-index coefficient palindromicity; the adjacent generic module supplies the complete linear-coefficient classifier. \\
\path{CentralQBinomialReduction.lean} & Exactly six public theorems and no definitions: sign pairing, even--odd dissection, naturality, integral and general division-free central reduction, and the field quotient under two nonvanishing hypotheses. \\
\path{CyclotomicFactorization.lean} & Exactly seven public theorems and no definitions: floor-quotient bounds, divisor reindexing, commutative-ring finite-product identities, interval extension, and Gaussian cyclotomic factorization over an integral domain. \\
\path{PrimitiveRootBlock.lean} & Exactly three public theorems and no definitions: \lean{gaussianBinomial_isPrimitiveRoot_eq_zero}, \lean{neg_one_pow_mul_pow_choose_two}, and \lean{finiteQPochhammerIn_isPrimitiveRoot}; in a commutative integral domain they give interior vanishing and, for positive primitive-root order, the top phase and complete block. \\
\path{QLucas.lean} & Exactly seven public theorems and no definitions: \lean{add_mul_add_sub_one}, \lean{choose_two_add}, \lean{coeff_finiteQPochhammerIn_neg_X}, \lean{finiteQPochhammerIn_neg_X_block}, \lean{coeff_block_pow_mul}, \lean{pow_choose_two_add_mul_eq}, and \lean{gaussianBinomial_q_lucas}.  The final theorem assumes only positive primitive-root order and the residue bounds \(b,s<d\), not \(r\le a\) or \(s\le b\).  Its local \path{two_mul_choose_two} is private; the public theorem belongs to \path{QChuVandermonde.lean}. \\
\path{TwoPhiOneReversal.lean} & Exactly two public definitions and twelve theorems: \lean{twoPhiOneFinite}, \lean{twoPhiOneReflection}, \lean{choose_two_add_succ_choose_two}, \lean{finiteQPochhammerIn_sub_eq}, \lean{finiteQPochhammerIn_reversal_ne_zero}, \lean{finiteQPochhammerIn_inv_pow_self}, \lean{twoPhiOneReflection_involutive}, \lean{twoPhiOneFinite_reversal}, \lean{twoPhiOneFinite_reversal_twice}, \lean{twoPhiOneFinite_eq_sum_twoPhiOneTerm}, \lean{twoPhiOne_eq_twoPhiOneFinite_inv_pow}, \lean{twoPhiOne_reversal}, \lean{twoPhiOne_reversal_twice}, and \lean{twoPhiOne_one_eq_twoPhiOneFinite_zero}. \\
\path{QChuVandermonde.lean} & Exactly ten public theorems and no definitions: \lean{two_mul_choose_two}, \lean{mul_sub_one_eq_mul_sub_add}, \lean{finiteQPochhammerIn_div_eq_sum_chu}, \lean{q_chu_vandermonde_first}, \lean{finiteQPochhammerIn_div_eq_sum_chu_second}, \lean{twoPhiOneFinite_mul_finiteQPochhammerIn_eq_chu_second}, \lean{q_chu_vandermonde_second}, \lean{q_chu_vandermonde_second_by_reversal}, \lean{twoPhiOne_q_chu_vandermonde_first}, and \lean{twoPhiOne_q_chu_vandermonde_second}.  Both full-domain formulas are exact; the named by-reversal route alone retains the additional \(C\ne0\) and \((A;q)_n\ne0\) hypotheses. \\
\path{JacobiTwoSquareCount.lean} & Exactly four public theorems and no definitions: \lean{sumSqRep_two_eq_four_mul_twoSquareDivisorSum}, \lean{sumSqRep_two_eq_four_mul_prod}, \lean{theta_sq_eq_chi4_lambert}, and \lean{theta_sq_eq_odd_lambert}.  The count holds for nonzero natural input, the product keeps its even-valuation hypothesis, and both Lambert identities require only \(\lVert q\rVert<1\) in a complete normed field. \\
\path{CyclotomicDivisibility.lean} & Exactly three public theorems and no definitions: \lean{cyclotomic_exponent_eq_one_iff}, \lean{cyclotomic_dvd_gaussianBinomial_iff}, and \lean{gaussianBinomial_mul_isPrimitiveRoot}.  The carry criterion assumes \(k\le n\), \(d>0\); divisibility is in \(\RationalNumbers[X]\); the primitive-root value needs no \(b\le a\). \\
\path{QCatalan.lean} & Exactly one public definition, \lean{qCatalan}, and eleven theorems: \lean{map_qInt}, \lean{qInt_X_monic}, \lean{qInt_X_natDegree}, \lean{X_sub_one_mul_qInt}, \lean{qInt_X_eq_prod_cyclotomic}, \lean{qInt_X_dvd_gaussianBinomial_rat}, \lean{qInt_X_dvd_gaussianBinomial_int}, \lean{qInt_X_mul_qCatalan}, \lean{qCatalan_natDegree}, \lean{qCatalan_eval_one_mul}, and \lean{qCatalan_eval_one}.  The integral \path{divByMonic} quotient is defined for every natural \(n\); coefficient positivity, palindromicity, and a combinatorial interpretation are not asserted. \\
\path{NewtonInterpolation.lean} & Exactly three definitions and nineteen theorems: \lean{newtonCoeff}, \lean{nodeNewtonPoly}, \lean{newtonInterpolant}; \lean{newtonCoeff_eq}, \lean{newtonCoeff_zero}, \lean{newtonCoeff_mul_prod}, \lean{nodeNewtonPoly_succ}, \lean{eval_nodeNewtonPoly}, \lean{degree_nodeNewtonPoly_lt}, \lean{nodeNewtonPoly_eq_interpolate}, \lean{eq_nodeNewtonPoly_of_eval_eq}, \lean{coeff_nodeNewtonPoly_self}, \lean{newtonCoeff_eq_sum}, \lean{nodal_range_pow}, \lean{prod_erase_pow_sub_pow}, \lean{newtonCoeff_pow_eq_sum}, \lean{newtonPoly_succ}, \lean{eval_newtonPoly}, \lean{degree_newtonPoly_lt}, \lean{newtonPoly_eq_interpolate}, \lean{eq_newtonPoly_of_eval_eq}, and \lean{coeff_newtonPoly_self}. \\
\path{QBetaIntegral.lean} & Exactly one definition and eight theorems: \lean{qBeta}; \lean{qNumber_pos}, \lean{qBeta_term_eq}, \lean{qBeta_eq_prod}, \lean{qBeta_eq_qGamma}, \lean{qBeta_comm}, \lean{qBeta_pos}, \lean{qBeta_add_one_left}, and \lean{qBeta_add_one_right}. \\
\path{QPochhammerElementaryIdentities.lean} & Exactly 13 public theorems: unit-valued ring reversals, root-safe terminating and adjacent cross identities, and field quotient forms with their exact nonzero-denominator hypotheses; the cross identities are total by zero extension, while the adjacent quotients retain \(k<n\). \\
\path{QDifferenceAnnihilation.lean} & Exactly four public theorems: the scaled inverse-Gaussian characteristic polynomial, its complete geometric monomial moments, scalar-extension top-coefficient extraction, and strict-degree \(q\)-difference annihilation; \cref{dyadicq:rem:scaled-qdiff}. \\
\path{GeometricQBinomialLagrange.lean} & Denominator-free geometric Lagrange numerators and residual moments, including the global identification of the numerator with the scaled inverse kernel at base and scale \(q\); normalized quotient results retain their distinct-node field hypotheses. \\
\path{QBinomialInversionSpecializations.lean} & Exactly two public definitions and four public theorems: the residual and reconstruction rows at Gaussian base \(q^2\), scale \(-q\), and their pointwise closed forms need only a ring; exactly the two total convolution inverses require a commutative ring. \\
\path{ThueMorseExponential.lean} & Complete-block exponential factorization, Prouhet cancellation, and translation of inner powers. \\
\path{ThueMorseComplexExponential.lean} & Weighted and signed finite complex-exponential block transforms, including arbitrary affine shifts and the empty block. \\
\path{FabiusQBinomialFormula.lean} & Centered and translated finite $q$-binomial formulas, arbitrary numerators, common-translation invariance, and invariance under equal rational dyadic representations. \\
\path{FabiusQBinomialTaylor.lean} & Identification of the finite $q$ expression with the Taylor reduction polynomial. \\
\path{FabiusQBinomialScalarFormula.lean} and \path{FabiusDyadicQBinomialScalar.lean} & Scalar normalizations, real/complex translated forms, and the constant-polynomial coefficient theorem. \\
\path{FabiusRawQBinomialFormula.lean} and \path{FabiusRawQBinomialScalar.lean} & Raw inverse-dyadic forms and their scalar versions. \\
\path{FabiusBinaryReductionSeries.lean} & Exact binary residual telescope, absolute convergence, and uniform tail bound. \\
\path{FabiusParityPowerSeries.lean} & Corrected scale-zero parity-power presentation. \\
\path{FabiusGlobalQBinomialSeries.lean} & Termwise substitution of finite $q$-identities into the binary telescope. \\
\path{FabiusComplexShiftSpline.lean} & Shifted finite splines, fixed-cutoff Taylor expansion, and uniform shift bound. \\
\path{FabiusDiscreteLimitToeplitz.lean} & Row mass, exact total variation, uniform variation bound, and finite-row Toeplitz theorem. \\
\path{FabiusDiscreteLimitIntegration.lean} & Reindexing of the discrete formula and convergence to the signed global function. \\
\path{TaylorReduction.lean} & Binary Taylor-reduction polynomials, their exact analytic identity, and finite termination on dyadic inputs. \\
\path{FabiusUniformSpline.lean} & Centered finite-power spline approximants and the uniform analytic error estimate used by the discrete rows. \\
\bottomrule
\end{longtable}
\endgroup

Except where explicitly marked below as frontier-dependent, the following statements are
direct paper-level corollaries or reorganizations of those formalized results:
\begin{itemize}
\item the expectation and two contour representations
      \eqref{dyadicweb:eq:expectation-Y}, \eqref{dyadicweb:eq:cauchy-X}, and \eqref{dyadicweb:eq:q-node-contour};
\item the composition and Bell-polynomial formulae for $d_n$;
\item the moment-corrected spline transform and its reciprocal
      \eqref{dyadicweb:eq:moment-corrected-spline}--\eqref{dyadicweb:eq:reciprocal-spline};
\item exact finite stabilization of the discrete row at dyadic arguments;
\item the row polynomial $W_n$, the shift-operator form, and the $q$-Richardson interpretation;
\item the explicit $\BigO_Q(4^{-n})$ bound \eqref{dyadicweb:eq:D-error}, which remains a
      frontier sharpening because it depends on the unformalized shifted-spline estimate
      \eqref{dyadicweb:eq:shifted-spline-bound-local}, whose proof obligation and edge-case
      warning are recorded in the Primary Exposition Gap Register;
\item the synchronized nearest-dyadic $q$- and moment-limits.
\end{itemize}
Apart from that frontier sharpening, the listed statements require no new unproved property
of the Fabius function; each direct corollary follows by finite algebra, coefficient
extraction, or the already formalized uniform spline estimates.  The direct corollaries are
marked here to distinguish conceptual reorganization from theorem names that already occur
verbatim in the repository.

\section{Conclusion}
\label{dyadicweb:sec:conclusion}

The exact dyadic, global-series, and discrete-limit formulae are best understood not as a
collection of special tricks but as a small number of reusable operators acting in different
variables.

The product $B_m\mathscr A$ separates numerator arithmetic from the universal probability
law.  Expanding its coefficient in centered moments gives the classical Thue--Morse formula;
reading it as an expectation gives a probabilistic representation; taking a Cauchy integral
gives an analytic one.  Introducing the geometric node $z$ turns the same coefficient into
the leading term of a degree-$n$ polynomial.  The Gaussian row at base $1/2$ is then simply
the barycentric divided difference on $-1,-2,\ldots,-2^n$, and the $q$-Pochhammer symbol is the
node polynomial of that grid.

The centered finite spline sits directly inside the exact arithmetic.  Its matched-grid value
is the zeroth centered-moment term, and the remaining moments form an exact, invertible defect
correction.  A second correction uses the normalized $q$-Pochhammer polynomial in the spline
degree: it preserves constants, annihilates geometric error modes, and yields exact recovery
at dyadic points.  The resulting discrete row is therefore naturally viewed as a
$q$-Richardson extrapolator, not merely as an opaque Toeplitz average.

Finally, the two arbitrary-argument $q$ constructions have distinct logical roles.  The global
$q$-series is the binary residual telescope with an alternative finite formula inserted into
each summand.  The discrete limit is a finite extrapolation row whose shift dependence
disappears only asymptotically.  Their meeting point is the dyadic grid: there the binary
series terminates, the discrete row stabilizes, the synchronized row is exact, and all roads
return to the same master coefficient.

\renewcommand{\refname}{Topic references}
\begin{thebibliography}{99}

\bibitem{dyadicweb:fabius1966}
J.~Fabius,
\newblock A probabilistic example of a nowhere analytic $C^\infty$-function,
\newblock \emph{Zeitschrift fuer Wahrscheinlichkeitstheorie und Verwandte Gebiete}
  \textbf{5} (1966), 173--174.
\newblock DOI: \nolinkurl{10.1007/BF00536652}.

\bibitem{dyadicweb:arias1982}
J.~Arias de Reyna,
\newblock Definicion y estudio de una funcion indefinidamente diferenciable de soporte
  compacto,
\newblock \emph{Revista de la Real Academia de Ciencias Exactas, Fisicas y Naturales de
  Madrid} \textbf{76} (1982), no.~1, 21--38.
\newblock English translation: \emph{An infinitely differentiable function with compact
  support: Definition and properties}, arXiv:\nolinkurl{1702.05442} (2017).

\bibitem{dyadicweb:arias2018}
J.~Arias de Reyna,
\newblock Arithmetic of the Fabius function,
\newblock \emph{INTEGERS} \textbf{18} (2018), Article A51;
  arXiv:\nolinkurl{1702.06487}, version~3.

\bibitem{dyadicweb:haugland2016}
J.~K.~Haugland,
\newblock Evaluating the Fabius function,
\newblock arXiv:\nolinkurl{1609.07999} (2016).

\bibitem{dyadicweb:allouche2003}
J.-P.~Allouche and J.~Shallit,
\newblock \emph{Automatic Sequences: Theory, Applications, Generalizations},
\newblock Cambridge University Press, 2003.

\bibitem{dyadicweb:andrews1976}
G.~E.~Andrews,
\newblock \emph{The Theory of Partitions},
\newblock Encyclopedia of Mathematics and its Applications~2,
  Addison--Wesley, 1976.

\bibitem{dyadicweb:gasper2004}
G.~Gasper and M.~Rahman,
\newblock \emph{Basic Hypergeometric Series}, second edition,
\newblock Encyclopedia of Mathematics and its Applications~96,
  Cambridge University Press, 2004.

\bibitem{dyadicweb:kac2002}
V.~Kac and P.~Cheung,
\newblock \emph{Quantum Calculus},
\newblock Universitext, Springer, 2002.

\bibitem{dyadicweb:dlmf17}
NIST Digital Library of Mathematical Functions, Chapter~17,
\newblock \emph{$q$-Hypergeometric and Related Functions},
\newblock \url{https://dlmf.nist.gov/17}.

\bibitem{dyadicweb:reshetnikovMSE}
V.~Reshetnikov,
\newblock A conjectured closed form for the Fabius function at dyadic rationals,
\newblock Mathematics Stack Exchange, question 3283519, 5 July 2019.

\bibitem{dyadicweb:proveit}
V.~Reshetnikov,
\newblock \emph{ProveIt: Analysis/FabiusFunction},
\newblock Lean~4 formal development, repository snapshot
  \path{75dac857e96d858140e7e5520172d0c659bde98e}, 24 August 2026.
\newblock \url{https://github.com/VladimirReshetnikov/ProveIt/tree/75dac857e96d858140e7e5520172d0c659bde98e/Analysis/FabiusFunction}.

\bibitem{dyadicweb:mainarticle}
V.~Reshetnikov,
\newblock \emph{The Fabius Function and Rvachev's Up-Function: Exact arithmetic,
  probability, Fourier analysis, and corrected sharp asymptotics},
\newblock manuscript in the \texttt{ProveIt} repository, snapshot 25 August 2026.
\newblock \url{https://github.com/VladimirReshetnikov/ProveIt/blob/main/Analysis/FabiusFunction/docs/Fabius_Function_and_Rvachev_Up/Fabius_Function_and_Rvachev_Up.tex}.

\end{thebibliography}


\renewcommand{\arraystretch}{1}
% END SYNTHESIZED TOPIC: exact dyadic calculus and alternative representations

% BEGIN SYNTHESIZED TOPIC: exact dyadic arithmetic to endpoint asymptotics
% Former source SHA-256 values:
% c5ad7c5298d958ab63f459a3e246c2293d3020525223bc05ef2d0e1edab58f10
% 6e98efd3afc402159a7f080e8604c951d5bc51be4a73383da67c1241d46d54ca
\clearpage
\part{From Exact Dyadic Formulae to the Small-Argument Asymptotic}
\label{part:dyadicasym}
\begin{boxedremark}
\textbf{Synthesized provenance and status.} This part merges the former
\nolinkurl{Fabius_Dyadic_Asymptotic_Bridge/Fabius_Dyadic_Asymptotic_Bridge.tex} and
\nolinkurl{Fabius_Dyadic_Formulae_to_Asymptotics/Fabius_Dyadic_Formulae_to_Asymptotics.tex}.
The latter supplies the canonical normalization, moment--Laplace route, exact-phase theorem,
status ledger, and Lean correspondence.  The former contributes the Mellin calculation,
vertical coefficient extraction, fixed and moderately growing numerator rays, exact
Bromwich-to-coefficient collapse, and continuum-boundary analysis.  Repeated recurrence,
composition, and $q$-binomial preliminaries now refer to \cref{part:dyadicweb}.  The vertical
and growing-prefix routes remain natural-language formalization obligations.\par\smallskip
Former source SHA-256 values:
\nolinkurl{c5ad7c5298d958ab63f459a3e246c2293d3020525223bc05ef2d0e1edab58f10} and
\nolinkurl{6e98efd3afc402159a7f080e8604c951d5bc51be4a73383da67c1241d46d54ca}.
\end{boxedremark}
\section*{Topic abstract}
The exact arithmetic of the Fabius function and its endpoint asymptotics are usually presented
as two nearly disjoint subjects.  At reciprocal powers of two one has rational recurrences,
composition sums, Thue--Morse power sums, and a finite half-base $q$-binomial formula; near the
origin one has a negative-Laplace product, a lower-Lambert saddle, and a logarithmically
periodic correction.  This report supplies the missing bridge.

The central exact observation is that the apparently formidable $q$-binomial numerator at
$F(2^{-n})$ is not a new asymptotic object.  After choosing a convenient scalar translation,
the inner Thue--Morse block becomes a coefficient of the ratio
$A(-2^kX)/A(-X)$, where
\[
 A(z)=\sum_{n\ge0}a_nz^n
     =\prod_{j\ge1}\frac{\EulerE^{z/2^j}-1}{z/2^j},
 \qquad
 F(2^{-n})=2^{-\binom n2}a_n.
\]
The Gaussian row at $q=1/2$ is exactly the divided-difference functional on the geometric
nodes $-1,-2,\ldots,-2^n$; it annihilates all lower powers and extracts the leading
coefficient.  Consequently the complete finite numerator collapses identically to
$\M_na_n$, where $\M_n=\prod_{j=1}^n(2^j-1)$.  This ``decompression'' proves that the
$q$-Pochhammer, $q$-binomial, Thue--Morse, recurrence, composition, and generating-product
formulae all encode the same coefficient sequence.

Two asymptotic derivations are then developed.  The fully rigorous real-variable route writes
the half moment as $d_n=\Expectation(1-X)^n$, compares it through second order with the negative
Laplace transform $A(-n)=\Expectation\EulerE^{-nX}$, and uses the exact quadratic-plus-periodic
decomposition of $\log A(-s)$.  A complementary coefficient-saddle route applies Cauchy's
formula directly to $a_n=[z^n]A(z)$; it reaches the same lower-Lambert coordinate and provides
an independent explanation of every cancellation.  If $L=\log2$ and $\lambda_n$ is the
large solution of
\[
 \lambda_n-\frac{\log\lambda_n}{L}=n,
 \qquad
 \lambda_n=-\frac1L W_{-1}(-L2^{-n}),
\]
then the resulting sharp dyadic formula is
\[
 \log F(2^{-n})=
 -\frac L2\lambda_n^2+
 \left(1+\frac L2\right)\lambda_n-\frac12\log\lambda_n
 +C_F+\Psi(\lambda_n)+\BigO(n^{-1}),
\]
where $\Psi$ is a smooth, nonconstant, mean-zero, one-periodic function with explicit
Gamma--zeta Fourier coefficients.  Expanding the Lambert phase yields
\[
 \log F(2^{-n})=
 -\frac L2n^2-n\log n+\left(1+\frac L2\right)n
 -\frac{\log^2n}{2L}+C_F-
 \frac{\log^2n}{2L^2n}+\Psi(\lambda_n)+\BigO(n^{-1}).
\]
The exact phase is essential: replacing it by $\log_2n$ leaves a generally larger
$(\log n)/n$ term.  Corollaries give equally sharp asymptotics for the recurrence sequence,
the half moments, the composition sum, and the literal $q$-binomial numerator.  Numerical
checks against exact rational values confirm the constant, phase, and first saddle
correction.


\section{The missing bridge}
\label{dyadicasym:sec:introduction}

The Fabius function admits two descriptions that look, at first sight, almost antagonistic.
Its values at dyadic rationals are algebraic and finite: every value is rational; there are
terminating recurrences; the Thue--Morse sequence provides finite signed power sums; and a
half-base Gaussian-binomial row packages those sums into a compact formula.  Its endpoint
asymptotic is analytic and infinite: one studies a Laplace product, separates a periodic
renormalization of a dyadic harmonic sum, and evaluates a saddle determined by the lower real
branch of Lambert's $W$-function.

The two subjects are not merely compatible.  They are two coordinate systems on the same
object.  The exact formulae contain the asymptotic information in a highly compressed form,
and the purpose of this report is to make the decompression explicit.

\subsection{What is proved}

There are four levels of connection.

\begin{enumerate}
\item The reciprocal-power value $F(2^{-n})$ is exactly a half moment divided by a factorial
and a triangular power of two.  This already gives the quadratic logarithmic decay.

\item After factorial normalization, all reciprocal-power values are coefficients of one
entire function $A$.  The recurrence, composition sum, and infinite product are equivalent
presentations of $A$.

\item The finite $q$-binomial/Thue--Morse formula can be reduced identically to the same
coefficient $a_n=[z^n]A(z)$.  No limiting argument is involved.  The $q$-binomial row is a
geometric divided difference; its role is to undo the blockwise Thue--Morse cancellation.

\item The product for $A(-s)$ has an exact quadratic-plus-periodic logarithm.  Transferring
from $A(-n)$ to the endpoint moment, or applying a coefficient saddle directly to $A$, yields
the corrected lower-Lambert asymptotic.
\end{enumerate}

The first, second, third, and the full real-variable fourth level are represented in the Lean
formalization under \repo.  The coefficient-saddle calculation in \Cref{dyadicasym:sec:coefficient-saddle}
is included as an independent conventional-mathematics derivation; it is useful because it
starts from the coefficient formula itself and exposes the same cancellations with very
little probability language.

\subsection{Headline notation and result}

Throughout,
\begin{equation}
 L=\log2.                                                     \label{dyadicasym:eq:L}
\end{equation}
Let $P$ be the periodic correction in the exact negative-Laplace decomposition, let
\begin{equation}
 \overline P=\int_0^1P(u)\Differential u,
 \qquad \Psi(u)=P(u)-\overline P,                            \label{dyadicasym:eq:Psi-intro}
\end{equation}
and put
\begin{equation}
 A_0=\frac{6\gamma^2+12\gamma_1-\pi^2}{12},
 \qquad
 C_F=\frac{A_0}{L}-\frac{7L}{12}-\frac12\log\pi.            \label{dyadicasym:eq:CF-intro}
\end{equation}
Here $\gamma$ is Euler's constant and $\gamma_1$ is the first Stieltjes constant.
For every sufficiently large integer $n$, let $\lambda_n$ be the large positive solution of
\begin{equation}
 \lambda_n2^{-\lambda_n}=2^{-n},
 \qquad\text{equivalently}\qquad
 \lambda_n-\frac{\log\lambda_n}{L}=n.                       \label{dyadicasym:eq:lambda-n-intro}
\end{equation}

\begin{theorem}[Sharp reciprocal-power asymptotic]
\label{dyadicasym:thm:headline}
As $n\to\infty$,
\begin{equation}
 \boxed{
 \log F(2^{-n})=
 -\frac L2\lambda_n^2+
 \left(1+\frac L2\right)\lambda_n-\frac12\log\lambda_n
 +C_F+\Psi(\lambda_n)+\BigO(n^{-1}).}                       \label{dyadicasym:eq:headline-compact}
\end{equation}
Equivalently,
\begin{align}
 \log F(2^{-n})={}&-\frac L2n^2-n\log n+
 \left(1+\frac L2\right)n-\frac{\log^2n}{2L}+C_F           \notag\\
 &-\frac{\log^2n}{2L^2n}+\Psi(\lambda_n)+\BigO(n^{-1}).   \label{dyadicasym:eq:headline-expanded}
\end{align}
The function $\Psi$ is smooth, mean-zero, one-periodic, and nonconstant.  Its Fourier
coefficients are explicit.
\end{theorem}

The rest of the report derives \Cref{dyadicasym:thm:headline} from the exact arithmetic formulae rather
than taking it as an externally supplied endpoint theorem.

\begin{boxedremark}
The most important conceptual point is that the periodic term is already present in the exact
dyadic values.  It does not arise because one interpolates between dyadic points.  The
fractional parts of the exact saddle coordinates $\lambda_n$ drift through the logarithmic
period, so the single rational sequence $F(2^{-n})$ samples the entire periodic correction.
\end{boxedremark}

\section{Exact dyadic inputs from the preceding part}
\label{dyadicasym:sec:exact-inputs}

The exact algebra needed here is developed once in
\cref{part:dyadicweb}.  In the present notation,
\begin{equation}
 d_n=\Expectation X^n=\Expectation(1-X)^n,\qquad
 a_n=\frac{d_n}{n!},\qquad
 A(z)=\sum_{n\ge0}a_nz^n=\mathscr A(z),                  \label{dyadicasym:eq:an-A}
\end{equation}
with
\begin{align}
 F(2^{-n})&=2^{-\binom n2}a_n,                             \label{dyadicasym:eq:F-an}\\
 A(z)&=\prod_{j\ge1}\frac{\EulerE^{z/2^j}-1}{z/2^j},          \label{dyadicasym:eq:A-product}\\
 A(2z)&=\frac{\EulerE^z-1}{z}A(z),\qquad A(z)=\EulerE^zA(-z).     \label{dyadicasym:eq:A-reflection}
\end{align}
The first equality is \eqref{dyadicweb:eq:inverse-dyadic-input}; the product and
refinement are \eqref{dyadicweb:eq:A-product}--\eqref{dyadicweb:eq:A-refinement}.
The recurrence, ordered-composition formula, Mersenne normalization, and general-numerator
master coefficient are respectively
\eqref{dyadicweb:eq:d-recurrence-companion},
\eqref{dyadicweb:eq:ordered-composition},
\eqref{dyadicweb:eq:q-barycentric-normalization}, and
\eqref{dyadicweb:eq:master-coefficient}.

The normalized endpoint and generating-product identifications are the exact declarations
\path{fabius_inverse_two_pow_eq_recurrenceSequence},
\path{halfMomentGeneratingSeries_eq_fabiusRecurrenceSequence_series},
\path{complexGeneratingFunction_eq_fabiusRecurrenceSequence_series}, and
\path{fabiusRecurrenceSequence_series_neg_eq_tprod} in
\path{FabiusRecurrenceSequence.lean}.

For later asymptotic bookkeeping, let
\(\M_n=\prod_{j=1}^n(2^j-1)\).  The literal inverse-power half-base
$q$-numerator is the finite geometric extractor applied to the degree-$n$ coefficient
polynomial, and the exact collapse reads
\begin{equation}
 \mathcal N_n=\M_na_n.                                    \label{dyadicasym:eq:N-m-a}
\end{equation}
This is the $m=1$ specialization of the general formula in
\cref{dyadicweb:sec:q-extraction}; it is not a separate asymptotic input.  Thus every
recurrence, composition, moment, and literal $q$ consequence below is a corollary of one
coefficient sequence rather than a parallel derivation.

\section{The first asymptotic already visible in the exact values}
\label{dyadicasym:sec:coarse}

Before resolving the periodic fluctuation, the endpoint identity alone gives the principal
quadratic law.  Taking logarithms in \eqref{dyadicasym:eq:F-an},
\begin{equation}
 \log F(2^{-n})=\log d_n-\log(n!)-\binom n2L.               \label{dyadicasym:eq:log-F-exact}
\end{equation}
Because $0\le\RvachevUp\le1$,
\begin{equation}
 d_n\le n\int_0^1t^{n-1}\Differential t=1.                           \label{dyadicasym:eq:dn-upper}
\end{equation}
A fixed positive portion of the mass on $[0,1/2]$ gives, for $n\ge1$,
\begin{equation}
 2^{-(n+1)}\le d_n.                                         \label{dyadicasym:eq:dn-lower}
\end{equation}
Combining these bounds with elementary factorial estimates yields the explicit inequality
\begin{equation}
 \left|\log F(2^{-n})+\frac L2n^2\right|
 \le3n\log(n+1),\qquad n\ge1.                              \label{dyadicasym:eq:coarse-bound}
\end{equation}
Therefore
\begin{equation}
 \lim_{n\to\infty}\frac{\log F(2^{-n})}{n^2}=-\frac L2.    \label{dyadicasym:eq:coarse-limit}
\end{equation}

This is already a derivation of the leading endpoint asymptotic from exact arithmetic.  It
also explains why neither the factorial nor the half moment can change the coefficient of
$n^2$: $\log(n!)=\BigO(n\log n)$ and $\log d_n=\BigO(n)$ under the crude bounds.  To recover
the $-n\log n$ term, the $\log^2n$ term, the constant, and the periodic fluctuation, however,
one must understand $d_n$ much more sharply.

\section{The exact negative-Laplace decomposition}
\label{dyadicasym:sec:laplace-product}

The sharp asymptotic enters through the same function $A$.  For $s>0$ put
\begin{equation}
 B(s)=A(-s)=\Expectation\EulerE^{-sX},
 \qquad
 \Lam(s)=\log B(s).                                         \label{dyadicasym:eq:B-Lambda}
\end{equation}
From \eqref{dyadicasym:eq:A-product},
\begin{equation}
 B(s)=\prod_{j=1}^{\infty}
 \frac{1-\EulerE^{-s/2^j}}{s/2^j},                               \label{dyadicasym:eq:B-product}
\end{equation}
and hence
\begin{equation}
 \Lam(s)=\sum_{j=1}^{\infty}\kappa(s/2^j),
 \qquad
 \kappa(u)=\log\frac{1-\EulerE^{-u}}u.                          \label{dyadicasym:eq:Lambda-harmonic}
\end{equation}
This is a dyadic harmonic sum.  Its dilation equation is exact:
\begin{equation}
 \Lam(2s)=\Lam(s)+\log(1-\EulerE^{-s})-\log s.                  \label{dyadicasym:eq:Lambda-dilation}
\end{equation}

\subsection{Periodization of the dilation equation}

Define the forward logarithmic tail
\begin{equation}
 T(s)=\sum_{m=0}^{\infty}\log(1-\EulerE^{-s2^m})<0,
 \qquad
 \mathcal E(s)=-T(s)>0.                                    \label{dyadicasym:eq:tail}
\end{equation}
Since $2^m\ge m+1$,
\begin{equation}
 0\le\mathcal E(s)\le
 \frac{\EulerE^{-s}}{(1-\EulerE^{-s})^2},                            \label{dyadicasym:eq:tail-bound-report}
\end{equation}
and in particular $\mathcal E(s)\le4\EulerE^{-s}$ for $s\ge L$.
For $u\in\RealNumbers$, set
\begin{equation}
 P(u)=\Lam(2^u)+\frac L2(u^2-u)+T(2^u).                    \label{dyadicasym:eq:P-report}
\end{equation}

\begin{theorem}[Exact quadratic-plus-periodic decomposition]
\label{dyadicasym:thm:periodic-decomposition}
The function $P$ is smooth and one-periodic.  For every $s>0$,
\begin{equation}
 \boxed{
 \Lam(s)=-\frac{(\log s)^2}{2L}+\frac12\log s
 +P(\log_2s)+\mathcal E(s).}                               \label{dyadicasym:eq:Lambda-periodic}
\end{equation}
All derivatives of $P$ are bounded.
\end{theorem}

\begin{proof}
The forward tail satisfies
\begin{equation*}
 T(2s)=T(s)-\log(1-\EulerE^{-s}).
\end{equation*}
Under $u\mapsto u+1$, the change in the quadratic term of \eqref{dyadicasym:eq:P-report} is
\begin{equation*}
 \frac L2\bigl((u+1)^2-(u+1)-(u^2-u)\bigr)=Lu=\log s.
\end{equation*}
This cancels the changes in \eqref{dyadicasym:eq:Lambda-dilation} and in $T$, proving
$P(u+1)=P(u)$.  Rearrangement gives \eqref{dyadicasym:eq:Lambda-periodic}.  Local uniform convergence of
the defining series permits differentiation; periodicity then makes every derivative bounded.
\end{proof}

The drift $-(\log s)^2/(2L)$ is the analytic counterpart of the triangular power
$2^{-\binom n2}$ in the exact endpoint formula.  The periodic remainder is the residual
freedom in solving the dilation equation on logarithmic scale.

\subsection{Mean, Fourier modes, and the asymptotic constant}

Let $\overline P$ and $\Psi$ be as in \eqref{dyadicasym:eq:Psi-intro}.  The Mellin transform of the
kernel makes the normalization explicit.  For $-1<\Re z<0$,
\begin{align}
 \widetilde\kappa(z)
 &=\int_0^\infty\kappa(u)u^{z-1}\Differential u
   =-\Gamma(z)\zeta(1+z),
 \label{dyadicasym:eq:kappa-mellin}\\
 \widetilde\Lam(z)
 &=\frac{2^z}{1-2^z}\bigl(-\Gamma(z)\zeta(1+z)\bigr).
 \label{dyadicasym:eq:Lambda-mellin}
\end{align}
The first identity follows by integration by parts and the subtracted Bose integral.  The
geometric factor is the Mellin image of the same dyadic scale that produced the Mersenne
factors in the exact formulas.  Its poles at $2\pi\ImaginaryUnit k/L$ give the Fourier modes; the
finite part at zero gives the mean
\begin{equation}
 \overline P=\frac{A_0}{L}-\frac L{12},                     \label{dyadicasym:eq:P-mean}
\end{equation}
where $A_0$ is defined in \eqref{dyadicasym:eq:CF-intro}.  With
\begin{equation}
 \chi_k=\frac{2\pi\ImaginaryUnit k}{L},                                \label{dyadicasym:eq:chi}
\end{equation}
the Fourier coefficients of the centered function are
\begin{equation}
 \FourierSeriesCoefficient{\Psi}{0}=0,
 \qquad
 \boxed{
 \FourierSeriesCoefficient{\Psi}{k}=-\frac{\Gamma(-\chi_k)\zeta(1-\chi_k)}L}
 \quad(k\ne0).                                             \label{dyadicasym:eq:Psi-Fourier-report}
\end{equation}
The Gamma factor has no zeros, and $\zeta$ has no zeros on $\Re z=1$ away from its pole at
$1$, so every nonzero Fourier coefficient is nonzero.  Thus $\Psi$ is genuinely nonconstant.
Stirling's vertical estimate for $\Gamma$ makes the coefficients exponentially small; in
fact the oscillation has amplitude only a few parts in $10^6$, but it cannot be deleted at an
error scale tending to zero.

The Gaussian normalization later contributes $-\tfrac12\log(2\pi)$.  Therefore
\begin{equation}
 \overline P-\frac12\log(2\pi)
 =\frac{A_0}{L}-\frac{7L}{12}-\frac12\log\pi
 =C_F.                                                       \label{dyadicasym:eq:CF-from-P}
\end{equation}
Numerically,
\begin{equation}
 \overline P=-1.10904536543418668\ldots,
 \qquad
 C_F=-2.02798389863885942\ldots.                            \label{dyadicasym:eq:constants-numeric}
\end{equation}

\section{Sharp route I: endpoint moments versus Laplace moments}
\label{dyadicasym:sec:endpoint-laplace}

This route begins with the exact rational object $d_n=\Expectation(1-X)^n$ and converts it into the
negative-Laplace product at the natural scale $s=n$.  It is the route closest to the current
Lean bridge modules.

\subsection{Second-order kernel comparison}

For $0\le x\le1$ and $n\ge1$, the logarithmic expansion
\begin{equation*}
 n\log(1-x)=-nx-\frac n2x^2+\BigO(nx^3)
\end{equation*}
is useful only after the large-$x$ region is controlled uniformly.  A convenient global
form is
\begin{equation}
 \left|(1-x)^n-\EulerE^{-nx}\left(1-\frac n2x^2\right)\right|
 \le16\EulerE^{-nx}(nx^3+n^2x^4).                               \label{dyadicasym:eq:kernel-comparison}
\end{equation}
Integrating against the law of $X$ gives
\begin{align}
 d_n={}&B(n)-\frac n2\Expectation\bigl(X^2\EulerE^{-nX}\bigr)+R_n,        \label{dyadicasym:eq:dn-B-second}\\
 |R_n|\le{}&16\left(
 n\Expectation(X^3\EulerE^{-nX})+n^2\Expectation(X^4\EulerE^{-nX})\right).              \label{dyadicasym:eq:Rn-bound}
\end{align}

Let $q(s)=\Lam(s)=\log B(s)$.  The normalized tilted second moment is
\begin{equation}
 \frac{\Expectation(X^2\EulerE^{-sX})}{B(s)}=q''(s)+q'(s)^2.               \label{dyadicasym:eq:tilted-second}
\end{equation}
The corresponding third- and fourth-moment identities are polynomials in
$q',q'',q^{(3)},q^{(4)}$.  Differentiation of \eqref{dyadicasym:eq:Lambda-periodic}, together with the
exponentially small tail, shows that these normalized moments have precisely the inverse
powers needed to make \eqref{dyadicasym:eq:Rn-bound} a relative $\BigO(n^{-1})$ error after the second
term is retained.

\begin{proposition}[Endpoint--Laplace logarithmic transfer]
\label{dyadicasym:prop:endpoint-laplace-transfer}
As $n\to\infty$,
\begin{equation}
 \boxed{
 \log d_n=q(n)-\frac n2\bigl(q''(n)+q'(n)^2\bigr)
 +\BigO(n^{-1}).}                                          \label{dyadicasym:eq:log-d-transfer}
\end{equation}
\end{proposition}

\begin{proofidea}
Divide \eqref{dyadicasym:eq:dn-B-second} by $B(n)$.  The second-order displacement is
\begin{equation*}
 \alpha_n=\frac n2\bigl(q''(n)+q'(n)^2\bigr)
          =\BigO(\log^2n/n),
\end{equation*}
and the normalized remainder is $\BigO(n^{-1})$.  Thus
$d_n/B(n)=1-\alpha_n+\BigO(n^{-1})$.  The local estimate
$\log(1-y)=-y+\BigO(y^2)$ applies; since
$\alpha_n^2=\BigO(\log^4n/n^2)=\BigO(n^{-1})$, it gives
\eqref{dyadicasym:eq:log-d-transfer}.
\end{proofidea}

\subsection{Differentiating the exact periodic decomposition}

Write
\begin{equation}
 \ell=\log s,
 \qquad
 u=\log_2s=\frac\ell L.                                    \label{dyadicasym:eq:ell-u}
\end{equation}
Ignoring an exponentially small differentiated tail, \eqref{dyadicasym:eq:Lambda-periodic} gives
\begin{equation}
 q'(s)=\frac{-\ell/L+1/2+\Psi'(u)/L}{s}+\BigO(\EulerE^{-s}),     \label{dyadicasym:eq:qprime}
\end{equation}
and
\begin{equation}
 q''(s)=\frac{\ell/L-1/2-1/L-\Psi'(u)/L+\Psi''(u)/L^2}{s^2}
 +\BigO(\EulerE^{-s}).                                           \label{dyadicasym:eq:qsecond}
\end{equation}
Squaring \eqref{dyadicasym:eq:qprime} and adding \eqref{dyadicasym:eq:qsecond}, the terms linear in $\ell$ that do
not involve $\Psi'$ cancel.  Since $\Psi'$ and $\Psi''$ are bounded,
\begin{equation}
 \frac s2\bigl(q''(s)+q'(s)^2\bigr)
 =\frac{\ell^2}{2L^2s}-
   \frac{\ell\,\Psi'(u)}{L^2s}+\BigO(s^{-1}).              \label{dyadicasym:eq:cumulant-expansion}
\end{equation}
This cancellation is the analytic reason that the first nonperiodic endpoint correction is
$-\log^2n/(2L^2n)$ rather than a mixture of $\log^2n/n$ and $\log n/n$ terms.

Substituting \eqref{dyadicasym:eq:Lambda-periodic} and \eqref{dyadicasym:eq:cumulant-expansion} into
\eqref{dyadicasym:eq:log-d-transfer} yields the first sharp form.

\begin{theorem}[Half moments at the logarithmic phase]
\label{dyadicasym:thm:dn-logphase}
Let $u_n=\log_2n$.  Then
\begin{align}
 \log d_n={}&-\frac{\log^2n}{2L}+\frac12\log n
 +\overline P+\Psi(u_n)                                    \notag\\
 &-\frac{\log^2n}{2L^2n}
 +\frac{\log n}{L^2n}\Psi'(u_n)+\BigO(n^{-1}).             \label{dyadicasym:eq:dn-logphase}
\end{align}
\end{theorem}

Combining \eqref{dyadicasym:eq:dn-logphase} with Stirling's formula and
\eqref{dyadicasym:eq:log-F-exact} gives
\begin{align}
 \log F(2^{-n})={}&-\frac L2n^2-n\log n+
 \left(1+\frac L2\right)n-\frac{\log^2n}{2L}+C_F           \notag\\
 &-\frac{\log^2n}{2L^2n}+\Psi(u_n)
 +\frac{\log n}{L^2n}\Psi'(u_n)+\BigO(n^{-1}).             \label{dyadicasym:eq:F-periodic-derivative}
\end{align}
Formula \eqref{dyadicasym:eq:F-periodic-derivative} has already recovered the constant and the periodic
function from the exact endpoint values.  The derivative term is not an extra oscillation; it
is the first-order displacement from the naive phase $u_n$ to the exact Lambert phase.

\section{The lower-Lambert phase and the derivative cancellation}
\label{dyadicasym:sec:lambert-phase}

\subsection{The exact phase}

The natural saddle coordinate at $x=2^{-n}$ is
\begin{equation}
 \lambda_n=-\frac1L W_{-1}(-L2^{-n}),                       \label{dyadicasym:eq:lambda-W}
\end{equation}
where $W_{-1}$ is the lower real branch of Lambert's function
\cite{dyadicasym:corless1996}.  It is the large solution of \eqref{dyadicasym:eq:lambda-n-intro}.  Its fixed-point
form is
\begin{equation}
 \lambda_n=n+\log_2\lambda_n.                               \label{dyadicasym:eq:lambda-fixed}
\end{equation}
The standard lower-Lambert expansion gives
\begin{align}
 \lambda_n={}&n+\frac{\log n}{L}+\frac{\log n}{L^2n}
 +\frac{\log n-\tfrac12\log^2n}{L^3n^2}                    \notag\\
 &+\frac{\log n-\tfrac32\log^2n+\tfrac13\log^3n}{L^4n^3}
 +\BigO\!\left(\frac{\log^4n}{n^4}\right).               \label{dyadicasym:eq:lambda-expansion}
\end{align}
For the periodic phase transfer only the first displacement is needed.  Define
\begin{equation}
 \delta_n=\log_2\lambda_n-\log_2n.                          \label{dyadicasym:eq:delta}
\end{equation}
Then
\begin{equation}
 \delta_n=\frac{\log n}{L^2n}
 +\BigO\!\left(\frac{\log^2n}{n^2}\right).                \label{dyadicasym:eq:delta-expansion}
\end{equation}

\subsection{Why \texorpdfstring{$\Psi$}{Psi} is evaluated at \texorpdfstring{$\lambda_n$}{lambda n}}

Since $n$ is an integer, periodicity and \eqref{dyadicasym:eq:lambda-fixed} imply the exact identity
\begin{equation}
 \Psi(\lambda_n)=\Psi(\lambda_n-n)
 =\Psi(\log_2\lambda_n).                                   \label{dyadicasym:eq:Psi-lambda-exact}
\end{equation}
The fixed-point step is the exact declaration
\path{dyadicLambertPhase_fixedPoint} in \path{FabiusLambertPhase.lean}; the periodic rewrite
uses \path{negativeLaplacePsi_periodic} in \path{PeriodicMean.lean}.
Taylor expansion at $u_n=\log_2n$, with uniformly bounded second derivative, gives
\begin{align}
 \Psi(\lambda_n)
 &=\Psi(u_n+\delta_n)                                      \notag\\
 &=\Psi(u_n)+\frac{\log n}{L^2n}\Psi'(u_n)
   +\BigO(n^{-1}).                                          \label{dyadicasym:eq:Psi-phase-transfer}
\end{align}
Indeed, the error in \eqref{dyadicasym:eq:delta-expansion} contributes
$\BigO(\log^2n/n^2)=\BigO(n^{-1})$, and the quadratic Taylor remainder contributes the same
or less.  Substitution of \eqref{dyadicasym:eq:Psi-phase-transfer} into
\eqref{dyadicasym:eq:F-periodic-derivative} proves \eqref{dyadicasym:eq:headline-expanded}.

This calculation explains exactly why the phase may not be replaced by $\log_2n$ at the
claimed precision.  One would omit the term
\begin{equation}
 \frac{\log n}{L^2n}\Psi'(\log_2n),                         \label{dyadicasym:eq:missing-phase-term}
\end{equation}
which is generally of order $(\log n)/n$, not $1/n$.  The tiny amplitude of $\Psi$ affects
numerical visibility, but not the asymptotic order.

\begin{warning}[Phase precision is amplified]
The leading exponent contains $-L\lambda^2/2$.  An error $\Delta\lambda$ in the phase changes
the exponent by approximately $-L\lambda\Delta\lambda$.  Therefore deriving an elementary
expansion with remainder $\BigO(n^{-1})$ requires the Lambert phase to one order beyond a
naive term count.  Keeping $\lambda_n$ exact inside $\Psi$ avoids this bookkeeping hazard.
\end{warning}

\subsection{Compact versus expanded form}

Define
\begin{equation}
 \mathcal M_n=-\frac L2\lambda_n^2+
 \left(1+\frac L2\right)\lambda_n-\frac12\log\lambda_n
 +C_F+\Psi(\lambda_n).                                     \label{dyadicasym:eq:Mn}
\end{equation}
Substitution of \eqref{dyadicasym:eq:lambda-expansion} and collection of terms gives
\begin{align}
 \mathcal M_n={}&-\frac L2n^2-n\log n+
 \left(1+\frac L2\right)n-\frac{\log^2n}{2L}+C_F           \notag\\
 &-\frac{\log^2n}{2L^2n}+\Psi(\lambda_n)+\BigO(n^{-1}).   \label{dyadicasym:eq:compact-to-expanded}
\end{align}
Thus \eqref{dyadicasym:eq:headline-compact} and \eqref{dyadicasym:eq:headline-expanded} are equivalent at the stated
precision.  The compact formula is structurally cleaner and is the one that survives
unchanged for arbitrary small real arguments.

\section{Sharp route II: a coefficient saddle from the exact generating product}
\label{dyadicasym:sec:coefficient-saddle}

The endpoint--Laplace route starts from $d_n=\Expectation(1-X)^n$.  There is a second route beginning
directly with the recurrence coefficient $a_n=[z^n]A(z)$.  It is particularly well suited to
the question posed here because $a_n$ is exactly what remains after the $q$-binomial formula
is decompressed.

\subsection{Positive-axis growth of \texorpdfstring{$A$}{A}}

The reflection identity \eqref{dyadicasym:eq:A-reflection} and
\Cref{dyadicasym:thm:periodic-decomposition} give, for $r>0$,
\begin{align}
 K(r):=\log A(r)
 ={}&r+\Lam(r)                                               \notag\\
 ={}&r-\frac{\log^2r}{2L}+\frac12\log r
 +P(\log_2r)+\mathcal E(r).                                \label{dyadicasym:eq:K-positive}
\end{align}
The reference saddle is chosen by ignoring the bounded derivative of the periodic term:
\begin{equation}
 r-\frac{\log r}{L}=n.                                      \label{dyadicasym:eq:reference-saddle}
\end{equation}
Its large solution is exactly $r=\lambda_n$.  The exact coefficient saddle for
$rK'(r)=n$ differs from $\lambda_n$ by only $\BigO(1)$, because
\begin{equation}
 rK'(r)=r-\frac{\log r}{L}+\frac12+\frac1L P'(\log_2r)
 +\BigO(r\EulerE^{-r}).                                          \label{dyadicasym:eq:exact-saddle-shift}
\end{equation}
An $\BigO(1)$ displacement of a saddle of curvature $\asymp1/r$ changes the logarithmic
coefficient estimate by only $\BigO(1/r)$.

\subsection{Vertical extraction and the normalized saddle mass}

The probability representation gives a stronger contour than the circular sketch.  For
$n\ge1$ and every $c>0$,
\begin{equation}
 a_n=\frac1{2\pi\ImaginaryUnit}\int_{c-\ImaginaryUnit\infty}^{c+\ImaginaryUnit\infty}
 \frac{A(z)}{z^{n+1}}\Differential z.                               \label{dyadicasym:eq:vertical-extraction}
\end{equation}
The integral is absolutely convergent because
$|A(c+\ImaginaryUnit t)|\le A(c)$ and $|c+\ImaginaryUnit t|^{-n-1}$ is integrable.  Fubini and the
inverse-Laplace kernel recover $\Expectation X^n/n!=a_n$, so this is an exact consequence of the
moment law rather than an inversion assumption for $F$.

Set $c=\lambda_n$ and write $z=\lambda_n(1+\ImaginaryUnit\theta)$.  Define
\begin{equation}
 \mathcal C_n=\sqrt{\frac{\lambda_n}{2\pi}}
 \int_{-\infty}^{\infty}
 \frac{A(\lambda_n(1+\ImaginaryUnit\theta))}{A(\lambda_n)}
 (1+\ImaginaryUnit\theta)^{-n-1}\Differential\theta.                           \label{dyadicasym:eq:Cn-def}
\end{equation}
Then the exact coefficient identity is
\begin{equation}
 F(2^{-n})=
 \frac{2^{-\binom n2}A(\lambda_n)\lambda_n^{-n}}
 {\sqrt{2\pi\lambda_n}}\,\mathcal C_n.                   \label{dyadicasym:eq:exact-coefficient-saddle}
\end{equation}
Substituting the periodic decomposition and using
$\log_2\lambda_n=\lambda_n-n$ gives the exact separation
\begin{equation}
 \boxed{\log F(2^{-n})-\mathcal M_n
 =\mathcal E(\lambda_n)+\log\mathcal C_n.}                \label{dyadicasym:eq:exact-dyadic-separation}
\end{equation}
Thus all algebraic normalization, the constant, and the periodic phase have been removed
before any approximation is made.

\subsection{Local phase, domination, and all orders}

Put $u=\PrincipalLogarithm(1+\ImaginaryUnit\theta)$.  On a fixed central neighborhood, the exponent in
\eqref{dyadicasym:eq:Cn-def} satisfies
\begin{align}
 \Phi_{\lambda}^{\mathrm{coef}}(\theta)
 :={}&\log\frac{A(\lambda(1+\ImaginaryUnit\theta))}{A(\lambda)}-(n+1)u\notag\\
 ={}&\lambda(\ImaginaryUnit\theta-u)-\frac u2-\frac{u^2}{2L}
 +\Psi\!\left(\lambda+\frac uL\right)-\Psi(\lambda)
 +\mathcal E_{\lambda}^{\mathrm{coef}}(\theta),           \label{dyadicasym:eq:coefficient-local-phase}
\end{align}
where every fixed derivative of the last term is $O(\EulerE^{-\lambda})$.  This is the same
scaled local phase as in the continuous Bromwich analysis.  Globally,
\begin{equation}
 \left|\frac{A(\lambda(1+\ImaginaryUnit\theta))}{A(\lambda)}
 (1+\ImaginaryUnit\theta)^{-n-1}\right|
 \le(1+\theta^2)^{-(n+1)/2},                              \label{dyadicasym:eq:coefficient-tail-bound}
\end{equation}
so the vertical contour supplies its own Gaussian central bound and integrable tails.

With $\theta=v/\sqrt\lambda$, the phase begins with $-v^2/2$; the
$\lambda^{-1/2}$ contribution is odd and integrates to zero.  Gaussian contraction gives
\begin{equation}
 \log\mathcal C_n
 =\frac{A_1(\lambda_n)}{\lambda_n}+O(\lambda_n^{-2}),      \label{dyadicasym:eq:Cn-first}
\end{equation}
where $A_1$ is displayed in \eqref{dyadicasym:eq:A1-report}.  More generally, expanding
\eqref{dyadicasym:eq:coefficient-local-phase}, exponentiating, and integrating its even
Gaussian moments yields, for every fixed $N$,
\begin{equation}
 \log\mathcal C_n
 =\sum_{j=1}^{N-1}\frac{A_j(\lambda_n)}{\lambda_n^j}
 +O_N(\lambda_n^{-N}).                                    \label{dyadicasym:eq:Cn-all-orders}
\end{equation}
The coefficient functions are the same periodic differential polynomials as in the
continuous saddle because the local phases agree.  In particular,
\begin{equation}
 \log F(2^{-n})=
 -\frac L2\lambda_n^2+
 \left(1+\frac L2\right)\lambda_n-\frac12\log\lambda_n
 +C_F+\Psi(\lambda_n)+O(n^{-1}),                           \label{dyadicasym:eq:saddle-final}
\end{equation}
which is \eqref{dyadicasym:eq:headline-compact}.

\begin{boxedremark}
The coefficient saddle provides a particularly direct answer to the original question.
Starting from the exact $q$-binomial value, \eqref{dyadicasym:eq:N-m-a} isolates $a_n$; coefficient
extraction of that coefficient from the exact product produces the lower-Lambert asymptotic.
The endpoint--Laplace route and the coefficient route are not the same proof in disguise:
one compares $(1-X)^n$ with $\EulerE^{-nX}$ on the real line, while the other analyzes a complex
coefficient integral.  Their agreement fixes the normalization and phase very rigidly.
\end{boxedremark}

\section{Fixed and moderately growing dyadic numerators}
\label{dyadicasym:sec:fixed-numerator}

The master coefficient extends Route~II from reciprocal powers to rays $x=m2^{-n}$.
Factor
\begin{equation}
 B_m(z)=\EulerE^{(m-1)z}\mathcal S_m(z),\qquad
 \mathcal S_m(z)=\sum_{h=0}^{m-1}\ThueMorseSignSymbol_h\EulerE^{-hz}.
 \label{dyadicasym:eq:fixed-prefix-factor}
\end{equation}
For $\Re z=r>0$,
$|\mathcal S_m(z)-1|\le\EulerE^{-r}/(1-\EulerE^{-r})$; thus a fixed prefix is dominated
uniformly on vertical lines by its first exponential.

Let $m\ge1$ be fixed, put $x=m2^{-n}$, let $\lambda=\lambda(x)$ satisfy
$x=\lambda2^{-\lambda}$, and set $r=\lambda/m$.  Then
\begin{equation}
 n=mr-\log_2r,\qquad \log_2r=\lambda-n,
 \qquad \Psi(\log_2r)=\Psi(\lambda).
 \label{dyadicasym:eq:fixed-ray-phase-lock}
\end{equation}
Define
\begin{equation}
 \mathcal M(x)=-\frac L2\lambda^2+
 \left(1+\frac L2\right)\lambda-\frac12\log\lambda
 +C_F+\Psi(\lambda).
 \label{dyadicasym:eq:general-main}
\end{equation}
For fixed \(m\), the estimate above gives \(\mathcal S_m(r)>0\) once \(n\) is
sufficiently large, so all real logarithms below are defined.  Vertical extraction of
\([z^n](B_mA)\) at \(z=r(1+\ImaginaryUnit\theta)\) gives the normalized mass
\begin{equation}
 \mathcal C_{m,n}
 =\sqrt{\frac{\lambda}{2\pi}}
 \int_{-\infty}^{\infty}
 \frac{B_m(r(1+\ImaginaryUnit\theta))A(r(1+\ImaginaryUnit\theta))}
      {B_m(r)A(r)}
 (1+\ImaginaryUnit\theta)^{-n-1}\Differential\theta.
 \label{dyadicasym:eq:Cmn-def}
\end{equation}
The coefficient identity identifies \(\mathcal C_{m,n}\) with a positive normalization of
\(F(m2^{-n})\), so it is positive in the same eventual range.  Taking real logarithms then
gives the exact separation
\begin{equation}
 \boxed{
 \log F(m2^{-n})-\mathcal M(m2^{-n})
 =\mathcal E(r)+\log\mathcal S_m(r)+\log\mathcal C_{m,n}.}
 \label{dyadicasym:eq:fixed-ray-separation}
\end{equation}
On the central window, the prefix estimate reduces the exponent to
\eqref{dyadicasym:eq:coefficient-local-phase} with an error whose fixed derivatives are
$O_m(\EulerE^{-\lambda/m})$.  The same vertical-tail bound therefore yields
\begin{equation}
 \log\mathcal C_{m,n}
 =\sum_{j=1}^{N-1}\frac{A_j(\lambda)}{\lambda^j}
 +O_{m,N}(\lambda^{-N}).
 \label{dyadicasym:eq:fixed-ray-all-orders}
\end{equation}
Hence the full coefficient expansion uses the same periodic $A_j$ on every fixed-numerator
ray.

The proof remains uniform for a growing numerator $m=m(n)$ whenever
\begin{equation}
 \frac{\lambda/m}{\log\lambda}\longrightarrow\infty,
 \qquad\text{equivalently }m=o(\lambda/\log\lambda).
 \label{dyadicasym:eq:moderate-numerator-condition}
\end{equation}
Indeed, each fixed family of prefix derivatives is
$O_N(r^{K_N}\EulerE^{-r})$, smaller than every required power of $1/\lambda$ in this
regime.  This does not reach the large numerators needed for sharp interpolation of every
real argument; that two-parameter boundary is described below.

\section{Exact Bromwich-to-coefficient collapse}
\label{dyadicasym:sec:bromwich-collapse}

The continuous endpoint integral and the coefficient integral are connected before any
asymptotic approximation.  Iterating
\begin{equation}
 A(-2z)=\frac{1-\EulerE^{-z}}zA(-z)
 \label{dyadicasym:eq:negative-refinement}
\end{equation}
gives
\begin{equation}
 A(-2^nz)=2^{-\binom n2}z^{-n}A(-z)
 \prod_{j=0}^{n-1}(1-\EulerE^{-2^jz}).
 \label{dyadicasym:eq:negative-refinement-iterated}
\end{equation}
The finite product is the complete Thue--Morse block
\begin{equation}
 P_n(z)=\prod_{j=0}^{n-1}(1-\EulerE^{-2^jz})
 =\sum_{h=0}^{2^n-1}\ThueMorseSignSymbol_h\EulerE^{-hz}.
 \label{dyadicasym:eq:Bromwich-TM-product}
\end{equation}
Starting from the Laplace--Stieltjes inversion formula for $F$, putting
$x=2^{-n}$, substituting $s=2^nz$, and using $A(z)=\EulerE^zA(-z)$ yields
\begin{equation}
 F(2^{-n})=\frac{2^{-\binom n2}}{2\pi\ImaginaryUnit}
 \int\frac{A(z)P_n(z)}{z^{n+1}}\Differential z.
 \label{dyadicasym:eq:Bromwich-rescaled}
\end{equation}
With $A(z)=\Expectation\EulerE^{zX}$, inverse Laplace transforms this exactly into
\begin{equation}
 F(2^{-n})=2^{-\binom n2}\Expectation\left[
 \sum_{h=0}^{2^n-1}\ThueMorseSignSymbol_h\frac{(X-h)_+^n}{n!}\right].
 \label{dyadicasym:eq:positive-part-collapse}
\end{equation}
Because $0\le X\le1$, every $h\ge1$ term vanishes almost surely and the remaining term is
$2^{-\binom n2}a_n$.  On the saddle line $\Re z=\lambda_n$ one also has
$P_n(z)=1+O(\EulerE^{-\lambda_n})$, while the continuous radius
$2^n\lambda_n=2^{\lambda_n}$ maps exactly to the coefficient radius $\lambda_n$.
Thus the rescaled continuous integrand and the coefficient integrand differ only by a
beyond-all-orders finite block.

\section{Consequences for every exact arithmetic representation}
\label{dyadicasym:sec:corollaries}

\subsection{The normalized recurrence and composition sum}

Since $a_n=2^{\binom n2}F(2^{-n})$, \eqref{dyadicasym:eq:headline-expanded} immediately gives
\begin{align}
 \log a_n={}&-n\log n+n-\frac{\log^2n}{2L}+C_F             \notag\\
 &-\frac{\log^2n}{2L^2n}+\Psi(\lambda_n)+\BigO(n^{-1}).   \label{dyadicasym:eq:an-asymptotic}
\end{align}
Thus the exact composition sum \eqref{dyadicweb:eq:ordered-composition} has the asymptotic
\begin{align}
 &\sum_{\rho\vDash n}
 \prod_{j=1}^{\ell(\rho)}
 \frac{1}{(2^{s_j}-1)(r_j+1)!}                             \notag\\
 &\quad=\exp\!\left(
 -n\log n+n-\frac{\log^2n}{2L}+C_F
 -\frac{\log^2n}{2L^2n}+\Psi(\lambda_n)
 +\BigO(n^{-1})\right).                                    \label{dyadicasym:eq:composition-asymptotic}
\end{align}
The recurrence therefore has a moving periodic multiplier; no ansatz with a single fixed
constant can be sharp.

\subsection{The half moments}

Adding Stirling's expansion to \eqref{dyadicasym:eq:an-asymptotic},
\begin{align}
 \log d_n={}&\frac12\log(2\pi n)+C_F+\Psi(\lambda_n)
 -\frac{\log^2n}{2L}                                      \notag\\
 &-\frac{\log^2n}{2L^2n}+\BigO(n^{-1}).                   \label{dyadicasym:eq:dn-asymptotic}
\end{align}
In particular,
\begin{equation}
 d_n\sim\sqrt{2\pi n}\,
 \exp\!\left(C_F+\Psi(\lambda_n)-\frac{\log^2n}{2L}\right).
                                                                    \label{dyadicasym:eq:dn-equivalent}
\end{equation}
The notation ``$\sim$'' is legitimate because the omitted logarithmic error tends to zero.
The periodic factor $\exp\Psi(\lambda_n)$ remains part of the equivalent.

\subsection{The literal \texorpdfstring{$q$}{q}-binomial numerator}

By \eqref{dyadicasym:eq:N-m-a},
\begin{equation}
 \mathcal N_n=\M_na_n.                                      \label{dyadicasym:eq:N-product}
\end{equation}
Furthermore,
\begin{equation}
 \M_n=2^{\binom{n+1}{2}}(1/2;1/2)_n                       \label{dyadicasym:eq:M-q}
\end{equation}
and
\begin{equation}
 \log(1/2;1/2)_n=
 \log(1/2;1/2)_\infty+\BigO(2^{-n}).                       \label{dyadicasym:eq:qpoch-limit}
\end{equation}
Consequently
\begin{align}
 \log\mathcal N_n={}&\frac L2n^2-n\log n+
 \left(1+\frac L2\right)n-\frac{\log^2n}{2L}              \notag\\
 &+C_F+\log(1/2;1/2)_\infty
 -\frac{\log^2n}{2L^2n}+\Psi(\lambda_n)
 +\BigO(n^{-1}).                                            \label{dyadicasym:eq:N-asymptotic}
\end{align}
This formula is the sharp asymptotic of the finite signed $q$-binomial/Thue--Morse sum itself.
It makes the scale cancellation visible: the numerator grows like
$\exp(Ln^2/2)$, while the explicit denominator in \eqref{dyadicweb:eq:input-q-dyadic} contains
$2^{n^2}$.  The residual triangular factor is exactly the decay
$2^{-\binom n2}$ in \eqref{dyadicasym:eq:F-an}.

\subsection{A first all-orders corollary}

The general small-argument saddle expansion supplies a first correction
\begin{equation}
 A_1(u)=\frac1{24}+\frac1{2L}
 -\frac{\Psi''(u)+\Psi'(u)^2}{2L^2}.                        \label{dyadicasym:eq:A1-report}
\end{equation}
At reciprocal powers this gives
\begin{equation}
 \log F(2^{-n})=\mathcal M_n+
 \frac{A_1(\lambda_n)}{\lambda_n}+
 \BigO(\lambda_n^{-2}).                                    \label{dyadicasym:eq:first-correction-dyadic}
\end{equation}
The same correction can be transported to $a_n$, $d_n$, the composition sum, and
$\mathcal N_n$ by the exact multiplicative normalizations above.  Higher orders have the same
structure: a Poincar\'e series in inverse powers of $\lambda_n$ with smooth periodic
coefficient functions evaluated at the exact phase.

\section{What each exact formula contributes asymptotically}
\label{dyadicasym:sec:formula-audit}

It is useful to separate exact equivalence from asymptotic convenience.  The formulae compute
the same numbers, but they expose different structures.

\begingroup\small
\begin{longtable}{@{}
  >{\raggedright\arraybackslash}p{0.25\textwidth}
  >{\raggedright\arraybackslash}p{0.31\textwidth}
  >{\raggedright\arraybackslash}p{0.36\textwidth}@{}}
\toprule
Exact representation & What it exposes & Best asymptotic use\\
\midrule
\endhead
Endpoint identity
$F(2^{-n})=d_n/(n!2^{\binom n2})$ &
The triangular dyadic factor, factorial drift, and one unknown half moment &
Immediately proves the coefficient $-L/2$ of $n^2$; becomes sharp after $d_n$ is compared
with $A(-n)$\\
\addlinespace
Recurrence for $a_n$ &
A positive triangular coefficient recursion &
Reassemble into $A(2z)=E(z)A(z)$; direct iteration term by term obscures the moving saddle\\
\addlinespace
Composition sum &
A positive sum over all ordered descents of the recurrence &
Identifies $a_n$ exactly; asymptotics are best obtained after resummation into $A$, not by
selecting a few compositions\\
\addlinespace
Entire product $A(z)=\prod E(z/2^j)$ &
The dyadic scale at every analytic level &
Primary source for both the negative-Laplace harmonic sum and the coefficient saddle\\
\addlinespace
Thue--Morse block &
A zero of exact order $k$ and a finite product over dyadic scales &
Converts the inner signed power sum into a coefficient of $A(-2^kX)/A(-X)$\\
\addlinespace
Half-base $q$-binomial row &
A divided difference on $-1,-2,\ldots,-2^n$ &
Extracts the leading coefficient $a_n$ after the block cancellation; termwise asymptotics
are badly conditioned\\
\addlinespace
$q$-Pochhammer normalization &
The Mersenne product $\M_n$ in normalized form &
Contributes the elementary limit
$(1/2;1/2)_n\to(1/2;1/2)_\infty$ and cancels $\M_n$ exactly\\
\addlinespace
Denominator and valuation formulae &
Prime-adic arithmetic of the exact rationals &
Strong arithmetic information, but by themselves insufficient to determine Archimedean
magnitude or the periodic saddle phase\\
\bottomrule
\end{longtable}
\endgroup

\subsection{Why direct termwise analysis of the \texorpdfstring{$q$}{q}-sum is misleading}

The formula
\begin{equation*}
 F(2^{-n})=\frac{\mathcal N_n}{2^{n^2}(1/2;1/2)_n}
\end{equation*}
contains an explicit logarithmic contribution $-Ln^2$, whereas the final leading term is only
$-Ln^2/2$.  The missing $+Ln^2/2$ is generated inside the alternating finite numerator.
There is no contradiction: the outer Gaussian-binomial row is designed to annihilate all
lower scale powers exactly, and the surviving leading coefficient is exponentially larger
than a naive bounded-summand estimate would suggest.

The cancellation can be quantified without inspecting any individual summand.  From
\eqref{dyadicasym:eq:N-asymptotic},
\begin{equation}
 \log\mathcal N_n=\frac L2n^2+\BigO(n\log n),               \label{dyadicasym:eq:N-leading}
\end{equation}
while
\begin{equation}
 -\log\bigl(2^{n^2}(1/2;1/2)_n\bigr)
 =-Ln^2-\log(1/2;1/2)_\infty+\LittleO(1).                   \label{dyadicasym:eq:q-den-leading}
\end{equation}
Their sum is $-Ln^2/2+\BigO(n\log n)$, as required.

\begin{remark}[Numerical conditioning]
A floating-point implementation of the literal signed formula is normally inferior to the
positive recurrence.  It must resolve the two exact annihilations before dividing by a factor
of size $2^{n^2}$.  The decompressed identity $\mathcal N_n=\M_na_n$ is therefore not only a
proof device; it is the stable evaluation strategy.
\end{remark}

\section{How much of the full small-argument asymptotic is encoded by inverse dyadics?}
\label{dyadicasym:sec:general-x}

\subsection{The arbitrary-argument formula}

For $x\downarrow0$, let
\begin{equation}
 \lambda(x)=-\frac1L W_{-1}(-Lx),
 \qquad
 \lambda(x)2^{-\lambda(x)}=x.                              \label{dyadicasym:eq:lambda-x}
\end{equation}
The full endpoint saddle has the same compact main term:
\begin{equation}
 \log F(x)=
 -\frac L2\lambda(x)^2+
 \left(1+\frac L2\right)\lambda(x)-\frac12\log\lambda(x)
 +C_F+\Psi(\lambda(x))+\BigO(\lambda(x)^{-1}).              \label{dyadicasym:eq:general-x-main}
\end{equation}
At $x=2^{-n}$ this becomes \eqref{dyadicasym:eq:headline-compact}.  Thus the exact dyadic arithmetic
recovers the restriction of the entire arbitrary-argument asymptotic, including the same
constant and the same periodic function at the same exact phase.

\subsection{Dense sampling of the periodic correction}

The sequence of inverse dyadics does not freeze the periodic phase.  Put
\begin{equation}
 y_n=\lambda_n-n=\log_2\lambda_n.                           \label{dyadicasym:eq:yn}
\end{equation}
Then $y_n$ is increasing and unbounded, while
\begin{equation}
 y_{n+1}-y_n\longrightarrow0.                               \label{dyadicasym:eq:yn-increments}
\end{equation}
Indeed, $\lambda_n\sim n$ and
$y_{n+1}-y_n=\log_2(\lambda_{n+1}/\lambda_n)=\BigO(n^{-1})$.

\begin{lemma}[Density of the dyadic Lambert phases]
\label{dyadicasym:lem:phase-density}
The fractional parts $\{\lambda_n\}=\{y_n\}$ are dense in $[0,1]$.
\end{lemma}

\begin{proof}
Fix $\theta\in[0,1]$.  For every sufficiently large integer $m$, let $n(m)$ be the first
index for which $y_{n(m)}\ge m+\theta$.  Then
\begin{equation*}
 0\le y_{n(m)}-(m+\theta)
 <y_{n(m)}-y_{n(m)-1}.
\end{equation*}
The right side tends to zero by \eqref{dyadicasym:eq:yn-increments}.  Hence the fractional parts of
$y_{n(m)}$ tend to $\theta$.
\end{proof}

After removing the nonperiodic terms in \eqref{dyadicasym:eq:headline-expanded}, the exact rational
sequence therefore determines $\Psi$ uniquely.  More precisely, if
\begin{align}
 R_n={}&\log F(2^{-n})+\frac L2n^2+n\log n-
 \left(1+\frac L2\right)n+\frac{\log^2n}{2L}               \notag\\
 &-C_F+\frac{\log^2n}{2L^2n},                              \label{dyadicasym:eq:renormalized-Rn}
\end{align}
then
\begin{equation}
 R_n=\Psi(\lambda_n)+\BigO(n^{-1}).                         \label{dyadicasym:eq:Rn-Psi}
\end{equation}
For any $u\in[0,1]$, choose a subsequence with
$\{\lambda_{n_j}\}\to u$; continuity gives $R_{n_j}\to\Psi(u)$.  In this precise sense the
inverse-power values encode the whole periodic correction, not just a single phase sample.

\subsection{What inverse powers alone do not prove}

The inverse-power sequence does recover the universal leading quadratic law.  If
$x=2^{-t}$ and $n\le t<n+1$, monotonicity gives
$F(2^{-(n+1)})\le F(x)\le F(2^{-n})$; dividing the coarse endpoint estimates by
$t^2$ yields
\begin{equation}
 \log F(2^{-t})\sim-\frac L2t^2.
 \label{dyadicasym:eq:monotone-leading-law}
\end{equation}
It does not, however, extend \eqref{dyadicasym:eq:headline-compact} uniformly to all
$x\downarrow0$.  The same bracket places a point
$2^{-(n+1)}\le x\le2^{-n}$ between two endpoint values whose logarithms differ by order $n$:
\begin{equation}
 \log F(2^{-n})-\log F(2^{-(n+1)})=Ln+\BigO(\log n).        \label{dyadicasym:eq:dyadic-gap}
\end{equation}
That uncertainty is enormous compared with the desired $\BigO(n^{-1})$ remainder.

The global formula \eqref{dyadicweb:eq:input-q-dyadic} contains every dyadic rational and could, in principle,
be used on meshes much finer than one reciprocal-power interval.  To derive
\eqref{dyadicasym:eq:general-x-main} by that route one would need a two-parameter asymptotic uniform in
both numerator and denominator.  The arbitrary-argument Bromwich saddle avoids that difficult
finite-sum uniformity.  Accordingly, the exact inverse-power formulae recover the full phase
function and all constants, but a genuinely uniform analytic argument remains necessary to
pass from the dense arithmetic set to a sharp theorem on a continuum.

More concretely, for $x_n=m_n2^{-n}$ the master coefficient contains the prefix factor
\begin{equation}
 \mathcal S_{m_n}(r_n)=\sum_{h<m_n}\ThueMorseSignSymbol_h\EulerE^{-hr_n},
 \qquad r_n=\lambda(x_n)/m_n.                              \label{dyadicasym:eq:large-prefix-problem}
\end{equation}
The fixed/moderate analysis requires $r_n$ large, whereas a mesh fine enough to control
an additive $o(1)$ error in $\log F$ needs $m_n\gg\lambda$.  In that latter regime the
prefix cannot be replaced by one and depends on a long stretch of numerator digits.  A
uniform theorem there is a two-parameter de-Poissonization of the exact dyadic formula,
essentially the continuous Bromwich problem in another coordinate system.

\section{Numerical verification from exact rational values}
\label{dyadicasym:sec:numerics}

The recurrence \eqref{dyadicweb:eq:d-recurrence-companion} computes $a_n=d_n/n!$, hence $F(2^{-n})$, in exact rational
arithmetic.  For example,
\begin{equation}
 F(2^{-10})=
 \frac{134926369}{1246394851358539387238350848000},         \label{dyadicasym:eq:F10}
\end{equation}
so
\begin{equation*}
 \log F(2^{-10})=-50.5775682823421608125\ldots.
\end{equation*}

Let
\begin{equation}
 \rho_n=\log F(2^{-n})-\mathcal M_n.                        \label{dyadicasym:eq:rho}
\end{equation}
The leading theorem predicts $\rho_n=\BigO(1/\lambda_n)$, and the first correction predicts
$\lambda_n\rho_n=A_1(\lambda_n)+\BigO(1/\lambda_n)$.  The table was computed by exact rational
recurrence followed by one high-precision logarithm; $P$ and its derivatives were evaluated
from the absolutely convergent product and forward-tail series.

\begin{table}[htbp]
\centering
\small
\begin{tabular}{@{}rrrrrr@{}}
\toprule
$n$ & $\lambda_n$ & $\rho_n$ & $\lambda_n\rho_n$
& $A_1(\lambda_n)$ & $\lambda_n^2(\rho_n-A_1/\lambda_n)$\\
\midrule
10  & 13.78503056  & 0.05801333439 & 0.7997155875 & 0.7629678468 & 0.5065687288\\
20  & 24.62186833  & 0.03182756966 & 0.7836542295 & 0.7629268731 & 0.5103462408\\
40  & 45.50804986  & 0.01701388954 & 0.7742689337 & 0.7629490545 & 0.5151456258\\
80  & 86.43351899  & 0.008896709914& 0.7689739453 & 0.7629823191 & 0.5178773369\\
160 & 167.3870441   & 0.004576872612& 0.7661091776 & 0.7630070725 & 0.5192522062\\
\bottomrule
\end{tabular}
\caption{Exact reciprocal-power values against the compact Lambert main and first correction.}
\label{dyadicasym:tab:numerics}
\end{table}

The fourth and fifth columns approach one another, while the last column remains bounded and
appears to settle, exactly as \eqref{dyadicasym:eq:first-correction-dyadic} predicts.

The fixed-ray coefficient theorem has an independent exact-rational check.  With
$\rho_{m,n}=\log F(m2^{-n})-\mathcal M(m2^{-n})$, the following rows retain the
nontrivial numerators from the original bridge calculation.
\begin{table}[htbp]
\centering
\small
\begin{tabular}{@{}rrrrr@{}}
\toprule
$m$ & $n$ & $\lambda$ & $\lambda\rho_{m,n}$ &
$\lambda^2(\rho_{m,n}-A_1(\lambda)/\lambda)$\\
\midrule
3 & 20 & 22.93448405 & 0.7855611616 & 0.5164\\
3 & 40 & 43.87020711 & 0.7748252957 & 0.5182\\
3 & 80 & 84.82139378 & 0.7691025485 & 0.5188\\
5 & 20 & 22.14711904 & 0.7864070706 & 0.5162\\
5 & 40 & 43.10795409 & 0.7751701258 & 0.5203\\
5 & 80 & 84.07161885 & 0.7692992153 & 0.5214\\
\bottomrule
\end{tabular}
\caption{Exact coefficient checks on the fixed rays $m=3,5$.}
\label{dyadicasym:tab:fixed-ray-numerics}
\end{table}

A second check displays the cancellation in the literal $q$-formula.  Since
\begin{equation*}
 \log F(2^{-n})=\log\mathcal N_n-n^2L-\log(1/2;1/2)_n,
\end{equation*}
one obtains:

\begin{table}[htbp]
\centering
\small
\begin{tabular}{@{}rrrrr@{}}
\toprule
$n$ & $\log\mathcal N_n$ & $n^2L$ & $\log(1/2;1/2)_n$ & $\log F(2^{-n})$\\
\midrule
10 & 17.4960644003 & 69.3147180560 & -1.2410853733 & -50.5775682823\\
20 & 95.4684669707 & 277.2588722240 & -1.2420611411 & -180.5483441121\\
40 & 447.4055728988& 1109.0354888959& -1.2420620948 & -660.3878539023\\
80 & 1957.8744511397&4436.1419555836& -1.2420620948 & -2477.0254423491\\
\bottomrule
\end{tabular}
\caption{Logarithmic cancellation in the finite $q$-binomial formula.}
\label{dyadicasym:tab:q-cancellation}
\end{table}

The Pochhammer logarithm has already converged to
\begin{equation}
 \log(1/2;1/2)_\infty=-1.2420620948124149458\ldots.         \label{dyadicasym:eq:qpoch-numeric}
\end{equation}
The nontrivial compensation is between $\log\mathcal N_n\sim Ln^2/2$ and the explicit
$-Ln^2$ denominator.

\clearpage
\section{Correspondence with the Lean development}
\label{dyadicasym:sec:lean-map}

The relevant modules in the 25 August 2026 repository snapshot \cite{dyadicasym:proveit2026}
form a fairly literal proof graph.  The table records the role of each group rather than every import edge.

\begingroup\small
\begin{longtable}{@{}
  >{\raggedright\arraybackslash}p{0.33\textwidth}
  >{\raggedright\arraybackslash}p{0.59\textwidth}@{}}
\toprule
Module & Role in the bridge\\
\midrule
\endhead
\path{FabiusRecurrenceSequence.lean} &
Defines $a_n=d_n/n!$, proves its recurrence, identifies
$F(2^{-n})=2^{-\binom n2}a_n$, and connects its ordinary series to the analytic generating
function and dyadic product.\\
\addlinespace
\path{FabiusInverseDyadicClosedForm.lean} &
Unfolds the inverse-dyadic recurrence into the exact composition formula.\\
\addlinespace
\path{FabiusRawQBinomialFormula.lean},
\path{FabiusQBinomialFormula.lean},
\path{FabiusQBinomialScalarFormula.lean},
\path{FabiusDyadicQBinomialScalar.lean} &
Develop the Thue--Morse block factorization, half-base Gaussian interpolation, scalar
translation invariance, and the global exact dyadic formula.  The coefficient-extraction
argument in \cref{dyadicweb:sec:q-extraction} is the ordinary-mathematics unpacking of this layer.\\
\addlinespace
\path{FabiusDyadicLogBounds.lean} &
Combines the exact endpoint identity with elementary half-moment bounds to prove the
$-Ln^2/2$ law and an explicit $3n\log(n+1)$ error.\\
\addlinespace
\path{NegativeLaplace.lean},
\path{PeriodicCorrection.lean},
\path{PeriodicFourier.lean} &
Construct the negative-Laplace logarithm, its exact quadratic-plus-periodic decomposition,
the forward-tail bound, regularity, mean normalization, and Fourier data.\\
\addlinespace
\path{EndpointLaplaceComparison.lean} &
Proves the second-order pointwise and integrated comparison between $(1-x)^n$ and
$\EulerE^{-nx}(1-nx^2/2)$ and transfers the estimate through the logarithm.\\
\addlinespace
\path{DyadicSharpDecomposition.lean},
\path{FabiusDyadicSharpCumulant.lean} &
Assemble the exact logarithmic decomposition of $F(2^{-n})$, including Stirling, the periodic
term, the endpoint/Laplace error, and the exact second tilted cumulant.\\
\addlinespace
\path{LaplacePeriodicSecondOrder.lean} &
Differentiates the periodic Laplace decomposition and proves the explicit
$\Psi'(\log_2n)$ correction with an $\BigO(1/n)$ remainder.\\
\addlinespace
\path{FabiusLambertPhase.lean},
\path{FabiusDyadicSharpAsymptotic.lean} &
Construct the exact lower-Lambert phase, transfer the periodic derivative term to
$\Psi(\lambda_n)$, and prove the corrected sharp dyadic theorem.\\
\addlinespace
\path{FabiusSharpLambertMain.lean},
\path{FabiusSharpAsymptotic.lean},
\path{FabiusFullAsymptoticExpansion.lean} &
Prove the arbitrary-small-argument Lambert main and continue the saddle expansion to all
orders.\\
\bottomrule
\end{longtable}
\endgroup

The coefficient-saddle derivation in \Cref{dyadicasym:sec:coefficient-saddle} is deliberately marked as
an alternative exposition.  The formal dyadic proof uses the endpoint--Laplace comparison;
the formal arbitrary-argument proof uses the Bromwich/saddle kernel.  No claim is made that
the Cauchy coefficient route appears as a separate theorem under that exact name.

\section{Conclusions}
\label{dyadicasym:sec:conclusion}

The conspicuous gap between exact dyadic evaluation and endpoint asymptotics can be closed
without adding a new special function or an independent heuristic ansatz.

First, the exact reciprocal-power values are controlled by a single coefficient sequence:
\begin{equation*}
 F(2^{-n})=2^{-\binom n2}a_n,
 \qquad
 A(z)=\sum_{n\ge0}a_nz^n=\prod_{j\ge1}E(z/2^j).
\end{equation*}
The recurrence and composition formula are simply two expansions of this identity.

Second, the $q$-binomial formula reduces exactly to the same sequence.  The inner Thue--Morse
sum is a coefficient of $A(-2^kX)/A(-X)$, and the outer half-base Gaussian row is the
leading-coefficient functional on the nodes $-2^k$.  Therefore
\begin{equation*}
 \mathcal N_n=\M_na_n
\end{equation*}
with no approximation.  This is the direct algebraic connection that makes asymptotic
analysis of the finite $q$-formula feasible.

Third, the negative-axis product of $A$ satisfies an exact dyadic dilation equation.  Removing
its quadratic drift leaves a smooth periodic function.  A second-order comparison between the
endpoint kernel $(1-X)^n$ and the Laplace kernel $\EulerE^{-nX}$ supplies the first cumulant
correction; differentiating the exact periodic decomposition identifies that correction; and
the lower-Lambert fixed point moves it to the exact phase.  The result is
\begin{equation*}
 \log F(2^{-n})=
 -\frac L2\lambda_n^2+
 \left(1+\frac L2\right)\lambda_n-\frac12\log\lambda_n
 +C_F+\Psi(\lambda_n)+\BigO(n^{-1}).
\end{equation*}

Finally, the inverse dyadic sequence is richer than a sparse one-phase sample: the fractional
parts of $\lambda_n$ are dense, so the exact rational values determine the entire continuous
periodic correction after renormalization.  What they do not provide by monotonicity alone is
a uniform continuum theorem at the sharp error scale; that step still requires a uniform
saddle argument or a two-parameter analysis of the general dyadic formula.

The answer to the motivating question is therefore affirmative in a strong sense.  The
small-argument asymptotic can be re-derived from the exact calculation formulae, including the
constant, the logarithmic correction, the exact lower-Lambert phase, and the nonconstant
periodic fluctuation.  The only essential move is to respect the exact cancellations built
into those formulae rather than attempting to estimate their finite summands independently.

\clearpage
\section*{Supplementary material from this synthesis}
\addcontentsline{toc}{section}{Supplementary material from this synthesis}
\section{Algebra behind the expanded dyadic formula}
\label{dyadicasym:app:algebra}

This appendix records the expansion needed to pass from the compact Lambert form to the
explicit expression in $n$.  Put $\ell=\log n$ and seek
\begin{equation}
 \lambda_n=n+\frac\ell L+\frac{c_1}{n}+\frac{c_2}{n^2}
 +\BigO(\ell^3/n^3).                                       \label{dyadicasym:eq:lambda-ansatz}
\end{equation}
The fixed-point equation is
\begin{equation}
 \lambda_n=n+\frac1L\log\lambda_n.                         \label{dyadicasym:eq:lambda-fixed-app}
\end{equation}
Writing
\begin{equation*}
 \lambda_n=n\left(1+\frac\ell{Ln}+\frac{c_1}{n^2}
 +\frac{c_2}{n^3}+\cdots\right)
\end{equation*}
and expanding the logarithm gives
\begin{equation*}
 \log\lambda_n=
 \ell+\frac\ell{Ln}+\frac{c_1-\ell^2/(2L^2)}{n^2}
 +\BigO(\ell^3/n^3).
\end{equation*}
Comparison in \eqref{dyadicasym:eq:lambda-fixed-app} yields
\begin{equation}
 c_1=\frac\ell{L^2},
 \qquad
 c_2=\frac{\ell-\ell^2/2}{L^3},                            \label{dyadicasym:eq:c1-c2}
\end{equation}
which gives the first four terms of \eqref{dyadicasym:eq:lambda-expansion}.

Now set
\begin{equation*}
 H(\lambda)=-\frac L2\lambda^2+
 \left(1+\frac L2\right)\lambda-\frac12\log\lambda.
\end{equation*}
Substitution of \eqref{dyadicasym:eq:lambda-ansatz} with \eqref{dyadicasym:eq:c1-c2} and collection through order
$1/n$ gives
\begin{align}
 H(\lambda_n)={}&-\frac L2n^2-n\ell+
 \left(1+\frac L2\right)n-\frac{\ell^2}{2L}                \notag\\
 &-\frac{\ell^2}{2L^2n}+\BigO(n^{-1}).                    \label{dyadicasym:eq:H-expanded}
\end{align}
The notation $\BigO(n^{-1})$ in \eqref{dyadicasym:eq:H-expanded} includes bounded multiples of $1/n$;
the displayed $\ell^2/n$ term is separated because it is asymptotically larger.  Adding
$C_F+\Psi(\lambda_n)$ proves \eqref{dyadicasym:eq:compact-to-expanded}.

\section{A stable numerical recipe}
\label{dyadicasym:app:numerical}

The following procedure reproduces \Cref{dyadicasym:tab:numerics} without evaluating the unstable signed
$q$-sum.

\begin{algorithm}[Exact arithmetic plus convergent products]
\label{dyadicasym:alg:numerics}
\leavevmode
\begin{enumerate}
\item Compute $a_0=1$ and, for $1\le j\le n$,
\begin{equation*}
 a_j=\frac1{2^j-1}\sum_{k<j}\frac{a_k}{(j-k+1)!}
\end{equation*}
in rational arithmetic.

\item Form the exact rational
\begin{equation*}
 F(2^{-n})=2^{-\binom n2}a_n
\end{equation*}
and take its logarithm only after conversion to the desired floating-point precision.

\item Solve $\lambda-\log\lambda/L=n$ on the large branch, either by
$\lambda=-W_{-1}(-L2^{-n})/L$ or by Newton iteration initialized with
$n+\log n/L$.

\item Reduce $u=\lambda-\Floor{\lambda}$ and put $s=2^u\in[1,2)$.  Evaluate
\begin{align*}
 P(u)={}&\sum_{j\ge1}\log\frac{1-\EulerE^{-s/2^j}}{s/2^j}
 +\frac L2(u^2-u)\\
 &+\sum_{m\ge0}\log(1-\EulerE^{-s2^m}).
\end{align*}
The first series has a geometric small-argument tail and the second a doubly exponential tail.
Use stable primitives such as \texttt{expm1} for the first factors.

\item Subtract $\overline P$ from \eqref{dyadicasym:eq:P-mean} to obtain $\Psi(u)$.  Differentiate the
locally uniformly convergent series termwise, or use the rapidly convergent Fourier series
\eqref{dyadicasym:eq:Psi-Fourier-report}, to obtain $\Psi'$ and $\Psi''$.

\item Evaluate $\mathcal M_n$ from \eqref{dyadicasym:eq:Mn} and, when desired, $A_1$ from
\eqref{dyadicasym:eq:A1-report}.
\end{enumerate}
\end{algorithm}

For a direct verification of the $q$-formula one may compute
$\mathcal N_n=\M_na_n$ exactly and compare it with the finite signed expression.  The equality
is a better test than comparing floating-point approximations of the two sides because it
checks both cancellation layers over the rationals.

\section{Notation cross-reference}
\label{dyadicasym:app:notation}

\begingroup\small
\begin{longtable}{@{}ll@{}}
\toprule
Symbol & Meaning\\
\midrule
\endhead
$L$ & $\log2$\\
$F$ & bounded Fabius function / distribution function on $[0,1]$\\
$\RvachevUp$ & Rvachev up-function, with $\RvachevUp(t)=F(1-|t|)$ on $[-1,1]$\\
$X$ & random variable with distribution function $F$\\
$d_n$ & half moment $\Expectation(1-X)^n$\\
$a_n$ & factorially normalized half moment $d_n/n!$\\
$A(z)$ & $\sum a_nz^n=\Expectation\EulerE^{zX}$\\
$E(z)$ & $(\EulerE^z-1)/z$\\
$\M_n$ & $\prod_{j=1}^n(2^j-1)$\\
$\mathcal N_n$ & finite $q$-binomial/Thue--Morse numerator for $F(2^{-n})$\\
$B(s)$ & $A(-s)=\Expectation\EulerE^{-sX}$\\
$\Lam(s)$ & $\log B(s)$\\
$P$ & exact one-periodic correction in $\Lam$\\
$\Psi$ & $P-\overline P$, mean-zero periodic correction\\
$C_F$ & constant $\overline P-\tfrac12\log(2\pi)$\\
$\lambda_n$ & large solution of $\lambda-\log\lambda/L=n$\\
$\mathcal M_n$ & compact Lambert main term for $\log F(2^{-n})$\\
\bottomrule
\end{longtable}
\endgroup

\renewcommand{\refname}{Topic references}
\begin{thebibliography}{99}

\bibitem{dyadicasym:arias1982}
J.~Arias de Reyna,
\newblock Definici\'on y estudio de una funci\'on indefinidamente diferenciable de soporte
  compacto,
\newblock \emph{Revista de la Real Academia de Ciencias Exactas, F\'isicas y Naturales de
  Madrid} \textbf{76} (1982), no.~1, 21--38;
  English translation, arXiv:\nolinkurl{1702.05442} (2017).

\bibitem{dyadicasym:arias2018}
J.~Arias de Reyna,
\newblock Arithmetic of the Fabius function,
\newblock \emph{INTEGERS} \textbf{18} (2018), Article A51;
  arXiv:\nolinkurl{1702.06487}, version~3.

\bibitem{dyadicasym:andrews1976}
G.~E.~Andrews,
\newblock \emph{The Theory of Partitions},
\newblock Encyclopedia of Mathematics and its Applications~2,
  Addison--Wesley, 1976.

\bibitem{dyadicasym:allouche2003}
J.-P.~Allouche and J.~Shallit,
\newblock \emph{Automatic Sequences: Theory, Applications, Generalizations},
\newblock Cambridge University Press, 2003.

\bibitem{dyadicasym:corless1996}
R.~M.~Corless, G.~H.~Gonnet, D.~E.~G.~Hare, D.~J.~Jeffrey, and D.~E.~Knuth,
\newblock On the Lambert $W$ function,
\newblock \emph{Advances in Computational Mathematics} \textbf{5} (1996), 329--359.

\bibitem{dyadicasym:dlmf17}
NIST Digital Library of Mathematical Functions, Chapter~17,
\newblock \emph{$q$-Hypergeometric and Related Functions},
\newblock \url{https://dlmf.nist.gov/17}.

\bibitem{dyadicasym:flajolet1995}
P.~Flajolet, X.~Gourdon, and P.~Dumas,
\newblock Mellin transforms and asymptotics: harmonic sums,
\newblock \emph{Theoretical Computer Science} \textbf{144} (1995), 3--58.

\bibitem{dyadicasym:flajolet2009}
P.~Flajolet and R.~Sedgewick,
\newblock \emph{Analytic Combinatorics},
\newblock Cambridge University Press, 2009.

\bibitem{dyadicasym:gasper2004}
G.~Gasper and M.~Rahman,
\newblock \emph{Basic Hypergeometric Series}, second edition,
\newblock Encyclopedia of Mathematics and its Applications~96,
  Cambridge University Press, 2004.

\bibitem{dyadicasym:hayman1956}
W.~K.~Hayman,
\newblock A generalisation of Stirling's formula,
\newblock \emph{Journal f\"ur die reine und angewandte Mathematik}
  \textbf{196} (1956), 67--95.

\bibitem{dyadicasym:haugland2016}
J.~K.~Haugland,
\newblock Evaluating the Fabius function,
\newblock arXiv:\nolinkurl{1609.07999} (2016).

\bibitem{dyadicasym:kac2002}
V.~Kac and P.~Cheung,
\newblock \emph{Quantum Calculus},
\newblock Universitext, Springer, 2002.

\bibitem{dyadicasym:reshetnikovMSE}
V.~Reshetnikov,
\newblock A conjectured closed form for the Fabius function at dyadic rationals,
\newblock Mathematics Stack Exchange, question 3283519, 5 July 2019.

\bibitem{dyadicasym:proveit2026}
V.~Reshetnikov,
\newblock \emph{ProveIt: Analysis/FabiusFunction},
\newblock Lean~4 formal development, repository snapshot
  \path{a24af8347aba5d38b8febf2cf9a19eeef6aba18a}, 25 August 2026.
\newblock
\url{https://github.com/VladimirReshetnikov/ProveIt/tree/a24af8347aba5d38b8febf2cf9a19eeef6aba18a/Analysis/FabiusFunction}.

\bibitem{dyadicasym:fabiusExposition2026}
V.~Reshetnikov,
\newblock \emph{The Fabius Function and Rvachev's Up-Function},
\newblock technical exposition in the same repository snapshot, 25 August 2026.

\end{thebibliography}


\renewcommand{\arraystretch}{1}
% END SYNTHESIZED TOPIC: exact dyadic arithmetic to endpoint asymptotics





% BEGIN SYNTHESIZED TOPIC: quantitative Thue--Morse convergence
% Former source SHA-256 values:
% 577c5f3426def68a774fedf3fce61552d32200c8da526ace44178f2a8995a6d3
% 384d69b461cb94af33f1c080703e269ea77104405d1940413f1944172b3312c0
\clearpage
\part{Signed Thue--Morse Prefix Sums and the Fabius Function}
\label{part:tmrates}
\begin{boxedremark}
\textbf{Synthesized provenance and status.} This part merges the former
\nolinkurl{Fabius_Thue_Morse_Convergence_Rate/Fabius_Thue_Morse_Convergence_Rate.tex}
and
\nolinkurl{Fabius_Thue_Morse_Convergence_Rate-2/Fabius_Thue_Morse_Convergence_Rates.tex}.
The Fourier/deconvolution source supplies the canonical sharp and all-orders results;
the elementary source contributes zero-run identities, an explicit finite-level
curvature certificate, binary-phase and endpoint formulas.  The qualitative bridge
reuses formal inputs, but all quantitative rates, constants, deconvolution expansions,
and acceleration results remain natural-language proofs and formal obligations.\par\smallskip
Former source SHA-256 values:
\nolinkurl{577c5f3426def68a774fedf3fce61552d32200c8da526ace44178f2a8995a6d3} and
\nolinkurl{384d69b461cb94af33f1c080703e269ea77104405d1940413f1944172b3312c0}.
\end{boxedremark}
\section*{Topic abstract}
Let $\varepsilon_m=(-1)^{\BinaryDigitSum(m)}$ be the signed Thue--Morse sequence and let
$\sigma_{k+1}(m)=\sum_{j=0}^{m}\sigma_k(j)$ be its inclusive iterated prefix sums.  The
formal development in \repo\ proves that a shifted and normalized order-$(n+1)$ prefix row
is exactly the coefficient row of a finite product of discrete uniform laws, and that its
associated histogram converges pointwise to Rvachev's compactly supported function $\RvachevUp$.
Sampling the same row at $\Floor{2^n x}$ consequently converges to the bounded Fabius
function $F(x)$.  The existing argument is deliberately qualitative and supplies no rate.

This report makes the convergence quantitative.  At the natural cell centers, the normalized
prefix coefficients admit a uniform deconvolution expansion
\[
 H_n(\lambda)=\RvachevUp(\lambda)
 -\frac{3n+4}{72\,4^n}\RvachevUp''(\lambda)
 +\left(\frac{15n+16}{43200\,16^n}
       +\frac{(3n+4)^2}{10368\,16^n}\right)\RvachevUp^{(4)}(\lambda)
 +\BigO\!\left(\frac{n^3}{64^n}\right).
\]
It follows that the sharp uniform center-grid error is
$\frac13 n4^{-n}(1+o(1))$; the same rate holds after piecewise-linear interpolation.
The piecewise-constant histogram is less accurate because of cell quantization:
$\|\mathcal W_n-\RvachevUp\|_\infty=2^{-n}+\BigO(n4^{-n})$.
Most strikingly, the literal dyadic-floor approximation
\[
 A_n(x)=\frac{\sigma_{n+1}(\Floor{2^n x})}{2^{\binom n2}}
\]
has a still larger, entirely geometric first-order error.  Uniformly on $[0,1]$,
\[
 A_n(x)=F(x)+\frac{n+2-2\fract{2^nx}}{2^{n+1}}F'(x)
          +\BigO\!\left(\frac{n^2}{4^n}\right),
\]
and therefore
\[
 \left\|\frac{2^{n+1}}n(A_n-F)-F'\right\|_\infty\longrightarrow0,
 \qquad
 \frac{2^n}{n}\|A_n-F\|_\infty\longrightarrow1.
\]
The factor $n$ comes from the exact inward displacement of the finite support, not from slow
weak convergence.  At $x=1/2$ the phenomenon has the closed form
$A_n(1/2)-1/2=(n+2)2^{-n}-2^{-\binom n2}$.

The bridge identities used as input are formalized in Lean.  The Fourier asymptotic expansion,
its sharp constants, and the accelerated finite-difference constructions developed here are
new analytic deductions and are not claimed to have been formalized in the repository.


\section{Scope, conventions, and the answer in one table}
\label{tmrates:sec:scope}

The source article \cite{tmrates:reshetnikovExposition} isolates the connection between iterated
prefix sums and $\RvachevUp$, corrects two index/normalization defects in a naive formulation, and
proves pointwise convergence.  It also says explicitly that its monotonicity-and-continuity
argument gives no uniform rate.  The purpose of this supplement is to determine that rate,
including its leading constant, and to explain why several superficially similar readings of
the same integer row have different orders of accuracy.

We use the source article's normalization throughout:
\begin{itemize}
\item $F$ is the bounded distribution function, equal to $0$ on $(-\infty,0]$ and to $1$ on
      $[1,\infty)$;
\item $\RvachevUp$ is the even compactly supported density on $[-1,1]$, with
      $\RvachevUp(x-1)=F(x)$ for $0\le x\le1$;
\item the Fourier transform is
      $\FourierTwoPi{f}(\xi)=\int_{\RealNumbers}
      f(x)\EulerE^{-2\pi\ImaginaryUnit x\xi}\Differential x$.
\end{itemize}
The sharp derivative theorem from the source gives
\begin{equation}
 \|F^{(r)}\|_\infty=\|\RvachevUp^{(r)}\|_\infty
 =2^{\binom{r+1}{2}}.                                      \label{tmrates:eq:sharp-derivatives}
\end{equation}
In particular $\|F'\|_\infty=\|\RvachevUp'\|_\infty=2$ and
$\|\RvachevUp''\|_\infty=8$.

\begin{boxedremark}
\textbf{The three natural rates.}
Put $h_n=2^{-n}$.  The same normalized prefix row produces three different errors.
\begin{center}
\begin{tabularx}{0.95\textwidth}{@{}>{\raggedright\arraybackslash}p{0.23\textwidth}
>{\raggedright\arraybackslash}p{0.30\textwidth}X@{}}
\toprule
Reading of the row & Sharp uniform error & Dominant mechanism\\
\midrule
Natural center values, or their linear interpolant
& $\displaystyle \frac13 n h_n^2(1+o(1))$
& Deconvolution of the omitted mean-zero Bernoulli tail; its variance is of order $nh_n^2$.\\[0.35em]
Piecewise-constant histogram $\mathcal W_n$
& $\displaystyle h_n+\BigO(nh_n^2)$
& Horizontal quantization by half a cell, amplified by the optimal slope $2$.\\[0.35em]
Dyadic-floor Fabius sample $A_n$
& $\displaystyle n h_n(1+o(1))$
& The finite support is shorter than $[-1,1]$ by $(n+2)h_n/2$ at each end, producing an inward drift.\\
\bottomrule
\end{tabularx}
\end{center}
Thus the pointwise floor formula is not the most accurate reconstruction of the function from
the prefix row.  It is worse than the histogram by a factor asymptotic to $n$, while a
centered linear reconstruction is better than the histogram by a factor asymptotic to
$3\cdot2^n/n$.
\end{boxedremark}

\subsection{What is formalized and what is new here}

The exact prefix convolution law, Prouhet cancellation, zero runs, generating-series identity,
coefficient bridge, corrected cell centers, and pointwise limit are developed in
\path{ThueMorsePrefix.lean} and \path{ThueMorseApproximation.lean}; the finite atomic and step
approximants are developed in \path{EarlyApproximants.lean},
\path{EarlyMeasureBridge.lean}, \path{StepMeasureBridge.lean}, and
\path{StepApproximationLimit.lean} \cite{tmrates:proveit}.  These are the algebraic and
measure-theoretic inputs used below.

The new ingredient is a quantitative inversion of the finite cosine product on its natural
Nyquist interval.  It yields a complete asymptotic expansion in even derivatives of $\RvachevUp$.
The first two nonzero terms are enough to determine all sharp rates stated in the abstract.
No part of that new expansion is represented here as already machine-checked.

\section{The exact discrete bridge}
\label{tmrates:sec:bridge}

\subsection{Signed Thue--Morse prefixes}

Let
\begin{equation}
 \varepsilon_m=(-1)^{\BinaryDigitSum(m)},\qquad m\ge0,                 \label{tmrates:eq:tm-sign}
\end{equation}
where $\BinaryDigitSum(m)$ is the number of ones in the binary expansion of $m$.  Define inclusive
iterated prefix sums by
\begin{equation}
 \sigma_0(m)=\varepsilon_m,
 \qquad
 \sigma_{k+1}(m)=\sum_{j=0}^{m}\sigma_k(j).                \label{tmrates:eq:prefix-def}
\end{equation}
The word \emph{inclusive} matters: one has
\begin{equation}
 \sigma_{k+1}(m+1)-\sigma_{k+1}(m)=\sigma_k(m+1),          \label{tmrates:eq:forward-difference}
\end{equation}
with the lower-order value at the new index.

Repeated summation is convolution with the discrete truncated-power kernel.  For $k\ge1$,
\begin{equation}
 \sigma_k(m)=\sum_{r=0}^{m}\binom{m-r+k-1}{k-1}\varepsilon_r. \label{tmrates:eq:prefix-convolution}
\end{equation}
The signed generating series is
\begin{equation}
 \sum_{m\ge0}\varepsilon_m z^m=\prod_{j\ge0}(1-z^{2^j}),
 \qquad
 \sum_{m\ge0}\sigma_k(m)z^m
 =\frac{\prod_{j\ge0}(1-z^{2^j})}{(1-z)^k}.               \label{tmrates:eq:prefix-generating}
\end{equation}
All identities in this paragraph are formal power-series identities over $\IntegerNumbers$.

The order-$r$ zero of a complete block at $z=1$ gives a useful exact boundary
certificate.  For $1\le k\le r$,
\begin{align}
 \sigma_k(m)&=0 &&(2^r-k\le m\le2^r-1),
 \label{tmrates:eq:zero-run}\\
 \sigma_k(2^r-k-1)&=(-1)^{r-k},\qquad
 \sigma_k(2^r)=-1.
 \label{tmrates:eq:zero-run-boundary}
\end{align}
These dyadic zero runs explain both the long flat stretches of the prefix rows and the
exceptional right endpoint of the floor sample.

\subsection{The finite polynomial and the one-order shift}

For $n\ge0$ put
\begin{equation}
 P_n(z)=\prod_{i=1}^{n}(1+z+\cdots+z^{2^i-1}),
 \qquad
 g_n=\deg P_n=2^{n+1}-n-2.                                \label{tmrates:eq:Pn}
\end{equation}
The cancellation of the first $n+1$ factors in \eqref{tmrates:eq:prefix-generating} gives the exact
coefficient identity
\begin{equation}
 [z^m]P_n(z)=\sigma_{n+1}(m),
 \qquad 0\le m<2^{n+1}.                                   \label{tmrates:eq:coefficient-bridge}
\end{equation}
This is the source of the order shift: the polynomial indexed by $n$ corresponds to the
$(n+1)$st, not the $n$th, prefix row.

Every factor of $P_n$ is palindromic, hence
\begin{equation}
 [z^m]P_n=[z^{g_n-m}]P_n,
 \qquad
 P_n(1)=2^{\binom{n+1}{2}}.                               \label{tmrates:eq:Pn-symmetry-mass}
\end{equation}
The coefficients are nonnegative integers.  Thus, after normalization, they form a symmetric
probability mass function.

\subsection{Natural centers, atomic law, and histogram}

Set
\begin{equation}
 h_n=2^{-n},
 \qquad
 \lambda_{n,m}=\frac{2m-g_n}{2^{n+1}}
 =h_n\left(m-\frac{g_n}{2}\right).                        \label{tmrates:eq:center}
\end{equation}
The points $\lambda_{n,m}$ lie on a translate of the lattice $h_n\IntegerNumbers$.  Define
\begin{equation}
 p_{n,m}=2^{-\binom{n+1}{2}}[z^m]P_n,
 \qquad
 \mu_n=\sum_{m\in\IntegerNumbers}p_{n,m}\,\delta_{\lambda_{n,m}},    \label{tmrates:eq:atomic-law}
\end{equation}
where coefficients outside $0\le m\le g_n$ are understood to be zero.  Then $\mu_n$ is a
probability measure.  Its mass density relative to the lattice spacing is
\begin{equation}
 H_n(\lambda_{n,m})=\frac{p_{n,m}}{h_n}
 =\frac{[z^m]P_n}{2^{\binom n2}}.                          \label{tmrates:eq:Hn}
\end{equation}
By \eqref{tmrates:eq:coefficient-bridge}, this is exactly
\begin{equation}
 H_n(\lambda_{n,m})
 =\frac{\sigma_{n+1}(m)}{2^{\binom n2}},
 \qquad 0\le m<2^{n+1}.                                  \label{tmrates:eq:Hn-prefix}
\end{equation}

Let $I_{n,m}$ be the cell of width $h_n$ centered at $\lambda_{n,m}$.  Spreading each atom of
$\mu_n$ uniformly over its cell gives the step density
\begin{equation}
 \mathcal W_n(x)=\sum_m H_n(\lambda_{n,m})\IndicatorOf{I_{n,m}}(x), \label{tmrates:eq:Wn}
\end{equation}
with half weight from each adjacent cell at a shared endpoint.  In the interior of a cell,
$\mathcal W_n$ is simply the corresponding normalized prefix coefficient.  The formal
repository proves
\begin{equation}
 \mathcal W_n(x)\longrightarrow\RvachevUp(x)                     \label{tmrates:eq:qualitative-W}
\end{equation}
pointwise.  The rest of this report quantifies \eqref{tmrates:eq:qualitative-W} and the moving-center
sample derived from it.

\section{The hidden Bernoulli tail}
\label{tmrates:sec:tail}

\subsection{Rademacher factorization}

Let $R_{\ell,r}$ be independent Rademacher variables, each taking $\pm1$ with probability
$1/2$.  The binary factorization of every discrete uniform variable in \eqref{tmrates:eq:Pn} gives
an equivalent representation of \eqref{tmrates:eq:atomic-law}:
\begin{equation}
 X_n=\sum_{\ell=1}^{n}\sum_{r=1}^{\ell}
       2^{-\ell-1}R_{\ell,r},
 \qquad \mathcal L(X_n)=\mu_n.                            \label{tmrates:eq:Xn}
\end{equation}
Consequently
\begin{equation}
 \FourierTwoPi{\mu_n}(\xi)
 =\prod_{\ell=1}^{n}
  \left(\cos\frac{\pi\xi}{2^\ell}\right)^\ell.          \label{tmrates:eq:finite-cosine}
\end{equation}
The complete series
\begin{equation}
 Y=\sum_{\ell=1}^{\infty}\sum_{r=1}^{\ell}
       2^{-\ell-1}R_{\ell,r}                              \label{tmrates:eq:Y}
\end{equation}
converges absolutely, and its density is $\RvachevUp$.  Equivalently,
\begin{equation}
 \FourierTwoPi{\RvachevUp}(\xi)
 =\prod_{\ell=1}^{\infty}
  \left(\cos\frac{\pi\xi}{2^\ell}\right)^\ell.          \label{tmrates:eq:infinite-cosine}
\end{equation}
This cosine representation is equivalent to the more familiar infinite sinc product.

Write
\begin{equation}
 T_n=Y-X_n
 =\sum_{\ell>n}\sum_{r=1}^{\ell}2^{-\ell-1}R_{\ell,r}.    \label{tmrates:eq:Tn}
\end{equation}
Then $X_n$ and $T_n$ are independent and
\begin{equation}
 \RvachevUp(x)\Differential x=\mu_n*\mathcal L(T_n).                        \label{tmrates:eq:tail-convolution}
\end{equation}
Thus $\mu_n$ is obtained by deleting a small, symmetric smoothing tail from the target law.
The center-grid asymptotics are an exact lattice version of deconvolving that tail.

\subsection{Exact support radius and variance}

\begin{proposition}[Geometry and second moment of the omitted tail]
\label{tmrates:prop:tail-moments}
The tail $T_n$ has mean zero, support contained in $[-r_n,r_n]$, and variance $v_n$, where
\begin{align}
 r_n&=\sum_{\ell>n}\ell2^{-\ell-1}
     =\frac{n+2}{2^{n+1}},                                  \label{tmrates:eq:rn}\\
 v_n&=\sum_{\ell>n}\ell2^{-2\ell-2}
     =\frac{3n+4}{36\,4^n}.                                 \label{tmrates:eq:vn}
\end{align}
Both support endpoints occur.
\end{proposition}

\begin{proof}
The mean vanishes by symmetry.  The maximal possible absolute value is obtained by taking
all signs equal, so it is the first displayed sum.  The geometric-tail identity
\[
 \sum_{\ell=n+1}^{\infty}\ell q^\ell
 =q^{n+1}\frac{(n+1)-nq}{(1-q)^2}
\]
with $q=1/2$ gives \eqref{tmrates:eq:rn}.  Independence and $\Variance(R_{\ell,r})=1$ give the variance
sum, and the same identity with $q=1/4$, followed by the extra factor $1/4$, gives
\eqref{tmrates:eq:vn}.
\end{proof}

The support of $\mu_n$ is therefore
\begin{equation}
 \SupportOperator\mu_n=[-1+r_n,1-r_n]\cap\bigl(h_n(\IntegerNumbers-g_n/2)\bigr),  \label{tmrates:eq:support-mun}
\end{equation}
because $g_n/2^{n+1}=1-r_n$.  Two scales should be distinguished:
\begin{equation}
 r_n\sim\frac n2h_n,
 \qquad
 \sqrt{v_n}\sim\sqrt{\frac n{12}}\,h_n.                  \label{tmrates:eq:two-tail-scales}
\end{equation}
The support radius is a worst-case all-signs quantity, while the standard deviation is the
typical scale.  The floor-indexed approximation will inherit the former; the correctly
centered deconvolution error will inherit the variance itself.

\subsection{Why the leading center correction is second order}

If a smooth density $f$ is convolved with a small mean-zero law of variance $v$, then formally
\[
 f*\nu=f+\frac v2f''+\cdots.
\]
Solving backward gives $f-\frac v2f''+\cdots$.  Since \eqref{tmrates:eq:tail-convolution} smooths the
atomic row by the tail law, one should expect
\begin{equation}
 H_n(\lambda)-\RvachevUp(\lambda)
 \approx-\frac{v_n}{2}\RvachevUp''(\lambda)
 =-\frac{3n+4}{72\,4^n}\RvachevUp''(\lambda).                    \label{tmrates:eq:heuristic-leading}
\end{equation}
The sign is that of deconvolution: in a convex region the unsmoothed profile lies below the
smoothed one.  The next sections make \eqref{tmrates:eq:heuristic-leading} uniform and rigorous and
compute the fourth-derivative term.

\section{Fourier inversion on the moving lattice}
\label{tmrates:sec:fourier}

The main technical point is that one must invert the finite atomic law on exactly one period
of its lattice Fourier series.  Ordinary Fourier inversion of a density cannot be applied
directly to $\mu_n$.

\subsection{A shifted-lattice inversion formula}

\begin{lemma}[Nyquist inversion for a shifted lattice]
\label{tmrates:lem:lattice-inversion}
Let $\mu=\sum_{k\in\IntegerNumbers}p_k\delta_{a+hk}$ be a finite measure on a translate of $h\IntegerNumbers$, and let
$\FourierTwoPi{\mu}(\xi)=\sum_kp_k\EulerE^{-2\pi\ImaginaryUnit(a+hk)\xi}$.
Then, for every $m\in\IntegerNumbers$,
\begin{equation}
 \frac{p_m}{h}
 =\int_{-1/(2h)}^{1/(2h)}
   \FourierTwoPi{\mu}(\xi)\EulerE^{2\pi\ImaginaryUnit(a+hm)\xi}\Differential\xi.          \label{tmrates:eq:lattice-inversion}
\end{equation}
\end{lemma}

\begin{proof}
Insert the finite sum defining $\FourierTwoPi{\mu}$ and integrate term by term.  The
$k$th term is
\[
 p_k\int_{-1/(2h)}^{1/(2h)}\EulerE^{2\pi\ImaginaryUnit h(m-k)\xi}\Differential\xi,
\]
which is $p_m/h$ when $k=m$ and zero otherwise.  The shift $a$ cancels from the phase.
\end{proof}

For $\mu_n$, the interval in \eqref{tmrates:eq:lattice-inversion} is
\begin{equation}
 B_n=[-2^{n-1},2^{n-1}].                                   \label{tmrates:eq:Bn}
\end{equation}
Combining \eqref{tmrates:eq:Hn} and \eqref{tmrates:eq:lattice-inversion} gives
\begin{equation}
 H_n(\lambda_{n,m})
 =\int_{B_n}\FourierTwoPi{\mu_n}(\xi)
   \EulerE^{2\pi\ImaginaryUnit\lambda_{n,m}\xi}\Differential\xi.                  \label{tmrates:eq:Hn-inversion}
\end{equation}

\subsection{The exact deconvolution multiplier}

On $B_n$, all factors in the omitted cosine tail are positive: for $\ell>n$,
$|\pi\xi/2^\ell|\le\pi/4$.  Hence one may divide \eqref{tmrates:eq:infinite-cosine} by
\eqref{tmrates:eq:finite-cosine} and write
\begin{equation}
 \FourierTwoPi{\mu_n}(\xi)=\FourierTwoPi{\RvachevUp}(\xi)R_n(\xi),
 \qquad
 R_n(\xi)=\prod_{\ell>n}
  \left(\sec\frac{\pi\xi}{2^\ell}\right)^\ell,
 \qquad \xi\in B_n.                                      \label{tmrates:eq:Rn}
\end{equation}
The multiplier $R_n$ is the reciprocal characteristic function of the omitted tail on the
frequency range where lattice inversion needs it.

There is a useful global bound.  Put $t=\pi\xi/2^n$, so $|t|\le\pi/2$.  Reindexing
$\ell=n+r$ and using Vi\`ete's product
$\prod_{r\ge1}\cos(t/2^r)=\sin t/t$ gives
\begin{equation}
 R_n(\xi)
 =\left(\frac{t}{\sin t}\right)^n
   \prod_{r=1}^{\infty}\sec(t/2^r)^r.                     \label{tmrates:eq:Rn-viete}
\end{equation}
The second product is bounded uniformly for $|t|\le\pi/2$, while
$t/\sin t\le\pi/2$.  Thus there is an absolute $C_*$ such that
\begin{equation}
 1\le R_n(\xi)\le C_*\left(\frac\pi2\right)^n,
 \qquad \xi\in B_n.                                      \label{tmrates:eq:Rn-global-bound}
\end{equation}
This exponential growth is harmless because $\FourierTwoPi{\RvachevUp}$ decays faster
than every power.
Indeed, compact support and smoothness give, for each $M\ge1$,
\begin{equation}
 |\FourierTwoPi{\RvachevUp}(\xi)|
 \le\frac{\|\RvachevUp^{(M)}\|_1}{(2\pi|\xi|)^M}
 \le\frac{2^{1+\binom{M+1}{2}}}{(2\pi|\xi|)^M}
 \qquad(\xi\ne0),                                        \label{tmrates:eq:fourier-rapid}
\end{equation}
where \eqref{tmrates:eq:sharp-derivatives} was used in the second inequality.

\section{Sharp center-grid asymptotics}
\label{tmrates:sec:center-rate}

\subsection{Expansion of the reciprocal tail product}

For $|u|<\pi/2$,
\begin{equation}
 \log\sec u=\frac{u^2}{2}+\frac{u^4}{12}+\BigO(u^6).       \label{tmrates:eq:logsec}
\end{equation}
The two geometric tails needed below are
\begin{align}
 \frac12\sum_{\ell>n}\ell4^{-\ell}
 &=\frac{3n+4}{18\,4^n}=:\alpha_n,                        \label{tmrates:eq:alpha-n}\\
 \frac1{12}\sum_{\ell>n}\ell16^{-\ell}
 &=\frac{15n+16}{2700\,16^n}=: \beta_n.                  \label{tmrates:eq:beta-n}
\end{align}
Therefore, on a low-frequency range growing sufficiently slowly with $n$,
\begin{equation}
 \log R_n(\xi)
 =\alpha_n(\pi\xi)^2+\beta_n(\pi\xi)^4
  +\BigO\!\left(\frac{n}{64^n}|\xi|^6\right),            \label{tmrates:eq:log-Rn}
\end{equation}
and exponentiation gives
\begin{equation}
 R_n(\xi)=1+\alpha_n(\pi\xi)^2
 +\left(\beta_n+\frac{\alpha_n^2}{2}\right)(\pi\xi)^4
 +\BigO\!\left(\frac{n^3}{64^n}|\xi|^6\right).          \label{tmrates:eq:Rn-expansion}
\end{equation}

\subsection{The uniform expansion}

\begin{theorem}[Two-term lattice deconvolution expansion]
\label{tmrates:thm:center-expansion}
Uniformly for all lattice centers $\lambda=\lambda_{n,m}$,
\begin{align}
 H_n(\lambda)
={}&\RvachevUp(\lambda)
 -\frac{3n+4}{72\,4^n}\RvachevUp''(\lambda) \notag\\
 &+\left(
   \frac{15n+16}{43200\,16^n}
  +\frac{(3n+4)^2}{10368\,16^n}
  \right)\RvachevUp^{(4)}(\lambda)
 +\BigO\!\left(\frac{n^3}{64^n}\right).                 \label{tmrates:eq:center-expansion}
\end{align}
The implicit constant is absolute and effective.  In particular,
\begin{equation}
 H_n(\lambda)-\RvachevUp(\lambda)
 =-\frac{3n+4}{72\,4^n}\RvachevUp''(\lambda)
  +\BigO\!\left(\frac{n^2}{16^n}\right)                 \label{tmrates:eq:center-leading}
\end{equation}
uniformly on the moving lattice.
\end{theorem}

\begin{proof}
Use \eqref{tmrates:eq:Hn-inversion} and \eqref{tmrates:eq:Rn} and split the integral at
$|\xi|=2^{n/4}$.  On the inner interval, \eqref{tmrates:eq:Rn-expansion} is uniform.  Since
$\int_{\RealNumbers}|\xi|^6|\FourierTwoPi{\RvachevUp}(\xi)|\Differential\xi<\infty$,
its integrated remainder is
$\BigO(n^3 64^{-n})$ uniformly in the phase $\lambda$.

On the complementary part of $B_n$, combine \eqref{tmrates:eq:Rn-global-bound} with
\eqref{tmrates:eq:fourier-rapid}.  Taking, for example, $M=32$ makes the resulting integral
$o(64^{-n})$.  The tails of the three polynomial Fourier integrals outside $B_n$ are even
smaller.  We may therefore replace the truncated integrals of the displayed polynomial terms
by integrals over all of $\RealNumbers$.

With the chosen Fourier convention,
\begin{equation*}
 \int_{\RealNumbers}(\pi\xi)^2\FourierTwoPi{\RvachevUp}(\xi)
 \EulerE^{2\pi\ImaginaryUnit\lambda\xi}\Differential\xi
 =-\frac14\RvachevUp''(\lambda),
 \qquad
 \int_{\RealNumbers}(\pi\xi)^4\FourierTwoPi{\RvachevUp}(\xi)
 \EulerE^{2\pi\ImaginaryUnit\lambda\xi}\Differential\xi
 =\frac1{16}\RvachevUp^{(4)}(\lambda).
\end{equation*}
Substituting \eqref{tmrates:eq:alpha-n} and \eqref{tmrates:eq:beta-n} gives
\eqref{tmrates:eq:center-expansion}.  Omitting the fourth-derivative term and using
\eqref{tmrates:eq:sharp-derivatives} gives \eqref{tmrates:eq:center-leading}.
\end{proof}

\begin{remark}[Variance form]
Because $\alpha_n=2v_n$, the leading term is exactly
\begin{equation}
 H_n=\RvachevUp-\frac{v_n}{2}\RvachevUp''+\cdots,                        \label{tmrates:eq:variance-form}
\end{equation}
which confirms the deconvolution heuristic \eqref{tmrates:eq:heuristic-leading}.
\end{remark}

\subsection{Sharp norm constant}

\begin{corollary}[Sharp center-grid rate]
\label{tmrates:cor:center-sharp}
One has the uniform profile limit
\begin{equation}
 \sup_{m\in\IntegerNumbers}
 \left|
  \frac{4^n}{n}\bigl(H_n(\lambda_{n,m})-\RvachevUp(\lambda_{n,m})\bigr)
  +\frac1{24}\RvachevUp''(\lambda_{n,m})
 \right|\longrightarrow0.                                \label{tmrates:eq:center-profile-limit}
\end{equation}
Consequently,
\begin{equation}
 \lim_{n\to\infty}\frac{4^n}{n}
 \sup_{m\in\IntegerNumbers}|H_n(\lambda_{n,m})-\RvachevUp(\lambda_{n,m})|
 =\frac13.                                                 \label{tmrates:eq:center-sharp-norm}
\end{equation}
\end{corollary}

\begin{proof}
The coefficient $(3n+4)/(72n)$ tends to $1/24$, while the remainder in
\eqref{tmrates:eq:center-leading}, after multiplication by $4^n/n$, tends uniformly to zero.  The
lattices become dense, and $\RvachevUp''$ is continuous with exact supremum $8$.  This supremum
is attained; for example, the attained-point form of the sharp derivative theorem gives
$\RvachevUp''(-3/4)=F''(1/4)=8$.  Hence the norm limit is $8/24=1/3$.
\end{proof}

\subsection{A continuous reconstruction with the same rate}

Let $L_n$ be the piecewise-linear interpolant through the points
$(\lambda_{n,m},H_n(\lambda_{n,m}))$, with zero values continued outside the finite support.
The ordinary interpolation error for a $C^2$ function is at most
$h_n^2\|\RvachevUp''\|_\infty/8=h_n^2$.  This is smaller by a factor $1/n$ than the leading
center bias.

\begin{corollary}[Sharp continuous linear reconstruction]
\label{tmrates:cor:linear-rate}
Uniformly on $\RealNumbers$,
\begin{equation}
 \frac{4^n}{n}(L_n-\RvachevUp)\longrightarrow-\frac1{24}\RvachevUp'',   \label{tmrates:eq:linear-profile}
\end{equation}
in the supremum norm.  In particular,
\begin{equation}
 \|L_n-\RvachevUp\|_\infty=\frac n{3\,4^n}(1+o(1)).              \label{tmrates:eq:linear-sharp}
\end{equation}
\end{corollary}

Thus iterated prefix sums do not merely converge to $\RvachevUp$ after a qualitative passage through
histograms.  Read at their actual centers and interpolated linearly, they give a continuous
approximation with a fully determined, curvature-shaped leading error.

\subsection{The all-orders pattern}

The same proof is not tied to fourth order.  Write
\begin{equation}
 \log\sec z=\sum_{r\ge1}c_rz^{2r}
 \quad(c_1=1/2,\ c_2=1/12,\ c_3=1/45,\ldots)              \label{tmrates:eq:logsec-full}
\end{equation}
and put
\begin{equation}
 s_{r,n}=c_r\sum_{\ell>n}\ell2^{-2r\ell}.                 \label{tmrates:eq:srn}
\end{equation}
For a fixed $M$, expand
\begin{equation}
 \exp\left(\sum_{r=1}^{M-1}s_{r,n}z^{2r}\right)
 =\sum_{j=0}^{M-1}b_{j,n}z^{2j}+\BigO(z^{2M}).             \label{tmrates:eq:bjn}
\end{equation}
Then the lattice inversion argument gives
\begin{equation}
 H_n(\lambda)
 =\sum_{j=0}^{M-1}\frac{(-1)^j b_{j,n}}{4^j}
    \RvachevUp^{(2j)}(\lambda)
 +\BigO_M\!\left(n^M4^{-Mn}\right),                      \label{tmrates:eq:all-orders}
\end{equation}
uniformly on the center lattice.  Each $b_{j,n}$ is an explicit polynomial in the geometric
tails $s_{r,n}$.  Thus the connection admits a full asymptotic differential expansion, not
just a big-$O$ convergence estimate.

\subsection{An elementary finite-level certificate}
\label{tmrates:sec:effective-certificate}

The Fourier expansion above determines the sharp asymptotics.  A complementary argument
gives explicit, if nonsharp, bounds at every level $n\ge3$.  Let
$\beta_h=h^{-1}\IndicatorOf{[-h/2,h/2]}$ and define the mass-preserving dilation
\[
 (D_h\RvachevUp)(x)=h^{-1}\RvachevUp(x/h).
\]
Put \(q_n=\beta_{h_n}^{*n}*D_{h_n}\RvachevUp\), and rename the cardinal weights
\begin{equation}
 \pi_{n,k}=h_nq_n(kh_n),\qquad k\in\IntegerNumbers.
 \label{tmrates:eq:cardinal-weights}
\end{equation}
They form a symmetric probability distribution and satisfy the exact identities
\begin{align}
 \RvachevUp(\lambda_{n,m})
 &=\sum_{k\in\IntegerNumbers}\pi_{n,k}H_n(\lambda_{n,m-k}),
 \label{tmrates:eq:cardinal-average}\\
 \sum_{k\in\IntegerNumbers}k^2\pi_{n,k}&=\frac n{12}+\frac19.
 \label{tmrates:eq:cardinal-second-moment}
\end{align}
The second formula follows by differentiating the Poisson identity twice: the
$n$ box factors kill every nonzero integer frequency, leaving the variance
$n/12+1/9$ at the origin.

Writing $w_{n,m}=H_n(\lambda_{n,m})$, the factorization
$(1-z)^2P_n=(1-z^{2^{n-1}})(1-z^{2^n})P_{n-2}$ and the coefficient bound for
$P_{n-2}$ give
\begin{equation}
 |w_{n,m}-2w_{n,m-1}+w_{n,m-2}|\le32\,4^{-n}.
 \label{tmrates:eq:discrete-curvature}
\end{equation}
A telescoping symmetric-difference argument upgrades this to
$|w_{n,m+k}+w_{n,m-k}-2w_{n,m}|\le32\,4^{-n}k^2$.
Combining it with \eqref{tmrates:eq:cardinal-average}--
\eqref{tmrates:eq:cardinal-second-moment} yields the explicit certificate
\begin{equation}
 |H_n(\lambda_{n,m})-\RvachevUp(\lambda_{n,m})|
 \le\left(\frac{4n}{3}+\frac{16}{9}\right)4^{-n}.
 \label{tmrates:eq:effective-center-bound}
\end{equation}
Consequently
\begin{align}
 \|\mathcal W_n-\RvachevUp\|_\infty
 &\le2^{-n}+\left(\frac{4n}{3}+\frac{16}{9}\right)4^{-n},
 \label{tmrates:eq:effective-histogram-bound}\\
 \|L_n-\RvachevUp\|_\infty
 &\le\left(\frac{4n}{3}+\frac{25}{9}\right)4^{-n}.
 \label{tmrates:eq:effective-linear-bound}
\end{align}
These are effective envelopes, not replacements for the sharper limiting constants proved
above.

\section{The histogram: a sharp first-order quantization error}
\label{tmrates:sec:histogram-rate}

The center values are exceptionally accurate, but $\mathcal W_n$ holds each value constant
across a cell.  That reconstruction injects a larger first-order horizontal error.

For a point $x$ in the interior of a cell, let $c_n(x)$ denote that cell's center and define
the sawtooth phase
\begin{equation}
 \rho_n(x)=\frac{c_n(x)-x}{h_n},
 \qquad -\frac12<\rho_n(x)<\frac12.                       \label{tmrates:eq:rho}
\end{equation}
Then \eqref{tmrates:eq:center-leading} and Taylor's theorem give, uniformly away from the shared cell
endpoints,
\begin{equation}
 \mathcal W_n(x)-\RvachevUp(x)
 =h_n\rho_n(x)\RvachevUp'(x)+\BigO(nh_n^2).                      \label{tmrates:eq:hist-local}
\end{equation}
At a shared endpoint, the half-weight convention averages the two adjacent center values and
cancels the first-order odd displacement, leaving only $\BigO(nh_n^2)$ there.

\begin{theorem}[Sharp histogram rate]
\label{tmrates:thm:hist-sharp}
The step approximants satisfy
\begin{equation}
 \|\mathcal W_n-\RvachevUp\|_\infty
 =2^{-n}+\BigO\!\left(\frac n{4^n}\right).                \label{tmrates:eq:hist-sharp}
\end{equation}
Equivalently,
\begin{equation}
 \lim_{n\to\infty}2^n\|\mathcal W_n-\RvachevUp\|_\infty=1.      \label{tmrates:eq:hist-limit-constant}
\end{equation}
\end{theorem}

\begin{proof}
For an interior point, combine \eqref{tmrates:eq:center-leading} with the optimal Lipschitz bound
$\|\RvachevUp'\|_\infty=2$:
\[
 |\mathcal W_n(x)-\RvachevUp(x)|
 \le \BigO(nh_n^2)+2|c_n(x)-x|
 \le h_n+\BigO(nh_n^2).
\]
At a shared endpoint the averaged value obeys the same bound.  There remains only the thin
uncovered strip between the support of $\mathcal W_n$ and an endpoint of $[-1,1]$.
Its width is at most $r_n+h_n/2=\BigO(nh_n)$.  Since $\RvachevUp$ and all of its
derivatives vanish at $\pm1$, Taylor's theorem with the sharp bounded third derivative gives
\[
 |\RvachevUp(\pm1\mp s)|\le \frac{\|\RvachevUp^{(3)}\|_\infty}{6}s^3
 =\BigO(n^3h_n^3)=\LittleO(nh_n^2)
 \qquad(0\le s\le r_n+h_n/2).
\]
Outside $[-1,1]$ both functions vanish.  This completes the upper estimate globally.

For the lower estimate, choose cells whose centers tend to $-1/2$, where
$\RvachevUp'(-1/2)=2$, and approach one edge of each chosen cell from its interior.  The center error
is $\BigO(nh_n^2)$, while the change of $\RvachevUp$ across half a cell is
$h_n+\BigO(h_n^2)$.  Taking the supremum gives
$\|\mathcal W_n-\RvachevUp\|_\infty\ge h_n-\BigO(nh_n^2)$.
\end{proof}

\begin{remark}
The pointwise proof in the source article could not reveal \eqref{tmrates:eq:hist-sharp}: weak
convergence plus unimodality controls values through fixed comparison points but does not
track a moving point at distance $h_n/2$ from a cell center.  The rate is a local geometric
phenomenon, and the sharp constant comes from the sharp derivative bound.
\end{remark}

\section{The floor-prefix approximation to the bounded Fabius function}
\label{tmrates:sec:floor-rate}

\subsection{Definition and exact moving-center displacement}

The source theorem obtains $F$ from the prefix row by evaluating its $\RvachevUp$ approximation at
half the argument.  Define
\begin{equation}
 A_n(x)=\mathcal J_n(x/2)
 =\frac{\sigma_{n+1}(\Floor{2^nx})}{2^{\binom n2}},
 \qquad 0\le x\le1.                                      \label{tmrates:eq:An}
\end{equation}
Let
\begin{equation}
 \theta_n(x)=\fract{2^nx}=2^nx-\Floor{2^nx}.        \label{tmrates:eq:theta}
\end{equation}
The normalized prefix value in \eqref{tmrates:eq:An} is not located at $x-1$.  Its actual center is
\begin{align}
 \lambda_n(x)
 &=\lambda_{n,\Floor{2^nx}} \notag\\
 &=x-1+\delta_n(x),                                        \label{tmrates:eq:lambda-drift}\\
 \delta_n(x)
 &=\frac{n+2-2\theta_n(x)}{2^{n+1}}.                       \label{tmrates:eq:delta}
\end{align}
This identity, not merely the fact that $\lambda_n(x)\to x-1$, determines the dominant error.
It also gives
\begin{equation}
 \frac{n}{2^{n+1}}<\delta_n(x)\le\frac{n+2}{2^{n+1}}.     \label{tmrates:eq:delta-range}
\end{equation}
In terms of the omitted-tail support radius \eqref{tmrates:eq:rn},
\begin{equation}
 \delta_n(x)=r_n-\theta_n(x)h_n.                           \label{tmrates:eq:delta-rn}
\end{equation}
Thus the index $\Floor{2^nx}$ compensates for the ordinary fractional cell position but
not for the shortening of the finite support.  The residual is roughly $r_n\sim nh_n/2$,
which is much larger than either the cell width error of the histogram or the variance-scale
center bias.

\subsection{First- and second-order expansions}

\begin{theorem}[Uniform floor-prefix expansion]
\label{tmrates:thm:floor-expansion}
Uniformly for $x\in[0,1]$,
\begin{equation}
 A_n(x)
 =F(x)+\delta_n(x)F'(x)
  +\BigO\!\left(\frac{n^2}{4^n}\right).                  \label{tmrates:eq:floor-first}
\end{equation}
More precisely,
\begin{align}
 A_n(x)
={}&F(x)+\delta_n(x)F'(x) \notag\\
 &+\left(
    \frac{\delta_n(x)^2}{2}
    -\frac{3n+4}{72\,4^n}
   \right)F''(x)
 +\BigO\!\left(\frac{n^3}{8^n}\right).                  \label{tmrates:eq:floor-second}
\end{align}
All remainder constants are independent of $x$ and $n$.
\end{theorem}

\begin{proof}
By the exact coefficient identification,
$A_n(x)=H_n(\lambda_n(x))$.  Apply \eqref{tmrates:eq:center-leading} at
$\lambda_n(x)=x-1+\delta_n(x)$, use $\RvachevUp(x-1)=F(x)$ and the same identity for derivatives,
and Taylor-expand around $x-1$.  Since $\delta_n=\BigO(n2^{-n})$, the first omitted Taylor
term in \eqref{tmrates:eq:floor-first} is $\BigO(n^2 4^{-n})$.

For \eqref{tmrates:eq:floor-second}, retain the quadratic Taylor term and the leading center-bias
term.  The remainders are
$\BigO(\delta_n^3)$, $\BigO(n4^{-n}\delta_n)$, and
$\BigO(n^2 16^{-n})$, all absorbed by $\BigO(n^3 8^{-n})$.
\end{proof}

The earlier elementary route supplies constants at finite level rather than only uniform
asymptotic notation.  They are weaker than the Fourier expansion above but useful as a
direct certificate.

\begin{proposition}[Effective floor-profile certificate]
\label{tmrates:prop:effective-floor-certificate}
For every \(n\geq3\) and \(x\in[0,1]\),
\begin{align}
 A_n(x)&=F(x)+F'(x)\delta_n(x)+\frac12F''(x)\delta_n(x)^2
          +R^{(2)}_n(x),\label{tmrates:eq:effective-second-order}\\
 |R^{(2)}_n(x)|&\leq
 \left(\frac{4n}{3}+\frac{16}{9}\right)4^{-n}
 +\frac43(n+2)^3 8^{-n},\label{tmrates:eq:effective-second-order-remainder}\\
 A_n(x)&=F(x)+F'(x)\delta_n(x)+R^{(1)}_n(x),\label{tmrates:eq:effective-first-order}\\
 |R^{(1)}_n(x)|&\leq
 \left(n^2+\frac{16}{3}n+\frac{52}{9}\right)4^{-n}
 <(n+3)^2 4^{-n}.\label{tmrates:eq:effective-first-order-remainder}
\end{align}
Consequently the normalized profile has the explicit uniform bound
\begin{equation}
 \sup_{x\in[0,1]}
 \left|\frac{2^n}{n}\bigl(A_n(x)-F(x)\bigr)-\frac12F'(x)\right|
 \leq \frac2n+
 \left(n+\frac{16}{3}+\frac{52}{9n}\right)2^{-n}.
 \label{tmrates:eq:effective-uniform-profile-bound}
\end{equation}
\end{proposition}

\begin{proof}
Replace \(A_n(x)=H_n(x-1+\delta_n(x))\) by \(\RvachevUp\) using
\eqref{tmrates:eq:effective-center-bound}, then apply Taylor's theorem with the sharp
bounds \(\|F''\|_\infty\leq8\) and \(\|F'''\|_\infty\leq64\) and
\(\delta_n(x)\leq(n+2)2^{-n-1}\).  This gives the two displayed remainders.  Finally use
\(2^n\delta_n/n=1/2+(1-\theta_n)/n\), \(0<1-\theta_n\leq1\), and
\(0\leq F'\leq2\) to obtain \eqref{tmrates:eq:effective-uniform-profile-bound}.
\end{proof}

Expanding the exact displacement
$\delta_n=(n/2+1-\theta_n)2^{-n}$ makes the binary phase visible.
\begin{corollary}[Three-scale binary-phase expansion]
\label{tmrates:cor:binary-phase}
Uniformly for $x\in[0,1]$,
\begin{align}
 A_n(x)-F(x)
 &=\frac12F'(x)n2^{-n}
   +F'(x)\bigl(1-\theta_n(x)\bigr)2^{-n}\notag\\
 &\quad+\frac18F''(x)n^2 4^{-n}+\BigO(n4^{-n}).
 \label{tmrates:eq:binary-phase-expansion}
\end{align}
For dyadic $x$, the phase $\theta_n(x)$ vanishes eventually.  For a general real
$x$, it is the shifted binary tail rather than a periodic function of $n$.
\end{corollary}

\begin{figure}[htbp]
\centering
\begin{tikzpicture}
\begin{axis}[
  width=0.86\textwidth,
  height=0.46\textwidth,
  ymode=log,
  xmin=3,xmax=20,
  xlabel={$n$},
  ylabel={error scale},
  grid=both,
  legend style={font=\small,at={(0.98,0.98)},anchor=north east},
  tick label style={font=\small},
  label style={font=\small}
]
\addplot+[mark=*] coordinates {
 (3,0.375) (4,0.25) (5,0.15625) (6,0.09375) (7,0.0546875)
 (8,0.03125) (9,0.017578125) (10,0.009765625) (11,0.00537109375)
 (12,0.0029296875) (13,0.0015869140625) (14,0.0008544921875)
 (15,0.000457763671875) (16,0.000244140625) (17,0.00012969970703125)
 (18,0.00006866455078125) (19,0.0000362396240234375) (20,0.000019073486328125)
};
\addlegendentry{$n2^{-n}$}
\addplot+[mark=square*] coordinates {
 (3,0.125) (4,0.0625) (5,0.03125) (6,0.015625) (7,0.0078125)
 (8,0.00390625) (9,0.001953125) (10,0.0009765625) (11,0.00048828125)
 (12,0.000244140625) (13,0.0001220703125) (14,0.00006103515625)
 (15,0.000030517578125) (16,0.0000152587890625) (17,0.00000762939453125)
 (18,0.000003814697265625) (19,0.0000019073486328125) (20,0.00000095367431640625)
};
\addlegendentry{$2^{-n}$}
\addplot+[mark=triangle*] coordinates {
 (3,0.140625) (4,0.0625) (5,0.0244140625) (6,0.0087890625)
 (7,0.00299072265625) (8,0.0009765625) (9,0.000308990478515625)
 (10,0.000095367431640625) (11,0.0000288486480712890625)
 (12,0.00000858306884765625) (13,0.000002518296241760254)
 (14,0.0000007301568984985352) (15,0.00000020954757928848267)
 (16,0.000000059604644775390625) (17,0.000000016821548342704773)
 (18,0.00000000471482053399086) (19,0.0000000013133103493601084)
 (20,0.0000000003637978807091713)
};
\addlegendentry{$n^2 4^{-n}$}
\end{axis}
\end{tikzpicture}
\caption{The three scales in \eqref{tmrates:eq:binary-phase-expansion}.  The deterministic
center drift $n2^{-n}$ dominates the one-cell scale $2^{-n}$ and the quadratic
Taylor scale $n^2 4^{-n}$.}
\label{tmrates:fig:three-scales-supplement}
\end{figure}


\begin{corollary}[Uniform profile and sharp norm rate]
\label{tmrates:cor:floor-sharp}
The entire rescaled error profile converges uniformly:
\begin{equation}
 \left\|\frac{2^{n+1}}n(A_n-F)-F'\right\|_\infty
 \longrightarrow0.                                       \label{tmrates:eq:floor-profile}
\end{equation}
Consequently,
\begin{equation}
 \lim_{n\to\infty}\frac{2^n}{n}\|A_n-F\|_\infty=1.       \label{tmrates:eq:floor-sharp}
\end{equation}
For each fixed $x$,
\begin{equation}
 \lim_{n\to\infty}\frac{2^{n+1}}n\bigl(A_n(x)-F(x)\bigr)
 =F'(x).                                                   \label{tmrates:eq:floor-pointwise-profile}
\end{equation}
\end{corollary}

\begin{proof}
By \eqref{tmrates:eq:delta},
\[
 \frac{2^{n+1}}n\delta_n(x)
 =1+\frac{2-2\theta_n(x)}n\longrightarrow1
\]
uniformly in $x$.  Multiplying the remainder in \eqref{tmrates:eq:floor-first} by
$2^{n+1}/n$ gives $\BigO(n2^{-n})$.  This proves \eqref{tmrates:eq:floor-profile}.  Taking norms and
using $\|F'\|_\infty=2$ gives \eqref{tmrates:eq:floor-sharp}.
\end{proof}

\subsection{The equivalent statement for the \texorpdfstring{$\RvachevUp$}{Up} sample}

The source notation uses
\begin{equation}
 \mathcal J_n(u)
 =\frac{\sigma_{n+1}(\Floor{2^{n+1}u})}{2^{\binom n2}}.
\end{equation}
For $0\le u<1$ put $\vartheta_n(u)=\fract{2^{n+1}u}$ and
\begin{equation}
 d_n(u)=\frac{n+2-2\vartheta_n(u)}{2^{n+1}}.
\end{equation}
Then
\begin{equation}
 \mathcal J_n(u)
 =\RvachevUp(2u-1)+d_n(u)\RvachevUp'(2u-1)
  +\BigO\!\left(\frac{n^2}{4^n}\right).                  \label{tmrates:eq:J-expansion}
\end{equation}
At $u=1$, the prefix coefficient identity is outside its range and the exact zero-run value is
$-2^{-\binom n2}$; since $\RvachevUp(1)=\RvachevUp'(1)=0$, this exceptional endpoint is much smaller than
the uniform principal scale and does not change the norm asymptotic.
At the symmetric center one has the exact value
\begin{equation}
 \mathcal J_n(1/2)=1-2^{-\binom n2}.
 \label{tmrates:eq:J-half}
\end{equation}

The bounded floor sample has superexponentially accurate endpoints:
\begin{align}
 A_n(0)&=2^{-\binom n2}, &
 A_n(1)&=1-2^{-\binom n2},
 \label{tmrates:eq:endpoint-values}\\
 A_n(0)-F(0)&=2^{-\binom n2}, &
 A_n(1)-F(1)&=-2^{-\binom n2}.
 \label{tmrates:eq:endpoint-errors}
\end{align}
The first identity is the constant coefficient of $P_n$; the second is the sharp
zero-run boundary value \eqref{tmrates:eq:zero-run-boundary}.

\section{An exact midpoint identity}
\label{tmrates:sec:midpoint}

The sharp constant in \eqref{tmrates:eq:floor-sharp} can be seen without asymptotic analysis at the
single point where $F'$ attains its maximum.

\begin{theorem}[Exact floor-prefix error at $x=1/2$]
\label{tmrates:thm:midpoint-exact}
For every $n\ge2$,
\begin{equation}
 A_n\!\left(\frac12\right)
 =\frac12+\frac{n+2}{2^n}-2^{-\binom n2}.                 \label{tmrates:eq:midpoint-exact}
\end{equation}
Equivalently,
\begin{equation}
 \frac{2^n}{n}
 \left(A_n\!\left(\frac12\right)-\frac12\right)
 =1+\frac2n-\frac{2^{n-\binom n2}}n.                     \label{tmrates:eq:midpoint-scaled}
\end{equation}
\end{theorem}

\begin{proof}
Put $M=2^{n-1}$ and write $Q=P_{n-1}=\sum_{r=0}^{d}q_rz^r$, where
$d=g_{n-1}=2^n-n-1$, extending $q_r$ by zero outside $0\le r\le d$.  Since $M<2^n$,
multiplication by the last factor of $P_n$ gives
\begin{equation}
 [z^M]P_n=\sum_{r=0}^{M}q_r=:C.                            \label{tmrates:eq:midpoint-C}
\end{equation}
The polynomial $Q$ is palindromic and has total mass
\begin{equation}
 T=Q(1)=2^{\binom n2}.                                     \label{tmrates:eq:midpoint-T}
\end{equation}
By symmetry,
\begin{equation}
 2C-T=\sum_{r=d-M}^{M}q_r.                                 \label{tmrates:eq:midpoint-band}
\end{equation}
The band in \eqref{tmrates:eq:midpoint-band} has $n+2$ coefficients.  To evaluate them, write
$Q=P_{n-2}(1+\cdots+z^{2^{n-1}-1})$.  Put
$T_0=P_{n-2}(1)=2^{\binom{n-1}{2}}$.  The $n$ interior coefficients of the band are all
$T_0$, because the convolution window contains all coefficients of $P_{n-2}$.  Each of the
two boundary coefficients is $T_0-1$: its window omits only the monic leading coefficient of
$P_{n-2}$, and palindromicity gives the same value at the opposite boundary.  Hence
\begin{equation}
 2C-T=(n+2)T_0-2.                                         \label{tmrates:eq:midpoint-band-value}
\end{equation}
Since $T/T_0=2^{n-1}$,
\[
 \frac CT=\frac12+\frac{n+2}{2^n}-\frac1T.
\]
Finally, \eqref{tmrates:eq:coefficient-bridge} and \eqref{tmrates:eq:An} identify $C/T$ with
$A_n(1/2)$, proving \eqref{tmrates:eq:midpoint-exact}.
\end{proof}

Combining this exact lower witness with the elementary center certificate and Taylor's
theorem gives a completely explicit finite-level bracket, complementary to the sharp
asymptotic profile:
\begin{equation}
 -2^{-\binom n2}
 \le \|A_n-F\|_\infty-(n+2)2^{-n}
 \le (n+3)^2 4^{-n}\qquad(n\ge3).
 \label{tmrates:eq:uniform-finite-bracket}
\end{equation}

\begin{table}[htbp]
\centering
\caption{Exact midpoint error and its sharp-rate normalization.}
\label{tmrates:tab:midpoint}
\begin{tabular}{@{}rcc@{}}
\toprule
$n$ & $A_n(1/2)-1/2$ & $2^n\bigl(A_n(1/2)-1/2\bigr)/n$\\
\midrule
2  & $0.5$                 & $1.000000000000$\\
3  & $0.5$                 & $1.333333333333$\\
4  & $0.359375$            & $1.437500000000$\\
5  & $0.2177734375$        & $1.393750000000$\\
6  & $0.124969482421875$   & $1.333007812500$\\
8  & $0.0390624962747097$  & $1.249999880791$\\
10 & $0.0117187499999716$  & $1.199999999997$\\
12 & $0.0034179687500000$  & $1.166666666667$\\
16 & $0.0002746582031250$  & $1.125000000000$\\
20 & $0.00002098083496094$ & $1.100000000000$\\
\bottomrule
\end{tabular}
\end{table}

The exact identity also explains an otherwise puzzling numerical observation: after a few
levels, the error is almost indistinguishable from the tangent displacement
$2\delta_n(1/2)=(n+2)2^{-n}$.  The remaining defect is the super-exponentially small term
$2^{-\binom n2}$.

\subsection{Correct centering can be exact at the same point}

The large error in \eqref{tmrates:eq:midpoint-exact} is not an unavoidable defect of the prefix row.
It is caused by reading the wrong index as the point $-1/2$ on the $\RvachevUp$ axis.  When $n$ is
even, $-1/2$ itself is a lattice center, with index
\begin{equation}
 m_n^*=2^{n-1}-\frac{n+2}{2}.                              \label{tmrates:eq:midpoint-correct-index}
\end{equation}
For this index, $m_n^*=(g_{n-1}-1)/2$.  Since $P_{n-1}$ is palindromic of odd degree, exactly
half of its total coefficient mass lies at indices at most $m_n^*$.  The last-factor prefix
identity therefore gives
\begin{equation}
 H_n(-1/2)=\frac12=\RvachevUp(-1/2) \qquad(n\ \text{even}).       \label{tmrates:eq:midpoint-centered-exact}
\end{equation}
The floor rule instead uses $M=2^{n-1}$, which is $(n+2)/2$ whole cells to the right of
$m_n^*$.  This is the discrete form of the support drift.

\section{Bias removal and accelerated prefix reconstructions}
\label{tmrates:sec:acceleration}

The expansion \eqref{tmrates:eq:center-expansion} can be used constructively.  It tells us how to
remove the leading deconvolution bias using only neighboring prefix coefficients.

For a function on the center lattice, define the centered second difference
\begin{equation}
 \Deltah^2H(\lambda)
 =\frac{H(\lambda+h_n)-2H(\lambda)+H(\lambda-h_n)}{h_n^2}. \label{tmrates:eq:discrete-laplacian}
\end{equation}
Put
\begin{equation}
 c_n=\frac{3n+4}{72}h_n^2
\end{equation}
and define the corrected grid
\begin{equation}
 H_n^{[1]}(\lambda)=H_n(\lambda)+c_n\Deltah^2H_n(\lambda). \label{tmrates:eq:corrected-grid}
\end{equation}
Every value in \eqref{tmrates:eq:corrected-grid} is a rational linear combination of three normalized
iterated-prefix values.

\begin{theorem}[One-step deconvolution acceleration]
\label{tmrates:thm:accelerated}
Uniformly on the center lattice,
\begin{equation}
 H_n^{[1]}(\lambda)-\RvachevUp(\lambda)
 =-\frac{n^2}{1152\,16^n}\RvachevUp^{(4)}(\lambda)
  +\LittleO\!\left(\frac{n^2}{16^n}\right).                \label{tmrates:eq:accelerated-profile}
\end{equation}
Consequently,
\begin{equation}
 \lim_{n\to\infty}\frac{16^n}{n^2}
 \sup_{\lambda\in h_n(\IntegerNumbers-g_n/2)}
 |H_n^{[1]}(\lambda)-\RvachevUp(\lambda)|
 =\frac89.                                                 \label{tmrates:eq:accelerated-norm}
\end{equation}
\end{theorem}

\begin{proof}
For a smooth function,
\begin{equation}
 \Deltah^2f=f''+\frac{h_n^2}{12}f^{(4)}+\BigO(h_n^4).      \label{tmrates:eq:second-difference-expansion}
\end{equation}
Write \eqref{tmrates:eq:center-expansion} as
$H_n=\RvachevUp-c_n\RvachevUp''+d_n\RvachevUp^{(4)}+\BigO(n^3h_n^6)$, where
\[
 d_n=\left(\frac{15n+16}{43200}
          +\frac{(3n+4)^2}{10368}\right)h_n^4.
\]
Substitute this into \eqref{tmrates:eq:corrected-grid} and use
\eqref{tmrates:eq:second-difference-expansion}.  If the uniform remainder is denoted by $R_n$, then
$\|\Deltah^2R_n\|_\infty\le4h_n^{-2}\|R_n\|_\infty$; hence
$c_n\Deltah^2R_n=\BigO(n^4h_n^6)=\LittleO(n^2h_n^4)$.  Thus the remainder remains negligible
at the claimed scale.  The second-derivative terms cancel, and the coefficient of
$\RvachevUp^{(4)}$ is
\begin{equation*}
 d_n-c_n^2+\frac{c_nh_n^2}{12}
 =-\frac{n^2}{1152}h_n^4+\BigO(nh_n^4).
\end{equation*}
This proves \eqref{tmrates:eq:accelerated-profile}.  By \eqref{tmrates:eq:sharp-derivatives},
$\|\RvachevUp^{(4)}\|_\infty=2^{10}=1024$, so the sharp norm constant is
$1024/1152=8/9$.
\end{proof}

Higher corrections follow by replacing the reciprocal-tail multiplier with successively
longer polynomials in the lattice Laplacian.  The all-orders expansion
\eqref{tmrates:eq:all-orders} supplies the coefficients recursively.  To preserve the accelerated
rate in a continuous reconstruction, one should use interpolation of sufficiently high order:
piecewise-linear interpolation has an $\BigO(h_n^2)$ geometric error and would mask the
$\BigO(n^2h_n^4)$ corrected lattice error, whereas local cubic interpolation has
$\BigO(h_n^4)$ error and retains the latter as the dominant term.

\section{Conceptual synthesis}
\label{tmrates:sec:synthesis}

\subsection{One row, three coordinate choices}

The normalized integer row
\begin{equation}
 m\longmapsto\frac{\sigma_{n+1}(m)}{2^{\binom n2}}        \label{tmrates:eq:one-row}
\end{equation}
contains highly accurate samples, but the meaning of $m$ must be specified.

\begin{enumerate}
\item The intrinsic coordinate is the center
$\lambda_{n,m}=h_n(m-g_n/2)$.  At that coordinate the row differs from $\RvachevUp$ by a curvature
term of order $nh_n^2$.
\item Holding a center value constant across its width-$h_n$ cell changes the coordinate by
at most $h_n/2$.  The target has maximal slope $2$, so the histogram error is $h_n$.
\item Replacing the intrinsic coordinate by the naive dyadic coordinate implicit in
$m=\Floor{2^nx}$ overlooks the support deficit $r_n$.  It changes the coordinate by
roughly $nh_n/2$, so the Fabius error is roughly $nh_nF'(x)/2$.
\end{enumerate}
The three rates are therefore compatible, not contradictory.  They measure distinct
reconstructions of the same discrete data.

\subsection{Variance versus support radius}

The center expansion sees the omitted tail only through its even cumulants.  Its mean is zero,
so the first nonzero effect is its variance $v_n\asymp nh_n^2$.  The floor-coordinate drift,
by contrast, sees the extreme support loss $r_n\asymp nh_n$.  This is why a centered analytic
argument is exponentially more accurate than a support-aligned but miscentered index rule.

This distinction is common in lattice approximation.  Weak convergence controls averages,
and moment matching controls smooth centered observables, but an endpoint-derived coordinate
can retain a worst-case support defect long after the bulk distribution has become extremely
close.

\subsection{Why the \texorpdfstring{$n$}{n} is structural}

The degree
\begin{equation}
 g_n=2^{n+1}-n-2
\end{equation}
is shorter than the naive value $2^{n+1}$ by exactly $n+2$.  After scaling by
$2^{n+1}$ and centering, half of this deficit appears at each endpoint.  Hence the factor $n$
in the floor rate is already encoded in the elementary degree formula for $P_n$.  No delicate
property of $F$ is needed to predict its presence; the analytic work is needed to turn the
horizontal deficit into the sharp vertical profile $F'$ and to control the smaller terms.

\subsection{A practical recommendation}

For numerical reconstruction from a computed prefix row:
\begin{itemize}
\item do not use $m/2^n$ as the physical coordinate;
\item attach each value to $\lambda_{n,m}$ from \eqref{tmrates:eq:center};
\item use at least linear interpolation rather than a histogram;
\item if another factor of roughly $4^n/n$ is valuable, apply the three-point correction
      \eqref{tmrates:eq:corrected-grid} and use cubic interpolation.
\end{itemize}
The coefficients remain exact rationals until the final interpolation or evaluation stage.

\section{A Lean formalization roadmap}
\label{tmrates:sec:lean-roadmap}

The new proof separates naturally into reusable lemmas.  A formalization need not begin with
the full asymptotic theorem.

\begin{enumerate}
\item \textbf{Shifted-lattice Fourier inversion.}  Formalize \cref{tmrates:lem:lattice-inversion} for
a finitely supported measure on $a+h\IntegerNumbers$.  This is a finite-sum orthogonality calculation.
\item \textbf{Rademacher realization.}  Refine the existing finite-measure bridge to the
triangular Rademacher array \eqref{tmrates:eq:Xn}; obtain \eqref{tmrates:eq:finite-cosine} and the independent
tail decomposition \eqref{tmrates:eq:Tn}.
\item \textbf{Exact geometric tails.}  Add the closed forms \eqref{tmrates:eq:rn}, \eqref{tmrates:eq:vn},
\eqref{tmrates:eq:alpha-n}, and \eqref{tmrates:eq:beta-n}.  These reduce to finite manipulations of geometric
series.
\item \textbf{Uniform $\log\sec$ remainder.}  Prove a sixth-order Taylor bound on
$[-1/4,1/4]$ and use it scale by scale.  The low-frequency cutoff can be chosen as
$2^{n/4}$ or replaced by any convenient dyadic integer bound.
\item \textbf{Rapid Fourier decay.}  Package repeated integration by parts for compactly
supported $C^\infty$ functions and instantiate it with the already formalized sharp derivative
bounds.
\item \textbf{Center expansion.}  Combine the preceding lemmas with Fourier derivative
identities.  The first-order version \eqref{tmrates:eq:center-leading} is substantially shorter than
the fourth-order theorem and already proves all three basic rates.
\item \textbf{Exact drift.}  Formalize \eqref{tmrates:eq:lambda-drift} by integer-floor arithmetic.
Then \eqref{tmrates:eq:floor-first} is a direct Taylor estimate.
\item \textbf{Midpoint identity.}  Prove \cref{tmrates:thm:midpoint-exact} independently of Fourier
analysis from palindromicity and the central coefficient plateau.  It is an especially useful
regression test for index conventions.
\end{enumerate}

The formal theorem statements should distinguish three domains: values at the center lattice,
the step function on all reals, and the floor sample on $[0,1]$.  Combining them into one
statement would obscure exactly the coordinate issue responsible for the different rates.

\section{Conclusion}
\label{tmrates:sec:conclusion}

Iterated summation of the signed Thue--Morse sequence does more than produce a sequence that
happens to resemble the Fabius function.  After the correct one-order shift, its finite row is
an exact lattice deconvolution of $\RvachevUp$.  The omitted part is an explicit triangular Bernoulli
tail with support radius $(n+2)2^{-n-1}$ and variance $(3n+4)/(36\,4^n)$.  Those two exact
numbers govern two different approximation regimes.

At the natural centers, the mean-zero tail leaves a variance-scale curvature bias, yielding
the sharp rate $n/(3\,4^n)$.  Holding the values constant across cells introduces the larger
sharp rate $2^{-n}$.  Reading the same row at the naive dyadic floor leaves the full support
deficit as a coordinate drift and yields the still larger sharp rate $n2^{-n}$.  The latter
has the uniform error profile $F'$ and an exact midpoint formula.

The main conceptual correction is therefore simple but consequential: the integer index is
not itself the continuum coordinate.  Once the genuine center map is retained, the
Thue--Morse prefix construction is far more accurate than its original pointwise formulation
suggests, and its error admits a complete, explicitly computable differential expansion.

\clearpage
\section*{Supplementary material from this synthesis}
\addcontentsline{toc}{section}{Supplementary material from this synthesis}
\section{Elementary geometric-tail formulas}
\label{tmrates:app:tails}

For $|q|<1$ and $n\ge0$,
\begin{equation}
 \sum_{\ell=n+1}^{\infty}\ell q^\ell
 =q\frac{\Differential}{\Differential q}\left(\frac{q^{n+1}}{1-q}\right)
 =q^{n+1}\frac{(n+1)-nq}{(1-q)^2}.                        \label{tmrates:eq:geometric-tail}
\end{equation}
The specializations used in the text are
\begin{align}
 \sum_{\ell>n}\ell2^{-\ell}&=(n+2)2^{-n},\\
 \sum_{\ell>n}\ell4^{-\ell}&=\frac{3n+4}{9\,4^n},\\
 \sum_{\ell>n}\ell16^{-\ell}&=\frac{15n+16}{225\,16^n},\\
 \sum_{\ell>n}\ell64^{-\ell}&=\frac{63n+64}{3969\,64^n}.
\end{align}
These identities make every coefficient in the all-orders expansion elementary: no
unevaluated infinite tail sum is necessary.

\section{A short exact-arithmetic algorithm}
\label{tmrates:app:algorithm}

The coefficient row $[z^m]P_n$ can be built without polynomial multiplication.  Multiplication
by $1+\cdots+z^{2^i-1}$ is a sliding-window sum.  The following pseudocode maintains exact
integers and costs $\BigO(2^n)$ arithmetic operations per level.

\begin{lstlisting}[style=pseudo]
coeff := [1]
for i := 1,...,n:
    width := 2^i
    nextLength := length(coeff) + width - 1
    next := array(nextLength, 0)
    window := 0
    for m := 0,...,nextLength-1:
        if m < length(coeff):
            window := window + coeff[m]
        if m >= width:
            window := window - coeff[m-width]
        next[m] := window
    coeff := next

# Normalized prefix sample:
H_n(m) := coeff[m] / 2^(n(n-1)/2)
# Physical coordinate:
lambda_n(m) := (2m - (2^(n+1)-n-2)) / 2^(n+1)
\end{lstlisting}

The midpoint identity \eqref{tmrates:eq:midpoint-exact}, symmetry
$\mathtt{coeff[m]}=\mathtt{coeff[g_n-m]}$, the total mass
$2^{\binom{n+1}{2}}$, and the central maximum $2^{\binom n2}$ provide independent exact
checks of an implementation.

\renewcommand{\refname}{Topic references}
\begin{thebibliography}{99}

\bibitem{tmrates:fabius1966}
J.~Fabius,
\newblock A probabilistic example of a nowhere analytic $C^\infty$-function,
\newblock \emph{Zeitschrift f\"ur Wahrscheinlichkeitstheorie und Verwandte Gebiete}
  \textbf{5} (1966), 173--174.
\newblock DOI: \nolinkurl{10.1007/BF00536652}.

\bibitem{tmrates:arias1982}
J.~Arias de Reyna,
\newblock Definici\'on y estudio de una funci\'on indefinidamente diferenciable de soporte
  compacto,
\newblock \emph{Revista de la Real Academia de Ciencias Exactas, F\'isicas y Naturales de
  Madrid} \textbf{76} (1982), no.~1, 21--38.
\newblock English translation: \emph{An infinitely differentiable function with compact
support: Definition and properties}, arXiv:\nolinkurl{1702.05442} (2017).

\bibitem{tmrates:allouche2003}
J.-P.~Allouche and J.~Shallit,
\newblock \emph{Automatic Sequences: Theory, Applications, Generalizations},
\newblock Cambridge University Press, 2003.

\bibitem{tmrates:reshetnikovExposition}
V.~Reshetnikov,
\newblock \emph{The Fabius Function and Rvachev's Up-Function: Exact arithmetic,
probability, Fourier analysis, and corrected sharp asymptotics},
\newblock human-readable synthesis in the \texttt{ProveIt} repository, snapshot
25 August 2026.
\newblock \url{https://github.com/VladimirReshetnikov/ProveIt/blob/e4aa81a0b6b7e4f9943c3a5789c4a781ee120be4/Analysis/FabiusFunction/docs/Fabius_Function_and_Rvachev_Up/Fabius_Function_and_Rvachev_Up.tex}.

\bibitem{tmrates:proveit}
V.~Reshetnikov,
\newblock \emph{ProveIt: Analysis/FabiusFunction},
\newblock Lean~4 formal development, repository snapshot
\path{e4aa81a0b6b7e4f9943c3a5789c4a781ee120be4}, 25 August 2026.
\newblock \url{https://github.com/VladimirReshetnikov/ProveIt/tree/e4aa81a0b6b7e4f9943c3a5789c4a781ee120be4/Analysis/FabiusFunction}.

\end{thebibliography}


\renewcommand{\arraystretch}{1}
% END SYNTHESIZED TOPIC: quantitative Thue--Morse convergence


% BEGIN SYNTHESIZED TOPIC: inverse and small-argument saddle analysis
% Former source provenance values:
% Historically recorded SHA-256 (not reproduced by a reachable source blob):
% f9d8605761aaaa1b2c2af83e3c5c55dcd6acfc847402163458b03b68c7b35ff8
% Authoritative archived standalone Git blob: 7b5cd9f7887de95236045de21c176d77962721bb
% 85f51a20fc7b6bdf3b1d049ec4506f508aee3c4cd70554e6a68cbcc30977cb0b
\clearpage
\part{Inverse and Small-Argument Saddle Analysis}
\label{part:inverse}
\begin{boxedremark}
\textbf{Source provenance and status.} This topic synthesizes the former
\nolinkurl{Fabius_Inverse_and_Saddle_Research_Frontiers/Fabius_Inverse_and_Saddle_Research_Frontiers.tex}
(historically recorded SHA-256 \nolinkurl{f9d8605761aaaa1b2c2af83e3c5c55dcd6acfc847402163458b03b68c7b35ff8};
the audit did not reproduce it from a reachable whole-file blob, whose authoritative archived
Git identifier is \nolinkurl{7b5cd9f7887de95236045de21c176d77962721bb})
and \nolinkurl{Small_Argument_Asymptotics/Small_Argument_Asymptotics.tex}
(SHA-256 \nolinkurl{85f51a20fc7b6bdf3b1d049ec4506f508aee3c4cd70554e6a68cbcc30977cb0b}).
The inverse dossier supplies the canonical mathematical backbone; the specialist
phase-polynomial and right-endpoint transfers, exact coefficient means, higher coefficients,
Fourier-amplitude analysis, numerical checks, traps, and open obligations are retained from
the small-argument notebook.  Formalized inputs are distinguished from
derived and numerical consequences throughout.
\end{boxedremark}
\section*{Topic abstract}
The Fabius function is unusually rich: a single compactly supported
$C^\infty$ density produces a self-similar distribution function, exact
dyadic arithmetic, Fourier products, a probability model, and a
small-argument expansion whose coefficients are smooth periodic functions.
The companion article \emph{The Fabius Function and Rvachev's Up-Function}
develops most of that landscape.  The present paper is deliberately
different.  It does not repeat the introductory theory.  Instead it fills
two identifiable gaps in the exposition by turning recently formalized
material into ordinary mathematics.

The first gap is the inverse.  We construct the totalized inverse
$\FabiusClampedQuantile=F^{-1}$ on the whole real line, prove its global monotonicity,
continuity, reflection law, comparison calculus, and exact values at all
outputs of the dyadic evaluator, and show how every flatness estimate for
$F$ becomes an effective steepness estimate for $\FabiusClampedQuantile$.  We then derive
the classical interior calculus of the inverse and invert the sharp
Lambert-phase asymptotic.  The result is the explicit endpoint law
\[
 \FabiusClampedQuantile(y)\sim
 \frac{\sqrt{\log(1/y)}}{\EulerE\sqrt{\log 2}}\,
 \exp\!\left(-\sqrt{2\log 2\,\log(1/y)}\right)
 \qquad (y\downarrow0),
\]
together with its reflected analogue at $1$.  This explains in one formula
why the inverse is steeper than every positive power of $y$, although it
still tends to zero faster than every negative power of $\log(1/y)$.

The second gap concerns the internal machinery of the all-orders saddle
expansion.  We give complete derivations of the closed periodic jet
formula, its shifted signed-Stirling conversion, the collapse of the
centered exponent coefficients, the formal mass-to-logarithm recursions,
the leading-derivative law, and the full second logarithmic correction.
Finally, we give a proof-level account of the arbitrary-order analytic
remainder: the order-dependent central window, vertical Taylor control,
finite exponential substitution, Gaussian contraction, polynomial-tail
replacement, minor-arc estimates, logarithmic transfer, and return from
the dyadic logarithmic scale to $x\downarrow0$.

Statements explicitly represented by theorem declarations in the Lean
development are marked \Formalized.  Consequences proved here from those
statements by standard real or analytic arguments are marked \Derived.


\begin{boxedremark}
\textbf{Formalization boundary.}
This document is a research-frontier notebook, not part of the project's
formalization-backed primary exposition.  It deliberately preserves a mixture of proved
Lean inputs, ordinary mathematical deductions, symbolic coefficient calculations, and
numerical evidence so that the unformalized steps are not lost.  A statement here is a
project theorem only when an exact current Lean declaration is named for the full statement;
the labels \Formalized\ and \Derived\ record the draft's provenance but do not themselves
certify that every displayed consequence has been formalized.  In particular, H\"older
consequences beyond the now-formalized inverse-curvature identities, inverse endpoint
asymptotics, the ordered-composition coefficient formulas and the logarithmic
highest-derivative law remain research obligations.
\end{boxedremark}

\section{Purpose, scope, and status conventions}

\subsection{A supplement rather than a second survey}

Let $F$ denote the bounded, totalized Fabius function: $F(x)=0$ for
$x\leq0$, $F(x)=1$ for $x\geq1$, and on $[0,1]$ it is the unique strictly
increasing $C^\infty$ function associated with Rvachev's compactly
supported up-function.  We use only a small portion of the foundational
theory:
\begin{align}
 F(0)&=0,&F(1)&=1,&F(1-x)&=1-F(x), \label{inverse:eq:F-basic}\\
 F'(x)&=2\RvachevUp(2x-1),&&&0<x<1, \label{inverse:eq:Fprime-up}
\end{align}
where $\RvachevUp$ is positive on $(-1,1)$ and $\RvachevUp(0)=1$.  Consequently $F$ is a
continuous strictly increasing bijection of $[0,1]$ onto itself.
The exact second derivative is now also a named formal input:
\[
 F''(x)=8\bigl(\RvachevUp(4x-1)-\RvachevUp(4x-3)\bigr).
\]
Its positive, negative, and zero loci are respectively \((0,1/2)\),
\((1/2,1)\), and \(\RealNumbers\setminus(0,1)\cup\{1/2\}\), by
\path{deriv_deriv_fabiusReal_pos_iff},
\path{deriv_deriv_fabiusReal_neg_iff}, and
\path{deriv_deriv_fabiusReal_eq_zero_iff}.

The main companion article already explains the random-series model,
infinite convolutions, the Fourier transform, moment recurrences,
Thue--Morse identities, dyadic rationality, and the broad shape of the
Lambert saddle analysis.  Repeating those chapters would obscure the
missing material.  The division of labor adopted here is summarized in
\cref{inverse:tab:audit}.

\begin{longtable}{>{\raggedright\arraybackslash}p{0.24\textwidth}
                  >{\raggedright\arraybackslash}p{0.35\textwidth}
                  >{\raggedright\arraybackslash}p{0.31\textwidth}}
\caption{Expository audit and the role of this companion.}
\label{inverse:tab:audit}\\
\toprule
Topic & Situation in the main article & Treatment here\\
\midrule
\endfirsthead
\toprule
Topic & Situation in the main article & Treatment here\\
\midrule
\endhead
Inverse function & The current article summarizes the totalized inverse and its endpoint
estimates. & Full construction, exact identities, dyadic inversion,
effective steepness, endpoint limits, regularity, and inverse
asymptotics.\\
Closed saddle jets & Closed formulas are displayed, but the operator
calculus and the shifted Stirling conversion are compressed. & Complete
induction, forcing-term propagation, coefficient table, and the reason for
the index shift $s(n+1,m+1)$.\\
Mass and logarithmic coefficients & The all-orders result is stated and
some coefficients are listed. & Formal exponential recursion, Gaussian
contraction, composition formulas, structural bounds, and the leading
derivative law.\\
Second correction & The final polynomial appears, while much of the
order-four algebra is easy to miss. & A line-by-line calculation from
$E_1,\dots,E_4$ through $B_4$, Gaussian moments, $a_2$, and $A_2$.\\
Arbitrary-order remainder & The main article gives a high-level proof
architecture. & The order-dependent contour window and every analytic
error mechanism are separated and justified.\\
Probability, Fourier theory, dyadic arithmetic & Already developed in
substantial depth. & Used as input only; not duplicated.\\
\bottomrule
\end{longtable}

\subsection{What the labels mean}

A paragraph headed \Formalized\ restates a mathematical theorem whose
content is represented directly in the Lean files under
\begin{center}
\path{Analysis/FabiusFunction/Lean/FabiusFunction}.
\end{center}
A paragraph headed \Derived\ gives a paper-level consequence.  Its proof
uses formalized results plus a standard theorem such as the inverse
function theorem, ordinary asymptotic inversion, or elementary formal
power-series algebra.  The distinction matters especially in
\cref{inverse:sec:inverse-calculus,inverse:sec:coefficient-formulas}: those sections contain
new synthesis, but they do not pretend that every displayed corollary has
a declaration of its own.

\begin{boxedremark}
The exact Lambert phase is never silently replaced inside a periodic
coefficient.  If $\lambda=\lambda(x)$ is the exact lower solution of
$x=\lambda2^{-\lambda}$, then a term such as $A_j(\lambda)$ remains at that
exact phase.  Replacing $\lambda$ by an asymptotic approximation inside
$A_j$ is a separate composition problem, because a bounded phase error
does not become small modulo the period.
\end{boxedremark}

\subsection{Local saddle notation and its crosswalk}
\label{inverse:sec:notation-crosswalk}
\label{smallarg:sec:crosswalk}

Three namings for the same principal saddle objects appear across the specialist source, the
Wikipedia draft, and the Lean development:
\begin{center}
\begin{tabular}{@{}lll@{}}
\toprule
this synthesis / the primary article & Wikipedia draft & Lean object\\
\midrule
\(\lambda(x)\) & \(\lambda\) & \nolinkurl{fabiusLambertPhase x}\\
\(\Psi\) & \(\Psi\) & \nolinkurl{negativeLaplacePsi}\\
\(\mathcal M(x)\) & main term & \nolinkurl{fabiusSharpLambertMain x}\\
\(A_j\) & \(P_j\) & \nolinkurl{fabiusSaddleLogCoefficient j}\\
\(a_j\) & \(b_j\) & \nolinkurl{fabiusSaddleMassCoefficient j}\\
\(A_1\) & \(P_1\) & \nolinkurl{fabiusFirstSaddleCorrection}\\
\(A_2\) & \(P_2\) & \nolinkurl{fabiusSecondSaddleCorrection}\\
\(C_F\) & \(C_F\) & \nolinkurl{fabiusSharpAsymptoticConstant}\\
\bottomrule
\end{tabular}
\end{center}

The Wikipedia draft calls the coefficient functions merely continuous.  The direct
regularity result \cref{inverse:cor:A-structure} makes them smooth, one-periodic, and
bounded, while \cref{inverse:sec:closed-jets,inverse:sec:coefficient-formulas} constructs
them from smooth periodic jets.

The local jets and exponent coefficients have no counterpart in the Wikipedia draft.  The
primary article gives them different letters because \(p\), \(d\), and \(E\) are already
used there: \(p_j\) is the Lambert phase polynomial and \(p_m\) the binary prefix, \(d_n\)
is the half-moment sequence, and \(E\) denotes both the tail error and the exponential
quotient.  The exact correspondence is
\begin{center}
\begin{tabular}{@{}lll@{}}
\toprule
Local symbol & Main-exposition symbol & Lean object \\
\midrule
\(p_n\) & \(Z_n\) & \nolinkurl{negativeLaplacePeriodicJet} \\
\(d_n\) & \(\widetilde Z_n\) & \nolinkurl{negativeLaplaceBoundedExponentJet} \\
\(E_m\) & \(\Xi_m\) & \nolinkurl{negativeLaplaceExponentCoefficient} \\
\(\Differential/\Differential t\) & \(\mathrm D\) & --- \\
\bottomrule
\end{tabular}
\end{center}

The tilde in \(\widetilde Z_n\) is apt: the periodic and bounded jet families differ by the
single constant \((-1)^{n+1}n!\), by \eqref{inverse:eq:d-from-p}.  In the derivative row,
\(t\) is the independent variable used by the former specialist source; the consolidated
synthesis otherwise writes that variable as \(u\) and defines \(D=\Differential/\Differential u\).  In this
final part, \(d_n\) is a bounded saddle jet and is not the half-moment sequence
\(d_n=\Expectation X^n\) used in the probability and dyadic parts.  When numerical work later returns
to half moments, it writes them as \(h_n\) to keep the two roles visibly separate.

\section{The inverse as a global mathematical object}
\label{inverse:sec:inverse}

\subsection{The order isomorphism on the unit interval}

Set
\[
 I=[0,1],\qquad
 \clamp(y)=\min\{1,\max\{0,y\}\}.
\]
Because $F|_I$ is a continuous strictly increasing bijection $I\to I$, it
is an order isomorphism.  Its ordinary inverse will first be denoted by
$F_I^{-1}:I\to I$.

\begin{definition}[Totalized inverse]
\label{inverse:def:totalized-inverse}
The totalized inverse of the Fabius function is
\[
 \FabiusClampedQuantile(y)=F_I^{-1}\bigl(\clamp(y)\bigr),\qquad y\in\RealNumbers.
\]
Thus the inverse is extended by constants outside its natural range,
exactly as $F$ itself is extended by constants outside its natural domain.
\end{definition}

\begin{theorem}[Global inverse package; \Formalized]
\label{inverse:thm:inverse-package}
The function $\FabiusClampedQuantile$ has the following properties.
\begin{enumerate}
\item $0\leq \FabiusClampedQuantile(y)\leq1$ for every $y\in\RealNumbers$.
\item If $0\leq y\leq1$, then
\[
 F(\FabiusClampedQuantile(y))=y.
\]
If $0\leq x\leq1$, then
\[
 \FabiusClampedQuantile(F(x))=x.
\]
\item $\FabiusClampedQuantile$ is strictly increasing on $[0,1]$, monotone on $\RealNumbers$, and
continuous on $\RealNumbers$.
\item The restriction $\FabiusClampedQuantile|_I$ is a bijection $I\to I$.
\end{enumerate}
\end{theorem}

\begin{proof}
The clamp maps $\RealNumbers$ continuously and monotonically into $I$.  The inverse
of an order isomorphism between compact real intervals is itself an order
isomorphism, hence is continuous and strictly increasing.  Their
composition is therefore continuous and monotone on the whole line and
takes values in $I$.  If $y\in I$, then $\clamp(y)=y$, so the usual
right-inverse identity for $F_I^{-1}$ gives $F(\FabiusClampedQuantile(y))=y$.  Likewise,
$F(x)\in I$ for $x\in I$, and the left-inverse identity gives
$\FabiusClampedQuantile(F(x))=x$.  Strict increase and bijectivity on $I$ are inherited
from $F_I^{-1}$.
\end{proof}

The totalization is not cosmetic.  It produces a globally continuous
monotone function and makes the reflection law valid without restricting
the variable to $[0,1]$.

\subsection{A four-way comparison calculus}

The order-isomorphism viewpoint turns inverse inequalities into direct
inequalities without any numerical inversion.

\begin{theorem}[Comparison equivalences; \Formalized]
\label{inverse:thm:inverse-comparisons}
For $x,y\in[0,1]$,
\begin{align}
 \FabiusClampedQuantile(y)\leq x&\iff y\leq F(x), \label{inverse:eq:inv-le}\\
 F(x)\leq y&\iff x\leq \FabiusClampedQuantile(y), \label{inverse:eq:F-le}\\
 \FabiusClampedQuantile(y)<x&\iff y<F(x), \label{inverse:eq:inv-lt}\\
 F(x)<y&\iff x<\FabiusClampedQuantile(y). \label{inverse:eq:F-lt}
\end{align}
\end{theorem}

\begin{proof}
Apply the strictly increasing function $F$ to the two sides of an
inequality involving $\FabiusClampedQuantile(y)$, then use $F(\FabiusClampedQuantile(y))=y$.  For example,
\[
 \FabiusClampedQuantile(y)\leq x
 \iff F(\FabiusClampedQuantile(y))\leq F(x)
 \iff y\leq F(x).
\]
The other three equivalences are identical, with the order reversed in
placement but not in direction because $F$ and $\FabiusClampedQuantile$ are increasing.
\end{proof}

\begin{remark}
The equivalences are often more useful than an explicit inverse formula.
They let one convert a known dyadic or asymptotic bound for $F$ into a
location bound for $\FabiusClampedQuantile$ with no loss in constants.
\end{remark}

\subsection{Tails, endpoints, and reflection}

\begin{theorem}[Constant tails and symmetry; \Formalized]
\label{inverse:thm:inverse-symmetry}
For every $y\in\RealNumbers$,
\begin{align}
 y\leq0&\Longrightarrow \FabiusClampedQuantile(y)=0,\\
 y\geq1&\Longrightarrow \FabiusClampedQuantile(y)=1,\\
 \FabiusClampedQuantile(1-y)&=1-\FabiusClampedQuantile(y). \label{inverse:eq:inv-reflection}
\end{align}
In particular,
\[
 \FabiusClampedQuantile(0)=0,\qquad \FabiusClampedQuantile(1)=1,\qquad \FabiusClampedQuantile(1/2)=1/2.
\]
\end{theorem}

\begin{proof}
The two tail identities follow immediately from the clamp.  Suppose first
that $0\leq y\leq1$.  Then $\FabiusClampedQuantile(y)\in[0,1]$, and by the symmetry of $F$,
\[
 F(1-\FabiusClampedQuantile(y))
   =1-F(\FabiusClampedQuantile(y))
   =1-y.
\]
The point $1-\FabiusClampedQuantile(y)$ lies in $[0,1]$, so uniqueness of the inverse gives
$\FabiusClampedQuantile(1-y)=1-\FabiusClampedQuantile(y)$.  If $y\leq0$ or $y\geq1$, the same equation follows
from the two constant tails.  Setting $y=0,1,$ and $1/2$ gives the stated
special values.
\end{proof}

\section{Exact inversion of the dyadic evaluator}
\label{inverse:sec:dyadic-inverse}

\subsection{The exact statement}

Let $\mathsf F_n(a)$ denote the exact bounded dyadic evaluator from the
formal development: for $0\leq a\leq2^n$,
\[
 \mathsf F_n(a)=F\!\left(\frac{a}{2^n}\right),
\]
and for larger $a$ it is clamped at the endpoint value $1$.

\begin{theorem}[Dyadic inversion; \Formalized]
\label{inverse:thm:dyadic-inversion}
For all $n,a\in\NaturalNumbers$,
\[
 \FabiusClampedQuantile\bigl(\mathsf F_n(a)\bigr)
 =\frac{\min\{a,2^n\}}{2^n}.
\]
In particular, if $0\leq a\leq2^n$, then
\[
 \FabiusClampedQuantile\!\left(F\!\left(\frac a{2^n}\right)\right)=\frac a{2^n}.
\]
\end{theorem}

\begin{proof}
When $a\leq2^n$, the argument $a/2^n$ lies in $I$, so the left-inverse
identity in \cref{inverse:thm:inverse-package} applies.  When $a\geq2^n$, the
bounded evaluator equals $1$, the minimum equals $2^n$, and both sides are
$1$.
\end{proof}

The theorem is stronger than a table of examples: every exact rational
produced by the dyadic algorithm immediately gives an exact algebraic
input for the inverse.

\subsection{A small exact table}

\begin{table}[ht]
\centering
\small
\begin{tabular}{@{}ccc@{}}
\toprule
$x$ & $F(x)$ & exact inverse statement\\
\midrule
$0$ & $0$ & $\FabiusClampedQuantile(0)=0$\\
$\dfrac1{16}$ & $\dfrac{143}{2073600}$
  & $\FabiusClampedQuantile\!\left(\dfrac{143}{2073600}\right)=\dfrac1{16}$\\[0.55em]
$\dfrac18$ & $\dfrac1{288}$
  & $\FabiusClampedQuantile\!\left(\dfrac1{288}\right)=\dfrac18$\\[0.55em]
$\dfrac14$ & $\dfrac5{72}$
  & $\FabiusClampedQuantile\!\left(\dfrac5{72}\right)=\dfrac14$\\[0.55em]
$\dfrac38$ & $\dfrac{73}{288}$
  & $\FabiusClampedQuantile\!\left(\dfrac{73}{288}\right)=\dfrac38$\\[0.55em]
$\dfrac12$ & $\dfrac12$
  & $\FabiusClampedQuantile\!\left(\dfrac12\right)=\dfrac12$\\[0.55em]
$\dfrac58$ & $\dfrac{215}{288}$
  & $\FabiusClampedQuantile\!\left(\dfrac{215}{288}\right)=\dfrac58$\\[0.55em]
$\dfrac34$ & $\dfrac{67}{72}$
  & $\FabiusClampedQuantile\!\left(\dfrac{67}{72}\right)=\dfrac34$\\[0.55em]
$\dfrac78$ & $\dfrac{287}{288}$
  & $\FabiusClampedQuantile\!\left(\dfrac{287}{288}\right)=\dfrac78$\\[0.55em]
$\dfrac{15}{16}$ & $\dfrac{2073457}{2073600}$
  & $\FabiusClampedQuantile\!\left(\dfrac{2073457}{2073600}\right)=\dfrac{15}{16}$\\
\bottomrule
\end{tabular}
\caption{Exact inverse values obtained from the dyadic evaluator and
reflection.}
\label{inverse:tab:inverse-dyadic}
\end{table}

\begin{remark}
The threshold $F(1/4)=5/72$ will reappear in the effective endpoint
steepness theorem.  Its occurrence is structural rather than arbitrary:
$1/4=2^{-2}$ is precisely the scale on which the quadratic flatness bound
is valid.
\end{remark}

\section{\texorpdfstring{Flatness of \(F\) becomes steepness of \(\FabiusClampedQuantile\)}{Flatness of F becomes steepness of its inverse}}
\label{inverse:sec:inverse-steepness}

\subsection{The transported power inequalities}

Write
\[
 C_n=2^{\binom{n+1}{2}}.
\]
The effective flatness theorem for $F$ says, in one of its standard forms,
that
\[
 0\leq x\leq2^{-n}
 \quad\Longrightarrow\quad
 F(x)\leq C_nx^n. \label{inverse:eq:flatness-input}
\]
The sharp version improves the right-hand side by the factor $n!$.

\begin{theorem}[Inverted effective and sharp flatness; \Formalized]
\label{inverse:thm:inverted-flatness}
Let $n\in\NaturalNumbers$ and $0\leq y\leq1$.  If
\[
 2^n\FabiusClampedQuantile(y)\leq1,
\]
then
\begin{align}
 y&\leq C_n\FabiusClampedQuantile(y)^n, \label{inverse:eq:inverse-effective}\\
 y&\leq \frac{C_n}{n!}\FabiusClampedQuantile(y)^n. \label{inverse:eq:inverse-sharp}
\end{align}
Moreover, the scale condition follows from the argument condition
\[
 0\leq y\leq F(2^{-n}).
\]
\end{theorem}

\begin{proof}
Put $x=\FabiusClampedQuantile(y)$.  The scale hypothesis says $0\leq x\leq2^{-n}$, while
$F(x)=y$.  Substituting $x$ into the effective and sharp flatness bounds
for $F$ gives \eqref{inverse:eq:inverse-effective} and
\eqref{inverse:eq:inverse-sharp}.

Now suppose $y\leq F(2^{-n})$.  By
\eqref{inverse:eq:inv-le}, $\FabiusClampedQuantile(y)\leq2^{-n}$, which is exactly the required
scale condition.
\end{proof}

For $n\geq1$, taking positive $n$th roots yields the more geometric form.

\begin{corollary}[Root lower bounds; \Formalized\ input, \Derived\ form]
\label{inverse:cor:root-lower}
If $n\geq1$, $0<y\leq F(2^{-n})$, then
\begin{align}
 \FabiusClampedQuantile(y)&\geq C_n^{-1/n}y^{1/n}, \label{inverse:eq:root-effective}\\
 \FabiusClampedQuantile(y)&\geq
 \left(\frac{n!}{C_n}\right)^{1/n}y^{1/n}. \label{inverse:eq:root-sharp}
\end{align}
\end{corollary}

\begin{proof}
All quantities are nonnegative.  Rearrange
\cref{inverse:eq:inverse-effective,inverse:eq:inverse-sharp} and apply the monotonicity of
the positive $n$th-root function.
\end{proof}

\subsection{An effective bound against every line}

The quadratic case already detects infinite endpoint slope.

\begin{theorem}[Effective linear steepness; \Formalized]
\label{inverse:thm:effective-linear}
Let $M,y\in\RealNumbers$.  If
\[
 0<y,\qquad y\leq F(1/4)=\frac5{72},\qquad
 8M^2y<1,
\]
then
\[
 My<\FabiusClampedQuantile(y).
\]
\end{theorem}

\begin{proof}
The comparison theorem gives $\FabiusClampedQuantile(y)\leq1/4$, so the quadratic inverse
flatness estimate applies:
\[
 y\leq8\FabiusClampedQuantile(y)^2. \label{inverse:eq:quadratic-inv}
\]
If $M\leq0$, then $My\leq0<\FabiusClampedQuantile(y)$ because $y>0$ and $\FabiusClampedQuantile$ is strictly
increasing on $[0,1]$.  Suppose $M>0$ and, toward a contradiction,
$\FabiusClampedQuantile(y)\leq My$.  Then
\[
 y\leq8\FabiusClampedQuantile(y)^2\leq8M^2y^2.
\]
Dividing by $y>0$ gives $1\leq8M^2y$, contrary to the last hypothesis.
\end{proof}

\subsection{Infinite one-sided slope}

\begin{theorem}[Endpoint quotient divergence; \Formalized]
\label{inverse:thm:inverse-quotient}
As $y\downarrow0$,
\[
 \frac{\FabiusClampedQuantile(y)}{y}\longrightarrow+\infty.
\]
By reflection, as $y\uparrow1$,
\[
 \frac{1-\FabiusClampedQuantile(y)}{1-y}\longrightarrow+\infty.
\]
\end{theorem}

\begin{proof}
Fix a real target $A$.  Put $B=|A|+1$.  For all sufficiently small
positive $y$ we have simultaneously
\[
 y<F(1/4),\qquad 8B^2y<1.
\]
By \cref{inverse:thm:effective-linear}, $By<\FabiusClampedQuantile(y)$, and hence
\[
 A\leq B<\frac{\FabiusClampedQuantile(y)}y.
\]
This is exactly the filter definition of divergence to $+\infty$.

For the right endpoint use \eqref{inverse:eq:inv-reflection}:
\[
 \frac{1-\FabiusClampedQuantile(y)}{1-y}
 =\frac{\FabiusClampedQuantile(1-y)}{1-y},
\]
and let $1-y\downarrow0$.
\end{proof}

\subsection{Failure of every positive Hölder estimate}

\begin{theorem}[Stronger than non-Lipschitz; \Derived]
\label{inverse:thm:no-holder}
For every $\alpha>0$,
\[
 \frac{\FabiusClampedQuantile(y)}{y^\alpha}\longrightarrow+\infty
 \qquad (y\downarrow0).
\]
Consequently there are no $C,\delta>0$ for which
\[
 \FabiusClampedQuantile(y)\leq Cy^\alpha\qquad(0\leq y\leq\delta).
\]
The reflected assertion holds at $1$:
\[
 \frac{1-\FabiusClampedQuantile(y)}{(1-y)^\alpha}\longrightarrow+\infty
 \qquad (y\uparrow1).
\]
\end{theorem}

\begin{proof}
Choose an integer $n>1/\alpha$.  On the window
$0<y\leq F(2^{-n})$, \cref{inverse:eq:root-effective} gives
\[
 \frac{\FabiusClampedQuantile(y)}{y^\alpha}
 \geq C_n^{-1/n}y^{1/n-\alpha}.
\]
The exponent $1/n-\alpha$ is negative, so the right-hand side tends to
$+\infty$.  A local Hölder upper bound would keep the quotient bounded and
is therefore impossible.  The statement at $1$ follows from reflection.
\end{proof}

\begin{boxedremark}
Each fixed $n$ supplies a root lower bound only on the shrinking window
$y\leq F(2^{-n})$.  There is no bound uniform in $n$.  This is not a
defect: it is exactly the scale dependence needed to prove that
$\FabiusClampedQuantile$ outruns \emph{every fixed} positive power.
\end{boxedremark}


\section{Classical calculus and the exact regularity of the inverse}
\label{inverse:sec:inverse-calculus}

This section uses the formalized inverse package together with standard
one-variable analysis.  Smoothness, the reciprocal first-derivative formula,
the exact second-derivative formula, and the strict closed-half shape are now
marked \Formalized; the general higher inverse-derivative formulas remain
\Derived.

\subsection{Smoothness and differential identities in the interior}

\begin{theorem}[Interior inverse calculus through first order; \Formalized]
\label{inverse:thm:inverse-calculus}
The inverse $\FabiusClampedQuantile$ is $C^\infty$ on $(0,1)$.  For $0<y<1$,
\begin{equation}
 \FabiusClampedQuantile'(y)
   =\frac1{F'(\FabiusClampedQuantile(y))}
     =\frac1{2\RvachevUp(2\FabiusClampedQuantile(y)-1)}>0.
 \label{inverse:eq:Gprime}
\end{equation}
\end{theorem}

\begin{proof}
For $x\in(0,1)$, \eqref{inverse:eq:Fprime-up} and positivity of $\RvachevUp$ on
$(-1,1)$ imply $F'(x)>0$.  The $C^\infty$ inverse function theorem
therefore applies at every interior point and gives a smooth local inverse.
Uniqueness identifies all these local inverses with $\FabiusClampedQuantile$.  Differentiating
$F(\FabiusClampedQuantile(y))=y$ and substituting \eqref{inverse:eq:Fprime-up} proves
\eqref{inverse:eq:Gprime}.
\end{proof}

The corresponding Lean declarations are \path{fabiusInv_contDiffOn_Ioo},
\path{fabiusInv_hasDerivAt},
\path{deriv_fabiusInv_eq_inv_two_mul_rvachevUp}, and
\path{deriv_fabiusInv_pos}.

\begin{remark}[Current global smoothness locus; \Formalized]
The totalized inverse has locally constant tails.  The newer declarations
\path{fabiusInv_differentiableAt_iff} and
\path{fabiusInv_contDiffAt_infty_iff} sharpen the interior statement: the
inverse is differentiable, and indeed \(C^\infty\), exactly at the points
\(y\ne0,1\).  At each clamping endpoint it is not differentiable.  More
generally, \path{fabiusInv_contDiffAt_iff} says that order-zero smoothness is
global, whereas every positive finite order and \(C^\infty\) have precisely
those two exceptions.  Its order side condition deliberately excludes
Mathlib's stronger analytic order; the analytic locus is formalized separately by
\path{fabiusInv_analyticAt_iff} in \path{InverseNotElementary.lean}.
\end{remark}

\begin{proposition}[Second inverse derivative; \Formalized]
\label{inverse:prop:inverse-second-derivative}
For $0<y<1$,
\begin{equation}
 \FabiusClampedQuantile''(y)
   =-\frac{F''(\FabiusClampedQuantile(y))}
           {F'(\FabiusClampedQuantile(y))^3}.
 \label{inverse:eq:Gsecond}
\end{equation}
\end{proposition}

\begin{proof}
Differentiating the identity behind \eqref{inverse:eq:Gprime} once more gives
\[
 F''(\FabiusClampedQuantile(y))\FabiusClampedQuantile'(y)^2+
 F'(\FabiusClampedQuantile(y))\FabiusClampedQuantile''(y)=0.
\]
Substitution of \eqref{inverse:eq:Gprime} yields \eqref{inverse:eq:Gsecond}.  This is
\path{deriv_fabiusInv_hasDerivAt} and its pointwise corollary
\path{deriv_deriv_fabiusInv}.
\end{proof}

\begin{remark}[Higher inverse derivatives; \Derived]
Repeated differentiation expresses every higher derivative of $\FabiusClampedQuantile$ as a
rational polynomial in the derivatives of $F$ evaluated at $\FabiusClampedQuantile(y)$, with a
power of $F'(\FabiusClampedQuantile(y))$ in the denominator.  These general Fa\`a di Bruno
inverse-derivative polynomials are not separately packaged in the Lean API.
\end{remark}

\begin{corollary}[Midpoint differential data; \Formalized]
\label{inverse:cor:midpoint-inverse}
At the fixed point $1/2$,
\[
 \FabiusClampedQuantile(1/2)=1/2,\qquad
 \FabiusClampedQuantile'(1/2)=\frac12,\qquad
 \FabiusClampedQuantile''(1/2)=0.
\]
\end{corollary}

\begin{proof}
The value follows from reflection.  Since $\RvachevUp(0)=1$,
\eqref{inverse:eq:Gprime} gives $\FabiusClampedQuantile'(1/2)=1/2$.  The up-function is even, so
$\RvachevUp'(0)=0$, hence $F''(1/2)=4\RvachevUp'(0)=0$; now use
\eqref{inverse:eq:Gsecond}.  The three exact Lean declarations are
\path{fabiusInv_half}, \path{deriv_fabiusInv_half}, and
\path{deriv_deriv_fabiusInv_half}.
\end{proof}

\begin{proposition}[Exact interior curvature loci; \Formalized]
For \(0<y<1\),
\[
 \FabiusClampedQuantile''(y)<0\iff0<y<1/2,\qquad
 \FabiusClampedQuantile''(y)>0\iff1/2<y<1,
\]
and
\[
 \FabiusClampedQuantile''(y)=0\iff y=1/2.
\]
These are \path{deriv_deriv_fabiusInv_neg_iff},
\path{deriv_deriv_fabiusInv_pos_iff}, and
\path{deriv_deriv_fabiusInv_eq_zero_iff}.  Their interior hypothesis is
part of the statement: \(\FabiusClampedQuantile\) is not differentiable at the clamping points.
\end{proposition}

\subsection{Shape reversal under inversion}

The formal development proves that $F$ is strictly convex on
$(0,1/2)$ and strictly concave on $(1/2,1)$.  Inversion reverses these
two curvature signs.

\begin{proposition}[Concavity and convexity of the inverse; \Formalized]
\label{inverse:prop:inverse-shape}
The function $\FabiusClampedQuantile$ is strictly concave on $[0,1/2]$ and strictly convex
on $[1/2,1]$.
\end{proposition}

\begin{proof}
The inverse is strictly increasing and preserves the midpoint.  It therefore
maps the first half into itself, where $F'$ is strictly increasing and
positive.  Its reciprocal in \eqref{inverse:eq:Gprime} is strictly decreasing, so the
derivative criterion and continuity give strict concavity on the closed first
half.  On the second half $F'$ is strictly decreasing and positive, hence the
reciprocal derivative of the inverse is strictly increasing and the same
criterion gives strict convexity.

The closed-interval conclusions are
\path{strictConcaveOn_fabiusInv_firstHalf} and
\path{strictConvexOn_fabiusInv_secondHalf}.  Their derivative criteria inspect
only the two open interval interiors, while global continuity supplies the
endpoints; no finite endpoint derivative is asserted.
\end{proof}

\begin{remark}[Interior derivative blow-up; \Formalized]
The secant slopes diverge by \cref{inverse:thm:inverse-quotient}.  More
strongly, the positive interior derivative satisfies
\[
 \FabiusClampedQuantile'(y)\longrightarrow+\infty\quad(y\downarrow0),\qquad
 \FabiusClampedQuantile'(y)\longrightarrow+\infty\quad(y\uparrow1).
\]
These are \path{tendsto_deriv_fabiusInv_atTop_at_zero_right} and
\path{tendsto_deriv_fabiusInv_atTop_at_one_left}.  They are one-sided limits
through the open interval, not derivative values at the nondifferentiable
clamping points.
\end{remark}

\subsection{\texorpdfstring{The \(C^\infty\) and analytic loci}{The smooth and analytic loci}}

\begin{theorem}[Exact regularity locus; \Formalized]
\label{inverse:thm:inverse-regularity}
The totalized inverse has the following regularity.
\begin{enumerate}
\item $\FabiusClampedQuantile$ is $C^\infty$ precisely on
\[
 \RealNumbers\setminus\{0,1\}.
\]
\item It has no finite derivative at $0$ or at $1$.
\item It is real analytic precisely on
\[
 \RealNumbers\setminus[0,1].
\]
\end{enumerate}
\end{theorem}

\begin{proof}
On $(-\infty,0)$ and $(1,\infty)$ the function is constant, hence smooth
and analytic.  On $(0,1)$ smoothness follows from
\cref{inverse:thm:inverse-calculus}.  At $0$, the left derivative is $0$ because
the function is constant to the left, whereas the right difference
quotient $\FabiusClampedQuantile(y)/y$ tends to $+\infty$.  Thus no finite derivative
exists.  Reflection gives the same conclusion at $1$.

It remains to exclude analyticity in the open unit interval.  Suppose
$\FabiusClampedQuantile$ were analytic near some $y_0\in(0,1)$.  Its derivative there is
nonzero by \eqref{inverse:eq:Gprime}; the analytic inverse function theorem would
then make its local inverse $F$ analytic near $x_0=\FabiusClampedQuantile(y_0)$.  But the
Fabius function is nowhere analytic on the interior of its transition
interval.  This contradiction proves the claim.

The exact smoothness and differentiability classifications are
\path{fabiusInv_contDiffAt_infty_iff} and
\path{fabiusInv_differentiableAt_iff}; the analytic classification is
\path{fabiusInv_analyticAt_iff}.
\end{proof}

\subsection{Analytic-density and non-elementarity consequences}

The generic analytic-locus obstruction is proved in
\cref{integration:sec:elementary,integration:sec:not-elementary}; applying it to the exact
locus above gives a fully formal inverse conclusion.  The theorem
\path{not_eqOn_fabiusInv_of_dense_analyticLocus} excludes agreement with \(\FabiusClampedQuantile\) on any
subset of \([0,1]\) having nonempty interior whenever the comparison function has dense
analytic locus.  Thus \path{fabiusInv_not_isElementary} proves that the totalized inverse
is not elementary.  The stronger
\path{not_isElementaryOrInverse_eqOn_fabiusInv} and
\path{fabiusInv_not_isElementaryOrInverse} reach the same conclusion for the class generated
by elementary functions and admissible continuous inverse branches.  The branch constructor
requires analyticity of the chosen total extension on the interior of the complement of its
open branch domain; its conditional Lambert-style criterion is not an instantiated Mathlib
Lambert \(W\) declaration.  These domain qualifications are essential at branch boundaries.

\section{Inverting the sharp small-argument asymptotic}
\label{inverse:sec:inverse-asymptotic}

The power inequalities show qualitative infinite steepness.  The sharp
Lambert expansion gives the actual scale.  This section extracts that
scale without needing the detailed value of the periodic fluctuation.

\subsection{The exact lower Lambert phase}

Put
\[
 L=\log2.
\]
For sufficiently small $x>0$, let $\lambda=\lambda(x)$ be the large
positive solution of
\begin{equation}
 x=\lambda2^{-\lambda}
   =\lambda\EulerE^{-L\lambda}. \label{inverse:eq:lambert-phase}
\end{equation}
Equivalently,
\[
 \lambda(x)=-\frac1L W_{-1}(-Lx),
\]
but \eqref{inverse:eq:lambert-phase} is the identity used in the proofs.

The sharp expansion established for $F$ has the form
\begin{equation}
 \log F(x)
 =-\frac L2\lambda^2+
   \left(1+\frac L2\right)\lambda
   -\frac12\log\lambda
 +C_F+\Psi(\lambda)
   +\BigO(\lambda^{-1}), \label{inverse:eq:sharp-input}
\end{equation}
This is the general-real endpoint statement
\eqref{dyadicasym:eq:general-x-main}, restated here in the inverse notation.
Here $\Psi$ is smooth, bounded, and one-periodic.  The exact value of
$C_F$ and the Fourier description of $\Psi$ are important for $F$ itself,
but surprisingly they do not affect the leading equivalent of the
inverse.

\subsection{Asymptotic solution for the phase}

Let $y=F(x)$ and
\[
 T=\log\frac1y.
\]
Negating \eqref{inverse:eq:sharp-input} gives
\begin{equation}
 T=\frac L2\lambda^2-B\lambda+\frac12\log\lambda
     -C_F-\Psi(\lambda)+\BigO(\lambda^{-1}),
 \qquad B=1+\frac L2. \label{inverse:eq:T-lambda}
\end{equation}

\begin{lemma}[Asymptotic inversion of the quadratic phase; \Derived]
\label{inverse:lem:lambda-from-T}
As $T\to+\infty$,
\[
 \lambda
 =\sqrt{\frac{2T}{L}}+\frac BL
   +\BigO\!\left(\frac{\log T}{\sqrt T}\right).
\]
\end{lemma}

\begin{proof}
Since $\Psi$ is bounded, \eqref{inverse:eq:T-lambda} first gives
\[
 T=\frac L2\lambda^2+\BigO(\lambda+\log\lambda).
\]
Hence $\lambda\to\infty$ and, with
\[
 q=\sqrt{\frac{2T}{L}},
\]
we have $\lambda/q\to1$.  Dividing
\[
 T-\frac L2\lambda^2
 =-\frac L2(\lambda-q)(\lambda+q)
\]
by $\lambda+q\asymp q$ and using the preceding error estimate improves
this to $\lambda=q+\BigO(1)$.

Write
\[
 \lambda=q+\frac BL+\rho.
\]
Insert this into \eqref{inverse:eq:T-lambda}.  The terms proportional to $q$
cancel because $L(B/L)q=Bq$.  What remains is
\[
 Lq\rho
 =-\frac12\log\!\left(q+\frac BL+\rho\right)+\BigO(1)
   +\BigO(\rho^2).
\]
The already known bound $\rho=\BigO(1)$ makes the final term bounded.
Since $q\asymp\sqrt T$ and $\log q\asymp\log T$, division by $Lq$ gives
\[
 \rho=\BigO\!\left(\frac{\log T}{\sqrt T}\right).
\]
\end{proof}

\subsection{The endpoint equivalent for the inverse}

\begin{theorem}[Sharp inverse asymptotic; \Derived]
\label{inverse:thm:sharp-inverse-asymptotic}
As $y\downarrow0$,
\begin{equation}
 \boxed{\;
 \FabiusClampedQuantile(y)\sim
 \frac{\sqrt{\log(1/y)}}{\EulerE\sqrt{\log2}}\,
 \exp\!\left[-\sqrt{2\log2\,\log(1/y)}\right].
 \;} \label{inverse:eq:sharp-inverse}
\end{equation}
Equivalently,
\begin{equation}
 \log\FabiusClampedQuantile(y)
 =-\sqrt{2L\log(1/y)}
  +\frac12\log\log(1/y)
  -1-\frac12\log L+o(1). \label{inverse:eq:log-sharp-inverse}
\end{equation}
At the right endpoint,
\begin{equation}
 1-\FabiusClampedQuantile(y)\sim
 \frac{\sqrt{\log(1/(1-y))}}{\EulerE\sqrt{\log2}}\,
 \exp\!\left[-\sqrt{2\log2\,\log(1/(1-y))}\right]
 \quad(y\uparrow1). \label{inverse:eq:right-inverse-asymptotic}
\end{equation}
\end{theorem}

\begin{proof}
Let $x=\FabiusClampedQuantile(y)$.  Then $y=F(x)$ and $x$ is related to its exact phase by
$x=\lambda\EulerE^{-L\lambda}$.  Therefore
\[
 \log x=\log\lambda-L\lambda.
\]
With $T=\log(1/y)$, \cref{inverse:lem:lambda-from-T} gives
\[
 \lambda=q+\frac BL+o(1),\qquad
 q=\sqrt{\frac{2T}{L}},
\]
because $(\log T)/\sqrt T=o(1)$.  Also $\log\lambda=\log q+o(1)$.
Consequently
\begin{align*}
 \log x
 &=\log q-Lq-B+o(1)\\
 &=-\sqrt{2LT}
   +\frac12\log T+\frac12\log\frac2L
   -1-\frac L2+o(1)\\
 &=-\sqrt{2LT}
   +\frac12\log T-1-\frac12\log L+o(1),
\end{align*}
because $L=\log2$.  This is \eqref{inverse:eq:log-sharp-inverse}.
Exponentiating proves \eqref{inverse:eq:sharp-inverse}.  Finally,
\eqref{inverse:eq:inv-reflection} gives
$1-\FabiusClampedQuantile(y)=\FabiusClampedQuantile(1-y)$ and hence
\eqref{inverse:eq:right-inverse-asymptotic}.
\end{proof}

\begin{remark}[Why the periodic fluctuation disappears at leading order]
The bounded term $\Psi(\lambda)$ changes the solution $\lambda$ of
\eqref{inverse:eq:T-lambda} only by $\BigO(1/\lambda)$.  Since
$\log x=\log\lambda-L\lambda$, this changes $\log x$ by $o(1)$ and
therefore changes $x$ by a factor tending to $1$.  Periodicity returns in
the next relative correction, but not in the leading equivalent
\eqref{inverse:eq:sharp-inverse}.
\end{remark}

\subsection{The complete hierarchy of elementary scales}

\begin{corollary}[Between powers and logarithms; \Derived]
\label{inverse:cor:inverse-scale-hierarchy}
For every $\alpha>0$ and every $m>0$,
\begin{align}
 y^\alpha&=o(\FabiusClampedQuantile(y)), \label{inverse:eq:power-below-G}\\
 \FabiusClampedQuantile(y)&=o\!\left((\log(1/y))^{-m}\right)
 \label{inverse:eq:G-below-log}
\end{align}
as $y\downarrow0$.
\end{corollary}

\begin{proof}
Set $T=\log(1/y)$.  Formula \eqref{inverse:eq:log-sharp-inverse} gives
\[
 \log\frac{\FabiusClampedQuantile(y)}{y^\alpha}
 =\alpha T-\sqrt{2LT}+\BigO(\log T)\longrightarrow+\infty,
\]
which proves \eqref{inverse:eq:power-below-G}.  Likewise,
\[
 \log\!\left(\FabiusClampedQuantile(y)T^m\right)
 =-\sqrt{2LT}+\BigO(\log T)\longrightarrow-\infty,
\]
which proves \eqref{inverse:eq:G-below-log}.
\end{proof}

\begin{boxedremark}
The inverse endpoint is therefore neither power-like nor logarithmic.
It lies in the stretched-logarithmic regime
$\exp(-c\sqrt{\log(1/y)})$ with an essential
$\sqrt{\log(1/y)}$ prefactor.
\end{boxedremark}

\section{Closed periodic saddle jets}
\label{inverse:sec:closed-jets}

We now turn from the inverse to the coefficient machinery hidden inside
the all-orders endpoint expansion.  Throughout this part,
\[
 L=\log2,\qquad D=\frac{\Differential}{\Differential u},
\]
and $\Psi$ is the smooth one-periodic function occurring in the sharp
Lambert main term.

\subsection{The recursive jets}

Define the periodic jets $p_n$ by
\begin{align}
 p_0(u)&=\frac12+\frac{\Psi'(u)}L, \label{inverse:eq:p0}\\
 p_{n+1}(u)
 &=\left(\frac DL-(n+1)\right)p_n(u)
   +\frac{(-1)^{n+1}n!}{L}. \label{inverse:eq:p-recurrence}
\end{align}
The bounded exponent jets are
\begin{equation}
 d_n(u)=p_n(u)+(-1)^{n+1}n!. \label{inverse:eq:d-from-p}
\end{equation}
The subtraction in \eqref{inverse:eq:d-from-p} removes the factorial part that
would otherwise obscure the centered saddle expansion.

For $n\geq0$, write
\[
 H_n=\sum_{k=1}^n\frac1k,\qquad H_0=0,
\]
and define integers $c(n,m)$ by
\begin{equation}
 \prod_{k=1}^n(z-k)=\sum_{m=0}^n c(n,m)z^m.
 \label{inverse:eq:c-definition}
\end{equation}

\begin{theorem}[Closed periodic jets; \Formalized]
\label{inverse:thm:closed-jets}
For every $n\geq0$,
\begin{align}
 p_n(u)
 &=
 (-1)^n n!\left(\frac12+\frac{H_n}{L}\right)
 +\sum_{m=0}^n
   c(n,m)\frac{\Psi^{(m+1)}(u)}{L^{m+1}},
 \label{inverse:eq:p-closed}\\
 d_n(u)
 &=
 (-1)^n n!\left(\frac{H_n}{L}-\frac12\right)
 +\sum_{m=0}^n
   c(n,m)\frac{\Psi^{(m+1)}(u)}{L^{m+1}}.
 \label{inverse:eq:d-closed}
\end{align}
\end{theorem}

\begin{proof}
Let
\[
 T_k=\frac DL-k.
\]
The operators $T_k$ commute because each is a polynomial in $D$.  Iterating
\eqref{inverse:eq:p-recurrence} gives
\begin{equation}
 p_n=T_nT_{n-1}\cdots T_1p_0
 +\sum_{j=0}^{n-1}
   T_nT_{n-1}\cdots T_{j+2}
   \left(\frac{(-1)^{j+1}j!}{L}\right),
 \label{inverse:eq:p-iteration}
\end{equation}
where the product is empty when $j=n-1$.

First consider the homogeneous term.  Since
\[
 T_n\cdots T_1
 =\prod_{k=1}^n\left(\frac DL-k\right),
\]
its action on the constant $1/2$ is
\[
 \frac12\prod_{k=1}^n(-k)=\frac{(-1)^nn!}{2}.
\]
Its action on $\Psi'/L$ is found by expanding the operator polynomial with
\eqref{inverse:eq:c-definition}:
\[
 \prod_{k=1}^n\left(\frac DL-k\right)\frac{\Psi'}L
 =\sum_{m=0}^n
   c(n,m)\frac{\Psi^{(m+1)}}{L^{m+1}}.
\]

Now propagate the forcing term inserted at stage $j$.  It is constant, so
each subsequent $T_k$ acts on it as multiplication by $-k$.  Hence its
contribution at stage $n$ is
\begin{align*}
 \frac{(-1)^{j+1}j!}{L}
 \prod_{k=j+2}^n(-k)
 &=
 \frac{(-1)^{j+1}j!}{L}
 \frac{(-1)^{n-j-1}n!}{(j+1)!}\\
 &=\frac{(-1)^nn!}{(j+1)L}.
\end{align*}
Summing over $j=0,\dots,n-1$ gives
\[
 \frac{(-1)^nn!}{L}\sum_{j=0}^{n-1}\frac1{j+1}
 =\frac{(-1)^nn!H_n}{L}.
\]
Combining the constant, derivative, and forcing contributions proves
\eqref{inverse:eq:p-closed}.  Finally,
\[
 p_n+(-1)^{n+1}n!
 =(-1)^nn!\left(\frac12+\frac{H_n}{L}-1\right)
   +\text{derivative part},
\]
which is exactly \eqref{inverse:eq:d-closed}.
\end{proof}

\subsection{Why the Stirling indices are shifted}

Let $s(n,k)$ be the signed Stirling number of the first kind, defined by
\[
 z(z-1)\cdots(z-n+1)=\sum_{k=0}^n s(n,k)z^k.
\]

\begin{proposition}[Shifted Stirling conversion; \Formalized]
\label{inverse:prop:stirling-shift}
For $0\leq m\leq n$,
\[
 c(n,m)=s(n+1,m+1).
\]
\end{proposition}

\begin{proof}
Apply the defining identity with $n+1$:
\[
 z(z-1)\cdots(z-n)
 =\sum_{k=0}^{n+1}s(n+1,k)z^k.
\]
Both sides have zero constant term, so divide by $z$:
\[
 (z-1)\cdots(z-n)
 =\sum_{m=0}^n s(n+1,m+1)z^m.
\]
Comparison with \eqref{inverse:eq:c-definition} proves the formula.
\end{proof}

\begin{warning}
The coefficients are not $s(n,m)$.  The product starts with $(z-1)$
rather than $z$, and this forces the simultaneous shift
$(n,m)\mapsto(n+1,m+1)$.  In particular, the nonzero constant column
$c(n,0)=(-1)^nn!$ is essential.
\end{warning}

\subsection{Coefficient rows and the first jets}

The first rows of $c(n,m)$, listed from $m=0$ to $m=n$, are
\[
\begin{array}{c|rrrrrr}
n&c(n,0)&c(n,1)&c(n,2)&c(n,3)&c(n,4)&c(n,5)\\
\hline
0&1\\
1&-1&1\\
2&2&-3&1\\
3&-6&11&-6&1\\
4&24&-50&35&-10&1\\
5&-120&274&-225&85&-15&1
\end{array}
\]
and \eqref{inverse:eq:d-closed} gives
\begin{align}
 d_0={}&-\frac12+\frac{\Psi'}L, \label{inverse:eq:d0-explicit}\\
 d_1={}&\frac12-\frac1L-\frac{\Psi'}L+\frac{\Psi''}{L^2},
 \label{inverse:eq:d1-explicit}\\
 d_2={}&-1+\frac3L+\frac{2\Psi'}L
       -\frac{3\Psi''}{L^2}+\frac{\Psi'''}{L^3},
 \label{inverse:eq:d2-explicit}\\
 d_3={}&3-\frac{11}{L}-\frac{6\Psi'}L
       +\frac{11\Psi''}{L^2}-\frac{6\Psi'''}{L^3}
       +\frac{\Psi^{(4)}}{L^4}. \label{inverse:eq:d3-explicit}
\end{align}

\begin{corollary}[Structural properties; \Formalized]
\label{inverse:cor:jet-structure}
Every $p_n$ and $d_n$ is smooth, one-periodic, and bounded on $\RealNumbers$.
Moreover, $d_n$ involves derivatives of $\Psi$ only through order $n+1$,
and the coefficient of the highest derivative is $L^{-(n+1)}$.
\end{corollary}

\begin{proof}
The closed formulas are finite linear combinations of derivatives of the
smooth periodic function $\Psi$ and constants.  Smoothness, periodicity,
and boundedness follow.  Since $c(n,n)=1$, the coefficient of
$\Psi^{(n+1)}$ is $L^{-(n+1)}$.
\end{proof}


\section{The centered exponent and its algebraic collapse}
\label{inverse:sec:centered-exponent}

\subsection{From a long Taylor expansion to one short formula}

At the saddle, introduce the half-power parameter
\[
 \varepsilon=\lambda^{-1/2}
\]
and the Gaussian variable $v$.  After removing the exact quadratic term,
the centered exponent has a formal expansion
\[
 \mathcal E(u,v;\varepsilon)
 =\sum_{m\geq1}E_m(u,v)\varepsilon^m.
\]
The original construction of $E_m$ contains separate contributions from
the vertical negative-Laplace jet and from the Taylor expansion of
$\log(1+\ImaginaryUnit\varepsilon v)$.  The bounded exponent jets make these two
pieces collapse.

\begin{theorem}[Closed exponent coefficient; \Formalized]
\label{inverse:thm:closed-exponent}
For every $m\geq1$,
\begin{equation}
 E_m(u,v)=\ImaginaryUnit^m\left(
   \frac{d_{m-1}(u)}{m!}v^m
   +\frac{(-1)^{m+1}}{m+2}v^{m+2}
 \right). \label{inverse:eq:E-closed}
\end{equation}
In particular,
\begin{align}
 E_1&=\ImaginaryUnit v\,\frac{3d_0+v^2}{3}, \label{inverse:eq:E1}\\
 E_2&=v^2\,\frac{-2d_1+v^2}{4}, \label{inverse:eq:E2}\\
 E_3&=\ImaginaryUnit v^3\left(-\frac{d_2}{6}-\frac{v^2}{5}\right), \label{inverse:eq:E3}\\
 E_4&=\frac{v^4}{24}(d_3-4v^2). \label{inverse:eq:E4}
\end{align}
\end{theorem}

\begin{proof}
Before centering, the coefficient at order $m$ has two contributions.
The vertical jet contributes $\ImaginaryUnit^m p_{m-1}(u)v^m/m!$, while the
logarithmic saddle coordinate contributes a factorial term to the same
monomial and a universal $v^{m+2}$ term.  Using
\[
 d_{m-1}=p_{m-1}+(-1)^m(m-1)!
\]
combines the two $v^m$ coefficients into
$\ImaginaryUnit^m d_{m-1}/m!$.  The remaining logarithmic term is
\[
 \frac{\ImaginaryUnit^{m+2}(-1)^m}{m+2}v^{m+2}
 =\ImaginaryUnit^m\frac{(-1)^{m+1}}{m+2}v^{m+2},
\]
because $\ImaginaryUnit^2=-1$.  This is \eqref{inverse:eq:E-closed}.  Substitution of
$m=1,2,3,4$ gives \eqref{inverse:eq:E1}--\eqref{inverse:eq:E4}.
\end{proof}

\begin{corollary}[Parity; \Formalized]
\label{inverse:cor:E-parity}
For every $m\geq1$,
\[
 E_m(u,-v)=(-1)^mE_m(u,v).
\]
Thus the odd-indexed $E_m$ are odd polynomials in $v$ and the
even-indexed $E_m$ are even.
\end{corollary}

\begin{proof}
Both monomials in \eqref{inverse:eq:E-closed} have parity $m$, because $m+2$ has
the same parity as $m$.
\end{proof}

\section{Formal exponential, Gaussian contraction, and logarithm}
\label{inverse:sec:coefficient-formulas}

The generic exponential, logarithmic, and Gaussian engines---including their uniqueness,
naturality, parity, and finite-remainder theorems---are developed in
\cref{integration:sec:exp-coeff,integration:sec:log-coeff,integration:sec:gaussian-contraction}.
This section records only their Fabius-specific instantiation and the ordered-composition
formulas needed for the later explicit coefficients.

\subsection{The exponential recursion}

Define polynomials $B_n(u,v)$ by the formal identity
\begin{equation}
 \exp\!\left(\sum_{m\geq1}E_m(u,v)\varepsilon^m\right)
 =\sum_{n\geq0}B_n(u,v)\varepsilon^n.
 \label{inverse:eq:B-definition}
\end{equation}
Thus $B_0=1$.

\begin{proposition}[Finite exponential recursion; \Formalized]
\label{inverse:prop:exp-recursion}
For $n\geq1$,
\begin{equation}
 nB_n=\sum_{m=1}^n mE_mB_{n-m}. \label{inverse:eq:B-recursion}
\end{equation}
\end{proposition}

\begin{proof}
Apply the generic coefficient identity from \cref{integration:sec:exp-coeff} to the
Fabius exponent sequence \(E_m(u,v)\).
\end{proof}

\begin{corollary}[Ordered-composition exponential formula; \Derived]
\label{inverse:cor:exp-composition}
For \(n\geq1\),
\[
 B_n=\sum_{\substack{r\geq1\\m_1+\cdots+m_r=n}}
 \frac1{r!}E_{m_1}\cdots E_{m_r},
\]
where the inner sum is over ordered compositions of $n$ into $r$ positive
parts.
\end{corollary}

\begin{proof}[Proof of \cref{inverse:cor:exp-composition}]
Expand the finite coefficient of the formal exponential by the number of factors.  This is
an algebraic consequence of the recursion in each degree; no analytic convergence is used.
The repository does not currently package this ordered-composition statement as one general
Lean declaration.
\end{proof}

\begin{corollary}[Parity of the exponential coefficients; \Formalized]
\label{inverse:cor:B-parity}
For all $n\geq0$,
\[
 B_n(u,-v)=(-1)^nB_n(u,v).
\]
\end{corollary}

\begin{proof}
Every product contributing to $B_n$ has indices
$m_1+\cdots+m_r=n$.  By \cref{inverse:cor:E-parity}, its parity is
$(-1)^{m_1+\cdots+m_r}=(-1)^n$.
\end{proof}

\subsection{Gaussian contraction and mass coefficients}

For a polynomial $P$ define normalized Gaussian contraction by
\[
 \Gauss[P]
 =\frac1{\sqrt{2\pi}}\int_{\RealNumbers}\EulerE^{-v^2/2}P(v)\,\Differential v.
\]
It is determined by
\begin{equation}
 \Gauss[v^{2k}]=(2k-1)!!,\qquad
 \Gauss[v^{2k+1}]=0. \label{inverse:eq:gaussian-moments}
\end{equation}
The mass coefficients are
\begin{equation}
 a_j(u)=\Gauss[B_{2j}(u,\cdot)],\qquad j\geq0.
 \label{inverse:eq:mass-coeff-definition}
\end{equation}
In particular $a_0=1$.  The odd half-powers vanish automatically:
\[
 \Gauss[B_{2j+1}]=0
\]
by \cref{inverse:cor:B-parity}.

\begin{theorem}[Closed composition formula for \(a_j\); \Derived]
\label{inverse:thm:mass-composition}
For $j\geq1$,
\begin{align}
 a_j(u)
 ={}&(-1)^j
 \sum_{r=1}^{2j}\frac1{r!}
 \sum_{\substack{m_1+\cdots+m_r=2j\\m_i\geq1}}
 \sum_{S\subseteq\{1,\dots,r\}}
 (2j+2|S|-1)!! \notag\\
 &\qquad{}\times
 \prod_{i\in S}\frac{(-1)^{m_i+1}}{m_i+2}
 \prod_{i\notin S}\frac{d_{m_i-1}(u)}{m_i!}.
 \label{inverse:eq:mass-composition}
\end{align}
Consequently $a_j$ is a polynomial with rational coefficients in
$d_0,\dots,d_{2j-1}$.
\end{theorem}

\begin{proof}
Use the ordered-composition formula in
\cref{inverse:cor:exp-composition} for $B_{2j}$.  Each factor
\[
 E_{m_i}
 =\ImaginaryUnit^{m_i}\left(
   \frac{d_{m_i-1}}{m_i!}v^{m_i}
   +\frac{(-1)^{m_i+1}}{m_i+2}v^{m_i+2}
 \right)
\]
offers two choices.  Let $S$ be the set of factors from which the second,
universal monomial is chosen.  Since $\sum_i m_i=2j$, the product of the
powers of $\ImaginaryUnit$ is
\[
 \ImaginaryUnit^{2j}=(-1)^j,
\]
and the total power of $v$ is $2j+2|S|$.  Its Gaussian contraction is
$(2j+2|S|-1)!!$.  Multiplying the remaining scalar factors gives
\eqref{inverse:eq:mass-composition}.  The largest possible jet index is
$m_i-1\leq2j-1$.
\end{proof}

\begin{boxedremark}
Formula \eqref{inverse:eq:mass-composition} is a human-readable algebraic
consequence of the formal exponent and exponential-recursion theorems.
The repository machine-checks the underlying recursion, the structural
top-jet theorem, and explicit low orders.  It does not package this exact
threefold summation as a single general Lean declaration.
\end{boxedremark}

\subsection{Logarithmic coefficients}

Define $A_j(u)$ by
\begin{equation}
 1+\sum_{j\geq1}a_j(u)z^j
 =\exp\!\left(\sum_{j\geq1}A_j(u)z^j\right).
 \label{inverse:eq:log-coeff-definition}
\end{equation}
These are the coefficients that occur in $\log F$.

\begin{proposition}[Mass-to-logarithm conversion; \Formalized]
\label{inverse:prop:log-recursion}
For $n\geq1$,
\begin{equation}
 A_n
 =a_n-\frac1n\sum_{k=1}^{n-1}kA_k\,a_{n-k}.
 \label{inverse:eq:A-recursion}
\end{equation}
\end{proposition}

\begin{proof}
Specialize the mutually inverse formal logarithm/exponential transforms of
\cref{integration:sec:log-coeff}.  The recursive form is the coefficient identity for
\(A'=(1+\sum a_jz^j)'/(1+\sum a_jz^j)\).
\end{proof}

\begin{corollary}[Ordered-composition logarithm formula; \Derived]
\label{inverse:cor:log-composition}
For \(n\geq1\),
\begin{equation}
 A_n
 =\sum_{r=1}^n\frac{(-1)^{r-1}}r
   \sum_{\substack{q_1+\cdots+q_r=n\\q_i\geq1}}
   a_{q_1}\cdots a_{q_r}. \label{inverse:eq:A-compositions}
\end{equation}
\end{corollary}

\begin{proof}[Proof of \cref{inverse:cor:log-composition}]
Take the finite coefficient of \(\log(1+w)\).  This is an algebraic identity in each fixed
degree, but the displayed general composition sum is not currently exported as a single
Lean theorem.
\end{proof}

\begin{corollary}[Regularity and derivative order; \Formalized\ input,
\Derived\ conclusion]
\label{inverse:cor:A-structure}
For every $j$, the functions $a_j$ and $A_j$ are smooth, one-periodic,
and bounded.  Neither involves a derivative of $\Psi$ of order greater
than $2j$.
\end{corollary}

\begin{proof}
The function $a_j$ is a polynomial in
$d_0,\dots,d_{2j-1}$, and $d_n$ is smooth, periodic, and depends on
$\Psi$ only through order $n+1$.  Thus $a_j$ has the stated properties
and uses derivatives only through order $2j$.  The recursion
\eqref{inverse:eq:A-recursion} transfers the same assertions to $A_j$.
\end{proof}

\subsection{The top jet and the highest derivative}

The most complicated-looking coefficient has a very simple dependence on
its newest jet.

\begin{theorem}[Leading mass-jet law; \Formalized{} input, \Derived{} logarithmic consequence]
\label{inverse:thm:top-jet}
For $j\geq1$, the mass coefficient $a_j$ is affine in $d_{2j-1}$, and
the coefficient of that jet is
\begin{equation}
 \frac{(-1)^j}{2^j j!}. \label{inverse:eq:top-jet-coefficient}
\end{equation}
The logarithmic coefficient $A_j$ is derived here to have the same top-jet coefficient; this transfer is not asserted by the cited mass-coefficient theorem.
Consequently the coefficient of $\Psi^{(2j)}(u)$ in $A_j(u)$ is
\begin{equation}
 \boxed{\frac{(-1)^j}{2^j j!L^{2j}}.}
 \label{inverse:eq:highest-derivative}
\end{equation}
\end{theorem}

\begin{proof}
The jet $d_{2j-1}$ occurs only in $E_{2j}$.  Therefore the only term of
$B_{2j}$ that contains it is the linear term $E_{2j}$ itself.  From
\eqref{inverse:eq:E-closed}, its jet part is
\[
 \ImaginaryUnit^{2j}\frac{d_{2j-1}}{(2j)!}v^{2j}
 =(-1)^j\frac{d_{2j-1}}{(2j)!}v^{2j}.
\]
Gaussian contraction multiplies this by $(2j-1)!!$.  Since
\[
 (2j)!=2^jj!(2j-1)!!,
\]
the coefficient is \eqref{inverse:eq:top-jet-coefficient}.

In \eqref{inverse:eq:A-recursion}, every correction term
$A_ka_{j-k}$ has $k<j$ and $j-k<j$; it involves jets only through
$d_{2j-3}$ and therefore cannot alter the coefficient of $d_{2j-1}$.
Thus $A_j$ has the same top-jet coefficient.  Finally,
\cref{inverse:cor:jet-structure} says that the coefficient of
$\Psi^{(2j)}$ in $d_{2j-1}$ is $L^{-2j}$, giving
\eqref{inverse:eq:highest-derivative}.
\end{proof}

\begin{remark}
The theorem provides a quick consistency check for every explicit
coefficient.  For example, $A_1$ must contain
$-\Psi''/(2L^2)$, while $A_2$ must contain
$+\Psi^{(4)}/(8L^4)$.
\end{remark}

\section{The first and second logarithmic corrections}
\label{inverse:sec:second-correction}

\subsection{Order two: the first correction}

From \cref{inverse:eq:E1,inverse:eq:E2},
\[
 B_2=E_2+\frac12E_1^2.
\]
A direct expansion gives
\[
 B_2
 =-\frac{d_0^2+d_1}{2}v^2
   +\left(\frac14-\frac{d_0}{3}\right)v^4
   -\frac1{18}v^6.
\]
Using $\Gauss[v^2]=1$, $\Gauss[v^4]=3$, and
$\Gauss[v^6]=15$ yields
\begin{equation}
 a_1=A_1
 =-\frac{d_0^2}{2}-d_0-\frac{d_1}{2}-\frac1{12}.
 \label{inverse:eq:A1-jets}
\end{equation}
Equivalently,
\[
 A_1=-\frac{6d_0^2+12d_0+6d_1+1}{12}.
\]
Substituting \cref{inverse:eq:d0-explicit,inverse:eq:d1-explicit} recovers the first
periodic correction in terms of $\Psi'$ and $\Psi''$.

\subsection{Order four of the formal exponential}

The fourth exponential coefficient is
\begin{equation}
 B_4
 =E_4+E_3E_1+\frac12E_2^2
   +\frac12E_2E_1^2+\frac1{24}E_1^4.
 \label{inverse:eq:B4-identity}
\end{equation}
This follows either from the composition formula or from the recurrence
\eqref{inverse:eq:B-recursion}.  Substitution of
\cref{inverse:eq:E1,inverse:eq:E2,inverse:eq:E3,inverse:eq:E4} and collection of even powers of $v$
gives the following polynomial.

\begin{proposition}[The order-four reference polynomial; \Formalized]
\label{inverse:prop:B4-polynomial}
Writing $d_k=d_k(u)$,
\begin{align}
 B_4(u,v)={}&
 \left(\frac{d_0^4}{24}+\frac{d_0^2d_1}{4}
       +\frac{d_0d_2}{6}+\frac{d_1^2}{8}
       +\frac{d_3}{24}\right)v^4 \notag\\
 &+\left(\frac{d_0^3}{18}-\frac{d_0^2}{8}
       +\frac{d_0d_1}{6}+\frac{d_0}{5}
       -\frac{d_1}{8}+\frac{d_2}{18}
       -\frac16\right)v^6 \notag\\
 &+\left(\frac{d_0^2}{36}-\frac{d_0}{12}
       +\frac{d_1}{36}+\frac{47}{480}\right)v^8 \notag\\
 &+\left(\frac{d_0}{162}-\frac1{72}\right)v^{10}
   +\frac1{1944}v^{12}. \label{inverse:eq:B4-polynomial}
\end{align}
\end{proposition}

\begin{proof}
Let
\[
 P_1=d_0v+\frac{v^3}{3},\qquad
 P_3=-\frac{d_2v^3}{6}-\frac{v^5}{5},
\]
so that $E_1=\ImaginaryUnit P_1$ and $E_3=\ImaginaryUnit P_3$.  Then
$E_3E_1=-P_3P_1$ and $E_1^4=P_1^4$.  The even terms $E_2,E_4$ are already
real.  Insert these expressions in \eqref{inverse:eq:B4-identity}.  Each summand
has degree at most $12$ and is even.  Multiplication gives
\[
\begin{array}{c|ccccc}
 &v^4&v^6&v^8&v^{10}&v^{12}\\
\hline
E_4&
 d_3/24&-1/6&0&0&0\\
E_3E_1&
 d_0d_2/6&d_2/18+d_0/5&1/15&0&0\\
E_2^2/2&
 d_1^2/8&-d_1/8&1/32&0&0\\
E_2E_1^2/2&
 d_0^2d_1/4&
 d_0d_1/6-d_0^2/8&
 d_1/36-d_0/12&
 -1/72&0\\
E_1^4/24&
 d_0^4/24&d_0^3/18&d_0^2/36&d_0/162&1/1944
\end{array}
\]
and the two universal $v^8$ contributions satisfy
$1/15+1/32=47/480$.  Summing the columns proves
\eqref{inverse:eq:B4-polynomial}.
\end{proof}

\subsection{Gaussian contraction and the second mass coefficient}

The relevant Gaussian moments are
\[
 \Gauss[v^4]=3,\quad
 \Gauss[v^6]=15,\quad
 \Gauss[v^8]=105,\quad
 \Gauss[v^{10}]=945,\quad
 \Gauss[v^{12}]=10395.
\]

\begin{theorem}[Second mass coefficient; \Formalized]
\label{inverse:thm:a2}
Gaussian contraction of \eqref{inverse:eq:B4-polynomial} gives
\begin{align}
 a_2={}&
 \frac{d_0^4}{8}+\frac{5d_0^3}{6}
 +\frac{3d_0^2d_1}{4}+\frac{25d_0^2}{24}
 +\frac{5d_0d_1}{2}+\frac{d_0d_2}{2}
 +\frac{d_0}{12} \notag\\
 &+\frac{3d_1^2}{8}+\frac{25d_1}{24}
 +\frac{5d_2}{6}+\frac{d_3}{8}
 +\frac1{288}. \label{inverse:eq:a2}
\end{align}
\end{theorem}

\begin{proof}
Multiply the coefficients of $v^4,v^6,v^8,v^{10},v^{12}$ in
\eqref{inverse:eq:B4-polynomial} by $3,15,105,945,10395$, respectively, and
collect like monomials.  For illustration, the coefficient of $d_0$ is
\[
 15\cdot\frac15
 +105\left(-\frac1{12}\right)
 +945\cdot\frac1{162}
 =3-\frac{35}{4}+\frac{35}{6}
 =\frac1{12},
\]
and the universal constant is
\[
 15\left(-\frac16\right)
 +105\frac{47}{480}
 +945\left(-\frac1{72}\right)
 +10395\frac1{1944}
 =\frac1{288}.
\]
The remaining coefficients follow by the same arithmetic.
\end{proof}

\subsection{The second logarithmic coefficient}

Since
\[
 \log(1+a_1z+a_2z^2+\cdots)
 =a_1z+\left(a_2-\frac12a_1^2\right)z^2+\cdots,
\]
we have $A_2=a_2-a_1^2/2$.

\begin{theorem}[Second periodic saddle correction; \Formalized]
\label{inverse:thm:A2}
The second logarithmic coefficient is
\begin{equation}
 \boxed{
 A_2=
 \frac{
  8d_0^3+12d_0^2d_1+12d_0^2+48d_0d_1+12d_0d_2
  +6d_1^2+24d_1+20d_2+3d_3
 }{24}.}
 \label{inverse:eq:A2}
\end{equation}
It is smooth, one-periodic, and bounded.
\end{theorem}

\begin{proof}
Insert \eqref{inverse:eq:A1-jets} and \eqref{inverse:eq:a2} into
$A_2=a_2-a_1^2/2$.  The quartic term $d_0^4/8$ cancels, as do the
terms $d_0/12$ and $1/288$ after the square is expanded.  The remaining
terms reduce to \eqref{inverse:eq:A2}.  Smoothness, periodicity, and boundedness
follow because $A_2$ is a polynomial in $d_0,d_1,d_2,d_3$.
\end{proof}

\begin{corollary}[Two explicit correction orders; \Formalized]
\label{inverse:cor:two-corrections}
Let $\lambda=\lambda(x)$ be the exact lower Lambert phase.  Then, as
$x\downarrow0$,
\begin{equation}
 \log F(x)
 =\Main(x)+\frac{A_1(\lambda)}{\lambda}
          +\frac{A_2(\lambda)}{\lambda^2}
          +\BigO(\lambda^{-3}), \label{inverse:eq:two-correction-expansion}
\end{equation}
where $\Main(x)$ is the exact sharp Lambert main term.  Since
$\lambda\asymp-\log x$, the remainder is also
$\BigO((-\log x)^{-3})$.
\end{corollary}

\begin{warning}
The coefficients in \eqref{inverse:eq:two-correction-expansion} are evaluated at
the exact phase $\lambda(x)$.  Substituting a truncated asymptotic formula
for $\lambda$ into $A_1$ or $A_2$ is not licensed merely by
$\lambda-\lambda_{\rm approx}=\BigO(1)$: periodic functions are sensitive
to bounded phase shifts.
\end{warning}


\section{Why the all-orders remainder is actually rigorous}
\label{inverse:sec:all-orders-proof}

The formal coefficient algebra by itself does not prove an asymptotic
expansion.  One must control the exact contour integral uniformly on a
window that grows with the requested order, prove that exponentiation does
not magnify the Taylor error, replace the finite Gaussian window by the
whole line, and show that the discarded contour is smaller than every
retained inverse power.  This section spells out that analytic chain.

\subsection{The exact normalized saddle mass}

For $x>0$ small, let $\lambda=\lambda(x)$ be the exact phase
\eqref{inverse:eq:lambert-phase}.  The exact saddle reduction can be organized as
\begin{equation}
 \log F(x)
 =\Main(x)+\log\mathcal G(x)+\mathcal T(x),
 \label{inverse:eq:exact-saddle-reduction}
\end{equation}
where:
\begin{itemize}
\item $\Main(x)$ is the exact sharp Lambert main term, containing the
quadratic phase, the Gaussian normalization, the periodic finite part, and
the fixed constant;
\item $\mathcal G(x)>0$ is the normalized saddle-kernel mass and
$\mathcal G(x)\to1$;
\item $\mathcal T(x)$ is the forward negative-Laplace tail, which is flat
against every inverse power of $\lambda$.
\end{itemize}
Thus all algebraic corrections beyond $\Main$ come from $\log\mathcal G$.

It is technically convenient to begin on the real dyadic logarithmic
scale
\[
 x=2^{-t},\qquad t\to+\infty,
\]
and to write
\[
 b=\lambda(2^{-t}),\qquad \varepsilon=b^{-1/2}.
\]
This is not a restriction to integer $t$.  Every sufficiently small
positive $x$ has the unique representation $x=2^{-t}$ with
$t=-\log_2x$.

After the saddle contour is centered and rescaled by
$v=\sqrt b\,\theta$, its central kernel has the normalized form
\begin{equation}
 \frac1{\sqrt{2\pi}}\EulerE^{-v^2/2}
 \exp R_b(v), \label{inverse:eq:normalized-central-kernel}
\end{equation}
where $R_b$ is the exact centered exponent.  The formal expansion of
$R_b$ is
\begin{equation}
 R_b(v)\sim\sum_{m\geq1}E_m(b,v)b^{-m/2}.
 \label{inverse:eq:R-formal}
\end{equation}
The first variable of $E_m$ is evaluated at $b$ because all coefficient
functions are periodic in the exact phase.

\subsection{The order-dependent central window}

For a requested order $N\geq0$, define
\begin{equation}
 A_N(b)=\sqrt{32(N+1)\log b}. \label{inverse:eq:central-radius}
\end{equation}

\begin{lemma}[Scale separation; \Formalized]
\label{inverse:lem:central-scale}
For every fixed $N$,
\[
 A_N(b)\longrightarrow+\infty,\qquad
 \varepsilon A_N(b)\longrightarrow0
 \qquad(b\to+\infty).
\]
Moreover, for every fixed $r>0$ and $\delta>0$,
\[
 A_N(b)^r=o(b^\delta).
\]
\end{lemma}

\begin{proof}
The first limit is immediate from $\log b\to\infty$.  The second follows
from
\[
 \varepsilon A_N(b)
 =\sqrt{\frac{32(N+1)\log b}{b}}\longrightarrow0.
\]
Finally, every fixed power of $\log b$ grows more slowly than every
positive power of $b$.
\end{proof}

The window has two simultaneous purposes.  It grows, so replacing a
Gaussian integral over it by a whole-line Gaussian integral creates a
tiny tail.  Yet its angular width
$A_N(b)/\sqrt b$ shrinks to zero, so all local Taylor estimates remain
uniform.

\subsection{All-order control of the exact centered exponent}

The exact negative-Laplace factor contains a finite periodic part plus
endpoint and forward-product corrections.  On the vertical line through
the saddle, the proof decomposes it into:
\begin{enumerate}
\item a Taylor polynomial in $\theta=v/\sqrt b$ whose coefficients are
the periodic jets $p_n$;
\item the Taylor polynomial of $\log(1+\ImaginaryUnit\theta)$;
\item two boundary-jet remainders from the finite part;
\item forward-product jets whose size is exponentially small in $b$;
\item explicit integral remainders for the vertical derivatives.
\end{enumerate}
The factorial pieces in the first two items combine into the bounded
jets $d_n$, producing \eqref{inverse:eq:E-closed}.

\begin{proposition}[Integrated exponent truncation; \Formalized]
\label{inverse:prop:exponent-truncation}
Fix $N$.  The Taylor depth can be chosen, depending only on $N$, so that
on $|v|\leq A_N(b)$ the exact exponent may be replaced by
\[
 R_{b,N}(v)
 =\sum_{m=1}^{2N+1}E_m(b,v)b^{-m/2}
\]
with an error which, after multiplication by the Gaussian weight and
integration over the central window, is
\[
 \BigO(b^{-N-1}).
\]
The implied constant may depend on $N$ but not on $b$.
\end{proposition}

\begin{proof}
All vertical derivatives of the periodic finite part are uniformly
bounded because the periodic jets are smooth and bounded.  Taylor's
theorem therefore bounds a vertical remainder of order $M$ by
\[
 C_M|\theta|^{M+1}
 =C_M b^{-(M+1)/2}|v|^{M+1}.
\]
The logarithmic coordinate has an analogous remainder on the shrinking
angular window, because $\varepsilon A_N(b)\to0$ keeps
$1+\ImaginaryUnit\theta$ uniformly away from zero.  Boundary jets satisfy the same
kind of estimate, while the forward-product terms are exponentially
small and hence smaller than any prescribed inverse power.

Choose $M$ with enough surplus beyond $2N+1$.  On the central window,
powers of $|v|$ are bounded by powers of $A_N(b)$, hence by powers of
$\log b$.  By \cref{inverse:lem:central-scale}, the spare powers of $b^{-1/2}$
absorb every such logarithmic loss.  The pointwise exponent error is
therefore integrably $\BigO(b^{-N-1})$ against
$\EulerE^{-v^2/2}\Differential v$.  Removing terms above order $2N+1$ from the finite
Taylor polynomial costs the same order because each carries at least the
corresponding spare half-power.  This leaves $R_{b,N}$.
\end{proof}

\subsection{Finite exponentiation without an infinite-series argument}

A common informal step is to exponentiate \eqref{inverse:eq:R-formal} term by
term.  The formal development instead uses a finite Taylor polynomial for
the exponential and an explicit remainder.

Let
\[
 \mathcal P_N(b,v)
 =\sum_{\ell=0}^{2N+1}B_\ell(b,v)b^{-\ell/2}.
 \label{inverse:eq:reference-polynomial}
\]
It is the reference polynomial obtained from the formal recursion in
\cref{inverse:prop:exp-recursion}.

\begin{proposition}[Finite exponential substitution; \Formalized]
\label{inverse:prop:finite-exp-substitution}
For every fixed $N$,
\begin{equation}
 \frac1{\sqrt{2\pi}}
 \int_{|v|\leq A_N(b)}
 \EulerE^{-v^2/2}
 \left|\exp R_b(v)-\mathcal P_N(b,v)\right|\Differential v
 =\BigO(b^{-N-1}). \label{inverse:eq:central-reference-error}
\end{equation}
\end{proposition}

\begin{proof}
First replace $R_b$ by $R_{b,N}$ using
\cref{inverse:prop:exponent-truncation}.  On the central window the truncated
exponent tends uniformly to zero after the exact quadratic term has been
removed, and it is bounded by a fixed polynomial in $v$ times
$b^{-1/2}$.  Apply Taylor's formula
\[
 \EulerE^z=\sum_{r=0}^{L-1}\frac{z^r}{r!}
      +\frac{z^L}{(L-1)!}\int_0^1(1-s)^{L-1}\EulerE^{sz}\,\Differential s
\]
with $L=2(N+1)$.  Uniform control of the real part of $R_{b,N}$ bounds the
integral remainder by a Gaussian-integrable polynomial times
$b^{-N-1}$.

Now expand the finite Taylor polynomial in the finitely many powers
$b^{-m/2}$.  The coefficient of $b^{-\ell/2}$ for
$\ell<L$ is exactly $B_\ell$, by the same recursion that defines the
formal exponential coefficients.  Algebraically, the difference between
the finite exponential Taylor polynomial and
$\mathcal P_N$ factors as
\[
 b^{-L/2}\times
 \{\text{a finite quotient polynomial in }b^{-1/2},v,
   E_1,\dots,E_{2N+1}\}.
\]
Because $L/2=N+1$ and the coefficient functions are bounded, Gaussian
integration gives \eqref{inverse:eq:central-reference-error}.
\end{proof}

\subsection{Gaussian contraction on the whole line}

By \cref{inverse:cor:B-parity}, odd $\ell$ contribute odd polynomials.  Therefore
\begin{equation}
 \Gauss[\mathcal P_N(b,\cdot)]
 =\sum_{j=0}^{N}a_j(b)b^{-j}. \label{inverse:eq:reference-contraction}
\end{equation}
The exact central integral, however, only runs over
$|v|\leq A_N(b)$.  We must show that the polynomial reference can be
extended to the whole real line at the required order.

\begin{lemma}[Gaussian polynomial tail; \Formalized]
\label{inverse:lem:gaussian-polynomial-tail}
For every fixed polynomial $P$ and every fixed $N$,
\[
 \int_{|v|>A_N(b)}\EulerE^{-v^2/2}|P(v)|\,\Differential v
 =\BigO(b^{-N-1}).
\]
The same assertion holds for a finite family of polynomials whose
coefficients are uniformly bounded in $b$.
\end{lemma}

\begin{proof}
For every integer $q\geq0$ there is a constant $C_q$ such that
\[
 \int_A^\infty v^q\EulerE^{-v^2/2}\,\Differential v
 \leq C_q(1+A^q)\EulerE^{-A^2/4}
 \qquad(A\geq1).
\]
At $A=A_N(b)$,
\[
 \EulerE^{-A_N(b)^2/4}
 =\EulerE^{-8(N+1)\log b}
 =b^{-8(N+1)}.
\]
The polynomial factor in $A_N(b)$ is a power of $\log b$ and is absorbed
by the enormous surplus of inverse powers.  Sum over the finitely many
monomials of $P$.
\end{proof}

Combining \cref{inverse:prop:finite-exp-substitution,inverse:lem:gaussian-polynomial-tail}
with \eqref{inverse:eq:reference-contraction}, the central contour contributes
\begin{equation}
 \sum_{j=0}^{N}a_j(b)b^{-j}+\BigO(b^{-N-1}).
 \label{inverse:eq:central-mass-expansion}
\end{equation}

\subsection{The complementary contour}

The exact contour outside the central window is split into an intermediate
region and a far region.  The two estimates use different mechanisms.

\begin{proposition}[All-order minor-arc bounds; \Formalized]
\label{inverse:prop:minor-arcs}
For every fixed $N$, the normalized contour outside
$|v|\leq A_N(b)$ contributes $\BigO(b^{-N-1})$.
More precisely:
\begin{enumerate}
\item the intermediate region has a bound with fourfold power slack,
of order at most a constant times $b^{-4(N+1)}$ once the relevant
dyadic index is at least $b/4$;
\item the far region is bounded by a constant times
\[
 b\exp\!\left(-\frac{\log2}{4}b\right),
\]
which is flat against every inverse power of $b$.
\end{enumerate}
\end{proposition}

\begin{proof}
On the intermediate region the modulus estimate for the saddle kernel
contains a decay factor whose exponent is quadratic in the enlarged
radius.  Substitution of \eqref{inverse:eq:central-radius} yields
\[
 \exp\{-4(N+1)\log b\}=b^{-4(N+1)}
\]
after the standard comparison between the dyadic index and $b$.  The
remaining geometric and polynomial factors consume far less than the
three spare blocks of inverse powers.

In the far region the dyadic product supplies geometric decay in the
index.  Summing the geometric tail gives a bound of the displayed
exponential form.  For every $K$,
\[
 b^K\,b\EulerE^{-(\log2)b/4}\longrightarrow0,
\]
so this contribution is $\BigO(b^{-N-1})$ for the requested fixed $N$.
\end{proof}

\subsection{The all-orders mass expansion}

\begin{theorem}[Normalized mass expansion; \Formalized]
\label{inverse:thm:mass-all-orders}
For every fixed $N\geq0$, as $t\to+\infty$ with
$b=\lambda(2^{-t})$,
\begin{equation}
 \mathcal G(2^{-t})
 =\sum_{j=0}^{N}\frac{a_j(b)}{b^j}
  +\BigO(b^{-N-1}). \label{inverse:eq:mass-all-orders}
\end{equation}
The coefficient functions are evaluated at the exact phase $b$.
\end{theorem}

\begin{proof}
The central part is \eqref{inverse:eq:central-mass-expansion}.  Add the
intermediate and far parts controlled by \cref{inverse:prop:minor-arcs}.  All
three errors are $\BigO(b^{-N-1})$.
\end{proof}

\subsection{Transfer through the logarithm}

\begin{theorem}[Logarithmic expansion on the dyadic real scale;
\Formalized]
\label{inverse:thm:log-dyadic-all-orders}
For every fixed $N\geq0$,
\begin{equation}
 \log\mathcal G(2^{-t})
 =\sum_{j=1}^{N}\frac{A_j(b)}{b^j}
  +\BigO(b^{-N-1}),
 \qquad b=\lambda(2^{-t}). \label{inverse:eq:log-mass-all-orders}
\end{equation}
\end{theorem}

\begin{proof}
The boundedness of the $a_j$ and \eqref{inverse:eq:mass-all-orders} with $N=0$
give $\mathcal G(2^{-t})=1+\BigO(b^{-1})$, hence
$\mathcal G>0$ and $|\mathcal G-1|<1/2$ eventually.  Taylor's theorem for
$\log(1+z)$ is therefore uniform.  Substitute
\eqref{inverse:eq:mass-all-orders} into the finite Taylor polynomial for the
logarithm.  The coefficient of $b^{-j}$ is exactly $A_j$, either by
\eqref{inverse:eq:A-compositions} or by the recursion
\eqref{inverse:eq:A-recursion}.  The logarithmic Taylor remainder and the
original mass remainder are both $\BigO(b^{-N-1})$.
\end{proof}

\subsection{Adding the flat tail and returning to \texorpdfstring{\(x\downarrow0\)}{x to 0}}

The tail $\mathcal T$ in \eqref{inverse:eq:exact-saddle-reduction} satisfies, for
every $K$,
\[
 \mathcal T(x)=\BigO(\lambda(x)^{-K}).
\]
It therefore does not change any coefficient in a Poincar\'e expansion.

\begin{theorem}[Full exact-phase expansion; \Formalized]
\label{inverse:thm:full-all-orders}
For every $N\geq0$,
\begin{equation}
 \log F(x)
 =\Main(x)+
   \sum_{j=1}^{N}\frac{A_j(\lambda(x))}{\lambda(x)^j}
   +\BigO(\lambda(x)^{-N-1})
 \qquad(x\downarrow0). \label{inverse:eq:full-all-orders}
\end{equation}
Equivalently, using a truncation with indices $j<N$,
\[
 \log F(x)-\Main(x)-
 \sum_{j=1}^{N-1}\frac{A_j(\lambda(x))}{\lambda(x)^j}
 =\BigO(\lambda(x)^{-N}).
\]
\end{theorem}

\begin{proof}
First combine \eqref{inverse:eq:exact-saddle-reduction},
\eqref{inverse:eq:log-mass-all-orders}, and the flatness of $\mathcal T$ on the
scale $x=2^{-t}$.  Next use the exact change of variables
$t=-\log_2x$.  As $x\downarrow0$, $t\to+\infty$ and
$\lambda(2^{-t})=\lambda(x)$, so the dyadic-real estimate transports
without approximation.  This yields \eqref{inverse:eq:full-all-orders}.
\end{proof}

\begin{warning}[Poincar\'e does not mean convergent]
The theorem is fixed-order.  For each $N$ there is an error constant
depending on $N$, and the central radius itself depends on $N$.  No
uniform control as $N\to\infty$ is asserted, and no convergence of the
infinite coefficient series is implied.
\end{warning}

\subsection{Exact transfer to the right endpoint}

\begin{corollary}[Right-endpoint deficiency; \Derived{} from formalized identities]
\label{inverse:cor:right-endpoint-all-orders}
For every \(N\geq0\), as \(x\downarrow0\),
\begin{equation}
 \log\bigl(1-F(1-x)\bigr)
 =\Main(x)+
   \sum_{j=1}^{N}\frac{A_j(\lambda(x))}{\lambda(x)^j}
   +\BigO(\lambda(x)^{-N-1}).
 \label{inverse:eq:right-endpoint-all-orders}
\end{equation}
\end{corollary}

\begin{proof}
The reflection identity gives \(1-F(1-x)=F(x)\) exactly.  Substitute
\eqref{inverse:eq:full-all-orders}; no asymptotic error is introduced by the transfer.
\end{proof}

\subsection{The separate elementary expansion of the Lambert phase}

The exact-phase theorem above does not require replacing \(\lambda(x)\).  For comparison,
put \(S=-\log x\) and \(z=\log S-\log L\).  There are unique polynomials
\(\ell_j\), with \(\deg\ell_j=j\) and \(\ell_j(0)=0\), such that for each fixed
\(N\geq1\)
\begin{equation}
 L\lambda(x)=S+z+\sum_{j=1}^{N-1}\frac{\ell_j(z)}{S^j}
 +\BigO\!\left(\frac{(1+|z|)^N}{S^N}\right).
 \label{inverse:eq:lambert-phase-all-orders}
\end{equation}
The letter \(\ell\) avoids collision with the periodic saddle jets \(p_n\) of
\cref{inverse:sec:closed-jets}.  The first terms are
\[
 \ell_1(z)=z,\qquad
 \ell_2(z)=z-\frac{z^2}{2},\qquad
 \ell_3(z)=z-\frac{3z^2}{2}+\frac{z^3}{3}.
\]
If \(\mathfrak s(j,k)\) denotes an \emph{unsigned} Stirling number of the first kind, then
\begin{equation}
 \ell_j(z)=\sum_{m=1}^{j}
 \frac{(-1)^{m+1}}{m!}\,
 \mathfrak s(j,j-m+1)z^m.
 \label{inverse:eq:lambert-phase-stirling}
\end{equation}
Equivalently, one may use signed first-kind Stirling numbers and absorb the sign into that
convention; mixing the signed convention with the displayed exterior sign would be wrong.
Substituting a finite truncation of \eqref{inverse:eq:lambert-phase-all-orders} into the
periodic coefficients is a separate composition problem: the quadratic main exponent
amplifies phase error, and even a bounded error moves the argument modulo one.  No such
all-orders substitution is claimed here.

\section{Higher explicit corrections and their Fourier signatures}
\label{inverse:sec:higher-explicit}

\begin{boxedremark}
\textbf{Merged specialist source.}
This section preserves the independent symbolic and numerical work formerly stored at
\nolinkurl{Small_Argument_Asymptotics/Small_Argument_Asymptotics.tex}, SHA-256
\nolinkurl{85f51a20fc7b6bdf3b1d049ec4506f508aee3c4cd70554e6a68cbcc30977cb0b}.
Its repeated derivations of the lower-Lambert phase, periodic main term, closed jets,
general coefficient recursions, and first two corrections have been removed in favor of
\cref{part:dyadicasym,inverse:sec:closed-jets,inverse:sec:coefficient-formulas,inverse:sec:second-correction}.
What remains is the material genuinely added by that specialist calculation: exact mean
identities, expanded third and fourth corrections, derivative/Fourier forms, amplitude
estimates, high-precision checks, and numerical traps.  The displayed \(A_3\) and \(A_4\) formulas and
the numerical evidence are not promoted to Lean theorems by this merge.
\end{boxedremark}

\subsection{The third and fourth coefficients in bounded-jet coordinates}
The first two jet formulas are \cref{inverse:eq:A1-jets,inverse:eq:A2}; continuing the same
symbolic exponential--Gaussian--logarithm pipeline gives the following two research-frontier
coefficients.

\begin{equation}\label{smallarg:eq:A3jets}
\begin{aligned}
  A_3=-\frac{1}{720}\bigl(&180d_0^4+720d_0^3d_1+120d_0^3d_2+240d_0^3
    +360d_0^2d_1^2+2520d_0^2d_1\\
   &+1440d_0^2d_2+180d_0^2d_3+2160d_0d_1^2+720d_0d_1d_2+1440d_0d_1\\
   &+2520d_0d_2+960d_0d_3+90d_0d_4+120d_1^3+1980d_1^2+1800d_1d_2\\
   &+180d_1d_3+150d_2^2+720d_2+780d_3+210d_4+15d_5-2\bigr).
\end{aligned}
\end{equation}

One order further, in the same coordinates,
\begin{equation}\label{smallarg:eq:A4jets}
\begin{aligned}
  A_4=\frac{1}{5760}\bigl(&1152d_0^5+8640d_0^4d_1+2880d_0^4d_2+240d_0^4d_3+1440d_0^4\\
   &+11520d_0^3d_1^2+2880d_0^3d_1d_2+28800d_0^3d_1+25920d_0^3d_2\\
   &+6720d_0^3d_3+480d_0^3d_4+2880d_0^2d_1^3+63360d_0^2d_1^2\\
   &+46080d_0^2d_1d_2+4320d_0^2d_1d_3+17280d_0^2d_1+2880d_0^2d_2^2\\
   &+46080d_0^2d_2+29280d_0^2d_3+6000d_0^2d_4+360d_0^2d_5\\
   &+23040d_0d_1^3+8640d_0d_1^2d_2+74880d_0d_1^2+123840d_0d_1d_2\\
   &+30720d_0d_1d_3+2160d_0d_1d_4+21600d_0d_2^2+3840d_0d_2d_3\\
   &+17280d_0d_2+30720d_0d_3+14640d_0d_4+2400d_0d_5+120d_0d_6\\
   &+720d_1^4+30720d_1^3+28800d_1^2d_2+2160d_1^2d_3+14400d_1^2\\
   &+3600d_1d_2^2+66240d_1d_2+36960d_1d_3+6720d_1d_4+360d_1d_5\\
   &+25200d_2^2+12000d_2d_3+840d_2d_4+480d_3^2+5760d_3\\
   &+7392d_4+2720d_5+360d_6+15d_7\bigr).
\end{aligned}
\end{equation}

Its top term is $15d_7/5760=d_7/384$, and $(-1)^4/(2^4\cdot4!)=1/384$: a
fourth independent check of \cref{inverse:thm:top-jet}.


The occurrence of \(d_7/384\) in \(A_4\) is consistent with the formalized mass top-jet
coefficient in \cref{inverse:thm:top-jet}.  Its transfer through the logarithm, and hence the
general highest-derivative conclusion, remains a derived calculation as stated there.

\subsection{The same coefficients in the derivatives of \texorpdfstring{$\Psi$}{Psi}}

Write $P_k=\Psi^{(k)}(u)$.  Then
\[
  A_1(u)=\frac1{24}+\frac1{2L}-\frac{P_2+P_1^2}{2L^2},
\]
\begin{equation}\label{smallarg:eq:A2psi}
\begin{aligned}
  A_2(u)={}&\frac{P_1}{24L}
    +\frac{2P_1+1}{4L^2}
    -\frac{P_1^3+3P_1^2+3P_1P_2+3P_2+P_3}{6L^3}\\
   &+\frac{4P_1^2P_2+4P_1P_3+2P_2^2+P_4}{8L^4}.
\end{aligned}
\end{equation}
The corresponding form of $A_3$ is long.  Its \emph{linear part} is the
following; the omitted terms are quadratic and higher in the derivatives of
$\Psi$ and reach $8.9\times10^{-7}$ over a period, which is four decades
below the $6.7\times10^{-2}$ swing of $A_3$ but is not negligible if six
digits are wanted.
\begin{equation}\label{smallarg:eq:A3psi}
\begin{aligned}
  A_3(u)={}&-\frac7{2880}-\frac1{24L}-\frac1{4L^2}+\frac1{6L^3}
    +\frac{P_1+P_2}{24L^2}
    +\frac{48P_1+24P_2-P_3}{48L^3}\\
   &-\frac{6P_2+9P_3+P_4}{12L^4}
    +\frac{3P_4+P_5}{12L^5}
    -\frac{P_6}{48L^6}
    +(\text{quadratic and higher in the }P_k).
\end{aligned}
\end{equation}

\subsection{The \texorpdfstring{$\Psi$}{Psi}-free parts}

\begin{center}
\begin{tabular}{@{}clr@{}}
\toprule
$j$ & $\Psi$-free part of $A_j$ & value\\
\midrule
$1$ & $\dfrac1{24}+\dfrac1{2L}$ & $0.7630141871$\\[0.6em]
$2$ & $\dfrac1{4L^2}$ & $0.5203422453$\\[0.6em]
$3$ & $-\dfrac7{2880}-\dfrac1{24L}-\dfrac1{4L^2}+\dfrac1{6L^3}$ & $-0.0824216430$\\
$4$ & $\dfrac{-L^3-72L^2-816L+144}{1152L^4}$ & $-1.7167954639$\\
\bottomrule
\end{tabular}
\end{center}

The order-two value is remarkably clean: the $\Psi$-free part of $A_2$ is
exactly $1/(4\log^2 2)$, with no rational term and no $1/L$ term at all.

\subsection{Exact period means}

The \(\Psi\)-free part is not literally the period mean, because nonlinear products of
derivatives can have nonzero mean.  Write
\(\langle H\rangle=\int_0^1H(u)\,\Differential u\).  Fourier multiplication gives the exact identity
\begin{equation}
 \bigl\langle\Psi^{(a)}\Psi^{(b)}\bigr\rangle
 =\sum_{k\ne0}(2\pi\ImaginaryUnit k)^a(-2\pi\ImaginaryUnit k)^b
   |\FourierSeriesCoefficient{\Psi}{k}|^2.
 \label{smallarg:eq:derivative-product-mean}
\end{equation}
Integration by parts over one period gives, in particular,
\[
 \langle\Psi'\Psi''\rangle=0,\qquad
 \langle(\Psi')^2\Psi''\rangle=0,\qquad
 \langle\Psi'\Psi'''\rangle=-\langle(\Psi'')^2\rangle.
\]
Substitution in the exact formulas for \(A_1\) and \(A_2\) leaves
\begin{align}
 \langle A_1\rangle
 &=\frac1{24}+\frac1{2L}
   -\frac{\langle(\Psi')^2\rangle}{2L^2},
   \label{smallarg:eq:mean-A1}\\
 \langle A_2\rangle
 &=\frac1{4L^2}
   -\frac{\langle(\Psi')^2\rangle}{2L^3}
   -\frac{\langle(\Psi'')^2\rangle}{4L^4}
   -\frac{\langle(\Psi')^3\rangle}{6L^3}.
   \label{smallarg:eq:mean-A2}
\end{align}
\begin{samepage}
These are identities, not truncated estimates.  Numerically,
\[
 \langle(\Psi')^2\rangle=8.917038\times10^{-11},\quad
 \langle(\Psi'')^2\rangle=3.520305\times10^{-9},\quad
 \langle(\Psi')^3\rangle=1.115006\times10^{-21},
\]
and hence
\[
 \langle A_1\rangle=0.7630141870184,\qquad
 \langle A_2\rangle=0.5203422413049.
\]
\end{samepage}
Their gaps from the \(\Psi\)-free constants are respectively about
\(-9.3\times10^{-11}\) and \(-3.9\times10^{-9}\).  The cubic term in
\eqref{smallarg:eq:mean-A2} is tiny but retained because deleting it would turn an exact
identity into an approximation.

The Gamma--zeta formula \eqref{dyadicasym:eq:Psi-Fourier-report} gives, numerically,
\[
\begin{aligned}
  \FourierSeriesCoefficient{\Psi}{1}
    &=(7.55864754098-7.47010355898\,\ImaginaryUnit)\,10^{-7},
    &|\FourierSeriesCoefficient{\Psi}{1}|&=1.06271162519\,10^{-6},\\
  \FourierSeriesCoefficient{\Psi}{2}
    &=(3.24812347723+5.85175898712\,\ImaginaryUnit)\,10^{-13},
    &|\FourierSeriesCoefficient{\Psi}{2}|&=6.69278636792\,10^{-13},\\
  \FourierSeriesCoefficient{\Psi}{3}
    &=(-2.81236523199-2.58674397916\,\ImaginaryUnit)\,10^{-19},
    &|\FourierSeriesCoefficient{\Psi}{3}|&=3.82107872359\,10^{-19}.
\end{aligned}
\]
Thus \(\Psi\) itself oscillates only by a few parts in \(10^6\), even though repeated
differentiation greatly amplifies that oscillation in the higher saddle coefficients.

\section{Fourier form and the amplitude law}\label{smallarg:sec:fourier}

Because
$\Psi^{(m)}(u)=\sum_{k\ne0}(2\pi\ImaginaryUnit k)^m
\FourierSeriesCoefficient{\Psi}{k}\EulerE^{2\pi\ImaginaryUnit ku}$,
\cref{inverse:sec:coefficient-formulas} together with \cref{dyadicasym:eq:Psi-Fourier-report} turns every coefficient into
an explicit Gamma--zeta Fourier object.  Discarding the terms of degree two
and higher in the derivatives of $\Psi$ leaves the linearized Fourier form
\[
  A_j(u)=A_j^{(0)}
   +2\operatorname{Re}\Bigl(\sum_{k\ge1}c_j(k)
   \FourierSeriesCoefficient{\Psi}{k}\EulerE^{2\pi\ImaginaryUnit ku}\Bigr)
   +R_j(u),
\]
where \(A_j^{(0)}\) is the \(\Psi\)-free constant and \(R_j\) is the raw sum of
terms of degree at least two in derivatives of \(\Psi\).  Thus
\(\langle A_j\rangle=A_j^{(0)}+\langle R_j\rangle\), as quantified above.  The
discarded remainder is measured at
$\max_u|R_j(u)|=1.6\times10^{-8}$ for $j=2$ and $8.9\times10^{-7}$ for
$j=3$, growing with $j$ like everything else here.  It is negligible against
the oscillation it sits inside, and it is not negligible in absolute terms.
Here $c_j(k)$ is the polynomial in $w=2\pi\ImaginaryUnit k$ read off from
\cref{smallarg:eq:A2psi,smallarg:eq:A3psi}.  For the first two,
\[
  c_1(k)=-\frac{w^2}{2L^2},
  \qquad
  c_2(k)=\frac{w}{24L}+\frac{w}{2L^2}-\frac{3w^2+w^3}{6L^3}+\frac{w^4}{8L^4}.
\]
The resulting data for $A_2$ are all confirmed numerically in
\cref{smallarg:sec:numerics}.
For the first Fourier mode, the linearized amplitude is
$2|c_2(1)\FourierSeriesCoefficient{\Psi}{1}|$.  The left column is what the linearized
form gives, the right what the exact
$A_2$ gives; they are quoted separately because the difference is real, if
small, and quoting ten digits of a linearization would be an overclaim.
\[
\begin{aligned}
  \text{\(\Psi\)-free constant / mean}&&\tfrac1{4L^2}=0.5203422452514&\qquad
    \langle A_2\rangle=0.5203422413049\\
  k=1\ \text{amplitude}&&1.9398782\times10^{-3}&\qquad
    1.9398782\times10^{-3}\\
  \text{maximum at}&&\{\lambda\}=0.1011297&\qquad\{\lambda\}=0.1011281\\
  \text{minimum at}&&\{\lambda\}=0.6011297&\qquad\{\lambda\}=0.6011313
\end{aligned}
\]
The exact peak-to-peak is $3.8797564\times10^{-3}$, twice the linearized
$k=1$ amplitude to the displayed precision, and the extrema agree to six.  The
maximum of the exact $A_2$ is $0.522282117$ and its minimum is
$0.518402361$.

\subsection{The amplitude law}

The top term of \cref{inverse:thm:top-jet} dominates the oscillation, and it multiplies
the $k$-th Fourier mode by $(2\pi k)^{2j}$.  Taking $k=1$,
\begin{equation}\label{smallarg:eq:amplitude}
  \text{amplitude of }A_j\ \approx\
  2|\FourierSeriesCoefficient{\Psi}{1}|\,\frac{(2\pi^2/L^2)^j}{j!},
  \qquad \frac{2\pi^2}{L^2}=41.0845769104 .
\end{equation}

\begin{center}
\begin{tabular}{@{}cr@{}}
\toprule
$j$ & predicted $k=1$ amplitude\\
\midrule
$1$ & $8.73\times10^{-5}$\\
$2$ & $1.79\times10^{-3}$\\
$3$ & $2.46\times10^{-2}$\\
$4$ & $2.52\times10^{-1}$\\
$5$ & $2.07$\\
$6$ & $1.42\times10^{1}$\\
$7$ & $8.33\times10^{1}$\\
\bottomrule
\end{tabular}
\end{center}

These are estimates from the top term alone.  Comparing them with the truth
requires care about a convention: an \emph{amplitude} is half a
peak-to-peak swing, and mixing the two silently doubles every second
entry.  Everything below is an amplitude.

\begin{center}
\begin{tabular}{@{}crrr@{}}
\toprule
$j$ & \cref{smallarg:eq:amplitude} & actual & ratio\\
\midrule
$1$ & $8.73\times10^{-5}$ & $8.73\times10^{-5}$ & $1.00$\\
$2$ & $1.79\times10^{-3}$ & $1.94\times10^{-3}$ & $1.08$\\
$3$ & $2.46\times10^{-2}$ & $\approx3\times10^{-2}$ & $\approx1.2$\\
$4$ & $2.52\times10^{-1}$ & $\approx0.41$ & $\approx1.6$\\
\bottomrule
\end{tabular}
\end{center}

At $j=1$ the top derivative term dominates to the displayed precision, but it is not literally the only periodic term: $A_1$ also contains $-(\Psi')^2/(2L^2)$.  The equality in the table is therefore a very accurate linearization, not an exact amplitude identity.  At
$j=2$ the $w^3$ term of $c_2$ adds in phase and lifts the amplitude above
the estimate; the entries at $j=3$ and $j=4$ are halves of the swings
measured in \cref{smallarg:sec:numerics}.  The ratio drifts upward because the
subleading terms of $c_j$ contribute more as $j$ grows.  So
\cref{smallarg:eq:amplitude} gives the growth rate and not the constant, and should
never be used to size a particular coefficient.  So the periodic parts of
the coefficients grow rapidly with $j$, even though \(\Psi\) itself has
amplitude only about \(2\times10^{-6}\) (roughly \(4\times10^{-6}\) peak to
peak): differentiation progressively cancels the smallness of
$\FourierSeriesCoefficient{\Psi}{1}$.  Any numerical extraction of a coefficient beyond
the first must
therefore control the phase; see \cref{smallarg:sec:traps}.

\subsection{Why the expansion is only asymptotic}

$|\FourierSeriesCoefficient{\Psi}{k}|$ decays like
$\exp(-\pi|\chi_k|/2)=\exp(-\pi^2k/L)$, that is,
about $\EulerE^{-14.24k}$, while mode $k$ is amplified by
$(2\pi k)^{2j}/(2^jj!\,L^{2j})$.  The $k$-th contribution therefore behaves like
$\exp(2j\log k-14.24k)$, maximized near $k^\ast=2j/14.24=0.14j$.  For $j$ beyond
about $8$ the higher modes dominate, and $\sup|A_j|$ grows faster than
geometrically.

What is worth noticing is \emph{which} divergence this is.  A saddle-point
expansion normally diverges factorially, because the Gaussian moments grow
like $(2j-1)!!$ against a fixed exponent.  Here that growth is exactly
cancelled: the $(2j)!$ in $E_{2j}$ absorbs it, which is the content of
\cref{inverse:thm:top-jet} and the reason the coefficient of the top jet is the mild
$1/(2^jj!)$.  The divergence that remains is of a different origin
altogether --- it is driven by the high-frequency content of $\Psi$, through
the $(2\pi k)^{2j}$ that differentiation applies to the $k$-th mode.  A
reader who expects the usual factorial mechanism will look for it in the
Gaussian integration, where it is not.  (This framing is due to the parallel
write-up of the main article, where it also appears.)

\begin{remark}
The expansion is a Poincar\'e asymptotic expansion: each truncation is valid
with an $\BigO(\lambda^{-N})$ remainder for fixed $N$, and there is no claim,
here or in the formalization, that the series converges for any fixed $x$.
The estimate in this subsection is a heuristic size estimate, not a theorem.
\end{remark}

\section{Numerical verification}\label{smallarg:sec:numerics}

$F(2^{-n})$ is computed \emph{exactly} as a rational number from the
inverse-power recursion used by the Lean evaluator
(\nolinkurl{Arithmetic.lean}, \lstinline[style=pseudo]{fabiusInversePowTwoTable}):
\[
  v_0=1,\qquad
  v_{n+1}=\frac{1}{2^{\,n+1}-1}
    \sum_{k=0}^{n}\frac{v_k}{2^{\binom{n+1}{2}-\binom{k}{2}}\,(n-k+2)!},
  \qquad F(2^{-n})=v_n.
\]
Sanity checks: $v_1=1/2=F(1/2)$ and $v_2=5/72=F(1/4)$.  All computations below
use $260$-digit working precision, with $\Psi$ and its derivatives summed over
six Gamma--zeta modes and $\zeta$ evaluated by Euler--Maclaurin
(\cref{smallarg:sec:traps}).

Define the successive residuals
\[
  r_1=\lambda\bigl(\log F-\mathcal M\bigr),\quad
  r_2=\lambda^2\Bigl(\log F-\mathcal M-\frac{A_1}{\lambda}\Bigr),\quad
  r_3=\lambda^3\Bigl(\log F-\mathcal M-\frac{A_1}{\lambda}-\frac{A_2}{\lambda^2}\Bigr).
\]

\begin{center}\small
\resizebox{\textwidth}{!}{%
\begin{tabular}{@{}rrrrrrrr@{}}
\toprule
$n$ & $\lambda$ & $r_1$ & $A_1(\lambda)$ & $r_2$ & $A_2(\lambda)$ & $r_3$ & $A_3(\lambda)$\\
\midrule
$10$ & $13.78503056$ & $0.7997155875$ & $0.7629678468$ & $0.5065687288$ & $0.5195595404$ & $-0.1790787344$ & $-0.0868938679$\\
$20$ & $24.62186833$ & $0.7836542295$ & $0.7629268731$ & $0.5103462408$ & $0.5184188039$ & $-0.1987615843$ & $-0.1111875606$\\
$40$ & $45.50804986$ & $0.7742689337$ & $0.7629490545$ & $0.5151456258$ & $0.5187247759$ & $-0.1628801414$ & $-0.1121601297$\\
$60$ & $66.04538587$ & $0.7709808331$ & $0.7630910552$ & $0.5210834261$ & $0.5221643489$ & $-0.0713899656$ & $-0.0512594569$\\
$80$ & $86.43351899$ & $0.7689739453$ & $0.7629823191$ & $0.5178773369$ & $0.5193822939$ & $-0.1300787245$ & $-0.1046130809$\\
$120$ & $126.9885547$ & $0.7671769004$ & $0.7630717251$ & $0.5213102874$ & $0.5218167515$ & $-0.0643151389$ & $-0.0540815552$\\
$160$ & $167.3870441$ & $0.7661091776$ & $0.7630070725$ & $0.5192522062$ & $0.5199082072$ & $-0.1098060709$ & $-0.0973365576$\\
\bottomrule
\end{tabular}}
\end{center}

The decisive test is that each residual, minus its predicted coefficient, is
$\BigO(1/\lambda)$.  Writing $e_j=r_j-A_j$:

\begin{center}\small
\begin{tabular}{@{}rrrrr@{}}
\toprule
$n$ & $\{\lambda\}$ & $\lambda e_1$ & $\lambda e_2$ & $\lambda e_3$\\
\midrule
$100$ & $0.737929$ & $0.51802529$ & $-0.11213991$ & $-1.7523898$\\
$112$ & $0.893526$ & $0.52019699$ & $-0.07798492$ & $-1.3773605$\\
$124$ & $0.033795$ & $0.52164121$ & $-0.06156984$ & $-1.3168268$\\
$136$ & $0.161500$ & $0.52166737$ & $-0.06826635$ & $-1.5329675$\\
$148$ & $0.278716$ & $0.52062406$ & $-0.08861040$ & $-1.8452719$\\
$160$ & $0.387044$ & $0.51925221$ & $-0.10980607$ & $-2.0872350$\\
\bottomrule
\end{tabular}
\end{center}

$\lambda e_1$ reproduces $A_2$ with mean $0.5203$; $\lambda e_2$ reproduces
$A_3$ with mean $-0.0824$; and $\lambda e_3$ is bounded, locating $A_4$ between
about $-1.3$ and $-2.1$.  Each column stays bounded and oscillates with the
phase rather than drifting, which is precisely the assertion that the
corresponding coefficient is correct and the next one is $\BigO(1)$.

Pointwise, $A_2$ is confirmed at the $10^{-5}$ level once the next order is
included.  For $140\le n\le160$:

\begin{center}\small
\begin{tabular}{@{}rrrrr@{}}
\toprule
$n$ & $\{\lambda\}$ & $A_2(\lambda)$ & $r_2$ & $r_2-A_3(\lambda)/\lambda$\\
\midrule
$140$ & $0.201650$ & $0.5219078895$ & $0.5214039042$ & $0.5218323910$\\
$145$ & $0.250301$ & $0.5214906179$ & $0.5209456288$ & $0.5214143255$\\
$150$ & $0.297351$ & $0.5209853273$ & $0.5203979284$ & $0.5209087825$\\
$155$ & $0.342900$ & $0.5204425007$ & $0.5198167900$ & $0.5203665207$\\
$160$ & $0.387044$ & $0.5199082072$ & $0.5192522062$ & $0.5198337121$\\
\bottomrule
\end{tabular}
\end{center}

The last two columns agree to about $7.5\times10^{-5}$, which is
$A_4/\lambda^2$ with $A_4\approx-2$.  Note that the \emph{variation} of $A_2$
over this window is $2.0\times10^{-3}$, twenty-six times the residual, so this
is a genuine pointwise verification of the periodic shape and not merely of a
mean.

\begin{remark}
An independent high-precision computation performed in a separate worktree,
using half-moments to $n=400$ at $400$-digit precision and a different route to
$\lambda$, located the minimum of the order-$\lambda^{-2}$ residual at
$\{\lambda\}=0.602$ and measured a peak-to-peak spread of about
$3\times10^{-3}$; \cref{smallarg:eq:A2psi} predicts a minimum at $0.6011313$ and a
peak-to-peak of $3.8798\times10^{-3}$ over the full period.  That is a phase
prediction, not a magnitude fit, and it was made before the comparison.
\end{remark}

\subsection{The fourth coefficient}

The same apparatus confirms \cref{smallarg:eq:A4jets}.  Writing
$r_4=\lambda^4(\log F-\mathcal M-A_1/\lambda-A_2/\lambda^2-A_3/\lambda^3)$,
at $300$-digit precision with eight Gamma--zeta modes:

\begin{center}\small
\begin{tabular}{@{}rrrr@{}}
\toprule
$n$ & $\{\lambda\}$ & $r_4$ & $A_4(\lambda)$\\
\midrule
$120$ & $0.988555$ & $-1.299548008$ & $-1.305123624$\\
$136$ & $0.161500$ & $-1.532967481$ & $-1.506670619$\\
$152$ & $0.315745$ & $-1.939809221$ & $-1.890185255$\\
$160$ & $0.387044$ & $-2.087234974$ & $-2.035031382$\\
$176$ & $0.519792$ & $-2.166743696$ & $-2.124617975$\\
$192$ & $0.641266$ & $-1.992935176$ & $-1.970225155$\\
$200$ & $0.698346$ & $-1.853347748$ & $-1.840001848$\\
\bottomrule
\end{tabular}
\end{center}

$\lambda(r_4-A_4)$ stays inside $[-8.8,0.8]$ over $120\le n\le200$ with no
drift, which places $A_5$ and confirms $A_4$.  The $\Psi$-free part
$-1.7167954639$ sits inside the measured range, and the swing of the
$A_4$ column, $0.82$ peak to peak, is $1.6$ times the leading-term estimate
$0.50$ of \cref{smallarg:eq:amplitude} --- the same upward drift of that ratio seen
at $j=2$ ($1.08$) and $j=3$ ($1.24$), because the subleading terms of
$c_j(1)$ contribute more as $j$ grows.  \Cref{smallarg:eq:amplitude} is a
leading-term estimate and should be read as one.

\subsection{A worked example}\label{smallarg:sec:worked}

It is worth seeing the expansion used once, on a value that can be written
down exactly.  At $x=2^{-10}$ the evaluator gives
\[
\begin{aligned}
  F(2^{-10})
    &=\frac{134\,926\,369}
      {1\,246\,394\,851\,358\,539\,387\,238\,350\,848\,000},\\
  \log F(2^{-10})&=-50.57756828234216081254\dots .
\end{aligned}
\]
and $\lambda=13.7850305606$, so $\{\lambda\}=0.785031$.  The successive
truncations miss $\log F$ by

\begin{center}\small
\begin{tabular}{@{}lrr@{}}
\toprule
truncation & error at $n=10$ & error at $n=160$\\
\midrule
$\mathcal M(x)$ & $5.80\times10^{-2}$ & $4.58\times10^{-3}$\\
${}+A_1/\lambda$ & $2.67\times10^{-3}$ & $1.85\times10^{-5}$\\
${}+A_2/\lambda^2$ & $-6.84\times10^{-5}$ & $-2.34\times10^{-8}$\\
${}+A_3/\lambda^3$ & $-3.52\times10^{-5}$ & $-2.66\times10^{-9}$\\
\bottomrule
\end{tabular}
\end{center}

where $n=160$ has $\lambda=167.387044$ and
$\log F(2^{-160})=-9489.629874649676879065\dots$

Two things are visible.  At $n=160$ each order gains roughly a factor
$\lambda$, as it should: $250$, then $790$.  The last gain is only $8.8$,
and that is not a defect --- the third-order error is already dominated by
$A_4/\lambda^4$, which \cref{smallarg:eq:A4jets} puts at
$-1.72/167.387^4=-2.19\times10^{-9}$ against a measured
$-2.66\times10^{-9}$.  So the table is measuring $A_4$, not a failure.

At $n=10$, by contrast, the fourth term barely helps: the error falls only
from $6.8\times10^{-5}$ to $3.5\times10^{-5}$.  With $\lambda=13.8$ and
$A_4\approx-1.7$, the term $A_4/\lambda^4$ is itself $4.7\times10^{-5}$,
comparable to the error it is meant to reduce.  This is the ordinary
behaviour of a Poincar\'e expansion near the end of its useful range, and it
is the practical reason \cref{smallarg:sec:fourier} declines to claim convergence:
the coefficients grow, so for a fixed $x$ there is a best truncation, and
going past it makes the approximation worse.  At $x=2^{-10}$ that optimum is
already at hand around the third term.

\subsection{An independent confirmation to \texorpdfstring{$n=400$}{n=400}}

The following was produced in a separate worktree, by a different agent,
from a different representation of $F$, and is reproduced with its
permission.  To avoid collision with the bounded saddle jets \(d_n\), write \(h_n\) for
the half moments in this paragraph.  They were computed to $n=400$ at $400$-digit working
precision from
\[
  h_n=\frac1{(n+1)(2^n-1)}\sum_{k<n}\binom{n+1}{k}h_k,
\]
whose terms are all positive, so there is no cancellation; the values agree
with the exact rationals to $400$ digits at $n=5,20,60$.  Then
$\log F(2^{-n})=\log h_n-\log n!-\binom n2\log2$, with $\lambda$ from the
lower Lambert branch (satisfying $\lambda2^{-\lambda}=x$ to $10^{-21}$) and
$\Psi$ from ten Gamma--zeta modes.  Nothing in this table shares an
implementation with the preceding numerical tables.

\begin{center}\small
\begin{tabular}{@{}rrrrrr@{}}
\toprule
$n$ & $\lambda$ & $\{\lambda\}$ & $r_2$ measured & $A_2(\lambda)$ & $r_3$\\
\midrule
$100$ & $106.73793$ & $0.7379$ & $0.51802529$ & $0.51907590$ & $-0.11214$\\
$125$ & $132.04488$ & $0.0449$ & $0.52169792$ & $0.52216224$ & $-0.06131$\\
$150$ & $157.29735$ & $0.2974$ & $0.52039793$ & $0.52098533$ & $-0.09240$\\
$175$ & $182.51185$ & $0.5119$ & $0.51801879$ & $0.51869969$ & $-0.12427$\\
$200$ & $207.69835$ & $0.6984$ & $0.51821775$ & $0.51875313$ & $-0.11120$\\
$225$ & $232.86334$ & $0.8633$ & $0.52015631$ & $0.52049092$ & $-0.07792$\\
$250$ & $258.01129$ & $0.0113$ & $0.52175813$ & $0.52198119$ & $-0.05755$\\
$275$ & $283.14540$ & $0.1454$ & $0.52199208$ & $0.52220755$ & $-0.06101$\\
$300$ & $308.26804$ & $0.2680$ & $0.52104845$ & $0.52130958$ & $-0.08050$\\
$325$ & $333.38103$ & $0.3810$ & $0.51967269$ & $0.51997997$ & $-0.10244$\\
$350$ & $358.48577$ & $0.4858$ & $0.51856496$ & $0.51889001$ & $-0.11653$\\
$375$ & $383.58340$ & $0.5834$ & $0.51810490$ & $0.51841439$ & $-0.11872$\\
$400$ & $408.67481$ & $0.6748$ & $0.51833670$ & $0.51860654$ & $-0.11028$\\
\bottomrule
\end{tabular}
\end{center}

The decisive column is the last.  Over this range $\lambda$ quadruples, so
an error of even $10^{-4}$ in $A_2$ would appear as a drift in $r_3$ of
order $3\times10^{-2}$ with a definite sign.  Instead $r_3$ is $-0.112$ at
$n=100$ and $-0.110$ at $n=400$, never leaving $[-0.125,-0.058]$ in
between, and oscillating rather than drifting --- which is what a periodic
$A_3$ requires.  Its sample mean over the thirteen rows is about $-0.088$,
against the $\Psi$-free part $-0.0824216430$ of $A_3$; the sample is not a
uniform phase average and carries an $A_4/\lambda$ term, so the gap is
expected.

The $A_2(\lambda)$ column also straddles the predicted mean
$1/(4L^2)=0.5203422452514019$, and the row-by-row differences
$r_2-A_2$ match the prediction $A_3(\lambda)/\lambda$ to about $10^{-6}$:
at $\{\lambda\}=0.2680$ the predicted difference is $-2.62\times10^{-4}$
against a measured $-2.611\times10^{-4}$, and at $\{\lambda\}=0.6748$,
$-2.71\times10^{-4}$ against $-2.698\times10^{-4}$.

Neither this table nor the preceding numerical calculations is machine-checked.  What is
machine-checked is listed in \cref{inverse:sec:module-concordance}.

\section{Two traps}\label{smallarg:sec:traps}

\begin{warning}[The phase trap]
Sampling $n=10,20,40,80,160,320$ to estimate a coefficient is invalid.  Doubling
$n$ adds about $1$ to $\log_2\lambda$, so $\{\lambda\}$ returns to nearly the
same place and a geometric sample holds the phase almost fixed.  Since every
coefficient is one-periodic, such a sample cannot see the oscillation at all;
it shows only the $\BigO(1/\lambda)$ drift in magnitude, and looks like
convergence to a constant.  This mistake was made independently on this project
before it was caught.  Always sweep $n$ densely and report against
$\{\lambda\}$, as in \cref{smallarg:sec:numerics}.
\end{warning}

\begin{warning}[The zeta trap]
Many $\zeta$ implementations route through the alternating Dirichlet eta
series, dividing by $1-2^{1-s}$.  At $s=1-\chi_k$ one has
$2^{1-s}=\EulerE^{2\pi\ImaginaryUnit k}=1$ exactly, so that prefactor has a pole at precisely
the arguments \cref{dyadicasym:eq:Psi-Fourier-report} needs.  The singularity is removable in $\zeta$
but not in that formula, and the result is either an overflow or a plausible
wrong number.  Concretely, \texttt{mpmath.zeta} at these arguments returns
values whose seventh significant digit drifts with the working precision:
\[
\begin{aligned}
  \texttt{mp.dps}=30:&\quad \FourierSeriesCoefficient{\Psi}{1}
    =7.55865005\times10^{-7}-7.47011001\times10^{-7}\,\ImaginaryUnit,\\
  \texttt{mp.dps}=50:&\quad \FourierSeriesCoefficient{\Psi}{1}
    =7.55864674\times10^{-7}-7.47009636\times10^{-7}\,\ImaginaryUnit,\\
  \texttt{mp.dps}=80:&\quad \FourierSeriesCoefficient{\Psi}{1}
    =7.55865060\times10^{-7}-7.47008702\times10^{-7}\,\ImaginaryUnit,
\end{aligned}
\]
whereas a hand-rolled Euler--Maclaurin $\zeta$ ($N=300$, $25$ Bernoulli terms)
is precision-stable and returns
$7.55864754097769\times10^{-7}-7.47010355898008\times10^{-7}\,\ImaginaryUnit$ at every
working precision.  Use Euler--Maclaurin or Hurwitz, and always check at two
precisions.
\end{warning}


\section{Specialist open items}
\label{inverse:sec:specialist-open}
\label{smallarg:sec:open}
The synthesis leaves the following claim-level obligations explicit:
\begin{enumerate}
\item Package \cref{inverse:thm:mass-composition,inverse:cor:log-composition} for the generic
  \texttt{expCoeff} and \texttt{logCoeff} recursions as exact general Lean declarations.
  This is the remaining formal composition layer, so it is worth recording what the
  implementation requires.  Mathlib has \texttt{Composition n} with a \texttt{Fintype}
  instance.  Writing \(\mathcal C(N)\) for the finite type of ordered compositions of \(N\)
  into positive parts, the missing ingredient is a first-block decomposition equivalence
  \[
   \mathcal C(n+1)\simeq
   \coprod_{b=1}^{n+1}\mathcal C(n+1-b),
  \]
  which is the form needed for induction against the recursion.

  The route that avoids \texttt{Composition} is to work directly with the \(k\)-fold
  finite-convolution powers \(E^{*k}\), as plain finite sums, and prove
  \[
   \operatorname{expCoeff}(E)_n
   =\sum_{k\le n}\frac1{k!}(E^{*k})_n
   \qquad(E_0=0).
  \]
  At coefficient level the product rule is immediate: writing
  \((\mathrm DA)_n=nA_n\),
  \[
   \mathrm D(A*B)=(\mathrm DA)*B+A*(\mathrm DB)
  \]
  is the identity \(n=j+(n-j)\) inside the convolution sum.  Promoting it to
  \[
   \mathrm D(E^{*(k+1)})=(k+1)\,E^{*k}*\mathrm DE
  \]
  also needs associativity of convolution, a double-sum reindexing not presently packaged
  for bare sequences.  Either route is a self-contained finite-sum combinatorics project;
  neither bridge is supplied merely by the displayed natural-language formulas.
\item formalize the transfer of the mass top-jet coefficient in
  \cref{inverse:thm:top-jet} through the logarithm and the resulting
  highest-\(\Psi\)-derivative coefficient at every order;
\item formalize or independently certify the displayed \(A_3\) and \(A_4\) formulas and
  their numerical checks;
\item surface effective constants in the fixed-order \(O(\lambda^{-N})\) remainders; and
\item prove or refute convergence, Borel summability, or a sharper optimal-truncation law
  for the formal correction series discussed in \cref{smallarg:sec:fourier}.
\end{enumerate}
No substitution of a truncated elementary approximation for the exact phase inside the
periodic coefficients is asserted.  No claim is made that an individual \(A_j\) is nowhere
zero, and all numerical tables remain evidence rather than proof.

\section{A proof-oriented concordance with the Lean modules}
\label{inverse:sec:module-concordance}

The following table records where the two missing clusters are proved.
Paths are relative to
\begin{center}
\path{Analysis/FabiusFunction/Lean/FabiusFunction/}.
\end{center}

\begingroup
\small
\begin{longtable}{>{\raggedright\arraybackslash}p{0.31\textwidth}
                  >{\raggedright\arraybackslash}p{0.49\textwidth}
                  >{\raggedright\arraybackslash}p{0.10\textwidth}}
\caption{Module concordance for this companion.}
\label{inverse:tab:module-concordance}\\
\toprule
Module & Principal role & Status\\
\midrule
\endfirsthead
\toprule
Module & Principal role & Status\\
\midrule
\endhead
\texttt{FabiusInverse.lean} &
Order isomorphism of $[0,1]$, totalized inverse, global continuity and
monotonicity, comparison equivalences, tails, reflection, exact dyadic
inversion, reciprocal first and second derivatives, interior smoothness,
strict closed-half curvature, transported flatness, effective steepness,
exact global differentiability and $C^\infty$ loci, and endpoint quotient divergence. & \Formalized\\
\texttt{InverseNotElementary.lean} &
Exact analytic locus of the totalized inverse and the analytic-density
obstructions used in its non-elementarity theorem. & \Formalized\\
\texttt{EffectiveMonotoneInverse.lean} &
Exact two-definition/six-theorem generic surface for the computable unit
clamp, tolerant difference certificates, fixed-depth dyadic bisection, and
sequential inversion on \([0,1]\) from a supplied positive computable inverse
modulus; the following module constructs such a certificate from rational gap
bounds. & \Formalized\\
\texttt{EffectiveGapInverse.lean} &
Exact eight-declaration surface deriving an inverse modulus, sequential
computability, and effective uniform continuity from computable positive
rational dyadic-gap lower bounds.  Its total theorem certifies precisely the
clamped extension, not arbitrary off-interval values of the inverse. & \Formalized\\
\texttt{FabiusInverseComputable.lean} &
The zero-definition/one-theorem specialization proving
\texttt{fabiusInv\_isComputableRealFunction} for the total clamped inverse. &
\Formalized\\
\texttt{FabiusSaddleJet}\\\texttt{ClosedForm.lean} &
Operator solution of the recursive periodic jets and the harmonic-number
constant term. & \Formalized\\
\texttt{FabiusSaddleJet}\\\texttt{Stirling.lean} &
Identification of the operator coefficients with
$s(n+1,m+1)$ and the corresponding closed derivative formula. &
\Formalized\\
\texttt{FabiusSaddleExponent}\\\texttt{ClosedForm.lean} &
Collapse of each centered exponent coefficient to one bounded jet
monomial plus one universal monomial; parity. & \Formalized\\
\texttt{FabiusSaddleExpansion}\\\texttt{Coefficients.lean} &
Formal exponential coefficients, Gaussian contraction, mass
coefficients, logarithmic transfer, periodicity, smoothness, and
boundedness. & \Formalized\\
\texttt{FabiusSaddleLeading}\\\texttt{Coefficient.lean} &
Affine dependence of $a_j$ on the newest jet and coefficient
$(-1)^j/(2^jj!)$. & \Formalized\\
\texttt{FabiusSecondSaddle}\\\texttt{Correction.lean} &
Order-four exponential identity, explicit $B_4$, Gaussian contraction,
$a_2$, $A_2$, and regularity of the second correction. & \Formalized\\
\texttt{FabiusSecondSaddle}\\\texttt{Expansion.lean} &
Insertion of $A_2$ into the exact-phase asymptotic with a
third-order remainder. & \Formalized\\
\texttt{FabiusSaddleCentral}\\\texttt{AllOrders.lean} &
Exact central decomposition, finite reference polynomial, odd-term
cancellation, and integrated central remainders. & \Formalized\\
\texttt{FabiusSaddleTail}\\\texttt{AllOrders.lean} &
Order-dependent radius and complementary intermediate/far contour bounds
of arbitrary fixed inverse-power order. & \Formalized\\
\texttt{FabiusSaddleMass}\\\texttt{AllOrders.lean} &
Assembly of central Taylor control, finite exponential substitution,
whole-line Gaussian contraction, and minor arcs into the complete mass
expansion. & \Formalized\\
\texttt{FabiusFullAsymptotic}\\\texttt{Expansion.lean} &
Transfer from mass to logarithm, addition of the flat forward tail, and
transport from $x=2^{-t}$ to the full limit $x\downarrow0$. & \Formalized\\
\bottomrule
\end{longtable}
\endgroup

\subsection{Named declarations worth knowing}

The structural uniform-modulus layer is the exhaustive public surface of
\path{InverseModulus.lean}: one definition,
\path{Fabius.fabiusIntervalMass}, and thirty-one theorems,
\path{Fabius.fabiusIntervalMass_reflect},
\path{Fabius.fabiusIntervalMass_eq_zero_of_add_nonpos},
\path{Fabius.fabiusIntervalMass_eq_zero_of_one_le},
\path{Fabius.monotoneOn_fabiusIntervalMass_firstHalf},
\path{Fabius.antitoneOn_fabiusIntervalMass_secondHalf},
\path{Fabius.strictMonoOn_fabiusIntervalMass_firstHalf},
\path{Fabius.strictAntiOn_fabiusIntervalMass_secondHalf},
\path{Fabius.fabiusReal_sub_le_sub},
\path{Fabius.fabiusReal_lt_fabiusIntervalMass_of_mem_Ioo},
\path{Fabius.fabiusIntervalMass_eq_fabiusReal_iff},
\path{Fabius.fabiusReal_sub_lt_sub},
\path{Fabius.fabiusReal_sub_eq_sub_iff},
\path{Fabius.fabiusReal_add_le},
\path{Fabius.fabiusReal_add_eq_iff},
\path{Fabius.fabiusInv_sub_le_sub_of_mem_Icc},
\path{Fabius.fabiusInv_sub_eq_sub_iff_of_mem_Icc},
\path{Fabius.fabiusInv_sub_le_sub_of_le},
\path{Fabius.fabiusInv_sub_le_sub},
\path{Fabius.fabiusInv_add_le},
\path{Fabius.abs_fabiusInv_sub_le},
\path{Fabius.abs_fabiusInv_sub_eq_iff_of_mem_Icc},
\path{Fabius.dist_fabiusInv_le},
\path{Fabius.fabiusInv_min_one},
\path{Fabius.isGreatest_abs_fabiusInv_sub},
\path{Fabius.sSup_abs_fabiusInv_sub_eq},
\path{Fabius.isGreatest_abs_fabiusInv_sub_Icc},
\path{Fabius.sSup_abs_fabiusInv_sub_Icc_eq},
\path{Fabius.abs_fabiusInv_sub_lt_of_abs_sub_lt_fabiusReal},
\path{Fabius.abs_fabiusInv_sub_le_of_abs_sub_le_fabiusReal},
\path{Fabius.fabiusReal_le_abs_sub_of_le_abs_fabiusInv_sub}, and
\path{Fabius.forall_abs_fabiusInv_sub_lt_iff}.
For nonnegative length the fixed-length Fabius increment is reflection-invariant,
vanishes on both constant tails, and has the global weak increase--decrease shape.  When
\(0<h\leq1\), the maximal support-meeting branches \([-h,(1-h)/2]\) and
\([(1-h)/2,1]\) are respectively strictly increasing and strictly decreasing.  On the
unit interval, the endpoint minimum has exact strict-interior and degenerate/endpoint
equality loci, and constrained superadditivity has its full equality classification.  The
ordered and order-free inverse-gap bounds include complete ordered and absolute equality
loci on \([0,1]\), while clamping creates additional cases and therefore precludes the
corresponding global equality claims.  Subadditivity, absolute and metric self-moduli,
saturation at one, and attained exact \path{IsGreatest}/\path{sSup} moduli hold globally
and on the unit interval.  The final four declarations are the strict and closed
effective-injectivity bounds, their contrapositive separation form, and the uniform
strict-threshold biconditional.  This module does not supply the report's closed
\(\Delta_r\) lower bound, recursive-modulus wrapper, tolerant-bisection realizer,
sequential computability theorem, combined computable-real-function theorem, or input-bit
asymptotics.

Those realizer conclusions are supplied downstream.  The exhaustive public
surface of \path{EffectiveMonotoneInverse.lean} is the two definitions
\path{Fabius.SequentiallyComputableOn} and \path{Fabius.unitClamp}, and the
six theorems
\path{Fabius.unitClamp_sequentiallyComputable},
\path{Fabius.tolerantDifference_error},
\path{Fabius.tolerantDifference_safe_updates},
\path{Fabius.tolerantDifference_inconclusive},
\path{Fabius.tolerantBisection_correct}, and
\path{Fabius.effectiveInversionOn_Icc}.  Its natural recursion has the
requested output precision \(p\) as a fixed iteration bound.  A live bracket
index at depth \(s\) names the adjacent dyadic interval of width \(2^{-s}\);
a certified sign halves it, whereas an inconclusive comparison freezes a
certified midpoint and doubles its numerator through the remaining grids.
The final code is therefore on the \(2^{-p}\) grid whether or not the success
branch was taken.  The comparison scale \(q=p+d(p)+3+s\) absorbs both oracle
errors, and the four-unit signed-numerator margin implements the three-way
comparison without deciding equality of represented reals.

The hypotheses of \path{Fabius.effectiveInversionOn_Icc} include the
computable positive denominator \(d\) and the strict inverse-modulus
implication at \(d(p)^{-1}\).  This reusable theorem consumes that certificate.
The exhaustive four-definition/four-theorem public surface of
\path{EffectiveGapInverse.lean} is
\path{Fabius.EffectivelyUniformContinuousOn}, the structure
\path{Fabius.ComputablePositiveRationalSequence},
\path{Fabius.ComputablePositiveRationalSequence.value},
\path{Fabius.ComputablePositiveRationalSequence.reciprocalDenominator},
\path{Fabius.ComputablePositiveRationalSequence.reciprocalDenominator_spec},
\path{Fabius.inverseModulus_of_positiveRationalGap},
\path{Fabius.effectiveInversionOn_Icc_of_computablePositiveRationalGap}, and
\path{Fabius.clampedEffectiveInversion_of_computablePositiveRationalGap}.
The structure supplies computable positive natural numerators and denominators,
and the reciprocal-denominator specification gives a computable positive
natural denominator with reciprocal strictly below the represented value.  A
strict increasing inverse pair on \([0,1]\) must satisfy
\[
 \alpha(p)\le f(x+2^{-p})-f(x)
 \qquad\text{for every }x\in[0,1-2^{-p}].
\]
The first theorem converts that bound into the inverse modulus.  With a
computable dyadic approximation of \(f\) and both interval maps, the next
theorem proves sequential computability and effective uniform continuity of
\(g\) on the unit interval.  The last proves
\path{IsComputableRealFunction} only for
\(x\mapsto g(\lean{unitClamp}(x))\), which agrees with \(g\) on \([0,1]\) but
does not constrain its unclamped values outside that interval.  The separate
\path{FabiusInverseComputable.lean} has zero public definitions and the single
public theorem \path{Fabius.fabiusInv_isComputableRealFunction}.  It
instantiates the generic theorem with the centered-spline evaluator and Delta
denominator, clamps arbitrary computable input names, and imports the
logarithmic-Delta effective-continuity theorem.  Hence the conclusion is the
full \path{IsComputableRealFunction} statement for the totalized inverse.
The fixed-depth module and Fabius specialization expose exactly nine
declarations; the rational-gap module adds the eight declarations above.  No
input-bit or running-time asymptotic is asserted.

The inverse interface is centered around
\begin{lstlisting}[style=pseudo]
fabiusOrderIso
fabiusReal_fabiusInv
fabiusInv_fabiusReal
strictMonoOn_fabiusInv
monotone_fabiusInv
continuous_fabiusInv
fabiusInv_le_iff_le_fabiusReal
fabiusReal_le_iff_le_fabiusInv
fabiusInv_lt_iff_lt_fabiusReal
fabiusReal_lt_iff_lt_fabiusInv
fabiusInv_one_sub
fabiusInv_fabiusDyadicUnit
fabiusInv_differentiableAt_iff
fabiusInv_contDiffAt_infty_iff
fabiusInv_analyticAt_iff
fabiusInv_not_isElementaryOrInverse
le_two_pow_mul_fabiusInv_pow
le_two_pow_div_factorial_mul_fabiusInv_pow
mul_lt_fabiusInv
tendsto_fabiusInv_div_atTop
tendsto_one_sub_fabiusInv_div_one_sub_atTop
\end{lstlisting}
The all-orders endpoint interface culminates in
\begin{lstlisting}[style=pseudo]
fabiusSaddleKernelMass_dyadicLambert_hasAsymptoticExpansion
fabiusSaddleRatio_dyadicLambert_hasAsymptoticExpansion
log_fabius_dyadicReal_sub_sharpLambertMain_hasAsymptoticExpansion
log_fabius_sub_sharpLambertMain_hasAsymptoticExpansion
log_fabius_sub_sharpLambertExpansion_isBigO
\end{lstlisting}

These names are not needed to read the mathematics, but they make the
article useful as a map back into the source tree.

\section{What is formalized and what is newly synthesized}
\label{inverse:sec:status-ledger}

For clarity, the principal statements can be divided into three groups.

\subsection{Directly formalized}

The following content is represented directly by declarations in the
repository:
\begin{itemize}
\item the global totalized inverse package and all four comparison
equivalences;
\item the constant tails, reflection, midpoint, and exact dyadic
inversion;
\item interior $C^\infty$ regularity, the exact formulas for $\FabiusClampedQuantile'$ and
$\FabiusClampedQuantile''$, and strict concavity on $[0,1/2]$ and strict convexity on
$[1/2,1]$;
\item the effective and sharp inverse power inequalities, effective
linear steepness, and the two endpoint quotient limits;
\item the exact global differentiability, $C^\infty$, and analytic loci
of the totalized inverse;
\item the closed jet formulas, shifted Stirling conversion, and closed
exponent coefficients;
\item the formal exponential and logarithmic recurrence machinery,
the mass top-jet coefficient, and the first and second corrections;
\item the arbitrary-order central, tail, mass, and logarithmic
asymptotic theorems.
\end{itemize}

\subsection{Classical consequences derived in this article}

The following statements are proved here from the formalized input:
\begin{itemize}
\item divergence of $\FabiusClampedQuantile(y)/y^\alpha$ for every $\alpha>0$;
\item the sharp equivalent \eqref{inverse:eq:sharp-inverse} and the scale
hierarchy \eqref{inverse:eq:power-below-G}--\eqref{inverse:eq:G-below-log};
\item the exact right-endpoint transfer
\eqref{inverse:eq:right-endpoint-all-orders}, the separate Lambert-phase expansion
\eqref{inverse:eq:lambert-phase-all-orders}, and the period-mean identities
\eqref{smallarg:eq:mean-A1}--\eqref{smallarg:eq:mean-A2};
\item the explicit all-$j$ composition sum \eqref{inverse:eq:mass-composition}
and the highest-$\Psi$-derivative corollary
\eqref{inverse:eq:highest-derivative}.
\end{itemize}

\subsection{Symbolic and numerical frontier}

The expanded \(A_3\) and \(A_4\) formulas, their derivative forms and
\(\Psi\)-free parts, the Fourier-amplitude estimates, and the residual tables in
\cref{inverse:sec:higher-explicit,smallarg:sec:fourier,smallarg:sec:numerics} are retained as
symbolic or numerical evidence.  They are calibrated against the formalized recursions and
the first two corrections, but no exact current Lean declaration certifies the displayed
higher formulas or their numerical verification.

\subsection{Claims deliberately not made}

This companion does \emph{not} claim:
\begin{itemize}
\item convergence or Borel summability of the all-orders series;
\item an all-orders substitution of an approximate Lambert phase inside
the periodic coefficients;
\item a closed form for the inverse at arbitrary rational arguments;
\item analyticity of either $F$ or $\FabiusClampedQuantile$ inside the transition interval;
\item uniformity of the remainder constants as the truncation order tends
to infinity.
\end{itemize}

\section{Concluding synthesis}

The inverse function and the saddle coefficient machinery illuminate
opposite sides of the same endpoint phenomenon.

On the geometric side, $F$ approaches zero more rapidly than every
power.  Its inverse must therefore leave zero more rapidly than every
power, and the formalized comparison calculus turns this slogan into
effective inequalities with exact constants.  The sharp Lambert
expansion goes further: the inverse lives on the
stretched-logarithmic scale
\[
 \FabiusClampedQuantile(y)\asymp
 \sqrt{\log(1/y)}
 \exp\!\left(-\sqrt{2\log2\,\log(1/y)}\right),
\]
with the exact leading constant $1/(\EulerE\sqrt{\log2})$.  Reflection transfers
the entire picture to the endpoint $1$.

On the asymptotic side, the apparent complexity of the periodic
corrections is organized by three simple structures.  First, the
recursive jets are solved by a product of commuting first-order
operators; the forcing terms sum to a harmonic number and the derivative
coefficients are shifted signed Stirling numbers.  Second, centering
collapses every exponent coefficient to two monomials, forcing parity and
eliminating all odd half-powers under Gaussian contraction.  Third, the
newest jet can enter only linearly, so the coefficient of the highest
derivative $\Psi^{(2j)}$ is known at every order before the lower-order
algebra is performed.

The analytic remainder proof then explains why these formal coefficients
are not merely symbolic.  The growing window
$A_N(b)=\sqrt{32(N+1)\log b}$ separates the local and global tasks:
locally it is still narrow enough for uniform Taylor expansion, while
globally it is wide enough to make both Gaussian and contour tails smaller
than the requested inverse power.  Finite exponential substitution and
logarithmic transfer complete a genuine Poincar\'e expansion at the exact
Lambert phase.

Together, these missing parts close two expository loops.  Flatness of
$F$ is now matched by a precise theory of the steep inverse, and the
headline all-orders expansion is now matched by a transparent account of
the coefficient algebra and the analytic estimates that make it true.

\clearpage
\section*{Supplementary material from this synthesis}
\addcontentsline{toc}{section}{Supplementary material from this synthesis}
\section{Repository paths used in the audit}
\label{inverse:app:paths}

\begin{lstlisting}[style=pseudo]
Analysis/FabiusFunction/
  Lean/FabiusFunction/FabiusInverse.lean
  Lean/FabiusFunction/FabiusSaddleJetClosedForm.lean
  Lean/FabiusFunction/FabiusSaddleJetStirling.lean
  Lean/FabiusFunction/FabiusSaddleExponentClosedForm.lean
  Lean/FabiusFunction/FabiusSaddleLeadingCoefficient.lean
  Lean/FabiusFunction/FabiusSecondSaddleCorrection.lean
  Lean/FabiusFunction/FabiusSecondSaddleExpansion.lean
  Lean/FabiusFunction/FabiusSaddleCentralAllOrders.lean
  Lean/FabiusFunction/FabiusSaddleTailAllOrders.lean
  Lean/FabiusFunction/FabiusSaddleMassAllOrders.lean
  Lean/FabiusFunction/FabiusFullAsymptoticExpansion.lean
  docs/Fabius_Function_and_Rvachev_Up/
    Fabius_Function_and_Rvachev_Up.tex
\end{lstlisting}

\renewcommand{\refname}{Topic references}
\begin{thebibliography}{99}

\bibitem{inverse:fabius1966}
J.~Fabius,
\newblock A probabilistic example of a nowhere analytic $C^\infty$-function,
\newblock \emph{Zeitschrift f\"ur Wahrscheinlichkeitstheorie und Verwandte
Gebiete} \textbf{5} (1966), 173--174.
\newblock DOI: \nolinkurl{10.1007/BF00536652}.

\bibitem{inverse:rvachev1990}
V.~A.~Rvachev,
\newblock Compactly supported solutions of functional-differential equations
and their applications,
\newblock \emph{Russian Mathematical Surveys} \textbf{45} (1990), no.~1,
87--120.

\bibitem{inverse:arias1982}
J.~Arias de Reyna,
\newblock Definici\'on y estudio de una funci\'on indefinidamente
diferenciable de soporte compacto,
\newblock \emph{Revista de la Real Academia de Ciencias Exactas, F\'isicas y
Naturales de Madrid} \textbf{76} (1982), no.~1, 21--38.
\newblock English translation:
\emph{An infinitely differentiable function with compact support:
Definition and properties}, arXiv:\nolinkurl{1702.05442} (2017).

\bibitem{inverse:arias2018}
J.~Arias de Reyna,
\newblock Arithmetic of the Fabius function,
\newblock \emph{INTEGERS} \textbf{18} (2018), Article A51;
arXiv:\nolinkurl{1702.06487}, version~3.

\bibitem{inverse:corless1996}
R.~M.~Corless, G.~H.~Gonnet, D.~E.~G.~Hare, D.~J.~Jeffrey, and D.~E.~Knuth,
\newblock On the Lambert $W$ function,
\newblock \emph{Advances in Computational Mathematics}
\textbf{5} (1996), 329--359.

\bibitem{inverse:proveit}
V.~Reshetnikov,
\newblock \emph{ProveIt: Analysis/FabiusFunction},
\newblock Lean~4 formal development, repository snapshot 25 August 2026.
\newblock
\url{https://github.com/VladimirReshetnikov/ProveIt/tree/main/Analysis/FabiusFunction}.

\bibitem{inverse:maincompanion}
V.~Reshetnikov,
\newblock \emph{The Fabius Function and Rvachev's Up-Function},
\newblock repository document in
\path{Analysis/FabiusFunction/docs/Fabius_Function_and_Rvachev_Up},
snapshot 25 August 2026.

\end{thebibliography}


\renewcommand{\arraystretch}{1}
% END SYNTHESIZED TOPIC: inverse and small-argument saddle analysis

% BEGIN FRONTIER PART: general-nome value identity (added 2026-09-02)
\clearpage
\part{The General-Nome Value Identity for the Geometric Fabius Recursion}

\begin{boxedremark}
\textbf{Status.}  Everything in this part is a frontier claim: numerically
verified in exact rational arithmetic, not formalized.  The two negative
results of \Cref{genq:sec:negative} \emph{are} settled, in the sense that a
single exact counterexample refutes each; they are recorded here because both
refuted statements are the ones a reader would guess.  Added 2 September 2026.
\end{boxedremark}

\section{The dyadic identity and the shape of its deformation}
\label{genq:sec:statement}

The primary exposition's central dyadic result evaluates the Fabius function at
the points \(2^{-n}\) through a \(q\)-binomial--Thue--Morse sum at the fixed
nome \(q=1/2\):
\begin{equation}
\label{genq:eq:dyadic}
  F\!\left(2^{-n}\right)
  =\frac{1}{2^{n^{2}}\QPochhammer{1/2}{1/2}{n}}
   \sum_{k=0}^{n}
     \frac{\GaussianBinomial nk{1/2}}{4^{\Btri k}\,(n+k)!}
     \sum_{r<2^{k}}\ThueMorseSign r\,\bigl(r-2^{k}\bigr)^{n+k}.
\end{equation}
This is formalized: \lean{fabiusAtInverseTwoPow_eq_qBinomialThueMorse_sum} in
\path{FabiusQBinomialFormula.lean}.  Its prefactor is now also identified with a
group order, \(2^{n^{2}}\QPochhammer{1/2}{1/2}{n}=\lvert\mathrm{GL}_{n}(\mathbb
F_{2})\rvert\), by \lean{two_pow_nat_sq_mul_qPochhammer_half_eq_card_GL} in
\path{FabiusGeneralLinearDenominator.lean}.

The nome \(1/2\) is written into \eqref{genq:eq:dyadic} in three separate
places, and it is not obvious which of them are essential.  Two of the three are
not.  The outer Gaussian weights are governed by the finite \(q\)-binomial
theorem, which the atlas proves over an arbitrary commutative ring
(\lean{finite_qBinomial_theorem}, \path{FiniteQBinomialCore.lean}); and the
inner signed block is governed by a powerset annihilation engine stated for an
\emph{arbitrary} family of step weights over an arbitrary commutative ring
(\lean{sum_powerset_neg_one_pow_eval} and
\lean{sum_powerset_neg_one_pow_pow_card}, \path{ThueMorseSparseProuhet.lean}),
in which the dyadic steps \(2^{j}\) are never used.  Only the left-hand side is
genuinely base-two: \lean{halfMoment} hard-codes the denominator
\((n+2)(2^{n+1}-1)\), and the formal-series engine of the proof rescales by
\(2\) throughout.

The deformation below replaces the steps \(2^{j}\) by \(q^{-j}\) and the
half-moment recursion by its \(q\)-analogue.

\begin{definition}[The \(q\)-deformed data]
\label{genq:def:data}
Let \(K\) be a field of characteristic zero and let \(q\in K\) satisfy
\(q^{m}\neq1\) for \(1\le m\le n\); equivalently
\(\QPochhammer qqn\neq0\), which in particular forces \(q\neq0,1\).  Define
\begin{align}
  d_{q}(0)&=1, &
  d_{q}(n+1)&=\frac{\displaystyle\sum_{k=0}^{n}\binom{n+2}{k}d_{q}(k)}
                  {(n+2)\bigl(q^{-(n+1)}-1\bigr)}, &
  a_{q}(n)&=\frac{d_{q}(n)}{n!},
\end{align}
and, for \(k,d\in\NaturalNumbers\) and \(x\in K\), the \emph{\(q\)-Prouhet
block}
\begin{equation}
\label{genq:eq:block}
  P_{q}(k,d,x)=\sum_{T\subseteq\{0,\dots,k-1\}}(-1)^{\lvert T\rvert}
    \Bigl(x+\sum_{j\in T}q^{-j}\Bigr)^{d}.
\end{equation}
\end{definition}

At \(q=1/2\) the Mersenne factor reappears, since
\(q^{-(n+1)}-1=2^{n+1}-1\); this is precisely why
\(\QPochhammer{1/2}{1/2}{n}\) occurs in \eqref{genq:eq:dyadic}, and
\(d_{1/2}=\lean{halfMoment}\), \(a_{1/2}=\lean{fabiusRecurrenceSequence}\).

\begin{candidate}[General-nome value identity]
\label{genq:cand:main}
With the data of \Cref{genq:def:data}, for every \(n\in\NaturalNumbers\),
\begin{equation}
\label{genq:eq:general}
  q^{\Btri n}\,a_{q}(n)
  =\frac{1}{q^{-n^{2}}\QPochhammer qqn}
   \sum_{k=0}^{n}
     \frac{\GaussianBinomial nkq}{q^{-2\Btri k}\,(n+k)!}\,
     P_{q}\Bigl(k,\;n+k,\;-\sum_{j<k}q^{-j}\Bigr).
\end{equation}
\end{candidate}

Note the exponent \(q^{+\Btri n}\) on the left, and that the inner block is
centred at \(-\sum_{j<k}q^{-j}\), the negative of the largest node of the
block, rather than at any single power of \(q\).  Both points are sharp; see
\Cref{genq:sec:negative}.

\section{Numerical evidence}
\label{genq:sec:evidence}

\Cref{genq:cand:main} was checked in exact rational arithmetic --- no floating
point --- for \(q\in\{1/2,\;1/3,\;2/5,\;3/4,\;1/4\}\) and \(0\le n\le6\), both
sides agreeing as rationals in all thirty-five cases.  At \(q=1/2\) the left
side reproduces the Fabius dyadic values
\[
  1,\quad \tfrac12,\quad \tfrac5{72},\quad \tfrac1{288},\quad
  \tfrac{143}{2073600},\quad \tfrac{19}{33177600},
\]
which are exactly the values \eqref{genq:eq:dyadic} produces, so
\eqref{genq:eq:general} contains the formalized dyadic identity as its
\(q=1/2\) specialization.

That specialization is a corollary and not a restatement.  Under the Boolean
cube reindexing \(T\mapsto r=\sum_{j\in T}2^{j}\)
(\lean{sum_powerset_two_pow}), \(P_{1/2}(k,d,x)\) becomes
\(\sum_{r<2^{k}}\ThueMorseSign r\,(x+r)^{d}\) with \(x=-(2^{k}-1)\), whereas
\eqref{genq:eq:dyadic} uses \(x=-2^{k}\).  The two centrings differ term by
term and agree only after the full \(\GaussianBinomial nkq\)-weighted sum over
\(k\) has been taken; a centring-shift lemma is therefore part of the
obligation below.

\section{Two refutations}
\label{genq:sec:negative}

Both statements refuted here are the natural guesses, which is why they are
recorded rather than passed over.

\begin{warning}[The naive centring is wrong off the dyadic point]
\label{genq:warn:centring}
Replacing the centring \(-\sum_{j<k}q^{-j}\) in \eqref{genq:eq:general} by the
single power \(-q^{-k}\) --- the literal transcription of the \(-2^{k}\) of
\eqref{genq:eq:dyadic} --- yields a statement that is true for every \(n\) at
\(q=1/2\) and false at every \(n\ge1\) for \(q=1/3\) and \(q=3/4\), by exact
computation.  The dyadic case is thus accidental: at \(q=1/2\),
\(\sum_{j<k}2^{j}=2^{k}-1\) and the discrepancy is absorbed by the outer sum.
\end{warning}

\begin{warning}[The left-hand side is not a distribution function]
\label{genq:warn:cdf}
It is tempting to read \(q^{\Btri n}a_{q}(n)\) as the geometric-uniform
distribution function \lean{geometricUniformCDF} evaluated at \(q^{n}\),
because that reading is correct at \(q=1/2\), where it is the Fabius
CDF.  It is false in general: at \(q=3/4\) and \(n=1\) the recursion gives
\(d_{q}(1)=\tfrac32\) and \(\Btri1=0\), so the left side of
\eqref{genq:eq:general} equals \(3/2>1\) and cannot be the value of any
distribution function.  Whether some renormalization of \(a_{q}\) admits a
probabilistic reading is open; the naive one does not.
\end{warning}

\Cref{genq:warn:cdf} is worth stating explicitly because the corresponding
dyadic statement is both true and formalized
(\lean{geometricUniformCDF_one_half_eq_fabiusReal},
\path{ProbabilityRepresentation.lean}), so the false generalization inherits an
unusually strong appearance of support.

\section{What remains to be formalized}
\label{genq:sec:obligation}

\begin{obligation}[General-nome value identity]
\label{genq:obl:main}
Prove \eqref{genq:eq:general} in Lean.  Three components are already available
at the required generality and need only to be instantiated:
\begin{enumerate}[label=(\roman*)]
  \item the outer Gaussian weights, by \lean{finite_qBinomial_theorem} over an
        arbitrary commutative ring;
  \item the vanishing of \(P_{q}(k,d,x)\) for \(d<k\) and its sharp value at
        \(d=k\), by \lean{sum_powerset_neg_one_pow_eval} and
        \lean{sum_powerset_neg_one_pow_pow_card} at the step weights
        \(w_{j}=q^{-j}\), whose product \(\prod_{j<k}q^{-j}=q^{-\Btri k}\)
        supplies the denominator \(q^{-2\Btri k}\) after squaring;
  \item the centring-shift lemma of \Cref{genq:sec:evidence}.
\end{enumerate}
The genuinely missing components are:
\begin{enumerate}[label=(\alph*)]
  \item a definition of \(d_{q}\) and \(a_{q}\) --- the atlas has only the
        base-two \lean{halfMoment} and \lean{fabiusRecurrenceSequence}, whose
        denominator \((n+2)(2^{n+1}-1)\) is hard-coded; and
  \item a nome-general replacement for the formal-series engine of
        \path{FabiusQBinomialFormula.lean}, whose
        \lean{recurrenceSeries_rescale_two} rescales by \(2\) and whose
        \lean{refinementFactorSeries} is a product over
        \(\mathrm{rescale}(-2^{j})\).  The abstract coefficient extractor
        \lean{weighted_factor_coefficient} in the same file is already stated
        for an abstract weight, node and annihilation family, so its
        \emph{statement} should transfer unchanged.
\end{enumerate}

\begin{warning}[Three of those declarations are \texttt{private}]
\label{genq:warn:private}
\lean{recurrenceSeries_rescale_two}, \lean{refinementFactorSeries} and
\lean{weighted_factor_coefficient} are all declared \texttt{private} in
\path{FabiusQBinomialFormula.lean}, so none of them is accessible from another
module as things stand.  Component (b) therefore begins with a step this
register originally omitted: making those three public, or relocating the
abstract extractor to a module the general-nome development can import.  That
is bookkeeping rather than mathematics, but it is not nothing --- the extractor
is the piece the whole route was supposed to reuse, and "reuse it unchanged" is
not available until it is visible.
\end{warning}
\end{obligation}

\begin{obligation}[Probabilistic reading, if any]
\label{genq:obl:prob}
Decide whether \(a_{q}\) admits a probabilistic interpretation generalizing the
Fabius one.  \Cref{genq:warn:cdf} refutes the naive candidate, so this
obligation is to find the correct normalization or to show that none exists.
Note that \path{GeometricUniformCDF.lean} already develops
\lean{geometricUniformCDF} at a general nome with its density, support and
smoothness, so only the closed form at the grid points \(q^{n}\) is missing.
\end{obligation}
% END FRONTIER PART: general-nome value identity

% BEGIN POST-CONSOLIDATION REGISTER: Primary_Exposition_Gap_Register/Primary_Exposition_Gap_Register.tex
% Original source SHA-256: 06c7b888d9601b67ad7a5c0aee3f087d44d9ecaaa9abfb3bf3edfda4bd29c0d1
\clearpage
\part{Primary Exposition Gap Register}
\begin{boxedremark}
\textbf{Source provenance and status.}
Former path:
\nolinkurl{Primary_Exposition_Gap_Register/Primary_Exposition_Gap_Register.tex}.
This post-consolidation register preserves smaller claims removed during the primary
exposition's exact-counterpart audit.  Candidate statements below are of two kinds.  A
candidate marked \emph{discharged} carries a named Lean citation in the surrounding text:
a public declaration now matches its hypotheses, normalization, domain and conclusion, and
its former formalization obligation has been deleted.  A candidate carrying no such
citation is still open; it is not claimed as a Lean theorem, and it still needs a public
theorem with the stated hypotheses and conclusion.  Declarations identified as formal
background support an open candidate without stating it.

This register is audited against a library that keeps moving.  Six of its original
obligations were discharged by formalization work that landed within days of the audit date:
one on the audit date itself, four the day after, one the day after that.  The entries below
record that outcome.
The remaining open entries carry the same exposure: they are accurate as of the audit date
recorded under \emph{Provenance} at the end of this part and can be overtaken in the same
way, so re-check them against the current library before relying on them.
\end{boxedremark}

\section{Fourier-decay normalization}

The source proves polynomially weighted decay.  The conventional shifted-weight form was
removed from the primary exposition for want of a named counterpart; the library has since
supplied one, so the statement is recorded here as proved rather than as a gap:
\begin{candidate}[Discharged]
For every \(N\in\NaturalNumbers\), there is \(C_N>0\) such that
\[
 |\FourierTwoPi{\RvachevUp}(\xi)|\le C_N(1+|\xi|)^{-N}
 \qquad(\xi\in\RealNumbers).
\]
\end{candidate}
It is standardly equivalent to the proved polynomially weighted estimate after a
bounded-region argument and a change of constant, and it is now also a literal named Lean
theorem: \lean{rvachevFourier_real_shiftedDecay} in \path{PoissonSummation.lean} states
exactly the displayed bound for every exponent, universally quantified over all real
arguments, so the range \(|\xi|\le1\) is covered by the statement itself rather than by a
separate argument.  The formal library proves strictly more than the candidate:
\lean{rvachevFourier_real_iteratedDeriv_shiftedDecay} in the same module gives the same
shifted-weight majorant for every iterated derivative of the transform, and the displayed
estimate is its order-zero case.

\section{Dyadic specializations and evaluator cost}

\subsection{Concrete rational values}

Let \(F\) denote the bounded Fabius function and \(\RvachevUp\) Rvachev's compactly supported
up-function.  Write \(c_n=\int_{-1}^{1}x^{2n}\RvachevUp(x)\Differential x\).  The general dyadic evaluator
is proved correct, and these specializations now have named theorems:
\begin{candidate}[Discharged]
\[
 c_1=\frac19,\qquad
 F(1/4)=\frac5{72},\qquad
 F(1/8)=\frac1{288},\qquad
 F(3/8)=\frac{73}{288}.
\]
\end{candidate}
\begin{obligation}[Discharged]
Prove named specializations by reducing the exact evaluator and bridge each result to the
analytic function.
\end{obligation}
All four values now have named theorems in \path{DyadicSpecializations.lean}, each in the
two requested forms.  The executable reductions are \lean{moment_one} (`\(\mu_1=1/9\)' from
the defining recurrence), \lean{fabiusDyadic_two_one}, \lean{fabiusDyadic_three_one}, and
\lean{fabiusDyadic_three_three}, closed by rational normalization of the dyadic formula.
The analytic bridges are \lean{integral_sq_mul_rvachev_eq_one_ninth} together with its
interval form \lean{intervalIntegral_sq_mul_rvachev_eq_one_ninth} (the display's
\(c_1\)), and \lean{fabiusReal_one_quarter}, \lean{fabiusReal_one_eighth},
\lean{fabiusReal_three_eighths}, valid for every bounded Fabius function via
\lean{fabiusDyadic_cast} and \lean{integral_even_pow_mul_rvachev_eq_moment}.

\subsection{Popcount complexity}

Termination is formalized, but the removed prose did not distinguish the following two
possible complexity claims.  For \(e,a\in\NaturalNumbers\) with \(0<a<2^e\), let \(D(e,a)\) be the
number of recursive calls that enter the evaluator's interior branch.
\begin{candidate}[Discharged]
The exact-count candidate is \(D(e,a)=\BinaryDigitSum(a)\).
\end{candidate}
\begin{candidate}[Discharged]
The weaker upper-bound candidate is \(D(e,a)\le\BinaryDigitSum(a)\).
\end{candidate}
\begin{obligation}[Discharged]
Define \(D\) in Lean, settle which statement matches the executable recursion, and prove
the resulting relation to binary weight, including endpoint and saturation branches.
\end{obligation}
Both candidates are settled in \path{EvaluatorPopcount.lean}.  The count is
\lean{dyadicInteriorCalls}, defined by literally the branch structure of
\lean{fabiusDyadicUnitAux} (the same guards and the same leading-bit-stripping recursion
argument, so each interior entry is counted once).  The verdict: \emph{the exact count
matches the executable recursion} — \lean{dyadicInteriorCalls_eq_binaryWeight} proves
\(D(e,a)=\BinaryDigitSum(a)\) for every \(a<2^e\), the endpoint \(a=0\) included, since each interior
step strips exactly the leading binary digit
(\lean{binaryWeight_add_pow_two}).  On the saturation branch \(2^e\le a\) the evaluator
answers without recursing, \(D(e,a)=0\)
(\lean{dyadicInteriorCalls_of_saturated}), and only there does the exact statement degrade
to the inequality; the weaker candidate is the unconditional corollary
\lean{dyadicInteriorCalls_le_binaryWeight}.  The all-ones worst case
\(D(e,2^e-1)=e\) is \lean{dyadicInteriorCalls_two_pow_sub_one}, via
\lean{binaryWeight_two_pow_sub_one}.

\subsection{Independent analytic derivations}

Earlier prose derived the dyadic formula through Taylor integrals, block substitutions,
and the exact remainder \(-F(y)\).  The final rational identities are formal, but those
intermediate analytic equations are not public theorem-for-theorem counterparts.
\begin{obligation}[Discharged]
Formalize named Taylor integral identities for the bounded and signed global functions,
including the blockwise changes of variables and the exact remainder.
\end{obligation}
The intermediate analytic equations are now public theorems.  In
\path{DyadicAnalytic.lean}: the dyadic Taylor sum \lean{analyticTaylorSum} with coefficient
data \lean{analyticInverseValue} at the scales \lean{inverseTwoPowReal}; the Taylor-level
Fabius differential equation \lean{analyticTaylorSum\_derivative\_identity}; the exact-remainder
identity \lean{analytic\_scale\_recurrence} (\(F(2^{-m}+y)=\mathrm{Taylor}_m(y)-F(y)\) for
\(0\le y\le2^{-m}\), the register's remainder \(-F(y)\) with no error term); and the
blockwise change of variables \lean{fabiusReal\_dyadic\_bit\_recurrence}, the identity behind
every interior step of the exact evaluator.  The signed global function satisfies the same
law on the same window: \lean{extendedFabius\_scale\_recurrence} in \path{GlobalDyadic.lean},
via agreement with the bounded function on the unit interval.

\subsection{Global block fold for the signed evaluator}

Let \(n\in\NaturalNumbers\), put \(s=2^n\), and for \(a\in\NaturalNumbers\) write
\[
 a=2sB+R,\qquad 0\le R<2s,\qquad
 C=\begin{cases}R,&R\le s,\\2s-R,&R>s.\end{cases}
\]
The executable definition performs exactly this quotient--remainder fold, but the following
analytic formula is not exposed as a public theorem with these variables and hypotheses.
\begin{candidate}[Triangular block reduction]
\[
 0\le C\le s,\qquad
 \FabiusGlobal\!\left(\frac a{2^n}\right)
 =(-1)^{\BinaryDigitSum(B)}F\!\left(\frac C{2^n}\right).
\]
\end{candidate}
\begin{obligation}[Discharged]
Package the quotient--remainder decomposition, folded numerator, and Thue--Morse block sign
as a public theorem about \lean{extendedFabiusDyadicValue}; then compose it with
\lean{extendedFabiusDyadicValue_cast}.  The evaluator definition, its cast theorem, and the
private first-block fold used inside the proof do not themselves state the displayed result.
\end{obligation}
Both levels are now public in \path{GlobalBlockFold.lean}, with the register's explicit
variables and hypotheses \(a=2sB+R\), \(0<a\), \(R<2s\): the executable fold is
\lean{extendedFabiusDyadicValue\_block\_fold}, the folded-numerator bound \(0\le C\le s\)
is \lean{blockFoldCore\_le}, and the displayed analytic formula
\(\FabiusGlobal(a/2^n)=(-1)^{\BinaryDigitSum(B)}F(C/2^n)\), for every bounded Fabius function, is
\lean{extendedFabius\_block\_fold}, composed through
\lean{extendedFabiusDyadicValue\_cast}, \lean{fabiusDyadicUnit\_eq\_fabiusDyadic}, and
\lean{fabiusDyadic\_cast}.

\subsection{Least-common-denominator valuation}

Let \(\Delta_n\) be the least common multiple of the reduced denominators of
\(\RvachevUp(a/2^n)\) over odd support numerators \(a\).  Lean proves that the reduced denominator
of \(F(2^{-n})\) divides \(\Delta_n\), and separately computes the two-adic valuation of
\(F(2^{-n})\).  The deduction combining those facts is now a named theorem.
\begin{candidate}[Discharged]
For every \(n\ge1\),
\[
 \TwoAdicValuation(\Delta_n)\ge \binom n2+1+\TwoAdicValuation(n!).
\]
\end{candidate}
This is \lean{dyadicDenominator_padicVal_two_lower_bound} in \path{PaperStatements.lean},
which states the displayed inequality with exactly the hypothesis \(1\le n\).  Its proof is
the requested bridge: it passes from the divisibility supplied by \lean{proposition_fifteen}
to the natural two-adic valuation of \(\Delta_n\).

\subsection{Rvachev rational-input wrapper}

The formal API has exact reduced-rational \texttt{Option} wrappers for the bounded function
and the signed global extension.  For \(\RvachevUp\) it once had only an already decoded
exponent--integer-numerator evaluator; the matching wrapper now exists.
\begin{candidate}[Discharged]
Define an exact rational-input evaluator for \(\RvachevUp\) that returns ``not dyadic'' exactly when
the reduced denominator is not a power of two, and otherwise returns the rational value
computed by \lean{rvachevDyadic}.
\end{candidate}
The wrapper is \lean{evalRvachevDyadic} in \path{Arithmetic.lean}, built on
\lean{dyadicExponent?} in the same shape as \lean{evalFabiusDyadic} and
\lean{evalExtendedFabiusDyadic}.  Both requested theorems are public: the recognition
theorem \lean{evalRvachevDyadic_eq_none_iff} in \path{DyadicCorrectness.lean}, and the
correctness theorems \lean{evalRvachevDyadic_eq_some_correct} and
\lean{evalRvachevDyadic_complete_correct} in \path{GlobalDyadic.lean}, the latter producing
the rational value and its identification with analytic \(\RvachevUp\) for every dyadic rational
input.  The formalization is more general than the candidate asked for: the recognition
theorem is an instance of a single evaluator-generic lemma quantified over an arbitrary
\(\NaturalNumbers\to\IntegerNumbers\to\RationalNumbers\) evaluator, so the same lemma also yields
\lean{evalFabiusDyadic_eq_none_iff} and \lean{evalExtendedFabiusDyadic_eq_none_iff}.

\section{Centered finite splines and local finiteness}

\subsection{Global positive-part form and a rejected raw-spline recurrence}

Let \(S_p\) denote the finite-cutoff centered spline from
\eqref{dyadicweb:eq:centered-spline-local}.  The following globally stated positive-part
formula now has an exact public counterpart.
\begin{candidate}[Discharged]
For \(p\ge1\) and \(x\in\RealNumbers\),
\[
 S_p(x)=\frac1{2^{\binom p2}p!}
 \sum_{r\ge0}(-1)^{\BinaryDigitSum(r)}
 \left(2^px-\frac12-r\right)_+^p,
\]
where the sum is pointwise finite.
\end{candidate}
\begin{warning}[The formerly proposed raw recurrence is false globally]
The identity
\[
 S_{p+1}(x)\stackrel{?}{=}\int_0^1S_p(2x-u)\,\Differential u
\]
is not a formalization target as stated.  Substitution in the displayed positive-part
formula at \(p=1\) and \(x=1\) gives \(S_2(1)=1\), whereas the proposed integral is
\(1/2\).  The public Lean recurrence concerns the centered finite CDF, not this all-real raw
spline.
\end{warning}
\begin{obligation}[Discharged]
Expose a public all-real theorem for the positive-part formula with its exact normalization,
domain, and pointwise-finiteness interpretation.  If a raw-spline smoothing law is desired,
first formulate the correct identity and domain rather than reusing the refuted display
above.  The public declarations
\lean{uniformCenteredPartialCDF_eq_integral} and
\lean{fabiusUniformSpline_eq_centeredPartialCDF_all} are formal background; private
fixed-range positive-spline lemmas and a bridge valid only on \([0,1]\) do not state these
global claims.
\end{obligation}
The all-real theorem is now public in \path{FabiusUniformSpline.lean}:
\lean{fabiusUniformSpline_eq_tsum_positivePart} states, for \(p\ge1\) and every real \(x\),
\[
 S_p(x)=\bigl(2^{\binom p2}\,p!\bigr)^{-1}
 \sum_{r\ge0}(-1)^{\BinaryDigitSum(r)}\left(2^px-\tfrac12-r\right)_+^{p},
\]
with \(\sum_{r\ge0}\) read as a \lean{tsum} and the positive part as
\(\max(\,\cdot\,,0)^p\); the pointwise-finiteness interpretation is its companion
\lean{positivePart_summand_support_finite}, which exhibits the summand's support inside
\(\{r<\Ceiling{2^px}\}\).  The proof needs no fixed-range detour: left of the support both
sides vanish, and for \(x\ge0\) the positive parts cut the \lean{tsum} down to exactly the
definitional prefix of the centered spline, the reversed power contributing the sign
\((-1)^p\) absorbed by the normalizing constant.  No raw-spline smoothing law is asserted,
in accordance with the warning above.

\subsection{What ``a factor of two'' is intended to compare}

Let \(U_p(x)=\Pr(T_p\le x)\) be the uncentered finite CDF and
\(C_p(x)=U_p(x-2^{-p-1})\) its midpoint correction.
\begin{candidate}[Discharged]
For every \(p\in\NaturalNumbers\) and \(x\in\RealNumbers\), the respective sandwich-and-Lipschitz arguments certify
\[
 |U_p(x)-F(x)|\le 2^{1-p},
 \qquad
 |C_p(x)-F(x)|\le 2^{-p}.
\]
Thus the centered construction halves this particular certified upper bound.
\end{candidate}
The two inequalities are \lean{abs_uniformPartialCDF_sub_fabiusReal_le} and
\lean{abs_uniformCenteredPartialCDF_sub_fabiusReal_le} in \path{FabiusComputability.lean},
each stated for every \(p\in\NaturalNumbers\) and every real \(x\), and each deduced there from the
sandwich theorems \lean{uniformPartialCDF_sandwich} and
\lean{uniformCenteredPartialCDF_sandwich} together with the global Lipschitz bound.  The
factor-two comparison is
\lean{uniformCenteredPartialCDF_error_bound_eq_half_uniformPartialCDF_error_bound} in the
same module, whose documentation states the intended reading verbatim: it compares the
bounds, not the actual pointwise errors, and makes no optimality claim.

\subsection{Uniform CDF control and weak convergence}

The centered-spline error estimate and pointwise CDF convergence are public Lean results.
The removed prose also compared the uniform estimate with weak convergence for the same
finite laws:
\begin{candidate}[Discharged]
For \(p\ge1\), let \(\mu_p\) be the law of the finite weighted uniform-coordinate sum,
translated to the right by \(2^{-p-1}\), and let \(\mu\) be the
weighted-infinite-sum law.  Their CDFs are respectively \(C_p\) and \(F\), and
\[
 \sup_{x\in\RealNumbers}|C_p(x)-F(x)|\le 2^{-p}.
\]
On \([0,1]\), the identification \(C_p=S_p\) recovers the centered-spline estimate.
Consequently \(\mu_p\) converges weakly to \(\mu\), with the displayed uniform-CDF rate
providing quantitative information beyond the resulting unquantified convergence statement.
\end{candidate}
Every part of this candidate is now public Lean.  The translated finite law is
\lean{uniformCenteredPartialLaw} in \path{UniformSplineWeakConvergence.lean}, a probability
measure by instance, with CDF identified as the midpoint-corrected finite CDF by
\lean{cdf_uniformCenteredPartialLaw}; the displayed all-real bound is
\lean{abs_uniformCenteredPartialCDF_sub_fabiusReal_le} in \path{FabiusComputability.lean};
and the weak convergence \(\mu_p\to\mu\) in the probability-measure topology is
\lean{uniformCenteredPartialLaw_tendsto}, proved by L\'evy's continuity theorem from the
pointwise characteristic-function convergence
\lean{tendsto_charFun_uniformCenteredPartialLaw}, itself a dominated-convergence consequence
of the pointwise partial-sum convergence \lean{tendsto_uniformPartialSum}.  As instructed
below, the proof does not route through \lean{finiteConvolutionProbability_tendsto}; the
phrase \emph{strictly stronger} is rendered as the neutral comparison above, so no
nonconverse is required.
\begin{obligation}[Discharged]
Bundle the translated finite laws and the limiting law as probability measures; identify
their CDFs; combine the existing interval estimate with support facts to obtain the displayed
all-real bound; and prove that the bound implies convergence in the probability-measure weak
topology.  If the phrase \emph{strictly stronger} is retained, define the intended comparison
and prove the required nonconverse as well.  Do not use
\lean{finiteConvolutionProbability_tendsto} as this bridge: that theorem concerns a different
sequence of symmetric two-atom convolution laws and the limiting measure with density
\(\RvachevUp\), not the continuous uniform-coordinate partial laws whose centered CDF is \(C_p\).
\end{obligation}

\subsection{Local finiteness of the signed translate family}

For fixed \(x\in\RealNumbers\), consider
\[
 n\longmapsto(-1)^{\BinaryDigitSum(n)}\RvachevUp(x-2n-1).
\]
\begin{candidate}[Discharged]
Its support in \(\NaturalNumbers\) is finite; equivalently, only finitely many translates in the defining
series for \(\FabiusGlobal(x)\) are nonzero.
\end{candidate}
The finite-support statement is public as
\lean{extendedFabius_summand_hasFiniteSupport} in \path{GlobalExtension.lean}: for each
fixed real \(x\) it asserts \texttt{Function.HasFiniteSupport} of the signed translate
sequence itself, with local finiteness as its conclusion rather than as a step inside
another proof.  The public summability theorem \lean{extendedFabius_summable} is now a
corollary of it.

\section{Step cells and corrected-prefix conventions}

\subsection{Exact histogram-cell tiling}

Put
\[
 g_n=2^{n+1}-n-2,\qquad
 \lambda_{n,m}=\frac{2m-g_n}{2^{n+1}},\qquad
 I_{n,m}=\left[\lambda_{n,m}-2^{-n-1},
                  \lambda_{n,m}+2^{-n-1}\right].
\]
\begin{candidate}[Discharged]
Consecutive cells meet only at their common endpoint, and
\[
 \bigcup_{m=0}^{g_n}I_{n,m}
 =\left[-1+\frac{n+1}{2^{n+1}},
          1-\frac{n+1}{2^{n+1}}\right].
\]
\end{candidate}
\begin{obligation}[Discharged]
Package adjacency, finite interval coverage, and both endpoint calculations as a public Lean
theorem.  If ``support'' means \(\operatorname{Function.support}\) of the step approximant
rather than the union of its defining cells, also prove positivity of every in-range
coefficient: the current support theorem is only an enclosure.
\end{obligation}
The package is public in \path{EarlyApproximants.lean}.  Adjacency is
\lean{stepIntervalRight_eq_left_succ} together with
\lean{stepInterval_inter_succ}, which computes the intersection of consecutive cells as
exactly the shared-endpoint singleton; coverage is \lean{stepInterval_biUnion} (abutting
cells tile one closed interval, for any top index); the endpoint calculations are
\lean{stepIntervalLeft_zero_eq} and \lean{stepIntervalRight_degree_eq}; and the displayed
tiling identity, packaged, is \lean{stepInterval_tiling}.  The conditional coefficient
half is also done: \lean{approximationPolynomial_coeff_pos} proves every in-range
coefficient positive (each coefficient of the product of all-ones geometric factors
receives at least one contributing monomial), upgrading the support enclosure to an exact
description.

\subsection{Why half-endpoint cells are needed}

The public step approximant uses half weight at each cell endpoint.  The elementary
counterexample that motivated that convention was removed because it has not itself been
packaged as a theorem:
\begin{candidate}[Discharged]
Modify the level-one step approximant by replacing its half-endpoint representatives with
ordinary closed-interval indicators, leaving its coefficients and normalization unchanged.
At \(x=0\), the cells \([-1/2,0]\) and \([0,1/2]\) then each contribute one, so the modified
value is \(2\), rather than the intended value \(1\).
\end{candidate}
\begin{obligation}[Discharged]
Define the ordinary-closed variant and prove its level-one value at zero.  The public Lean API
proves the endpoint values and adjacency law in
\lean{halfEndpointIntervalIndicator_self_left},
\lean{halfEndpointIntervalIndicator_self_right}, and
\lean{halfEndpointIntervalIndicator_add_adjacent}, as well as the corrected value in
\lean{stepApproximant_one_zero}; it neither defines the ordinary-closed variant nor states
the displayed double-counting result.  A source comment explaining why the convention is
useful is not a theorem.
\end{obligation}
The counterexample is now a theorem in \path{EarlyApproximants.lean}:
\lean{closedIntervalIndicator} is the naive representative,
\lean{naiveStepApproximant} the modified approximant with unchanged coefficients and
normalization, and \lean{naiveStepApproximant_one_zero} proves the displayed
double-counting value \(2\) at the origin for \(n=1\), directly against
\lean{stepApproximant_one_zero}'s corrected value \(1\).

\subsection{Independent roles of the order and coordinate corrections}

The quantitative Thue--Morse synthesis already preserves both mechanisms in detail: the
one-order coefficient shift and the centered affine coordinate.  The stronger statement that
the two repairs are independently necessary still lacks one exact theorem.
\begin{candidate}
For \(k\ge1\) and \(0\le m<2^k\), the order correction matches polynomial index \(k-1\)
with prefix row \(k\):
\[
 [z^m]P_{k-1}(z)=\sigma_k(m),
\]
whereas the centered coordinate
\(\lambda_{n,m}=(2m-g_n)/2^{n+1}\) independently repairs the geometric interpretation of that
row.  Retaining either naive convention leaves a defect that the other correction does not
repair.
\end{candidate}
\begin{obligation}[Discharged]
Package two separate failure witnesses---an explicit counterexample to matching \(P_k\)
with prefix row \(k\), and the endpoint obstruction for the naive grid coordinate---and
combine them into a theorem saying that neither correction alone gives the corrected
coefficient-and-coordinate interpretation.
\end{obligation}
The witnesses and their combination are now theorems in
\path{ThueMorseApproximation.lean}, beside the corrected identity.  The naive-order
witness is \lean{iteratedPrefix\_ne\_approximationPolynomial\_coeff\_same\_index}: at \(k=1\),
\(m=1\), inside the dyadic cutoff, the row value is \(0\) while \([z^1]P_1=1\) --- a defect
purely in the coefficients, untouched by any coordinate convention.  The naive-coordinate
witness is \lean{naiveGridCoordinate\_ne\_stepIntervalLeft}: the uncentered grid
(\lean{naiveGridCoordinate}) puts the first level-one cell's left endpoint at \(-1/4\)
while the corrected tiling requires \(-1/2\) --- a defect purely in the geometry, untouched
by any order convention.  The combined statement is
\lean{order\_and\_coordinate\_corrections\_independent}.

\section{Already-preserved audit dispositions}

Two removed primary claims need no duplicate candidate here.  Outcome-wise absolute
convergence of the random coordinate series is preserved, with its geometric majorant and
stronger uniform-convergence argument, in
\eqref{integration:eq:weighted-coordinate-sum}--\eqref{integration:eq:coordinate-majorant}.
Finite-\(Q\) dependence of the discrete rows and the warning that those rows are not partial
sums of the binary or global \(q\)-series are preserved in
\cref{dyadicweb:sec:discrete,dyadicweb:tab:two-infinite}.  Their common-limit theorem is
formalized; the finite-stage dependence and non-equality statements remain quarantined here.

\section{Complex-shift spline estimate}

The estimate cited by the primary exposition has the factor \(\exp(|Q-\tfrac12|)\).  Put
\[
 N_p(x)=\max\!\left(0,\Floor{2^px+\frac12}\right),\qquad
 S_p^{(Q)}(x)=\frac{(-1)^p}{2^{\binom p2}p!}
 \sum_{r=0}^{N_p(x)-1}(-1)^{\BinaryDigitSum(r)}(r-2^px+Q)^p,
\]
with an empty sum when \(N_p(x)=0\).  The proposed sharpening was removed from the primary
exposition, and has since been formalized:
\begin{candidate}[Discharged]
For every \(p\in\NaturalNumbers\) with \(p\ge1\), every \(Q\in\ComplexNumbers\), and every \(x\in\RealNumbers\),
\[
 \left|S_p^{(Q)}(x)-S_p^{(1/2)}(x)\right|
 \le2^{-(p-1)}
 \left(\exp\!\left(\left|Q-\frac12\right|\right)-1\right).
\]
\end{candidate}
\begin{warning}
The accompanying assertion that the relevant triangular exponent attains equality
\emph{exactly in the top term} is false when \(p=2\): both endpoint indices attain the same
exponent.  The formalization avoids the refuted claim by proving the exponent comparison as
an inequality, from monotonicity of \(d\mapsto\binom d2\), and never asserts a unique
maximizing index.
\end{warning}
The public \(\exp-1\) estimate is
\lean{norm_fabiusComplexShiftSpline_sub_center_le_half_pow_mul_exp_sub_one_all} in
\path{FabiusComplexShiftSpline.lean}, and it is more general than the candidate: it carries
no \(p\ge1\) hypothesis, holding for every \(p\in\NaturalNumbers\) and every real \(x\), with the
naturally truncated exponent making the degenerate degree an instance rather than an
exception.  Its finite Taylor majorant is formalized separately, as the exact branch
expansion \lean{fabiusComplexShiftSpline_sub_center_eq_taylorBranches} in the same module
together with the reflected factorial sum bounded in
\path{FabiusDiscreteLimitComplexShift.lean}.  The weaker \(\exp\) form
\lean{norm_fabiusComplexShiftSpline_sub_center_le_half_pow_mul_exp_all} is derived from it.

\section{Prefix-kernel interpretations}

The binomial kernel and the resulting prefix-sum formula are formalized.  These two
interpretations are not:
\begin{candidate}
\begin{itemize}
\item for \(i,k\in\NaturalNumbers\) with \(k\ge1\), \(\binom{i+k-1}{k-1}\) counts ordered \(k\)-tuples
      of nonnegative integers summing to \(i\);
\item for \(k\in\NaturalNumbers\) with \(k\ge1\), the kernel is the \(k\)-fold discrete convolution
      power of the constant kernel.
\end{itemize}
\end{candidate}
\begin{obligation}[Discharged]
Choose a finite-support or locally finite convolution model, prove the counting bijection,
and connect it to the existing \texttt{iteratedPrefixKernel}.
\end{obligation}
Both interpretations are now theorems in \path{PrefixKernelCounting.lean}, in the finite-support model: the
tuples form \lean{Finset.Nat.antidiagonalTuple} and the convolution is the finite
antidiagonal sum \lean{natConv}.  The counting bijection is
\lean{card_antidiagonalTuple} (stars and bars for weak compositions: \(\binom{n+k}{k}\)
ordered \((k+1)\)-tuples of naturals summing to \(n\)), proved from the
first-coordinate decomposition \lean{antidiagonalTuple_succ_eq} and the hockey-stick
identity; \lean{constKernelConvPow_eq_card} identifies the \(k\)-fold convolution power
of the constant kernel with the tuple count; and the connections to the kernel are
\lean{iteratedPrefixKernel_eq_card_antidiagonalTuple} and
\lean{iteratedPrefixKernel_eq_constKernelConvPow}, for every \(k\ge1\).

\section{An actual infinite formal product}

Lean proves coefficient stabilization of the finite products
\(\prod_{j<r}(1-z^{2^j})\).  The stronger construction formerly asserted was:
\begin{candidate}
Construct \(\prod_{j\ge0}(1-z^{2^j})\) as a formal infinite product and prove that it equals
the Thue--Morse power series.
\end{candidate}
\begin{obligation}[Discharged]
Select an appropriate formal-power-series topology or summability API, define the product,
and prove equality with the coefficient-stabilized series.  Until then, the notation is
only shorthand for finite stabilization.
\end{obligation}
The product is now genuine, in \path{ThueMorseFormalProduct.lean}.  The chosen API is
Mathlib's coefficientwise topology on \(\IntegerNumbers[[X]]\)
(\lean{PowerSeries.WithPiTopology}), in which convergence is exactly coefficient
stabilization.  \lean{hasProd\_one\_sub\_X\_two\_pow} proves that the partial products
over \emph{arbitrary} finite sets of factors converge to the Thue--Morse series,
\lean{multipliable\_one\_sub\_X\_two\_pow} packages multipliability, and
\lean{tprod\_one\_sub\_X\_two\_pow} is the display:
\(\prod'_{j\ge0}(1-X^{2^j})\) \emph{equals} the Thue--Morse series as a \lean{tprod}.
The engine \lean{coeff\_prod\_one\_sub\_X\_pow} is proved over an arbitrary commutative
ring and arbitrary exponent families: a finite product of factors \(1-X^{m_j}\) with all
exponents above \(d\) is indistinguishable from \(1\) in every coefficient up to \(d\).

\section{Formal-background crosswalk}

The following declarations support the proved background cited above.  None of them is
itself a counterpart for the adjacent candidate statement; where a candidate has since
acquired an exact counterpart, that declaration is named inline with the candidate and is
not repeated in this table.

\begingroup
\footnotesize
\setlength{\tabcolsep}{0.6em}
\noindent\begin{tabularx}{\textwidth}{@{}>{\raggedright\arraybackslash}p{0.29\textwidth}
  >{\raggedright\arraybackslash}X@{}}
\toprule
Background & Lean declaration and module\\
\midrule
Polynomially weighted Fourier decay &
\lean{rvachevFourier_real_rapidDecay} in \path{PoissonSummation.lean}\\
Even moments and exact dyadic evaluation &
\lean{moment}, \lean{fabiusDyadicValue}, \lean{fabiusDyadicValue_cast},
\lean{fabiusDyadic_cast} in \path{Arithmetic.lean}, \path{PaperStatements.lean}, and
\path{DyadicAnalytic.lean}\\
Complex-shift exponential bound &
\lean{norm_fabiusComplexShiftSpline_sub_center_le_half_pow_mul_exp} in
\path{FabiusComplexShiftSpline.lean}\\
Prefix kernel and convolution formula &
\lean{iteratedPrefixKernel} and \lean{iteratedPrefix_eq_sum_kernel} in
\path{ThueMorsePrefix.lean}\\
Even/odd Thue--Morse recurrences &
\lean{binaryWeight_two_mul} and \lean{binaryWeight_two_mul_add_one} in
\path{DyadicClosedForm.lean}\\
Finite block product and coefficient stabilization &
\lean{thueMorseBlockPolynomial_eq_product} and \lean{coeff_finite_thueMorse_product} in
\path{ThueMorseGenerating.lean}\\
Signed dyadic evaluator and analytic cast &
\lean{extendedFabiusDyadicValue} and \lean{extendedFabiusDyadicValue_cast} in
\path{Arithmetic.lean} and \path{GlobalDyadic.lean}\\
Dyadic denominator and two-adic inputs &
\lean{proposition_fifteen}, \lean{dyadicDenominator_eq_fabiusDyadicDenominator}, and
\lean{fabiusAtInverseTwoPow_padicVal_two_closed} in \path{PaperStatements.lean} and
\path{TwoAdic.lean}\\
Dyadic recognition and decoded Rvachev evaluation &
\lean{dyadicExponent?_exists_iff} and \lean{rvachevDyadic_cast_global} in
\path{DyadicCorrectness.lean} and \path{OriginalPaperSupplement.lean}\\
Centered finite-CDF integral and spline bridge &
\lean{uniformCenteredPartialCDF_eq_integral} and
\lean{fabiusUniformSpline_eq_centeredPartialCDF_all} in
\path{FabiusUniformSpline.lean}\\
Uncentered and centered probability sandwiches &
\lean{uniformPartialCDF_sandwich} and \lean{uniformCenteredPartialCDF_sandwich} in
\path{FabiusUniformSpline.lean}\\
Centered uniform-CDF control and pointwise convergence &
\lean{abs_fabiusUniformSpline_sub_fabiusReal_le_all},
\lean{uniformCenteredPartialCDF_tendsto_weightedSumCDF}, and
\lean{weightedSumCDF_eq_fabiusReal} in \path{FabiusComputability.lean},
\path{FabiusUniformSpline.lean}, and \path{ProbabilityRepresentation.lean}\\
Separate symmetric-atomic-law weak limit &
\lean{finiteConvolutionProbability_tendsto} and
\lean{finiteConvolutionProbability_tendsto_fabius} in \path{WeakConvergence.lean}\\
Signed translate summability and one-block evaluation &
\lean{extendedFabius_summable} and \lean{extendedFabius_eq_single_translate} in
\path{GlobalExtension.lean}\\
Step-approximant support enclosure and cell coordinates &
\lean{support_stepApproximant_subset}, \lean{polynomialAtomLocation}, and
\lean{stepApproximant_at_polynomialAtomLocation} in \path{StepApproximationLimit.lean},
\path{EarlyApproximants.lean}, and \path{ThueMorseApproximation.lean}\\
Half-endpoint cell gluing and the level-one value &
\lean{halfEndpointIntervalIndicator_self_left},
\lean{halfEndpointIntervalIndicator_self_right},
\lean{halfEndpointIntervalIndicator_add_adjacent}, and
\lean{stepIntervalRight_eq_left_succ}, together with
\lean{stepApproximant_one_zero}, in \path{EarlyApproximants.lean} and
\path{StepApproximationLimit.lean}\\
Corrected prefix order and centered sampling &
\lean{iteratedPrefix_eq_approximationPolynomial_coeff},
\lean{correctedPrefixCoefficient_eq_stepApproximant}, and
\lean{correctedPrefixGridSample_eq_stepApproximant} in
\path{ThueMorseApproximation.lean}\\
\bottomrule
\end{tabularx}
\endgroup

\section{Endpoint material in the consolidated volume}

The combined inverse and small-argument saddle synthesis preserves the removed raw
Mellin-inversion derivation, exponential Gamma--zeta decay estimates, numerical periodic profiles, proposed
ordered-composition formulas for saddle coefficients, explicit higher Lambert polynomials,
a proposed leading-derivative law, uniqueness and phase-density arguments, half-moment
asymptotics, and strengthened endpoint comparisons.

The primary article retains only the exact negative-Laplace decomposition, the proved
Gamma--zeta Fourier series, the evaluated constant, lower-Lambert reduction, Bromwich
inversion, the first two periodic saddle corrections, and the recursive all-orders theorem.

\section{Provenance}

This register was created from the claim-by-claim audit of the former drafts and primary
exposition on 25 August 2026.  Related numerical and exploratory material is preserved in
the dyadic, \(q\)-binomial, integration, repeated-integration, and combined inverse/small-argument
syntheses of this same research-frontier volume.

% END POST-CONSOLIDATION REGISTER: Primary_Exposition_Gap_Register/Primary_Exposition_Gap_Register.tex

\clearpage
\section*{Consolidated promotion rule}
\addcontentsline{toc}{section}{Consolidated promotion rule}
When a frontier claim is formalized, promotion is claim-by-claim rather than
source-by-source.  The required evidence is an exact current Lean declaration whose
hypotheses, normalization, domain, and conclusion match the mathematical sentence.  The
primary exposition should cite that name and module, while this volume should remove or
relabel the discharged obligation without erasing any still-unformalized surrounding
argument.  Until that audit is complete, the claim remains here.

\end{document}
